\documentclass[11pt,a4paper]{report}
\usepackage{fontspec,amsmath,amssymb,amsthm,microtype,xcolor}
\setmainfont{lmroman10-regular.otf}[BoldFont=lmroman10-bold.otf,ItalicFont=lmroman10-italic.otf,BoldItalicFont=lmroman10-bolditalic.otf]
\usepackage[margin=25mm]{geometry}
\usepackage[unicode,hidelinks]{hyperref}
\setlength{\emergencystretch}{3em}
\setlength{\parskip}{0.4em}
\setlength{\parindent}{0pt}
\setcounter{tocdepth}{1}
\makeatletter
\renewcommand*\l@section{\@dottedtocline{1}{1.5em}{3.8em}}
\renewcommand*\@pnumwidth{2.5em}
\renewcommand*\@tocrmarg{3.5em}
\makeatother
\setcounter{secnumdepth}{2}
\allowdisplaybreaks[2]
\raggedbottom
\title{\texorpdfstring{Kinetic Interacting Particle Langevin Monte Carlo}{Kinetic Interacting Particle Langevin Monte Carlo}\\[0.6em]\large Mathematical proofs}
\author{Omer Deniz Akyildiz\\[0.4em]\normalsize Prepared with AI assistance}
\date{}
\begin{document}
\pagenumbering{roman}
\maketitle
\chapter*{About this document}
\addcontentsline{toc}{chapter}{About this document}
These draft notes present the assumptions, convergence theorems and detailed mathematical proofs for this development. The supporting arguments are arranged so that prerequisites precede their applications. They use ordinary mathematical notation throughout.

Prepared by Omer Deniz Akyildiz with AI assistance, these notes form part of an ongoing effort to formalize Langevin algorithms comprehensively. This is draft work and may contain mistakes.

\section*{Literature and scope}
\begin{enumerate}
\item Paul Valsecchi Oliva and Ö. Deniz Akyildiz (2026). \emph{Kinetic interacting particle Langevin Monte Carlo}. Statistics and Computing 36, article 221. \url{https://doi.org/10.1007/s11222-026-10980-z}.
\item Arnak S. Dalalyan and Lionel Riou-Durand (2020). \emph{On sampling from a log-concave density using kinetic Langevin diffusions}. Bernoulli 26(3), 1956–1988. \url{https://doi.org/10.3150/19-BEJ1178}.
\item Pierre Monmarché (2021). \emph{High-dimensional MCMC with a standard splitting scheme for the underdamped Langevin diffusion}. Electronic Journal of Statistics 15(2), 4117–4166. \url{https://arxiv.org/abs/2007.05455v4}.
\end{enumerate}
\tableofcontents
\clearpage
\pagenumbering{arabic}
\chapter{Assumptions and main results}
\section{\texorpdfstring{The actual Brownian KIPLD parameter convergence proposition}{The actual Brownian KIPLD parameter convergence proposition}}


Let \(N\ge1\), and suppose the joint potential \(U(\theta,x)\) is differentiable, \(m\)-strongly convex, and has \(L\)-Lipschitz gradient, with \(0<m\le L\). Initialize the scaled position \(Z_0\sim\nu\) with finite second moment. Independently initialize scaled velocity as \(\mathcal N(0,N^{-1}I)\), and use the independent Brownian source. If \(\gamma^2\ge m+L\), the constructed physical KIPLD parameter law satisfies


\[\boxed{W_2(\mathcal L(\Theta_t),\pi_\Theta)\le\sqrt2e^{-mt/\gamma}\bigl(\mathbb E\|Z_0-\bar Z_*\|^2\bigr)^{1/2}.}\]


Here \(\bar Z_*\sim\pi_Z\) is independent of \(Z_0\), and \(\pi_Z(dz)\propto e^{-V_N(z)}dz\). The target parameter marginal is exactly \(\pi_\Theta(d\theta)\propto\ell(\theta)^N d\theta\). This is the supplied manuscript’s continuous convergence proposition.




The proof uses \(c=\sqrt N\), unit-temperature friction \(c\gamma\), and time \(t/c\). The sum potential has curvature \(Nm\), so its mixing exponent becomes \((Nm)/(c\gamma)\,(t/c)=mt/\gamma\). The flow is rescaled in velocity only; its position endpoint is unchanged. A separate theorem verifies its full physical integral equation and its independent thermal initial law. Brownian scaling proves that the transformed forcing has amplitude \(\sqrt{2\gamma/N}\). Projection to the parameter cannot increase a coupling’s cost, and the independent initial coupling gives the displayed constant.



The two terms have different interpretations. The exponential term measures dynamical distance from the equilibrium position law; the statistical radius measures concentration of that equilibrium parameter marginal around its unique maximizer. Taking time large removes the first term, while increasing particle count is needed to reduce the second.


{\small\color{gray}Supporting proof: Section~\ref{proof-module-KIPLMC-Continuous-lean}.\par}
\section{\texorpdfstring{KIPLMC1: the complete parameter RMS theorem}{KIPLMC1: the complete parameter RMS theorem}}


Use the same assumptions and thermal initialization as above. Let \(\theta_*\) be the unique marginal-likelihood maximizer. KIPLMC1 calls the shared Dalalyan–Riou-Durand exponential step with the correlated Gaussian velocity–position innovations. If


\[\gamma^2\ge m+L,\qquad 0<h,\qquad4\gamma Lh\le m,\]


then its parameter iterates obey


\[\boxed{\bigl(\mathbb E\|\Theta_n-\theta_*\|^2\bigr)^{1/2}\le C_1q^n+hC_2+\sqrt{\frac{d_\theta}{Nm}},}\]
\[q=1-\frac{3mh}{4\gamma},\qquad C_1=\sqrt2\bigl(\mathbb E\|Z_0-\bar Z_*\|^2\bigr)^{1/2},\]
\[C_2=\frac{\sqrt2L}{m}\sqrt{\frac{d_\theta+Nd_x}{N}}.\]


The three terms have different meanings. The first remembers initialization and vanishes geometrically. The second is the numerical bias and decreases with the step. The third is the statistical concentration radius and decreases with particle count. Increasing the number of iterations cannot remove either persistent floor.




The proof transfers the exact Gaussian Markov law through the clock and velocity conjugacy, applies the shared first-order KLMC theorem, and checks every power of \(N\). The factor \(NL\) from smoothness, the transformed step \(h/\sqrt N\), and curvature \(Nm\) leave precisely \(h(L/m)\sqrt{2(d_\theta+Nd_x)/N}\). Parameter projection and the L\ensuremath{{}^{2}} triangle inequality add the sharp equilibrium radius. The independent initial coupling bounds the true initial Wasserstein distance. Gaussian laws, their independence, target invariance, and finite transport moments are proved internally.



The numerical floor depends on the intensive dimension \(d_x+d_\theta/N\). The equilibrium parameter contribution decreases as \(N^{-1/2}\), while the latent part of the numerical floor remains at fixed step. Consequently increasing particle count alone cannot drive the total parameter error to zero.


{\small\color{gray}Supporting proof: Section~\ref{proof-module-KIPLMC-CanonicalExponential-lean}.\par}
\section{\texorpdfstring{KIPLMC2: the supplied rate with an explicit uniform constant}{KIPLMC2: the supplied rate with an explicit uniform constant}}


KIPLMC2 uses the shared O-B-A-B-O composition with independent Gaussian thermostat vectors. Under the same joint potential assumptions, suppose


\[\gamma>0,\qquad \gamma^2\ge4L,\qquad0<h,\qquad33\gamma^3h\le m.\]


Put \(d_z=d_\theta+Nd_x\) and \(D_0=W_2(\nu,\pi_Z)\). The supplied manuscript’s bound holds with the following explicit constant, independent of \(N,n,h,d_\theta,d_x\):


\[C=\sqrt2+\frac{6\sqrt2\gamma}{m}\left(4\sqrt L+\frac{Lm}{66\gamma^3}\right)+\frac1{\sqrt m}.\]
\[\boxed{\begin{aligned}(\mathbb E\|\Theta_n-\theta_*\|^2)^{1/2}\le C\bigl[&(\sqrt N D_0+hN\sqrt{d_z})e^{-mhn/(6\gamma)}\\&+h\sqrt{Nd_z}+\sqrt{d_\theta/N}\bigr].\end{aligned}}\]


The proof first obtains a stronger direct bound. In the weighted phase metric, the two-force OBABO step contracts by \(r=\sqrt{1-mh/(3\gamma)}\). Compare one such step with an exact O–Hamiltonian–O step, which preserves the proved Gibbs target. The common-noise local defect is quadratic in \(h\); summing the contractive perturbations gives the persistent position bias


\[\frac{6\sqrt2\gamma}{m}\,h\sqrt{\frac{d_z}{N}}\left(4\sqrt L+\frac{Lh}{2}\right).\]


This argument proves sampling from the Gaussian chain directly. It requires no assumed discrete invariant law or external OBABO error estimate. Exact clock scaling and \(r^n\le e^{-mhn/(6\gamma)}\) turn that stronger result into the supplied coarse rate, retaining both of its transient terms. The final concentration term is the same sharp statistical contribution as for KIPLMC1. The accuracy corollaries expose positive step choices, natural particle counts, and rounded logarithmic iteration counts.



The constant is uniform in particle count and dimensions because the admissible step cap bounds its remaining step dependence using only curvature and friction. The displayed coarse particle dependence is a consequence of enlarging the sharper intensive bound. Its additional transient contribution decays and does not alter the equilibrium concentration term.


{\small\color{gray}Supporting proof: Section~\ref{proof-module-KIPLMC-OBABOPaperRate-lean}.\par}
\section{\texorpdfstring{Sharp concentration of the powered parameter marginal}{Sharp concentration of the powered parameter marginal}}


Let \(N\ge1\) and let the joint potential
\(U:\mathbb R^{d_\theta+d_x}\to\mathbb R\) be differentiable and
\(m\)-strongly convex, with \(m>0\). The dimensions \(d_\theta,d_x\)
may differ and may be zero. Define


\[\ell(\theta)=\int e^{-U(\theta,x)}dx,\]
\[\pi_\Theta(d\theta)\propto\ell(\theta)^N d\theta.\]


The normalization is positive and finite. There is a unique maximizer
\(\theta_\star\) of \(\ell\), and its powered marginal satisfies


\[\boxed{\int\|\theta-\theta_\star\|^2d\pi_\Theta(\theta)
\le\frac{d_\theta}{Nm}.}\]


This gives the sharp concentration contribution
\(\sqrt{d_\theta/(Nm)}\) in both particle bounds. The exact powered Gibbs target, its kinetic invariance, and both complete numerical error bounds are proved in the supporting modules.



The negative log marginal is strongly convex after latent integration, and powering the marginal likelihood multiplies this curvature by particle count. The moment estimate then has numerator exactly the parameter dimension. This sharp coefficient is carried unchanged into both numerical particle error theorems.


{\small\color{gray}Supporting proof: Section~\ref{proof-module-KIPLMC-Target-lean}.\par}
\section{\texorpdfstring{KIPLMC1 reuses exponential integration exactly}{KIPLMC1 reuses exponential integration exactly}}


Write \(z=(\theta,x^1/\sqrt N,\ldots,x^N/\sqrt N)\),
and scale velocities the same way. Put
\(V_N(z)=\sum_iU(\theta,\sqrt N z_i)\) and
\(F_N=\nabla V_N/N\). KIPLMC1 is exactly the shared exponential update
with force \(F_N\) and phase noise amplitude \(\sqrt{2\gamma/N}\):


\[\begin{aligned}
v^+&=\psi_0(h)v-\psi_1(h)F_N(z)+\sqrt{2\gamma/N}\,\varepsilon,\\
z^+&=z+\psi_1(h)v-\psi_2(h)F_N(z)+\sqrt{2\gamma/N}\,\varepsilon'.
\end{aligned}\]


After multiplying a latent coordinate by \(\sqrt N\), its velocity update
is precisely


\[V^{x_i,+}=\psi_0(h)V^{x_i}-\psi_1(h)\nabla_xU(\theta,x^i)
+\sqrt{2\gamma}\,\varepsilon^i.\]


The coordinate identity follows by differentiating the lifted potential and undoing the latent spatial scale. The same response coefficients govern both parameter and latent coordinates.



The parameter counterpart uses the particle-average gradient and thermal amplitude \(\sqrt{2\gamma/N}\). Both parameter and latent position updates integrate their corresponding velocity responses using the same correlated innovation block. These coordinate identities therefore specify the full common exponential algorithm in physical variables.


{\small\color{gray}Supporting proof: Section~\ref{proof-module-KIPLMC-Exponential-lean}.\par}
\section{\texorpdfstring{KIPLMC2 is the exact OBABO composition}{KIPLMC2 is the exact OBABO composition}}


Set \(\delta=e^{-\gamma h/2}\),
\(\sigma=\sqrt{(1-\delta^2)/N}\), and take separate noise vectors
\(\xi,\xi'\) for thermostat half-steps. KIPLMC2 is exactly


\[v_{1/2}=\delta v+\sigma\xi-\tfrac h2F_N(z),\qquad
z^+=z+hv_{1/2},\qquad
v^+=\delta\bigl(v_{1/2}-\tfrac h2F_N(z^+)\bigr)+\sigma\xi'.\]


The new position is used when recomputing the second force. The five subflows compose in the displayed order. The parameter coordinates use the average pair gradient, while undoing latent scaling gives the individual pair gradients and ordinary latent thermal amplitude.



The two thermostat vectors are independent standard Gaussians and independent of the current state. Their amplitudes give covariance \(I/N\) in scaled equilibrium velocity. Multiplying latent position and velocity by \(\sqrt N\) recovers the individual pair-gradient kick and unit-temperature latent refresh.


{\small\color{gray}Supporting proof: Section~\ref{proof-module-KIPLMC-Splitting-lean}.\par}
\chapter{Mathematical argument}
\section{\texorpdfstring{Use interacting particles to integrate out latent variables}{Use interacting particles to integrate out latent variables}}


The observation \(y\) is fixed. A model gives a joint density
\(p_\theta(x,y)\), but the quantity to optimize is its integral over an
unobserved variable:


\[p_\theta(y)=\int p_\theta(x,y)dx,
\qquad\theta_*=\arg\max_\theta\log p_\theta(y).\]


Computing this integral and its gradient exactly can be difficult. The
particle method retains \(N\) latent states \(X^1,\ldots,X^N\), coupled
through the same parameter \(\theta\). The force on the parameter is the
average of their individual forces, whereas each latent particle sees its own
force:


\[F^\theta=\frac1N\sum_{i=1}^N\nabla_\theta U(\theta,X^i),
\qquad F^{x_i}=\nabla_xU(\theta,X^i).\]


Adding a parameter noise of size \(N^{-1/2}\) makes the whole system a
Langevin system with a useful explicit equilibrium. The particle number then
has two roles: it provides interacting latent samples and determines how
strongly the stationary parameter distribution concentrates near the optimizer.



The common parameter is represented once while each latent variable has its own coordinate. Conditional on that parameter, the equilibrium latent integrations factor into identical marginal likelihood integrals. Increasing particle count therefore sharpens the parameter marginal through a power of the likelihood, while increasing the computational state dimension.

\section{\texorpdfstring{The assumptions concern the whole statistical model}{The assumptions concern the whole statistical model}}


The theory assumes \(U\in C^2\), \(\mu>0\), \(L>0\), and


\[\langle z-z',\nabla U(z)-\nabla U(z')\rangle
\ge\mu\|z-z'\|^2,\qquad
\|\nabla U(z)-\nabla U(z')\|\le L\|z-z'\|,\]


for \(z=(\theta,x)\). Equivalently, in finite-dimensional Euclidean space,
\(\mu I\preceq\nabla^2U(z)\preceq LI\). Thus \(\mu\le L\) in positive
dimension. Joint curvature is essential: even a well-behaved posterior in
\(x\), with \(\theta\) held fixed, does not automatically control movement
along directions in which both variables change together.




The paper's unregularized logistic example illustrates the distinction.
Its Gaussian prior contributes a quadratic in \(x-\theta\), which leaves
common shifts partly uncontrolled; logistic curvature can vanish far from the
origin. Spanning covariates give strict convexity, while a uniform positive
strong-convexity constant requires additional regularization. The source itself
separates these experiments from its strong-convexity theorem.



The required lower curvature concerns arbitrary joint directions \((u,w)\), including directions that move parameter and latent variables together. Positive curvature in either coordinate with the other fixed is insufficient. A direction in which both move can remove a prior penalty and expose weak or vanishing curvature in the likelihood.

\section{\texorpdfstring{Give the parameter and particles their own velocities}{Give the parameter and particles their own velocities}}


KIPLD adds velocity blocks \(V^\theta,V^{x_i}\) and uses the same
friction in each:


\[d\theta=V^\theta dt,\qquad dX^i=V^{x_i}dt,\]
\[dV^\theta=-(\gamma V^\theta+F^\theta)dt
+\sqrt{2\gamma/N}\,dB^0,
\quad dV^{x_i}=-(\gamma V^{x_i}+F^{x_i})dt+\sqrt{2\gamma}\,dB^i.\]


The Brownian blocks are independent. The parameter velocity has equilibrium
covariance \(N^{-1}I\); every latent velocity has covariance \(I\). These
scales are chosen together with the average parameter force. Changing one
without changing the others generally changes the equilibrium target.




The noise acts on velocity, while positions are integrals of velocities.
This creates the smoother position paths responsible for the exponential
integrator's improved discretization dependence. The target remains a
statistical estimation problem: the final parameter distribution should
concentrate around \(\theta_*\) as the particle number grows.



The parameter noise is reduced by the square root of particle count, matching its averaged force and concentrated equilibrium marginal. Each physical latent equation retains unit temperature. Scaling all latent positions and velocities by \(N^{-1/2}\) turns these different physical amplitudes into one uniform phase-space temperature.

\section{\texorpdfstring{One rescaling reveals an ordinary kinetic diffusion}{One rescaling reveals an ordinary kinetic diffusion}}


The particle system looks anisotropic, but it becomes a standard kinetic
system in \(d_z=d_\theta+Nd_x\) dimensions. At rescaled time \(t\), define


\[\widetilde Z_t=(\theta_{\sqrt Nt},N^{-1/2}X^1_{\sqrt Nt},\ldots),
\qquad\widetilde V_t=(\sqrt N V^\theta_{\sqrt Nt},V^{x_1}_{\sqrt Nt},\ldots).\]


Set


\[\bar U_N(\theta,z_1,\ldots,z_N)=\frac1N\sum_{i=1}^NU(\theta,\sqrt N z_i),
\qquad f_N=N\bar U_N,\qquad\widetilde\gamma=\sqrt N\gamma.\]


Brownian scaling produces independent standard blocks
\(\widetilde B_t^i=N^{-1/4}B^i_{\sqrt Nt}\). Direct substitution yields


\[d\widetilde Z_t=\widetilde V_tdt,\qquad
d\widetilde V_t=-(\widetilde\gamma\widetilde V_t+\nabla f_N(\widetilde Z_t))dt
+\sqrt{2\widetilde\gamma}\,d\widetilde B_t.\]


This is Proposition A.1. The parameter position stays unchanged under the
spatial scaling, which will make projection of the final transport estimate
particularly simple. An original step \(\eta\) corresponds to a rescaled
step \(\widetilde\eta=\eta/\sqrt N\).



Position is fixed by the velocity-only transformation, so the parameter sample is unchanged. Friction and step are changed inversely, preserving their dimensionless product. The normalized Gaussian velocity and transformed Brownian amplitude must be together; a deterministic clock identity alone would not identify the sampler’s probability law.

\section{\texorpdfstring{Averaging preserves curvature after scaling the latent blocks}{Averaging preserves curvature after scaling the latent blocks}}


The gradients of the averaged potential are


\[\nabla_\theta\bar U_N=\frac1N\sum_i\nabla_\theta U(\theta,\sqrt Nz_i),
\qquad\nabla_{z_i}\bar U_N=\frac1{\sqrt N}\nabla_xU(\theta,\sqrt Nz_i).\]


Apply joint strong convexity to each pair
\((\theta,\sqrt Nz_i),(\theta',\sqrt Nz'_i)\), divide by \(N\), and sum.
The repeated parameter term and the scaled latent terms combine exactly:


\[\langle Z-Z',\nabla\bar U_N(Z)-\nabla\bar U_N(Z')\rangle
\ge\mu\left(\|\theta-\theta'\|^2+\sum_i\|z_i-z'_i\|^2\right).\]


For smoothness, Jensen bounds the squared norm of the averaged parameter
gradient by the average of squared norms. Adding the latent-gradient squares
then permits one joint \(L\)-Lipschitz estimate per particle, giving
\(\|\nabla\bar U_N(Z)-\nabla\bar U_N(Z')\|\le L\|Z-Z'\|\).




Thus the standard potential \(f_N=N\bar U_N\) has constants
\(m_N=N\mu\), \(M_N=NL\). Their ratio is \(L/\mu\), independent of
\(N\). This exact matching is what makes the KIPLMC1 error bound scale well
with the particle number.



The lifted sum potential has curvature and gradient Lipschitz constants multiplied by \(N\). Dividing the potential by \(N\) cancels both factors in physical time. As a result the synchronous rate depends on the original pair constants, while total coordinate dimension appears only through thermal and numerical moments.

\section{\texorpdfstring{Integrating out particles raises the likelihood to the Nth power}{Integrating out particles raises the likelihood to the Nth power}}


The rescaled kinetic Gibbs density is


\[\widetilde\Pi(dz,dv)\propto
\exp\left[-\sum_iU(\theta,\sqrt Nz_i)-\frac12\|v\|^2\right]dz\,dv.\]


Integrating the Gaussian velocity gives a constant. For each latent block,
substitute \(x_i=\sqrt Nz_i\). Its Jacobian is independent of \(\theta\),
so it is absorbed in normalization. The remaining product separates:


\[\pi_\Theta(\theta)\propto
\prod_{i=1}^N\int e^{-U(\theta,x_i)}dx_i
=p_\theta(y)^N=e^{-N\kappa(\theta)},\quad
\kappa=-\log p_\theta(y).\]


At stationarity, more particles therefore sharpen the likelihood landscape
rather than changing its maximizer. This is the central statistical mechanism.
It relies on the parameter noise being present at the prescribed scale; a
noise-free parameter gradient flow has a different invariant structure.




The density computation identifies the candidate marginal. A complete
continuous-time theorem additionally proves normalization, moments, and
invariance of the kinetic solution.



The latent variables become independent under the conditional Gibbs density for a fixed parameter. Fubini multiplies their identical integrals into \(\ell(\theta)^N\). Dividing by the finite positive full partition function gives the powered marginal, so the concentration argument applies to a normalized probability measure.

\section{\texorpdfstring{Convert marginal strong convexity into a sharp equilibrium radius}{Convert marginal strong convexity into a sharp equilibrium radius}}


Strong Prékopa–Leindler preserves the parameter curvature when the latent
variable is integrated out: \(\kappa\) is \(\mu\)-strongly convex. Thus
\(\theta_*\), its minimizer, is unique, and the potential of
\(\pi_\Theta\) has curvature \(N\mu\).




A useful sharp moment argument integrates the divergence of
\((\theta-\theta_*)e^{-N\kappa(\theta)}\). After justifying the boundary
limit with the Gaussian tails, it gives


\[\mathbb E_{\pi_\Theta}
\langle\theta-\theta_*,N\nabla\kappa(\theta)\rangle=d_\theta.\]


Since \(\nabla\kappa(\theta_*)=0\), strong monotonicity bounds the left
side below by \(N\mu\mathbb E\|\theta-\theta_*\|^2\). Hence


\[W_2(\pi_\Theta,\delta_{\theta_*})
=\bigl(\mathbb E_{\pi_\Theta}\|\theta-\theta_*\|^2\bigr)^{1/2}
\le\sqrt{\frac{d_\theta}{N\mu}}.\]


The coefficient is sharp: an isotropic Gaussian with covariance
\((N\mu)^{-1}I\) attains equality. This explains why a proof with an extra
\(\sqrt2\) coefficient, although valid, would not reproduce Proposition 2
exactly.



The radius is a root-mean-square statement around the maximizer, rather than a covariance statement around the marginal mean. Its explicit coefficient is important for the final triangle inequality. No equality of the mean and maximizer is required to add this equilibrium contribution to the numerical sampling error.

\section{\texorpdfstring{Separate approaching equilibrium from its concentration}{Separate approaching equilibrium from its concentration}}


The only coupling of a probability law with a point mass fixes the second
coordinate, so Wasserstein distance to a point is exactly root mean square
error. The triangle inequality gives


\[\bigl(\mathbb E\|\theta_n-\theta_*\|^2\bigr)^{1/2}
\le W_2(\operatorname{Law}(\theta_n),\pi_\Theta)
+\sqrt{\frac{d_\theta}{N\mu}}.\]


The first term is a numerical sampling error; the second is the equilibrium
statistical radius. Increasing \(n\) reduces the transient sampling error.
Decreasing the step reduces discretization bias. Increasing \(N\) reduces
the concentration term. These operations affect different pieces of the bound.




The map from full rescaled position, or full phase space, to its parameter
block is 1-Lipschitz. Applying the same projection to any coupling cannot
increase its quadratic transport cost. Therefore a bound on the ordinary
kinetic sampler immediately controls the parameter sampling term once the
algorithm-law correspondence has been proved.



For any coupling of the numerical position and an equilibrium position, project both to parameters and use Minkowski’s inequality around the maximizer. The projected discrepancy is bounded by the full scaled-state discrepancy because the projection is nonexpansive. Infimizing the latter cost yields the sampling term plus the equilibrium parameter radius.

\section{\texorpdfstring{KIPLMC1 is the same exponential step in particle coordinates}{KIPLMC1 is the same exponential step in particle coordinates}}


Reuse \(\psi_0,\psi_1,\psi_2\) from the kinetic integrator. With fresh
correlated Gaussian pairs of covariance \(C\), KIPLMC1 updates


\[\theta^+=\theta+\psi_1(\eta)V^\theta-\psi_2(\eta)F^\theta
+\sqrt{2\gamma/N}\,\varepsilon^{0\prime},\]
\[V^{\theta+}=\psi_0(\eta)V^\theta-\psi_1(\eta)F^\theta
+\sqrt{2\gamma/N}\,\varepsilon^0.\]


Each latent block uses the same formulas with \(X^i,V^{x_i},F^{x_i}\)
and noise scale \(\sqrt{2\gamma}\). This is Algorithm 1.




Proposition B.1 proves the correspondence by changing variables in the
ordinary exponential step with \(\widetilde\gamma=\sqrt N\gamma\) and
\(\widetilde\eta=\eta/\sqrt N\). The coefficients obey


\[\widetilde\psi_0(\widetilde\eta)=\psi_0(\eta),\quad
\widetilde\psi_1(\widetilde\eta)=N^{-1/2}\psi_1(\eta),\quad
\widetilde\psi_2(\widetilde\eta)=N^{-1}\psi_2(\eta).\]


The velocity and position Gaussian integrals scale by \(N^{-1/4}\) and
\(N^{-3/4}\), respectively. Their covariance entries therefore scale by
\(N^{-1/2},N^{-1},N^{-3/2}\). Matching these laws, in addition to matching
the deterministic terms, permits reuse of the same kinetic convergence proof.



The parameter force contains a particle mean because the common parameter occurs in every potential summand. A latent force belongs to only one summand. Undoing its spatial scaling cancels the averaging factor, explaining why the physical latent update uses its individual gradient and ordinary thermal amplitude.

\section{\texorpdfstring{The N factors cancel in the exponential-\allowbreak{}integrator bound}{The N factors cancel in the exponential-integrator bound}}


Apply the kinetic KLMC theorem with
\(m_N=N\mu\), \(M_N=NL\), \(h_N=\eta/\sqrt N\), and
\(\gamma_N=\sqrt N\gamma\). Then


\[\frac{m_Nh_N}{\gamma_N}=\frac{\mu\eta}{\gamma},\qquad
\frac{M_Nh_N}{m_N}\sqrt{2d_z}
=\sqrt2\frac L\mu\eta\sqrt{\frac{d_\theta}N+d_x}.\]


The friction and step restrictions reduce exactly to
\(\gamma\ge\sqrt{\mu+L}\) and \(\eta\le\mu/(4\gamma L)\).
Projection and concentration give


\[\mathcal E_n\le\sqrt2D_0
\left(1-\frac{3\mu\eta}{4\gamma}\right)^n
+\sqrt2\frac L\mu\eta\sqrt{\frac{d_\theta}N+d_x}
+\sqrt{\frac{d_\theta}{N\mu}}.\]


The true initial position Wasserstein distance is a valid \(D_0\).
The paper instead upper-bounds it by the root mean square distance to an
independently sampled stationary position. Initialization must include mutually
independent Gaussian velocity blocks, independent of position, so that the
rescaled velocity is precisely the standard Gaussian product needed by the
shared theorem.




Choosing \(N\) of order \(d_\theta/\varepsilon^2\) and \(\eta\) of
order \(\varepsilon/\sqrt{d_x}\), subject to stability, makes both floors
small. The transient then needs order \(\sqrt{d_x}/\varepsilon\) steps,
up to logarithmic and curvature factors.



The total scaled dimension is \(D=d_\theta+Nd_x\), and thermal variance per coordinate is \(1/N\). Thus the stationary velocity radius is \(\sqrt{D/N}=\sqrt{d_\theta/N+d_x}\). This identity is the probabilistic meaning of the particle-factor cancellation in the first-order numerical floor.

\section{\texorpdfstring{OBABO composes five exactly solvable pieces}{OBABO composes five exactly solvable pieces}}


KIPLMC2 uses a symmetric splitting of the same kinetic diffusion. Write
\(O\) for friction plus noise, \(B\) for the force kick, and \(A\) for free
transport. The composition is


\[O(\eta/2)\;B(\eta/2)\;A(\eta)\;B(\eta/2)\;O(\eta/2).\]


Let \(\delta=e^{-\gamma\eta/2}\),
\(s=\sqrt{1-\delta^2}\). For a standard velocity block, an \(O\) half-step
is \(v\mapsto\delta v+s\xi\); for the parameter block it is
\(V^\theta\mapsto\delta V^\theta+s\xi/\sqrt N\).
A \(B\) half-step subtracts \(\eta F/2\), and the full \(A\) step sends
\(x\mapsto x+\eta v\).




The second \(B\) kick must use the force at the newly computed position.
The last \(O\) step uses an independent Gaussian vector. Symmetry gives the
method its second-order splitting structure, while only first derivatives of
\(U\) are evaluated. This is an entirely different construction from the
Hessian-corrected KLMC2 in the first paper.



Each thermostat half step separately preserves the appropriate Gaussian velocity law, each force kick is a velocity translation, and each drift is a position translation. Their order determines which refreshed velocity moves position and where the second nonlinear force is evaluated.

\section{\texorpdfstring{Expand the five substeps to inspect every coefficient}{Expand the five substeps to inspect every coefficient}}


For one standard block with force \(g\), combine the first \(O\) and
\(B\) operations:


\[v_{1/2}=\delta v+s\xi-\frac\eta2g(x),
\qquad x^+=x+\eta v_{1/2}.\]


Apply the final force kick and OU refresh:


\[v^+=\delta\left(v_{1/2}-\frac\eta2g(x^+)\right)+s\xi'
=\delta^2v-\frac{\eta\delta}2\bigl(g(x)+g(x^+)\bigr)
+s(\delta\xi+\xi').\]


In the parameter block replace \(g\) by the average \(F^\theta\), and
both Gaussian amplitudes by \(s/\sqrt N\). The latent blocks use individual
forces and amplitude \(s\). These identities reproduce Algorithm 2 and
Proposition B.2.




The resulting law is determined by two independent standard Gaussian
vectors, independent of the current state. Proving the displayed algebra for
arbitrary supplied noise vectors certifies the deterministic update;
proving the stochastic theorem additionally constructs this Gaussian kernel,
its iterates and the required moment estimates.



The final combined noise is \(s(\delta\xi+\xi^{\prime})\), but position depends only on the first refresh through the half velocity. Therefore the full phase output has a nontrivial shared-noise structure even though the two input thermostat vectors are independent. The expanded formula retains this structure exactly.

\section{\texorpdfstring{The discrete invariant law contributes a separate bias}{The discrete invariant law contributes a separate bias}}


A finite-step OBABO kernel \(P_h\) generally has its own invariant law
\(\Pi_h\), different from the kinetic Gibbs law \(\Pi\). The cited
Monmarché analysis supplies phase-space contraction and invariant bias bounds
of the form


\[W_2(\rho P_h^n,\sigma P_h^n)\le A q^{n/2}W_2(\rho,\sigma),
\qquad b=W_2(\Pi_h,\Pi)\le B.\]


Apply the triangle inequality twice, using invariance of \(\Pi_h\):


\[W_2(\rho P_h^n,\Pi)
\le A q^{n/2}\bigl(W_2(\rho,\Pi)+B\bigr)+B.\]


This explains two different occurrences of the numerical bias: a persistent
\(B\) and a transient \(Aq^{n/2}B\). Dropping the latter would lose a term
present in the supplied paper's revised theorem.




These inequalities are substantial results about the OBABO kernel,
including existence of \(\Pi_h\). A The proof theorem that takes them as hypotheses
checks the transport combination but does not prove the cited OBABO analysis.
The coverage audit distinguishes that conditional use from a completed
argument.



The development gives a second route that avoids constructing the discrete invariant law. An exact O–Hamiltonian–O step preserves the continuous Gibbs target. Compare it with the OBABO update on the same initial state and the same Gaussian vectors. The Verlet local defect has L\ensuremath{{}^{2}} cost of order \(h^2\sqrt{d_z}\). Contractive perturbation sums this cost directly over the chain. The resulting intensive bound is stronger than the supplied coarse estimate, so it also implies every term retained in the printed theorem.

\section{\texorpdfstring{Track the OBABO constants before hiding them}{Track the OBABO constants before hiding them}}


For the rescaled potential, the source uses


\[A_N=\sqrt3\max\{\sqrt{NL},(NL)^{-1/2}\},\qquad
q=1-\frac{\mu\eta}{3\gamma}.\]


Under \(\gamma\ge2\sqrt L\) and
\(0<\eta\le\eta_0=\mu/(33\gamma^3)\), one has \(0<q<1\).
Write \(M_L=\max\{\sqrt L,L^{-1/2}\}\). Because \(N\ge1\),
\(A_N\le\sqrt{3N}M_L\). The detailed bias expression in the source
can be bounded by


\[B_N\le\sqrt3G_*\eta\sqrt{Nd_z},\]
\[G_*=\frac{6\gamma L}{\mu}M_L
\left(1+\frac{\eta_0L}{2\sqrt\mu}\right)
\left(1+e^{L\eta_0^2}
\left(\frac{\eta_0}{6}+\frac{\eta_0^2L}{24}\right)\right).\]


All factors increase with a nonnegative step, so evaluating at the fixed
upper limit \(\eta_0\) removes hidden step dependence. This makes \(G_*\)
a constant depending only on \(\mu,L,\gamma\).




The product \(A_NB_N\) is proportional to \(\eta N\sqrt{d_z}\), while
\(B_N\) itself is proportional to \(\eta\sqrt{Nd_z}\). The two terms have
different powers of \(N\), which is why the source tracks them separately
before introducing the implicit constant.



The clock changes both curvature constants and friction, while leaving the physical contraction ratio unchanged. Simplifying the bias only after substituting every factor avoids introducing spurious growing powers of \(N\). The resulting intensive estimate is stronger than the coarser printed envelope.

\section{\texorpdfstring{The revised KIPLMC2 theorem keeps both transient terms}{The revised KIPLMC2 theorem keeps both transient terms}}


Independent Gaussian initialization gives the same velocity marginal as
\(\widetilde\Pi\). Coupling these equal velocity blocks and using position
projection proves the exact identity


\[W_2(\rho_0,\widetilde\Pi)
=W_2(\nu,\widetilde\pi_Z)=D_0.\]


Insert the scaled \(A_N,B_N\) into the discrete equilibrium comparison,
project to \(\theta\), and add the concentration radius. With a constant
\(C\) depending only on \(\mu,L,\gamma\), the supplied manuscript obtains


\[\mathcal E_n\le C\left[
\bigl(\sqrt N D_0+\eta N\sqrt{d_\theta+Nd_x}\bigr)
e^{-\mu\eta n/(6\gamma)}
+\eta\sqrt{N(d_\theta+Nd_x)}
+\sqrt{d_\theta/N}\right].\]


The exponential follows from
\((1-\mu\eta/(3\gamma))^{n/2}\le e^{-\mu\eta n/(6\gamma)}\).
The paper gives an explicit common coefficient by taking the maximum of
\(\sqrt3M_L,\sqrt3G_*,3M_LG_*\) and absorbing the concentration factor
\(\mu^{-1/2}\).




This bound has worse particle dependence than the KIPLMC1 bound. The paper
regards that dependence as a limitation of the available estimate and compares
it with numerical behavior; the experiments do not themselves remove the
factor from theorem.



The extra transient term is harmless for validity but matters when reproducing the displayed manuscript formula. It decays with the same rate as the initial transport discrepancy. The sharper invariant-comparator estimate bounds the chain directly and implies this coarser statement by nonnegative termwise enlargements.

\section{\texorpdfstring{Separate guarantees, complexity and experimental observations}{Separate guarantees, complexity and experimental observations}}


For KIPLMC2, choosing \(N\asymp d_\theta/\varepsilon^2\) gives


\[\sqrt{N(d_\theta+Nd_x)}
\asymp\frac{d_\theta}{\varepsilon}
\sqrt{1+\frac{d_x}{\varepsilon^2}}.\]


To make the persistent numerical error of order \(\varepsilon\), take
the minimum of the stability threshold and
\(\varepsilon/\sqrt{N(d_\theta+Nd_x)}\). In the small-error regime this
is a step of order \(\varepsilon^3/(d_\theta\sqrt{d_x})\). Both transient
terms can then be controlled after
\(\widetilde O(d_\theta\sqrt{d_x}/\varepsilon^3)\) steps, provided the
initial transport error grows at most polynomially in particle number and
dimension so that its logarithm remains in the suppressed factor.




The manuscript also studies logistic regression, Wisconsin cancer data, a
Bayesian neural network, a multi-well model and runtime comparisons. Its signed
weighted area between error curves is
\(\sum_{i=1}^M2i(c^{\mathrm{IPLA}}(i)-c^{\mathrm{KIPLMC2}}(i))/(M(M+1))\).
Positive values favor KIPLMC2 by this metric. Its predictive log score averages
the log probability assigned to each correct test label.




These are empirical comparisons with concrete data and numerical settings.
The strong-convexity convergence theorems, their asymptotic complexity
corollaries, the algorithm definitions and the experimental metrics are
different kinds of content. The logistic example is globally smooth and can be strictly convex under a full-span feature condition, while common-translation directions still prevent a positive uniform strong-convexity constant without additional regularization.



The statistical choice of particle count first fixes the dimension-dependent numerical coefficient. The bias budget then fixes a step, and the contraction rate fixes the iteration count. Empirical advantages at particular model settings remain separate from these global sufficient conditions, especially when a model lacks uniform joint strong convexity.

\chapter{Supporting proofs}
We now supply the detailed supporting arguments. Each section introduces the objects and estimates needed by the results that follow it.

\section{\texorpdfstring{Scaled particle geometry and exact coordinate updates}{Scaled particle geometry and exact coordinate updates}}
\label{proof-module-IPLA-Geometry-lean}


The scaling is chosen so that one Euclidean norm simultaneously measures the common parameter and the average physical particle energy: \(\|z\|^2=\|\theta\|^2+N^{-1}\sum_i\|x_i\|^2\). This identity explains both the common particle factor in the potential's regularity constants and its cancellation against the Euler step. All coordinate calculations below use orthogonal parameter and latent blocks.


\subsection{\texorpdfstring{The coordinate spaces and their dimensions}{The coordinate spaces and their dimensions}}


The scaled state is \(\mathbb R^{d_\theta}\times(\mathbb R^{d_x})^N\), and a joint pair belongs to \(\mathbb R^{d_\theta}\times\mathbb R^{d_x}\). Counting these orthogonal coordinates proves

\[\dim(\mathrm{State})=d_\theta+Nd_x,\qquad \dim(\mathrm{Joint})=d_\theta+d_x.\]

The block map extracts \(z_i\), while \(A_i z=\operatorname{joint}_i(z)=(\theta,\sqrt N z_i)\) restores the physical particle. The two coordinate evaluation identities follow directly from this definition. The definition of \(V_N\) then sums the original potential over these pairs.





In ordinary vector notation the state is \(\mathbb R^{d_\theta}\times(\mathbb R^{d_x})^N\). The latent coordinates in this space are scaled, rather than physical, particle locations. In particular, moving only the parameter changes every joint pair, whereas moving one latent block changes only one pair. This difference accounts for the two gradient formulas derived below and prevents an incorrect extra factor of \(N\) in a physical particle update.

\subsection{\texorpdfstring{Replication preserves a precise quadratic identity}{Replication preserves a precise quadratic identity}}


Each \(A_i\) is linear: checking a parameter coordinate and a latent coordinate separately proves preservation of addition, scalar multiplication and subtraction. Every parameter coordinate occurs \(N\) times, while each latent coordinate is multiplied by \(\sqrt N\) once. Therefore

\[\sum_i\|A_i z\|^2=N\|\theta\|^2+N\sum_i\|z_i\|^2=N\|z\|^2.\]

The proof expands squared Euclidean norms into finite sums and uses \((\sqrt N)^2=N\), which is valid also when \(N=0\). This exact identity supplies the constants in both the convexity and smoothness arguments.





The identity also applies to differences, since every extraction map is linear: \(\sum_i\|A_i(z-w)\|^2=N\|z-w\|^2\). It is essential to sum squared norms before taking a square root. Bounding every individual extraction by its largest possible operator norm and then summing would lose the precise particle scaling. The exact sum retains the geometry of the shared parameter and disjoint latent blocks.

\subsection{\texorpdfstring{Joint strong convexity lifts to N times the constant}{Joint strong convexity lifts to N times the constant}}


If \(U\) is \(m\)-strongly convex, then for \(a,b\ge0\), \(a+b=1\), each pair satisfies

\[U(aA_ix+bA_iy)\le aU(A_ix)+bU(A_iy)-\frac{mab}{2}\|A_i(x-y)\|^2.\]

Sum over \(i\), use linearity to identify the arguments with \(A_i(ax+by)\), and use the quadratic replication identity on \(x-y\). The result is the strong-convexity inequality for \(V_N\) with constant \(Nm\). The continuous linear map \(A_i\) is obtained from its finite-dimensional linear map; its evaluation theorem is definitional. If \(U\) is differentiable, the finite sum of \(U\circ A_i\) is differentiable by the chain rule.





The deficit in the strong-convexity inequality is additive. Its sum is \((mab/2)\sum_i\|A_i(x-y)\|^2\), hence exactly \((Nmab/2)\|x-y\|^2\). Thus replication strengthens convexity in the scaled state by \(N\). This conclusion follows for all convex-combination weights, including the endpoints, where the deficit vanishes. Differentiability is needed later for gradients, but not for this convex-combination argument itself.

\subsection{\texorpdfstring{The gradient chain rule and the lifted Lipschitz constant}{The gradient chain rule and the lifted Lipschitz constant}}


Differentiating in a direction \(w\) gives

\[\langle\nabla V_N(z),w\rangle=\sum_i\langle\nabla U(A_i z),A_iw\rangle.\]

Cauchy–Schwarz in the particle index, followed by the quadratic replication identity, proves \(\sum_i\|A_iv\|\|A_iw\|\le N\|v\|\|w\|\). If \(\nabla U\) is \(L\)-Lipschitz, subtract the two gradient identities to obtain

\[\langle\nabla V_N(x)-\nabla V_N(y),w\rangle\le NL\|x-y\|\|w\|.\]

Choose \(w=\nabla V_N(x)-\nabla V_N(y)\). If its norm is zero the desired bound is immediate; otherwise cancel its positive norm. Thus \(\nabla V_N\) is \(NL\)-Lipschitz, with the same particle factor as the convexity constant.






For completeness, take absolute values in the directional identity before applying Cauchy–Schwarz. The resulting bound is \(|\langle\nabla V_N(x)-\nabla V_N(y),w\rangle|\le NL\|x-y\|\|w\|\). Duality of the Euclidean norm then gives the Lipschitz estimate directly. Because both regularity constants contain the same factor \(N\), the condition number of the lifted potential is \(L/m\), the original joint condition number.

\subsection{\texorpdfstring{The two gradient blocks}{The two gradient blocks}}


For a unit parameter basis vector \(e_j\), all \(A_i e_j\) equal the matching joint parameter vector. For a latent basis vector belonging to particle \(i\), \(A_qe_{i,j}=0\) for \(q\ne i\), and \(A_i e_{i,j}=\sqrt N e_j\). Insert these basis vectors into the directional gradient formula. The Euclidean inner product reads off coordinates, yielding

\[\nabla_\theta V_N(z)=\sum_i\nabla_\theta U(\theta,\sqrt Nz_i),\qquad (\nabla V_N(z))_i=\sqrt N\nabla_xU(\theta,\sqrt Nz_i).\]

This calculation supplies the parameter average and ordinary physical-particle drift after the step size is divided by \(N\).





The parameter derivative is a sum because the same parameter appears in every term of \(V_N\). The latent derivative has only one contributing summand, but differentiation of \(x_i=\sqrt N z_i\) contributes its factor \(\sqrt N\). No derivative of a normalized empirical potential is being taken here: the lifted potential is the full sum. Averaging enters only when this gradient is multiplied by the Euler step \(\gamma/N\).

\subsection{\texorpdfstring{The parameter projection is a contraction}{The parameter projection is a contraction}}


The continuous linear map \(p_\theta z=\theta\) simply selects the parameter coordinates. Expanding squared norms gives \(\|p_\theta z\|^2\le\|z\|^2\), since the omitted latent squared coordinates are nonnegative. Both norms are nonnegative, so taking square roots proves \(\|p_\theta z\|\le\|z\|\). The definition of operator norm then gives \(\|p_\theta\|\le1\). The coordinate evaluation theorem identifies each output coordinate directly. These inequalities permit quadratic transport to control the parameter RMS.





For a difference of states, \(\|p_\theta z-p_\theta w\|\le\|z-w\|\) follows by linearity. Therefore any coupling of full states also couples their parameter projections with no greater quadratic cost. The same contraction applies pointwise inside an \(L^2\) norm, which will allow a full-state transport estimate to control the parameter error about a deterministic statistical center.

\subsection{\texorpdfstring{One ULA step is exactly one physical IPLA step}{One ULA step is exactly one physical IPLA step}}


Define the scaled update with physical step \(h\) by

\[z'=z-\frac hN\nabla V_N(z)+\sqrt{2h/N}\,\xi.\]

Its parameter coordinate is \(\theta'=\theta-(h/N)\sum_i\nabla_\theta U(\theta,x_i)+\sqrt{2h/N}\,\xi_\theta\). For \(N>0\), multiply the \(i\)-th scaled latent coordinate by \(\sqrt N\). The drift coefficient becomes \((\sqrt N)^2h/N=h\), and \(\sqrt N\sqrt{2h/N}=\sqrt{2h}\). Consequently

\[x_i'=x_i-h\nabla_xU(\theta,x_i)+\sqrt{2h}\,\xi_i.\]


The score-block lemmas are obtained by testing the derivative against single basis vectors: a parameter basis vector is replicated across all pairs, while a latent basis vector appears in one pair with factor \(\sqrt N\). In the particle update, multiplying the scaled coordinate by \(\sqrt N\) turns the drift coefficient \((\gamma/N)\sqrt N\sqrt N\) into \(\gamma\), and the noise coefficient into \(\sqrt{2\gamma}\). The parameter block keeps the averaged force and noise coefficient \(\sqrt{2\gamma/N}\). Positive \(N\) protects exactly these unscaling cancellations.





In the probabilistic algorithm take \(\gamma>0\), so every displayed noise coefficient is an ordinary positive square root. The parameter variance per iteration is \(2\gamma/N\) in each coordinate, while the variance of each physical latent coordinate is \(2\gamma\). This asymmetry is a consequence of the scaled Euclidean representation. It recovers the intended interacting algorithm exactly, before making any approximation or convergence estimate.

\section{\texorpdfstring{The scaled particle force and its geometry}{The scaled particle force and its geometry}}
\label{proof-module-KIPLMC-Geometry-lean}


The shared IPLA state stores \(z=(\theta,X^1/\sqrt N,\ldots,X^N/\sqrt N)\) in a Euclidean coordinate space with independent parameter and latent dimensions. This section keeps physical time, scaling velocities by the same spatial map. In this convention the kinetic temperature is \(1/N\), the potential is the particle average \(\bar U_N\), and the force is its gradient. It is equivalent to the manuscript's additional time change but avoids reimplementing the shared kinetic step.



The scaled average potential keeps curvature and smoothness independent of particle count in physical time. Its coordinate derivatives explain why the parameter force is averaged while physical latent forces are unaveraged.

\subsection{\texorpdfstring{Define the average and prove that its gradient is the force}{Define the average and prove that its gradient is the force}}


the scaled particle state space reuses the scaled Euclidean IPLA state. The existing lifted potential is \(V_N(z)=\sum_iU(\theta,\sqrt Nz_i)\). Define \(\bar U_N=V_N/N\) and \(F=\nabla V_N/N\). To prove \(\nabla\bar U_N=F\), compare their inner products against an arbitrary vector. The derivative rule for a scalar multiple gives \(D\bar U_N[w]=N^{-1}DV_N[w]\), and the inner-product characterization of gradients identifies the two vectors. The argument uses equality of all inner products, the proved differentiability of the lifted potential, and the derivative rule for a constant multiple. For \(N>0\), cancellation of the nonzero particle number gives \(V_N=N\bar U_N\).



The spatial scaling places the common parameter and all latent particles in one Euclidean space while preserving the physical latent coordinates through multiplication by \(\sqrt N\). Dividing the lifted potential by \(N\) makes its gradient the physical-time force. The potential and temperature must be scaled consistently.

\subsection{\texorpdfstring{Average the lifted strong convexity}{Average the lifted strong convexity}}


The existing particle geometry proves \(V_N\) is \(Nm\)-strongly convex whenever the pair potential is \(m\)-strongly convex. Multiply its convex-combination inequality by \(N^{-1}>0\): the function terms become \(\bar U_N\), and the quadratic correction \(N^{-1}(Nm)\|x-y\|^2/2\) becomes \(m\|x-y\|^2/2\). The argument constructs the strong-convexity structure directly, obtains the lifted inequality, multiplies it by the nonnegative inverse particle number, and uses cancellation of nonzero denominators and expansion of the scalar polynomials to cancel constants. The proof requires positive particle count but imposes no artificial equality of the parameter and latent dimensions.



The quadratic correction in strong convexity scales by exactly the same scalar as the potential. Since the lifted potential’s curvature is \(Nm\), averaging removes the particle factor. The resulting force geometry therefore has the original pair curvature, regardless of the total particle-space dimension.

\subsection{\texorpdfstring{Recover the mean parameter force and unaveraged latent force}{Recover the mean parameter force and unaveraged latent force}}


The parameter block of \(\nabla V_N\) is the sum of pair parameter gradients, so \(F_\theta=N^{-1}\sum_i\nabla_\theta U(\theta,X^i)\). Each latent block of \(\nabla V_N\) is \(\sqrt N\nabla_xU\). Dividing by \(N\) and undoing the spatial scaling gives \(\sqrt N F_i=\nabla_xU(\theta,X^i)\). The parameter proof simplifies with the existing gradient-coordinate theorem. The latent proof expands coordinate scalar multiplication, groups the coefficient into \((\sqrt N)^2N^{-1}\), rewrites \((\sqrt N)^2=N\), and cancels the nonzero \(N\). These identities verify the force used by both particle algorithms.



The parameter appears in every summand, so its derivative is a particle average. Each latent coordinate appears in just one summand, and undoing the coordinate scaling recovers its unaveraged derivative. These different factors explain the apparent asymmetry between parameter and latent equations.

\subsection{\texorpdfstring{The force retains the pair smoothness constant}{The force retains the pair smoothness constant}}


The shared lifted gradient has Lipschitz constant \(NL\). Since \(F=\nabla V_N/N\), its difference satisfies \(\|F(x)-F(y)\|\le N^{-1}(NL)\|x-y\|=L\|x-y\|\). The argument starts from the existing lifted Lipschitz estimate, rewrites norms of scalar multiples with the nonnegative inverse particle count, multiplies the inequality by that inverse, and cancels \(N\). This proves that the physical-time scaled force has a smoothness constant independent of the number of particles, a key feature of the favorable KIPLMC1 scaling.



The average-potential normalization cancels the lifted gradient’s factor \(N\). Consequently neither the strong-convexity nor smoothness constant grows with particle count in physical time. The dimension still affects thermal moments, but the deterministic contraction rate is particle independent.

\section{\texorpdfstring{Exponential response kernels and valid Gaussian covariance matrices}{Exponential response kernels and valid Gaussian covariance matrices}}
\label{proof-module-KineticLangevin-Coefficients-lean}


The functions \(\psi_0(t)=e^{-\gamma t}\), \(\psi_{k+1}(t)=\int_0^t\psi_k(s)ds\) encode exactly integrated friction. The Hessian method also uses \(\varphi_{k+1}(t)=\int_0^te^{-\gamma(t-s)}\psi_k(s)ds\). Integral definitions avoid division by friction, so zero friction is handled exactly. The module proves their calculus and constructs positive semidefinite Gram covariances; it does not yet construct the corresponding Gaussian random variables.



These kernels encode time integration of friction without approximation. Their primitive identities and covariance relations are the common algebraic input for both the first-order and Hessian-corrected numerical schemes.

\subsection{\texorpdfstring{Define iterated response coefficients and prove their calculus}{Define iterated response coefficients and prove their calculus}}


The recursive definition chooses an exponential at index zero and an interval integral at a successor index; the two the defining formulas lemmas expose these cases to rewriting. Continuity follows by induction: exponential continuity handles zero, while continuity of the primitive handles the successor. The fundamental theorem of calculus then proves \(\psi'_{k+1}(t)=\psi_k(t)\), using interval integrability and continuity of the integrand. The separate zero-index derivative is \(\psi'_0=-\gamma\psi_0\), obtained from the exponential chain rule; algebraic rewriting and simplification reconcile its scalar expression with the recursive notation.



The first three responses have a direct dynamical interpretation. The exponential propagates an initial velocity, its first primitive propagates velocity into position, and its second primitive propagates a constant force into position. The recursive notation also makes all higher-order time integrations unambiguous.

\subsection{\texorpdfstring{Derive friction balance before dividing by friction}{Derive friction balance before dividing by friction}}


Differentiate \(1-\psi_0(t)\): its derivative is \(\gamma\psi_0(t)\), and its value at zero is zero. Integrating gives \(\gamma\psi_1(t)=1-\psi_0(t)\). Integrating this identity once more gives \(\gamma\psi_2(t)=t-\psi_1(t)\). First, the argument applies the integral form of the fundamental theorem, then commutes a scalar through the integral. The second proof integrates equality of two functions, distributes the integral over subtraction with its integrability proofs, and evaluates the constant integral. Only after establishing these division-free identities do the closed formula for \(\psi_1\) and the closed formula for \(\psi_2\) use \(\gamma\ne0\) to derive the familiar quotient formulas.



For positive friction the explicit expressions are \(\psi_1(t)=(1-e^{-\gamma t})/\gamma\) and \(\psi_2(t)=t/\gamma-(1-e^{-\gamma t})/\gamma^2\). The integral identities should remain primary: they remain well conditioned mathematically at zero friction, where the quotient expressions have removable singularities.

\subsection{\texorpdfstring{Prove exact time scaling of every primitive}{Prove exact time scaling of every primitive}}


For \(c\ne0\), substitution \(r=cs\) gives \(\psi_{c\gamma,k}(t)=c^{-k}\psi_{\gamma,k}(ct)\). At index zero this is the scalar identity \(-(c\gamma)t=-\gamma(ct)\). At a successor, use the induction hypothesis inside the integral, pull out \(c^{-k}\), and change variables in the remaining primitive; one additional inverse factor appears. The argument uses the time-substitution rule, so oriented interval integrals also cover negative \(c\). Substituting \(t=h/c\) yields \(\psi_{c\gamma,k}(h/c)=\psi_{\gamma,k}(h)/c^k\). This is the exact coefficient identity used in mass and particle rescalings.



Each additional time integral contributes one factor of the inverse clock speed. This explains the power \(c^{-k}\), while the exponential itself has index zero and no such factor. The identity must scale friction and time together; changing only one does not preserve the response.

\subsection{\texorpdfstring{Establish initial values, signs and polynomial upper bounds}{Establish initial values, signs and polynomial upper bounds}}


At time zero the exponential coefficient is one and every primitive is zero. For \(t\ge0\), induction and nonnegative integration show every coefficient is nonnegative, even without a sign restriction on \(\gamma\). If also \(\gamma\ge0\), then \(e^{-\gamma t}\le1\). Integrating once gives \(\psi_1(t)\le t\), and integrating again gives \(\psi_2(t)\le t^2/2\). The integral comparison theorem receives continuous, hence interval-integrable, functions and the pointwise inequality on the interval. These elementary estimates control the local time-step error without approximating the exponential by an unproved series.



When \(\gamma\ge0\), the bound \(0<e^{-\gamma s}\le1\) on positive times is the starting point. Integrating it repeatedly gives \(0\le\psi_k(t)\le t^k/k!\). These bounds retain the correct small-step power without introducing inverse-friction constants.

\subsection{\texorpdfstring{Check the zero-\allowbreak{}friction limit exactly}{Check the zero-friction limit exactly}}


Putting \(\gamma=0\) in the integral definition yields \(\psi_0=1\), \(\psi_1=t\), and \(\psi_2=t^2/2\). The first two proofs simplify the exponential and constant integral. The third uses the integral of the first power. These are identities at zero, not limits inferred from quotient formulas, so no removable-singularity argument is needed. They later turn the kinetic exponential step into the constant-acceleration position–velocity update.



At zero friction the recursion begins with the constant one. Repeated integration gives \(\psi_k(t)=t^k/k!\) exactly. Thus the frictionless limits of the numerical formulas follow by evaluating the definitions, rather than by an unproved interchange of limits and integrals.

\subsection{\texorpdfstring{Define the Hessian response and its useful factorization}{Define the Hessian response and its useful factorization}}


the damped response \(\varphi\) assigns zero to index zero as an explicit convention and defines positive indices by exponential convolution. Its rewrite lemmas expose those cases. Nonnegativity follows from positivity of the exponential and nonnegativity of \(\psi_k\). Factoring \(e^{-\gamma t}\) out of the integral gives \(\varphi_{k+1}(t)=e^{-\gamma t}\int_0^te^{\gamma s}\psi_k(s)ds\). One proves the integrands agree by the exponential addition identity and scalar arithmetic. This factorization expresses the convolution as a product of a continuous exponential and a continuous primitive, proving continuity even though the upper limit and the kernel both depend on time. Every \(\varphi_k(0)\) is zero.



The extra exponential factor represents damping after the force perturbation has acted. The distinction between \(\psi\) and \(\varphi\) matters in a Hessian correction because one response comes from integrating a prior response, while the other comes from applying friction to that response.

\subsection{\texorpdfstring{Compare convolution coefficients in the correct direction}{Compare convolution coefficients in the correct direction}}


For \(\gamma,t\ge0\), the factor \(e^{-\gamma(t-s)}\le1\) on \(0\le s\le t\), and \(\psi_k(s)\ge0\). Thus \(\varphi_{k+1}(t)\le\int_0^t\psi_k(s)ds=\psi_{k+1}(t)\). The proof explicitly establishes integrability of the convolution and primitive before applying integral comparison. At zero friction the exponential factor is one, so equality holds. In particular \(\varphi_2\le\psi_2\), correcting the opposite ordering printed in the supplied Proposition 3; both coefficients can still be bounded above by \(t^2/2\).



These comparisons use nonnegative kernels. If \(0\le k_1\le k_2\) and the response being integrated is nonnegative, then integrating against either kernel preserves the order. The sign hypotheses must be retained: convolution does not preserve such inequalities for arbitrary signed integrands.

\subsection{\texorpdfstring{A Gram quadratic form is an integral of a square}{A Gram quadratic form is an integral of a square}}


Given continuous scalar kernels \(k_i\), define \(C_{ij}=\int_0^hk_i(t)k_j(t)dt\). Commutativity proves symmetry. For a finite coefficient vector \(z\), expand the quadratic form and interchange its finite sums with integrals:

\[\sum_{i,j}z_iC_{ij}z_j=\int_0^h\left(\sum_i z_ik_i(t)\right)^2dt.\]

The argument supplies interval-integrability witnesses for each product, pulls scalar factors inside, applies interchange of an integral and a finite sum twice, and normalizes the finite-sum expansion. For \(h\ge0\) the integral is nonnegative. Combining symmetry with this quadratic-form bound proves positive semidefiniteness. This is a matrix property, not yet an identification with Itô covariance.



The covariance matrix can therefore be recognized without computing its individual entries. For any coefficient vector, its quadratic form is a squared \(L^2[0,h]\) norm. Symmetry follows from commutativity of scalar multiplication, and nonnegativity follows from the square under the integral.

\subsection{\texorpdfstring{Instantiate the first-\allowbreak{} and second-\allowbreak{}order noise covariances}{Instantiate the first- and second-order noise covariances}}


The first integrator uses kernels \((\psi_0,\psi_1)\); the Hessian integrator uses \((\psi_0,\psi_1,\varphi_2,\varphi_3)\). the two response kernels and the four response kernels list them on \(\mathrm{Fin}(2)\) and \(\mathrm{Fin}(4)\). Their covariance definitions apply the same Gram construction. Finite case splits reduce kernel continuity to the proved coefficient continuity. Each final positive-semidefinite theorem is then a direct application of the general Gram result. In the stochastic algorithm the noise is additionally required to be centered jointly Gaussian with these covariances, independent across time and coordinates; those law properties are not encoded by matrix positivity alone.



Each block corresponds to one deterministic Brownian response kernel. The first-order scheme needs two kernels, while the Hessian-corrected scheme needs four on the same interval. Their off-diagonal entries encode shared Brownian forcing and generally cannot be set to zero.

\subsection{\texorpdfstring{Strict positivity of every positive-\allowbreak{}time response}{Strict positivity of every positive-time response}}


For \(t>0\), every response coefficient \(\psi_k(t)\) is strictly positive. The exponential base case is positive at every time. At a successor, the integrand is continuous and nonnegative on the interval, and the induction hypothesis makes it strictly positive at the right endpoint. the positive-integral theorem then gives a positive primitive. The proof explicitly supplies both the pointwise nonnegative bound and an endpoint witnessing positivity. No positive-friction hypothesis is needed because an exponential remains positive even when its exponent grows.



Strict positivity is stronger than a nonnegative upper bound: integrating a continuous positive response over an interval of positive length gives a positive number. This fact later excludes a degenerate covariance direction and permits division by the diagonal pivots in a Cholesky factorization.

\subsection{\texorpdfstring{All higher responses are bounded by their frictionless Taylor coefficient}{All higher responses are bounded by their frictionless Taylor coefficient}}


At zero friction induction gives \(\psi_{0,k}(t)=t^k/k!\). Integrating the induction hypothesis uses \(\int_0^ts^kds=t^{k+1}/(k+1)\), and the factorial recursion \((k+1)!=(k+1)k!\) matches the denominator. For nonnegative friction and time, compare each \(\psi_{\gamma,k}\) with its zero-friction counterpart. The exponential comparison starts the induction, and integral monotonicity preserves it at each successor. Substitution of the exact zero-friction identity yields \(0\le\psi_{\gamma,k}(t)\le t^k/k!\). These estimates provide every higher-order coefficient bound without a formal power-series approximation.



The induction integrates the preceding polynomial envelope: \(\int_0^t s^k/k!\,ds=t^{k+1}/(k+1)!\). This simple observation is the source of the exact factorials in local-error estimates. Replacing all responses by \(t\) would lose the order of the scheme.

\subsection{\texorpdfstring{Evaluate the first Hessian response without division}{Evaluate the first Hessian response without division}}


The convolution of the two exponential responses has a constant integrand: \(e^{-\gamma(t-s)}\psi_0(s)=e^{-\gamma t}\). Hence \(\varphi_1(t)=t\psi_0(t)\). Reversing the integration variable also gives \(\int_0^te^{-\gamma(t-s)}ds=\psi_1(t)\). To compute the next balance, multiply the \(\varphi_2\) integrand by \(\gamma\), substitute \(\gamma\psi_1(s)=1-\psi_0(s)\), and simplify the exponential product. Integrating yields \(\gamma\varphi_2(t)=\psi_1(t)-t\psi_0(t)\). One proves equality of the integrand functions, commutes constants through integrals, and evaluates the constant integral. These identities remain valid at zero friction.



The division-free expression is particularly useful inside integrals. It can be substituted before simplifying friction factors, so a cancellation remains valid uniformly as friction tends to zero. This avoids introducing an artificial positive-friction hypothesis into identities that hold for every real friction.

\subsection{\texorpdfstring{Compute covariance entries and the exact friction–noise balance}{Compute covariance entries and the exact friction–noise balance}}


The velocity variance coefficient is \(C_{00}=\int_0^he^{-2\gamma t}dt=\psi_{2\gamma,1}(h)\), proved by combining exponentials. Since \((\psi_1^2/2)'=\psi_0\psi_1\) and \(\psi_1(0)=0\), the fundamental theorem gives \(C_{01}=\psi_1(h)^2/2\). Finally the friction identity at rate \(2\gamma\) yields \(2\gamma C_{00}=1-e^{-2\gamma h}\). Therefore \(\psi_0(h)^2+2\gamma C_{00}=1\), the exact stationary velocity variance balance. One proves the exponential square identity and reduces the final expression with expansion of the scalar polynomials. An independent unit-variance input velocity and the specified Gaussian noise consequently retain unit variance after the linear thermostat part.



The fluctuation–dissipation relation is a covariance identity, not an independence statement. An initial Gaussian velocity damped by \(e^{-\gamma h}\), plus an independent response of the stated variance, has exactly its original covariance. This is the algebra behind the invariant thermostat.

\subsection{\texorpdfstring{Relate the second-\allowbreak{}order coefficients with an exact integrated balance}{Relate the second-order coefficients with an exact integrated balance}}


The derivative of \((t-s)\psi_1(s)+\psi_2(s)\), with \(t\) fixed, is \((t-s)\psi_0(s)\). Integrating and reversing the variable yields \(\int_0^te^{-\gamma(t-s)}s\,ds=\psi_2(t)\). Next put the defining integrals for \(\varphi_2\) and \(\varphi_3\) under one integral. The identity \(\psi_1(s)+\gamma\psi_2(s)=s\) simplifies their integrand, giving \(\varphi_2(t)+\gamma\varphi_3(t)=\psi_2(t)\). The argument explicitly supplies interval-integrability before combining the integrals and uses a linear combination of the coefficient identities for the scalar integrand equality. For nonnegative friction and time, \(\gamma\varphi_3(t)\ge0\), hence \(\varphi_2(t)\le\psi_2(t)\). This provides a direct corrected coefficient relation, rather than the reversed intermediate inequality printed in the supplied manuscript.



The integrated balance allows terms arising from different correction kernels to be combined before norms are taken. That order is important: cancellation at the level of the vector identity can be destroyed by bounding the terms separately. The later stochastic correction estimate relies on this exact relation.

\section{\texorpdfstring{One shared exponential integrator and the Hessian correction}{One shared exponential integrator and the Hessian correction}}
\label{proof-module-KineticLangevin-Integrator-lean}


The deterministic maps accept a force vector field and supplied noise vectors. Keeping the force explicit permits particle and inverse-mass reuse. Throughout the section, the first phase coordinate is velocity and the second is position. The response-kernel construction specifies the covariance matrices required by the probabilistic interpretation; these definitions by themselves make no assertion that arbitrary noise inputs have those laws.



The update is an exact solution of the linear frozen-force problem. Local numerical error therefore comes only from the nonlinear force variation, while all friction and Gaussian response coefficients remain exact.

\subsection{\texorpdfstring{Define phase space and the frozen-\allowbreak{}force step}{Define phase space and the frozen-force step}}


the velocity–position space is the pair \(E\times E\). For supplied additive innovation \((n_v,n_x)\), define

\[v^+=\psi_0v-\psi_1g(x)+n_v,\qquad x^+=x+\psi_1v-\psi_2g(x)+n_x.\]

the exponential update implements these two expressions. the KLMC update passes innovation \(\sqrt{2\gamma}(\xi,\xi')\), matching the paper's noise normalization. The velocity and position projection theorems are definitional equalities: the defining formulas exposes the corresponding pair coordinate. This separation lets KIPLMC1 use exactly this realization with its own temperature amplitude rather than copying the formula.



Conditioning on the starting phase freezes the force at one position. The resulting affine linear system can be solved by an integrating factor for velocity and then a time integral for position. The only nonlinear approximation in this step is keeping that force constant.

\subsection{\texorpdfstring{Check zero time and cancel shared additive noise}{Check zero time and cancel shared additive noise}}


At zero time, \(\psi_0=1\), \(\psi_1=\psi_2=0\); with zero innovation the step is the identity. One proves equality of pairs by coordinate extensionality and simplification. For two starting phases driven by the same innovation, subtraction cancels each additive noise term before any estimate:

\[\Delta v^+=\psi_0\Delta v-\psi_1\Delta g,\quad
\Delta x^+=\Delta x+\psi_1\Delta v-\psi_2\Delta g.\]

The proof expands pair subtraction, distributes scalar multiplication across differences, and uses additive cancellation for the additive group rearrangement. The KLMC cancellation theorem reuses this exact identity with the scaled noise argument. It does not need Gaussianity.



For two states using the same additive innovation, subtraction removes the innovation exactly. The remaining difference is a deterministic expression involving the two starting velocities, positions, and force values. This permits the numerical comparison to be analyzed pathwise before any expectation is taken.

\subsection{\texorpdfstring{Include every Hessian correction, including noise multipliers}{Include every Hessian correction, including noise multipliers}}


With \(H=\nabla^2f(x)\), the second-order map subtracts \(\varphi_2Hv\) from velocity and \(\varphi_3Hv\) from position. Its innovations are \(\sqrt{2\gamma}(\xi_0-H\xi_2)\) and \(\sqrt{2\gamma}(\xi_1-H\xi_3)\). The definition the Hessian-corrected update includes all four terms. With zero Hessian, the deterministic and random corrections vanish and the map becomes the original KLMC map using \((\xi_0,\xi_1)\). Pair extensionality and simplification prove this reduction. When comparing two nonzero Hessians, the noise terms differ too; this explicit map prevents silently discarding that contribution.



The four innovations are correlated because all are deterministic responses to the same Brownian interval. A Hessian multiplies the appropriate noise responses as well as the deterministic response. Omitting those stochastic correction terms changes the map and invalidates a second-order comparison.

\subsection{\texorpdfstring{The frictionless map is constant acceleration}{The frictionless map is constant acceleration}}


At \(\gamma=0\), the coefficient identities give \(\psi_0=1\), \(\psi_1=h\), \(\psi_2=h^2/2\), while \(\sqrt{2\gamma}=0\). Therefore \(v^+=v-hg(x)\) and \(x^+=x+hv-h^2g(x)/2\). The proof simplifies the map with the exact zero-friction coefficient lemmas. It asserts exact integration of a frozen force, not exact integration of a spatially varying Hamiltonian force along the whole interval.



When friction is zero, \(\psi_0=1\), \(\psi_1=h\), and \(\psi_2=h^2/2\). Substitution gives the elementary constant-acceleration formulas. This is a useful consistency check on signs, on the order of the phase coordinates, and on the factor one half in position.

\subsection{\texorpdfstring{Iterate a prescribed stream of innovations}{Iterate a prescribed stream of innovations}}


the KLMC sequence starts at the supplied phase and recursively applies the KLMC update with innovation number \(k\). Its zero and successor equations are definitional equalities, providing reliable rewrite points for later induction and law identification. the Hessian-corrected sequence uses the same recursion pattern with the four-block innovation required by the Hessian-corrected update. A prescribed noise stream is a pathwise object: assigning independent Gaussian coordinate laws and proving finite moments of this recursion is the next probabilistic task, not a hidden assumption in the definitions.



The recursion is a deterministic measurable function of the initial phase and the full innovation sequence. Probabilistic assumptions are imposed on these inputs afterward. In particular, independence across time is a property of the prescribed stream and is not supplied by recursion alone.

\subsection{\texorpdfstring{Verify that the frozen deterministic response solves its differential equations}{Verify that the frozen deterministic response solves its differential equations}}


For constant force \(g\), define \(v_t=\psi_0(t)v_0-\psi_1(t)g\) and \(x_t=x_0+\psi_1(t)v_0-\psi_2(t)g\). The coefficient derivative identities immediately give \(x'_t=v_t\). For velocity they give \(v'_t=-\gamma\psi_0(t)v_0-\psi_0(t)g\). Substituting \(\gamma\psi_1=1-\psi_0\) proves this equals \(-\gamma v_t-g\). The proof builds the vector derivatives from the scalar coefficient derivatives and rewrites the friction balance before normalizing scalar actions. the exact frozen-force response identity is then a definitional equality: the deterministic step evaluates these exact frozen responses at \(h\) and adds the supplied noise. Thus the term “exponential integrator” is supported by an ODE identity.



The deterministic response has initial velocity \(v\) and acceleration \(-\gamma v_t-g(x)\). Its integrated position starts at \(x\). Verifying these equations and the initial conditions identifies the response uniquely, providing a precise reference for the local force-freezing error.

\section{\texorpdfstring{KIPLMC1 reuses the shared exponential integrator}{KIPLMC1 reuses the shared exponential integrator}}
\label{proof-module-KIPLMC-Exponential-lean}


There is one exponential-step realization: the exponential update. KIPLMC1 supplies the proved averaged particle force and the common scaled-coordinate thermal amplitude \(\sqrt{2\gamma/N}\). The coordinate identities recover the paper's Algorithm 1 directly, including its separate velocity and position innovations. Noise covariance and independent-Gaussian law conditions must still be established for any stochastic convergence theorem.



The particle method uses the common exponential response with thermal covariance scaled by the inverse particle count. Exact coordinate identities recover the parameter and latent updates while preserving correlated position–velocity noise.

\subsection{\texorpdfstring{Supply the particle temperature to the shared map}{Supply the particle temperature to the shared map}}


the kinetic noise amplitude is \(a_N=\sqrt{2\gamma/N}\). the KIPLMC1 update calls the shared exponential map with force \(F\) and additive phase innovation \(a_N\,\text{noise}\). Consequently every block uses exactly the same friction response functions. In physical time and scaled coordinates this is an inverse temperature \(N\) kinetic system. The definition accepts arbitrary innovations; when used as the algorithm their velocity–position pair must have the common Gram covariance from the kinetic library.



In scaled coordinates the equilibrium velocity covariance is \(I/N\). The fluctuation–dissipation amplitude \(\sqrt{2\gamma/N}\) corresponds to that temperature. The response covariance itself is the common exponential Gram matrix; the amplitude is applied afterward to the entire correlated pair.

\subsection{\texorpdfstring{Recover both parameter updates}{Recover both parameter updates}}


Parameter coordinates are unchanged by spatial rescaling. Substitute \(F_\theta=N^{-1}\sum_i\nabla_\theta U_i\) in the two shared expressions:

\[V^{\theta+}=\psi_0V^\theta-\frac{\psi_1}N\sum_i\nabla_\theta U_i+a_N\xi_\theta,\quad
\theta^+=\theta+\psi_1V^\theta-\frac{\psi_2}N\sum_i\nabla_\theta U_i+a_N\xi'_\theta.\]

The proofs unfold the shared map, evaluate scalar multiplication in Euclidean coordinates, rewrite the proved force identity, and normalize division as multiplication by an inverse. The distinct projections of the innovation pair retain the correlated but generally unequal velocity and position increments.



The parameter position receives an integrated version of the same velocity noise. Thus the two displayed innovations are correlated even though they are written in different equations. Keeping their common block law is necessary for the algorithm’s correct transition measure.

\subsection{\texorpdfstring{Undo the latent temperature and force scaling}{Undo the latent temperature and force scaling}}


For \(N>0\), \(\sqrt N\sqrt{2\gamma/N}=\sqrt{2\gamma}\). The proof combines square roots using nonnegativity of \(N\), then cancels the positive particle number inside the root. Multiply each latent coordinate of the shared step by \(\sqrt N\). The force identity gives \(\sqrt NF_i=\nabla_xU_i\); the amplitude identity restores \(\sqrt{2\gamma}\) noise. Thus the physical latent velocity and position updates have the unaveraged pair force and the paper's latent thermal scale. The two coordinate proofs distribute multiplication, rewrite these exact identities and finish scalar rearrangement with expansion of the scalar polynomials.



The spatial factor \(\sqrt N\) simultaneously restores the unaveraged latent force and the unit-temperature latent noise. This explains why physical latent updates have familiar single-particle coefficients while the parameter noise remains reduced by particle averaging.

\subsection{\texorpdfstring{Construct the numerical recursion from a supplied innovation stream}{Construct the numerical recursion from a supplied innovation stream}}


the KIPLMC1 sequence starts from the given phase and uses innovation number \(n\) in the successor update. Each step is the shared exponential map specialized to the particle force. Recursive construction proves there is an explicit sequence for every prescribed stream; it is not an assumed sequence satisfying a desired estimate. To obtain the paper's random chain, one must additionally equip this stream with independent centered Gaussian pairs of the right covariance and prove its independence from initialization, measurability and finite moments.



Once an initial phase and the full innovation stream are supplied, the recursion determines every state. If that stream has the independent correlated-Gaussian block law and is independent of initialization, pushing the inputs through the recursion gives the Gaussian Markov sampler.

\section{\texorpdfstring{The OBABO splitting and its exact pathwise identities}{The OBABO splitting and its exact pathwise identities}}
\label{proof-module-KineticLangevin-OBABO-lean}


OBABO brackets a kick–drift–kick Hamiltonian step between two Ornstein–Uhlenbeck half steps. The module proves the composition and its reversal and cancellation identities for arbitrary supplied vectors. The Gaussian-law and convergence analysis is separate. Its thermostat coefficient is chosen so that the damping and fresh noise exactly preserve a target scalar velocity variance.



The five-stage chronology determines both deterministic reversibility and the placement of the two independent refresh noises. Its expanded identities are the starting point for the weighted contraction calculation.

\subsection{\texorpdfstring{Define each subflow and compose O–B–A–B–O}{Define each subflow and compose O–B–A–B–O}}


In velocity–position ordering, \(O(v,x)=(\delta v+\sigma\xi,x)\), \(B_h(v,x)=(v-hg(x),x)\), and \(A_h(v,x)=(v,x+hv)\). the OBABO update composes \(O,B_{h/2},A_h,B_{h/2},O\), using distinct noise arguments in the two thermostat steps. Expanding the nested definitions gives \(v_{1/2}=\delta v+\sigma\xi-hg(x)/2\), \(x^+=x+hv_{1/2}\), and \(v^+=\delta(v_{1/2}-hg(x^+)/2)+\sigma\xi'\). the expanded five-stage OBABO formula is proved by the defining formulas because this formula follows directly from the composition; the new-position force is structurally enforced.



The first and last thermostat stages are half steps, as are the two force kicks; the position drift is a full step. Keeping this chronology explicit matters because the second force evaluation occurs after the position update. The composition is symmetric in its deterministic stages.

\subsection{\texorpdfstring{Reverse the deterministic kick and drift}{Reverse the deterministic kick and drift}}


The kick leaves position fixed, so subtracting \(hg(x)\) and then adding it cancels exactly. The drift leaves velocity fixed, so adding \(hv\) and then subtracting it also cancels. Their inverse formulas hold for arbitrary real \(h\) and arbitrary force, with no smoothness assumption. velocity reversal negates velocity and keeps position; applying it twice is the identity. Conjugating a kick or drift by this reversal changes \(h\) to \(-h\). The proofs use pair extensionality and simplification of scalar signs, making the time-reversal algebra explicit.



A kick changes velocity while holding position fixed, so reversing its step size gives its inverse. A drift changes position while holding velocity fixed and has the same property. These triangular maps are the elementary reversible pieces used to build the Hamiltonian proposal.

\subsection{\texorpdfstring{The palindromic Hamiltonian composition has an exact inverse}{The palindromic Hamiltonian composition has an exact inverse}}


the velocity-Verlet map is \(B_{h/2}A_hB_{h/2}\). To invert it, reverse the order and invert each factor. Since the composition is palindromic, this is precisely \(B_{-h/2}A_{-h}B_{-h/2}\). The proof expands the two compositions, rewrites \((-h)/2=-(h/2)\), and cancels the three adjacent inverse pairs with the preceding lemmas. The OBABO zero-time theorem separately checks that \(h=0,\delta=1,\sigma=0\) makes every substep the identity. Reversal of the deterministic map alone does not assert reversibility of a stochastic transition with a particular probability density.



Velocity reversal changes the sign of the drift direction and interchanges forward and backward kicks. The palindrome of the Verlet composition therefore yields an exact algebraic inverse. This reversibility identity holds for any force; conservation of the true Hamiltonian is a separate question.

\subsection{\texorpdfstring{Synchronous noises cancel while the new forces remain}{Synchronous noises cancel while the new forces remain}}


Apply two OBABO steps with the same \(\xi,\xi'\). The difference of half velocities is \(a=\delta\Delta v-h(g(x)-g(y))/2\). Thus \(\Delta x^+=\Delta x+ha\). The final shared refresh cancels too, leaving \(\Delta v^+=\delta(a-h(g(x^+)-g(y^+))/2)\). The proof unfolds all subflows, proves the velocity and position identities as a conjunction, and normalizes vector expressions with vector-space algebra. Although additive noises disappear explicitly from the difference, the new positions still depend on the common first noise. The identity preserves their force evaluations instead of treating them as old-position gradients.



The two trajectories use the same first thermostat innovation and the same second thermostat innovation. Their differences still contain forces at their respective updated positions. Thus cancellation of noise simplifies the comparison but does not linearize the nonlinear force.

\subsection{\texorpdfstring{Choose damping and prove the thermal variance identity}{Choose damping and prove the thermal variance identity}}


The OU half-step damping is \(\delta=e^{-\gamma h/2}>0\). The exponential addition identity gives \(\delta^2=e^{-\gamma h}\). With \(\gamma,h\ge0\), this is at most one, so \(s^2=1-\delta^2\ge0\). For a target velocity variance \(u\ge0\), set \(\sigma=\sqrt{u(1-\delta^2)}\). Then \(\delta^2u+\sigma^2=u\). One proves positivity directly from positivity of the exponential, rewrites the square with the exponential addition formula, uses the nonpositive exponent to bound the exponential, and finally applies the square-root identity followed by expansion of the scalar polynomials. For independent centered input velocity and fresh Gaussian noise this is the scalar covariance balance; Gaussian-law preservation additionally uses those distributional facts.



With \(\delta=e^{-\gamma h/2}\), the deterministic variance contribution is \(\delta^2\). An independent standard Gaussian scaled by \(\sqrt{1-\delta^2}\) supplies the missing variance. Nonnegative \(\gamma h\) ensures that the square root is real and that thermostat is contractive.

\section{\texorpdfstring{KIPLMC2 is exactly the shared OBABO composition}{KIPLMC2 is exactly the shared OBABO composition}}
\label{proof-module-KIPLMC-Splitting-lean}


In scaled physical-time particle coordinates every velocity block has temperature \(1/N\). The parameter coordinates stay unchanged. Multiplying a latent coordinate by \(\sqrt N\) recovers the physical latent position and velocity and the unit-temperature latent noise used in Algorithm 2. The module proves these coordinate identities for the five-subflow map; no OBABO contraction theorem is hidden in the definitions.



The particle OBABO map retains the exact five-stage chronology in scaled coordinates. Undoing latent scaling restores the physical thermostat temperature and the original pair-gradient half kicks.

\subsection{\texorpdfstring{Define the scaled thermostat and half velocity}{Define the scaled thermostat and half velocity}}


Let \(\delta=e^{-\gamma h/2}\) and \(\sigma_N=\sqrt{(1-\delta^2)/N}\). thermostat amplitude defines \(\sigma_N\), the KIPLMC2 update passes it and the force \(F\) to the shared OBABO map, and the half-step velocity records \(v_{1/2}=\delta v+\sigma_N\xi-hF(x)/2\). The algorithm uses this newly formed half velocity for its full position drift. The second force kick is evaluated at the new position, and the final OU half step receives the distinct noise \(\xi'\). Both the position and velocity projection formulas are definitional equalities, proved by the defining formulas.



The first thermostat refresh and first force kick create the velocity used by the drift. The final force is evaluated at the newly updated position before the second thermostat refresh. Reordering either evaluation would produce a different splitting method.

\subsection{\texorpdfstring{The parameter half velocity contains the average gradient}{The parameter half velocity contains the average gradient}}


Evaluating the half velocity on parameter index \(j\) gives

\[(v_{1/2})_{\theta,j}=\delta v_{\theta,j}+\sigma_N\xi_{\theta,j}
-\frac{h}{2N}\sum_i\nabla_\theta U(\theta,X^i)_j.\]

The proof unfolds the half-velocity definition, distributes coordinate evaluation over vector addition, subtraction and scalar multiplication, then substitutes the parameter-average force identity. The remaining equality is scalar rational algebra, discharged by expansion of the scalar polynomials. This proof keeps the particle mean in the force rather than adding an assumed algorithm-equivalence premise.



The average parameter gradient is inherited from the averaged potential. It is multiplied by the half-step length because this is one of the two B stages. The full parameter force action over a step is distributed between the old and new positions.

\subsection{\texorpdfstring{Undo latent scaling in the first O and B substeps}{Undo latent scaling in the first O and B substeps}}


The thermostat identity is \(\sqrt N\sigma_N=\sqrt{1-\delta^2}\). Its proof combines square roots and cancels the positive \(N\) inside the argument. Multiplying the latent half velocity by \(\sqrt N\) then gives

\[\sqrt N(v_{1/2})_i=\delta(\sqrt Nv_i)+\sqrt{1-\delta^2}\xi_i-\frac h2\nabla_xU(\theta,X^i).\]

The coordinate proof expands multiplication through the three terms, applies the proved thermostat and force-unscaling identities, and normalizes multiplication order with expansion of the scalar polynomials. This yields exactly the physical latent OU refresh and unaveraged force kick.



The latent scaling cancels the temperature reduction in thermostat amplitude. Its physical noise is therefore the standard unit-temperature OU noise. The same scaling cancels the averaging in the latent force, giving the physical pair-gradient kick.

\subsection{\texorpdfstring{Recover physical parameter and latent position drifts}{Recover physical parameter and latent position drifts}}


The common position formula is \(z^+=z+h v_{1/2}\). Its parameter coordinate reads \(\theta^+=\theta+hV^\theta_{1/2}\) directly because that coordinate was not scaled; its theorem is definitional. In a latent coordinate, multiply both sides by \(\sqrt N\) and distribute: \(X^{i+}=X^i+hV^{x_i}_{1/2}\). The argument rewrites the shared position identity, simplifies Euclidean coordinate operations and closes the scalar distributive equality with expansion of the scalar polynomials. Together with the half-velocity and final-velocity identities, these are the paper's complete pointwise OBABO stages. Gaussian transition laws and their numerical bias remain separate analytic assertions.



Position does not receive an independent direct noise in this splitting. Its randomness is transported by the refreshed half velocity. The parameter and latent position formulas therefore follow from the same deterministic drift stage after their respective coordinate scalings.

\section{\texorpdfstring{Exact velocity and clock conjugacies}{Exact velocity and clock conjugacies}}
\label{proof-module-KIPLMC-Rescaling-lean}


The physical-time particle realization can be expressed in unit-temperature kinetic coordinates by scaling velocity and the clock. This section proves the deterministic conjugacy of the exponential map and the damping identity. Gaussian covariance scaling and Brownian time-change laws remain distinct probabilistic results.



The deterministic velocity-and-clock conjugacy converts the physical average force to the unit-temperature sum-potential force. Position is fixed, while response and damping identities control the remaining phase coefficients.

\subsection{\texorpdfstring{Invert the velocity scale}{Invert the velocity scale}}


the velocity rescaling sends \((v,x)\) to \((cv,x)\). If \(c\ne0\), applying the inverse scale yields \((c^{-1}cv,x)=(v,x)\). The argument simplifies the nested scalar action and cancels the nonzero factor. Position is unchanged, so a parameter-position projection commutes with this part of the change of variables.



The velocity-only map is invertible for a nonzero scale and fixes every position observable. Hence it changes the phase coordinates without changing the parameter sample itself. Later law identities exploit exactly this unchanged position projection.

\subsection{\texorpdfstring{Conjugate the shared exponential step}{Conjugate the shared exponential step}}


The exact coefficient scaling gives \(\psi_{c\gamma,k}(h/c)=\psi_{\gamma,k}(h)/c^k\). Apply a step with force \(g\), friction \(c\gamma\), duration \(h/c\), initial velocity \(cv\), and additive innovation whose velocity component is also scaled by \(c\). Its result equals the velocity rescaling of a physical-time step with force \(g/c^2\), friction \(\gamma\), duration \(h\). One proves the two scalar cancellations for the first primitive, expands the map coordinatewise, rewrites all coefficient scales, and normalizes the resulting vector identities with vector-space algebra. The innovation here is already additive; obtaining its prescribed Gaussian law under this conjugacy requires a separate covariance computation.



The response with one time primitive acquires one inverse clock factor; the position-force response acquires two. Combining those factors with the scaled velocity and force gives the deterministic step conjugacy. The noise law must be transformed separately because its two response coordinates scale differently.

\subsection{\texorpdfstring{Check the invariant friction–step product and particle scale}{Check the invariant friction–step product and particle scale}}


The OU damping depends on the product \(\gamma h\), so \(e^{-(c\gamma)(h/c)/2}=e^{-\gamma h/2}\) for \(c\ne0\). The proof unfolds the damping definition, reduces equality of exponentials to equality of exponents, and cancels the nonzero factor with cancellation of nonzero denominators. For the particle choice \(c=\sqrt N\), the coefficient identity \(((\sqrt N)^2)^{-1}=N^{-1}\) follows from the nonnegative square-root square theorem. Consequently the sum-potential force in rescaled time corresponds to the average-potential force in physical time.



Thermostat damping depends only on the dimensionless product of friction and step, so that product is invariant under the clock change. Choosing the scale \(\sqrt N\) converts the average force into the sum-potential force and simultaneously normalizes thermal velocity.

\section{\texorpdfstring{The quadratic upper model}{The quadratic upper model}}
\label{proof-module-UnadjustedLangevin-Analysis-Smooth-lean}


We next obtain the complementary upper bound from a Lipschitz gradient. This argument uses only first derivatives: restrict the potential to a segment, subtract its affine approximation and the proposed quadratic error, and show that the resulting scalar function decreases.



The quadratic upper model concerns the gradient of the same differentiable potential occurring in its function values. The argument uses only first derivatives and a Lipschitz gradient; the scalar line restriction and its interval integral provide the quadratic remainder, with no Hessian assumption.



The upper quadratic model is a first-order statement. Along a segment, the variation of the gradient contributes at most \(Lt\|y-x\|^2\); its integral over \([0,1]\) is \(L\|y-x\|^2/2\). This explains the factor one half and why differentiability with Lipschitz gradient suffices in place of a Hessian bound. At \(x=y\), the segment is constant and every remainder term vanishes.

\subsection{\texorpdfstring{The descent lemma}{The descent lemma}}


Let \(H\) be a real Hilbert space and \(L\ge0\). Suppose \(f:H\to\mathbb R\) is differentiable and

\[\|\nabla f(u)-\nabla f(v)\|\le L\|u-v\|\quad(u,v\in H).\]

Then for every \(x,y\),

\[f(y)\le f(x)+\langle\nabla f(x),y-x\rangle+\frac L2\|y-x\|^2.\]

Set \(d=y-x\), \(a=\langle\nabla f(x),d\rangle\), \(b=\|d\|^2\), and introduce the corrected segment function

\[q(t)=f(x+td)-f(x)-ta-\frac L2t^2b.\]

It suffices to show \(q(1)\le q(0)=0\).



The segment has derivative \(d\). The chain rule therefore gives, for every real \(t\),

\[\frac{d}{dt}f(x+td)=\langle\nabla f(x+td),d\rangle.\]

Subtract the derivatives of the affine and quadratic terms:

\[q'(t)=\langle\nabla f(x+td),d\rangle-a-Ltb.\]

For \(t\) in the interior of \([0,1]\), we have \(t\ge0\). Lipschitz continuity gives

\[\|\nabla f(x+td)-\nabla f(x)\|\le L\|td\|=Lt\|d\|.\]

By Cauchy–Schwarz,

\[\langle\nabla f(x+td)-\nabla f(x),d\rangle\le\|\nabla f(x+td)-\nabla f(x)\|\,\|d\|\le Lt\|d\|^2.\]

Substituting \(a\) and \(b\) shows that \(q'(t)\le0\).



The derivative formula also makes \(q\) continuous on \([0,1]\) and differentiable in its interior. The derivative criterion for monotonicity now implies that \(q\) is nonincreasing on this interval.



Compare the two endpoints:

\[q(1)\le q(0),\qquad f(x+d)-f(x)-a-\frac L2b\le0.\]

Finally \(x+d=x+(y-x)=y\). Substitute the definitions of \(a,b,d\) and rearrange to obtain the quadratic upper bound.



Equivalently, integrating the segment derivative gives \(f(y)-f(x)-\langle\nabla f(x),d\rangle=\int_0^1\langle\nabla f(x+td)-\nabla f(x),d\rangle dt\). The integrand is bounded above by \(Lt\|d\|^2\). Its integral is \(L\|d\|^2/2\), which independently exposes the constant in the upper model. The monotonicity proof above realizes the same calculation without assuming a second derivative.

\section{\texorpdfstring{Supporting planes and strong monotonicity}{Supporting planes and strong monotonicity}}
\label{proof-module-UnadjustedLangevin-StrongConvexity-lean}


We begin with the deterministic geometry behind synchronous Langevin coupling. Throughout this part, \(\langle\cdot,\cdot\rangle\) and \(\|\cdot\|\) are the real inner product and its norm. Strong convexity uses the normalization \(m/2\): the function \(v\mapsto f(v)-\frac m2\|v\|^2\) is convex.



The normalization is \(m\)-strong convexity of \(f\) on the whole space, meaning convexity after subtracting \(m\|x\|^2/2\). \(Df(x)\) is the continuous linear derivative; \(\nabla f(x)\) is its Hilbert-space representative, so \(Df(x)[v]=\langle\nabla f(x),v\rangle\). The proofs explicitly move between these two representations.



The normalization of strong convexity matters: subtracting \(m\|x\|^2/2\) removes exactly the curvature \(m\), and its gradient is \(mx\). Thus the supporting-plane argument applies to a genuinely convex function even when the original function is presented through a strong-convexity assumption. The two-point conclusions concern arbitrary points of the whole space; no minimizer or second derivative is needed.

\subsection{\texorpdfstring{The supporting-\allowbreak{}plane inequality}{The supporting-plane inequality}}


Let \(H\) be a real inner-product space and let \(f:H\to\mathbb R\) be differentiable and convex on all of \(H\). For any \(x,y\in H\),

\[f(x)+Df(x)[y-x]\le f(y).\]

Restrict \(f\) to the affine line \(\ell(t)=x+t(y-x)\). Since affine composition preserves convexity, \(\varphi=f\circ\ell\) is convex on \(\mathbb R\), with \(\varphi(0)=f(x)\) and \(\varphi(1)=f(y)\).



The derivative of the line is its constant direction:

\[\ell'(0)=y-x.\]

By the chain rule, the derivative of the restricted function is

\[\varphi'(0)=Df(x)[y-x].\]

For a differentiable convex function, the derivative at the left endpoint is bounded by the secant slope to a later point. Applying this at \(0<1\) gives

\[Df(x)[y-x]\le\frac{\varphi(1)-\varphi(0)}{1-0}.\]

Evaluate the two endpoints and rearrange:

\[Df(x)[y-x]\le f(y)-f(x),\qquad f(x)+Df(x)[y-x]\le f(y).\]

The line argument uses only values on the segment joining the two points. For \(0<t\le1\), convexity gives \((\varphi(t)-\varphi(0))/t\le\varphi(1)-\varphi(0)\). Sending \(t\downarrow0\) identifies the left side with the directional derivative. This also explains why the derivative is evaluated at the base point rather than at an unspecified point on the segment.

\subsection{\texorpdfstring{Strong monotonicity of the gradient}{Strong monotonicity of the gradient}}


Let \(H\) be a real Hilbert space. If \(f:H\to\mathbb R\) is differentiable and \(m\)-strongly convex, where \(m\in\mathbb R\), then for every \(x,y\in H\),

\[m\|x-y\|^2\le\langle\nabla f(x)-\nabla f(y),x-y\rangle.\]

Remove the strong-convexity quadratic. Set

\[g(v)=f(v)-\frac m2\|v\|^2.\]

The definition of strong convexity says precisely that \(g\) is convex.



The squared norm has derivative \(D\|v\|^2[u]=2\langle v,u\rangle\). Consequently \(g\) is differentiable and

\[Dg(v)[u]=Df(v)[u]-m\langle v,u\rangle.\]

Apply the supporting-plane inequality in both directions, and substitute the derivative just computed:

\[g(x)+Dg(x)[y-x]\le g(y),\qquad g(y)+Dg(y)[x-y]\le g(x).\]

Express the derivative of \(f\) using its gradient. The two inequalities become

\[\begin{aligned}f(x)-\tfrac m2\|x\|^2+\langle\nabla f(x),y-x\rangle-m\langle x,y-x\rangle&\le f(y)-\tfrac m2\|y\|^2,\\ f(y)-\tfrac m2\|y\|^2+\langle\nabla f(y),x-y\rangle-m\langle y,x-y\rangle&\le f(x)-\tfrac m2\|x\|^2.\end{aligned}\]

Add these inequalities. The function values and quadratic values cancel. Using \(\langle x-y,x-y\rangle=\|x-y\|^2\), what remains is

\[0\le\langle\nabla f(x)-\nabla f(y),x-y\rangle-m\|x-y\|^2.\]

This is the required monotonicity inequality.



Reversing the pair changes \(y-x\) into \(-(x-y)\). Therefore adding the two supporting inequalities makes the function values cancel and leaves a gradient difference paired with a displacement of the same orientation. The sign of this pairing is essential: it gives monotonicity of \(\nabla f\), rather than the opposite inequality. No Hessian is introduced.

\subsection{\texorpdfstring{The quadratic lower model}{The quadratic lower model}}


Let \(H\) be a real Hilbert space, \(f:H\to\mathbb R\) differentiable, and \(m\in\mathbb R\). If \(f\) is \(m\)-strongly convex on all of \(H\), then for every \(x,y\in H\),

\[f(x)+\langle\nabla f(x),y-x\rangle+\frac m2\|y-x\|^2\le f(y).\]

Set \(g(v)=f(v)-\frac m2\|v\|^2\). Strong convexity makes \(g\) convex.



Differentiate this shifted potential:

\[Dg(v)[u]=Df(v)[u]-m\langle v,u\rangle.\]

In particular, \(g\) is differentiable on \(H\).



Its supporting-plane inequality at \(x\), together with \(Df(x)[u]=\langle\nabla f(x),u\rangle\), yields

\[f(x)-\frac m2\|x\|^2+\langle\nabla f(x),y-x\rangle-m\langle x,y-x\rangle\le f(y)-\frac m2\|y\|^2.\]

Collect the quadratic correction by the identity

\[\|y\|^2-\|x\|^2-2\langle x,y-x\rangle=\|y-x\|^2.\]

Moving terms gives the desired lower model.



The correction is the squared displacement, not the difference of the two squared norms. The cross term from differentiating the subtracted quadratic is exactly what converts that difference into \(\|y-x\|^2\). When \(\nabla f(x)=0\), this inequality specializes to \(f(y)-f(x)\ge m\|y-x\|^2/2\), explaining its later use for Gaussian tail bounds about a minimizer.

\section{\texorpdfstring{A first Euler-\allowbreak{}step estimate}{A first Euler-step estimate}}
\label{proof-module-UnadjustedLangevin-GradientStep-lean}


Let \(d=x-y\) and \(v=g_x-g_y\). The identity \(\|d-hv\|^2=\|d\|^2-2h\langle v,d\rangle+h^2\|v\|^2\) combines strong monotonicity and Lipschitz control into the elementary Euler bound. The vectors \(g_x,g_y\) may be arbitrary vectors satisfying these inequalities; substitution of the potential's gradients is a later specialization.



There are two competing terms in the Euler displacement: the inner product with the gradient difference decreases the squared distance by \(2hm\|x-y\|^2\), whereas the squared drift increment increases it by at most \(h^2L^2\|x-y\|^2\). Nonnegativity of \(h\) is essential when inserting the monotonicity estimate into the negative cross term. If \(x=y\), Lipschitz continuity forces the gradient difference to vanish, so the zero-distance case requires no division.

\subsection{\texorpdfstring{The squared Euler-\allowbreak{}step bound}{The squared Euler-step bound}}


Let \(x,y,g_x,g_y\) lie in a real inner-product space, let \(m\in\mathbb R\), \(L\ge0\), and \(h\ge0\). Assume

\[m\|x-y\|^2\le\langle g_x-g_y,x-y\rangle,\qquad\|g_x-g_y\|\le L\|x-y\|.\]

Then

\[\|(x-hg_x)-(y-hg_y)\|^2\le(1-2hm+h^2L^2)\|x-y\|^2.\]

Both sides of the Lipschitz estimate are nonnegative, so squaring preserves the inequality:

\[\|g_x-g_y\|^2\le L^2\|x-y\|^2.\]

Multiply strong monotonicity by \(2h\ge0\), and multiply the squared Lipschitz estimate by \(h^2\ge0\):

\[2hm\|x-y\|^2\le2h\langle g_x-g_y,x-y\rangle,\qquad h^2\|g_x-g_y\|^2\le h^2L^2\|x-y\|^2.\]

Write \(d=x-y\), \(v=g_x-g_y\). Expanding the squared norm and applying the preceding estimates gives

\[\begin{aligned}\|(x-hg_x)-(y-hg_y)\|^2&=\|d-hv\|^2\\&=\|d\|^2-2h\langle v,d\rangle+h^2\|v\|^2\\&\le(1-2hm+h^2L^2)\|d\|^2.\end{aligned}\]

The expansion retains the negative mixed term until strong monotonicity has been used. Bounding that term by its absolute value would discard the dissipative contribution and leave only \((1+hL)^2\). The sharper expression \(1-2hm+h^2L^2\) therefore comes from the orientation of the gradient, while the positive quadratic term measures the cost of explicit discretization.

\subsection{\texorpdfstring{Taking the square root}{Taking the square root}}


Let \(H\) be a real inner-product space, \(x,y,g_x,g_y\in H\), \(m\in\mathbb R\), \(L\ge0\), and \(h\ge0\). Assume

\[m\|x-y\|^2\le\langle g_x-g_y,x-y\rangle,\qquad\|g_x-g_y\|\le L\|x-y\|.\]

Then

\[\|(x-hg_x)-(y-hg_y)\|\le\sqrt{1-2hm+h^2L^2}\,\|x-y\|.\]

No separate sign assumption on \(1-2hm+h^2L^2\) is required; the real square root is interpreted as zero on nonpositive inputs.



Apply the preceding squared estimate. With \(A=\|(x-hg_x)-(y-hg_y)\|\), \(r=\|x-y\|\), and \(c=1-2hm+h^2L^2\), it says \(A^2\le cr^2\).



Since \(A\ge0\), take square roots and use \(r^2\ge0\):

\[A\le\sqrt{cr^2}=\sqrt c\,\sqrt{r^2}=\sqrt c\,r.\]

The square-root product identity is valid here even without a sign hypothesis on \(c\), because the second factor \(r^2\) is nonnegative.



If \(r=0\), the Lipschitz estimate implies \(g_x=g_y\), and both sides vanish. If \(r>0\), the squared estimate forces \(c\ge0\), so ordinary square-root monotonicity applies. Thus the absence of a separate sign premise is harmless: a negative coefficient can only occur together with the zero-distance case under the displayed assumptions.

\section{\texorpdfstring{The Langevin update and cancellation under shared noise}{The Langevin update and cancellation under shared noise}}
\label{proof-module-UnadjustedLangevin-ULAStep-lean}


We now work in finite-dimensional Euclidean space. A ULA step adds the same scaled noise vector to a gradient step. When two starting points use this same noise, it cancels exactly from their difference. The deterministic estimates can therefore be applied pointwise, before introducing any probability law.



the Gaussian Euler update is the pointwise Euler map. Its definition accepts arbitrary vector noise; Gaussian laws enter later in realization theorems. the common-noise cancellation identity uses the same innovation on both sides, so that contribution cancels exactly. The monotonicity of the actual gradient is supplied by the strong-convexity theorem when applying the deterministic step bound.



Here the noise vector is an input to a deterministic map. Gaussianity and independence enter only when laws of successive steps are considered. Consequently a comparison with a common noise vector holds for every realization and may later be integrated under any joint distribution of the initial pair and that noise. The diffusion coefficient \(\sqrt{2h}\) corresponds to the continuous generator \(\Delta-\langle\nabla U,\nabla\rangle\).

\subsection{\texorpdfstring{The Euclidean state space}{The Euclidean state space}}


For a natural number \(d\), write \(E_d=\mathbb R^d\) with its Euclidean inner product and norm. Vectors may be represented as real-valued functions on the finite index set \(\{0,\ldots,d-1\}\); when \(d=0\), this index set is empty.



Introduce the abbreviation

\[E_d:=\mathbb R^d.\]

This equips the coordinate functions with the Euclidean norm \(\|x\|=(\sum_{i=0}^{d-1}x_i^2)^{1/2}\), so the subsequent gradient and distance estimates use the usual geometry of \(\mathbb R^d\).



The coordinate expression for the inner product is \(\langle x,y\rangle=\sum_{i=0}^{d-1}x_iy_i\). In particular, coordinatewise addition and scalar multiplication agree with the Euclidean vector operations used by the update. When \(d=0\), the empty coordinate list gives one state and a zero norm; later nondegenerate arguments separate this case explicitly.

\subsection{\texorpdfstring{One unadjusted Langevin step}{One unadjusted Langevin step}}


For a potential \(U:E_d\to\mathbb R\), a real step size \(h\), a state \(x\in E_d\), and a noise vector \(z\in E_d\), define

\[T_h^z(x)=x-h\nabla U(x)+\sqrt{2h}\,z.\]

The definition itself imposes no sign condition on \(h\); later contraction results assume \(h\ge0\).



The update consists of the current state, the drift increment \(-h\nabla U(x)\), and the noise increment \(\sqrt{2h}\,z\).



For \(h>0\) and standard Gaussian \(z\), the conditional covariance of the noise increment is \(2hI_d\). The deterministic center is \(x-h\nabla U(x)\), so the update is the Euler discretization of drift \(-\nabla U\) and diffusion \(\sqrt2\). The definition remains a deterministic function until a law for its inputs is specified.

\subsection{\texorpdfstring{Shared noise cancels}{Shared noise cancels}}


For every potential \(U:E_d\to\mathbb R\), \(h\in\mathbb R\), and vectors \(x,y,z\in E_d\),

\[T_h^z(x)-T_h^z(y)=(x-h\nabla U(x))-(y-h\nabla U(y)).\]

Expand the two updates. The two copies of \(\sqrt{2h}\,z\) subtract to zero:

\[(x-h\nabla U(x)+\sqrt{2h}\,z)-(y-h\nabla U(y)+\sqrt{2h}\,z)=(x-h\nabla U(x))-(y-h\nabla U(y)).\]

The cancellation uses the same realization of \(z\), rather than merely two noises with the same distribution. If the noises differed, an additional term \(\sqrt{2h}(z-z')\) would remain and would need its own moment estimate. Synchronous coupling deliberately removes that term, allowing the deterministic gradient geometry to control the whole update difference.

\subsection{\texorpdfstring{The initial shared-\allowbreak{}noise estimate}{The initial shared-noise estimate}}


Suppose \(U:E_d\to\mathbb R\) is differentiable and \(m\)-strongly convex, \(\nabla U\) is \(L\)-Lipschitz with \(L\ge0\), and \(h\ge0\). For every \(x,y,z\in E_d\),

\[\|T_h^z(x)-T_h^z(y)\|\le\sqrt{1-2hm+h^2L^2}\,\|x-y\|.\]

For \(h>0\) this factor is strictly below one whenever \(hL^2<2m\).



Cancel the common noise. The left-hand side becomes the distance between the two deterministic gradient steps.



Strong convexity gives

\[m\|x-y\|^2\le\langle\nabla U(x)-\nabla U(y),x-y\rangle,\]

and the Lipschitz assumption gives \(\|\nabla U(x)-\nabla U(y)\|\le L\|x-y\|\). Insert these two inequalities into the Euler-step square-root bound.



The coefficient is independent of the noise distribution because cancellation occurs before taking a norm. Hence the inequality holds pointwise and may be integrated against any coupled starting law. A strict contraction requires \(1-2hm+h^2L^2<1\); for \(h>0\), this is the stated comparison of the linear dissipative contribution with the quadratic discretization contribution.

\section{\texorpdfstring{The sharp deterministic contraction}{The sharp deterministic contraction}}
\label{proof-module-UnadjustedLangevin-SharpContraction-lean}


The preceding square-root factor uses strong monotonicity and Lipschitz continuity separately. The quadratic lower and upper models contain more information. We combine them into the interpolation inequality, then interpret that inequality as a ball containing the gradient difference. This yields the two sharp factors \(1-mh\) and \(Lh-1\), with the change of regime at \(h=2/(m+L)\).



The proof separates the two regimes for the maximum contraction factor \(\max(|1-mh|,|1-Lh|)\). Each scalar bound follows after the strong monotonicity and cocoercivity estimates have supplied the needed inner-product control. The final norm theorem removes the square by nonnegativity, retaining the precise step-regime hypotheses.



The strong and smooth quadratic models constrain the gradient difference more sharply than separate monotonicity and Lipschitz inequalities do. For \(d=x-y\) and \(v=\nabla U(x)-\nabla U(y)\), they locate \(v\) in the ball centered at \((m+L)d/2\) with radius \((L-m)\|d\|/2\). Applying \(v\mapsto d-hv\) to this ball explains the transition at \(h=2/(m+L)\). The two regimes meet at the common factor \((L-m)/(L+m)\).

\subsection{\texorpdfstring{Cocoercivity from two quadratic models}{Cocoercivity from two quadratic models}}


Let \(H\) be a real inner-product space, \(f:H\to\mathbb R\), \(g:H\to H\), and \(K\ge0\). Suppose that for all \(a,b\in H\),

\[f(a)+\langle g(a),b-a\rangle\le f(b)\le f(a)+\langle g(a),b-a\rangle+\frac K2\|b-a\|^2.\]

Then for all \(x,y\),

\[\|g(x)-g(y)\|^2\le K\langle g(x)-g(y),x-y\rangle.\]

The proof includes \(K=0\), and does not separately assume that \(g\) is a gradient.



Fix \(a,b\), \(t\ge0\), and set \(v=g(a)-g(b)\), \(z=a-tv\). Compare the lower model based at \(b\) with the upper model based at \(a\), both evaluated at \(z\):

\[f(b)+\langle g(b),z-b\rangle\le f(z)\le f(a)+\langle g(a),z-a\rangle+\frac K2\|z-a\|^2.\]

The displacements are \(z-a=-tv\) and \(z-b=(a-b)-tv\). Substitute them into the two models. Since \(t\ge0\), the quadratic term is \(\|z-a\|^2=t^2\|v\|^2\).



The difference of the linear terms satisfies

\[\langle g(a),v\rangle-\langle g(b),v\rangle=\|v\|^2.\]

Eliminating \(f(z)\) therefore gives

\[\left(t-\frac K2t^2\right)\|g(a)-g(b)\|^2\le f(a)-f(b)-\langle g(b),a-b\rangle.\]

Apply this estimate to \((a,b)=(x,y)\) and \((y,x)\), with the same \(t\ge0\). Adding will cancel the function values.



The two gradient norms are equal and the reversed displacement changes the sign of its inner product. Thus, writing \(A=\|g(x)-g(y)\|^2\) and \(B=\langle g(x)-g(y),x-y\rangle\), we have for every \(t\ge0\),

\[(2t-Kt^2)A\le B.\]

First consider \(K=0\). Taking \(t=1\) gives \(2A\le B\). As \(A\ge0\), this also shows \(B\ge0\).



If the desired conclusion \(A\le0\) were false, then \(A>0\). We could choose

\[t=\frac{B+1}{2A}\ge0.\]

The estimate \(2tA\le B\) would then force \(B+1\le B\).



To justify that substitution, \(A>0\) implies \(\|g(x)-g(y)\|\ne0\), so its square can be cancelled from the denominator. Hence indeed

\[2\frac{B+1}{2A}A=B+1.\]

The contradiction \(B+1\le B\) excludes \(A>0\). Thus \(A=0\), proving the claim when \(K=0\).



Now suppose \(K>0\). Choose \(t=K^{-1}\) in the common estimate:

\[\frac AK\le B.\]

Multiplying by \(K\) gives \(A\le KB\), as required.



For positive \(K\), the choice \(t=1/K\) maximizes the quadratic coefficient \(2t-Kt^2\). This is why the reciprocal smoothness constant enters the estimate. At \(K=0\), no reciprocal is available; allowing arbitrarily large nonnegative \(t\) instead forces the squared gradient difference to vanish. The two cases together give a boundary-valid statement.

\subsection{\texorpdfstring{The strong-\allowbreak{}smooth interpolation inequality}{The strong-smooth interpolation inequality}}


Let \(H\) be a real Hilbert space, \(f:H\to\mathbb R\) differentiable and \(m\)-strongly convex, and let \(\nabla f\) be \(L\)-Lipschitz, where \(L\ge0\) and \(m\le L\). Then for every \(x,y\),

\[mL\|x-y\|^2+\|\nabla f(x)-\nabla f(y)\|^2\le(m+L)\langle\nabla f(x)-\nabla f(y),x-y\rangle.\]

The inequality holds also at the boundary \(m=L\).



Subtract the strong-convexity quadratic and its gradient:

\[F(v)=f(v)-\frac m2\|v\|^2,\qquad G(v)=\nabla f(v)-mv.\]

The quadratic lower model for \(f\) becomes an ordinary supporting plane for \(F\):

\[F(a)+\langle G(a),b-a\rangle\le F(b).\]

This follows by expanding the squared distance and collecting the terms involving \(m\).



Apply the smooth upper model to \(f\), and substitute the same definitions of \(F\) and \(G\).



Expanding the quadratic terms shows that the upper curvature decreases from \(L\) to \(L-m\):

\[F(b)\le F(a)+\langle G(a),b-a\rangle+\frac{L-m}{2}\|b-a\|^2.\]

Since \(L-m\ge0\), cocoercivity for these two models gives

\[\|G(x)-G(y)\|^2\le(L-m)\langle G(x)-G(y),x-y\rangle.\]

Write \(d=x-y\) and \(v=\nabla f(x)-\nabla f(y)\). The shifted gradient difference is exactly

\[G(x)-G(y)=v-md.\]

Expand the cocoercivity inequality using this identity:

\[\|v\|^2-2m\langle v,d\rangle+m^2\|d\|^2\le(L-m)\bigl(\langle v,d\rangle-m\langle d,d\rangle\bigr).\]

Use \(\langle d,d\rangle=\|d\|^2\) and collect terms. This gives

\[\|v\|^2+mL\|d\|^2\le(m+L)\langle v,d\rangle.\]

When \(m=L\), the cocoercivity constant above is zero, forcing \(G(x)=G(y)\); consequently \(\nabla f(x)-\nabla f(y)=m(x-y)\). Thus no division by \(L-m\) or limiting argument is needed.



The shifted upper curvature \(L-m\) is nonnegative by assumption. Expanding \(\|v-md\|^2\le(L-m)\langle v-md,d\rangle\) gives coefficients \(-2m\) and \(-(L-m)\) on the mixed term; collecting them produces \(m+L\). The coefficient of \(\|d\|^2\) becomes \(mL\), which is the extra information missing from separate monotonicity and Lipschitz bounds.

\subsection{\texorpdfstring{The gradient difference lies in a ball}{The gradient difference lies in a ball}}


Let \(H\) be a real inner-product space, \(x,y,g_x,g_y\in H\), and \(m,L\in\mathbb R\) with \(m\le L\). Suppose

\[mL\|x-y\|^2+\|g_x-g_y\|^2\le(m+L)\langle g_x-g_y,x-y\rangle.\]

Then

\[\left\|(g_x-g_y)-\frac{m+L}{2}(x-y)\right\|\le\frac{L-m}{2}\|x-y\|.\]

The proposed radius is nonnegative because \(m\le L\). Both sides of the desired norm bound are therefore nonnegative, so it is enough to prove the corresponding squared inequality.



Set \(c=(m+L)/2\), \(r=(L-m)/2\), \(d=x-y\), and \(v=g_x-g_y\). Completing the square in the interpolation inequality gives

\[\begin{aligned}\|v-cd\|^2&=\|v\|^2-(m+L)\langle v,d\rangle+c^2\|d\|^2\\&\le(c^2-mL)\|d\|^2=r^2\|d\|^2.\end{aligned}\]

Taking nonnegative square roots proves the claim. If \(m=L\), the radius is zero and \(v=md\).



The algebra uses \(((m+L)/2)^2-mL=((L-m)/2)^2\). Thus the interpolation inequality describes a closed ball in the gradient-difference variable for each fixed displacement. Its radius vanishes at \(m=L\), identifying the gradient difference exactly with \(m(x-y)\). This geometric form avoids dividing by either curvature difference or point separation.

\subsection{\texorpdfstring{The centered Euler-\allowbreak{}step bound}{The centered Euler-step bound}}


Let \(H\) be a real inner-product space, \(m,L,h\in\mathbb R\) with \(m\le L\) and \(h\ge0\), and let \(x,y,g_x,g_y\in H\) satisfy

\[mL\|x-y\|^2+\|g_x-g_y\|^2\le(m+L)\langle g_x-g_y,x-y\rangle.\]

Then

\[\|(x-hg_x)-(y-hg_y)\|\le\left(\left|1-h\frac{m+L}{2}\right|+h\frac{L-m}{2}\right)\|x-y\|.\]

The preceding ball bound gives \(\|v-cd\|\le r\|d\|\), where \(d=x-y\), \(v=g_x-g_y\), \(c=(m+L)/2\), and \(r=(L-m)/2\).



Decompose the Euler difference into its centered part and its deviation:

\[d-hv=(1-hc)d-h(v-cd).\]

The triangle inequality and homogeneity of the norm give

\[\|d-hv\|\le|1-hc|\,\|d\|+h\|v-cd\|,\]

where \(|-h|=h\) because \(h\ge0\).



Multiply the ball bound by \(h\ge0\) and substitute it:

\[\|d-hv\|\le\bigl(|1-hc|+hr\bigr)\|d\|.\]

For \(m=L\), this reduces to the exact scalar factor \(|1-mh|\).



The two summands in the decomposition are parallel to the displacement and the deviation from the ball center, respectively. Their norm bound is the distance of the affine image of that ball from the origin. Because \(h\ge0\), scaling its radius contributes \(h(L-m)/2\). Keeping the absolute value on the center yields one expression valid on both sides of the step threshold.

\subsection{\texorpdfstring{The sharp small-\allowbreak{}step factor}{The sharp small-step factor}}


Let \(H\) be a real inner-product space, \(x,y,g_x,g_y\in H\), and \(m,L,h\in\mathbb R\). Suppose \(m\le L\), \(m+L>0\), \(0\le h\le2/(m+L)\), and

\[mL\|x-y\|^2+\|g_x-g_y\|^2\le(m+L)\langle g_x-g_y,x-y\rangle.\]

Then

\[\|(x-hg_x)-(y-hg_y)\|\le(1-mh)\|x-y\|.\]

Apply the centered bound. The step-size assumption and \(m+L>0\) imply \(h(m+L)\le2\), hence \(1-h(m+L)/2\ge0\).



Remove the absolute value with this sign and simplify its coefficient:

\[\left|1-h\frac{m+L}{2}\right|+h\frac{L-m}{2}=1-h\frac{m+L}{2}+h\frac{L-m}{2}=1-mh.\]

The threshold ensures that the center coefficient is nonnegative. Moreover \(m\le L\) implies \(mh\le h(m+L)/2\le1\), so the simplified factor \(1-mh\) is itself nonnegative. This sign check makes the resulting pointwise estimate suitable for integration into an \(L^2\) inequality without any absolute value adjustment.

\subsection{\texorpdfstring{The sharp large-\allowbreak{}step factor}{The sharp large-step factor}}


Let \(H\) be a real inner-product space, \(x,y,g_x,g_y\in H\), and \(m,L,h\in\mathbb R\). Suppose \(m\le L\), \(m+L>0\), \(h\ge0\), \(2/(m+L)\le h\), and

\[mL\|x-y\|^2+\|g_x-g_y\|^2\le(m+L)\langle g_x-g_y,x-y\rangle.\]

Then

\[\|(x-hg_x)-(y-hg_y)\|\le(Lh-1)\|x-y\|.\]

This is a distance bound; it is a strict contraction only when its coefficient is below one.



Apply the centered estimate. Now the step-size assumption gives \(2\le h(m+L)\), so \(1-h(m+L)/2\le0\).



Replace the absolute value by the negative of its argument. The coefficient becomes

\[\left|1-h\frac{m+L}{2}\right|+h\frac{L-m}{2}=h\frac{m+L}{2}-1+h\frac{L-m}{2}=Lh-1.\]

Above the threshold the center has crossed the origin, so its absolute value changes branch. The factor \(Lh-1\) is nonnegative because \(h(m+L)\ge2\) and \(m\le L\) imply \(Lh\ge1\). Its being below one is the separate stability condition \(Lh<2\), which is needed later to sum the geometric error recurrence.

\subsection{\texorpdfstring{Small-\allowbreak{}step contraction for a potential}{Small-step contraction for a potential}}


Let \(f:H\to\mathbb R\) be differentiable and \(m\)-strongly convex on a real Hilbert space, with an \(L\)-Lipschitz gradient, \(L\ge0\), \(m\le L\), and \(m+L>0\). If \(0\le h\le2/(m+L)\), then for every \(x,y\),

\[\|(x-h\nabla f(x))-(y-h\nabla f(y))\|\le(1-mh)\|x-y\|.\]

The quadratic models imply the strong-smooth interpolation inequality for \(g_x=\nabla f(x)\) and \(g_y=\nabla f(y)\). Apply the small-step Euler bound to these two vectors. Its hypotheses are exactly the stated curvature and step-size assumptions.



Here the two vectors inserted into the abstract estimate are the actual gradients, so their interpolation inequality follows from the potential's lower and upper models. No additional relation between arbitrary vectors is assumed. The estimate works uniformly over every pair \(x,y\), which is what permits the same contraction constant to act on random states.

\subsection{\texorpdfstring{Large-\allowbreak{}step contraction for a potential}{Large-step contraction for a potential}}


Let \(H\) be a real Hilbert space and \(f:H\to\mathbb R\) differentiable and \(m\)-strongly convex, with an \(L\)-Lipschitz gradient. Suppose \(L\ge0\), \(m\le L\), \(m+L>0\), \(h\ge0\), and \(2/(m+L)\le h\). Then for every \(x,y\in H\),

\[\|(x-h\nabla f(x))-(y-h\nabla f(y))\|\le(Lh-1)\|x-y\|.\]

Use the potential's interpolation inequality in the large-step Euler bound. This gives the second sharp factor \(Lh-1\).



The potential assumptions verify the geometric interpolation inequality before the large-step algebra is used. The conclusion is a bound for every \(h\) in this branch, including unstable values with \(Lh\ge2\). Only when applying the bound as a contraction in a convergence theorem is the upper restriction \(Lh<2\) imposed.

\subsection{\texorpdfstring{The endpoint-\allowbreak{}factor form}{The endpoint-factor form}}


Let \(f\) be differentiable and \(m\)-strongly convex on a real Hilbert space, with an \(L\)-Lipschitz gradient, \(L\ge0\), \(m\le L\), and \(m+L>0\). For every \(h\ge0\) and every \(x,y\),

\[\|(x-h\nabla f(x))-(y-h\nabla f(y))\|\le\max\{|1-mh|,|1-Lh|\}\,\|x-y\|.\]

If \(0<m\le L\), this factor is below one for \(0<h<2/L\). At \(m=L\), the two endpoint factors coincide.



Split according to whether \(h\le2/(m+L)\). In the first case the small-step theorem bounds the distance by \((1-mh)\|x-y\|\).



Since \(1-mh\le|1-mh|\le\max\{|1-mh|,|1-Lh|\}\), multiplying by \(\|x-y\|\ge0\) gives the desired bound in this case.



In the other case \(h>2/(m+L)\), so the large-step theorem bounds the distance by \((Lh-1)\|x-y\|\).



The elementary absolute-value inequality yields

\[Lh-1=-(1-Lh)\le|1-Lh|.\]

Finally \(|1-Lh|\le\max\{|1-mh|,|1-Lh|\}\). Multiply by the nonnegative distance to finish the second case.



The maximum records the two endpoint curvatures of the admissible interval. In the small-step branch it bounds \(1-mh\), while in the large-step branch it bounds \(Lh-1\). For \(0<m\le L\), choosing \(h=2/(m+L)\) balances the two endpoint factors and gives \((L-m)/(L+m)\), explaining the common transition point.

\subsection{\texorpdfstring{Sharp shared-\allowbreak{}noise contraction: small steps}{Sharp shared-noise contraction: small steps}}


Let \(f:E_d\to\mathbb R\) be differentiable and \(m\)-strongly convex, with an \(L\)-Lipschitz gradient, \(L\ge0\), \(m\le L\), and \(m+L>0\). If \(0\le h\le2/(m+L)\), then for all starting states \(x,y\) and every common noise vector \(z\),

\[\|T_h^z(x)-T_h^z(y)\|\le(1-mh)\|x-y\|.\]

The common noise cancels exactly, leaving the distance between the gradient steps from \(x\) and \(y\).



Apply the potential-based small-step estimate. It gives the factor \(1-mh\) independently of the value or distribution of \(z\).



Apply cancellation first and then the deterministic sharp estimate. Since the resulting coefficient contains no \(z\), its square integrates under any joint law of the starting pair and noise. Thus the sharp deterministic rate survives synchronous randomization exactly; the Gaussian law is needed for the algorithm's marginal transition, not for this comparison.

\subsection{\texorpdfstring{Sharp shared-\allowbreak{}noise contraction: large steps}{Sharp shared-noise contraction: large steps}}


Let \(f:E_d\to\mathbb R\) be differentiable and \(m\)-strongly convex, with an \(L\)-Lipschitz gradient, \(L\ge0\), \(m\le L\), and \(m+L>0\). If \(h\ge0\) and \(2/(m+L)\le h\), then for every \(x,y,z\in E_d\),

\[\|T_h^z(x)-T_h^z(y)\|\le(Lh-1)\|x-y\|.\]

Cancel the shared noise, reducing the two ULA updates to their deterministic gradient-step difference.



Apply the potential-based large-step estimate. This yields the factor \(Lh-1\) for every common noise vector.



The same pointwise cancellation applies in the second regime, but the factor now increases with the step. This is why the eventual bias bound contains the denominator \(2-Lh\): it is precisely \(1-(Lh-1)\). As \(Lh\) approaches two, contraction is lost and the geometric accumulation of local defects can no longer be bounded uniformly.

\section{\texorpdfstring{The local Euler defect}{The local Euler defect}}
\label{proof-module-UnadjustedLangevin-LocalError-lean}


Let \(H\) be a real Banach space, with norm \(\|\cdot\|\). All Lipschitz constants \(L\) below are nonnegative. All time integrals are Bochner integrals. For the Langevin application we will take \(H=L^2(P;\mathbb R^d)\); the norms below then measure random variables in mean square. The scalar defect estimate permits a nonnegative real parameter \(d\); its application uses the natural-number dimension, viewed as a real number. Begin with a Lipschitz drift \(g\), a reference curve \(Y\), and its additive noise \(B\). We first derive the error of freezing the drift over one time step directly from the integral equation.



The analytic calculation is conditional on the displayed integral equation and norm bounds. It neither constructs a stationary Langevin diffusion nor proves its score moment. The Gaussian and Brownian noise moments used later are proved separately.



The local defect is a vector-valued Bochner integral of drift variation from time zero to time \(h\). Continuity or the stated integrability assumptions justify its existence and all subsequent norm and linearity operations. The integration is oriented when its endpoints are reversed; the error estimates use \(h\ge0\).



The defect measures the accumulated variation of the drift during one reference segment. Its sign follows from \(Y_h=Y_0-hg(Y_0)+B_h-D_h\): the Euler step omits \(D_h\), so subtracting the reference from the numerical update adds the defect. The drift part of the increment is of order \(t\), and the Brownian part is of order \(\sqrt t\). Their time integrals have orders \(h^2\) and \(h^{3/2}\), which is the origin of the eventual square-root-in-step bias after summing the recurrence.

\subsection{\texorpdfstring{Definition of the local defect}{Definition of the local defect}}


For \(g:H\to H\), \(Y:\mathbb R\to H\), and \(h\in\mathbb R\), define the drift omitted when Euler freezes \(g\) at the initial point:

\[D_h(g,Y):=\int_0^h\bigl(g(Y_t)-g(Y_0)\bigr)\,dt.\]

Integrate the difference between the evolving drift and its frozen value. This defines the defect itself; the estimates below will be consequences of this definition.



The integrand vanishes at time zero, so this defect measures drift variation rather than the drift magnitude itself. In a space of random variables, the integral is a Bochner integral in that space and its norm is an \(L^2\) error. Defining the defect in this form keeps the common additive noise outside the remainder.

\subsection{\texorpdfstring{Separate the exact and frozen drift}{Separate the exact and frozen drift}}


If \(t\mapsto g(Y_t)\) is integrable on the interval between \(0\) and \(h\), then

\[D_h(g,Y)=\int_0^h g(Y_t)\,dt-h\,g(Y_0).\]

Expand the definition. Both the actual drift and the constant function \(g(Y_0)\) are integrable, so linearity separates their difference into two integrals.



The integral of the constant vector is \((h-0)g(Y_0)=h\,g(Y_0)\), yielding the identity.



Subtracting the constant inside the integral is legitimate because a constant vector is integrable on a finite interval. For negative \(h\), the oriented integral changes sign consistently on both terms. Later error estimates assume \(h\ge0\) so that time-dependent upper bounds can be integrated in their usual order.

\subsection{\texorpdfstring{An exact step is Euler minus its defect}{An exact step is Euler minus its defect}}


Assume the preceding drift integrability and the endpoint integral equation

\[Y_h-Y_0=-\int_0^h g(Y_t)\,dt+B_h.\]

Then the exact reference step has the form

\[Y_h=Y_0-h\,g(Y_0)+B_h-D_h(g,Y).\]

Solve the assumed integral equation for \(Y_h\), then replace the defect by the difference between the exact drift integral and its frozen approximation.



Cancel the two frozen-drift terms. The resulting equality is precisely the Euler update with the signed remainder \(-D_h\).



Substitute \(\int_0^hg(Y_t)dt=hg(Y_0)+D_h\) into the endpoint equation. The minus sign in front of the original drift integral gives \(-hg(Y_0)-D_h\). This sign convention is retained when comparing the numerical and exact updates, where subtracting the reference changes \(-D_h\) into \(+D_h\).

\subsection{\texorpdfstring{Continuity supplies drift integrability}{Continuity supplies drift integrability}}


Let \(H\) be a normed additive group, \(g:H\to H\) be \(L\)-Lipschitz with \(L\ge0\), and \(Y:\mathbb R\to H\) be continuous. For every real \(h\), the function \(t\mapsto g(Y_t)\) is interval integrable from \(0\) to \(h\).



A Lipschitz map is continuous, so \(g\circ Y\) is continuous. A continuous function is integrable on every compact interval, proving the claim.



The image of a compact time interval under a continuous path is compact. A continuous drift is bounded on that image, so the composition has finite integral norm there. In the lifted random-vector setting the same argument uses continuity in \(L^2\), not a pathwise boundedness claim about every sample trajectory.

\subsection{\texorpdfstring{The square-\allowbreak{}root integral}{The square-root integral}}


For every \(h\ge0\),

\[\int_0^h\sqrt t\,dt=\frac23h\sqrt h.\]

Write \(\sqrt t=t^{1/2}\) and use the power-integral formula. The exponent satisfies \(1/2>-1\), so the formula applies at the zero endpoint.



The contribution at zero vanishes. Since \(h\ge0\), the other endpoint is \(h^{1/2+1}=h^{1/2}h=\sqrt h\,h\).



Divide by \(1/2+1=3/2\) and rearrange the factors to obtain \((2/3)h\sqrt h\).



The integrand has an integrable singularity in its derivative at zero, but the function \(\sqrt t\) itself is continuous there. The primitive \((2/3)t^{3/2}\) vanishes at zero. Consequently the Brownian contribution to an integrated Lipschitz increment is exactly proportional to \(h^{3/2}\), with the constant \(2/3\) retained for the final numerical estimate.

\subsection{\texorpdfstring{Bound a reference increment}{Bound a reference increment}}


Let \(h\ge0\) and suppose that, for every \(0\le s\le h\),

\[\|g(Y_s)\|\le G,\qquad \|B_s\|\le N\sqrt s,\qquad Y_s-Y_0=-\int_0^s g(Y_r)\,dr+B_s.\]

Then, for each \(0\le t\le h\),

\[\|Y_t-Y_0\|\le tG+N\sqrt t.\]

First bound the drift integral on \([0,t]\) by its constant norm bound:

\[\left\|\int_0^t g(Y_s)\,ds\right\|\le G\,|t-0|=Gt,\]

where \(t\ge0\) removes the absolute value.



To justify the constant bound under that integral, note that \(s\in(0,t]\) implies \(0\le s\le t\le h\). Hence the assumed bound on \([0,h]\) applies to every such \(s\).



Insert the integral equation and apply the triangle inequality:

\[\|Y_t-Y_0\|\le\left\|\int_0^t g(Y_s)\,ds\right\|+\|B_t\|.\]

Negation leaves the norm of the drift integral unchanged.



Apply the two bounds to obtain \(Gt+N\sqrt t\), and commute the factors \(G\) and \(t\).



For each \(t\), the drift integral runs only over \([0,t]\), which lies in the assumed control interval \([0,h]\). Its norm is at most \(\int_0^tGds=tG\). Adding the noise estimate gives the stated bound; no independence between drift and noise is used, because this is a norm inequality before taking any probabilistic expectation.

\subsection{\texorpdfstring{Integrate the Lipschitz increment bound}{Integrate the Lipschitz increment bound}}


Assume \(h\ge0\), \(g\) is \(L\)-Lipschitz with \(L\ge0\), and for \(0\le t\le h\) the same integral equation and bounds \(\|g(Y_t)\|\le G\), \(\|B_t\|\le N\sqrt t\) hold. Then

\[\|D_h(g,Y)\|\le L\left(\frac{h^2G}{2}+\frac23h\sqrt h\,N\right).\]

For \(0<t\le h\), Lipschitz continuity gives

\[\|g(Y_t)-g(Y_0)\|\le L\|Y_t-Y_0\|.\]

Multiply the increment estimate by \(L\ge0\) to obtain the pointwise majorant

\[\|g(Y_t)-g(Y_0)\|\le L(tG+N\sqrt t).\]

The scalar majorant \(t\mapsto L(tG+N\sqrt t)\) is continuous, hence integrable on \([0,h]\).



Expand \(D_h\). The norm estimate for a Bochner integral, together with the integrable majorant, gives

\[\|D_h\|\le\int_0^h L(tG+N\sqrt t)\,dt.\]

Pull out constants and split the two integrable summands. Evaluate

\[\int_0^h t\,dt=\frac{h^2}{2},\qquad \int_0^h\sqrt t\,dt=\frac23h\sqrt h.\]

This leaves the two drift and noise contributions in the claimed expression.



Rearrange the scalar factors to place the drift term first and the noise term second.



The two terms have distinct origins and scales: \(Lh^2G/2\) comes from the deterministic accumulation of drift, and \((2/3)LNh^{3/2}\) comes from the square-root growth of the additive input. Retaining both terms until the stationary scales are substituted avoids losing the constant by an early coarse bound on the whole segment.

\subsection{\texorpdfstring{The exact rational constant}{The exact rational constant}}


The numerical factor in the corrected Dalalyan–Karagulyan estimate satisfies

\[\frac76\sqrt2\le\frac{33}{20}=1.65.\]

Use \((\sqrt2)^2=2\) and \(\sqrt2\ge0\). Comparing squares proves the bound, since \(49/18\le1089/400\).



Both sides are nonnegative, so their squares may be compared without changing the order. After squaring, the difference is the rational positive number \(1089/400-49/18\). This establishes the decimal constant as an upper bound, rather than relying on a rounded floating-point evaluation of \(\sqrt2\).

\subsection{\texorpdfstring{The local error with stationary moment scales}{The local error with stationary moment scales}}


Let \(h,d\in\mathbb R\) with \(h,d\ge0\), let \(g\) be \(L\)-Lipschitz with \(L\ge0\), and suppose \(Lh\le2\). Assume that, for every \(0\le t\le h\),

\[Y_t-Y_0=-\int_0^t g(Y_s)\,ds+B_t,\quad \|g(Y_t)\|\le\sqrt{Ld},\quad \|B_t\|\le\sqrt{2d}\sqrt t.\]

Then

\[\boxed{\|D_h(g,Y)\|\le\frac{33}{20}Lh\sqrt{hd}.}\]

In \(L^2\), the drift bound corresponds to a score second moment at most \(Ld\); here that bound remains a hypothesis.



Compare the nonnegative quantities \(h\sqrt{Ld}\) and \(\sqrt2\sqrt{hd}\) by squaring. It suffices to show \(h^2Ld\le2hd\).



Replace each squared square root by its nonnegative argument. Multiply \(Lh\le2\) by \(hd\ge0\), obtaining \(h^2Ld\le2hd\), and thus

\[h\sqrt{Ld}\le\sqrt2\sqrt{hd}.\]

Split products under the square roots and rearrange:

\[h\sqrt h\sqrt{2d}=h\sqrt2\sqrt{hd}.\]

The nonnegativity of \(h\) and \(2\) justifies the square-root product identities.



Apply the two-term local-defect bound with \(G=\sqrt{Ld}\) and \(N=\sqrt{2d}\):

\[\|D_h\|\le L\left(\frac{h^2\sqrt{Ld}}2+\frac23h\sqrt h\sqrt{2d}\right).\]

Bound the first term using \(h\sqrt{Ld}\le\sqrt2\sqrt{hd}\), and rewrite the second with the preceding identity. Since \(1/2+2/3=7/6\),

\[\|D_h\|\le L\frac76\sqrt2\,h\sqrt{hd}.\]

Multiply \((7/6)\sqrt2\le33/20\) by the nonnegative factor \(Lh\sqrt{hd}\). This gives the claimed constant without numerical rounding.



The condition \(Lh\le2\) converts the drift contribution, initially of order \(h^2\sqrt{Ld}\), into the same \(Lh\sqrt{hd}\) scale as the Brownian contribution. The sum of coefficients is \(\sqrt2(1/2+2/3)=(7/6)\sqrt2\). Replacing it by \(33/20\) is the only numerical enlargement in this local-defect estimate.

\section{\texorpdfstring{The noise-\allowbreak{}driven iterates}{The noise-driven iterates}}
\label{proof-module-UnadjustedLangevin-Dynamics-lean}


Write \(T_h(x)=x-h\nabla U(x)\), so one Euler update is \(T_h(x)+\sqrt{2h}\,z\). The noise vectors remain explicit. Common-noise algebra works for any supplied noise sequence; identifying the algorithm with Gaussian ULA also requires independent standard Gaussian noise hypotheses.



The Euler recursion applies the displayed update to the initial state and supplied innovation sequence. Its zero and successor identities follow directly from this recursive definition. Measurability follows by induction: the current iterate is measurable, the Lipschitz gradient is continuous and hence measurable, and the next innovation is measurable by assumption. Addition and scalar multiplication preserve these properties. This argument requires no independence or Gaussian-law assumption.



At a fixed index \(n\), the state is a function of the initial vector and the finite list \(\xi_0,\ldots,\xi_{n-1}\). The recursion first forms the drift from the current state and then uses the innovation with the same index. This indexing is what makes the next innovation independent of the current state under an independent-noise construction. Defining and measuring the recursion itself requires no distributional assumptions on the innovations.

\subsection{\texorpdfstring{Definition of the chain}{Definition of the chain}}


Given \(U:\mathbb R^d\to\mathbb R\), \(h\in\mathbb R\), an initial vector \(x_0\), and noise \((\xi_k)\), define

\[X_0=x_0,\qquad X_{k+1}=X_k-h\nabla U(X_k)+\sqrt{2h}\,\xi_k.\]

Start the recursion at the supplied initial vector: \(X_0=x_0\).



Apply the Euler update to \(X_k\) using the next noise vector \(\xi_k\). This determines \(X_{k+1}\).



The recursion determines a unique finite-time state from the initial vector and a finite prefix of innovations. The potential enters only through its gradient evaluated at already constructed states. This separates the algebraic definition of the method from later assumptions that make the gradient measurable, the noise Gaussian, or the transition contractive.

\subsection{\texorpdfstring{The initial iterate}{The initial iterate}}


For every potential, step, initial vector, and noise sequence, the recursively defined chain satisfies \(X_0=x_0\).



The equality is exactly the initial clause of the recursive definition.



At index zero the noise prefix is empty. Thus the initial law of the chain is exactly the pushforward law of the supplied initialization, without any transition or Gaussian convolution. This identity is the base case for the later induction identifying the noise-driven law with the iterated Markov transition.

\subsection{\texorpdfstring{The update identity}{The update identity}}


For every \(k\ge0\), the chain satisfies \(X_{k+1}=T_h(U,h,X_k,\xi_k)\).



The successor clause of the recursion gives the asserted update formula.



Writing the successor through the update map makes the dependence on the current state and fresh innovation explicit. The state on the right has index \(k\), and the noise has that same index; the result has index \(k+1\). This alignment is used in both the independence argument and the common-noise coupling comparison.

\subsection{\texorpdfstring{Measurability of every iterate}{Measurability of every iterate}}


Let \(\Omega\) be a measurable space. If \(\nabla U\) is \(L\)-Lipschitz, the initial vector \(X_0:\Omega\to\mathbb R^d\) is measurable, and every \(\xi_k:\Omega\to\mathbb R^d\) is measurable, then \(X_n\) is measurable for every \(n\).



A Lipschitz gradient is continuous and therefore measurable.



Induct on \(n\). The initial iterate is measurable by hypothesis. For the successor step, assume \(X_n\) is measurable.



The composition \(\nabla U\circ X_n\) is measurable. Subtracting its scalar multiple from \(X_n\), then adding \(\sqrt{2h}\,\xi_n\), preserves measurability. Hence \(X_{n+1}\) is measurable.



At the induction step, continuity of the gradient makes \(\nabla U(X_n)\) measurable, while measurable vector operations preserve the two scalar multiples and their sum. The proof concerns every fixed finite \(n\). Its conclusion permits the pushforward law of each iterate to be used in the distributional convergence statements.

\section{\texorpdfstring{The actual gradient on square-\allowbreak{}integrable random variables}{The actual gradient on square-integrable random variables}}
\label{proof-module-UnadjustedLangevin-LpDrift-lean}


Let \(P\) be a finite measure on \(\Omega\). The first five results need only a normed additive commutative group \(H\); all Lipschitz constants \(L\) are nonnegative. When \(H\) is a real Banach space, \(L^2(P;H)\) provides the Banach space needed for the local-defect estimate. For a Lipschitz vector field \(g:H\to H\), construct the operator \(\mathcal G[X]=[g\circ X]\), including the case \(g(0)\ne0\). Finiteness of \(P\) ensures the constant offset is square integrable.



The gradient results use a real Hilbert space \(H\). The two score-moment lemmas need only a real inner-product space, without completeness. Whenever \(g=\nabla U\), this is the potential’s actual gradient, and the resulting \(L^2\) recursion is proved to represent the original pointwise ULA chain.



The pointwise-composition operator lifts a Lipschitz pointwise map to \(L^2\), using its linear growth and the finite probability measure to prove integrability. Its representative theorem identifies the lifted map almost everywhere with ordinary composition. All subsequent Hilbert-space estimates concern these quotient classes and are invariant under representative choices.



Write \(\mathcal H=L^2(P;H)\) and \(\|X\|_2^2=\int\|X\|^2dP\). A Lipschitz drift has linear growth, \(\|g(x)\|\le L\|x\|+\|g(0)\|\), so it maps square-integrable random vectors into square-integrable random vectors whenever \(P\) is finite. The lifted drift is therefore an actual map on \(\mathcal H\). Almost-everywhere identities between representatives express equality of random variables; integrating their squared norms gives the corresponding Hilbert-space statements.

\subsection{\texorpdfstring{Center a Lipschitz vector field}{Center a Lipschitz vector field}}


In a normed additive group \(H\), if \(g:H\to H\) is \(L\)-Lipschitz, then so is

\[\widetilde g(x):=g(x)-g(0).\]

For arbitrary \(x,y\), the common offset cancels:

\[\|\widetilde g(x)-\widetilde g(y)\|=\|g(x)-g(y)\|\le L\|x-y\|.\]

This is precisely the required Lipschitz inequality.



Centering makes the vector field vanish at zero, so \(\|\widetilde g(x)\|=\|g(x)-g(0)\|\le L\|x\|\). This growth estimate is stronger than a bare continuity statement: squaring and integrating shows that composition with \(\widetilde g\) preserves \(L^2\). The constant offset is restored only after this integrable composition has been constructed.

\subsection{\texorpdfstring{Construct the lifted drift}{Construct the lifted drift}}


Let \(P\) be finite and \(g\) be \(L\)-Lipschitz. For \(X\in L^2(P;H)\), define

\[\mathcal G[X]:=[\widetilde g\circ X]+[g(0)].\]

The centered function satisfies \(\widetilde g(0)=0\) and \(\|\widetilde g(x)\|\le L\|x\|\), so composition produces an \(L^2\) class. Add the constant class \([g(0)]\), which belongs to \(L^2\) because \(P\) is finite.



The constant class has norm \(\|g(0)\|\sqrt{P(\Omega)}\), which is finite under the stated finite-measure hypothesis. The centered composition has norm at most \(L\|X\|_2\). Their sum consequently satisfies \(\|\mathcal G[X]\|_2\le L\|X\|_2+\|g(0)\|\sqrt{P(\Omega)}\), explicitly explaining why the lift is defined on every square-integrable input.

\subsection{\texorpdfstring{The lift agrees with the original vector field}{The lift agrees with the original vector field}}


For finite \(P\), Lipschitz \(g\), and \(X\in L^2(P;H)\),

\[\mathcal G[X](\omega)=g(X(\omega))\quad\text{for }P\text{-almost every }\omega.\]

Expand \(\mathcal G\) into the centered composition class plus the constant class.



Intersect the full-measure sets on which the representative of the sum, the composition, and the constant class agree with their pointwise formulas.



On this set the value is \((g(X(\omega))-g(0))+g(0)\). Cancel the offset to recover \(g(X(\omega))\).



Each constituent operation on equivalence classes has an almost-everywhere pointwise formula. A finite intersection of the corresponding full-measure sets still has full measure. On that set, the two copies of \(g(0)\) cancel exactly, identifying the constructed class with the intended drift rather than a centered approximation to it.

\subsection{\texorpdfstring{The actual pointwise drift is square integrable}{The actual pointwise drift is square integrable}}


For finite \(P\), Lipschitz \(g\), and \(X\in L^2(P;H)\),

\[g\circ X\in L^2(P;H).\]

The constructed class \(\mathcal G[X]\) has an \(L^2\) representative. Transfer its square integrability to \(g\circ X\) using their almost-everywhere equality.



The identity with \(\mathcal G[X]\) provides both measurability and finite square norm for the actual composition. This avoids assuming in advance that a nonlinear drift maps \(L^2\) to itself. The needed fact instead follows from its global Lipschitz growth and the finiteness of the underlying measure.

\subsection{\texorpdfstring{The drift keeps its Lipschitz constant in L\ensuremath{{}^{2}}}{The drift keeps its Lipschitz constant in L²}}


For finite \(P\) and \(L\)-Lipschitz \(g\), the lifted drift is \(L\)-Lipschitz:

\[\|\mathcal G[X]-\mathcal G[Y]\|_{L^2}\le L\|X-Y\|_{L^2}.\]

Fix \(X,Y\). The two constant classes cancel in their drift difference. For the remaining centered compositions, the pointwise Lipschitz estimate passes to the \(L^2\) norm, giving the same constant \(L\).



On a common full-measure set, \(\|g(X)-g(Y)\|^2\le L^2\|X-Y\|^2\). Integrating and taking square roots gives the lifted bound. The constant offset disappears from the difference, so there is no extra term proportional to \(\|g(0)\|\). This exact preservation of the Lipschitz constant is needed for the integral-equation fixed point.

\subsection{\texorpdfstring{Lift the small-\allowbreak{}step gradient contraction}{Lift the small-step gradient contraction}}


Let \(H\) be a real Hilbert space, \(P\) finite, and \(U:H\to\mathbb R\) differentiable and \(m\)-strongly convex on \(H\), with \(L\)-Lipschitz gradient. Assume \(m\le L\), \(m+L>0\), and \(0\le h\le2/(m+L)\). Write \(\mathcal G[X]=[\nabla U\circ X]\). Then, for all \(X,Y\in L^2(P;H)\),

\[\|(X-h\mathcal G[X])-(Y-h\mathcal G[Y])\|_{L^2}\le(1-mh)\|X-Y\|_{L^2}.\]

It suffices to prove the corresponding norm inequality almost everywhere between representatives: integrating the squared pointwise inequality gives the \(L^2\) bound.



Work on a common full-measure set where the outer difference and both Euler-step differences have their ordinary pointwise subtraction formulas.



Also require the scalar multiples, both lifted gradients, and \(X-Y\) to have their pointwise formulas. The two drift terms are then \(\nabla U(X(\omega))\) and \(\nabla U(Y(\omega))\).



Substitute these representative identities and apply the deterministic sharp contraction to \(x=X(\omega)\), \(y=Y(\omega)\):

\[\|(x-h\nabla U(x))-(y-h\nabla U(y))\|\le(1-mh)\|x-y\|.\]

The deterministic factor \(1-mh\) is nonnegative in the stated small-step regime. Its squared pointwise bound can therefore be integrated and square-rooted with that same factor. The resulting estimate acts on the actual nonlinear drift lift, so it can be inserted directly into the random-vector Euler comparison.

\subsection{\texorpdfstring{Lift the second step-\allowbreak{}size regime}{Lift the second step-size regime}}


Let \(H\) be a real Hilbert space, \(P\) finite, and \(U\) differentiable, \(m\)-strongly convex, with \(L\)-Lipschitz gradient. Assume \(m\le L\), \(m+L>0\), \(h\ge0\), and \(h\ge2/(m+L)\). Then

\[\|(X-h\mathcal G[X])-(Y-h\mathcal G[Y])\|_{L^2}\le(Lh-1)\|X-Y\|_{L^2}.\]

The inequality holds in this whole regime; it is a strict contraction when its factor is less than one.



Reduce the \(L^2\) inequality to its almost-everywhere pointwise counterpart with factor \(Lh-1\).



Choose a full-measure set on which the outer subtraction and each Euler-step subtraction agree with the pointwise vector formulas.



On the same set identify scalar multiplication, both lifted gradient representatives, and the representative of \(X-Y\).



After substitution, invoke the deterministic second-regime estimate at \(X(\omega),Y(\omega)\):

\[\|(x-h\nabla U(x))-(y-h\nabla U(y))\|\le(Lh-1)\|x-y\|.\]

In the second regime the deterministic coefficient is \(Lh-1\ge0\). Squaring the pointwise estimate and integrating retains this coefficient exactly. If \(Lh<2\), it lies below one; if not, the inequality remains true but no geometric contraction follows. The lift itself does not add a stability restriction beyond the deterministic one.

\subsection{\texorpdfstring{The local defect for the actual L\ensuremath{{}^{2}} gradient}{The local defect for the actual L² gradient}}


Let \(P\) be finite, \(H\) a real Hilbert space, and \(\nabla U\) be \(L\)-Lipschitz. Let \(Y,B:\mathbb R\to L^2(P;H)\), let \(h,d\in\mathbb R\), and assume \(h,d\ge0\), \(Lh\le2\). For every \(0\le t\le h\), suppose

\[\|\mathcal G[Y_t]\|_{L^2}\le\sqrt{Ld},\quad\|B_t\|_{L^2}\le\sqrt{2d}\sqrt t,\quad Y_t-Y_0=-\int_0^t\mathcal G[Y_s]\,ds+B_t.\]

Then

\[\|D_h(\mathcal G,Y)\|_{L^2}\le\frac{33}{20}Lh\sqrt{hd}.\]

Apply the Banach-space local-defect theorem to \(\mathcal G\). Its required Lipschitz property is the one just proved for the actual lifted gradient; the reference equation and moment-scale bounds are the displayed hypotheses.



The Banach-space defect estimate is applied with the ambient norm equal to \(\|\cdot\|_2\). Its drift moment and noise moment inputs are already norms in that space. Thus the resulting defect is a genuine root mean square local error, and no extra exchange of stochastic and time integrals is needed in this abstract calculation.

\subsection{\texorpdfstring{Identify the drift norm with the score moment}{Identify the drift norm with the score moment}}


For finite \(P\), a real inner-product space \(H\), Lipschitz \(g:H\to H\), and \(X\in L^2(P;H)\),

\[\|\mathcal G[X]\|_{L^2}^2=\int_\Omega\|g(X(\omega))\|^2\,dP(\omega).\]

Write the squared \(L^2\) norm as a self-inner-product and expand the \(L^2\) inner product into the integral of the pointwise inner product.



Almost everywhere, \(\mathcal G[X](\omega)=g(X(\omega))\). Replace the representative under the integral and use \(\langle v,v\rangle=\|v\|^2\).



The equality converts the abstract Hilbert-space norm of the drift into the usual probabilistic energy. When the law of \(X\) is known, the right side may be changed into an integral against that law. In particular, a stationary Gibbs law permits the Gibbs score estimate to control the drift norm along the whole reference.

\subsection{\texorpdfstring{A second-\allowbreak{}moment bound gives the drift bound}{A second-moment bound gives the drift bound}}


Under the preceding finite-measure and inner-product assumptions, if

\[\int_\Omega\|g(X(\omega))\|^2\,dP(\omega)\le C,\]

then

\[\|\mathcal G[X]\|_{L^2}\le\sqrt C.\]

For \(g=\nabla U\) and \(C=Ld\), the moment inequality is an explicit input, not a consequence of this lemma.



The norm-square identity turns the assumed moment bound into \(\|\mathcal G[X]\|^2\le C\). Taking the nonnegative square root yields the conclusion.



Because the left side is a square, the assumed upper bound necessarily gives \(C\ge0\) whenever the data are consistent. Monotonicity of the nonnegative square root then yields the desired norm bound. The lemma transfers an established energy estimate; it does not itself supply any probabilistic or functional inequality proving that estimate.

\subsection{\texorpdfstring{Define one ULA update in L\ensuremath{{}^{2}}}{Define one ULA update in L²}}


For finite \(P\), a real Hilbert space \(H\), a Lipschitz gradient \(\nabla U\), step \(h\in\mathbb R\), and \(X,Z\in L^2(P;H)\), define

\[\mathsf T_h(X,Z):=X-h\mathcal G[X]+\sqrt{2h}\,Z.\]

Form the drift step and add the scaled noise inside \(L^2\). All terms are already equivalence classes of square-integrable variables, so the vector-space operations produce another such class.



The lift makes the nonlinear gradient term an element of the same Hilbert space as the state and noise. Addition and scalar multiplication therefore define the whole step without leaving \(L^2\). For \(h>0\), the formula coincides with the ordinary ULA update once representatives are evaluated almost everywhere.

\subsection{\texorpdfstring{The L\ensuremath{{}^{2}} update is the pointwise ULA formula}{The L² update is the pointwise ULA formula}}


For finite \(P\), a real Hilbert space \(H\), Lipschitz \(\nabla U\), and arbitrary \(X,Z\in L^2(P;H)\),

\[\mathsf T_h(X,Z)(\omega)=X(\omega)-h\nabla U(X(\omega))+\sqrt{2h}\,Z(\omega)\quad\text{a.e.}\]

This identity needs no Gaussian or independence assumption on \(Z\).



Expand the definition of \(\mathsf T_h\) into its subtraction, drift, and scaled-noise terms.



Intersect the full-measure sets on which addition, subtraction, both scalar multiples, and the lifted gradient agree with their pointwise formulas.



Substitute these identities. The representative is exactly \(X(\omega)-h\nabla U(X(\omega))+\sqrt{2h}\,Z(\omega)\).



No conditional expectation appears in this representative identity. The lifted drift is simply the original gradient evaluated at the random state, while the noise class is multiplied by the deterministic scalar. Thus every subsequent pathwise cancellation can be carried out on a common full-measure set using the familiar Euclidean update formula.

\subsection{\texorpdfstring{Define the genuine L\ensuremath{{}^{2}} ULA chain}{Define the genuine L² ULA chain}}


For a finite measure \(P\), a Lipschitz gradient on \(E_d\), an initial class \(X_0\in L^2(P;E_d)\), and noise classes \(Z_n\in L^2(P;E_d)\), define recursively, with the supplied initial class at index zero,

\[X_{n+1}:=\mathsf T_h(X_n,Z_n).\]

The initial class is the supplied \(X_0\). At the successor index, apply the \(L^2\) ULA update to the preceding iterate and the next noise class \(Z_n\).



The recursion stays inside \(L^2\) at every finite index because its one-step map has already been defined on that whole space. In particular, no separate induction on second moments of the chain is needed for existence of these classes. Distributional properties still require the innovations to satisfy the separate Gaussian and independence hypotheses.

\subsection{\texorpdfstring{The initial iterate}{The initial iterate}}


The recursively defined chain satisfies \(X_0=\mathrm{initial}\).



This is the zero branch of the recursion, so both sides reduce to the same class.



The zero clause identifies the initial class exactly, including its law and its distance from a reference initial class. This matters when the recurrence is unwound: the initial term is \(q^n\|X_0-Y_0\|_2\), with no additional initialization error. The statement is independent of the later noise sequence.

\subsection{\texorpdfstring{The successor iterate}{The successor iterate}}


For every \(n\ge0\), the same finite-measure, Lipschitz-gradient, initial-class and noise-class data give

\[X_{n+1}=\mathsf T_h(X_n,Z_n).\]

This is the successor branch of the defining recursion; unfold that branch to obtain the equality.



Only the innovation at index \(n\) enters the update from \(n\) to \(n+1\). Exposing this clause is the algebraic step used before subtracting the reference update. The same noise is then present in both expressions and cancels, leaving the lifted deterministic Euler difference and the defect.

\subsection{\texorpdfstring{Identify the chain with the original algorithm}{Identify the chain with the original algorithm}}


For finite \(P\), Lipschitz gradient \(\nabla U:E_d\to E_d\), square-integrable initial and noise classes, and every \(n\ge0\),

\[X_n(\omega)=X\bigl(U,h,X_0(\omega),(Z_k(\omega))_{k\ge0},n\bigr)\quad\text{for }P\text{-a.e. }\omega.\]

Induct on \(n\). At \(n=0\), both recursions return the same initial representative, so equality holds for every outcome. For the successor case, assume the representatives agree at index \(n\).



On a common full-measure set, combine the one-step representative identity with the induction hypothesis. Replace \(X_n(\omega)\) by the pointwise chain’s preceding iterate and apply the two successor recursion equations. Both next iterates are the same ULA update.



For a fixed finite index, intersecting the previous equality set with the one-step representative equality set preserves full measure. This proves agreement of the recursively defined random vector with the class representative at that index. If desired, the countable intersection over all indices gives one full-measure set on which all discrete-time identities hold simultaneously.

\section{\texorpdfstring{Existence under the additive Langevin integral under equation}{Existence under the additive Langevin integral under equation}}
\label{proof-module-UnadjustedLangevin-LangevinIntegral-lean}


In a complete real normed space, a globally Lipschitz drift and a continuous additive path \(B\), \(B(0)=0\), determine a continuous solution of \(Y_t-x=-\int_0^tg(Y_s)ds+B_t\) on every finite nonnegative horizon. The construction is deterministic and later applies with the Banach space \(L^2(P)\).



For a continuous additive input \(B\) with \(B(0)=0\), solve \(Y(t)=x-\int_0^tg(Y(s))ds+B(t)\). Subtracting \(B\) transforms this into an ordinary integral equation with time-dependent Lipschitz drift; differentiability of \(B\) is unnecessary. A short-interval fixed point exists once the interval is small relative to the Lipschitz constant. Restarting at successive endpoints and shifting the input yields every finite horizon, with continuity and the integral identity preserved at each join.

\subsection{\texorpdfstring{Solving the integral equation on a short interval}{Solving the integral equation on a short interval}}


Remove the additive input by writing \(Y(t)=Z(t)+B(t)\). The new unknown satisfies \(Z(t)=x+\int_0^t f(s,Z(s))ds\), with \(f(s,z)=-g(z+B(s))\). This time-dependent drift is \(L\)-Lipschitz in \(z\) and continuous in time even though \(B\) need not be differentiable.



Bound \(\|B(t)\|\le C'\) on the compact interval and set \(G=\|g(x)\|+LC'\), radius \(a=2hG\), and drift bound \(M=G+La\). The condition \(Lh\le1/2\) proves \(Mh\le a\), giving the ball-invariance condition in the local contraction theorem for the integral equation. The proof checks these bounds and invokes the contraction fixed-point theorem.



The chosen fixed point gives a continuous \(Z\) with its integral identity. Add back \(B\), use \(B(0)=0\), and rewrite the negative drift integral to obtain \(Y(t)-x=-\int_0^tg(Y(s))ds+B(t)\). The solution is represented as a real-time function, with the required identity asserted on \([0,h]\).



The shifted unknown has an ordinary time integral even if the additive input has no derivative. The short-time radius and drift bounds ensure that the Picard map preserves its closed ball, while the Lipschitz estimate controls contraction there.

\subsection{\texorpdfstring{Splicing two integral-\allowbreak{}equation segments}{Splicing two integral-equation segments}}


Use \(Y(t)\) up to time \(a\), and \(Z(t-a)\) afterward. The equality \(Z(0)=Y(a)\) makes the pasted path continuous at the join. Split \(\int_0^t=\int_0^a+\int_a^t\), translate the second integral, and add the two supplied equations; the terms \(B(a)\) cancel, giving the same equation on \([0,a+b]\).



At the join, the two path values coincide by the restart condition. For later times split the drift integral at the joining time and translate the second part; the two input increments then add to the original input difference.

\subsection{\texorpdfstring{Existence on every finite horizon}{Existence on every finite horizon}}


Choose a positive segment length small enough that \(L\delta\le1/2\). Construct the first Picard segment, then repeatedly solve with initial state at the previous endpoint and shifted input \(B(a+s)-B(a)\). The splicing lemma combines finitely many segments covering the prescribed horizon. The local contraction restriction determines the construction mesh and imposes no restriction on the final horizon.



A finite horizon can be covered by finitely many equal short intervals. At each restart the shifted input is continuous and starts at zero, so the same local theorem applies repeatedly with a new initial point.

\subsection{\texorpdfstring{A short-\allowbreak{}interval solution in square mean}{A short-interval solution in square mean}}


Apply the Banach-space Picard theorem to \(H=L^2(P;K)\), \(g(X)=\mathcal G(\nabla U)(X)\). The lifted drift has the same Lipschitz constant as \(\nabla U\). Completeness of \(L^2\) and continuity of the supplied \(L^2\)-noise path are the exact fixed-point hypotheses.



The lifted gradient's global Lipschitz constant is unchanged, and the space of square-integrable vectors is complete. These facts verify the short-interval fixed-point assumptions in the random-vector Hilbert space.

\subsection{\texorpdfstring{A finite-\allowbreak{}horizon solution in square mean}{A finite-horizon solution in square mean}}


Use the finite-horizon Banach-space existence theorem with the lifted gradient drift. The conclusion is continuity in the \(L^2\) norm and an equality of \(L^2\) classes at every time in the horizon. The Bochner time integral therefore lives in the same Hilbert space; no pathwise differentiability of the random noise is asserted.



The finite-horizon construction produces one continuous \(L^2\)-valued curve. Its Bochner drift integral is well defined by continuity, and the integral identity is asserted at each time in the prescribed horizon.

\section{\texorpdfstring{Construct the kinetic integral equation by Picard iteration}{Construct the kinetic integral equation by Picard iteration}}
\label{proof-module-KineticLangevin-Dynamics-lean}


For existence, the phase product may use its convenient maximum norm. This norm is equivalent to the Euclidean phase norm but is used here only to establish a global Lipschitz vector field. The input noise is any continuous additive velocity path starting at zero. The theorem therefore constructs paths without postulating a reference kinetic process; a full Brownian solution theorem additionally identifies the stochastic noise and its laws.



The construction is pathwise and works for arbitrary continuous additive forcing. This provides the deterministic solution map on which the later Brownian law and coupling constructions are built.

\subsection{\texorpdfstring{Write the drift in the shared integral-\allowbreak{}equation sign convention}{Write the drift in the shared integral-equation sign convention}}


The kinetic drift is \((-\gamma v-ug(x),v)\). The shared Picard theorem solves an equation with a minus sign in front of its drift integral, so the negative drift stores \((\gamma v+ug(x),-v)\). This sign choice ensures that \(Y_t-Y_0=-\int_0^tD(Y_s)ds+(B_t,0)\) has velocity acceleration \(-\gamma v-ug(x)\) and position derivative \(v\). The second noise coordinate is exactly zero.



Writing the equation componentwise removes any ambiguity: \(v_t=v_0-\int_0^t(\gamma v_s+ug(x_s))ds+B_t\) and \(x_t=x_0+\int_0^tv_sds\). In particular the position has an ordinary time derivative even when the forcing is only continuous.

\subsection{\texorpdfstring{Prove a global Lipschitz constant for phase space}{Prove a global Lipschitz constant for phase space}}


If \(g\) is \(L\)-Lipschitz and \(\gamma,u,L\ge0\), write \(R=\|(v,x)-(w,y)\|_{\max}\). Each coordinate difference is at most \(R\), and \(\|g(x)-g(y)\|\le LR\). The velocity component of the negative drift is bounded by \((\gamma+uL)R\) using the triangle inequality. The position component is bounded by \(R\). Hence the maximum is bounded by \((\gamma+uL+1)R\). The argument expands the product norm maximum, proves the two coordinate bounds, rewrites the drift difference into scalar multiples of differences with vector-space algebra, and combines the nonnegative coefficients. The constant is deliberately sufficient for existence rather than optimized for contraction.



The maximum norm is useful because each coordinate projection has operator norm one. A sharper phase Lipschitz constant would also prove existence, but its value is immaterial for the later contraction rate: that rate is derived separately from curvature in the transformed Hilbert norm.

\subsection{\texorpdfstring{Invoke the proved continuous-\allowbreak{}path Picard construction}{Invoke the proved continuous-path Picard construction}}


Let the underlying normed space be complete, \(B\) continuous with \(B_0=0\), and \(T\ge0\). The shared theorem the continuous-input Picard existence theorem constructs a continuous phase path \(Y\), with \(Y_0=\text{initial}\), satisfying the kinetic integral identity for every \(0\le t\le T\). The call supplies the concrete negative drift and its proved Lipschitz constant. It packages the noise as \((B_t,0)\), whose continuity follows from product continuity and whose value at zero follows from the zero initial forcing condition. This is finite-horizon pathwise existence; Gaussian stationarity, a canonical Brownian law, and full probabilistic chain convergence are not conclusions of this theorem.



Picard iteration solves the integral equation on a short interval by contraction of the integral map. Global Lipschitz continuity allows successive intervals to be joined without finite-time blowup. On a prescribed finite horizon the resulting path is unique and depends continuously on the initial phase and the forcing.

\section{\texorpdfstring{The exact physical-\allowbreak{}time integral equation}{The exact physical-time integral equation}}
\label{proof-module-KIPLMC-KineticClock-lean}


Let \(Y_s\) solve the unit-temperature kinetic equation with sum force \(g\) and friction \(c\gamma\). The physical process is \(S_t=D_{c^{-1}}Y_{t/c}\), where \(D_a(v,x)=(av,x)\). This section proves the complete transformed integral equation in any complete normed vector space, including the L\ensuremath{{}^{2}} space of random variables used by the Brownian construction.



The time and velocity change is performed inside the integral equation using bounded linearity and substitution. This gives a physical-flow identity compatible with the already proved Brownian scaling.

\subsection{\texorpdfstring{Make velocity scaling a bounded linear map}{Make velocity scaling a bounded linear map}}


The velocity projection multiplied by \(c\), paired with the unchanged position projection, defines a continuous linear map. Its action is exactly the velocity rescaling. Linearity allows this map to pass through Bochner integrals; continuity ensures the transformed path remains continuous. Expanding the phase drift proves \(c^{-1}D_{c^{-1}}R_{g,c\gamma}(z)=R_{c^{-2}g,\gamma}(D_{c^{-1}}z)\). The two powers of \(c^{-1}\) in the force come from time and velocity separately.



Velocity scaling is a bounded linear map on the full phase space. It therefore commutes with Bochner integration and preserves continuity of the constructed paths. Its position component is the identity, so physical parameter observables are unaffected.

\subsection{\texorpdfstring{Change variables inside the integral equation}{Change variables inside the integral equation}}


For \(0\le t\le T\), positivity of \(c\) gives \(t/c\in[0,T/c]\). Apply the continuous linear map to the given solution identity at \(t/c\), pass it through the integral, and substitute \(r=s/c\). The drift identity identifies the transformed force and friction, and the additive source becomes \(c^{-1}B_{t/c}\). One proves equality of the complete vector-valued integrals, not merely equality of formal differential symbols. This theorem requires no derivative of the Brownian input.



Apply the map to the original solution equation and perform the positive time substitution inside its drift integral. The force, friction, and thermal input transform with their stated factors. This proves an identity between integral solutions, not merely between formal differential equations.

\section{\texorpdfstring{Construct correlated Gaussian blocks and an independent noise stream}{Construct correlated Gaussian blocks and an independent noise stream}}
\label{proof-module-KineticLangevin-Gaussian-lean}


A covariance matrix alone is insufficient to define a random algorithm. This section constructs centered finite-dimensional Gaussian probability measures with the Gram covariance and an infinite product giving one independent block per time and spatial coordinate. Correlation is retained within each block. The exact covariance theorem requires a nonnegative step and continuous kernels; the general Gaussian constructor remains a probability measure even outside that parameter regime.



The noise construction distinguishes independence of different blocks from correlation of responses inside a block. Both features are needed to identify the sampler law and its common-source coupling.

\subsection{\texorpdfstring{Construct the centered Gaussian probability measure and its vector moments}{Construct the centered Gaussian probability measure and its vector moments}}


For a finite kernel index set, the centered Gaussian response law is the multivariate Gaussian with mean zero and the defined Gram matrix. The named instances establish Gaussianity and total probability mass using the distribution constructor. The mean theorem computes its vector expectation as zero, and finite Gaussian response moments invokes the finite Gaussian second-moment theorem. These are derived properties of a specified probability measure, rather than moment fields of an abstract noise source. They hold even if the covariance argument is degenerate; identifying its entries with the requested matrix uses the positive-semidefinite theorem below.



A covariance Gram matrix determines a centered Gaussian law by applying a square-root linear map to a standard Gaussian. The resulting law is a probability measure even when the covariance is singular. Finite-dimensional Gaussian moments then provide the integrability needed for every coordinate calculation.

\subsection{\texorpdfstring{Evaluate coordinate means, moments and exact covariances}{Evaluate coordinate means, moments and exact covariances}}


Each coordinate projection is a continuous linear functional. Commuting it with the vector integral proves zero coordinate mean, and composing the finite-second-moment identity with that functional proves coordinate L\ensuremath{{}^{2}} integrability. For continuous kernels and \(h\ge0\), Gram positivity permits the multivariate-Gaussian covariance formula:

\[\operatorname{Cov}(X_i,X_j)=\int_0^hk_i(t)k_j(t)dt,\quad
\operatorname{Var}(X_i)=\int_0^hk_i(t)^2dt.\]

The two distribution theorems receive the already proved positive-semidefinite matrix witness. The variance proof only rewrites a product as a square. This identifies the exact correlated coordinate law required by the integrators, while its representation as Brownian stochastic integrals is a separate identity.



For centered variables, covariance is simply the expectation of their product. Thus each entry of the constructed matrix has an immediate probabilistic meaning. Diagonal entries give variances; off-diagonal entries describe the dependence between distinct response coordinates generated from the same noise.

\subsection{\texorpdfstring{Instantiate the two-\allowbreak{}block and four-\allowbreak{}block kinetic laws}{Instantiate the two-block and four-block kinetic laws}}


the two-response Gaussian law uses the kernel vector \((\psi_0,\psi_1)\); the four-response Gaussian law uses \((\psi_0,\psi_1,\varphi_2,\varphi_3)\). They are therefore concrete measures on two- and four-dimensional Euclidean space. Each is the linear pushforward of a normalized Gaussian measure and therefore has total mass one. The first law correlates a coordinate's velocity and position increments; the second also couples the Hessian-correction integrals. Separate spatial coordinates will use separate copies of these entire blocks, rather than separate copies of their individual components.



The two-block and four-block laws are different-dimensional Gaussian vectors with the response covariances specified earlier. They should not be confused with a collection of independent scalar draws. Independence can instead be used before applying the linear covariance factor.

\subsection{\texorpdfstring{Build one independent Gaussian block for each time and coordinate}{Build one independent Gaussian block for each time and coordinate}}


The path space is \(\Omega=(\mathbb R^{\iota})^{\mathbb N\times\mathrm{Fin}(d)}\). Define its law as the infinite product of the kernel Gaussian. Evaluation at \((n,i)\) is measurable and has exactly the original Gaussian law, by the infinite-product marginal theorem. The coordinate-evaluation family is jointly independent by independence of the product coordinates. Thus a time step retains all within-block covariance while distinct iteration–coordinate pairs are independent. The product's probability property establishes total mass one; the distributional identity proof records both measurability and the exact pushforward equality.



Use a product probability space indexed by time and spatial coordinate, with one Gaussian response block at each index. Product structure gives independence of different indices. Within each block the full response covariance is retained, so the correct velocity–position correlation survives spatial assembly.

\subsection{\texorpdfstring{Transfer block moments and assemble Euclidean innovations}{Transfer block moments and assemble Euclidean innovations}}


Evaluation on the infinite product is measure preserving, so composing the finite-block L\ensuremath{{}^{2}} identity with evaluation gives a finite second moment for each block. Its expectation is zero by the established the distributional identity integral transformation. the first-order Gaussian innovation collects component zero across spatial coordinates into velocity and component one into position. the second-order Gaussian innovation similarly collects each of the four components into a Euclidean vector. the Euclidean identification equips the coordinate functions with the Euclidean norm. Their measurable-map theorems unfold these finite coordinate constructions and apply compositional measurability. These maps supply concrete innovation streams for pathwise integrator recursions.



The Euclidean squared norm is the sum of the coordinate squares. Integrating this finite sum turns scalar moment identities into vector moment identities, with the spatial dimension appearing explicitly. This is the origin of the square-root-dimension factor in first-order error bounds.

\subsection{\texorpdfstring{Turn covariance into second and cross moments}{Turn covariance into second and cross moments}}


The Gaussian coordinates have zero mean. Thus their variance equals \(\int X_i^2\), and their covariance equals \(\int X_iX_j\). The second-moment proof starts from the exact variance formula, rewrites variance as the integral of the centered square, and substitutes the proved zero mean. The cross-moment proof unfolds covariance and substitutes both zero means. Each conclusion is an equality for the constructed Gaussian measure: \(\mathbb E X_i^2=\int_0^hk_i^2\) and \(\mathbb E X_iX_j=\int_0^hk_ik_j\). These formulas can be used in numerical-error calculations without assuming moments of an unspecified noise source.



Polarization supplies cross moments from covariance once the means are known to vanish. Integrability of each product follows from Cauchy–Schwarz and the finite second moments. Consequently the covariance calculation is a legitimate expectation identity, rather than a purely formal matrix comparison.

\subsection{\texorpdfstring{The scalar KLMC transition is an Gaussian pushforward}{The scalar KLMC transition is an Gaussian pushforward}}


the scalar exponential transition measure maps the constructed correlated two-dimensional Gaussian through the scalar KLMC step at a fixed phase state. The map is continuous in its innovation, so the pushforward remains a probability measure. Both output coordinates are a constant frozen-force response plus \(\sqrt{2\gamma}\) times a Gaussian coordinate. Their finite second moments follow by adding a constant L\ensuremath{{}^{2}} function to a scalar multiple of the proved Gaussian L\ensuremath{{}^{2}} coordinate. To compute the transition mean, change variables through the pushforward, integrate the two coordinates separately using their established integrability, and substitute zero Gaussian means. The result is exactly the frozen deterministic position–velocity pair. This is a one-step law and moment result, with no assumption that an abstract transition has the desired mean.



An affine image of a Gaussian remains Gaussian. Its mean is the deterministic frozen-force response, and its covariance is obtained by multiplying the response covariance on both sides by the noise amplitude. This identifies the transition measure of a single exponential step.

\subsection{\texorpdfstring{Correlated transition covariance and independent initialization}{Correlated transition covariance and independent initialization}}


The scalar phase transition retains covariance \(2\gamma C_{01}\) between its velocity and position coordinates. The proof pulls covariance through the transition pushforward, removes the deterministic means, moves both factors \(\sqrt{2\gamma}\) outside covariance, and uses the Gaussian Gram-covariance theorem; their product is \(2\gamma\) when friction is nonnegative. This checks the correlation in the update rather than replacing the two noises by independent draws.



To initialize independently, the canonical probability space is the product \(\mu_0\otimes Q\), where \(Q\) is the countable independent Gaussian-block path measure. The first projection has law \(\mu_0\), the second projection has law \(Q\), and the two projections are independent. The argument constructs this product explicitly, obtains probability from the two probability properties, and applies the product-marginal and independent-projection theorems. No later-time independence or convergence inequality is assumed by this construction.



Fresh noise must also be independent of the initial phase. Taking an independent product of initialization and the innovation stream gives this property at time zero and, by recursion, independence of the next block from the past. The resulting joint laws justify Markov iteration.

\section{\texorpdfstring{The correlated Gaussian particle innovation}{The correlated Gaussian particle innovation}}
\label{proof-module-KIPLMC-Noise-lean}


Velocity and position innovations within a coordinate share a Brownian source and are correlated. Distinct spatial coordinates are independent. This section constructs that joint probability measure using the kinetic library's two-dimensional Gram Gaussian, directly on the parameter-and-particle coordinate set. It works also when one of the dimensions is zero.



The innovation is built from independent spatial response blocks. Packing preserves each block’s correlation and gives the finite-dimensional Gaussian vectors used by the defined first particle algorithm.

\subsection{\texorpdfstring{Construct independent spatial Gaussian blocks and assemble phase vectors}{Construct independent spatial Gaussian blocks and assemble phase vectors}}


The coordinate set is \(I=\{1,\ldots,d_\theta\}\sqcup(\{1,\ldots,N\}\times\{1,\ldots,d_x\})\). For each \(i\in I\), take the shared kinetic Gaussian \((\varepsilon_i,\varepsilon_i\prime)\sim\mathcal N(0,C_{\gamma,h})\), independently across \(i\). Their finite product measure is the independent spatial Gaussian blocks. The assembly map collects the first entries into the velocity vector and the second entries into the position vector. The proof defines these two Euclidean vectors with the Euclidean identification and proves the map measurable coordinate by coordinate. Pushing the product measure forward gives the correlated exponential noise. Probability follows from the Gaussian probability properties and preservation of total mass under the measurable map; no normalization assumption is supplied.



For each scaled particle-space coordinate, draw one two-response Gaussian block. Independent products over coordinates preserve the within-block velocity–position covariance. Packing all first responses and all second responses produces the two full vector innovations used in the update.

\subsection{\texorpdfstring{Prove coordinate independence and retain the within-\allowbreak{}pair covariance}{Prove coordinate independence and retain the within-pair covariance}}


The product-space evaluation maps are independent, so the Gaussian blocks belonging to distinct coordinates are independent. Evaluating the assembled velocity and position at a fixed spatial coordinate recovers its original two-dimensional Gaussian block exactly. The proof composes the two pushforwards, proves equality of the composed function with the block evaluation map by checking the two possible indices, and uses the product-measure evaluation law. The covariance of the pair is therefore exactly the kinetic Gram matrix, already proved positive semidefinite in the shared library. Multiplication by \(\sqrt{2\gamma/N}\) is performed by the particle step, not by this base noise law.



Different spatial coordinate blocks are independent because they are projections of different product factors. The two entries of a single block generally remain correlated through their common response Gram matrix. Both properties are required to recover the intended particle transition law.

\section{\texorpdfstring{A noise-\allowbreak{}driven recursion has a law independent of its realization}{A noise-driven recursion has a law independent of its realization}}
\label{proof-module-UnadjustedLangevin-ChainLaw-lean}


For measurable \(F(x,z)\), the recursion \(X_{n+1}=F(X_n,Z_n)\) has a law determined by the initial law and the full noise law when the initialization and noise are independent.



At index \(k\), the state depends only on initialization and previous innovations. Under the stated independence hypotheses the pair formed by that state and the next innovation has product law. Pushing this pair through the update map yields the transition-law recursion. Induction consequently shows that any two realizations with the same initial law and innovation law have the same state law at every index, independently of their particular sample spaces.

\subsection{\texorpdfstring{The recursion driven by a noise sequence}{The recursion driven by a noise sequence}}


Define \(X_0=x\), \(X_{n+1}=F(X_n,w_n)\) for any update and noise path. No distribution assumptions are part of the recursion; they enter only when a law is assigned to its inputs.



For a fixed input \((x,w)\), the finite-time state uses only the first finitely many coordinates of the innovation path. No probability measure is needed to evaluate it. This separates the deterministic algorithm from its probabilistic realization and permits the same recursion to be applied on canonical product sources and on user-supplied probability spaces.

\subsection{\texorpdfstring{The noise-\allowbreak{}driven chain at zero}{The noise-driven chain at zero}}


The zero clause of the recursion evaluates to the supplied initial state. This defining identity is the base case for all law and measurability inductions.



Before any update, the output is exactly \(x\), independently of the innovation path. Consequently its law under an independent input product must be the initial marginal. The same identity also starts the induction proving that two deterministic recursions with the same update and inputs agree at every finite step.

\subsection{\texorpdfstring{The successor recursion}{The successor recursion}}


The successor clause unfolds to \(X_{n+1}=F(X_n,w_n)\). This exact recursion lets the generic construction later be identified with ULA by induction.



The successor equation identifies which innovation is used next: state \(n\) uses only indices below \(n\), and its successor consumes \(w_n\). Keeping this index convention explicit is important in the independence argument. A fresh-noise theorem later separates the prefix \((w_0,\ldots,w_{n-1})\) from this next coordinate.

\subsection{\texorpdfstring{Measurability of the noise-\allowbreak{}driven chain}{Measurability of the noise-driven chain}}


Induct on \(n\). The base map is the initial-state projection. At the successor, pair the measurable current state with the measurable coordinate evaluation \(w\mapsto w_n\), then compose with the measurable uncurried update. Thus each iterate is a measurable function of initialization and the full noise path.



The current state is a measurable function of the joint initialization/\allowbreak{}path input by induction. Coordinate evaluation on path space is measurable, so pairing it with that current state gives a measurable pair. Composing with the measurable update proves the successor claim. No regularity beyond measurability is needed for this law-level recursion theorem.

\subsection{\texorpdfstring{The marginal law of the recursion}{The marginal law of the recursion}}


Push \(\mu\otimes Q\) through the measurable recursion output at index \(n\). The product gives independent initialization and noise path. This canonical law depends only on the initial law, noise-path law, and update, not on a particular sample-space construction.



If \(F_n\) denotes the finite recursion map, define its canonical distribution by \((F_n)_\#(\mu\otimes Q)\). The product specifies independence of initialization and the entire innovation path while allowing arbitrary dependence within \(Q\). Thus this definition applies before specializing to iid noise and makes the exact inputs determining the output distribution explicit.

\subsection{\texorpdfstring{Its initial law}{Its initial law}}


At index zero the output map is first-coordinate projection. Since \(Q\) has mass one, the first marginal of \(\mu\otimes Q\) is exactly \(\mu\). This proves that the canonical law starts from the supplied initial law.



For a measurable initial-state set \(A\), the probability of the zero-step output lying in \(A\) is \((\mu\otimes Q)(A\times\Omega_Q)=\mu(A)Q(\Omega_Q)=\mu(A)\). The noise source's unit mass is the only normalization needed here. This direct marginal computation proves that the canonical construction preserves the specified initialization.

\subsection{\texorpdfstring{Law identification under independence}{Law identification under independence}}


Independence gives the joint input law \(P_{(X,W)}=P_X\otimes P_W\). Push that identity through the measurable recursion output and use pushforward composition. The resulting output law is exactly map w) n.



Independence means that the joint law of initialization and complete noise path factors into their two marginals. Applying the same measurable recursion to this equality yields equality of output laws. In particular the result does not require the input variables to have a concrete product-space representation; their factorized joint distribution is sufficient.

\subsection{\texorpdfstring{Law identification for independent identical innovations}{Law identification for independent identical innovations}}


Joint independence of the innovations and their common law \(\eta\) identify the entire noise-path law with \(\bigotimes_n\eta\). Substitute this identity into the independent-initialization law theorem. Countability of the time index supplies path-valued measurability.



For iid innovations with common law \(\eta\), every finite cylinder event factors as a product of \(\eta\)-probabilities. These cylinder probabilities identify the complete path law as \(\eta^{\otimes\mathbb N}\). Substituting this law in the independent-input result shows that the algorithm's distribution is determined by \(F\), \(\mu\) and \(\eta\), with no remaining source-specific information.

\subsection{\texorpdfstring{Uniqueness of the resulting marginal laws}{Uniqueness of the resulting marginal laws}}


Apply the canonical iid-law formula on both probability spaces. Equality of their initial laws and the shared innovation law makes the right sides identical, so their iterate laws agree. The theorem compares distributions without coupling the two constructions.



Both realizations are measurable images of the same abstract input probability \(\mu\otimes\eta^{\otimes\mathbb N}\). Hence each measurable output set receives the same probability in both. The conclusion is equality of their distributions at a fixed finite iteration, even though their random variables live on different spaces and no joint coupling between them was specified.

\subsection{\texorpdfstring{Equality with the ULA recursion}{Equality with the ULA recursion}}


Induct on the index with \(F(x,z)=x-h\nabla U(x)+\sqrt{2h}z\). At zero both recursions return the same state. If they agree at index \(n\), their successor equations apply the same function to the same state and innovation, so they agree at index \(n+1\). This identifies the generic noise-driven recursion with ULA pointwise.



With the Euler update \(F(x,z)=x-h\nabla U(x)+\sqrt{2h}\,z\), both pointwise constructions have the same initial clause and the same successor clause. Applying the induction hypothesis inside this common function proves equality at the next step. The argument is purely deterministic, so it can be reused with every supplied innovation sequence before assigning Gaussian laws.

\subsection{\texorpdfstring{Measurability of the ULA update}{Measurability of the ULA update}}


Lipschitzness makes \(\nabla U\) continuous. The joint ULA update is obtained from coordinate projections and that gradient by scalar multiplication and addition, hence is measurable. This discharges the generic recursion's update-measurability premise.



The map \((x,z)\mapsto x-h\nabla U(x)+\sqrt{2h}\,z\) is jointly continuous because the gradient is Lipschitz and hence continuous. Coordinate projections, addition and scalar multiplication preserve continuity. Thus its measurability does not require a step restriction, strong convexity, or Gaussianity; those hypotheses enter only the later moment and convergence estimates.

\subsection{\texorpdfstring{The ULA law under independent identical innovations}{The ULA law under independent identical innovations}}


Use measurability of the Gaussian Euler update in the generic iid law theorem, then rewrite the generic recursion using equality of the two Gaussian Euler recursions. The result identifies the actual ULA construction with its canonical product-noise law.



Apply the general iid realization formula to the measurable Euler update. The deterministic recursion equality identifies its output with the actual Euler chain. Its time-\(n\) law is therefore the canonical pushforward of the independent initialization and iid innovation path. This is the distributional bridge used to compare numerical chains constructed on different coupling extensions.

\subsection{\texorpdfstring{The square-\allowbreak{}integrable chain has the same law}{The square-integrable chain has the same law}}


The \(L^2\) chain representative agrees almost everywhere with the pointwise ULA recursion of the input representatives. Apply the pointwise law theorem and transfer pushforwards across that almost-everywhere equality. Quotient representative choices therefore do not change the algorithm law.



At every fixed finite index, the square-integrable recursion and its ordinary pointwise update agree almost surely. Almost-surely equal measurable random variables have the same pushforward law. The pointwise iid theorem therefore transfers to the square-integrable realization. This transfer ensures that analytic \(L^2\) estimates apply to the actual algorithm law, without depending on choices of representative random functions.

\section{\texorpdfstring{Identify recursive sampling laws with iterated Markov transitions}{Identify recursive sampling laws with iterated Markov transitions}}
\label{proof-module-UnadjustedLangevin-MarkovLaw-lean}


The path-space law of an iid-noise recursion agrees with the repeated one-step transition \(\mu\mapsto F_\#(\mu\otimes\eta)\). The proof explicitly establishes independence of the current state and the fresh innovation.



Write \(P_h\mu\) for the pushforward of \(\mu\otimes\gamma\) by the update map, where \(\gamma\) is the innovation law. For an independent-noise realization, the current state and fresh innovation have precisely that product law. Hence \(\rho_{k+1}=P_h\rho_k\), with \(\rho_0=\mu\). Prefix measurability and independence of a fresh coordinate justify the product identity, and pushforward probability preservation ensures all recursive laws have mass one.

\subsection{\texorpdfstring{The state as a function of a finite noise prefix}{The state as a function of a finite noise prefix}}


Evaluate the same recursion using only a finite prefix \((w_0,\ldots,w_{n-1})\). At the successor, recurse on the first \(n\) coordinates and use the final coordinate as the next innovation. This makes finite dependence explicit for the fresh-noise independence proof.



A time-\(n\) state can be evaluated from an input in \(E\times E^{\{0,\ldots,n-1\}}\). Restricting the prefix for the recursive call and reserving its final coordinate for the next update makes this finite dependence explicit. It provides the measurable factorization through past coordinates needed to prove independence from the next innovation.

\subsection{\texorpdfstring{Agreement with the full recursion}{Agreement with the full recursion}}


Induct on prefix length. Restricting an infinite path to its first \(n\) entries gives the same previous state, and the final finite coordinate is \(w_n\). The successor updates thus agree, identifying finite-prefix and infinite-path constructions pointwise.



At length zero both maps return the initial state. At length \(n+1\), their previous states agree by induction after restricting to the first \(n\) coordinates, and both next innovations are \(w_n\). Applying the common update proves equality. Thus replacing the infinite path by its finite prefix does not alter the algorithm output.

\subsection{\texorpdfstring{Measurability of the finite-\allowbreak{}prefix state}{Measurability of the finite-prefix state}}


The zero map is the initial projection. At the successor, restriction to the earlier finite coordinates and evaluation of the last coordinate are measurable; combine them with the measurable update and the induction hypothesis. This proves the state depends measurably on a finite prefix.



Restriction of a finite coordinate family and evaluation of its last coordinate are measurable product maps. The induction hypothesis gives measurability of the previous state from the shorter prefix. Pair these two measurable outputs and compose with the update. This proves finite-prefix measurability without any probabilistic assumption and allows independence to be passed through the finite recursion.

\subsection{\texorpdfstring{Independence of the prefix and the fresh innovation}{Independence of the prefix and the fresh innovation}}


The index sets \(\{0,\ldots,n-1\}\) and \(\{n\}\) are disjoint. Independence of the coordinate family on the iid product source therefore gives independence of the whole prefix vector and the next coordinate. This is stronger than pairwise independence of individual past entries.



On the iid path space, disjoint groups of coordinate evaluations generate independent sigma-algebras. The first group contains all indices less than \(n\), and the second contains only \(n\). Therefore the entire past prefix vector is independent of the fresh innovation. This joint statement is what permits an arbitrary measurable function of the prefix, including the current state, to inherit independence.

\subsection{\texorpdfstring{The transition acting on a law}{The transition acting on a law}}


Define \(M_\eta(\mu)=F_*(\mu\otimes\eta)\). Its product measure means the innovation is fresh and independent of the current state. The subsequent theorem derives that independence for the canonical recursion, rather than assuming it at each step.



The transition sends \(\mu\) to the law of \(F(X,Z)\) when \(X\sim\mu\), \(Z\sim\eta\), and \(X\perp Z\). Writing it as a product pushforward records these requirements exactly. Knowing only the two marginals would not determine the output if state and innovation were dependent, so the independence verification for the actual recursion is essential.

\subsection{\texorpdfstring{The iterated Markov laws}{The iterated Markov laws}}


Set \(\mu_0=\mu\), \(\mu_{n+1}=M_\eta(\mu_n)\). These are iterated independent-noise transition laws. The recursion alone does not yet assert that a previously defined sample-path algorithm realizes them.



The law recursion begins with \(\mu\) and repeatedly applies the same independent-noise transition. It defines a candidate sequence of algorithm distributions without constructing sample paths. The next results show that the measurable iid path recursion realizes precisely this sequence. This separates a transition-level definition from the theorem identifying it with an actual noise-driven realization.

\subsection{\texorpdfstring{The joint law of a state and fresh innovation}{The joint law of a state and fresh innovation}}


The current state is a measurable function of initialization and the first \(n\) noises. Product-source independence and prefix/\allowbreak{}fresh independence make it independent of \(w_n\). Its marginal is the output law of a noise-driven chain, and the innovation marginal is \(\eta\); their joint law is therefore the product of those actual marginals.



The current state factors through initialization and its finite innovation prefix. Initialization is independent of the entire path, and prefix/\allowbreak{}fresh independence separates the next coordinate from that past. Regrouping these facts shows independence of the current state and next innovation. Their joint law is consequently the product of the already defined current output law and \(\eta\).

\subsection{\texorpdfstring{The successor marginal transition}{The successor marginal transition}}


The pointwise successor is \(F(X_n,w_n)\). Push the proved product joint law of \((X_n,w_n)\) through \(F\); this yields \(\mu_{n+1}=M_\eta(\mu_n)\) for the canonical recursion laws. No later-time state/\allowbreak{}noise independence is assumed.



The successor state is the measurable image of \((X_n,w_n)\) under \(F\). The preceding product-law identity therefore gives its distribution as \(F_\#(\mu_n\otimes\eta)\). This proves the Markov successor relation from the supplied iid source. It does not assume conditional independence as a new property of the recursively constructed state.

\subsection{\texorpdfstring{Agreement of path-\allowbreak{}space and Markov laws}{Agreement of path-space and Markov laws}}


The canonical path law and the iterated Markov law have the same initial value and the same successor equation. Induction on \(n\) identifies them at every index. This closes the connection between sample-path and transition-law definitions.



Equality at time zero follows from the initial marginal computation. If the two law sequences agree at time \(n\), their common successor transition gives agreement at \(n+1\). This induction identifies the path-space and transition-space formulations for every finite iteration. All later kernel estimates can therefore be transferred to the original finite recursion.

\subsection{\texorpdfstring{Probability preservation by one transition}{Probability preservation by one transition}}


The independent product of two probability measures has mass one. A measurable pushforward preserves that total mass, so a transition acting on probability measures is again a probability measure.



The product of current-state and innovation probabilities has total mass one. The preimage of the entire output space under a measurable update is the entire product input space. Hence the pushforward transition also has mass one. This conclusion concerns normalization alone and does not assume that the transition preserves second moments or has an invariant distribution.

\subsection{\texorpdfstring{Probability preservation at every iterate}{Probability preservation at every iterate}}


Induct from the initial probability and apply probability preservation by a measurable Gaussian transition at each successor. This proves normalization of every transition iterate and provides the probability hypothesiss used in later law comparisons.



Assume the initial law is a probability. The one-step normalization theorem shows that each successor law is a probability whenever its predecessor is. Induction proves the property at every finite index. Keeping this fact available is necessary for later expectation inequalities, where the norm of a constant and the comparison between first and second moments use unit total mass.

\subsection{\texorpdfstring{The noise-\allowbreak{}driven chain has the iterated transition law}{The noise-driven chain has the iterated transition law}}


First identify the independent iid realization with the canonical product-path output law. Then apply the equality of the independent-noise marginal recursion and iterated transition laws. Every realization satisfying the original initialization/\allowbreak{}noise assumptions consequently has the same proved Markov iterate law.



An arbitrary admissible realization first has the same law as the canonical independent path recursion, by equality of its joint input distribution. The canonical path recursion then has the same law as the iterated transition. Combining these two identities yields the claimed Markov-law representation for the supplied process, including all initialization and innovation assumptions used to derive it.

\section{\texorpdfstring{Measurable transitions and iterated particle laws}{Measurable transitions and iterated particle laws}}
\label{proof-module-KIPLMC-Laws-lean}


The pathwise algorithm maps become probability transitions by integrating them against independent fresh-noise measures. This section reuses the generic Markov-law construction and independence argument already proved for ULA. For KIPLMC1 it accepts a specified correlated phase innovation law; KIPLMC2 uses an explicit pair of independent standard Gaussian vectors for its two OU half steps. Law equality is kept distinct from convergence rates.



Both explicit particle recursions are identified with their independent-noise Markov iterates. The exponential transition retains correlated response pairs, while OBABO retains two distinct independent thermostat noises.

\subsection{\texorpdfstring{Construct the OBABO noise law and prove both step maps measurable}{Construct the OBABO noise law and prove both step maps measurable}}


the two independent thermostat noises is the product of two standard Gaussian measures on the particle Euclidean space. Its product probability property gives total mass one. Lipschitzness of the particle force implies continuity. Unfolding the exponential or OBABO map shows that each is built from this continuous force, coordinate projections, addition and scalar multiplication, so each uncurried state–noise map is continuous and hence measurable. The argument supplies force continuity first, applies measurability of continuous maps, and discharges the finite composition with continuity and measurability.



The two thermostat noises form an independent product of Gaussian vector measures. Continuity of the force makes each finite-stage update measurable. Both particle algorithms therefore define probability transitions by pushing forward the state/\allowbreak{}noise product.

\subsection{\texorpdfstring{Define numerical laws by independent-\allowbreak{}noise pushforwards}{Define numerical laws by independent-noise pushforwards}}


the KIPLMC1 laws applies the shared Markov recursion with the exponential step, an initial phase measure \(\mu\), and a supplied innovation measure \(\eta\). the KIPLMC2 laws uses the OBABO step and its fixed two-Gaussian noise measure. Each successor law is the pushforward of the product of the current-state law and the fresh-noise law. The probability theorems invoke the proved step measurability and probability of the initial and innovation measures, so every iterate is a probability measure. For KIPLMC1 the covariance specification of \(\eta\) remains explicit and separate; probability alone does not identify the intended correlated Gaussian law.



At each step take a fresh innovation independent of the present state and apply the specified map. Integrating this update against its innovation measure defines one transition. Repeated pushforward gives the numerical laws without assuming an underlying chain with the desired estimates.

\subsection{\texorpdfstring{Prove the explicit exponential recursion has the Markov law}{Prove the explicit exponential recursion has the Markov law}}


First show that the KIPLMC1 sequence equals the generic noise-driven recursion. Induction is exact: both start at the same phase and their successor expressions use the same step and innovation number. The law theorem then assumes a measurable initial random phase, a jointly independent innovation sequence with common law \(\eta\), and independence of initialization from the whole sequence. The shared generic result derives the law at every step; no additional per-step state–noise independence premise is inserted. Rewriting the explicit recursion by the proved equality yields the pushforward identity with the KIPLMC1 laws. This theorem establishes actual-law correspondence, while a quantitative sampling estimate still requires moment and coupling analysis for the specified Gaussian innovation.



The explicit recursion is a measurable function of initialization and the innovation prefix. Independence of the next block gives the one-step transition identity. Induction then identifies its law with the iterated exponential transition measure.

\subsection{\texorpdfstring{The independent-\allowbreak{}Gaussian OBABO trajectory has the constructed transition law}{The independent-Gaussian OBABO trajectory has the constructed transition law}}


the KIPLMC2 sequence iterates the OBABO map with the two thermostat vectors from each supplied noise pair. If these pairs are independent across time, have the product-standard-Gaussian law the two independent thermostat noises, and are independent of initialization, then the pushforward of the pathwise recursion equals the KIPLMC2 laws at every step. The generic independence theorem derives independence of the current state from the fresh pair from these global assumptions. The proof is a direct reuse of that theorem with the established measurable OBABO map. Independence within each pair is encoded by its product Gaussian law; it is not confused with the independent sequence of pairs. This completes the algorithm-law correspondence, while contraction and the discrete invariant-measure bias remain separate quantitative obligations.



The same argument applies to the two-thermostat recursion, keeping the two noises distinct and independent. Their common block law is used at every step. Thus the trajectory’s pushforward agrees with the constructed OBABO Markov law.

\subsection{\texorpdfstring{Choose the constructed kinetic Gaussian for Algorithm 1}{Choose the constructed kinetic Gaussian for Algorithm 1}}


The generic law identity accepts a specified iid innovation measure. the Gaussian KIPLMC1 laws supplies the correlated Gaussian measure constructed in the explicit correlated particle-noise construction, whose spatial pairs have the shared kinetic Gram covariance. Thus the transition and its particle scaling are explicit. The earlier iid realization theorem specializes to this measure, with ordinary independence of initialization and fresh innovations. A convergence theorem must still derive a target comparison for these laws; the definition does not encode one.



For Algorithm 1 the innovation law is the constructed exponential response Gaussian, including its off-diagonal covariance. Substituting that law specializes the generic transition to the correlated-Gaussian sampler. Independent velocity and position noises would define a different law.

\section{\texorpdfstring{Couplings and quadratic transport}{Couplings and quadratic transport}}
\label{proof-module-UnadjustedLangevin-Coupling-lean}


Let \(H\) be a second-countable normed space with its Borel sigma-algebra. A coupling of \(\mu,\nu\) is a measure \(\gamma\) on \(H\times H\) with the prescribed two marginals. Its quadratic cost is \(c(\gamma)=(\int\|x-y\|^2d\gamma)^{1/2}\), and \(W_2(\mu,\nu)\) is the infimum of these costs. Costs are initially allowed to be infinite. Exhibiting one finite-cost coupling proves that the transport distance is finite before it is used in ordinary real inequalities.



For a joint law \(\gamma\), the quadratic cost is the root mean square separation of its two coordinates. A particular coupled pair provides an upper bound on the transport infimum, even when that pair is not optimal. The proof consequently proceeds through actual random variables and their \(L^2\) error. Finite error moments are established before any extended nonnegative quantity is interpreted as an ordinary real distance.

\subsection{\texorpdfstring{Definition of a coupling}{Definition of a coupling}}


A coupling of measures \(\mu,\nu\) on \(H\) is a joint measure \(\gamma\) on \(H\times H\) with those two marginals. Denote the family by \(\Gamma(\mu,\nu)\).



Choose the joint measure \(\gamma\) on pairs \((x,y)\).



Require \((\operatorname{pr}_1)_\#\gamma=\mu\) and \((\operatorname{pr}_2)_\#\gamma=\nu\).



The two marginal conditions say that integrating a test depending only on the first coordinate reproduces its integral under \(\mu\), and likewise for the second coordinate and \(\nu\). Dependence between the two coordinates is unrestricted. This freedom is precisely what lets a synchronous construction make their separation much smaller than under independent sampling.

\subsection{\texorpdfstring{Definition of quadratic coupling cost}{Definition of quadratic coupling cost}}


For a joint measure \(\gamma\), define

\[c(\gamma)=\left(\int\|x-y\|^2\,d\gamma(x,y)\right)^{1/2}\in[0,+\infty].\]

Take the extended \(L^2(\gamma)\) norm of the displacement map \((x,y)\mapsto x-y\).



The square root is taken after the squared separation has been integrated. This is a root mean square cost and differs from the mean distance. The extended value allows couplings with infinite second displacement moment; later finite-moment lemmas identify when this cost is an ordinary real number.

\subsection{\texorpdfstring{Definition of the transport infimum}{Definition of the transport infimum}}


Define the extended quadratic transport distance by

\[W_2(\mu,\nu)=\inf_{\gamma\in\Gamma(\mu,\nu)}c(\gamma).\]

Take the infimum over all couplings without selecting a minimizer.



The infimum optimizes over dependence while fixing both marginal laws. Every admissible coupling gives an upper bound, and nearly optimal couplings suffice for passing affine estimates to the infimum. The definition therefore does not require choosing an optimizer or proving that an optimal transport plan exists.

\subsection{\texorpdfstring{Real conversion of transport distance}{Real conversion of transport distance}}


The transport cost is initially allowed to lie in \([0,\infty]\). Whenever \(W_2(\mu,\nu)<\infty\), we use the same notation for its ordinary nonnegative real value. Every finite transport estimate below first proves this finiteness, usually by exhibiting a coupling whose displacement belongs to \(L^2\).



Indeed, if \(\gamma\in\Gamma(\mu,\nu)\) has finite cost \(c(\gamma)\), then \(0\le W_2(\mu,\nu)\le c(\gamma)<\infty\). All subsequent real inequalities involving this distance are consequently inequalities between finite numbers. Without the finite-cost hypothesis, the extended distance remains available, but it cannot be used in ordinary subtraction or division.

\subsection{\texorpdfstring{Every coupling supplies an upper bound}{Every coupling supplies an upper bound}}


For every \(\gamma\in\Gamma(\mu,\nu)\), \(W_2(\mu,\nu)\le c(\gamma)\).



An infimum is at most each member of the family. Apply this to the chosen coupling.



The comparison follows directly from the defining lower-bound property of an infimum. It is useful even for a coupling selected for analytical convenience rather than optimality. In particular, the joint law generated by running two algorithms on the same noise can be inserted immediately once its marginals are identified.

\subsection{\texorpdfstring{The joint law of two random vectors}{The joint law of two random vectors}}


If \(X,Y:\Omega\to H\) are measurable, then \((X,Y)_\#P\) is a coupling of \(X_\#P\) and \(Y_\#P\).



Set \(\gamma=(\omega\mapsto(X(\omega),Y(\omega)))_\#P\).



Compose the pushforward with the first projection. Since \(\operatorname{pr}_1(X(\omega),Y(\omega))=X(\omega)\), the first marginal is \(X_\#P\).



The same calculation with the second projection gives the second marginal \(Y_\#P\).



For a measurable test \(f\), \(\int f(x)d\gamma(x,y)=\mathbb Ef(X)\), while the corresponding second-coordinate identity is \(\int f(y)d\gamma=\mathbb Ef(Y)\). These are the two marginal statements in integral form. No independence between \(X\) and \(Y\) is imposed or needed.

\subsection{\texorpdfstring{Joint laws from almost everywhere measurable vectors}{Joint laws from almost everywhere measurable vectors}}


If \(X,Y\) are measurable almost everywhere with respect to \(P\), their joint law is still a coupling of their two laws.



Again define \(\gamma=(X,Y)_\#P\).



Apply pushforward composition for almost everywhere measurable maps to identify the first marginal with \(X_\#P\).



Apply its second-projection version to identify the second marginal with \(Y_\#P\).



Replace each almost-everywhere measurable variable by a measurable version on a full-measure set. Their joint pushforward is unchanged by this replacement, as are both marginals. This extends the coupling construction to representatives of \(L^2\) classes, whose natural pointwise descriptions need only hold almost everywhere.

\subsection{\texorpdfstring{Cost of the joint-\allowbreak{}law coupling}{Cost of the joint-law coupling}}


For measurable \(X,Y:\Omega\to H\),

\[c((X,Y)_\#P)=\|X-Y\|_{2,P}.\]

Expand the two definitions. The cost is the \(L^2\) norm of \((x,y)\mapsto x-y\) under the joint law.



Change variables through the pushforward. The displacement composed with the paired random vector is \(\omega\mapsto X(\omega)-Y(\omega)\), giving the identity.



The integrand under the joint law is the squared norm of coordinate subtraction. Composing that subtraction with \(\omega\mapsto(X(\omega),Y(\omega))\) produces exactly \(X-Y\). Taking the nonnegative square root after the pushforward integral gives equality of costs, rather than only a bound.

\subsection{\texorpdfstring{A coupled pair bounds transport}{A coupled pair bounds transport}}


For measurable \(X,Y\),

\[W_2(X_\#P,Y_\#P)\le\|X-Y\|_{2,P}.\]

Use the joint law as an admissible coupling, then replace its cost by the \(L^2\) distance just computed.



First establish that the joint law is admissible, then evaluate its cost, and only then compare with the infimum. This three-step argument isolates the distributional information in the marginals from the analytic information in \(\mathbb E\|X-Y\|^2\). It is the basic bridge from a common-noise error estimate to Wasserstein convergence.

\subsection{\texorpdfstring{The almost everywhere measurable transport bound}{The almost everywhere measurable transport bound}}


The same transport comparison holds for almost everywhere measurable \(X,Y\):

\[W_2(X_\#P,Y_\#P)\le\|X-Y\|_{2,P}.\]

The almost everywhere measurable joint-law construction supplies an admissible coupling, whose cost bounds the infimum.



The pushforward change of variables identifies that cost with \(\|X-Y\|_{2,P}\).



The coupling and cost are independent of changes on null sets, so the measurable-version argument carries the preceding proof over unchanged. This formulation is needed when the compared states are naturally constructed in a Hilbert space of random-variable classes rather than as globally specified measurable representatives.

\subsection{\texorpdfstring{Finite coupled cost gives a real bound}{Finite coupled cost gives a real bound}}


Let \(X,Y:\Omega\to H\) be measurable and assume \(X-Y\in L^2(P;H)\). Their joint law is a coupling of their marginal laws, with finite quadratic cost

\[\left(\int\|x-y\|^2d(X,Y)_\#P(x,y)\right)^{1/2}
=\left(\int\|X-Y\|^2dP\right)^{1/2}=\|X-Y\|_{2,P}<\infty.\]


The defining transport infimum is no larger than this particular cost. It is therefore finite, and the ordinary real estimate is

\[W_2(X_\#P,Y_\#P)\le\|X-Y\|_{2,P}.\]


Only the displacement moment is needed for this argument. Separate second moments of the two variables, or independence between them, are not required for the coupling-cost comparison.

\subsection{\texorpdfstring{The transport bound for L\ensuremath{{}^{2}} classes}{The transport bound for L² classes}}


For \(X,Y\in L^2(P;H)\), choose measurable representatives. Their joint law has the two representative laws as marginals. On a common full-measure set, the difference of the representatives equals a representative of the class \(X-Y\). Hence the joint-law displacement energy is exactly \(\|X-Y\|_{L^2}^2\).

\[W_2(X_\#P,Y_\#P)\le\|X-Y\|_{L^2}<\infty.\]


The finite upper bound also establishes finiteness of the transport distance. Changing either representative on a null set changes neither the law nor the norm, so this comparison is well defined for the original equivalence classes.

\subsection{\texorpdfstring{The real L\ensuremath{{}^{2}}-\allowbreak{}to-\allowbreak{}transport bridge}{The real L²-to-transport bridge}}


The preceding joint-law construction gives, for every pair \(X,Y\in L^2(P;H)\), the finite comparison

\[0\le W_2(X_\#P,Y_\#P)\le\|X-Y\|_{L^2}.\]


Its right side is an ordinary Hilbert norm, and its finiteness proves that the left side is an ordinary finite transport distance as well. Consequently any upper bound on the square-mean error of these two classes is also an upper bound on the distance of their laws.



This is the form used after the synchronous error recurrence has been summed. The reference marginal is then the target law, so the coupled Hilbert-space error directly controls convergence in Wasserstein distance.

\subsection{\texorpdfstring{Pointwise domination passes to L\ensuremath{{}^{2}}}{Pointwise domination passes to L²}}


Let \(q\ge0\), let \(A\) be almost everywhere strongly measurable, and suppose \(\|A(\omega)\|\le q\|D(\omega)\|\) almost everywhere. Then

\[\|A\|_{2,P}\le q\|D\|_{2,P}.\]

Apply monotonicity and positive homogeneity of the extended \(L^2\) seminorm to the almost everywhere domination. No independence is required.



Squaring the pointwise domination gives \(\|A\|^2\le q^2\|D\|^2\), since both norms and \(q\) are nonnegative. Integral monotonicity then yields the same relation for squared \(L^2\) norms. Taking square roots recovers the factor \(q\), with no assumption on the dependence of \(A\) and \(D\).

\subsection{\texorpdfstring{One-\allowbreak{}step error from contraction and a defect}{One-step error from contraction and a defect}}


Suppose \(q\ge0\), \(A\) is almost everywhere strongly measurable, and almost everywhere

\[X'-Y'=A-V,\qquad\|A\|\le q\|X-Y\|.\]

Then

\[\|X'-Y'\|_{2,P}\le q\|X-Y\|_{2,P}+\|V\|_{2,P}.\]

Replace \(X'-Y'\) by \(A-V\) in the seminorm using the almost everywhere error identity.



Minkowski gives \(\|A-V\|_2\le\|A\|_2+\|V\|_2\). Pointwise contraction bounds the first term by \(q\|X-Y\|_2\).



The defect sign does not affect its norm: \(\|-V\|_2=\|V\|_2\). Hence Minkowski adds the local defect size to the contracted previous error. This addition, rather than a squared-norm expansion, avoids requiring any orthogonality between the existing error and the defect.

\subsection{\texorpdfstring{One-\allowbreak{}step transport estimate}{One-step transport estimate}}


Under the preceding decomposition and contraction hypotheses, if \(X',Y'\) are measurable, then

\[W_2(X'_\#P,Y'_\#P)\le q\|X-Y\|_{2,P}+\|V\|_{2,P}.\]

The joint-law coupling bounds transport by \(\|X'-Y'\|_2\). Compose this with the synchronous \(L^2\) estimate.



The final pair itself supplies the admissible coupling for the next marginal laws. Therefore the one-step stochastic estimate need not compare two transition kernels abstractly; it can use the explicit shared-noise pair already constructed on the common space. The finite-moment version later ensures all displayed real quantities are legitimate.

\subsection{\texorpdfstring{The finite-\allowbreak{}moment one-\allowbreak{}step inequality}{The finite-moment one-step inequality}}


Assume \(q\ge0\), \(A\) is almost everywhere strongly measurable, and

\[X'-Y'=A-V,\qquad\|A\|\le q\|X-Y\|\quad\text{almost everywhere}.\]


If \(X-Y,V\in L^2(P;H)\), then the domination makes \(A\) square-integrable. Therefore \(A-V\), and hence \(X'-Y'\), is also square-integrable. All norms in the following estimate are consequently finite:

\[\begin{aligned}
\|X'-Y'\|_{2,P}&=\|A-V\|_{2,P}\\
&\le\|A\|_{2,P}+\|V\|_{2,P}\\
&\le q\|X-Y\|_{2,P}+\|V\|_{2,P}.
\end{aligned}\]


The first inequality is Minkowski, and the second follows by squaring the nonnegative pointwise domination, integrating, and taking square roots. This proves the numerical affine recurrence together with the finite moment needed to interpret it.

\section{\texorpdfstring{Passing from a chosen initial coupling to an infimum}{Passing from a chosen initial coupling to an infimum}}
\label{proof-module-UnadjustedLangevin-InitialCoupling-lean}


A chosen initial coupling gives an estimate in terms of its own cost. To obtain the initial Wasserstein distance, repeat the construction for every finite-cost coupling with the same output laws and the same constants. Approximate minimizers then let the chosen cost approach \(W_2(\mu,\nu)\). An independent product extension preserves each initial coupling while supplying the common noise; its numerical output law depends only on the initial marginal and the innovation law.



To replace a chosen initial error by a transport distance, the later construction must work for every admissible initial coupling with the same constants. For any \(\varepsilon>0\), a coupling can then be chosen with cost below \(W_2(\mu,\nu)+\varepsilon\). Apply the uniform affine estimate and let \(\varepsilon\downarrow0\). This argument uses only the defining property of an infimum; existence of an optimal coupling is unnecessary. Infinite-cost couplings require a separate check when the affine coefficient is zero.

\subsection{\texorpdfstring{A uniform affine estimate passes to the infimum}{A uniform affine estimate passes to the infimum}}


Let \(a\ge0\) be finite, \(b\in[0,+\infty]\), and \(\Gamma(\mu,\nu)\ne\varnothing\). If the fixed output measures satisfy

\[W_2(\mu',\nu')\le ac(\gamma)+b\quad\text{for every }\gamma\in\Gamma(\mu,\nu),\]

then \(W_2(\mu',\nu')\le aW_2(\mu,\nu)+b\).



Use the nonempty coupling family and expand transport distances as infima.



Multiplication by finite \(a\ge0\) and addition of \(b\) commute with the nonempty extended-real infimum, including \(a=0\). Thus \(aW_2(\mu,\nu)+b=\inf_\gamma(ac(\gamma)+b)\).



The output distance lies below every member of this last family, hence below its infimum.



For \(a>0\), choose couplings whose costs approach the infimum and apply the estimate to each one. Multiplication by \(a\) and addition of \(b\) commute with this monotone affine approximation. At \(a=0\), the conclusion is simply the common constant bound \(b\), so no optimizer or limiting transport cost is required.

\subsection{\texorpdfstring{A finite infimum supplies a finite-\allowbreak{}cost coupling}{A finite infimum supplies a finite-cost coupling}}


If \(W_2(\mu,\nu)<+\infty\), then some \(\gamma\in\Gamma(\mu,\nu)\) has \(c(\gamma)<+\infty\). Attainment of the infimum is not asserted.



The coupling-cost infimum is strictly below \(+\infty\) by the finiteness hypothesis.



An infimum strictly below a threshold has a member strictly below that threshold. Apply this with threshold \(+\infty\).



If every admissible coupling had infinite cost, the infimum of their costs would also be infinite. Thus a finite infimum provides at least one finite-cost plan. More quantitatively, an approximate minimizer with cost below the finite infimum plus one is already enough for all later finiteness arguments.

\subsection{\texorpdfstring{The real uniform-\allowbreak{}coupling estimate}{The real uniform-coupling estimate}}


Let \(a,b\ge0\) be finite and assume \(W_2(\mu,\nu)<+\infty\). If \(W_2(\mu',\nu')\le ac(\gamma)+b\) for every initial coupling, then

\[W_2(\mu',\nu')\le aW_2(\mu,\nu)+b.\]

A finite initial distance supplies a coupling. Apply the extended uniform-coupling lemma to obtain the desired inequality in extended reals.



The product \(aW_2(\mu,\nu)\) is finite.



Its sum with finite \(b\) is finite, so the extended output inequality permits monotone real conversion.



Distribute real conversion through the finite sum and product to obtain the real inequality.



The extended estimate is first applied while retaining finite transport and coupling values. Their embeddings and real values then agree, so the affine inequality may be read in ordinary real arithmetic. This separates the infimum argument from the finite-value conversion and avoids drawing information from an infinite real-converted cost.

\subsection{\texorpdfstring{Finite-\allowbreak{}cost initial couplings suffice}{Finite-cost initial couplings suffice}}


Let \(a,b\ge0\) be finite and assume \(W_2(\mu,\nu)<\infty\) and \(W_2(\mu',\nu')<\infty\). Suppose every finite-cost initial coupling satisfies

\[W_2(\mu',\nu')\le a c(\gamma)+b\qquad
(\gamma\in\Gamma(\mu,\nu),\ c(\gamma)<\infty).\]


For every \(\varepsilon>0\), finiteness of the initial infimum supplies a coupling \(\gamma_\varepsilon\) with

\[c(\gamma_\varepsilon)<W_2(\mu,\nu)+\varepsilon<\infty.\]


Insert this coupling into the uniform estimate. For \(a>0\), let \(\varepsilon\downarrow0\) to obtain

\[W_2(\mu',\nu')\le aW_2(\mu,\nu)+b.\]


If \(a=0\), any one of these finite-cost couplings already gives the conclusion \(W_2(\mu',\nu')\le b\). Thus no expression involving zero times infinite cost is needed. Infinite-cost couplings cannot improve a finite infimum, and no minimizing coupling is assumed to exist.

\section{\texorpdfstring{Iteration of the scalar error estimate}{Iteration of the scalar error estimate}}
\label{proof-module-UnadjustedLangevin-Recurrence-lean}


Once coupling gives \(e_{k+1}\le qe_k+b\), the remaining calculation is deterministic. We retain the exact geometric sum before obtaining the simpler bias-floor bound. The upper envelope approaches \(b/(1-q)\); this does not assert that the actual error converges to that value.



The scalar sequence need not have an algorithmic interpretation for these lemmas. The proofs separate the exact geometric sum from its uniform bias-floor envelope and explicitly require the contraction factor in \([0,1)\). Later the coupling construction proves that the actual ULA error satisfies this recurrence.



Each local error inserted \(j\) steps before the final time is damped by a factor \(q^j\). The finite geometric sum records these individual contributions, rather than replacing them immediately by an infinite sum. When \(0\le q<1\), its limit is \(1/(1-q)\); hence persistent defects create the floor \(b/(1-q)\). The homogeneous part tends to zero because \(q^n\to0\), and every tolerance strictly above the floor is eventually reached.

\subsection{\texorpdfstring{Iteration gives a geometric sum}{Iteration gives a geometric sum}}


If \(q\ge0\) and \(e_{k+1}\le qe_k+b\) for every \(k\), then

\[e_n\le q^ne_0+b\sum_{j=0}^{n-1}q^j.\]

Induct on \(n\). At zero the sum is empty and both sides equal \(e_0\). Assume the bound at \(n\).



Multiply the induction inequality by \(q\ge0\), preserving its direction.



Write \(S_n=\sum_{j=0}^{n-1}q^j\). The geometric identity gives \(qS_n+1=S_n+q^n\); multiply this equality by \(b\).



Combine the one-step bound with the multiplied induction bound. Since \(S_{n+1}=S_n+q^n\), obtain \(e_{n+1}\le q^{n+1}e_0+bS_{n+1}\).



The empty sum at \(n=0\) is zero, so the initial condition is preserved exactly. The next step multiplies every previous contribution by \(q\) and inserts one new defect of size \(b\). This shifts the powers in the geometric sum and explains the indexing \(j=0,\ldots,n-1\).

\subsection{\texorpdfstring{The exact finite-\allowbreak{}horizon envelope}{The exact finite-horizon envelope}}


If \(0\le q<1\) and \(e_{k+1}\le qe_k+b\), then

\[e_n\le q^ne_0+\frac{b(1-q^n)}{1-q}.\]

Apply the geometric-sum bound and evaluate the sum as \((q^n-1)/(q-1)\). Its denominator is nonzero because \(q<1\).



Change the signs of numerator and denominator simultaneously and put \(b\) in the numerator to obtain \(b(1-q^n)/(1-q)\).



The numerator \(1-q^n\) retains the finite-horizon effect of the defects: no defect has acted before time zero, and only \(n\) of them have appeared by index \(n\). Dividing by \(1-q>0\) evaluates the finite sum. The exact expression is useful when a finite iteration count is required rather than only an asymptotic floor.

\subsection{\texorpdfstring{A uniform bias-\allowbreak{}floor bound}{A uniform bias-floor bound}}


If \(0\le q<1\), \(b\ge0\), and \(e_{k+1}\le qe_k+b\), then

\[e_n\le q^ne_0+\frac b{1-q}.\]

Start from the exact finite-horizon bound. Its denominator \(1-q\) is positive.



Since \(q^n\ge0\) and \(b\ge0\), the numerator satisfies \(b(1-q^n)\le b\).



Divide by the positive denominator, add \(q^ne_0\), and compose with the exact bound.



Since \(0\le q^n\le1\), dropping \(1-q^n\) enlarges only the nonnegative forcing contribution. The initial term remains \(q^ne_0\). This yields a simpler envelope whose bias term is independent of the iteration count and whose transient term decreases geometrically.

\subsection{\texorpdfstring{Limit of the upper envelope}{Limit of the upper envelope}}


For \(0\le q<1\) and arbitrary real \(b,e_0\),

\[q^ne_0+\frac b{1-q}\longrightarrow\frac b{1-q}\quad(n\to\infty).\]

The geometric powers satisfy \(q^n\to0\).



Multiply this limit by \(e_0\) and add the constant \(b/(1-q)\).



The envelope is the sum of a constant and \(q^ne_0\). For \(0\le q<1\), the latter tends to zero for every fixed real initial value. Hence the envelope tends to \(b/(1-q)\). A bound on the sequence by this envelope controls its limiting upper behavior; it need not imply that the sequence itself converges.

\subsection{\texorpdfstring{Every accuracy above the floor is eventually attained}{Every accuracy above the floor is eventually attained}}


Assume \(0\le q<1\), \(b\ge0\), and \(e_{k+1}\le qe_k+b\). For every \(\varepsilon>b/(1-q)\), the errors satisfy \(e_n<\varepsilon\) for every sufficiently large \(n\).



The upper envelope tends to a value strictly below \(\varepsilon\), so it is eventually below \(\varepsilon\).



At each such index, place \(e_n\) below the envelope by the recurrence bound and then below \(\varepsilon\).



Given \(\varepsilon>b/(1-q)\), the positive gap to the floor can absorb the eventually small homogeneous term. Choose an index after which \(q^ne_0\) lies below that gap. The uniform recurrence bound then places every later error below \(\varepsilon\), which is the precise eventual-accuracy conclusion.

\section{\texorpdfstring{The shared Hilbert-\allowbreak{}space proof and both step regimes}{The shared Hilbert-space proof and both step regimes}}
\label{proof-module-UnadjustedLangevin-HilbertDalalyan-lean}


Work on a probability space in \(\mathcal H=L^2(P;H)\), with lifted gradient \(\mathcal G[X]=[\nabla U\circ X]\). Supply continuous reference segments satisfying the drift integral equation, joining at their endpoints, and having the stationary target law at grid times. Their endpoint innovations are the same as those driving the numerical chain. The drift and noise moment bounds yield the local defect; sharp deterministic Euler contraction and Minkowski then yield the error recurrence. At this stage the reference is supplied; the independent Brownian construction later supplies its existence and stationary laws.



The reference lives in the Hilbert space of square-integrable random vectors, so its drift equation uses a Bochner time integral in that space. Its grid endpoints share the numerical noise, have the prescribed stationary law, and join successive continuous segments. These assumptions supply the local defect estimate with constant \(33/20\). Summing the error recurrence gives the floor by dividing that defect by \(1-q\); the division is valid only after proving \(0\le q<1\).

\subsection{\texorpdfstring{The common L\ensuremath{{}^{2}} algorithm}{The common L² algorithm}}


For an initial \(X_0\in L^2(P;H)\) and innovations \(\xi_k\in L^2(P;H)\), set \(X_{k+1}=X_k-h\widehat{\nabla U}(X_k)+\sqrt{2h}\,\xi_k\).



The initial iterate is the supplied random vector.



Each successor uses the actual-gradient Euler update.



The drift lift evaluates the original gradient almost everywhere, so this recursion represents the ordinary ULA method on random inputs. Its Hilbert-space formulation is useful because the reference equation and defect integral then live in the same space. Each step stays square-integrable by the Lipschitz growth estimate for the gradient.

\subsection{\texorpdfstring{Unfolding one step}{Unfolding one step}}


The next iterate is exactly one L\ensuremath{{}^{2}} ULA update of the current iterate.



Evaluate the successor clause of the definition.



Exposing one successor step identifies which innovation must be shared with the reference endpoint. Subtracting a reference equation using that same innovation removes the stochastic increment exactly. This is the algebraic reason that the later error recurrence contains only the deterministic Euler difference and the local defect.

\subsection{\texorpdfstring{The continuous reference assumptions}{The continuous reference assumptions}}


For an \(L\)-Lipschitz gradient, step \(h\), noise vectors \(\xi_k\in\mathcal H\), and a measure \(\pi\), a reference consists of segments \(Y_{k,t}\in\mathcal H\) and increments \(B_{k,t}\in\mathcal H\) satisfying the following conditions.



Supply the two families \((Y_{k,t})\) and \((B_{k,t})\) in the Hilbert space \(\mathcal H\).



Require continuity of \(t\mapsto Y_{k,t}\) in the \(L^2\) norm for every \(k\).



Require compatible endpoints, common endpoint noise, and stationary grid laws:

\[Y_{k,h}=Y_{k+1,0},\qquad B_{k,h}=\sqrt{2h}\,\xi_k,\qquad(Y_{k,0})_\#P=\pi.\]

For \(0\le t\le h\), require

\[Y_{k,t}-Y_{k,0}=-\int_0^t g(Y_{k,s})\,ds+B_{k,t}.\]

Require the score moment bound \(\|g(Y_{k,t})\|_2\le\sqrt{Ld}\) throughout each closed step interval.



Require the increment moment bound \(\|B_{k,t}\|_2\le\sqrt{2d}\sqrt t\) on the same interval.



The segment-joining condition makes the local initial point of segment \(k+1\) equal to the endpoint of segment \(k\). The stationary grid-law condition identifies each comparison marginal with \(\pi\), while the endpoint-noise condition identifies the reference's final increment with the algorithm innovation. The two moment assumptions have different roles: drift controls the deterministic segment increment, and noise controls its square-root growth.

\subsection{\texorpdfstring{Full convergence calculation from the integral reference}{Full convergence calculation from the integral reference}}


Suppose a reference with the preceding properties is supplied, \(h\ge0\), and \(Lh\le2\). Let \(0\le q<1\), and assume the lifted Euler map contracts:

\[\|(X-hg(X))-(Y-hg(Y))\|_2\le q\|X-Y\|_2.\]

For any initial vector \(X_0\in\mathcal H\), the \(L^2\) chain driven by the reference endpoint noise satisfies

\[W_2((X_n)_\#P,\pi)\le q^n\|X_0-Y_{0,0}\|_2+\frac{(33/20)Lh\sqrt{hd}}{1-q}.\]

The statement is conditional on the supplied reference, and the initial term is the cost of the chosen coupling.



Set \(g=\widehat{\nabla U}\), name the actual iterates \(X_k\), and define

\[b=\frac{33}{20}Lh\sqrt{hd},\qquad D_k=\int_0^h\bigl(g(Y_{k,t})-g(Y_{k,0})\bigr)\,dt.\]

The bound \(b\) is nonnegative. Apply the local-defect theorem to the actual lifted gradient and the supplied equation and moment bounds, deriving \(\|D_k\|_2\le b\).



We claim the reference endpoints obey

\[Y_{k+1,0}=Y_{k,0}-hg(Y_{k,0})+\sqrt{2h}\,\xi_k-D_k.\]

The lifted gradient is Lipschitz and the reference segment is continuous, so their composition is continuous and Bochner integrable on \([0,h]\).



Apply the Euler-defect identity at \(t=h\), then substitute the segment join and endpoint noise identities. This proves the claimed reference update. Its defect has a minus sign.



Fix \(k\) and compare the actual update with the reference endpoint equation, both using the same \(\xi_k\).



Cancel their identical noise terms. Because the reference subtracts \(D_k\),

\[X_{k+1}-Y_{k+1,0}=\bigl(X_k-hg(X_k)\bigr)-\bigl(Y_{k,0}-hg(Y_{k,0})\bigr)+D_k.\]

Apply the triangle inequality, Euler contraction, and the defect bound. Setting \(e_k=\|X_k-Y_{k,0}\|_2\) gives \(e_{k+1}\le qe_k+b\).



The actual pair supplies a finite-cost coupling, so its transport distance is at most \(e_n\). The second marginal is \(\pi\) by the stationary grid-law assumption.



Apply the scalar recurrence bound to the derived inequality and compose it with the transport comparison.



The defect has size \(b=(33/20)Lh\sqrt{hd}\), and this size is uniform in \(k\) because the supplied moment bounds hold on every segment. The recurrence is consequently homogeneous in its constants. Its initial value remains the norm of the supplied pair, explaining why this theorem is a chosen-coupling statement before a later infimum argument is applied.

\subsection{\texorpdfstring{The small-\allowbreak{}step bound}{The small-step bound}}


Suppose \(U\) is differentiable and \(m\)-strongly convex on \(H\), \(\nabla U\) is \(L\)-Lipschitz, and \(0<m\le L\). Let \(0<h\le2/(m+L)\), and supply a reference with the preceding properties. For any initial vector \(X_0\in\mathcal H\),

\[W_2((X_n)_\#P,\pi)\le(1-mh)^n\|X_0-Y_{0,0}\|_2+\frac{33}{20}\frac Lm\sqrt{hd}.\]

This uses a chosen initial coupling and a supplied stationary reference. Replacing the initial cost by \(W_2(\mathcal L(X_0),\pi)\) needs the additional uniform family of reference couplings explained in the final section.



Since \(m+L>0\), clear the positive denominator in the step hypothesis to obtain \(h(m+L)\le2\).



Use \(m\le L\) to get \(2mh\le h(m+L)\le2\), hence \(mh\le1\).



The same threshold gives \(Lh\le2\). Thus \(q=1-mh\ge0\); since \(mh>0\), also \(q<1\).



The sharp small-step contraction of the lifted Euler map gives the contraction hypothesis with \(q=1-mh\). Apply the integral-reference theorem.



Because \(m,h>0\), simplify the bias term:

\[\frac{(33/20)Lh\sqrt{hd}}{1-(1-mh)}=\frac{(33/20)Lh\sqrt{hd}}{mh}=\frac{33}{20}\frac Lm\sqrt{hd}.\]

Substitute the bias identity into the integral-reference bound.



The factor \(q=1-mh\) lies in \([0,1)\) under the positive-curvature small-step assumptions. Its gap is exactly \(mh\), so dividing the local \(h^{3/2}\) error by that gap produces the global \(\sqrt h\) bias. The remaining factor \(L/m\) is the condition-number dependence of this estimate.

\subsection{\texorpdfstring{The large-\allowbreak{}step bound}{The large-step bound}}


Suppose \(U\) is differentiable and \(m\)-strongly convex on \(H\), \(\nabla U\) is \(L\)-Lipschitz, and \(0<m\le L\). Let \(h>0\), \(h\ge2/(m+L)\), and \(Lh<2\), and supply a reference with the stated properties. For any initial vector \(X_0\in\mathcal H\),

\[W_2((X_n)_\#P,\pi)\le(Lh-1)^n\|X_0-Y_{0,0}\|_2+\frac{33}{20}\frac{Lh}{2-Lh}\sqrt{hd}.\]

The denominator is positive by strict stability. The reference remains an input, and the initial term remains the chosen coupling cost.



Since \(m+L>0\), clear the denominator in the lower step threshold to obtain \(2\le h(m+L)\).



Use \(m\le L\) and \(h\ge0\): \(2\le h(m+L)\le2Lh\). Thus \(Lh\ge1\) and \(q=Lh-1\ge0\).



Strict stability \(Lh<2\) gives \(q=Lh-1<1\).



Apply the sharp large-step Euler contraction with \(q=Lh-1\), then invoke the integral-reference theorem.



Use \(1-(Lh-1)=2-Lh\) and factor out \(\sqrt{hd}\):

\[\frac{(33/20)Lh\sqrt{hd}}{1-(Lh-1)}=\left(\frac{(33/20)Lh}{2-Lh}\right)\sqrt{hd}.\]

Substitute this identity into the integral-reference bound to obtain the large-step formula.



Here \(q=Lh-1\) and the gap is \(1-q=2-Lh>0\). Dividing the same local-defect estimate by that gap gives the displayed large-step floor. The denominator records deterioration near the explicit-Euler stability boundary; the local error itself remains bounded, but too little contraction is available to damp its repeated insertion.

\section{\texorpdfstring{Synchronous Langevin coupling and geometric error bounds}{Synchronous Langevin coupling and geometric error bounds}}
\label{proof-module-IPLA-Dalalyan-lean}


The basic coupling estimate compares one Euler step with one stationary Langevin segment under the same noise. Deterministic gradient-step contraction controls the existing error, while a local drift integral controls the new error. Its finite geometric summation is the analytic source of the physical particle transport bound.


\subsection{\texorpdfstring{The shared synchronous coupling calculation}{The shared synchronous coupling calculation}}


The common Euler and diffusion noises cancel in the coupled difference. If \(e_k=\|X_k-Y_k\|_{L^2}\), contraction of the deterministic gradient step and the integrated drift-defect estimate give

\[e_{k+1}\le q e_k+b,\qquad b=\frac{33}{20}Lh\sqrt{hd}.\]


The local forcing is \(b=(33/20)Lh\sqrt{hd}\), so dividing by \(1-q=mh\) gives \((33/20)(L/m)\sqrt{hd}\). This division uses positive curvature and positive step. The factor \(q^k\) keeps the initial coupled discrepancy. To replace that discrepancy by Wasserstein distance, the construction must work uniformly over finite initial couplings and preserve a common algorithm marginal; those requirements are handled separately.

\section{\texorpdfstring{Statistical target, initialization and Gaussian innovations}{Statistical target, initialization and Gaussian innovations}}
\label{proof-module-IPLA-Model-lean}


The statistical center and the sampling reference are different objects: the former maximizes one latent-integrated likelihood, while the latter is a Gibbs probability on the enlarged particle state. Their connection is that the parameter marginal of the enlarged target has density proportional to \(M^N\). Establishing this connection requires the normalization and change-of-variables arguments developed later.


\subsection{\texorpdfstring{The marginal likelihood and its maximizer}{The marginal likelihood and its maximizer}}


The pairing map forms \(\operatorname{pair}(\theta,x)=(\theta,x)\) in the orthogonal joint space. The unnormalized parameter marginal is

\[M(\theta)=\int_{\mathbb R^{d_x}}e^{-U(\theta,x)}\,dx.\]

The definition uses the nonnegative extended integral, so it remains meaningful before finiteness has been established. A marginal maximizer is a point \(\theta_*\) such that \(M(\theta)\le M(\theta_*)\) for every \(\theta\). If \(U=-\log p_\theta(x,y)\), this is maximum marginal likelihood; a prior included in \(U\) gives the corresponding MAP interpretation. The deprecated name is an abbreviation for exactly the same predicate.





A maximizer is defined by the global comparison with every parameter value, so it does not depend on differentiability of the marginal likelihood. Positive finiteness will later permit the equivalent minimization problem for \(\Phi=-\log M\). At this stage, retaining the nonnegative integral avoids assuming those analytic facts in the definition. A subsequent existence theorem supplies an actual maximizer, and marginal strong convexity supplies uniqueness.

\subsection{\texorpdfstring{The actual Gibbs target and all Gaussian assumptions}{The actual Gibbs target and all Gaussian assumptions}}


For \(V_N\) defined above, the target is \(\Pi_N=Z_N^{-1}e^{-V_N}\,dz\), with \(Z_N=\int e^{-V_N}\,dz\). Both normalization and density are initially extended-valued; this definition alone asserts no probability property. The Gaussian innovation structure asks that every coordinate \(\xi_{k,a}\) have law \(\mathcal N(0,1)\), that all coordinates at all times be mutually independent, and that the entire noise path be independent of the initialization. These are explicit distributional hypotheses on the supplied random variables.




The Gaussian structure concerns the entire countable collection \((k,a)\mapsto\xi_{k,a}\), not only coordinate independence at one time. The separate initialization field states independence of \(Z_0\) from the whole innovation path. Together these conditions imply that every transition receives a fresh full state-space standard Gaussian. They are the probabilistic assumptions later used to equate a supplied realization with the canonical Markov laws; moment membership in \(L^2\) alone would not establish either the Gaussian distribution or independence.





The density in the scaled state is \(e^{-\sum_i U(\theta,\sqrt N z_i)}\). It couples the particle blocks through their common parameter, so it is generally not a product probability. Independence belongs to the innovation source, not to this Gibbs target or to the interacting iterates. In particular, all-time coordinate independence ensures that the next innovation is independent of every finite collection of earlier innovations.

\subsection{\texorpdfstring{The law of the L\ensuremath{{}^{2}} recursion}{The law of the L² recursion}}


A square-integrable initial state and square-integrable noise vectors lie in \(L^2(P)\). With \(h=\gamma/N\), define the random sequence through the shared ULA recursion

\[Z_{k+1}=Z_k-h\nabla V_N(Z_k)+\sqrt{2h}\,\xi_k,\qquad \rho_k=(Z_k)_\#P.\]



A globally Lipschitz gradient has linear growth: \(\|\nabla V_N(z)\|\le\|\nabla V_N(0)\|+NL\|z\|\). Consequently square-integrable current states and innovations yield square-integrable next states by the triangle inequality. Induction establishes the moment property for every finite iteration. The law \(\rho_k\) is the pushforward of the actual iterate, so later expectation statements refer to the algorithm's distribution rather than a chosen reference process.

\section{\texorpdfstring{Gibbs normalization and second moments}{Gibbs normalization and second moments}}
\label{proof-module-UnadjustedLangevin-Gibbs-lean}


On a finite-dimensional real Hilbert space, differentiability and \(m\)-strong convexity with \(m>0\) imply that \(\pi(dx)=Z^{-1}e^{-U(x)}dx\) is a probability measure with finite second moment. This establishes existence of the target measure used throughout the coupling proof.



Throughout, \(Z=\int e^{-U(x)}dx\) and \(\pi_U(dx)=Z^{-1}e^{-U(x)}dx\). Under positive strong convexity, a quadratic lower bound on \(U\) dominates the density and its polynomial weights by integrable Gaussian tails. Strict positivity of the exponential proves \(Z>0\), while domination proves \(Z<\infty\). Both facts are required to normalize the measure and to divide weighted moment integrals by \(Z\).

\subsection{\texorpdfstring{The unnormalized Gibbs weight}{The unnormalized Gibbs weight}}


The unnormalized density is \(w_U(x)=e^{-U(x)}\), embedded in \([0,\infty]\) with a nonnegative scalar. This choice lets the partition integral retain the value infinity; positivity of the exponential ensures that embedding does not discard any part of the density.



The exponential is strictly positive for every finite potential value. Thus this weight has no zero set, a fact used later both for positivity of the partition and for equivalence of Gibbs and Lebesgue null sets.

\subsection{\texorpdfstring{The partition function}{The partition function}}


The partition function is \(Z_U=\int e^{-U(x)}\,dx\), formed as a nonnegative extended integral. At this definition no finiteness assumption is made. Later normalization uses separate proofs of \(0<Z_U\) and \(Z_U<\infty\), which are needed to justify cancellation with its inverse.



The partition records the total unnormalized mass before division. Leaving finiteness open at this definition allows the later normalization theorem to state precisely which potential assumptions make a probability measure available.

\subsection{\texorpdfstring{The normalized Gibbs measure}{The normalized Gibbs measure}}


When \(0<Z_U<\infty\), define \(\pi_U(dx)=Z_U^{-1}e^{-U(x)}dx\). The nonnegative density is measurable and integrates to \(Z_U^{-1}Z_U=1\), so it defines a probability. The following strong-convexity estimates establish the required partition bounds before normalization is used.



For a measurable test \(f\), normalization should give \(\int f\,d\pi_U=Z^{-1}\int fe^{-U}dx\). This formula is used only after positivity and finiteness of \(Z\) have been established.

\subsection{\texorpdfstring{Integrability of a radial Gaussian}{Integrability of a radial Gaussian}}


For \(a>0\), \(e^{-a\|x\|^2}\) is Lebesgue integrable in finite dimension. The proof specializes the absolutely convergent Gaussian Fourier integral to real positive precision \(a\) and zero linear term, then takes the real part. The complex Gaussian integrand is the real embedding of the desired exponential, so the final simplification identifies the functions exactly. This also includes dimension zero.



The one-dimensional Gaussian integral factors over an orthonormal coordinate system. Positive precision \(a\) makes each scalar integral finite, so the product integral is finite in every finite dimension.

\subsection{\texorpdfstring{A quadratic lower bound on the potential}{A quadratic lower bound on the potential}}


Apply the strong supporting-plane inequality at zero: \(U(x)\ge U(0)+\langle\nabla U(0),x\rangle+\frac m2\|x\|^2\). Cauchy–Schwarz and \(ab\le a^2/m+(m/4)b^2\), valid for \(m>0\), absorb the linear term. The result is \(U(x)\ge U(0)-\|\nabla U(0)\|^2/m+(m/4)\|x\|^2\). Completing the scalar square justifies the linear-term absorption.



Completing a scalar square bounds \(-\|\nabla U(0)\|\|x\|+(m/2)\|x\|^2\) below by \(-\|\nabla U(0)\|^2/m+(m/4)\|x\|^2\). Retaining a positive quadratic tail is the purpose of this relaxation.

\subsection{\texorpdfstring{Integrability of the Gibbs weight}{Integrability of the Gibbs weight}}


Exponentiating the quadratic lower bound gives \(e^{-U(x)}\le e^{-U(0)+\|\nabla U(0)\|^2/m}e^{-(m/4)\|x\|^2}\). The right side is an integrable Gaussian multiple. Differentiability supplies continuity and measurability of the left side, so integrable domination transfers integrability from that explicit dominating function.



The prefactor is independent of \(x\); all spatial decay lies in the Gaussian factor. Hence its integrability transfers directly to the Gibbs weight by domination, with no assumption that the origin minimizes the potential.

\subsection{\texorpdfstring{The Gaussian weighted second moment}{The Gaussian weighted second moment}}


For \(a>0\), split the Gaussian decay into two equal factors. The scalar function \(s\mapsto s e^{-(a/2)s}\) is bounded on \([0,\infty)\), so \(\|x\|^2e^{-a\|x\|^2}\le C_a e^{-(a/2)\|x\|^2}\). The remaining Gaussian is integrable; the proof checks measurability and nonnegativity before applying domination.



For \(s\ge0\), the maximum of \(s e^{-as/2}\) is finite. Taking \(s=\|x\|^2\) leaves an integrable factor \(e^{-a\|x\|^2/2}\), proving the polynomially weighted Gaussian tail claim.

\subsection{\texorpdfstring{The weighted Gibbs second moment}{The weighted Gibbs second moment}}


Multiply the previous Gibbs Gaussian domination by \(\|x\|^2\). Its right side has the form \(C\|x\|^2e^{-(m/4)\|x\|^2}\), whose integrability was proved separately. The pointwise inequality uses \(\|x\|^2\ge0\), so multiplying it preserves order; this is the second-moment input for the normalized measure.



The squared norm is nonnegative, so multiplying the density domination preserves its direction. The resulting majorant is precisely the weighted Gaussian whose integrability has just been proved.

\subsection{\texorpdfstring{Strict positivity of the partition}{Strict positivity of the partition}}


The continuous function \(e^{-U}\) is strictly positive everywhere. Its extended integral cannot vanish: a zero integral of a nonnegative measurable function would force it to vanish almost everywhere, contradicting positivity on a nonzero-volume open set. This proves \(0<Z_U\) without requiring an upper bound on \(Z_U\).



Continuity and positivity at one point give a neighborhood on which the weight is bounded below by a positive constant. That neighborhood has positive Lebesgue volume, forcing the partition to be strictly positive.

\subsection{\texorpdfstring{Finiteness of the partition}{Finiteness of the partition}}


Strong convexity with \(m>0\) has already given ordinary integrability of \(e^{-U}\). For a nonnegative integrable function its extended integral is finite. Unfolding the partition function identifies that integral exactly with \(Z_U\), giving the guard needed for inverse cancellation.



For a nonnegative measurable function, ordinary integrability is equivalent to finiteness of its nonnegative extended integral. Applying this equivalence to the already dominated weight supplies the finite partition needed for normalization.

\subsection{\texorpdfstring{The ordinary integral formula for the partition}{The ordinary integral formula for the partition}}


Assume the Gibbs weight \(e^{-U}\) is measurable and integrable, as established from the preceding positive strong-convexity hypotheses. Since the weight is nonnegative, its nonnegative extended integral is finite and equals the ordinary Lebesgue integral. Thus the partition has the conventional formula

\[Z_U=\int_H e^{-U(x)}dx<\infty.\]


The equality is justified by the absolute integrability already proved, rather than by assigning a real value to a divergent integral. This same finite scalar can consequently be used in the normalized density and in ordinary weighted-integral formulas. Strict positivity, needed to divide by \(Z_U\), is supplied by the positive-weight argument.

\subsection{\texorpdfstring{Positivity of the ordinary normalizer}{Positivity of the ordinary normalizer}}


Suppose \(U\) is continuous and the preceding argument has established integrability of \(e^{-U}\). At any fixed point \(x_0\), the value \(e^{-U(x_0)}\) is strictly positive. Continuity therefore gives a neighborhood on which the weight is bounded below by a positive constant.



Such a neighborhood has positive Lebesgue volume, so its weighted integral is positive. Together with the finite partition formula this proves

\[0<\int_H e^{-U(x)}dx=Z_U<\infty.\]


Both positivity and finiteness are now available before taking the reciprocal. Hence \(Z_U^{-1}e^{-U(x)}\) is an ordinary normalized probability density, and every later division by the normalizer is legitimate.

\subsection{\texorpdfstring{Normalization of the Gibbs probability}{Normalization of the Gibbs probability}}


The mass of a density measure is its extended density integral. Pull the constant \(Z_U^{-1}\) out to get \(\pi_U(H)=Z_U^{-1}Z_U\). The proofs that \(Z_U\ne0\) and \(Z_U\ne\infty\) justify this product being one. This theorem constructs probability normalization rather than assuming normalization.



Evaluating the weighted measure on the whole space gives \(Z^{-1}Z\). Both factors are finite and \(Z>0\), so this product is one, proving normalization of the actual Gibbs measure.

\subsection{\texorpdfstring{The normalized Gibbs second moment}{The normalized Gibbs second moment}}


Changing measure converts integrability under \(\pi_U\) to integrability of \(Z_U^{-1}e^{-U(x)}\|x\|^2\) under volume. The weighted second moment is already integrable, and the real normalizing factor is a finite scalar. The weighted-density theorem supplies the almost-everywhere measurability and finiteness hypotheses for this conversion.



The weighted squared-norm integral is finite before dividing by the partition. Multiplication by the finite scalar \(Z^{-1}\) therefore gives a finite Gibbs second moment, establishing square-integrability of the state.

\subsection{\texorpdfstring{Square-\allowbreak{}integrability of the Gibbs state}{Square-integrability of the Gibbs state}}


For the identity random variable, membership in \(L^2(\pi_U;H)\) is equivalent to integrability of \(\|x\|^2\). The identity is measurable, and the preceding weighted-moment theorem supplies the integral. The construction uses the \(p=2\) characterization of square-integrability, retaining the actual function before passing to any quotient class.



Membership of the identity map in \(L^2(\pi_U)\) is exactly the finite second-moment statement. The identity is continuous and therefore measurable, supplying the other part of the membership criterion.

\subsection{\texorpdfstring{Square-\allowbreak{}integrability after translation}{Square-integrability after translation}}


Since \(\pi_U\) is a probability, the constant map \(x\mapsto z\) belongs to \(L^2\). Subtract it from the identity, already in \(L^2(\pi_U)\), to obtain \(x\mapsto x-z\in L^2(\pi_U)\). This is a moment assertion for every chosen center, without requiring that the center minimize the potential.



Use \(\|x-a\|^2\le2\|x\|^2+2\|a\|^2\). The first term has finite Gibbs integral and the second is a constant integrated against a probability measure, so every fixed translation is square-integrable.

\section{\texorpdfstring{Marginal partitions, existence and the exact parameter target}{Marginal partitions, existence and the exact parameter target}}
\label{proof-module-IPLA-MarginalConcentration-lean}


The main issue is to identify the parameter law of the normalized enlarged Gibbs distribution, with every volume constant accounted for. Uniform Gaussian domination makes the statistical marginal a positive continuous function decaying at infinity. Product integration and finite positive normalization then show that its \(N\)-th power is exactly the parameter density of the particle target.


\subsection{\texorpdfstring{Pair coordinates split the squared norm}{Pair coordinates split the squared norm}}


The pairing map has parameter coordinates \(\theta_j\) and latent coordinates \(x_j\). Expanding the Euclidean squared norm over the disjoint index sets gives

\[\|\operatorname{pair}(\theta,x)\|^2=\|\theta\|^2+\|x\|^2.\]

Every output coordinate of pairing is a continuous product projection; the finite coordinate-to-Euclidean identification is continuous as well. Hence \((\theta,x)\mapsto\operatorname{pair}(\theta,x)\) is jointly continuous. These elementary identities justify the orthogonal Gaussian domination and later dominated convergence.





The norm decomposition is orthogonal, rather than an inequality from the triangle inequality. Thus exponentials of a negative quadratic split into a product of a parameter Gaussian and a latent Gaussian. Joint continuity of the pairing map also means that any continuous joint potential has continuous slices in either block. These two facts are the geometric inputs to the ensuing domination and parameter-limit arguments.

\subsection{\texorpdfstring{Strong convexity supplies uniform Gaussian domination}{Strong convexity supplies uniform Gaussian domination}}


The shared strong-convexity lower bound, obtained by bounding the linear term at zero, is

\[U(v)\ge U(0)-\frac{\|\nabla U(0)\|^2}{m}+\frac m4\|v\|^2.\]

Set \(C=\exp(\|\nabla U(0)\|^2/m-U(0))\) and \(a=m/4>0\). Exponentiation and the norm splitting yield

\[e^{-U(\theta,x)}\le Ce^{-a\|\theta\|^2}e^{-a\|x\|^2}\le Ce^{-a\|x\|^2}.\]

The final majorant is independent of \(\theta\) and integrable over the latent Euclidean space. Continuity supplies measurability, so domination proves integrability of every latent slice. No bounded parameter neighborhood is needed.




Strong convexity and the affine score term at zero give a Gaussian majorant after completing a square. Split the joint squared norm as \(\|\operatorname{pair}(\theta,x)\|^2=\|\theta\|^2+\|x\|^2\). Dropping the nonnegative parameter contribution yields an integrable Gaussian bound in \(x\), uniform in \(\theta\) when needed for local parameter limits. This proves latent integrability before converting the extended marginal density to its real integral. Strict positivity follows because the exponential weight is positive everywhere.





Start from \(U(v)\ge U(0)+\langle\nabla U(0),v\rangle+(m/2)\|v\|^2\). Young's inequality gives \(\langle\nabla U(0),v\rangle\ge-\|\nabla U(0)\|^2/m-(m/4)\|v\|^2\). Combining the bounds leaves a strictly positive quadratic coefficient. The resulting Gaussian constant is finite, and its latent part is integrable in every finite dimension, including dimension zero.

\subsection{\texorpdfstring{The marginal is positive, finite and continuous}{The marginal is positive, finite and continuous}}


Let \(M(\theta)=\int e^{-U(\theta,x)}dx\). Slice integrability and nonnegativity identify the extended marginal with \(M(\theta)\). This excludes infinity. The integrand is strictly positive everywhere, and Lebesgue measure is nonzero, so \(M(\theta)>0\) and the extended marginal is strictly positive as well. If \(\theta_n\to\theta\), the integrands converge pointwise by continuity of \(U\), and the preceding common Gaussian majorant permits dominated convergence. Thus \(M\) is continuous as an ordinary positive real-valued marginal function.





Strict positivity of \(M(\theta)\) is not obtained by a lower Gaussian bound; it follows because a positive continuous density has positive integral on some ball of positive volume. Finiteness follows from the upper Gaussian bound. For continuity, every convergent parameter sequence can be passed through the latent integral using the same integrable majorant. Positivity and continuity will later justify taking a real logarithm without imposing extra hypotheses.

\subsection{\texorpdfstring{Gaussian parameter tails force existence of a maximizer}{Gaussian parameter tails force existence of a maximizer}}


Integrating the stronger Gaussian domination gives

\[0<M(\theta)\le C\left(\int e^{-a\|x\|^2}dx\right)e^{-a\|\theta\|^2}.\]

The finite constant is independent of \(\theta\), and the right-hand side tends to zero as \(\theta\) leaves compact sets. The squeeze theorem proves \(M(\theta)\to0\) at infinity. Since \(M(0)>0\), outside a sufficiently large compact region the values are below \(M(0)\). Continuity then attains a global maximum within a compact region. The same point therefore maximizes the original marginal likelihood. Uniqueness is proved later from Prékopa strong convexity.




Integrating the stronger Gaussian bound while retaining the parameter term gives \(M(\theta)\le C e^{-c\|\theta\|^2}\), for finite \(C\) and \(c>0\). Dominated convergence proves continuity of \(M\), and this bound proves \(M(\theta)\to0\) as \(\|\theta\|\to\infty\). Since \(M\) is positive at any chosen point, outside a sufficiently large ball it is smaller than that value. Compact attainment on the ball proves existence of the actual global marginal maximizer.





Choose \(R\) so that the Gaussian upper bound is below \(M(0)/2\) whenever \(\|\theta\|>R\). The continuous marginal attains a maximum on the closed ball of radius \(R\), and that maximum is at least \(M(0)\). No point outside the ball can improve it. This proves a global maximizer using only compact attainment and a quantitative tail estimate, without differentiating the marginal.

\subsection{\texorpdfstring{Product integration and the particle-\allowbreak{}scaling Jacobian}{Product integration and the particle-scaling Jacobian}}


Tonelli induction over \(n\) proves the extended-integral identity

\[\int\prod_{i=1}^nf(x_i)\,d\mu^{\otimes n}=\left(\int f\,d\mu\right)^n.\]

The base case is the mass-one empty product; the successor case separates one coordinate and multiplies the remaining integral. This statement records divergence correctly because no real conversion is used. For a nonzero real dilation \(r\) in dimension \(d\), Haar measure scaling gives \(\int f(rx)dx=|r^d|^{-1}\int f(x)dx\). Applying these two identities at \(r=\sqrt N\), with \(N>0\), yields

\[\int\prod_i e^{-U(\theta,\sqrt N z_i)}\,dz_1\cdots dz_N=\left(N^{-d_x/2}M(\theta)\right)^N.\]

The formal statement retains the absolute-value inverse-power Jacobian, which is equal to the displayed factor.




For a fixed \(\theta\), the lifted weight factors as \(\prod_i e^{-U(\theta,\sqrt N z_i)}\). Tonelli factors the nonnegative product integral into a product of identical latent integrals. The substitution \(x_i=\sqrt N z_i\) contributes \(N^{-d_x/2}\) per particle, so the fiber integral is \(N^{-Nd_x/2}M(\theta)^N\). This Jacobian is finite and strictly positive when \(N>0\), and does not depend on \(\theta\).





The product integral concerns nonnegative functions, so Tonelli applies before integrability has been established. Under the positive finite slice result, each factor is then an ordinary finite number. Dilation by \(\sqrt N\) in one latent block has determinant \(N^{d_x/2}\); the inverse determinant enters the integral. Repeating this for \(N\) blocks yields one common factor independent of the parameter, which later disappears on normalization.

\subsection{\texorpdfstring{A coordinate equivalence groups the state into \ensuremath{\theta} and particle blocks}{A coordinate equivalence groups the state into θ and particle blocks}}


The configuration map \(C(\theta,(z_i)_i)\) inserts the parameter coordinates first and then all particle coordinates. Its coordinate evaluation theorems follow directly from insertion. The inverse is \(z\mapsto(p_\theta z,(\operatorname{block}_iz)_i)\); checking each coordinate proves both inverse identities and linearity, hence a finite-dimensional continuous linear equivalence. The evaluation theorem identifies this equivalence with the configuration definition, while \(p_\theta C(\theta,z)=\theta\) and \(A_iC(\theta,z)=\operatorname{pair}(\theta,\sqrt N z_i)\). These identities give the exact arguments of the fiber integrals.





Grouping coordinates is a bijection, not an approximation or a restriction of the state space. The inverse recovers every original coordinate exactly, and both directions are linear. The potential evaluated in grouped coordinates therefore uses precisely the original parameter and physical particle values. This allows Fubini and density projection to be applied on a product space while still describing the same enlarged Euclidean Gibbs measure.

\subsection{\texorpdfstring{The weight factors and its parameter projection integrates each fiber}{The weight factors and its parameter projection integrates each fiber}}


In grouped coordinates, exponentiation converts the sum in \(V_N\) into a product:

\[e^{-V_N(C(\theta,z))}=\prod_i e^{-U(\theta,\sqrt N z_i)}.\]

Two measure identities justify using this pointwise formula. First, for a measurable equivalence \(e\), pushing \(\mu\) with weight \(w\circ e\) through \(e\) gives \((e_\#\mu)\) with weight \(w\); prove this by evaluating any measurable set and the pushforward integral formula. Second, projecting \(w(\theta,z)d\mu(\theta)d\nu(z)\) to its first coordinate gives density \(\int w(\theta,z)d\nu(z)\) relative to \(\mu\), by Tonelli on the indicator of an arbitrary measurable parameter set.





For a measurable parameter set \(A\), integrate \(\mathbf 1_A(\theta)e^{-V_N(\theta,z)}\) over both blocks. Tonelli first integrates the latent variables and produces a function of \(\theta\); integrating this function over \(A\) is exactly the projected unnormalized measure. Since equality holds for every measurable \(A\), it identifies the measure itself. The exponential product identity supplies its fiber density explicitly.

\subsection{\texorpdfstring{All finite positive volume constants cancel on normalization}{All finite positive volume constants cancel on normalization}}


A continuous linear equivalence carries Euclidean volume to a Haar measure, hence to \(c\) times the target volume for a finite positive \(c\). This proves the required coordinate-volume relation without assuming its scalar is one. For any measure \(\mu\) and \(0<c<\infty\), extended-real inverse multiplication gives

\[\frac{c\mu}{(c\mu)(E)}=\frac{\mu}{\mu(E)}.\]

Thus both the coordinate-grouping scalar and the particle-dilation Jacobian cancel in a normalized law. The formal normalization lemma proves this algebra for arbitrary \(\mu\), keeping the nonzero and finite conditions on \(c\) explicit.





Both changes of variables may introduce constants, but neither constant depends on \(\theta\). If an unnormalized measure is replaced by \(c\mu\), its total mass is replaced by \(c\mu(E)\), so the quotient by total mass is unchanged. Finiteness and positivity of \(c\) justify canceling it. Tracking these conditions avoids treating a zero or infinite scalar as an innocuous normalization factor.

\subsection{\texorpdfstring{The actual parameter marginal is the powered likelihood target}{The actual parameter marginal is the powered likelihood target}}


Define the normalized power target \(\pi_N^\theta(d\theta)=Z_\theta^{-1}M(\theta)^N d\theta\). To identify it, let \(\mu=e^{-V_N}dz\) and \(\rho=M(\theta)^Nd\theta\). Grouping coordinates changes volume by \(c\), and the fiber calculation contributes \(J^N\), \(J=N^{-d_x/2}\). The preceding pushforward identities prove \(c\,p_{\theta\#}\mu=J^N\rho\). Both scalars are finite and positive. Normalize this equality, cancel both scalars, and use preservation of total mass under measurable projection to conclude

\[p_{\theta\#}\Pi_N=\pi_N^\theta.\]

Every change of variables and normalization factor in this identity is proved; no formula for the marginal is assumed.




The configuration equivalence groups the orthogonal state coordinates into the parameter block and all latent blocks. The volume theorem tracks the finite-dimensional coordinate-change constant, while the density pushforward lemmas integrate each parameter fiber. The preceding factorization leaves a density proportional to \(M(\theta)^N\), with the scaling and coordinate-change factors independent of \(\theta\). Normalization cancels those positive finite factors exactly, yielding \[(p_\theta)_\#\Pi_N(d\theta)=\frac{M(\theta)^N\,d\theta}{\int M(u)^Ndu}.\] The final target probability and moment lemmas are consequences of this identity and the already proved lifted target facts.





Let \(J=N^{-Nd_x/2}\) denote the total latent dilation factor. Up to the fixed coordinate-volume constant, the projected unnormalized density is \(J M(\theta)^N\). Its integral over parameter space has the same prefactor. Dividing these two expressions proves the normalized density formula. In particular the parameter mode is a maximizer of \(M\), because raising a positive number to the positive integer \(N\) preserves order.

\subsection{\texorpdfstring{Normalize the power target and guard its moments}{Normalize the power target and guard its moments}}


The established projection identity and the probability property of \(\Pi_N\) prove that \(\pi_N^\theta\) is a probability measure. Its unnormalized density \(M^N\) is everywhere positive, so its extended partition is strictly positive. If that partition were infinite, the definition of the normalized target would be the zero measure, contradicting mass one; thus it is also finite. Finally the lifted Gibbs law has a finite second moment and \(p_\theta\) is a bounded linear map. The pushforward integral identity therefore gives \(\theta\mapsto\theta\in L^2(\pi_N^\theta)\). This supplies finite marginal moments before the later sharp bound is applied.





The lifted Gibbs law already has mass one, so its measurable projection does too. A hypothetical infinite marginal partition would make its normalized density zero and hence give total mass zero, a contradiction. Finally, \(\|p_\theta z\|^2\le\|z\|^2\) transfers the lifted second moment through the projection. This gives the finite marginal moments needed before a quantitative concentration radius is introduced.

\section{\texorpdfstring{One-\allowbreak{}dimensional Prékopa–Leindler by constructed monotone transport}{One-dimensional Prékopa–Leindler by constructed monotone transport}}
\label{proof-module-IPLA-PrekopaLeindler-lean}


For positive continuous densities in one dimension, cumulative distribution functions provide an explicit monotone change of variables. The proof uses the derivative of this change of variables and the weighted arithmetic–geometric mean inequality. It then removes normalization to obtain an inequality between arbitrary positive finite integrals, ready for the marginal-convexity argument.




The proof constructs transport directly from cumulative integrals of positive continuous densities. It derives the one-dimensional Prékopa–Leindler inequality without postulating a transport map or marginal log-concavity theorem.

\subsection{\texorpdfstring{Cumulative integrals are differentiable and strictly increasing}{Cumulative integrals are differentiable and strictly increasing}}


A density here is a continuous strictly positive real function \(f\), integrable on \(\mathbb R\), with \(\int f=1\). Define \(F(x)=\int_{(-\infty,x]}f\). The difference formula

\[F(y)-F(0)=\int_0^y f(t)dt\]

and the fundamental theorem of calculus give the strict derivative \(F'(x)=f(x)\). Differentiability implies continuity, and positivity of the derivative implies strict monotonicity. The density assumptions are precisely the inputs required for these assertions.





For \(y>x\), the increment \(F(y)-F(x)=\int_x^y f(t)dt\) is positive because the density is positive and continuous on a nonempty interval. Thus strict monotonicity also has a direct integral interpretation. The derivative identity follows locally from the fundamental theorem of calculus; the improper lower limit contributes only a constant and does not affect differentiation.

\subsection{\texorpdfstring{The cumulative range is exactly the open unit interval}{The cumulative range is exactly the open unit interval}}


Integrability implies \(F(x)\to0\) as \(x\to-\infty\). Splitting the total integral at \(x\) gives \(F(x)=1-\int_{(x,\infty)}f\), so \(F(x)\to1\) as \(x\to+\infty\). Nonnegativity gives \(0\le F\le1\); strict monotonicity, comparing \(x\) with \(x-1\) and \(x+1\), improves this to \(0<F(x)<1\). For every \(u\in(0,1)\), the endpoint limits supply points on either side of \(u\), and the intermediate value theorem supplies \(F(x)=u\). Therefore \(F(\mathbb R)=(0,1)\).




The endpoint limits identify the inverse's correct domain: neither zero nor one is attained at a finite point. For any interior probability level, choose sufficiently far left and right points using those limits, then use continuity to obtain a preimage. Strict monotonicity makes this preimage unique. Consequently the inverse cumulative map used in transport is well defined everywhere it is needed.

\subsection{\texorpdfstring{Construct and differentiate the monotone transport}{Construct and differentiate the monotone transport}}


For two densities with cumulative maps \(F,G\), set \(T=G^{-1}\circ F\), using the inverse function on its range. Since both ranges are \((0,1)\), \(G(T(x))=F(x)\) everywhere. Strict monotonicity makes \(G\) injective. The one-dimensional inverse function theorem applies because \(G'(T(x))=g(T(x))>0\), and the chain rule yields

\[T'(x)=\frac{f(x)}{g(T(x))}>0.\]

Thus \(T\) is strictly increasing, and the density ratio is a proved derivative of an explicitly constructed function.






The identity \(G(T(x))=F(x)\) matches the mass lying to the left of corresponding points. Differentiating gives \(g(T(x))T'(x)=f(x)\); positivity of the denominator allows division and shows that the transport preserves orientation. This also proves the density Jacobian relation needed in the integral calculation, without assuming a transport existence theorem as an extra analytic premise.

\subsection{\texorpdfstring{The normalized inequality is a change-\allowbreak{}of-\allowbreak{}variables calculation}{The normalized inequality is a change-of-variables calculation}}


Assume \(0\le t\le1\) and \(h((1-t)x+ty)\ge f(x)^{1-t}g(y)^t\). Put \(S(x)=(1-t)x+tT(x)\) and \(r=f(x)/g(T(x))>0\). Then \(S'=1-t+tr>0\), and convexity of the exponential gives \(r^t\le1-t+tr=S'\). At \(y=T(x)\), the pointwise premise therefore implies

\[f(x)=f(x)^{1-t}g(T(x))^t r^t\le h(S(x))S'(x).\]

Since \(S\) is strictly increasing, the one-dimensional Jacobian theorem changes the right-hand integral to the integral of \(h\) over \(S(\mathbb R)\). Nonnegativity then bounds it by the integral over all \(\mathbb R\), proving \(1\le\int h\). The formal conclusion uses an extended nonnegative integral, so \(h\) need not itself be integrable or continuous.




Set \(S(x)=(1-t)x+tT(x)\). The interpolation hypothesis bounds \(h(S(x))\) below by \(f(x)^{1-t}g(T(x))^t\). Weighted arithmetic–geometric mean gives \(S'(x)=(1-t)+tT'(x)\ge T'(x)^t\). Since \(T'=f/g(T)\), multiplication yields \(h(S(x))S'(x)\ge f(x)\). The strictly increasing change-of-variables formula then proves \(\int h\ge\int f=1\). For arbitrary positive finite masses, divide \(f,g\) by their integrals, apply the normalized statement, and multiply back; logarithmic mass terms recover \((\int f)^{1-t}(\int g)^t\).





The scalar inequality \((1-t)+tr\ge r^t\) holds for every \(r>0\), including \(t=0,1\). Substitution of \(r=f/g(T)\) converts the interpolation lower bound into an exact density comparison after multiplication by the Jacobian. Integrating over the image of the increasing map is enough: nonnegativity of \(h\) bounds that image integral by the integral on the entire line. Surjectivity of the interpolating map is therefore unnecessary.

\subsection{\texorpdfstring{Remove the unit-\allowbreak{}integral normalization}{Remove the unit-integral normalization}}


An everywhere positive integrable function on a nonzero measure space has strictly positive integral, since its support has positive measure. For the Prékopa weights write \(A=\int f>0\), \(B=\int g>0\), and \(C=A^{1-t}B^t>0\). Apply the normalized inequality to \(f/A\), \(g/B\), and \(h/C\). Expanding \(\log(f/A)=\log f-\log A\) verifies its pointwise premise. If \(h\) is nonnegative and integrable, the extended conclusion converts to \(1\le(\int h)/C\). Multiplication by \(C\) gives

\[\exp\!\left((1-t)\log\int f+t\log\int g\right)\le\int h.\]

The positivity lemma also applies to the zero-dimensional product measure needed later.





Scaling \(f\) and \(g\) to unit integral changes the interpolation factor by \(A^{1-t}B^t\). The normalized result gives a lower bound of one for the integral of the correspondingly scaled third function. Multiplying by this strictly positive finite factor restores the original masses. Writing powers as exponentials of logarithms makes the endpoint weights unambiguous and uses only logarithms of positive numbers.

\section{\texorpdfstring{Finite-\allowbreak{}dimensional Prékopa–Leindler with Gaussian tails}{Finite-dimensional Prékopa–Leindler with Gaussian tails}}
\label{proof-module-IPLA-PrekopaTensorization-lean}


The dimension induction requires a class of functions stable under marginalization. Continuous positive functions with Gaussian upper tails have exactly the needed properties: every coordinate slice is integrable, its integral varies continuously, and the resulting marginal again has Gaussian tails. This closes the induction without assuming marginal continuity or integrability independently.




Positive continuous weights dominated by an integrable Gaussian form a class closed under integrating out one coordinate. This supplies all continuity and integrability facts needed to iterate the one-dimensional transport proof.

\subsection{\texorpdfstring{Gaussian domination proves integrability and positive integral}{Gaussian domination proves integrability and positive integral}}


Write \(\mathrm{Coordinates}(n)=\mathbb R^n\). A Gaussian weight satisfies continuity, strict positivity, and \(f(x)\le Ce^{-a\sum_i x_i^2}\), with \(C,a>0\). The Euclidean Gaussian integral is finite, and the coordinate identification preserves volume, so domination proves \(f\in L^1\). Positivity on the whole nonzero volume space proves \(\int f>0\). Multiplying by \(c>0\) preserves these properties and changes the envelope constant to \(cC\).





A Gaussian envelope with coefficient \(a>0\) has finite integral on \(\mathbb R^n\), as follows by the product of one-dimensional Gaussian integrals. Since the function is nonnegative, domination proves its absolute integrability. Strict positivity then makes its integral positive. Multiplication by a positive scalar preserves all three properties and changes only the envelope constant, a fact needed when strong convexity supplies an exponential prefactor.

\subsection{\texorpdfstring{A coordinate slice and its marginal are integrable and continuous}{A coordinate slice and its marginal are integrable and continuous}}


The insertion map \((s,x)\mapsto(s,x_1,\ldots,x_n)\) is continuous coordinate by coordinate. Define \(f_1(x)=\int_\mathbb R f(s,x)ds\). Since \(\sum_i(s,x)_i^2=s^2+\sum_i x_i^2\ge s^2\), every slice is bounded by \(Ce^{-as^2}\), which is integrable independently of \(x\). Slice continuity supplies measurability and domination proves slice integrability. If \(x_j\to x\), joint continuity gives pointwise convergence in \(s\), and the same Gaussian bound allows dominated convergence. Hence \(f_1\) is continuous.




Write a point in \(\mathbb R^{n+1}\) as \(\operatorname{cons}(z,x)\), where \(z\in\mathbb R\) is the coordinate being integrated. The coordinate-insertion continuity theorem checks its first and remaining coordinates separately. A Gaussian weight has constants \(C,a>0\) satisfying

\[0<f(\operatorname{cons}(z,x))\le C\exp\!\left[-a\left(z^2+\sum_{i=1}^{n}x_i^2\right)\right]\le C e^{-az^2}.\]



The crucial majorant is independent of the remaining coordinates. Therefore a sequence of such coordinates changes only the integrand, not the integrable bound used in dominated convergence. Each slice is positive and continuous in the integrated coordinate, so its integral is positive and finite. The argument proves joint analytic control needed to regard the first marginal as a function in the same class.

\subsection{\texorpdfstring{Coordinate integration preserves Gaussian tails and total integral}{Coordinate integration preserves Gaussian tails and total integral}}


Retaining the full Gaussian envelope and factoring its exponent gives

\[f_1(x)\le C\left(\int_\mathbb R e^{-as^2}ds\right)e^{-a\sum_i x_i^2}.\]

The new constant is positive and finite, and integrating the positive slice makes \(f_1\) positive everywhere. Together with continuity, this proves that \(f_1\) is again a Gaussian weight. The coordinate splitting equivalence preserves product volume. Its integral formula and Fubini, justified by the Gaussian integrability already proved, then show \(\int f_1=\int f\).




Retaining the remaining-coordinate Gaussian factor gives a sharper conclusion than slice integrability alone. The new envelope constant is the old constant multiplied by \(\int e^{-as^2}ds\), while the decay coefficient \(a\) is unchanged. Fubini then identifies the marginal's total integral with the original integral because Gaussian domination has already established absolute integrability. These facts make repeated coordinate integration legitimate.

\subsection{\texorpdfstring{Induct on the dimension}{Induct on the dimension}}


For \(n=0\), the coordinate space is a singleton with volume one; each integral equals its function value and the result is the pointwise premise. For dimension \(n+1\), fix remaining-coordinate vectors \(x,y\). The one-dimensional theorem applied to the slices \(s\mapsto f(s,x)\), \(r\mapsto g(r,y)\), and \(u\mapsto h(u,(1-t)x+ty)\) gives the pointwise Prékopa inequality for \(f_1,g_1,h_1\). These marginals are Gaussian weights, so the induction hypothesis applies in dimension \(n\). Replace their integrals by those of the original weights to obtain

\[\left(\int f\right)^{1-t}\left(\int g\right)^t\le\int h.\]

The formal statement writes positive powers using exponentials and logarithms, including \(t=0\) and \(t=1\).




The induction splits \(\mathbb R^{d+1}\) into its first coordinate and \(d\) remaining coordinates. For fixed remaining coordinates, apply the proved one-dimensional inequality to the first-coordinate slices. Their integrals define marginal functions that satisfy the same interpolation inequality in \(d\) dimensions. Gaussian domination proves slice integrability, continuity of those marginal functions by dominated convergence, and Gaussian tails for them. The induction hypothesis applies, and Fubini identifies the iterated integrals with the original total integrals. Dimension zero is the one-point base case.





At the induction step the original pointwise hypothesis holds for every pair of first coordinates and every pair of remaining-coordinate vectors. Fixing the latter pair gives exactly the one-dimensional interpolation premise. Its conclusion is the corresponding premise for the marginal functions. Applying the induction hypothesis integrates the remaining coordinates, and the total-integral identities replace each iterated integral by its original multidimensional value.

\section{\texorpdfstring{Derived strong convexity of the actual marginal potential}{Derived strong convexity of the actual marginal potential}}
\label{proof-module-IPLA-StrongLogConcavity-lean}


Joint strong convexity is preserved when latent variables are integrated out, provided the parameter quadratic deficit is retained in the interpolation argument. The positive finite marginal integrals permit the Prékopa inequality to be converted into the strong-convexity inequality for \(-\log M\). This establishes both the correct concentration potential and uniqueness of the statistical maximizer.




Let \(M(\theta)=\int e^{-U(\theta,x)}dx\) and \(\Phi(\theta)=-\log M(\theta)\). The proved finite-dimensional Prékopa–Leindler theorem shows that \(\Phi\) retains the original joint strong-convexity constant \(m\).

\subsection{\texorpdfstring{Define a legitimate logarithmic marginal potential}{Define a legitimate logarithmic marginal potential}}


The real marginal integral \(M\) is defined from the original slice; the earlier domination theorem makes it finite and strictly positive. Thus \(\Phi=-\log M\) is an ordinary real-valued continuous function: continuity follows by composing the continuous marginal with logarithm away from zero. If \(\theta_*\) maximizes the extended marginal, its finite real representation and monotonicity of logarithm imply \(\Phi(\theta_*)\le\Phi(\theta)\) for every \(\theta\). No logarithm of zero or infinite partition is used.





The marginal potential is defined on all parameter values because \(0<M(\theta)<\infty\) everywhere. If \(\theta_*\) maximizes \(M\), monotonicity of logarithm gives \(\log M(\theta)\le\log M(\theta_*)\); negation gives global minimality of \(\Phi\). This equivalence concerns global extrema and requires no interchange of differentiation with integration.

\subsection{\texorpdfstring{Prepare Gaussian weights in product coordinates}{Prepare Gaussian weights in product coordinates}}


Define \(f_\theta(x)=e^{-U(\theta,x)}\) in ordinary finite-product coordinates, using the volume-preserving Euclidean identification for the latent block. The joint Gaussian domination gives \(f_\theta(x)\le Ce^{-m\|x\|^2/4}\), and continuity plus strict positivity make it a Gaussian weight. Its integral equals \(M(\theta)\) by the volume-preserving change of coordinates. Finally pairing preserves subtraction and affine combinations coordinatewise:

\[\operatorname{pair}(\theta,x)-\operatorname{pair}(\theta',y)=\operatorname{pair}(\theta-\theta',x-y),\]\[\operatorname{pair}(a\theta+b\theta',ax+by)=a\operatorname{pair}(\theta,x)+b\operatorname{pair}(\theta',y).\]



The latent coordinate identification is volume preserving, so it changes neither the marginal integral nor its dimension. The affine pairing identity allows joint strong convexity to be applied at \((\theta,x)\) and \((\theta',y)\). Its squared difference separates into the parameter and latent squares. This separation is what permits the latent penalty to be dropped while retaining the exact parameter modulus.

\subsection{\texorpdfstring{Apply Prékopa with the parameter convexity penalty}{Apply Prékopa with the parameter convexity penalty}}


Let \(a,b\ge0\), \(a+b=1\), \(\theta_r=a\theta+b\theta'\), and \(c=(mab/2)\|\theta-\theta'\|^2\). Strong convexity of \(U\), the affine pair identity and norm splitting imply

\[aU(\theta,x)+bU(\theta',y)\ge U(\theta_r,ax+by)+c,\]

after discarding the nonnegative latent penalty \((mab/2)\|x-y\|^2\). Exponentiate and apply Prékopa–Leindler with third weight \(e^{-c}f_{\theta_r}\), which remains a Gaussian weight. This gives \(M(\theta)^aM(\theta')^b\le e^{-c}M(\theta_r)\). All factors are positive, so take logarithms and negate to conclude

\[\Phi(a\theta+b\theta')\le a\Phi(\theta)+b\Phi(\theta')-\frac{mab}{2}\|\theta-\theta'\|^2.\]

Thus marginal strong convexity is a theorem of the joint assumptions.




For \(\theta_r=a\theta+b\vartheta\), \(a+b=1\), joint strong convexity yields \[U(\theta_r,ax+by)\le aU(\theta,x)+bU(\vartheta,y)-\tfrac{mab}{2}\|\theta-\vartheta\|^2,\] after dropping the additional nonnegative latent deficit. Thus the latent weights satisfy the Prékopa hypothesis with the multiplier \(e^{-c}\), \(c=mab\|\theta-\vartheta\|^2/2\), on the interpolated weight. Integrating gives \(M(\theta)^aM(\vartheta)^b\le e^{-c}M(\theta_r)\).





Let \(c=(mab/2)\|\theta-\theta'\|^2\). The joint inequality implies \(e^{-c}e^{-U(\theta_r,ax+by)}\ge e^{-aU(\theta,x)}e^{-bU(\theta',y)}\). Prékopa therefore gives \(e^{-c}M(\theta_r)\ge M(\theta)^aM(\theta')^b\). Taking logarithms yields \(\log M(\theta_r)\ge a\log M(\theta)+b\log M(\theta')+c\). Negating this is precisely strong convexity with modulus \(m\), including the quadratic factor \(1/2\).

\subsection{\texorpdfstring{The marginal maximizer exists uniquely}{The marginal maximizer exists uniquely}}


Because \(m>0\), the proved \(m\)-strong convexity makes \(\Phi\) strictly convex. Two marginal maximizers are both global minimizers of \(\Phi\); if distinct, the midpoint strict-convexity inequality would lie strictly below their common minimum, a contradiction. Hence they coincide. The earlier continuity-and-tail argument already supplied existence of a maximizer. Combining existence with this uniqueness yields \(\exists!\theta_*\) satisfying the original statistical maximizer predicate.





If two distinct minimizers existed, their midpoint would have value at most their common minimum minus \((m/8)\|\theta-\theta'\|^2\), strictly smaller because \(m>0\). This contradiction proves uniqueness. Existence came earlier from continuity and Gaussian decay, so the two arguments serve different purposes: coercive tails provide attainment, while the marginal interpolation inequality prevents multiple maximizing parameters.

\section{\texorpdfstring{Concentration around a minimizer without differentiability}{Concentration around a minimizer without differentiability}}
\label{proof-module-UnadjustedLangevin-GibbsConcentration-lean}


For a continuous \(q\)-strongly convex potential \(F\), \(q>0\), and an explicitly supplied minimizer \(z\), this module proves the equilibrium radius \(\int\|x-z\|^2d\pi_F\le2d/q\). The proof uses dilation and convexity, so it also applies to the nonsmooth potentials used by proximal algorithms.



Let \(x_\star\) be the stated mode and write the radial variable as \(r=\|x-x_\star\|\). Strong convexity gives \(U(x)-U(x_\star)\ge mr^2/2\). The comparison under a radial contraction retains the change-of-variables factor from \(d\)-dimensional volume; differentiating the resulting partition comparison controls the weighted squared radius. Dividing by the positive partition gives \(\int\|x-x_\star\|^2d\pi_U\le d/m\).

\subsection{\texorpdfstring{The quadratic potential bound at a mode}{The quadratic potential bound at a mode}}


Use strong convexity at the midpoint of \(z\) and \(x\), then use minimality \(F(z)\le F((z+x)/2)\). Rearranging yields \(F(x)\ge F(z)+(q/4)\|x-z\|^2\). This argument uses values of \(F\), not a derivative or a first-order optimality assertion at the minimizer.



At the mode the linear term vanishes. The strong lower model is therefore centered exactly at \(x_\star\), with no linear-term absorption or loss from \(m/2\) to \(m/4\).

\subsection{\texorpdfstring{Integrability around the mode}{Integrability around the mode}}


For \(q>0\), the midpoint bound gives \(e^{-F(x)}\le e^{-F(z)}e^{-(q/4)\|x-z\|^2}\). Translation invariance reduces integrability of the majorant to the centered Gaussian integral. Continuity makes the density measurable, so domination proves integrability without differentiability of \(F\).



Translate coordinates by \(x_\star\) and dominate the density by \(e^{-U(x_\star)}e^{-m\|x-x_\star\|^2/2}\). Lebesgue translation preserves the integral and the remaining Gaussian is integrable.

\subsection{\texorpdfstring{The weighted squared-\allowbreak{}distance moment}{The weighted squared-distance moment}}


Multiply the translated Gaussian bound by \(\|x-z\|^2\). Translation preserves Lebesgue measure, and a quadratic times a decaying Gaussian is integrable. Translation transports this integrability through \(x\mapsto x-z\) and then applies the nonnegative domination theorem to the desired weighted moment.



Multiplying that centered Gaussian majorant by squared distance keeps it integrable. This proves the numerator of the centered Gibbs second moment finite before the normalization comparison is used.

\subsection{\texorpdfstring{Contracting coordinates toward the mode}{Contracting coordinates toward the mode}}


For \(0<r<1\), compare \(F(z+r(x-z))\) with \(F(x)\) using strong convexity and \(F(z)\le F(x)\). Exponentiating and using \(e^u\ge1+u\) produces a linear correction in \(\|x-z\|^2\). The dilation \(x\mapsto z+r(x-z)\) multiplies volume by \(r^d\), hence \(Z+(q/2)r(1-r)M\le r^{-d}Z\), with \(M=\int\|x-z\|^2e^{-F(x)}dx\).



The radial map \(x\mapsto x_\star+t(x-x_\star)\) scales volume by \(t^d\). Strong convexity controls the potential change along this map, so the partition comparison retains both curvature and dimension.

\subsection{\texorpdfstring{A partition comparison for squared distance}{A partition comparison for squared distance}}


Subtract \(Z\) from the dilation inequality and divide by \(1-r>0\). As \(r\uparrow1\), \((r^{-d}-1)/(1-r)\to d\), while \((q/2)rM\to(q/2)M\). The order-closed limit theorem yields \(M\le(2d/q)Z\). The proof establishes the one-sided derivative limit of the scalar power before passing the inequality to the limit.



The squared-radius integral appears by varying the contraction parameter near one. Keeping its integrability explicit justifies passing the variation through the weighted integral and produces the factor \(d/m\).

\subsection{\texorpdfstring{Integrability of distance under Gibbs}{Integrability of distance under Gibbs}}


The density-measure integrability criterion changes the desired \(\pi_F\)-integral into a constant multiple of \(\int\|x-z\|^2e^{-F(x)}dx\). Its finiteness follows from the continuous, strongly convex mode estimate. No score moment or smoothness assumption is introduced during normalization.



The centered weighted moment is finite and the partition is finite and positive. Their quotient is consequently the ordinary Gibbs second moment, without any extended-value cancellation.

\subsection{\texorpdfstring{The Gibbs concentration bound around the mode}{The Gibbs concentration bound around the mode}}


First prove \(Z>0\), \(Z<\infty\), and integrability of the weighted distance. Then write \(\int\|x-z\|^2d\pi_F=M/Z\) and apply \(M\le(2d/q)Z\). Cancellation is justified by positivity of \(Z\); the result is \(\mathbb E_{\pi_F}\|X-z\|^2\le2d/q\), including \(d=0\).



The bound controls distance from the stated mode rather than the mean. This is the form needed for projected parameter concentration and for the initial moment scales in the coupling construction.

\section{\texorpdfstring{Uniform potential gaps and marginal-\allowbreak{}center displacement}{Uniform potential gaps and marginal-center displacement}}
\label{proof-module-ProximalLangevin-Perturbation-lean}


A one-sided uniform gap between potentials survives latent integration with the same gap size after taking negative logarithms. A smooth strongly convex comparison also supplies analytic control of the original marginal. Combining the marginal sandwich with strong quadratic growth proves an explicit displacement bound for the two actual centers.


\subsection{\texorpdfstring{Quadratic growth needs only strong convexity}{Quadratic growth needs only strong convexity}}


Let \(F\) be \(m\)-strongly convex with \(m\ge0\), and let \(z\) be a global minimizer. Fix any comparator \(x\), and set


\[
A=\frac m2\|x-z\|^2,\qquad G=F(x)-F(z).
\]


Both numbers are nonnegative. For \(0<t<1\), apply strong convexity along the segment from \(z\) to \(x\), and use global minimality at the intermediate point:


\[
F(z)\le F((1-t)z+tx)
\le (1-t)F(z)+tF(x)-t(1-t)A.
\]


Subtract \(F(z)\) and divide by the positive number \(t\). This gives \((1-t)A\le G\) for every such \(t\). To conclude using only a finite comparison, suppose \(G<A\). Then \(A>0\), and the choice


\[
t=\frac{A-G}{2A}\in(0,1/2]
\]


is admissible. But it makes \((1-t)A=(A+G)/2>G\), contradicting the preceding inequality. Therefore \(A\le G\), which is the sharp growth bound


\[
F(x)\ge F(z)+\frac m2\|x-z\|^2.
\]


No derivative of \(F\) is used: the argument needs only the convex-combination inequality and actual global minimality. In the marginal-center comparison, this lower bound turns a finite increase in marginal potential into an upper bound on the squared distance between centers.


\subsection{\texorpdfstring{Integrate the gap before taking logarithms}{Integrate the gap before taking logarithms}}


For any potentials \(V\le U\le V+\delta\), exponentiation and integration in the latent variable give


\[
e^{-\delta}k_V(\theta)\le k_U(\theta)\le k_V(\theta),
\qquad k_W(\theta)=\int e^{-W(\theta,x)}dx.
\]


The smooth strongly convex comparison potential \(V\) supplies positive finite integrals and an integrable Gaussian bound. The original marginal is continuous and decays at infinity; hence it has an actual maximizer. Taking negative logarithms, which is justified by positivity, gives


\[
F_V(\theta)\le F_U(\theta)\le F_V(\theta)+\delta,
\qquad F_W=-\log k_W.
\]


Let \(\theta_U\) minimize \(F_U\) and \(\theta_V\) minimize the \(m\)-strongly convex \(F_V\). Minimality and the sandwich imply \(F_V(\theta_U)-F_V(\theta_V)\le\delta\). Strong-convexity growth about its minimizer therefore yields


\[
\frac m2\|\theta_U-\theta_V\|^2\le\delta.
\]


This argument never differentiates the original marginal and never integrates an asymptotic remainder.



The sharp quadratic-growth lemma uses strong convexity and global minimality, without differentiability of the marginal potential. The comparison model supplies an integrable Gaussian bound for the original weights, giving continuity and existence of an original marginal maximizer by dominated convergence and decay at infinity. Finite positive marginal integrals are established before logarithms and real conversions.




For \(V\le U\le V+\delta\), exponentiation reverses the potential order: \(e^{-\delta}e^{-V}\le e^{-U}\le e^{-V}\). Integrating each latent fiber gives \(e^{-\delta}M_V(\theta)\le M_U(\theta)\le M_V(\theta)\). Positivity and finiteness of the reference marginal then imply the same for \(M_U\); only after that can one use monotonicity of the logarithm to obtain \(F_V\le F_U\le F_V+\delta\), where \(F=-\log M\). These guards exclude taking logarithms of zero or infinite marginal integrals.





Exponentiation reverses the potential order, while integration preserves the order of nonnegative functions. Because the reference marginal is positive and finite, the original marginal inherits both properties. Taking logarithms is now legitimate and gives the same finite additive gap \(\delta\). Performing integration before logarithms avoids any unsupported replacement of a logarithm of an integral by an integral of a logarithm.

\subsection{\texorpdfstring{Establish original-\allowbreak{}marginal existence}{Establish original-marginal existence}}


The inequality \(e^{-U}\le e^{-V}\) and strong-convexity Gaussian bounds prove integrability. Continuity follows by dominated convergence on parameter neighborhoods, using the comparison potential to provide an integrable latent majorant. The same Gaussian estimate proves decay at parameter infinity. A positive continuous function that vanishes at infinity attains its maximum on a sufficiently large compact ball.




The uniform gap dominates the original latent weight by an integrable reference weight, giving finite latent integrals. Locally uniform domination also permits passage to the limit in the parameter and proves continuity of the original marginal integral. The reference's Gaussian decay in the parameter then transfers through the sandwich, so the original marginal vanishes at infinity. A continuous positive marginal attaining a positive value must achieve its maximum on a sufficiently large compact ball. This is how existence of the original center follows from finite smoothing, rather than an added optimizer-existence assumption.





For a convergent parameter sequence, continuity of the original potential gives pointwise convergence of the latent weights. The comparison Gaussian gives an integrable dominating function uniform on that sequence. Dominated convergence proves marginal continuity. Retaining the Gaussian's parameter factor proves decay at infinity, so compact attainment supplies an actual original maximizer. These facts justify using real marginal potentials at every parameter value.

\subsection{\texorpdfstring{Strong growth controls the center shift}{Strong growth controls the center shift}}


Positive finite marginal densities make \(-\log\) order-reversing, so an original marginal maximizer minimizes \(F_U\). At minimizers \(\theta_U,\theta_V\), the sandwich and minimality give \(F_V(\theta_U)\le F_U(\theta_U)\le F_U(\theta_V)\le F_V(\theta_V)+\delta\). The preserved marginal strong convexity yields \(\|\theta_U-\theta_V\|^2\le2\delta/m\); taking the nonnegative square root gives the distance estimate.





The chain of comparisons uses minimality only once, at \(\theta_U\): \(F_U(\theta_U)\le F_U(\theta_V)\). The uniform gap handles the two other links. Thus the comparison-potential increase at \(\theta_U\) is at most \(\delta\), and strong growth gives distance at most \(\sqrt{2\delta/m}\). The original potential need not itself be differentiable or have a Lipschitz gradient.

\section{\texorpdfstring{Sharp Gibbs concentration by dilation}{Sharp Gibbs concentration by dilation}}
\label{proof-module-ProximalLangevin-SharpConcentration-lean}


A radial contraction of space compares an unnormalized Gibbs partition with its second moment. A one-sided scalar derivative at the identity dilation yields the sharp dimension-over-curvature constant. The argument differentiates an explicit polynomial in the dilation factor and requires only continuity and strong convexity of the potential itself.




Let \(F\) be continuous and \(q\)-strongly convex on a finite-dimensional Hilbert space of dimension \(d\), with \(q>0\) and global minimizer \(z\). Gaussian domination supplies the integrals used below.

\subsection{\texorpdfstring{A radial contraction compares partition integrals}{A radial contraction compares partition integrals}}


Strong convexity and minimality imply \(F(z+r(x-z))\le F(x)-\frac q2(1-r^2)\|x-z\|^2\) for \(0<r<1\). Exponentiate and use \(e^u\ge1+u\). Integrating and changing variables under \(x\mapsto z+r(x-z)\), whose Jacobian is \(r^d\), yields

\[Z+\frac q2(1-r^2)M\le r^{-d}Z,\qquad Z=\int e^{-F},\quad M=\int\|x-z\|^2e^{-F}.\]


Let \(w(x)=e^{-F(x)}\), \(Z=\int w\), \(M=\int\|x-z\|^2w(x)\,dx\), and \(d=\dim H\). Continuity and positive strong convexity first give integrability of both \(w\) and \(\|x-z\|^2w\). Strong convexity and the minimizer growth bound combine to give \[F(z+r(x-z))\le F(x)-\tfrac q2(1-r^2)\|x-z\|^2.\] Exponentiation and \(e^u\ge1+u\) imply \[w(z+r(x-z))\ge w(x)+\tfrac q2(1-r^2)\|x-z\|^2w(x).\]





Strong convexity along the segment from the mode to \(x\), together with quadratic growth from that mode, gives the radial deficit \(q(1-r^2)\|x-z\|^2/2\). Exponentiating and using \(e^u\ge1+u\) converts this into a density-plus-moment lower bound. Dilation changes volume by \(r^d\), so integrating the dilated density contributes \(r^{-d}Z\). Gaussian domination makes both sides finite.

\subsection{\texorpdfstring{Take the sharp limit and normalize}{Take the sharp limit and normalize}}


Multiply the dilation inequality by \(r^d\) and define \(\Phi(r)=r^d[Z+(q/2)(1-r^2)M]\). Then \(\Phi(r)\le Z=\Phi(1)\) for \(0<r<1\). Its left difference quotients at one are nonnegative. Differentiating only this scalar polynomial expression gives \(\Phi'(1)=dZ-qM\ge0\), so \(M\le dZ/q\). The partition function \(Z\) is finite and strictly positive, as proved by strong-convexity Gaussian domination. The actual Gibbs normalization gives

\[\int\|x-z\|^2\pi(dx)=M/Z\le d/q,\qquad\pi(dx)=Z^{-1}e^{-F(x)}dx.\]

The proof differentiates no potential and remains valid for continuous nonsmooth strongly convex \(F\).






After obtaining \(qM\le dZ\), strict positivity of \(Z\) permits division. The normalized Gibbs moment is exactly \(M/Z\), by its density definition and the finite integrability of the moment weight. Therefore the estimate is \(d/q\). Neither a derivative of the potential nor a boundary integration identity is needed, making the estimate directly applicable to continuous nonsmooth strongly convex energies.

\subsection{\texorpdfstring{The one-\allowbreak{}sided tangent calculation gives d/\allowbreak{}q}{The one-sided tangent calculation gives d/q}}


Retain \(Z,M,d\) from the dilation inequality and define \[\Phi(r)=r^d\bigl[Z+\tfrac q2(1-r^2)M\bigr].\] For \(0<r<1\), dilation proves \(\Phi(r)\le Z=\Phi(1)\). Thus for \(-1<u<0\), the quotient \((\Phi(1+u)-\Phi(1))/u\) is nonnegative: both numerator and denominator are nonpositive.





For negative increments \(u\) approaching zero, both \(\Phi(1+u)-\Phi(1)\) and \(u\) are nonpositive, so their quotient is nonnegative. Its limit equals \(dZ-qM\), computed by the product rule on the explicit scalar expression. The limit is therefore nonnegative. This sign argument explains why a bound valid only for contractions \(r<1\) still determines the derivative inequality at \(r=1\).

\section{\texorpdfstring{Pointwise particle recursion and square-\allowbreak{}integrable realizations}{Pointwise particle recursion and square-integrable realizations}}
\label{proof-module-IPLA-Dynamics-lean}


An analytic proof can work with random variables identified whenever they agree almost surely, but the algorithm itself is a pointwise recursion. These two descriptions must produce the same probability measures. The finite-time induction below supplies that bridge while keeping the physical parameter and latent updates explicit.


\subsection{\texorpdfstring{Recover the original pairs from scaled coordinates}{Recover the original pairs from scaled coordinates}}


For every particle \(i\), checking parameter and latent coordinates gives

\[A_i z=\operatorname{pair}(p_\theta z,\sqrt N\,\operatorname{block}_i z).\]

Substitution into the finite sum proves

\[V_N(z)=\sum_iU\bigl(p_\theta z,\sqrt N\,\operatorname{block}_iz\bigr).\]

This identifies the abstract lifted potential with the sum of the original pair potentials at the physical particles.





The equality of the original and lifted pair descriptions is coordinatewise. In a parameter coordinate both sides equal \(\theta\); in a latent coordinate both equal \(\sqrt N z_i\). Substituting these equalities in each potential term then identifies the finite sums. This elementary step fixes the interpretation of every force evaluation in the displayed physical recursion.

\subsection{\texorpdfstring{The explicit recursion and its two block formulas}{The explicit recursion and its two block formulas}}


Set \(Z_0=z_0\) and recursively let \(Z_{k+1}=T_\gamma(Z_k,\xi_k)\). The zero and successor identities are the two defining cases of this recursion. Applying the coordinate identities proved above at \(Z_k\) proves

\[\theta_{k+1}=\theta_k-\frac\gamma N\sum_i\nabla_\theta U(\theta_k,x_k^i)+\sqrt{2\gamma/N}\,\xi_k^\theta,\]\[x_{k+1}^i=x_k^i-\gamma\nabla_xU(\theta_k,x_k^i)+\sqrt{2\gamma}\,\xi_k^i,\qquad x_k^i=\sqrt N(Z_k)_i.\]

The particle formula requires \(N>0\), exactly where unscaling cancels \(N\).





For each step, first evaluate every pair force at the old parameter and old physical particles. Then update the common parameter using their average and update each latent particle using its own force. The formulas use the same innovation vector as the scaled state recursion. Thus the conversion between physical and scaled coordinates commutes with advancing the algorithm by one step.

\subsection{\texorpdfstring{Induction identifies the L\ensuremath{{}^{2}} and pointwise algorithms}{Induction identifies the L² and pointwise algorithms}}


At time zero, the representative of the \(L^2\) chain is the initialization. Suppose its representative equals the pointwise chain almost surely at time \(k\). The shared update representative theorem gives the Euler formula almost surely; intersect that full-measure event with the induction event, substitute the previous-state equality, and rewrite \(2(\gamma/N)=2\gamma/N\). This proves almost-sure equality at \(k+1\). Finally, almost-everywhere equal measurable maps have equal pushforward measures, hence

\[\rho_k=P\circ\bigl[\omega\mapsto Z_k^{\mathrm{pt}}(\omega)\bigr]^{-1}.\]




Only finitely many almost-sure equalities are needed in each induction step. Their intersection has probability one, and on this event the two successor formulas are literally equal. There is therefore no need to select representatives agreeing at every real time. For any measurable set \(A\), equality almost surely gives equal probabilities of membership in \(A\); this is exactly equality of the two pushforward laws.

\section{\texorpdfstring{Subgaussian Gibbs score and its sharp second moment}{Subgaussian Gibbs score and its sharp second moment}}
\label{proof-module-UnadjustedLangevin-GibbsScore-lean}


Differentiability, an \(L\)-Lipschitz gradient, and finite Gibbs normalization suffice for \(\int\|\nabla U\|^2d\pi\le Ld\). Strong convexity is used only by the final result to derive normalization.



The smooth upper model at \(x\), evaluated at \(x-v\), gives \(\langle\nabla U(x),v\rangle\le U(x)-U(x-v)+L\|v\|^2/2\). Multiplication by \(e^{-U(x)}\) and translation of the Lebesgue integral therefore bound the directional score moment generating function by \(e^{L\|v\|^2/2}\). Centering and the variance estimate follow by examining small scalar multiples of \(v\). Summing over an orthonormal basis gives the score-energy scale \(Ld\).

\subsection{\texorpdfstring{The exponential score comparison}{The exponential score comparison}}


The descent lemma gives \(U(x+v)\le U(x)+\langle\nabla U(x),v\rangle+(L/2)\|v\|^2\). Negate and exponentiate, then rearrange to \(e^{-U(x)}e^{-\langle\nabla U(x),v\rangle}\le e^{L\|v\|^2/2}e^{-U(x+v)}\). The inequality is pointwise and nonnegative, so it can then be integrated.



Evaluate the smooth upper model at \(x-v\), move the linear gradient term to the other side, and exponentiate. The factor \(e^{-U(x)}\) then turns the translated potential into a translated Gibbs weight.

\subsection{\texorpdfstring{The directional log moment-\allowbreak{}generating bound}{The directional log moment-generating bound}}


Integrate the pointwise exponential-score bound. Translation invariance gives \(\int e^{-U(x+v)}dx=Z_U\); normalization cancels this factor, leaving \(\mathbb E_{\pi_U}e^{-\langle\nabla U,v\rangle}\le e^{L\|v\|^2/2}\). Finiteness of \(Z_U\) is explicit because extended inverse cancellation needs it.



Translation invariance gives \(\int e^{-U(x-v)}dx=Z\). Dividing by the same positive partition cancels the unnormalized mass and leaves the directional exponential bound with variance proxy \(L\).

\subsection{\texorpdfstring{The sub-\allowbreak{}Gaussian directional score}{The sub-Gaussian directional score}}


Apply the preceding bound with the vector \(-tv\), for every real \(t\). The exponent becomes \(t\langle\nabla U,v\rangle\) and the right side is \(e^{t^2L\|v\|^2/2}\). This is exactly the directional sub-Gaussian moment bound with variance proxy \(L\|v\|^2\); the proof also establishes exponential integrability required by that structure.



Applying the exponential bound to positive and negative scalar multiples of a direction forces the first derivative at zero to vanish. The second-order comparison gives directional score variance at most \(L\|v\|^2\).

\subsection{\texorpdfstring{A square bounded by two exponentials}{A square bounded by two exponentials}}


The series \(\cosh a=\sum_{n\ge0}a^{2n}/(2n)!\) has nonnegative terms, so retaining its first two terms gives \(\cosh a\ge1+a^2/2\). Substitute \(2\cosh a=e^a+e^{-a}\) and rearrange to \(a^2\le e^a+e^{-a}-2\). The construction uses the power series for the hyperbolic cosine with the lower bound obtained by discarding nonnegative series terms, explicitly proving nonnegativity of every discarded even-power term.



The even combination \(e^s+e^{-s}-2\) has leading term \(s^2\) and nonnegative higher even powers. This supplies a pointwise domination of the square by exponential terms whose Gibbs integrals are controlled.

\subsection{\texorpdfstring{A second moment from exponential control}{A second moment from exponential control}}


Integrate \(t^2X^2\le e^{tX}+e^{-tX}-2\), giving \(t^2\mathbb EX^2\le2(e^{ct^2/2}-1)\). Subgaussianity provides the exponential integrability and hence the square integrability needed here. Divide by \(t^2\) for \(t\ne0\), and let \(t\to0\) to obtain \(\mathbb EX^2\le c\). The limit preserves the inequality because both sides converge.



Apply the scalar inequality to a small multiple of the score, integrate using the two exponential bounds, divide by the squared multiplier, and send it to zero. The limit yields the variance proxy as a true second-moment bound.

\subsection{\texorpdfstring{The full Gibbs gradient-\allowbreak{}energy bound}{The full Gibbs gradient-energy bound}}


Choose an orthonormal basis \((e_i)_{i=1}^d\) and use Parseval: \(\|\nabla U(x)\|^2=\sum_i\langle\nabla U(x),e_i\rangle^2\). Each coordinate is subgaussian with proxy \(L\), hence has second moment at most \(L\). Finite-sum integrability permits integration term by term, yielding \(\mathbb E_{\pi_U}\|\nabla U\|^2\le Ld\).



For an orthonormal basis, the squared gradient norm is the sum of the squared directional scores. Each integral is at most \(L\), so summing the \(d\) finite terms gives \(Ld\).

\subsection{\texorpdfstring{Square-\allowbreak{}integrability of the Gibbs gradient}{Square-integrability of the Gibbs gradient}}


The coordinate subgaussian bounds prove integrability of every squared gradient coordinate. Summing over the finite orthonormal basis proves integrability of \(\|\nabla U\|^2\). Lipschitz continuity provides strong measurability, and the square-integrability square criterion packages those facts as \(\nabla U\in L^2(\pi_U)\).



Continuity of the gradient supplies measurability and the preceding energy bound supplies the finite square integral. Together they establish the full \(L^2\) membership needed for the lifted reference drift.

\subsection{\texorpdfstring{The gradient-\allowbreak{}energy bound under strong convexity}{The gradient-energy bound under strong convexity}}


Strong convexity with \(m>0\) supplies the finite-partition premise of the score theorem. Substitute that proof into the first-order smoothness result to obtain \(\mathbb E_{\pi_U}\|\nabla U\|^2\le Ld\). The strong convexity constant is used for normalization here; it does not appear in the sharp score bound.



Strong convexity is used here to guarantee a normalized Gibbs probability and finite partition. The score-energy calculation itself comes from the smooth upper model and the translation comparison.

\section{\texorpdfstring{Finite quadratic transport transfers second moments}{Finite quadratic transport transfers second moments}}
\label{proof-module-UnadjustedLangevin-TransportMoments-lean}


Finite \(W_2\) distance to a finite-second-moment law supplies the other second moment automatically. This removes a redundant integrability premise when passing an initial law to the L\ensuremath{{}^{2}} formulation.



A finite-cost coupling transfers second moments through \(\|x\|^2\le2\|x-y\|^2+2\|y\|^2\). Thus if one marginal has finite second moment and the transport infimum is finite, a finite-cost coupling proves the other marginal has one as well. Conversely, the product coupling of two square-integrable probability measures has finite cost by the same inequality. Applying these statements to laws transfers the moment conclusion back to the underlying random vectors.

\subsection{\texorpdfstring{A finite transport distance supplies a second moment}{A finite transport distance supplies a second moment}}


Finite Wasserstein distance supplies a finite-cost coupling. On that coupling space the second coordinate is in \(L^2\) by the reference moment, and the coordinate difference is in \(L^2\) by finite coupling cost. Their sum is the first coordinate, so its law \(\mu\) has a finite second moment. No optimal coupling is selected.



Explicitly, \(\|x\|^2\le2\|x-y\|^2+2\|y\|^2\). Integrating on a finite-cost coupling bounds the first marginal's second moment by two finite quantities. Thus finiteness of transport to one finite-moment reference determines the ordinary second-moment class of the other marginal.

\subsection{\texorpdfstring{Second moments give a finite transport distance}{Second moments give a finite transport distance}}


Use the explicit independent coupling \(\mu\otimes\pi\). Each coordinate is in \(L^2\), hence their difference is in \(L^2\) and has finite cost. Wasserstein distance is at most this coupling cost, proving finiteness.



For independent \(X\sim\mu\) and \(Y\sim\pi\), the same squared triangle inequality gives \(\mathbb E\|X-Y\|^2\le2\mathbb E\|X\|^2+2\mathbb E\|Y\|^2\). Taking a square root supplies an explicit finite admissible cost, hence an upper bound for the transport infimum.

\subsection{\texorpdfstring{Transferring the moment statement through a law}{Transferring the moment statement through a law}}


The finite-transport theorem first proves \(\mathrm{id}\in L^2(\mu)\). Replace \(\mu\) by the supplied pushforward law of \(X\), then use the change-of-law \(L^2\) criterion to pull this moment back to \(X\) on its original sample space.



The law change is the identity \(\int\|x\|^2d\mu(x)=\mathbb E\|X\|^2\). It transfers the moment conclusion to the original random variable without constructing a new realization or asserting any independence of that realization from other variables on its space.

\subsection{\texorpdfstring{The moment implication for pushforward laws}{The moment implication for pushforward laws}}


Specialize the preceding theorem to \(\mu=P\circ X^{-1}\), whose law identity is reflexive. Finite transport of the actual pushforward to a finite-moment reference therefore gives \(X\in L^2(P)\).



The statement applies to the actual measured random vector, not merely to an abstract law with the same name. Measurability supplies the pushforward change of variables, while finite transport and the reference moment supply finiteness of the resulting squared-norm integral.

\section{\texorpdfstring{Normalization, moments and score moments of the lifted target}{Normalization, moments and score moments of the lifted target}}
\label{proof-module-IPLA-Foundations-lean}


Strong convexity supplies Gaussian domination of the enlarged Gibbs density, and the scaled geometry supplies its dimension and smoothness constants. These facts establish the normalization and finite moments needed for transport. They also give a quantitative equilibrium score moment, which controls the local drift defect in the shared Langevin coupling argument.




Apply the shared ULA Gibbs results to \(V_N\), whose differentiability and \(Nm\)-strong convexity were derived from \(U\). All normalization and moment facts are conclusions from the original joint assumptions.

\subsection{\texorpdfstring{The named target is the normalized Gibbs measure}{The named target is the normalized Gibbs measure}}


The definitions give \(\Pi_N(dz)=Z_N^{-1}e^{-V_N(z)}\,dz\). If \(N>0\) and \(m>0\), then \(Nm>0\). The lifted strong-convexity theorem and chain-rule differentiability therefore meet the shared Gibbs theorem's hypotheses. Its quadratic lower bound gives integrable Gaussian tails and a positive partition, proving \(0<Z_N<\infty\) and thus \(\Pi_N(\mathbb R^D)=1\). The same tails imply \(\int\|z\|^2d\Pi_N<\infty\), expressed as membership of the identity in \(L^2(\Pi_N)\).






The lower bound has the form \(V_N(z)\ge a\|z\|^2-b\) with \(a>0\). Therefore the density is bounded by \(e^b e^{-a\|z\|^2}\), and multiplication by \(\|z\|^2\) still leaves an integrable Gaussian majorant. Strict positivity of the exponential on a nonzero Euclidean volume space gives \(Z_N>0\). Dividing by this finite positive constant produces the stated probability and second moment.

\subsection{\texorpdfstring{Finite transport supplies square integrability of initialization}{Finite transport supplies square integrability of initialization}}


Suppose a measurable initialization \(X\) has \(W_2(X_\#P,\Pi_N)<\infty\). The target already has a finite second moment. For any finite-cost coupling, \(\|x\|^2\le2\|x-y\|^2+2\|y\|^2\); integrating transfers the target moment and coupling cost to the initial law. The shared transport-moment theorem formalizes this argument and proves \(X\in L^2(P)\). Thus this alternative initialization interface introduces no independent moment assumption.





A finite-cost coupling is enough; attainment of the Wasserstein infimum is unnecessary. Pull the coordinate functions onto its joint probability space and integrate the pointwise inequality. Its first marginal identifies the desired initial second moment, and its second marginal identifies the already finite Gibbs moment. The result also explains why a finite Wasserstein initialization condition can replace a separately stated square-integrability condition.

\subsection{\texorpdfstring{The lifted Gibbs score moment has the correct dimension}{The lifted Gibbs score moment has the correct dimension}}


The shared score estimate for a strongly convex potential with \(\ell\)-Lipschitz gradient is \(\int\|\nabla V\|^2d\pi\le\ell\dim H\). Substitute \(V=V_N\), \(\ell=NL\), and \(\dim H=d_\theta+Nd_x\) to obtain

\[\int\|\nabla V_N(z)\|^2d\Pi_N(z)\le NL(d_\theta+Nd_x).\]

The lifted smoothness and convexity hypotheses follow from the joint ones. The proof only rewrites the common Gibbs definition and exact dimension; it adds no score-moment or Hessian premise.






The bound involves the full lifted dimension, not merely the parameter dimension, because the score has both parameter and particle components. Its scaling is \(NL D\), where \(D=d_\theta+Nd_x\). When inserted into a physical step of length \(\gamma/N\), these factors combine with the step scaling. Keeping the score moment separate from the sharper parameter concentration estimate prevents confusing two different dimensions in the proof.

\section{\texorpdfstring{An abstract family of compatible finite-\allowbreak{}coupling references}{An abstract family of compatible finite-coupling references}}
\label{proof-module-UnadjustedLangevin-ReferenceFamily-lean}


This interface isolates the passage from a bound for each initial coupling to a theorem with the exact initial Wasserstein distance. A concrete product construction later supplies this interface without assuming an optimal plan.



The reference-family assumption quantifies over every finite-cost coupling of the initial and target laws. Each member supplies a common-noise reference with the same output law and constants, while its initial separation retains that coupling's exact cost. This quantifier order is necessary before taking a transport infimum. A reference for just one chosen initial pair gives a valid coupling bound, but does not by itself replace that initial cost by Wasserstein distance.

\subsection{\texorpdfstring{Couplings with finite quadratic cost}{Couplings with finite quadratic cost}}


Restrict the family of couplings to those with \(c\ne\infty\). This family retains the integrability needed to realize initial class distances and convert them to real values while taking the Wasserstein infimum.



Finite cost means \(\int\|x-y\|^2dc(x,y)<\infty\). It controls the difference between initial states even before separate moment estimates for each coordinate are used. Restricting to these couplings does not alter a finite transport infimum, since infinite costs cannot improve it.

\subsection{\texorpdfstring{A reference for every finite-\allowbreak{}cost initial coupling}{A reference for every finite-cost initial coupling}}


For each finite coupling, supply initial and noise \(L^2\) variables and an l2reference. The fields identify initial law, bound initial class distance by that coupling's cost, and require one common output law across the family. These are explicit construction obligations, with no convergence estimate among them.



For each initial coupling, extend the probability space with the same independent noise law. The numerical endpoint must depend only on its initial marginal and that noise, so its law is common to all extensions. The reference endpoint retains the target law; these law identities make comparison across couplings meaningful.

\subsection{\texorpdfstring{Finiteness of the output transport distance}{Finiteness of the output transport distance}}


Choose one finite initial coupling using the finite initial Wasserstein premise. Its supplied numerical/\allowbreak{}reference \(L^2\) variables yield a finite-cost output coupling. The output-law and stationary-reference fields identify its marginals, proving \(W_2(\nu_n,\pi)<\infty\).



One finite initial coupling is enough to exhibit a finite output coupling, since both endpoint random vectors then have a square-integrable difference. This preliminary finiteness result permits the later convergence inequality to be stated with ordinary real Wasserstein distances.

\subsection{\texorpdfstring{Passing the analytic bound through the initial infimum}{Passing the analytic bound through the initial infimum}}


Apply the analytic \(L^2\) estimate to every finite coupling realization in the family. Replace its initial error by the coupling cost and its numerical law by the common \(\nu_n\). Because \((1-mh)^n\ge0\), the affine finite-coupling infimum theorem gives the displayed bound with the true initial Wasserstein distance.



The coefficient of the initial cost is nonnegative, so taking its infimum preserves the affine bound. A negative coefficient would reverse the minimization and invalidate this step. The small-step restrictions guarantee the required sign as well as strict contraction.

\section{\texorpdfstring{The initial transport infimum}{The initial transport infimum}}
\label{proof-module-IPLA-ReferenceFamily-lean}


Transport distance is an infimum of coupling costs. To obtain an initial Wasserstein term from a synchronous process estimate, one must construct the reference for arbitrary finite-cost initial couplings and show that its algorithm marginal is independent of the chosen coupling. These are the two substantive requirements of the initial-coupling argument.


\subsection{\texorpdfstring{Take the infimum without assuming an optimal coupling}{Take the infimum without assuming an optimal coupling}}


A reference family supplies, for every finite-cost coupling \(c\) of \(\mu_0\) and \(\pi\), an \(L^2\) realization of the initial variables and common Euler noise. Its fields require that the initial \(L^2\) difference be at most \(\operatorname{cost}(c)\), and that every realization produce the same output law \(\mu_k\). The shared coupling theorem gives \(W_2(\mu_k,\pi)\le q^k\operatorname{cost}(c)+b\) for each such \(c\). Taking the cost infimum yields \(q^kW_2(\mu_0,\pi)+b\). A finite-cost coupling also proves output finiteness before real conversion. These fields are explicit construction obligations in the reusable interface; the public Gaussian theorem discharges them.




For every \(\varepsilon>0\), the definition of a finite infimum supplies a coupling with cost at most \(W_2(\mu_0,\pi)+\varepsilon\). Apply the same process estimate to it and let \(\varepsilon\) decrease to zero. Because the output law and additive bias are common to the entire family, the left side does not vary with this choice. This proves the transport bound without requiring existence of an optimal coupling.

\section{\texorpdfstring{Exact particle-\allowbreak{}number and time scaling in the transport bound}{Exact particle-number and time scaling in the transport bound}}
\label{proof-module-IPLA-Scaling-lean}


The particle theorem is obtained by exact substitution in a Euclidean Langevin estimate. The relevant effective dimension is \(D/N=d_x+d_\theta/N\), because the physical iteration uses a step divided by particle count. This reduction is algebraic and introduces no additional approximation or loss in the numerical constant.


\subsection{\texorpdfstring{Step restriction and square-\allowbreak{}root bias}{Step restriction and square-root bias}}


Since \(N>0\), dividing both sides by \(N\) preserves order, and \(Nm+NL=N(m+L)\). Therefore \(\gamma/N\le2/(Nm+NL)\) if and only if \(\gamma\le2/(m+L)\). The ratio \(NL/(Nm)\) cancels to \(L/m\), and

\[\frac\gamma N(d_\theta+Nd_x)=\gamma\left(d_x+\frac{d_\theta}{N}\right).\]

The intensive dimension in parentheses is nonnegative. The real square-root multiplication rule therefore factors the square root, giving the exact identity

\[\frac{33}{20}\frac{NL}{Nm}\sqrt{\frac\gamma N D}=\frac{33}{20}\frac Lm\sqrt{d_x+d_\theta/N}\sqrt\gamma.\]


The numerical substitutions are \(m_N=Nm\), \(L_N=NL\), \(h=\gamma/N\), \(D=d_\theta+Nd_x\). Positive \(N\) cancels the common factor in the step restriction. In the discretization term, \(L_N/m_N=L/m\) and \[hD=\gamma(d_x+d_\theta/N).\] Since the intensive dimension is nonnegative, the square-root product rule gives \(\sqrt{hD}=\sqrt{d_x+d_\theta/N}\sqrt\gamma\). These are exact identities, so this rewrite loses no constant.





All cancellations involving \(N\) occur under \(N>0\). The dimension identity separates the latent contribution, which remains of order \(d_x\), from the parameter contribution, which is reduced by \(N\). Thus increasing particle count reduces the parameter part of the discretization estimate but does not remove the latent Euler bias at a fixed physical step. The admissible physical step itself is independent of particle count.

\subsection{\texorpdfstring{Substitute in the Hilbert-\allowbreak{}space Dalalyan theorem}{Substitute in the Hilbert-space Dalalyan theorem}}


For an explicit finite-cost reference family and finite initial \(W_2\), the shared theorem gives

\[W_2(\nu_k,\pi)\le(1-m_Nh)^kW_2(\nu_0,\pi)+\frac{33}{20}\frac{L_N}{m_N}\sqrt{hD}.\]

The assumptions \(m\le L\), \(m>0\), \(\gamma>0\), and the base step restriction imply the corresponding lifted inequalities by multiplication or division by positive \(N\). Finally \(Nm(\gamma/N)=m\gamma\), and the preceding bias identity yields the displayed IPLA transport estimate. This reusable intermediate theorem accepts the reference family explicitly; the public Gaussian theorem constructs it through shared ULA results.






The transient factor rescales as \((1-m_Nh)^k=(1-m\gamma)^k\). The coupling theorem's initial term remains the Wasserstein distance in the scaled state geometry, so no additional factor should be attached to it. The only dimension-dependent simplification is in the discretization floor. Applying these identities to the entire inequality yields the particle transport bound with its original finite-time and step-size meaning.

\section{\texorpdfstring{Project transport convergence into parameter RMS error}{Project transport convergence into parameter RMS error}}
\label{proof-module-UnadjustedLangevin-Projection-lean}


For a continuous linear contraction \(T:H\to K\), reference point \(b\), and equilibrium projected radius \(c\), the parameter error obeys \((\int\|Tx-b\|^2d\mu)^{1/2}\le W_2(\mu,\pi)+c\). This derives the estimate directly from couplings.



For a 1-Lipschitz projection \(P\) and a target point \(a\), Minkowski gives \((\mathbb E\|P(X)-a\|^2)^{1/2}\le(\mathbb E\|X-Y\|^2)^{1/2}+(\mathbb E\|P(Y)-a\|^2)^{1/2}\). The first term may be minimized over couplings, and the second depends only on the reference marginal. This converts a full-state transport estimate plus a target concentration estimate into a projected root-mean-square error bound.

\subsection{\texorpdfstring{The projected root mean square error}{The projected root mean square error}}


Define the extended projected error \(R_{T,b}(\mu)=\|x\mapsto Tx-b\|_{L^2(\mu)}\). Using an extended norm retains infinity for a divergent second moment and allows the following transport comparison before any finiteness assumption.



This quantity is an error about the specified deterministic center \(b\), rather than a variance about the mean. It can therefore represent a parameter estimation error about a marginal mode. Keeping the extended integral initially allows the comparison to prove, rather than assume, its finiteness.

\subsection{\texorpdfstring{Its finite ordinary value}{Its finite ordinary value}}


Whenever the projected displacement has finite second moment, its projected root mean square radius is the ordinary finite number \(R_{T,b}(\mu)=(\int\|Tx-b\|^2d\mu)^{1/2}\). The transport and reference-radius bounds below establish this finiteness before the radius is used in a numerical inequality.



Once \(\int\|Tx-b\|^2d\mu<\infty\) is proved, its nonnegative square root is an ordinary real number. All later real-valued estimates use this finite case; a divergent second moment is never assigned a finite statistical-error interpretation.

\subsection{\texorpdfstring{The pointwise projected difference bound}{The pointwise projected difference bound}}


If \(\|T\|\le1\), then \(\|T(x-y)\|\le\|x-y\|\) pointwise. Monotonicity of the extended \(L^2\) norm bounds the projected difference norm by the original coupling cost.



Squaring the pointwise contraction, integrating under an arbitrary coupling, and taking nonnegative square roots gives \(\|T(X-Y)\|_{L^2}\le\|X-Y\|_{L^2}\). No relation between the projected coordinates and the unobserved coordinates is required.

\subsection{\texorpdfstring{Projected error bounded by a coupling}{Projected error bounded by a coupling}}


On a coupling write \(Tx-b=T(x-y)+(Ty-b)\). Minkowski bounds its \(L^2\) norm by the sum of the two norms; the contraction estimate controls the first by coupling cost. Marginal identities rewrite the outer norms as \(R_{T,b}(\mu)\) and \(R_{T,b}(\pi)\).



The decomposition is valid because \(T\) is linear. The second term depends only on the reference marginal, so its value is the same for every coupling. This uniformity is precisely what lets the next step minimize only the full-state coupling cost.

\subsection{\texorpdfstring{Projected error bounded by Wasserstein distance}{Projected error bounded by Wasserstein distance}}


The preceding bound holds for every coupling, with the same equilibrium radius. Take the infimum of coupling costs to obtain \(R_{T,b}(\mu)\le W_2(\mu,\pi)+R_{T,b}(\pi)\). Nonemptiness supplies the indexed-infimum identity; no minimizing plan is assumed.



For any \(\varepsilon>0\), choose a coupling whose cost is within \(\varepsilon\) of the finite transport infimum. The couplingwise inequality then bounds the projected error by \(W_2+\varepsilon\) plus the reference radius. Sending \(\varepsilon\downarrow0\) gives the claimed estimate without an optimizer.

\subsection{\texorpdfstring{Finiteness of projected error}{Finiteness of projected error}}


Finite transport guarantees a coupling exists and makes \(W_2(\mu,\pi)\) finite. The assumed equilibrium radius is bounded by a finite real embedding. The transport-radius inequality therefore bounds the current projected norm by a sum of two finite extended values, proving its finiteness.



The argument supplies the projected moment even if no full second-moment bound for the first law has been stated separately. It needs only finite quadratic transport and a finite reference projected radius. This is useful when the statistical quantity is lower-dimensional than the full sampling state.

\subsection{\texorpdfstring{The ordinary squared-\allowbreak{}moment formula}{The ordinary squared-moment formula}}


The finite projected square norm and measurability give \(x\mapsto Tx-b\in L^2(\mu)\). Its norm is therefore exactly \(\sqrt{\int\|Tx-b\|^2d\mu}\), a finite nonnegative number. Thus the projected radius estimate concerns an actual finite second moment.



Absolute integrability of the squared norm is established before using the ordinary integral. The identity then follows from \(\|F\|_{L^2}^2=\int\|F\|^2\), with the nonnegative square root selecting the norm rather than its negative.

\subsection{\texorpdfstring{Inserting a target second-\allowbreak{}moment bound}{Inserting a target second-moment bound}}


The supplied \(L^2\) membership makes the squared-norm integral genuine. Its bound by \(c^2\), together with \(c\ge0\), implies the root-mean-square norm is at most \(c\). Re-embed this finite real inequality to obtain the extended radius bound.



If the reference second moment is bounded by a radius squared, the radius must be nonnegative to conclude the desired norm bound after taking roots. The degenerate choice \(c=0\) forces the projection to equal the center almost everywhere.

\subsection{\texorpdfstring{The finite projected error estimate}{The finite projected error estimate}}


First prove current-radius finiteness using finite transport and the finite equilibrium bound. Then convert the extended transport-radius inequality to real values, using finiteness of every summand. The result is \(R_{T,b}^{\mathbb R}(\mu)\le W_2(\mu,\pi)+c\).



The final bound separates two sources of error: transport measures how far the numerical full-state law is from the target, while \(c\) measures target fluctuations around the chosen center. Only the first term can decay through additional sampling iterations at a fixed target.

\section{\texorpdfstring{Projection of quadratic transport to parameter error}{Projection of quadratic transport to parameter error}}
\label{proof-module-IPLA-Transport-lean}


Projection controls parameter RMS through full-state transport and an equilibrium parameter radius. The argument works on each coupling separately, using the ordinary triangle inequality in the Hilbert space of square-integrable random variables. Taking the coupling infimum then yields the desired Wasserstein contribution.


\subsection{\texorpdfstring{Minkowski on each coupling gives the projected RMS estimate}{Minkowski on each coupling gives the projected RMS estimate}}


For a linear contraction \(T\), any center \(b\), and a coupling \(c\) of \(\mu,\pi\), write \(Tx-b=T(x-y)+(Ty-b)\). Minkowski and \(\|T(x-y)\|\le\|x-y\|\) imply

\[\|Tx-b\|_{L^2(\mu)}\le\operatorname{cost}(c)+\|Ty-b\|_{L^2(\pi)}.\]

Taking the infimum over couplings proves the same inequality with \(W_2(\mu,\pi)\). Finite transport plus a finite equilibrium radius proves finiteness of the projected moment. Only then does the shared real conversion identify its value with \(\sqrt{\int\|Tx-b\|^2d\mu}\). A second-moment bound at equilibrium supplies the concentration radius. No optimal coupling or separate Wasserstein triangle inequality is needed.




With coupling variables \(X,Y\), the equality \(TX-b=T(X-Y)+(TY-b)\) separates sampling discrepancy from equilibrium spread. A bounded linear contraction gives \(\|T(X-Y)\|_2\le\|X-Y\|_2\), and Minkowski adds the second term. Both marginals identify the two RMS quantities. Finiteness of one coupling cost and the equilibrium radius first proves a finite sampling RMS, so the final square-root expectation has its usual meaning.

\section{\texorpdfstring{The dimension-\allowbreak{}sharp equilibrium parameter concentration}{The dimension-sharp equilibrium parameter concentration}}
\label{proof-module-IPLA-Concentration-lean}


The parameter target is a Gibbs law in parameter dimension alone, with potential \(N\Phi\). Its strong-convexity constant grows linearly with particle count. Applying the available continuous-potential concentration theorem and then transferring through the exact parameter projection produces the equilibrium term in the strongly convex particle error bound.




The exact parameter marginal has density proportional to \(M(\theta)^N=e^{-N\Phi(\theta)}\). Its potential is \(Nm\)-strongly convex in dimension \(d_\theta\), giving the equilibrium radius \(\sqrt{2d_\theta/(Nm)}\) without a separate concentration hypothesis.

\subsection{\texorpdfstring{Identify the scaled marginal potential and its Gibbs measure}{Identify the scaled marginal potential and its Gibbs measure}}


Multiply the strong-convexity inequality for \(\Phi\) by the nonnegative scalar \(N\). Every affine term and the quadratic penalty scale by \(N\), so \(N\Phi\) is \(Nm\)-strongly convex. Since \(M>0\),

\[e^{-N\Phi(\theta)}=e^{N\log M(\theta)}=M(\theta)^N.\]



The identity \(e^{-N\Phi}=M^N\) holds pointwise because \(M\) is finite and positive. Hence the two unnormalized measures, not merely their modes, coincide. Their normalizations coincide as well. This justifies applying a Gibbs concentration theorem directly to the statistical marginal target and explains why its potential has modulus \(Nm\).

\subsection{\texorpdfstring{A dilation estimate gives the sharp parameter moment}{A dilation estimate gives the sharp parameter moment}}


A marginal maximizer \(\theta_*\) minimizes \(\Phi\), hence also \(N\Phi\). The shared continuous strongly convex Gibbs theorem, which needs continuity and a mode rather than differentiability of \(\Phi\), bounds the squared distance by twice dimension over convexity:

\[\int\|\theta-\theta_*\|^2d\pi_N^\theta(\theta)\le\frac{2d_\theta}{Nm}.\]

Apply it to the continuous potential \(N\Phi\), use the proved \(Nm\)-strong convexity, and rewrite its Gibbs law by the preceding identity. Only the parameter dimension enters; latent dimension does not enlarge this equilibrium radius.





The estimate here is the established bound with factor two. It uses only continuity, positive strong convexity and an actual minimizing point of the marginal potential. It therefore remains applicable without proving that the marginal potential is differentiable. The latent variables have already been integrated away, so the dimension in the numerator is \(d_\theta\); replacing it by the enlarged state dimension would unnecessarily weaken the statistical concentration estimate.

\subsection{\texorpdfstring{Transfer the marginal moment to the lifted Gibbs target}{Transfer the marginal moment to the lifted Gibbs target}}


The exact projection identity \(p_{\theta\#}\Pi_N=\pi_N^\theta\) and the measurable pushforward integral formula yield

\[\int\|p_\theta z-\theta_*\|^2d\Pi_N(z)\le\frac{2d_\theta}{Nm}.\]

The projected Gibbs variable is already in \(L^2\), and subtracting the constant \(\theta_*\) preserves this. The shared projected-RMS theorem converts the real moment inequality to the extended RMS bound

\[\operatorname{RMS}_{\Pi_N}(p_\theta-\theta_*)\le\sqrt{2d_\theta/(Nm)}.\]

The radius is nonnegative and its square equals the moment bound, so the square-root step is legitimate. This term needs joint differentiability and strong convexity, with no gradient-Lipschitz premise.






For the nonnegative observable \(g(\theta)=\|\theta-\theta_*\|^2\), the pushforward identity gives \(\int g(p_\theta z)d\Pi_N=\int g(\theta)d\pi_N^\theta\). Thus the same finite upper bound is valid in the enlarged space without estimating latent deviations. Taking its nonnegative square root yields the equilibrium RMS radius. This radius can then be added to a full-state coupling error by Minkowski's inequality.

\section{\texorpdfstring{Realize every initial coupling independently of the noise}{Realize every initial coupling independently of the noise}}
\label{proof-module-UnadjustedLangevin-InitialExtension-lean}


Given any coupling \(\gamma\) of \(\mu\) and \(\pi\), extend it to \(\gamma\otimes Q\). This preserves its quadratic cost and supplies independent noise without requiring an optimal coupling.



Starting from a coupling \(\eta(dx,dy)\), adjoin an independent noise space by the product measure \(\eta\otimes Q\). The two coordinate projections retain their initial marginal laws and their joint law \(\eta\), while the noise retains law \(Q\) and is independent of the entire pair. Integrating a function of the initial pair therefore gives its integral under \(\eta\); in particular the initial \(L^2\) separation is exactly the quadratic coupling cost. This construction lets one repeat the same reference argument for every finite-cost initial coupling.

\subsection{\texorpdfstring{Normalization of an initial coupling}{Normalization of an initial coupling}}


The coupling's first marginal is the initial probability. Pushforward preserves total mass, so \(c(H\times H)=\mu(H)=1\). This constructs normalization of the joint measure using the defining marginal identity of the coupling.



For any joint coupling measure, total mass equals the total mass of either marginal. Its first marginal is the initial probability, so the joint measure has mass one. This observation establishes normalization of the coupling before forming a product with noise. It uses only the defining marginal equality and does not require any finite-cost premise.

\subsection{\texorpdfstring{Adjoining independent noise}{Adjoining independent noise}}


Extend an initial coupling \(c(dx,dy)\) by an independent noise sample \(s\sim Q\), using \(c\otimes Q\) on \((H\times H)\times S\). Both algorithms retain the prescribed initial joint law while sharing fresh noise independent of the pair.



The extended variable is \(((x,y),s)\), where \((x,y)\) retains the chosen initial coupling and \(s\) is sampled independently from the fixed source. Sharing this one source later synchronizes the numerical and reference noises. The construction changes neither initial marginal nor initial dependence, while making noise independence explicit at the level of a product probability.

\subsection{\texorpdfstring{Normalization of the product extension}{Normalization of the product extension}}


The coupling is a probability by its first marginal, and \(Q\) is a probability by hypothesis. Their product therefore has total mass one. This normalization supports all subsequent \(L^2\) constructions on the extension.



The product's total mass is \(c(E\times E)Q(S)=1\). Thus every extended source is a probability space, independently of which initial coupling was chosen. This normalization is used when transferring second moments, integrating constants and applying probability-space \(L^2\) estimates. Finite coupling cost is a separate condition introduced only for the moment constructions.

\subsection{\texorpdfstring{The retained joint initial law}{The retained joint initial law}}


Projecting \(c\otimes Q\) to the first factor integrates out the probability noise coordinate. Since \(Q(S)=1\), the resulting measure is exactly \(c\), preserving the entire initial dependence structure.



For a measurable set \(A\) of initial pairs, its extended preimage is \(A\times S\), with measure \(c(A)Q(S)=c(A)\). Therefore the joint initial law remains exactly \(c\), including its dependence structure. In particular an initial cost depending on both coordinates is preserved, not just the separate marginal distributions.

\subsection{\texorpdfstring{The retained numerical initial marginal}{The retained numerical initial marginal}}


Compose the joint projection with its first state projection. The extension's joint marginal is \(c\), whose first state marginal is \(\mu\). Pushforward composition therefore yields the specified algorithm initialization law.



Project first to the initial pair and then to its first coordinate. Pushforward composition and the first marginal condition give law \(\mu\). The auxiliary reference coordinate and fresh noise are both integrated out. This ensures the algorithm begins with the same prescribed law on every coupling extension, even when the couplings encode different initial dependence.

\subsection{\texorpdfstring{The retained reference initial marginal}{The retained reference initial marginal}}


Compose projection to the initial pair with its second state projection. The coupling's second marginal is \(\pi\), so the reference initialization on the extension has exactly law \(\pi\).



The second coordinate the stated pushforward law \(\pi\) by the coupling's second marginal identity. Extending by independent noise leaves this distribution unchanged. Thus a reference initialized from the second coordinate starts at the actual desired target law while remaining jointly realized with the algorithm initialization. No separate reference variable with only a matching marginal needs to be selected.

\subsection{\texorpdfstring{The retained noise marginal}{The retained noise marginal}}


The coupling measure has mass one. Projecting its product with \(Q\) to the noise factor therefore returns \(Q\). The sfinite hypothesis supplies the product marginal theorem's measure-theoretic condition.



For a measurable source set \(A\), its product preimage has mass \(c(E\times E)Q(A)=Q(A)\). Hence the noise projection preserves the source probability exactly. Every noise-defined random object, including a complete path, retains its original law after pullback. This distributional preservation is independent of the input coupling and its quadratic cost.

\subsection{\texorpdfstring{Independence of initial pair and noise}{Independence of initial pair and noise}}


The two coordinate projections of a product probability measure are independent. Here those coordinates are the full coupled initial pair and the fresh noise sample. Independence of either individual initial state from any noise-derived random object follows by composition.



Product factorization says that the full initial pair and the full noise source are independent. Composing with either state projection preserves independence, as does composing the source with a measurable path or innovation map. The conclusion therefore covers the whole innovation sequence at once, which is stronger than separately checking independence from each scalar innovation.

\subsection{\texorpdfstring{The joint initial-\allowbreak{}noise law}{The joint initial-noise law}}


Map \(((x,y),s)\) to \((x,s)\). Product pushforward and the first marginal identity of \(c\) give \(\mu\otimes Q\). Thus removing the reference coordinate leaves the same algorithm input law for every initial coupling.



For measurable sets \(A\) in the initial state and \(B\) in the source, the extended probability of \(x\in A,s\in B\) is \(c(A\times E)Q(B)=\mu(A)Q(B)\). This identifies the joint law as \(\mu\otimes Q\). The reference coordinate disappears from that law, explaining why the algorithm output can be common to all initial couplings.

\subsection{\texorpdfstring{Pushforward under the product observation}{Pushforward under the product observation}}


Factor the map through \(((x,y),s)\mapsto(x,s)\), whose law is \(\mu\otimes Q\). Pushforward composition proves that every measurable output depending only on algorithm initialization and noise has the coupling-independent law \(f_*(\mu\otimes Q)\).



Any measurable output using only \((x,s)\) factors through the projection whose law was just identified. Its pushforward is therefore obtained by applying the same output map to \(\mu\otimes Q\). This proves coupling independence of the output distribution before any numerical recursion is specified. The reference may still depend on \(y\), allowing a useful synchronous comparison.

\subsection{\texorpdfstring{The numerical output law}{The numerical output law}}


Apply the pushforward law under the independent product extension to the measurable recursion output. Its canonical pushforward is exactly the output law of a noise-driven chain, so the same output law is realized on every initial-coupling extension. This equality is derived rather than made a field of the coupling.



Take the output map to be the finite noise-driven recursion. Its measurability was proved independently of all coupling choices. Applying the preceding input-law identity gives exactly its canonical algorithm law. Thus an analytic estimate established on one initial-coupling extension always concerns the same numerical marginal, which is crucial when taking an infimum over initial coupling costs.

\subsection{\texorpdfstring{The innovation law after extension}{The innovation law after extension}}


First project the extension to its product-path noise law, then evaluate coordinate \(n\). The marginal of \(\bigotimes_k\eta\) at that coordinate is \(\eta\), giving the requested noise law on the extension.



If the source is an iid path product, evaluation at time \(n\) the stated pushforward law \(\eta\). The extension's source projection preserves that complete product law, so the pulled-back evaluation has the same marginal. This proves the innovation law for the concrete random functions used on the extended space, rather than for an unrelated auxiliary Gaussian variable.

\subsection{\texorpdfstring{Independence of the extended innovations}{Independence of the extended innovations}}


The iid path has independent coordinate evaluations. Pull this independence back along the extension's measure-preserving noise projection. The resulting innovation sequence remains jointly independent on the coupled initialization space.



Independence is a statement about the joint law of the innovation family. The noise projection preserves that entire law, not merely the individual coordinate distributions. Pulling back an independent path family therefore preserves all finite product factorizations. The innovations remain jointly independent on the same space that carries the initial coupling.

\subsection{\texorpdfstring{Independence of initialization and innovations}{Independence of initialization and innovations}}


The algorithm's initial state is a projection of the initial pair. Compose the pair/\allowbreak{}noise independence theorem with this state projection. This proves independence from the entire noise sample, not just from each individual noise vector.



Initialization depends only on the first coordinate of the coupling factor. The complete noise path depends only on the independent source factor. Composing the product-coordinate independence with these maps proves independence of initialization and whole path. This is the exact hypothesis needed for the canonical recursion law theorem and its fresh-noise consequences.

\subsection{\texorpdfstring{The initial displacement cost}{The initial displacement cost}}


The difference \(x-y\) depends only on the initial pair. Its \(L^2\) extended norm under \(c\otimes Q\) equals the same norm under \(c\), because the noise factor has mass one. The latter is exactly joint.



The initial difference \(x-y\) is constant in the added source variable. Integrating its squared norm over \(c\otimes Q\) multiplies the original coupling integral by \(Q(S)=1\). Taking square roots therefore preserves the initial quadratic cost exactly. Extending the probability space introduces no loss in the initial-error coefficient.

\subsection{\texorpdfstring{Square-\allowbreak{}integrability of the initial difference}{Square-integrability of the initial difference}}


The state-difference function is measurable. Its extended \(L^2\) norm equals the finite initial coupling cost by the previous theorem. These two facts prove square-integrability for the difference on the extension.



Measurability of coordinate subtraction gives measurability of the initial difference. The preserved cost identity gives a finite second moment when the coupling cost is finite. Together these facts place that difference in \(L^2\) on the extended space. Finiteness is therefore derived from the chosen transport coupling rather than imposed as a separate random-variable assumption.

\subsection{\texorpdfstring{Square-\allowbreak{}integrability of the reference initial state}{Square-integrability of the reference initial state}}


The reference coordinate the stated pushforward law \(\pi\). Pull its supplied \(L^2(\pi)\) identity moment back through this law to get square integrability on the extension. This establishes the reference initialization without choosing a separate realization.



The reference-coordinate law is \(\pi\), so changing variables gives \(\mathbb E\|Y_0\|^2=\int\|y\|^2\pi(dy)\). The assumed target second moment is finite. This provides an actual square-integrable reference variable on the coupling extension, preserving both its target law and its prescribed dependence with the first coordinate.

\subsection{\texorpdfstring{Square-\allowbreak{}integrability of the numerical initial state}{Square-integrability of the numerical initial state}}


Write \(x=(x-y)+y\). The difference is in \(L^2\) by finite coupling cost and the reference coordinate is in \(L^2\) by its law. Closure of \(L^2\) under addition gives square integrability of the initial algorithm state; no independent moment premise on \(\mu\) is introduced.



The identity \(X_0=(X_0-Y_0)+Y_0\) expresses initialization as a sum of two square-integrable variables. The triangle inequality gives \(\|X_0\|_2\le\|X_0-Y_0\|_2+\|Y_0\|_2\). Thus finite transport to a finite-second-moment target supplies the initial moment automatically. No extra initial second-moment hypothesis is needed in this coupling-based interface.

\subsection{\texorpdfstring{The numerical initial class}{The numerical initial class}}


Use square-integrability of the initial projection on the product space to convert the first state coordinate into an actual \(L^2\) class. The class is represented almost everywhere by \(((x,y),s)\mapsto x\), preserving its concrete relation to the coupling coordinates.



Once the first coordinate is proved square integrable, it determines an \(L^2\) state. Its random-function representative agrees almost surely with the coordinate \(x\). This construction fixes its joint realization with the reference coordinate and noise; it is not merely an abstract variable having law \(\mu\). The ensuing analytic recursion can therefore use the actual preserved coupling cost.

\subsection{\texorpdfstring{The reference initial class}{The reference initial class}}


Convert the square-integrable second state coordinate into an \(L^2\) class. Its representative equals \(((x,y),s)\mapsto y\) almost everywhere, with law \(\pi\), and it shares the same product probability space as the algorithm initialization.



The second-coordinate \(L^2\) state is formed on exactly the same extended source as the first. Almost surely its representative equals \(y\), whose law is \(\pi\). It can be inserted as the initial condition of the reference integral equation, while the first coordinate initializes the Euler chain and both use the independent source.

\subsection{\texorpdfstring{The law of the numerical initial class}{The law of the numerical initial class}}


The representative of the numerical initial class agrees almost everywhere with the first coordinate of the initial coupling. Its law is therefore the retained first marginal \(\mu\), because almost-everywhere changes leave measurable pushforwards unchanged.



Almost-sure equality implies equal probabilities of membership in every measurable state set. Therefore replacing the chosen square-integrable representative by the raw first coordinate leaves its law unchanged. The previously computed first marginal is \(\mu\), so the analytic initial state has precisely the prescribed algorithm initialization distribution.

\subsection{\texorpdfstring{The law of the reference initial class}{The law of the reference initial class}}


Replace the selected \(L^2\) representative by the second state projection almost everywhere. Its already proved marginal is \(\pi\), hence the class used as reference initialization has exactly that law.



The same representative argument applies to the second coordinate. Its target law was established before constructing its \(L^2\) state, so that construction cannot alter the reference marginal. Consequently stationarity results requiring target-distributed initialization apply to the actual state used in the coupling argument, with no further distributional identification left to assume.

\subsection{\texorpdfstring{The class distance equals the coupling cost}{The class distance equals the coupling cost}}


The representative of the class difference equals \(x-y\) almost everywhere. The norm of that \(L^2\) difference is the real value of its extended norm; the equality of product-space displacement energy and initial coupling cost identifies it with \((cc)\). The finite-cost premise justifies this real conversion.



The difference of the two constructed states represents \(x-y\) almost surely. Its squared \(L^2\) norm is therefore the original finite quadratic coupling integral, and its norm is the coupling's root mean-square cost. This exact equality identifies the initial analytic error with the quantity over which Wasserstein distance takes an infimum.

\section{\texorpdfstring{A reference for every initial coupling with one output law}{A reference for every initial coupling with one output law}}
\label{proof-module-UnadjustedLangevin-ProductReference-lean}


This construction turns the analytic reference estimate into an exact initial Wasserstein term by realizing each finite initial coupling on a product with one fixed noise source.



For each initial coupling, take the product with the prescribed independent noise source and run the actual pointwise update. The output law depends only on the first initial marginal and the innovation law, so it is the same across all couplings. The reference and the numerical chain then share a sample space while their initial \(L^2\) distance equals that coupling's cost. Uniform constants allow the final comparison to pass to the infimum over initial couplings.

\subsection{\texorpdfstring{The pointwise numerical update}{The pointwise numerical update}}


Define the ordinary Euler update \(F_h(x,z)=x-h\nabla U(x)+\sqrt{2h}z\) in an arbitrary Hilbert-space presentation. It is the pointwise function whose lifted composition gives the \(L^2\) numerical chain.



The update uses one current state and one innovation vector. Its additive Gaussian coefficient is \(\sqrt{2h}\), corresponding to covariance \(2hI\) when the innovation is standard. Defining this ordinary map before lifting it to \(L^2\) keeps the actual algorithm explicit and permits its law to be computed by measurable pushforwards.

\subsection{\texorpdfstring{Measurability of the update}{Measurability of the update}}


The Lipschitz gradient is continuous. Compose it with the first coordinate projection and combine the remaining continuous scalar and vector operations. The uncurried map \((x,z)\mapsto F_h(x,z)\) is therefore measurable.



Lipschitzness of the gradient gives continuity of \(x\mapsto\nabla U(x)\). Both coordinate projections on the state–innovation product are continuous, and fixed scalar combinations preserve continuity. Thus the entire update is measurable for any fixed step. No coupling or reference hypothesis enters this elementary regularity fact.

\subsection{\texorpdfstring{Square-\allowbreak{}integrability of the product noise}{Square-integrability of the product noise}}


The noise projection from \(c\otimes Q\) to \(Q\) is measure preserving because the coupling is a probability. Pull back the supplied \(L^2(Q)\) noise function along this projection. Its square integrability is unchanged.



For a source innovation \(W_n\), the extended variable is \(W_n(s)\). Its squared norm integral under \(c\otimes Q\) equals \(c(E\times E)\int\|W_n\|^2dQ\), hence the original finite second moment. The added initial coordinates therefore do not alter innovation integrability, just as they do not alter its distribution.

\subsection{\texorpdfstring{The product-\allowbreak{}space noise class}{The product-space noise class}}


Use the preceding square-integrability proof to form the class of \(((x,y),s)\mapsto W_n(s)\) on the concrete extension. Both the numerical chain and reference use this same \(L^2\) innovation.



The lifted innovation is the actual square-integrable state represented by the pulled-back source function. It is fixed before selecting any reference trajectory, so both numerical and reference constructions can be required to use precisely this innovation. This choice is what later makes the algorithm's output law independent of the second initial marginal and coupling dependence.

\subsection{\texorpdfstring{Its representative identity}{Its representative identity}}


The selected noise representative equals the pulled-back source innovation \(W_n(s)\) almost everywhere. Thus its arithmetic inside the square-integrable update agrees with the ordinary pointwise Gaussian update, and its coordinate laws are those of the original innovation.



Almost surely, evaluating the constructed innovation gives exactly the original source innovation evaluated at \(s\). Consequently its scalar arithmetic in the square-integrable Euler update agrees with the raw pointwise formula. These representative equalities also permit Gaussian coordinate laws and path independence to be transferred to the exact random variables used by the numerical chain.

\subsection{\texorpdfstring{The product numerical output law}{The product numerical output law}}


Push \(\mu\otimes Q\) through the pointwise noise-driven chain starting at \(x\) with source path \(W(s)\). This definition omits the reference coupling entirely, making explicit the common algorithm output law for all coupling realizations.



The canonical output depends on the initial law \(\mu\), fixed source probability \(Q\), and source innovation path \(W\). It is computed on \(\mu\otimes Q\), so no reference initialization appears in its definition. Every coupling extension will reproduce this same first-coordinate/\allowbreak{}source input law, establishing one common numerical output for the later transport infimum.

\subsection{\texorpdfstring{The representative of the product chain}{The representative of the product chain}}


Induct through the lifted chain and use the representative identities for initial state, product noise, lifted drift, addition, and scalar multiplication. At each fixed index the result agrees almost everywhere with the ordinary recursion driven by \(x\) and \(W(s)\).



At each successor, replace the lifted gradient, state, and innovation by their pointwise representatives on a common full-measure set. The resulting equation is exactly the ordinary Euler recursion. Induction proves agreement for every fixed finite index. This permits its law to be computed from the raw coordinates \((x,s)\), which have already been identified independently of the coupling.

\subsection{\texorpdfstring{The output law of the product chain}{The output law of the product chain}}


Replace the \(L^2\) chain by its almost-everywhere pointwise recursion. That recursion depends only on the algorithm's initial coordinate and noise sample, so the pushforward law under the independent product extension gives its coupling-independent law. The result is exactly the numerical law on the product construction.



The pointwise recursion ignores the reference coordinate \(y\). Under the extension, the law of the remaining pair \((x,s)\) is always \(\mu\otimes Q\). Pushforward composition therefore gives the same canonical output law for every finite initial coupling. Almost-sure replacement transfers this identity to the analytic chain whose errors are estimated in \(L^2\).

\subsection{\texorpdfstring{The product probability measure}{The product probability measure}}


A coupling of the initial probability with the reference measure is itself a probability, and its product with the fixed probability noise source is a probability. This normalization supplies normalization on every extension separately.



Every initial coupling has mass one by either marginal identity, and the source probability has mass one. Their product is consequently normalized. This fact is established separately on each extension and allows the same probability-space estimates to be reused uniformly over all couplings. It does not presume finite coupling cost or any invariant law of the numerical update.

\subsection{\texorpdfstring{A reference on the independent product space}{A reference on the independent product space}}


For every finite-cost initial coupling, the structure supplies a concrete l2reference on that coupling's product extension, with its initial state matching the second coordinate. Algorithm output laws and initial costs are proved outside the structure; no convergence estimate is a field.



The reference requirement is constructive: on each finite-cost extension, provide continuous joined segments with the integral equation, matching Euler noise and target-law moment bounds. The numerical law and initial cost are already consequences of the product construction. Thus the reference interface records analytic data, while the convergence estimate is proved from those data rather than included among them.

\subsection{\texorpdfstring{Finite output transport distance}{Finite output transport distance}}


Finite initial Wasserstein distance supplies at least one finite-cost initial coupling. On its concrete extension, the numerical and reference variables belong to \(L^2\), giving a finite-cost output coupling. Their established laws identify its marginals with the common output law and \(\pi\), hence the output Wasserstein distance is finite.



Finite initial transport guarantees the existence of at least one finite-cost initial coupling. Its analytic numerical and reference variables are in \(L^2\), so their difference has finite norm and their joint law has finite quadratic cost. Its first marginal is the common numerical output and its second is \(\pi\). Therefore output transport is finite before the real-valued numerical inequality is stated.

\subsection{\texorpdfstring{The small-\allowbreak{}step bound from a product reference}{The small-step bound from a product reference}}


For each finite-cost initial coupling \(c\), apply the analytic \(L^2\) Dalalyan theorem on \(c\otimes Q\). The initial class-distance theorem identifies its initial error exactly with the real coupling cost, and the common output-law theorem identifies its numerical law.



The resulting bound holds for every finite \(c\), with coefficient \((1-mh)^n\ge0\). Infimizing their costs gives the genuine \(W_2(\mu,\pi)\), without selecting an optimal coupling. Output finiteness is proved first so the final real-valued inequality cannot conceal infinity.



For each initial coupling cost \(C\), synchronous analysis gives output distance at most \((1-mh)^nC+(33/20)(L/m)\sqrt{hd}\). The multiplicative coefficient is nonnegative and the other quantities do not depend on the chosen coupling. Taking the cost infimum therefore replaces \(C\) by initial \(W_2\). Finiteness was established separately, so this operation applies to genuine finite real distances.

\subsection{\texorpdfstring{The large-\allowbreak{}step bound from a product reference}{The large-step bound from a product reference}}


Repeat the couplingwise analytic argument for the stable large-step branch, whose contraction factor is \(Lh-1\in[0,1)\). The geometric error sum has denominator \(2-Lh>0\). The finite-coupling infimum theorem replaces the initial class distance by the actual initial Wasserstein distance and yields the displayed large-step bias term.



In the second stable branch, the deterministic gradient-step factor is \(q=Lh-1\). Its geometric denominator is \(1-q=2-Lh>0\), producing floor \((33/20)Lh\sqrt{hd}/(2-Lh)\). The same product-law identity and coupling infimum apply because the constructed reference and common numerical output are unchanged. Only the scalar contraction branch and its denominator differ.

\section{\texorpdfstring{Gaussian and Brownian noise moments}{Gaussian and Brownian noise moments}}
\label{proof-module-UnadjustedLangevin-GaussianMoments-lean}


Fix a natural-number dimension \(d\in\mathbb N\), write \(E_d=\mathbb R^d\) with its Euclidean norm, and let \(P\) be a measure on \(\Omega\). A law statement \(X_i\sim\mathcal N(0,v)\) below means that the coordinate is almost everywhere measurable and has the stated pushforward measure. We prove the required \(L^2\) norms from these laws. Coordinate independence is unnecessary for second moments because \(\|X\|^2=\sum_iX_i^2\) has no cross terms.



For Brownian processes the input is the scalar Brownian increment laws, whose marginal and increment laws supply the Gaussian hypotheses. These facts establish the noise moments; existence and stationarity of a coupled Langevin reference require separate inputs.



The \(L^2\) norms and expectations below concern the actual measure \(P\), with square-integrability giving square integrability before passing to a quotient class. Laws are pushforward identities; independence is used only for cross-term cancellation, while marginal Gaussian second moments supply the dimension factor.



All second-moment calculations use the Euclidean identity \(\|v\|^2=\sum_{i=1}^d v_i^2\). Only the individual coordinate variances are needed for this identity: independence between coordinates is not required to sum their squares. Independence becomes a separate issue when identifying the joint Gaussian innovation law. For Brownian increments, replacing the starting time by \(a\) changes no second-moment constant because the variance is the elapsed time.

\subsection{\texorpdfstring{Second moment of a centered Gaussian}{Second moment of a centered Gaussian}}


For nonnegative variance \(v\), the centered Gaussian law has second moment

\[\int_{\mathbb R}x^2\,d\mathcal N(0,v)(x)=v.\]

The known Gaussian variance is \(v\). Expand variance as \(\int(x-\mathbb E[x])^2\), and insert the Gaussian mean \(\mathbb E[x]=0\).



Simplify \(x-0=x\). The variance integral is now exactly the required second moment.



The variance parameter is the second argument of \(\mathcal N(0,v)\), not its standard deviation. Centering therefore makes the variance integral exactly the squared-variable integral. At \(v=0\), the Gaussian is a point mass at zero and the same identity gives zero, so no positive-variance exception is needed.

\subsection{\texorpdfstring{Gaussian coordinates give a square-\allowbreak{}integrable vector}{Gaussian coordinates give a square-integrable vector}}


Let \(X:\Omega\to E_d\), and suppose every coordinate satisfies \(X_i\sim\mathcal N(0,v)\) under \(P\), with \(v\ge0\). Then \(X\in L^2(P;E_d)\).



In finite-dimensional Euclidean space, square integrability is equivalent to square integrability of each coordinate. Each coordinate belongs to \(L^2\) because its Gaussian law has a finite second moment; transfer that membership along the equality of laws.



A finite sum of integrable nonnegative coordinate squares is integrable. Conversely each coordinate square is bounded by the squared Euclidean norm. These two inequalities explain the equivalence between vector and coordinate \(L^2\) membership in finite dimension. Coordinate independence is not used for this membership statement.

\subsection{\texorpdfstring{Sum the coordinate moments}{Sum the coordinate moments}}


Under the same coordinate Gaussian laws,

\[\int_\Omega\|X(\omega)\|^2\,dP(\omega)=dv.\]

For each coordinate \(i\), retain the proved square integrability \(X_i\in L^2(P)\). It will justify integrating the finite sum of squared coordinates.



Change variables using the law of \(X_i\), then apply the scalar Gaussian moment:

\[\int_\Omega X_i^2\,dP=\int_{\mathbb R}x^2\,d\mathcal N(0,v)=v.\]

Expand the squared Euclidean norm and interchange the finite sum and integral:

\[\int_\Omega\|X\|^2\,dP=\sum_{i=0}^{d-1}\int_\Omega X_i^2\,dP.\]

Every summand is integrable by its \(L^2\) membership.



Each of the \(d\) summands equals \(v\); hence their sum is \(dv\).



No cross terms occur in the identity for the squared Euclidean norm. Thus possible correlations between distinct coordinates do not change this second moment; their individual variances determine it. The finite sum may be exchanged with the integral because each square has already been proved integrable.

\subsection{\texorpdfstring{The Gaussian vector as an L\ensuremath{{}^{2}} class}{The Gaussian vector as an L² class}}


Let \([X]\) be the \(L^2\) equivalence class constructed from \(X\), whose coordinates have laws \(\mathcal N(0,v)\). Then

\[\|[X]\|_{L^2}^2=dv.\]

Express the squared \(L^2\) norm as the self-inner-product, then use the definition of the \(L^2\) inner product as an integral.



The chosen representative of \([X]\) equals \(X\) almost everywhere. Therefore the integrand \(\langle[X](\omega),[X](\omega)\rangle\) agrees almost everywhere with \(\|X(\omega)\|^2\), so the two integrals agree.



Substitute the integral identity and the vector second-moment calculation \(\int\|X\|^2\,dP=dv\).



An \(L^2\) class identifies random vectors that agree almost everywhere. The norm is independent of the chosen representative, so the second-moment calculation for the original random vector computes the class norm as well. The inner-product formula makes that identification explicit without selecting pointwise values on exceptional sets.

\subsection{\texorpdfstring{The square-\allowbreak{}root Gaussian norm}{The square-root Gaussian norm}}


For the same Gaussian coordinate laws,

\[\|[X]\|_{L^2}=\sqrt{dv}.\]

This includes \(d=0\) and \(v=0\).



Replace \(dv\) by \(\|[X]\|^2\). Since norms are nonnegative, \(\sqrt{\|[X]\|^2}=\|[X]\|\).



The argument takes the unique nonnegative square root of \(dv\ge0\). A sign choice is therefore unnecessary, and the formula remains valid at both degenerate endpoints. This is the exact root mean square size used when placing Gaussian vectors inside the common \(L^2\) reference space.

\subsection{\texorpdfstring{Include the Langevin diffusion factor}{Include the Langevin diffusion factor}}


If every coordinate of \(X\) the stated pushforward law \(\mathcal N(0,v)\), then the Langevin-scaled vector satisfies

\[\|\sqrt2[X]\|_{L^2}=\sqrt{2d}\sqrt v.\]

Norm homogeneity gives \(\|\sqrt2[X]\|=\sqrt2\|[X]\|\), because \(\sqrt2\ge0\). Insert \(\|[X]\|=\sqrt{dv}\), then split the square roots of the nonnegative products.



Rearrange \(\sqrt2\sqrt d\sqrt v\) into \(\sqrt{2d}\sqrt v\).



Multiplying a random vector by \(\sqrt2\) doubles its second moment. Taking a square root changes \(dv\) into \(\sqrt{2dv}\), which is equal to the displayed product because the factors are nonnegative. This factor is the diffusion normalization producing a unit Laplacian in the Langevin generator.

\subsection{\texorpdfstring{Assemble Brownian coordinates}{Assemble Brownian coordinates}}


For scalar processes \(W_i:[0,\infty)\times\Omega\to\mathbb R\), define the Euclidean vector

\[W_t(\omega):=(W_0(t,\omega),\ldots,W_{d-1}(t,\omega))\in E_d.\]

Collect the coordinates at the fixed time \(t\) and outcome \(\omega\), using the finite-dimensional \(\ell^2\) norm.



Assembling the coordinates changes the codomain from a finite family of scalar variables to \(\mathbb R^d\). Evaluating coordinate \(i\) of the resulting vector simply recovers the original scalar variable at that index. This observation lets every coordinate law and measurability property transfer without a new probabilistic construction.

\subsection{\texorpdfstring{Brownian marginals are Gaussian}{Brownian marginals are Gaussian}}


Assume each scalar \(W_i\) satisfies \(\text{the Brownian increment laws}\) under \(P\). For each \(t\ge0\) and coordinate \(i\),

\[W_t^{(i)}\sim\mathcal N(0,t).\]

Apply the scalar Brownian marginal-law theorem at time \(t\). The \(i\)-th coordinate of the assembled vector is by definition the original scalar process.



For a pre-Brownian coordinate, the distribution at time \(t\) is already specified by its centered Gaussian marginal of variance \(t\). Selecting the coordinate of the assembled vector therefore gives the claimed law directly. This statement alone makes no assertion about independence between distinct spatial coordinates.

\subsection{\texorpdfstring{Brownian vectors belong to L\ensuremath{{}^{2}}}{Brownian vectors belong to L²}}


Under the scalar pre-Brownian assumptions, for every \(t\ge0\),

\[W_t\in L^2(P;E_d).\]

Apply Gaussian vector square integrability with variance \(v=t\), using the just-proved marginal laws of all coordinates.



Each coordinate has finite second moment \(t\). Summing the \(d\) coordinate integrals bounds the vector square and gives finite \(L^2\) membership. The conclusion is a fixed-time moment statement; continuity of the process as an \(L^2\) curve is established separately from increment laws.

\subsection{\texorpdfstring{The exact Brownian noise norm}{The exact Brownian noise norm}}


For scalar pre-Brownian coordinates and \(t\ge0\),

\[\|\sqrt2[W_t]\|_{L^2}=\sqrt{2d}\sqrt t.\]

Specialize the scaled Gaussian vector norm to the Brownian vector, whose coordinate variance is \(t\). This derives the square-root noise scale from the Brownian laws.



The squared norm is exactly \(2dt\), not merely bounded by it. Consequently the reference noise estimate uses the optimal moment scale supplied by the Gaussian law. At \(t=0\), the noise class has zero norm and hence equals zero in \(L^2\), agreeing with the integral-equation initialization.

\subsection{\texorpdfstring{Define an increment from any starting time}{Define an increment from any starting time}}


For \(a,t\ge0\), define the vector increment over the interval \([a,a+t]\) by

\[\Delta_aW_t(\omega):=W_{a+t}(\omega)-W_a(\omega).\]

Subtract the two Euclidean vectors at the interval endpoints. Each coordinate is the corresponding scalar Brownian increment.



The elapsed time parameter is kept separate from the starting time. This permits every reference step to be described on its local interval \([0,h]\), while its underlying Brownian input is taken from the global interval \([a,a+h]\). Vector subtraction is performed coordinatewise, so scalar increment laws apply immediately.

\subsection{\texorpdfstring{Increment variance is elapsed time}{Increment variance is elapsed time}}


Under the scalar pre-Brownian assumptions, for every \(a,t\ge0\) and every coordinate \(i\),

\[W_i(a+t)-W_i(a)\sim\mathcal N(0,t).\]

Expand the vector increment to its scalar coordinate. Apply the Brownian difference-law theorem at times \(a+t\) and \(a\); its variance is the nonnegative distance between those times.



It remains to identify that variance. As real numbers,

\[|(a+t)-a|=|t|=t,\]

because \(t\ge0\). Thus the increment law is exactly \(\mathcal N(0,t)\).



The distance of the two nonnegative time points is \(t\), because their difference is \((a+t)-a=t\ge0\). Thus the variance depends on elapsed time alone. This invariance of the increment law under a change of starting time is the probabilistic fact needed to give all reference segments the same error constant.

\subsection{\texorpdfstring{Increments are square integrable}{Increments are square integrable}}


For pre-Brownian coordinates and any \(a,t\ge0\),

\[\Delta_aW_t\in L^2(P;E_d).\]

Each increment coordinate has the centered Gaussian law of variance \(t\), so the Gaussian vector integrability theorem applies.



Apply the coordinate-to-vector moment criterion to the increment rather than to the two endpoints separately. Each increment coordinate has second moment equal to elapsed time, so the argument gives exactly the moment scale needed later. It does not require a bound growing with the absolute starting time \(a\).

\subsection{\texorpdfstring{The same noise norm on every reference segment}{The same noise norm on every reference segment}}


For scalar pre-Brownian coordinates and every \(a,t\ge0\),

\[\boxed{\|\sqrt2[W_{a+t}-W_a]\|_{L^2}=\sqrt{2d}\sqrt t.}\]

Apply the scaled Gaussian vector norm to the increment coordinate laws. The answer depends only on elapsed time \(t\), so every reference interval has the same moment bound.



The expression is independent of \(a\), which is crucial when applying the local-defect estimate at step indices \(k=0,1,\ldots\). Each segment has the same \(\sqrt t\) modulus in \(L^2\), so the resulting defect bound is uniform over the entire discrete trajectory. This uniformity supports one fixed forcing constant in the recurrence.

\section{\texorpdfstring{Coordinate Gaussian moments and L\ensuremath{{}^{2}} Brownian continuity}{Coordinate Gaussian moments and L² Brownian continuity}}
\label{proof-module-UnadjustedLangevin-GaussianCoordinates-lean}


Coordinatewise Gaussian laws suffice for the second moment \(\mathbb E\|X\|^2=dv\). This module carefully obtains Brownian \(L^2\) continuity and increment norms without assuming independence between spatial coordinates.



For coordinate variance \(t\), the Euclidean second moment is \(dt\), and multiplying the vector by \(\sqrt2\) gives squared \(L^2\) norm \(2dt\). For Brownian coordinates, the same argument applied to a difference uses variance \(|t-s|\), yielding \(\|\sqrt2(W_t-W_s)\|_2=\sqrt{2d}\sqrt{|t-s|}\). This quantitative increment formula proves continuity in \(L^2\) and supplies the exact noise scale on every shifted reference segment.

\subsection{\texorpdfstring{Square-\allowbreak{}integrability from Gaussian coordinate laws}{Square-integrability from Gaussian coordinate laws}}


Each centered Gaussian coordinate has a finite second moment. In finite Euclidean dimension, \(\|X\|^2=\sum_iX_i^2\), so the finite sum is integrable. Coordinatewise measurability also assembles to vector measurability, giving \(X\in L^2\). Coordinate independence is unnecessary for this moment statement.



The finite-dimensional norm square is a finite sum of coordinate squares. Each coordinate has a finite Gaussian second moment, proving vector square-integrability without requiring independence for this particular conclusion.

\subsection{\texorpdfstring{The Euclidean Gaussian second moment}{The Euclidean Gaussian second moment}}


Use \(\|X\|^2=\sum_iX_i^2\) and finite-sum integral linearity. The coordinate law \(X_i\sim\mathcal N(0,v)\) gives \(\mathbb EX_i^2=v\), hence \(\mathbb E\|X\|^2=|\iota|v\). The dimension factor is exactly the cardinality of the finite coordinate set.



Summing the coordinate second moments gives \(dv\). No mixed coordinate products occur, so the same formula remains valid for correlated coordinates having these marginal variances.

\subsection{\texorpdfstring{The squared norm of the Gaussian class}{The squared norm of the Gaussian class}}


For the given \(L^2\) class, its squared norm is the integral of the squared norm of its representative. Apply the preceding Gaussian coordinate moment formula to obtain \(\|X\|_2^2=|\iota|v\).



The class representative has the same squared-norm integral as the original vector. Thus the \(L^2\) norm square is exactly the already computed total coordinate variance.

\subsection{\texorpdfstring{The norm of the Gaussian class}{The norm of the Gaussian class}}


Take nonnegative square roots in \(\|X\|_2^2=|\iota|v\). Norm nonnegativity selects the positive root, giving \(\|X\|_2=\sqrt{|\iota|v}\), including zero variance and zero dimension.



Taking nonnegative square roots turns the second-moment identity into the root mean square norm. The formula covers zero variance and the empty coordinate set.

\subsection{\texorpdfstring{Coordinate laws after forming the class}{Coordinate laws after forming the class}}


A representative of the square-integrable class \([X]\) equals \(X\) almost everywhere. Coordinate evaluation preserves that equality. Since pushforward laws do not change under almost-everywhere equality, the coordinate Gaussian laws transfer to the class representative.



The coordinate of the selected class representative agrees almost everywhere with the original coordinate variable. Their pushforward laws therefore coincide with the stated scalar Gaussian law.

\subsection{\texorpdfstring{The norm is independent of the chosen representative}{The norm is independent of the chosen representative}}


The class \([X]\) has the transferred coordinate laws and the same total second moment as \(X\). Hence \(\|[X]\|_2^2=|\iota|v\). Both sides are nonnegative, so taking square roots gives the exact Gaussian class norm.



Changing representatives on a null set leaves the squared norm integral unchanged. The class norm is consequently the same \(\sqrt{dv}\) computed from the original coordinates.

\subsection{\texorpdfstring{The Langevin-\allowbreak{}scaled Gaussian norm}{The Langevin-scaled Gaussian norm}}


The scalar-norm identity gives \(\|\sqrt2[X]\|_2=\sqrt2\sqrt{|\iota|v}\). Nonnegativity permits the square-root product rules, producing \(\sqrt{2|\iota|}\sqrt v\). This is the precise Langevin noise constant used in the local error estimate.



Norm homogeneity multiplies the Gaussian root mean square norm by the absolute scale factor. For the Langevin scale \(\sqrt2\), the result is \(\sqrt{2d}\sqrt v\).

\subsection{\texorpdfstring{Assembling Brownian coordinates}{Assembling Brownian coordinates}}


Assemble the scalar coordinate processes into \(B_t(\omega)=(W_i(t,\omega))_{i\in\iota}\). Finite Euclidean coordinates determine the vector; no new random source or independence is introduced by this definition.



The vector at time \(t\) is assembled from the spatial scalar processes at that same time. Coordinate evaluation recovers each scalar process directly.

\subsection{\texorpdfstring{The coordinate marginal laws}{The coordinate marginal laws}}


At a fixed time \(t\ge0\), each scalar pre-Brownian coordinate the stated pushforward law \(\mathcal N(0,t)\). Unfolding the assembled vector shows that its \(i\)-th coordinate is exactly that scalar variable, giving the required coordinate law.



Each coordinate process has Brownian marginal law \(\mathcal N(0,t)\). The vector construction changes only the codomain, so these fixed-time coordinate laws are unchanged.

\subsection{\texorpdfstring{A vector Brownian increment}{A vector Brownian increment}}


Define the vector increment \(\Delta_aB_t=B_{a+t}-B_a\) coordinate by coordinate. This records both the starting time and elapsed nonnegative time, matching the local noise input on an arbitrary Langevin segment.



Subtracting the vector at the starting time yields the vector of scalar increments. This local-time parameterization is what permits reference segments to be shifted along a global process.

\subsection{\texorpdfstring{The coordinate increment laws}{The coordinate increment laws}}


The scalar Brownian increment law depends only on elapsed time: \(W_i(a+t)-W_i(a)\sim\mathcal N(0,t)\). Simplify \((a+t)-a=t\) in the variance and unfold the vector increment to obtain its coordinate law.



The scalar increment law has variance equal to elapsed time, independent of starting time. Thus the same coordinate Gaussian moment argument applies to every segment.

\subsection{\texorpdfstring{Square-\allowbreak{}integrability of the vector increment}{Square-integrability of the vector increment}}


Apply the finite sum of the Gaussian coordinate second moments to the increment's coordinate laws of variance \(t\). The assembled Brownian vector increment is square integrable even without independence between different coordinate processes.



Every increment coordinate has finite Gaussian second moment. The finite sum criterion then establishes that the whole vector increment belongs to \(L^2\).

\subsection{\texorpdfstring{The scaled vector-\allowbreak{}increment norm}{The scaled vector-increment norm}}


Apply the exact scaled Gaussian-coordinate norm theorem with variance \(t\). It yields \(\|\sqrt2[\Delta_aB_t]\|_2=\sqrt{2|\iota|}\sqrt t\), establishing the noise radius from Brownian laws alone.



Its squared Langevin-scaled norm is \(2dt\). Taking a square root gives a uniform elapsed-time modulus, with no dependence on the segment's absolute starting time.

\subsection{\texorpdfstring{Brownian values in the square-\allowbreak{}integrable space}{Brownian values in the square-integrable space}}


Convert the assembled Brownian vector at time \(t\) to its \(L^2\) class using its Gaussian-coordinate square-integrability proof. This places the time-indexed process in the Hilbert space where the deterministic integral equation is solved.



The vector at each nonnegative time determines a square-integrable class. This turns the Brownian coordinate construction into a Hilbert-space curve suited to a Bochner integral equation.

\subsection{\texorpdfstring{The Brownian representative identity}{The Brownian representative identity}}


For each fixed time, a representative of \([B_t]\) equals the original Brownian vector \(B_t\) almost everywhere. This identifies its coordinate laws while allowing continuity to be studied in the square-mean norm. The exceptional set may depend on that fixed time, which does not affect any marginal or moment statement.



The representative identity permits fixed-time laws and coordinate moments to be evaluated using the original assembled Brownian vector. All equality claims are almost everywhere at the specified time.

\subsection{\texorpdfstring{The coordinate law of a class difference}{The coordinate law of a class difference}}


The difference of two selected representatives agrees almost everywhere with \(B_t-B_s\). Brownian increment laws, including the sign reversal when \(t<s\), give a centered Gaussian with variance \(|t-s|\); centered Gaussian symmetry handles the reversed order.



The difference of two class values agrees almost everywhere with the vector of scalar Brownian increments. Its coordinate law therefore has variance equal to the distance between the two times.

\subsection{\texorpdfstring{The exact norm of a Brownian class difference}{The exact norm of a Brownian class difference}}


Apply the Gaussian \(L^2\) norm formula to the difference class, whose coordinates have variance \(|t-s|\). Thus \(\|[B_t]-[B_s]\|_2=\sqrt{|\iota|\,|t-s|}\). This exact increment norm is stronger than the continuity statement derived from it.



The norm of a class difference is the square root of the sum of coordinate increment variances. This gives \(\sqrt d\sqrt{|t-s|}\) and immediately tends to zero as \(s\to t\).

\subsection{\texorpdfstring{Continuity in the square-\allowbreak{}mean norm}{Continuity in the square-mean norm}}


As \(s\to t\), \(|t-s|\to0\), hence the exact norm formula gives \(\|[B_t]-[B_s]\|_2\to0\). The construction uses continuity of the square root and the metric characterization of continuity. This needs finite-dimensional Brownian laws, with no sample-path continuity premise.



The increment norm formula gives a quantitative modulus of continuity in \(L^2\). This supplies the continuous additive input required by the Hilbert-space Langevin existence theorem.

\subsection{\texorpdfstring{The Langevin-\allowbreak{}scaled increment norm}{The Langevin-scaled increment norm}}


Set the two times in the difference-norm identity to \(a+t\) and \(a\). Their distance is \(t\), and scalar multiplication by \(\sqrt2\) gives \(\sqrt{2|\iota|}\sqrt t\). This is the class-valued version of the earlier increment estimate.



Multiplication of the increment by \(\sqrt2\) produces exactly the diffusion scale of the Langevin equation. Its \(L^2\) norm is \(\sqrt{2d}\sqrt t\) on a segment of duration \(t\).

\subsection{\texorpdfstring{The Brownian class at zero}{The Brownian class at zero}}


At time zero every coordinate has the degenerate centered Gaussian law of variance zero. Equivalently its \(L^2\) norm is zero by the moment formula. Norm-zero equivalence proves that the whole class is zero, irrespective of the chosen pointwise representative.



At time zero, each scalar Brownian coordinate has the point-mass law at zero. The assembled vector class has zero \(L^2\) norm and is consequently the zero class.

\subsection{\texorpdfstring{A real-\allowbreak{}time version of the Brownian curve}{A real-time version of the Brownian curve}}


Extend the nonnegative-time \(L^2\) the process value at time \(t\)o real times by \(t\mapsto[B_{\max(t,0)}]\), implemented with restriction to nonnegative time. It is constant zero for negative times and agrees with the original class at nonnegative times.



Extending the nonnegative-time curve to a real interval uses the designated nonnegative time associated with each argument. On the actual nonnegative horizon this agrees with the original Brownian class.

\subsection{\texorpdfstring{Continuity of the real-\allowbreak{}time curve}{Continuity of the real-time curve}}


The clipping map \(t\mapsto\max(t,0)\) into nonnegative reals is continuous. Compose it with the established \(L^2\)-continuous Brownian map to obtain continuity of the real-time extension, including at zero.



Continuity of the time conversion and of the nonnegative-time \(L^2\) curve gives continuity of the extended curve. This is the real-time domain used by the interval integral.

\subsection{\texorpdfstring{The real-\allowbreak{}time curve at zero}{The real-time curve at zero}}


At real time zero the clipping map returns nonnegative zero. The previously proved zero-time \(L^2\) identity therefore simplifies the real-time extension to zero.



At the origin, the converted time is zero, so the extended noise class also vanishes. This verifies the initialization condition in the integral equation.

\subsection{\texorpdfstring{Its scaled increment norm}{Its scaled increment norm}}


For \(a,t\ge0\), clipping leaves \(a\) and \(a+t\) unchanged. Rewrite the real-time extension at both endpoints as the nonnegative-time Brownian classes, then apply their scaled increment formula. This supplies \(\sqrt{2d}\sqrt t\) in the real interval-integral formulation.



Within the nonnegative segment, the real-time version has exactly the original Brownian increment. Its Langevin-scaled norm therefore has the same elapsed-time formula and constant.

\section{\texorpdfstring{Construct the L\ensuremath{{}^{2}} Langevin equation from Brownian-\allowbreak{}coordinate laws}{Construct the L² Langevin equation from Brownian-coordinate laws}}
\label{proof-module-UnadjustedLangevin-BrownianLangevin-lean}


Coordinate pre-Brownian laws produce the continuous \(L^2\) additive noise needed by the deterministic integral-equation existence theorem. The exact noise norm is \(\sqrt{2d}\sqrt t\).



The additive input is the \(L^2\) class of \(\sqrt2W_t\). The Gaussian increment formula makes this input continuous in time, with value zero at zero and norm \(\sqrt{2dt}\). The lifted gradient is Lipschitz, so the deterministic Banach-space integral-equation theorem applies directly. The result is a continuous curve of square-integrable variables satisfying the Langevin drift equation as an equality in \(L^2\).

\subsection{\texorpdfstring{The Langevin-\allowbreak{}scaled Brownian input}{The Langevin-scaled Brownian input}}


Define \(N_a(t)=\sqrt2([B_{a+t}]-[B_a])\) in \(L^2\), using the real-time Brownian extension. The subtraction fixes the input at zero at local time zero, and the \(\sqrt2\) factor gives the overdamped Langevin diffusion normalization.



The initial absolute time \(a\) fixes which Brownian increment is used during this segment. For \(t\ge0\), the input is \(\sqrt2(B_{a+t}-B_a)\), so it has covariance \(2tI_d\). Subtracting the value at \(a\) is essential: a local integral equation starts with zero additive input even when the driving Brownian path has already evolved before the segment begins.

\subsection{\texorpdfstring{Continuity of the Brownian input in square mean}{Continuity of the Brownian input in square mean}}


The real-time Brownian map is continuous in \(L^2\) from its Gaussian increment moments. Composition with \(t\mapsto a+t\), subtraction of the fixed class \([B_a]\), and scalar multiplication preserve continuity. Thus the additive input satisfies the deterministic Picard theorem's continuity requirement.



Mean-square continuity is the relevant regularity for an integral equation in \(L^2\). The increment identity \(\mathbb E\|B_t-B_s\|^2=d|t-s|\) makes the norm difference tend to zero. Translating time, subtracting one fixed vector and multiplying by \(\sqrt2\) preserve this property. Thus continuity is derived from the actual Gaussian process rather than imposed as an independent reference assumption.

\subsection{\texorpdfstring{The Brownian input at zero}{The Brownian input at zero}}


At local time zero both Brownian terms are \([B_a]\), so they cancel and their scalar multiple is zero. The defining formula gives the definition and uses additive identities; no exceptional-set statement is needed because this is equality in \(L^2\).



A local solution is prescribed to have initial state \(Y_0\). Evaluating its integral equation at zero gives \(Y(0)-Y_0=N_a(0)\), since the zero-length integral vanishes. The equality \(N_a(0)=0\) therefore enforces the intended initial condition exactly. This cancellation holds for the square-integrable random vectors themselves, regardless of which pointwise versions are used.

\subsection{\texorpdfstring{The exact input norm}{The exact input norm}}


For \(a,t\ge0\), apply the exact real-time Brownian increment moment formula. It gives \(\|N_a(t)\|_2=\sqrt{2d}\sqrt t\), with \(d=|\iota|\). The Gaussian laws prove this value; it is not an additional noise-bound hypothesis.



There are \(d\) scalar increment coordinates, each with variance \(t\) and zero mean. Therefore \(\mathbb E\|N_a(t)\|^2=2\sum_{j=1}^d t=2dt\). Taking the nonnegative square root gives the exact norm. Independence is not needed to sum coordinate squares, although the Brownian vector construction separately supplies its Gaussian joint law.

\subsection{\texorpdfstring{A short-\allowbreak{}interval square-\allowbreak{}mean Langevin solution}{A short-interval square-mean Langevin solution}}


Supply the continuous, zero-starting \(N_a\) to the short-interval fixed-point theorem in the square-integrable Hilbert space. It constructs \(Y_t-Y_0=-\int_0^t\mathcal G(\nabla U)(Y_s)ds+N_a(t)\) for \(Lh\le1/2\). Attach the exact noise norm theorem to the existential package. Neither the integral equation nor its noise estimate is assumed.



On a short interval, the map taking a trial path to \(Y_0-\int_0^t\nabla U(Y_s)ds+N_a(t)\) has contraction constant at most \(Lh\) in the supremum of \(L^2\) norms. The stated restriction makes the fixed-point construction immediate. The noise is only added, so its continuity is sufficient; it need not have an ordinary time derivative.

\subsection{\texorpdfstring{A solution on every finite horizon}{A solution on every finite horizon}}


Use the unrestricted finite-horizon \(L^2\) existence theorem with the same constructed Brownian input. The resulting solution is continuous in \(L^2\), starts at the prescribed random variable, and satisfies the actual gradient integral equation. The exact \(\sqrt{2d}\sqrt t\) noise formula is included as a proved property.



The global finite-horizon construction joins sufficiently short fixed-point intervals, using each endpoint as the next initial state. Lipschitzness of the score preserves square integrability and gives uniqueness by Grönwall. The resulting equation is therefore an actual constructed gradient dynamics on the required horizon. The Brownian norm formula remains unchanged because it concerns the supplied additive input, not the unknown solution.

\section{\texorpdfstring{Transfer the Gibbs score estimate to an L\ensuremath{{}^{2}} reference}{Transfer the Gibbs score estimate to an L² reference}}
\label{proof-module-UnadjustedLangevin-StationaryMoments-lean}


A reference variable \(Y\) whose law is \(\pi\) inherits \(\|\nabla U(Y)\|_{L^2}\le\sqrt{Ld}\). The assertion that its law is Gibbs is a displayed premise here.



If \(Y_\#P=\pi_U\), change of variables gives \(\|\nabla U(Y)\|_2^2=\int\|\nabla U\|^2d\pi_U\). The Gibbs score estimate bounds the last expression by \(Ld\), independently of the time at which the stationary variable is observed. Taking the nonnegative square root supplies exactly the reference drift bound \(\sqrt{Ld}\) used in the local-error calculation. This uses the stationary law, rather than a pathwise bound on the drift.

\subsection{\texorpdfstring{The reference drift norm from its Gibbs law}{The reference drift norm from its Gibbs law}}


The representative of The pointwise-composition operator equals \(\nabla U(Y)\) almost everywhere. Change variables using \(P\circ Y^{-1}=\pi_U\), so its squared \(L^2(P)\) norm is \(\int\|\nabla U\|^2d\pi_U\le Ld\). Nonnegativity of the norm and of \(Ld\) then gives \(\|\nabla U(Y)\|_2\le\sqrt{Ld}\). The law identity, rather than independence, is the input to this transfer.



The equality of laws transports the squared gradient integral from the sample space to Gibbs. The lifted drift's norm-square identity then gives \(\|\mathcal G[Y]\|_2^2\le Ld\), and taking square roots yields the reference bound.

\subsection{\texorpdfstring{The drift moment under strong convexity}{The drift moment under strong convexity}}


Use strong convexity to prove \(Z_U<\infty\), and apply the previous law-transfer theorem with that proof. The conclusion is the same \(\sqrt{Ld}\) bound for the lifted drift of any \(L^2\) variable with Gibbs law. There is no separate assumed drift-moment bound hidden in this specialization.



Positive strong convexity supplies the normalization and moment properties needed to apply the preceding stationary-law calculation. The resulting bound is independent of the particular representation or coupling of the stationary random vector.

\section{\texorpdfstring{Assemble continuous Brownian reference segments}{Assemble continuous Brownian reference segments}}
\label{proof-module-UnadjustedLangevin-BrownianReference-lean}


At knot \(kh\), start the next integral-equation segment from the preceding endpoint and drive it by \(\sqrt2(W_{kh+t}-W_{kh})\). Its value at \(t=h\) becomes the next knot, so adjacent segments agree at their shared endpoint. The normalized endpoint increment is the numerical innovation. Once the stationary law and the score moment are supplied, these constructed segments satisfy exactly the reference assumptions used by the local-defect convergence theorem.

\subsection{\texorpdfstring{A Brownian-\allowbreak{}driven segment from a given initial class}{A Brownian-driven segment from a given initial class}}


Choose the integral solution supplied by finite-horizon existence for the integral equation driven by the continuous Brownian class for the segment starting at absolute time \(nh\), with the supplied initial \(L^2\) variable. Classical choice selects data from a proved existence theorem; its defining equation and continuity are recovered by the specification theorem.



Fix a segment length \(h>0\), an index \(n\), and an initial random vector \(K\). The local reference is a solution of \(Y_t-K=-\int_0^t\nabla U(Y_s)ds+\sqrt2(B_{nh+t}-B_{nh})\). Existence has already been proved for this input. Selecting such a solution introduces no extra equation or numerical estimate; all its analytic properties are inherited from that existence result.

\subsection{\texorpdfstring{The defining properties of the segment}{The defining properties of the segment}}


Apply the choice specification to obtain continuity, the time-zero initial value, the gradient integral equation on \([0,h]\), and the exact noise norm. These four properties describe the selected segment independently of its noncomputable construction.



The four properties have distinct later uses. Continuity makes the drift time integral meaningful. The initial-value identity permits adjoining intervals. The integral equation exposes the local Euler drift defect. Finally the exact Brownian norm controls reference increments quantitatively. Listing these consequences together ensures that selecting a solution does not discard any of the information needed in the coupling argument.

\subsection{\texorpdfstring{The successive reference knots}{The successive reference knots}}


Set the zeroth knot to the prescribed initial variable. Define the next knot to be the endpoint at local time \(h\) of the segment started at the current knot. This recursion ensures consecutive segments meet exactly as \(L^2\) classes.



Write \(K_0=Y^{\mathrm{in}}\) and \(K_{n+1}=Y^{(n)}_h\), where the \(n\)-th local solution starts at \(K_n\). This is an ordinary recursion in the space of square-integrable states. Each endpoint belongs to that space by the solution construction. It fixes all future segment initial values from the same original state and Brownian source.

\subsection{\texorpdfstring{The reference segments}{The reference segments}}


At index \(n\), select a restarted brownian-driven segment with the recursively constructed \(n\)-th knot as initial state. The absolute noise start is \(nh\), so all segments are driven by increments of the same Brownian process.



The segment index changes both its initial state and its absolute noise interval. Its equation uses increments beginning at \(nh\), rather than restarting an unrelated Brownian process at each iteration. Consequently the endpoint increment used for the Euler innovation is the same increment appearing in this reference segment. This agreement is the basis of synchronous noise cancellation.

\subsection{\texorpdfstring{The starting value of each segment}{The starting value of each segment}}


The chosen-segment specification states its value at zero equals its supplied initial state. Here that state is the \(n\)-th knot, so \(Y^{(n)}_0=K_n\) follows directly by specialization.



At local time zero the drift integral and Brownian increment both vanish. The segment equation thus gives \(Y^{(n)}_0=K_n\). Together with \(K_n=Y^{(n-1)}_h\) for \(n>0\), this proves exact joining. There is no accumulated boundary error between segments, a necessary fact when applying one identical local comparison repeatedly.

\subsection{\texorpdfstring{The initial reference knot}{The initial reference knot}}


Apply the segment-start identity at \(n=0\), then unfold the zeroth knot. The actual constructed reference therefore begins at exactly the prescribed random variable, preserving its chosen joint realization with the numerical initialization.



The first knot is the prescribed random vector, not merely a new variable with the same law. Hence the first reference segment retains the exact initial coupling with the numerical chain. This matters when the final estimate is proved for arbitrary initial couplings: their root mean-square separation must remain the initial error actually used by the synchronous argument.

\subsection{\texorpdfstring{The endpoint Euler innovations}{The endpoint Euler innovations}}


Normalize the actual segment endpoint noise by \(1/\sqrt{2h}\). This defines the \(L^2\) innovation used by the Euler chain, so it shares the exact Brownian increment driving the reference segment.



The endpoint input is \(N_{nh}(h)=\sqrt2(B_{(n+1)h}-B_{nh})\). Dividing it by \(\sqrt{2h}\) gives \((B_{(n+1)h}-B_{nh})/\sqrt h\), the standard Brownian increment normalization. Defining the innovation from the same input ensures that multiplication by the Euler amplitude recovers the reference noise exactly; its Gaussian law is verified separately.

\subsection{\texorpdfstring{Identification of endpoint noise}{Identification of endpoint noise}}


For \(h>0\), \(\sqrt{2h}\ne0\). Multiply the normalized noise class by \(\sqrt{2h}\) and cancel the inverse scalar. The resulting equality identifies the Euler noise contribution with the same Brownian endpoint input in the integral equation.



Positivity of \(h\) makes \(\sqrt{2h}\) strictly positive, so scalar division and multiplication cancel. Thus \(\sqrt{2h}\,\xi_n=N_{nh}(h)\) as square-integrable random vectors. This equality is stronger than agreement in distribution: it allows the two noise terms to disappear when subtracting the Euler update and the reference endpoint equation on their common space.

\subsection{\texorpdfstring{The stationary Brownian reference}{The stationary Brownian reference}}


Populate l2reference with the constructed segments. Their specification proves continuity and the integral equation; recursive knots prove endpoint compatibility; normalized noise proves endpoint matching; Gibbs-score moment transfer proves the drift bound. The supplied all-time Gibbs-law premise remains visible and is later discharged by the stationarity construction.



Assuming every reference time has Gibbs law, the score moment gives \(\|\nabla U(Y_t^{(n)})\|_2\le\sqrt{Ld}\). The Brownian construction supplies \(\|N_{nh}(t)\|_2=\sqrt{2d}\sqrt t\). Along with segment joining and endpoint noise matching, these bounds meet the analytic coupling hypotheses. Stationarity remains an explicit input at this intermediate stage, rather than a consequence of the integral equation alone.

\subsection{\texorpdfstring{The reference under strong convexity}{The reference under strong convexity}}


Positive strong convexity gives a finite positive Gibbs partition and finite target moments. Apply the reference construction using these target properties and the Brownian endpoint-noise identities. Its all-time stationary-law premise remains explicit at this stage and is supplied later by the invariant-approximation argument.



Positive strong convexity implies a finite positive Gibbs partition and a finite second moment. These normalization and integrability facts discharge the target-side requirements before the reference is assembled. The remaining law premise says that the constructed solution has this target distribution at every segment time. A later stationarity proof supplies that premise using invariant approximations and a continuous-limit argument.

\section{\texorpdfstring{Strong convexity forces m \ensuremath{\le} L}{Strong convexity forces m ≤ L}}
\label{proof-module-UnadjustedLangevin-PotentialAssumptions-lean}


On a nontrivial Hilbert space the numerical compatibility \(m\le L\) follows from strong convexity and a Lipschitz gradient of the same differentiable potential.



For distinct \(x,y\), strong monotonicity and Cauchy–Schwarz give \(m\|x-y\|^2\le\langle\nabla U(x)-\nabla U(y),x-y\rangle\le L\|x-y\|^2\). Cancelling the strictly positive squared distance proves \(m\le L\). This requires a state space containing distinct points. In a singleton space every distance vanishes, so the geometric convergence statements are instead resolved directly with zero transport cost.

\subsection{\texorpdfstring{Compare monotonicity with Cauchy–Schwarz}{Compare monotonicity with Cauchy–Schwarz}}


Choose \(x\ne0\), possible because the Hilbert space is nontrivial. Strong monotonicity gives \(m\|x\|^2\le\langle\nabla U(x)-\nabla U(0),x\rangle\), while Cauchy–Schwarz and Lipschitzness bound the same quantity by \(L\|x\|^2\). Cancel the strictly positive \(\|x\|^2\) to prove \(m\le L\). No positivity premise on \(m\) is needed here.



The nontriviality assumption supplies a direction with positive squared norm. Dividing the two bounds for the same pairing by that norm gives the curvature comparison. On a singleton space no such direction exists, so the inequality need not follow; all transport errors there are instead treated directly as zero.

\section{\texorpdfstring{Normalized increments are independent standard Gaussians}{Normalized increments are independent standard Gaussians}}
\label{proof-module-UnadjustedLangevin-BrownianInnovations-lean}


For a pre-Brownian coordinate and positive grid spacing \(h\), \(Z_n=(B_{(n+1)h}-B_{nh})/\sqrt h\) has standard normal law. Across independent coordinate paths, the entire coordinate-time family is independent.



At a positive step \(h\), set \(\xi_k=(W_{(k+1)h}-W_{kh})/\sqrt h\). Brownian increments have variance \(h\), so these normalized variables have the standard Gaussian law. Their time intervals are disjoint, which gives independence within each coordinate path. Independence of the coordinate paths then combines with that time independence to identify the entire array of scalar innovations, and hence the vector innovations, with the required product Gaussian law.

\subsection{\texorpdfstring{The normalized mesh increment}{The normalized mesh increment}}


Define \(\xi_n=(B_{(n+1)h}-B_{nh})/\sqrt h\). The denominator is interpreted as real scalar division; positive-step law results below explicitly require \(h>0\), so that the variance rescaling is legitimate.



Division by \(\sqrt h\) rescales an increment of variance \(h\) to variance one. The positive-step assumption is needed so that this denominator is nonzero.

\subsection{\texorpdfstring{The standard Gaussian law of a normalized increment}{The standard Gaussian law of a normalized increment}}


A Brownian increment over an interval of length \(h\) is \(\mathcal N(0,h)\). Scaling by \(1/\sqrt h\) multiplies its variance by \(1/h\), producing \(\mathcal N(0,1)\). The proof rewrites \((n+1)h-nh=h\) and uses \((\sqrt h)^2=h\), justified by \(h>0\).



A Brownian increment is centered Gaussian with variance \(h\). Scaling its law by \(1/\sqrt h\) gives \(\mathcal N(0,1)\), because the variance scales by \(1/h\).

\subsection{\texorpdfstring{Square-\allowbreak{}integrability of a normalized increment}{Square-integrability of a normalized increment}}


The exact standard Gaussian law provides every finite moment, in particular the second. Pull this square integrability back through the standard Gaussian law of the rescaled Brownian increment to prove \(\xi_n\in L^2(P)\).



The normalized variable has the standard Gaussian law and therefore finite second moment. Law transfer gives its \(L^2\) membership on the original Brownian source.

\subsection{\texorpdfstring{Its unit second moment}{Its unit second moment}}


Transfer the integral of \(x^2\) along the standard Gaussian law of \(\xi_n\). Its mean is zero and variance one, hence \(\mathbb E\xi_n^2=1\). This is an exact second-moment identity for the normalized increment.



Its second moment is exactly one, agreeing with the unit variance of the scalar innovation law. This supplies the diagonal energy used when assembling vector innovations.

\subsection{\texorpdfstring{Independence of disjoint normalized increments}{Independence of disjoint normalized increments}}


Brownian increments on the consecutive grid intervals are independent. Normalizing each by the same deterministic scalar is a coordinatewise measurable transformation and preserves independence. The formulation permits \(h=0\), where every increment is the constant zero and the independence assertion remains valid.



Different time cells are disjoint, so Brownian independent-increment properties apply. Rescaling each cell by a deterministic nonzero factor preserves independence.

\subsection{\texorpdfstring{Reindexing the independent scalar array}{Reindexing the independent scalar array}}


Assume the row-valued random maps \((X_{i,n})_n\) are independent and each row is internally independent. Replace the almost-everywhere measurable coordinates by measurable versions; countability of \(\mathbb N\) lets all replacements in a row hold on one full-measure set. Apply independence of independent families, then transfer the result back across the coordinatewise almost-everywhere equalities.



The scalar array is indexed jointly by coordinate and time. Reindexing it as one family retains independence because every finite subfamily corresponds to the same finite set of observations.

\subsection{\texorpdfstring{Independence across coordinates and time steps}{Independence across coordinates and time steps}}


Each grid increment depends measurably on two entries of its coordinate grid path. Independence of those paths gives independence across coordinate rows, while Brownian independent increments gives independence within each row. Apply regrouping the independent coordinate-time array to obtain joint independence of all \((i,n)\) innovations.



The within-coordinate independence of disjoint time increments combines with independence of the full spatial paths. Together they give independence of all normalized scalar entries in the array.

\subsection{\texorpdfstring{Independence inherited from spatial paths}{Independence inherited from spatial paths}}


Restrict each independent full coordinate path to the grid \(nh\), preserving independence under composition. The preceding theorem then yields joint coordinate-time independence of the normalized increments. This avoids requiring a separate measurable map into an uncountable path space.



Each spatial path is an independent source factor, and its normalized increments are measurable functions of it. This transfers independence to the corresponding innovation sequences.

\subsection{\texorpdfstring{The full product law of all innovations}{The full product law of all innovations}}


Every coordinate-time innovation the stated pushforward law \(\mathcal N(0,1)\), and the family is jointly independent. The countable product-law theorem therefore identifies the whole innovation map with \(\bigotimes_{(i,n)}\mathcal N(0,1)\). Countability supplies measurability of the path-valued map from coordinatewise almost-everywhere measurability.



The joint innovation law is the product of all standard Gaussian coordinate laws. Grouping entries by time therefore yields independent standard Gaussian vectors, which is stronger than merely Gaussian coordinate marginals.

\section{\texorpdfstring{Pull Brownian noise onto every coupled initialization}{Pull Brownian noise onto every coupled initialization}}
\label{proof-module-UnadjustedLangevin-ProductBrownian-lean}


Brownian-coordinate laws, innovation independence and their L\ensuremath{{}^{2}} realizations survive passage to the product \(\gamma\otimes Q\). This builds all reference fields except the explicitly displayed stationarity premise, discharged by the later canonical assembly.



The product construction pulls Brownian coordinate processes back along the noise projection. Because that projection preserves their probability measure, their Gaussian increment laws and independence remain valid. The initial pair lies in the other factor, so the entire noise family is independent of initialization. The resulting \(L^2\) reference has stationary Gibbs marginal, the exact Brownian increment norm, and a common numerical endpoint noise, supplying the ingredients for the uniform initial-coupling comparison.

\subsection{\texorpdfstring{Pulling Brownian coordinates back by a measure-\allowbreak{}preserving map}{Pulling Brownian coordinates back by a measure-preserving map}}


A measure-preserving map leaves the law of every measurable random variable unchanged after pullback. Apply this to every finite-dimensional Brownian vector. Their Gaussian laws and covariance data are therefore preserved, giving a pre-Brownian process on the new sample space.



If a map preserves the source probability, the law of any measurable finite collection of source evaluations is unchanged after composition. Apply this to the Brownian values at arbitrary finitely many times. Their centered Gaussian joint distribution and covariance \(\min(s,t)\) are preserved. Thus pulling the paths to the coupling extension retains the required Brownian finite-dimensional laws.

\subsection{\texorpdfstring{Pulling back the square-\allowbreak{}integrable Brownian vector}{Pulling back the square-integrable Brownian vector}}


Both sides represent the vector \((B_i(t,f(\omega)))_i\). Pulling back a full-measure representative equality along a measure-preserving map again gives a full-measure equality. The two constructions therefore determine the same square-integrable class.



Both constructions evaluate the same finite vector of scalar Brownian coordinates after the measure-preserving map. Their representatives agree almost surely because the source representative equality pulls back to a full-measure event. Equality almost surely identifies the two \(L^2\) vectors. This makes coordinate assembly compatible with moving between source and extended probability spaces.

\subsection{\texorpdfstring{Pulling back the endpoint innovation}{Pulling back the endpoint innovation}}


The \(L^2\) pullback is linear, hence commutes with subtraction and scalar multiplication. Apply the preceding Brownian-vector pullback identity to the two endpoints in the normalized increment. This identifies the pulled-back Euler noise with the noise constructed after pullback.



The normalized Euler noise is a fixed scalar multiple of a difference of two Brownian \(L^2\) vectors. A measure-preserving pullback is linear, so it commutes with that subtraction and scalar multiplication. Applying the coordinate-vector pullback identity at both endpoints yields equality of the two noise constructions. Therefore extending first or normalizing first produces the same actual innovation.

\subsection{\texorpdfstring{The source Euler innovation}{The source Euler innovation}}


Take the actual selected representative of the previously constructed \(L^2\) Brownian Euler noise on the source space. This is the random function supplied to the generic product algorithm before choosing any initial coupling.



Choose a measurable random-function version of each source Euler innovation. These functions supply the input path for the generic product-law construction. Their finite-time laws and independence will be proved from the Brownian formulas below. Fixing the functions on the source before choosing an initial coupling ensures that the numerical output can be represented by one common source law.

\subsection{\texorpdfstring{Square-\allowbreak{}integrability of the source innovation}{Square-integrability of the source innovation}}


Every representative of a given \(L^2\) class is square integrable by has a finite square moment. Apply that fact to the Brownian Euler noise class. This avoids assuming an extra noise-moment condition in the product construction.



An \(L^2\) random variable has finite second moment for every representative agreeing with it almost surely. Therefore the selected source innovation functions inherit the finite norm already established for the Brownian-normalized class. No additional moment condition is required merely because a pointwise version is chosen for the product recursion.

\subsection{\texorpdfstring{Its vector representative}{Its vector representative}}


Unfold the nested class-valued noise definition and apply representative equalities for scalar multiplication, subtraction, and the Brownian vector classes. Almost everywhere the result is \((\sqrt{2h})^{-1}\sqrt2(B_{nh+h}-B_{nh})\), with real-time clipping as in the original definition.



Expand the innovation into its two Brownian endpoint values and fixed normalization factors. Almost-sure identities for sums, differences and scalar multiples show that its pointwise version equals the normalized vector increment. The time extension only specifies how Brownian evaluation is interpreted outside nonnegative times; in the positive mesh regime every grid endpoint already lies in the ordinary Brownian domain.

\subsection{\texorpdfstring{Its scalar coordinate representative}{Its scalar coordinate representative}}


Evaluate the vector representative formula in coordinate \(i\). For \(h>0\), use \(\sqrt{2h}=\sqrt2\sqrt h\) and cancel the nonzero \(\sqrt2\), obtaining \((B_i((n+1)h)-B_i(nh))/\sqrt h\).



For \(h>0\), factor \(\sqrt{2h}=\sqrt2\sqrt h\) and cancel the positive \(\sqrt2\). In coordinate \(i\), the result is \([B_i((n+1)h)-B_i(nh)]/\sqrt h\). This explicit formula connects the exact innovation representative to the standard Gaussian increment law and fixes the variance normalization without an inequality.

\subsection{\texorpdfstring{The coordinate at nonnegative mesh times}{The coordinate at nonnegative mesh times}}


For nonnegative mesh times, restriction to nonnegative time leaves their real values unchanged. Rewrite both clipped grid endpoints into the corresponding nonnegative-real products. The coordinate formula becomes the ordinary Brownian grid increment divided by \(\sqrt h\).



All grid times \(nh\) and \((n+1)h\) are nonnegative when \(h>0\). Hence the real-time extension agrees there with the original nonnegative-time Brownian process. The coordinate identity is consequently the usual Brownian grid-increment formula. This removes time-domain bookkeeping before applying the Gaussian increment and independence results.

\subsection{\texorpdfstring{The normalized Brownian coordinate increment}{The normalized Brownian coordinate increment}}


The nonnegative-grid formula is exactly the definition of the normalized brownian mesh increment. After converting the positive real step to its nonnegative-real form, the two random functions agree almost everywhere. This gives a direct bridge to the established innovation-law theorems.



The innovation coordinate equals the established normalized Brownian increment almost surely at each fixed step and coordinate. Thus every distributional statement proved for the normalized increment applies to the actual source coordinate. This bridge is needed because the numerical chain uses selected \(L^2\) representatives, while the Brownian law theorem is initially stated for raw scalar increment functions.

\subsection{\texorpdfstring{The scalar innovation law}{The scalar innovation law}}


A normalized Brownian increment the stated pushforward law \(\mathcal N(0,1)\). Transfer this law across the equality between a scalar mesh innovation and its normalized Brownian increment using almost-everywhere invariance of laws. Hence the actual source noise representative has the required Gaussian coordinate law.



A length-\(h\) Brownian increment is centered Gaussian with variance \(h\). Scaling by \(h^{-1/2}\) multiplies its variance by \(h^{-1}\), giving variance one. Equality almost surely of the two coordinate versions preserves their pushforward distributions. This proves the standard Gaussian law of the precise source innovation used in the algorithm.

\subsection{\texorpdfstring{Independence of the source innovation array}{Independence of the source innovation array}}


The Brownian grid-path theorem gives independence of all normalized coordinates indexed by \((i,n)\). Reindex by the injective coordinate swap to \((n,i)\), then transfer across the almost-everywhere coordinate identities. This matches the index order expected by the vector innovation theorem.



Independent Brownian coordinates and independent disjoint increments give one jointly independent family indexed by coordinate and time interval. Swapping the order of these two indices is a bijection and preserves every finite factorization. Almost-sure coordinate replacement then transfers that full independence to the selected source innovations. Correct marginal laws alone would not establish this joint conclusion.

\subsection{\texorpdfstring{Brownian coordinates on the product space}{Brownian coordinates on the product space}}


Projection from an initial-coupling product extension to the noise source is measure preserving. Pull each scalar Brownian coordinate back along it using preservation of Brownian finite-dimensional laws under a measure-preserving pullback. The actual coordinate processes on the extension retain their Brownian finite-dimensional laws.



On \(c\otimes Q\), the second-factor projection is measure preserving because \(c\) has mass one. Pulling every scalar Brownian path through it leaves its finite-dimensional distributions unchanged. This constructs Brownian coordinate processes on the same space as the coupled initial states, with no dependence of their source laws on the particular initial coupling.

\subsection{\texorpdfstring{Identification of product and reference noise}{Identification of product and reference noise}}


The lifted source noise and the extension's constructed Brownian noise are both the measure-preserving pullback of the source \(L^2\) Brownian Euler noise. Use representative extensionality to identify the former with this pullback, then apply the equality of mesh increments before and after pullback to identify the latter.



The product innovation is the pullback of the fixed source innovation, whereas the reference noise is normalized after pulling back the Brownian coordinates. Linearity and the preceding pullback identities show that these are the same \(L^2\) vector. Consequently the numerical and reference constructions genuinely share endpoint noise, permitting exact cancellation in their coupled difference.

\subsection{\texorpdfstring{The product coordinate innovation law}{The product coordinate innovation law}}


Pull the source standard-Gaussian coordinate law through the extension's noise projection. The lifted product-noise representative agrees almost everywhere with the raw pulled-back source noise. Transfer the law across that equality to obtain the coordinate law for the exact classes used by the numerical chain.



The source projection preserves the standard Gaussian coordinate law, so its raw pullback remains standard Gaussian. The product innovation representative agrees almost surely with that pullback. Changing representatives therefore gives the same Gaussian marginal for the exact variables used by the numerical chain. This verifies the distributional premise without replacing the coupled construction by a separate Gaussian source.

\subsection{\texorpdfstring{Independence of product innovations}{Independence of product innovations}}


Identify the source coordinate-time path law as the product of standard Gaussians. Pull that law to the extension; its coordinate family is independent by the product-law independence criterion. Finally replace raw source pullbacks by the almost-everywhere equal product-noise representatives.



Pulling the complete source innovation path through a measure-preserving map preserves its whole joint product Gaussian law. Every finite subfamily of product innovation coordinates therefore remains independent. Countably many almost-sure coordinate equalities can be imposed simultaneously on a full-measure event, allowing the raw path law to be transferred to the chosen representative path.

\subsection{\texorpdfstring{Independence from numerical initialization}{Independence from numerical initialization}}


The raw first initial coordinate depends only on the coupling factor, while the whole source noise sequence depends only on the noise factor. Product independence therefore applies. Countability of the sequence lets representative equalities hold simultaneously at all indices; transport independence to the selected \(L^2\) initialization and noise representatives.



The first initial coordinate is measurable from the coupling factor, while the complete innovation path is measurable from the independent source factor. Product independence gives their factorized joint law. Simultaneous almost-sure representative identities preserve this independence for the \(L^2\) initialization and actual innovation path. Independence therefore applies to all future innovations together, as required by the recursion law theorem.

\subsection{\texorpdfstring{Independence from reference initialization}{Independence from reference initialization}}


Use the same product independence argument with the second initial coordinate. Both initialization and whole noise-path representative equalities are transferred almost everywhere, proving the actual reference class is independent of all fresh Euler noise.



The reference initialization uses the second coordinate of the same coupling factor and is likewise independent of the entire source. Replacing it and the innovation path by almost-surely equal representatives leaves the joint law unchanged. This independence is used when constructing a Langevin process started at Gibbs law with fresh Brownian forcing, rather than only in the numerical algorithm-law argument.

\subsection{\texorpdfstring{The Brownian product probability}{The Brownian product probability}}


The coupling is a probability by its marginal identity and the Brownian source measure is a probability. Their product extension is normalized. This normalization supplies the probability-space hypotheses of the Brownian pullback and \(L^2\) reference constructors.



The Brownian extension's normalization follows from \(c(E\times E)=1\) and \(Q(S)=1\). Thus all pulled-back Brownian and square-integrable constructions occur on a genuine probability space. Unit mass also ensures that adding the source does not change initial second moments, noise norms or the cost of the chosen initial coupling.

\subsection{\texorpdfstring{The Gibbs reference as a square-\allowbreak{}integrable class}{The Gibbs reference as a square-integrable class}}


Name the already proved finite-second-moment theorem for the strongly convex Gibbs law. The abbreviation supplies exactly the \(L^2(\pi_U)\) identity moment needed to construct the reference initialization on every coupling extension.



Positive strong convexity gives a Gaussian upper bound on the target density and its quadratic moment weight. The normalized Gibbs law consequently has finite second moment. This proved target fact supplies the reference-coordinate integrability on every finite-coupling extension, and finite coupling cost then supplies the algorithm-coordinate integrability as well.

\subsection{\texorpdfstring{The square-\allowbreak{}integrable Brownian product reference}{The square-integrable Brownian product reference}}


The pulled-back Brownian segments have continuous square-mean paths and satisfy the gradient integral equation. Their endpoint Brownian increments equal the common product innovations. Given the supplied stationary laws, the Gibbs score bound supplies the remaining drift-moment requirement. Thus every stated reference property is available, with stationarity retained as a premise of this intermediate result.



The constructed local segments supply continuity, joining and the gradient integral equation. Their Brownian endpoint noise equals the common product innovation by the preceding identity. Given their target laws, the Gibbs score moment supplies the uniform drift bound. The resulting analytic reference therefore has every required field, while segment stationarity remains an explicit premise at this intermediate stage.

\subsection{\texorpdfstring{The product reference supplied by the Brownian construction}{The product reference supplied by the Brownian construction}}


For each finite initial coupling, use the square-integrable brownian product reference. Its initial reference equals the second-coordinate \(L^2\) class by the segment-initial identity. These constructions fill the a reference on the independent product construction fields. Only the explicitly quantified segment-stationarity premise remains an input at this intermediate stage.



For every finite initial coupling, form the same source product and construct its Brownian reference. Its first segment starts at the actual second-coordinate reference state, preserving the initial cost. These references together satisfy the coupling-family requirement. The only remaining probabilistic obligation is the stated Gibbs law at each segment time; later canonical stationarity discharges it uniformly across these constructions.

\subsection{\texorpdfstring{The convergence bound from the constructed product Brownian source}{The convergence bound from the constructed product Brownian source}}


Derive \(m\le L\) from strong convexity and smoothness, construct the product reference from the stated stationary segment laws, and apply the small-step product-reference theorem. This intermediate endpoint has the genuine initial \(W_2\) term; the later independent Brownian construction discharges its stationarity premise.



Strong convexity and smoothness imply \(m\le L\) on a nontrivial space, placing the parameters in the deterministic contraction regime. Apply the product-reference theorem with the constructed Brownian family and supplied segment target laws. Taking the initial coupling infimum gives the genuine initial Wasserstein term. This endpoint is still conditional on segment stationarity, which the later complete construction proves.

\section{\texorpdfstring{Independent realization of an initial coupling and noise}{Independent realization of an initial coupling and noise}}
\label{proof-module-IPLA-ProductReference-lean}


The product of an arbitrary initial coupling and an independent noise source gives a common probability space for the algorithm and reference. Its product structure supplies initialization/\allowbreak{}noise independence and preserves each initial marginal. Crucially, changing the coupling changes only their dependence, not the law of the algorithm output.


\subsection{\texorpdfstring{Separate initialization from the noise source}{Separate initialization from the noise source}}


On \((H\times H)\times S\) with law \(c\otimes Q\), let the two initial variables be the first and second coordinates of \(c\), and let every innovation depend only on \(S\). Product independence makes both initial variables independent of the noise source. The first chain evolves by \(x\mapsto x-h\nabla V(x)+\sqrt{2h}\xi_k\), and its output law is the pushforward of \(\mu_0\otimes Q\) by the measurable recursion; hence it depends only on \(\mu_0\), not on the chosen coupling \(c\). This verifies the common-output obligation needed to take the coupling infimum. Finite coupling cost and Gibbs moments supply the \(L^2\) initial realization; the proved output finiteness justifies the subsequent real-valued transport bounds.




Let the coupling coordinates be \(X_0,Y_0\) and let \(S\) carry the noise. Under \(c\otimes Q\), the law of \((X_0,S)\) is always \(\mu_0\otimes Q\), regardless of the second coordinate or the coupling dependence. The measurable Euler recursion uses only this pair, so its output law is fixed. The reference may use \(Y_0\) and the same source, providing the required synchronous coupling for every initial cost.

\section{\texorpdfstring{Canonical Gaussian path law and realization independence}{Canonical Gaussian path law and realization independence}}
\label{proof-module-IPLA-AlgorithmLaw-lean}


Once the initial law and the independent Gaussian path law are fixed, the algorithm law is obtained by applying a deterministic measurable recursion to their product. This gives an realization-independent meaning to the finite-time output distribution. It also permits coupling constructions on auxiliary probability spaces without changing the law of the algorithm being analyzed.




The scalar Gaussian coordinate hypotheses uniquely determine the entire countable noise path. Independence of initialization then determines every output law, even across different underlying probability spaces.

\subsection{\texorpdfstring{Construct the independent Gaussian noise path}{Construct the independent Gaussian noise path}}


A noise coordinate is a pair \((k,a)\) consisting of iteration and state coordinate. Put the countable product law \(\bigotimes_{(k,a)}\mathcal N(0,1)\) on scalar paths \(w\). Packing sets \((\operatorname{pack}w)_k(a)=w(k,a)\); its measurability follows because every output coordinate is a measurable product projection, followed by the measurable Euclidean coordinate identification. The canonical vector path law is the pushforward of that scalar product law by packing.





A cylinder event specifies only finitely many scalar coordinates and has probability equal to the product of their standard Gaussian probabilities. These cylinder probabilities determine the countable product law. Grouping the coordinates into vectors changes their representation but does not change any event or distributional relationship. The resulting time-indexed vectors therefore retain both coordinate independence and independence between distinct time slices.

\subsection{\texorpdfstring{The algorithm hypotheses identify the entire path law}{The algorithm hypotheses identify the entire path law}}


Mutual independence of the scalar family implies that its joint map has law equal to the infinite product of its individual coordinate laws. Since each law is \(\mathcal N(0,1)\),

\[\mathcal L((\xi_{k,a})_{k,a})=\bigotimes_{k,a}\mathcal N(0,1).\]

Composing this map with measurable packing and using the composition rule for pushforwards gives \(\mathcal L((\xi_k)_k)=\mathcal G\). The argument proves independence of vector blocks from scalar independence and never assumes a grouping principle without justification.




The canonical flat source is the product of one scalar Gaussian for each time/\allowbreak{}coordinate pair. The packing map merely groups these coordinates into a vector at each time; coordinate evaluation proves its measurability. The supplied innovation structure has the same finite-coordinate Gaussian product laws by mutual independence, so the product-law uniqueness theorem identifies its whole flat path law. Applying the measurable packing map gives equality of the full vector path law, including all cross-time relations.





The joint distribution theorem for an independent countable family applies to all time-coordinate pairs at once. It is stronger than knowing the correct vector law separately at each time: separate laws alone would permit temporal dependence. After identifying this full scalar law, composition with coordinate grouping proves the vector-path identity. This is the exact distributional input used when iterating a Markov transition.

\subsection{\texorpdfstring{Identify the common update and recursion}{Identify the common update and recursion}}




Write the common update as \(T(z,u)=z-(\gamma/N)\nabla V_N(z)+\sqrt{2\gamma/N}\,u\). Equality of updates gives equality of the zero-step chains immediately. If two chains agree after \(k\) steps, applying the same function \(T\) to the equal states and the same \(k\)-th innovation proves agreement after \(k+1\) steps. No distributional premise is required for this deterministic induction.

\subsection{\texorpdfstring{A common initial law gives common output laws}{A common initial law gives common output laws}}


Let \(F_k(z_0,\xi)\) be the measurable map that iterates the update \(k\) times. Initialization–path independence gives the joint input law \(\mu_0\otimes\mathcal G\), where \(\mathcal G\) is the canonical Gaussian path law. Hence

\[\rho_k=(F_k)_\#(\mu_0\otimes\mathcal G).\]


The exact update identity matches the generic noise-driven recursion to the scaled IPLA recursion. Independence then makes the joint initialization/\allowbreak{}path law the product of the initial pushforward and the canonical noise path law. Push this common source through the measurable time-\(k\) recursion. Two realizations on different probability spaces therefore have the same output law whenever their initial laws agree; the conclusion is about the actual entire algorithm law at the specified time.





The map \(F_k\) depends only on the initialization and the first \(k\) vector innovations. Its measurability follows by composing the measurable update a finite number of times. Independence then factors the input law before any recursion is evaluated. Thus a coupling construction is free to alter the auxiliary probability space, provided it preserves these input laws and their independence; the algorithm's output marginal remains fixed.

\section{\texorpdfstring{Euler approximations converge to the integral solution}{Euler approximations converge to the integral solution}}
\label{proof-module-UnadjustedLangevin-EulerLimit-lean}


This module proves a deterministic convergence theorem for additive Euler discretization. With mesh \(h=T/n\), its terminal value converges to the continuous integral solution, including when the state space itself is \(L^2(P)\).



The numerical approximation uses the increments of the same additive input as the exact integral solution. Noise consequently cancels in the comparison, leaving only the error from freezing the drift on each mesh interval. Uniform continuity of the drift along the exact compact-time trajectory makes that local error smaller than \(h\varepsilon\) on a fine mesh. The bound \((1+Lh)^n\le e^{Lnh}\) keeps accumulation uniformly controlled over a fixed horizon, so the final error tends to zero.

\subsection{\texorpdfstring{The Euler chain for a fixed additive input}{The Euler chain for a fixed additive input}}


Define \(X_0=x\) and \(X_{k+1}=X_k-hg(X_k)+B((k+1)h)-B(kh)\). The additive input is one fixed function evaluated at mesh endpoints; consequently the numerical chain and the integral solution can be compared with exactly matching noise increments.



Using increments of one fixed input makes the discrete and continuous equations directly comparable. The numerical error is caused by drift freezing, rather than by approximating the noise itself.

\subsection{\texorpdfstring{Its initial value}{Its initial value}}


Evaluating the recursive definition at index zero gives \(X_0=x\). The proof establishes this by the defining identity; no continuity of the input or positivity of the mesh is required for this initialization identity.



The zero index uses no input increment, so the chain starts exactly at its prescribed point. This supplies the zero initial error when the exact solution is initialized at the same point.

\subsection{\texorpdfstring{Its successor update}{Its successor update}}


Unfold the successor clause once to expose \(X_{k+1}=X_k-hg(X_k)+\Delta B_k\). This defining identity is the equation used when subtracting the exact solution's increment and deriving the error recurrence.



Writing \(\Delta B_k=B((k+1)h)-B(kh)\) isolates the common additive increment. It cancels when this equation is subtracted from the exact integral increment over the same mesh interval.

\subsection{\texorpdfstring{The exact integral increment}{The exact integral increment}}


Subtract the integral equations at \(b\) and \(a\). Continuity of \(g\circ Y\) makes it interval integrable, so interval-integral additivity gives \(\int_0^b-\int_0^a=\int_a^b\). Thus \(Y(b)-Y(a)=-\int_a^b g(Y(s))ds+B(b)-B(a)\), including oriented intervals when \(b<a\).



The drift integral over \([a,b]\) is the difference of the two integrals starting at zero. This interval-additivity identity is the step that aligns the exact solution with the numerical mesh.

\subsection{\texorpdfstring{A finite-\allowbreak{}horizon recurrence estimate}{A finite-horizon recurrence estimate}}


From \(e_{k+1}\le qe_k+b\), \(e_0=0\), \(q\ge1\), and \(b\ge0\), induction proves \(e_n\le nbq^n\). In the successor step bound \(q\cdot kbq^k+b\) by \((k+1)bq^{k+1}\), using \(q^{k+1}\ge1\). This coarse bound avoids dividing by \(q-1\), so it also covers \(q=1\).



The bound \(nbq^n\) is deliberately coarse and remains valid at \(q=1\). It avoids a division by \(q-1\), making the Euler convergence estimate uniform as the mesh tends to zero.

\subsection{\texorpdfstring{Stability of the explicit Euler map}{Stability of the explicit Euler map}}


Write the difference as \((x-y)-h(g(x)-g(y))\). The triangle inequality and Lipschitzness give \(\|\cdot\|\le\|x-y\|+hL\|x-y\|=(1+Lh)\|x-y\|\). The hypothesis \(h\ge0\) removes the absolute value in the scalar norm.



This estimate uses only Lipschitzness and a triangle inequality, so its factor is \(1+Lh\). It controls finite-time stability of the approximation without requiring strong convexity or a contractive Euler step.

\subsection{\texorpdfstring{The error from a drift continuity modulus}{The error from a drift continuity modulus}}


At one step the common noise cancels. The Euler map contributes \((1+Lh)e_k\), and the exact integral remainder contributes at most \(h\varepsilon\), because the drift changes by at most \(\varepsilon\) within one mesh interval. Apply the finite recurrence bound and \((1+Lh)^n\le e^{Lnh}\), then use \(nh=T\) to obtain \(e_n\le T\varepsilon e^{LT}\).



Within each mesh cell, a drift variation at most \(\varepsilon\) contributes at most \(h\varepsilon\) to the local remainder. Multiplying by the finite-time stability factor and using \(nh=T\) gives the stated uniform horizon bound.

\subsection{\texorpdfstring{Convergence of the matching-\allowbreak{}input Euler approximation}{Convergence of the matching-input Euler approximation}}


The continuous drift path \(g\circ Y\) is uniformly continuous on \([0,T]\). Given any positive error tolerance, choose a drift modulus \(\varepsilon\) making \(T\varepsilon e^{LT}\) small, then take the mesh \(T/(n+1)\) below its continuity radius. The previous error bound gives endpoint convergence. The case \(T=0\) is separated so no division by the horizon is needed.



Uniform continuity of the exact drift curve supplies a single mesh size that controls every cell of the horizon. Sending that mesh size to zero makes the remainder tolerance arbitrarily small, and hence the endpoint error tends to zero.

\subsection{\texorpdfstring{Euler convergence in square mean}{Euler convergence in square mean}}


Apply the deterministic matching-input Euler convergence theorem in \(L^2(P;H)\), using the Lipschitz lifted gradient. The numerical endpoint converges in the \(L^2\) norm to the constructed solution at \(T\). This is the topology later used to pass smooth bounded tests and identify laws.



Convergence is in the norm of the random-vector Hilbert space. It therefore controls expectations of bounded Lipschitz tests by Cauchy–Schwarz, which is the form later used to identify the limiting marginal law.

\section{\texorpdfstring{A finite O(h\ensuremath{{}^{3/2}}) rejected-\allowbreak{}displacement bound}{A finite O(h³ᐟ²) rejected-displacement bound}}
\label{proof-module-UnadjustedLangevin-MALAError-lean}


The Metropolis-adjusted Langevin proposal provides an exactly invariant approximation to the diffusion. This module proves that its mean rejected displacement under Gibbs initialization is at most \(Ch\sqrt h\), with a finite constant derived from the potential assumptions.



For the Langevin proposal, the log Metropolis ratio compares the target-density ratio with the reversed Gaussian-proposal ratio. Expanding the two quadratic exponents exposes a trapezoid error in the potential and a difference of squared gradient norms. Lipschitz-gradient estimates control both by powers of the proposal displacement. The scalar bound \(1-\min(1,e^r)\le|r|\) converts this log-ratio control into a rejection bound; multiplying by displacement gives the expected correction size.

\subsection{\texorpdfstring{The Langevin Gaussian proposal}{The Langevin Gaussian proposal}}


Define \(T_h(x,\xi)=x-h\nabla U(x)+\sqrt{2h}\xi\). This is the same Euler proposal used by ULA; MALA applies an accept/\allowbreak{}reject rule to this proposed state.



Conditioned on the current state, the proposal has mean \(x-h\nabla U(x)\) and covariance \(2hI_d\). Thus its density is the same Gaussian kernel used by ULA. Only the subsequent decision to accept or remain at \(x\) distinguishes the two algorithms.

\subsection{\texorpdfstring{The logarithmic Metropolis ratio}{The logarithmic Metropolis ratio}}


The log acceptance ratio is \(R_h(x,y)=U(x)-U(y)+(\|y-x+h\nabla U(x)\|^2-\|x-y+h\nabla U(y)\|^2)/(4h)\). The common Gaussian normalizers cancel in the reverse/\allowbreak{}forward proposal density ratio.



The ratio is \(\pi(y)q_h(y,x)/(\pi(x)q_h(x,y))\). Taking its logarithm cancels the target partition function as well as the two Gaussian normalizers, leaving the displayed expression. Positive \(h\) gives nondegenerate proposal densities in both directions.

\subsection{\texorpdfstring{Expansion of the proposal ratio}{Expansion of the proposal ratio}}


Expand both squared norms using the real inner-product norm identity. The \(\|y-x\|^2\) terms cancel, and division by \(4h\) gives \(R_h=U(x)-U(y)+\frac12\langle\nabla U(x)+\nabla U(y),y-x\rangle-\frac h4(\|\nabla U(y)\|^2-\|\nabla U(x)\|^2)\). The premise \(h\ne0\) justifies denominator cancellation.



Writing \(d=y-x\), expand \(\|d+hg_x\|^2\) and \(\|-d+hg_y\|^2\) before subtracting. Their common \(\|d\|^2\) term disappears, and the surviving linear terms add to \(2h\langle g_x+g_y,d\rangle\). The remaining difference of squared gradients has coefficient \(-h/4\).

\subsection{\texorpdfstring{Bounding the difference of squared gradients}{Bounding the difference of squared gradients}}


Factor \(|\|g_y\|^2-\|g_x\|^2|=|\|g_y\|-\|g_x\||(\|g_y\|+\|g_x\|)\). The reverse triangle inequality bounds the first factor by \(\|g_y-g_x\|\), then Lipschitzness bounds it by \(L\|y-x\|\). The other factor is nonnegative, so multiplication preserves the estimate.



The estimate does not bound \(\|g_x\|\) or \(\|g_y\|\) uniformly over space. Instead it isolates their growth in one factor and uses \(L\|y-x\|\) for their difference. The resulting expression is compatible with polynomial moment bounds under Gibbs.

\subsection{\texorpdfstring{A trapezoid bound for the potential increment}{A trapezoid bound for the potential increment}}


Use the smooth upper and lower first-order models for the potential increment at the two endpoints. Compare the increment with the average endpoint directional derivative, and use the Lipschitz bound on the derivative difference. The resulting trapezoid error is at most \(L\|y-x\|^2\); no second derivative of \(U\) is assumed.



Restrict \(U\) to the segment \(x+td\). Its derivative is \(\langle\nabla U(x+td),d\rangle\), whose variation is at most \(L|t-s|\|d\|^2\). Integrating this variation bounds the difference between the true increment and the endpoint trapezoid, using only a Lipschitz first derivative.

\subsection{\texorpdfstring{An absolute log-\allowbreak{}ratio bound}{An absolute log-ratio bound}}


Insert the exact log-ratio expansion and use the triangle inequality. Bound its potential trapezoid term by \(L\|y-x\|^2\) and its squared-gradient difference by \(L\|y-x\|(\|g_x\|+\|g_y\|)\). Positivity of \(h\) gives the displayed coefficient \(hL/4\) without an absolute value on the step.



The potential term is quadratic in displacement, whereas the second term contains \(h\) times displacement and gradient size. When the Gaussian proposal displacement is of order \(\sqrt h\), both are small enough that multiplication by the rejection displacement produces an integrable error of order \(h^{3/2}\).

\subsection{\texorpdfstring{Rejection probability bounded by the log ratio}{Rejection probability bounded by the log ratio}}


If \(r\ge0\), \(e^r\ge1\), so rejection is zero. If \(r<0\), \(1-e^r\le-r=|r|\) by \(e^r\ge1+r\). Both cases also give \(0\le1-\min(1,e^r)\), establishing a nonnegative rejection probability and its log-ratio bound.



The exponential tangent inequality \(e^r\ge1+r\) is valid for every real \(r\). On the negative branch it turns a nonlinear rejection probability into a linear bound by \(-r\); on the nonnegative branch rejection is exactly zero. This avoids a Taylor remainder for acceptance.

\subsection{\texorpdfstring{The rejection-\allowbreak{}weighted displacement}{The rejection-weighted displacement}}


Multiply \(1-\min(1,e^r)\le|r|\) by the nonnegative displacement norm \(\|y-x\|\). This yields the rejection-weighted coupling error bound with no condition on the states.



The displacement factor must remain outside the acceptance comparison because the later coupling cost measures distance, rather than rejection probability alone. Integrating this product will control how much rejecting a proposal separates MALA from the common-noise ULA chain.

\subsection{\texorpdfstring{A proposal-\allowbreak{}amplitude envelope}{A proposal-amplitude envelope}}


Use the step-independent amplitude \(A(x,\xi)=\|\nabla U(x)\|+\sqrt2\|\xi\|\). For \(0\le h\le1\), it bounds the proposal displacement after extracting \(\sqrt h\); its cubic integrability will give a uniform rejection-error constant.



The amplitude contains the score at the initial point and the size of the independent innovation. Its law is fixed when \(h\) varies. This separation is what permits a single integrable majorant and a single rejection constant for an entire vanishing-mesh sequence.

\subsection{\texorpdfstring{The rejected-\allowbreak{}displacement correction}{The rejected-displacement correction}}


Define \(D_h(x,\xi)=\|T_h(x,\xi)-x\|(1-\min(1,e^{R_h(x,T_h(x,\xi))}))\). This is the acceptance-variable-averaged distance lost on rejection, so no additional uniform random variable appears in its integral.



For a fixed state and innovation, integrating over an independent uniform accept/\allowbreak{}reject variable would give exactly this expected separation between the proposal and accepted state. The weighted-measure definition performs that integration directly and keeps the same coupling cost.

\subsection{\texorpdfstring{Nonnegativity of the amplitude envelope}{Nonnegativity of the amplitude envelope}}


Both score and noise norms are nonnegative, and \(\sqrt2\ge0\). Their sum defining \(A\) is therefore nonnegative. This sign permits later powers, products, and integral comparisons without absolute-value ambiguity.



When both score and innovation vanish, the amplitude is zero and so is the proposal displacement. Otherwise its nonnegative powers can be used as majorants without introducing an additional absolute value, which simplifies the cubic-moment calculation.

\subsection{\texorpdfstring{Nonnegativity of the correction}{Nonnegativity of the correction}}


The proposal displacement norm is nonnegative and the rejection factor is nonnegative by the rejection estimate \(1-\min(1,e^r)\le|r|\). Their product is nonnegative, permitting use of extended integrals and domination by an integrable envelope.



This nonnegativity means the expected rejection displacement can be handled first as a nonnegative integral. Once an integrable cubic envelope is available, that same quantity is a finite ordinary expectation and can enter the characteristic-function comparison.

\subsection{\texorpdfstring{A bound on proposal displacement}{A bound on proposal displacement}}


The triangle inequality gives \(\|T_h-x\|\le h\|g_x\|+\sqrt{2h}\|\xi\|\). For \(0\le h\le1\), \(h\le\sqrt h\) and \(\sqrt{2h}=\sqrt2\sqrt h\). Factor out \(\sqrt h\) to obtain \(\|T_h-x\|\le\sqrt h A(x,\xi)\).



For \(0\le h\le1\), the elementary inequality \(h^2\le h\) gives \(h\le\sqrt h\). Thus even the deterministic drift displacement can be placed on the same \(\sqrt h\) scale as the Gaussian displacement, with the drift magnitude retained inside the amplitude.

\subsection{\texorpdfstring{A pointwise bound on rejected displacement}{A pointwise bound on rejected displacement}}


Use rejection \(\le|R_h|\), the log-ratio estimate, and \(\|g_{T_h}\|\le\|g_x\|+L\|T_h-x\|\). Replace the displacement by \(\sqrt hA\), and absorb the remaining step factors with \(h\le1\). The result is \(D_h\le h\sqrt h\,L(2+L)A^3\), whose envelope has no step-size dependence.



The estimate has a constant independent of \(h\): all score and innovation dependence is in \(A^3\). In particular, averaging under Gibbs and Gaussian noise gives an \(O(h^{3/2})\) defect, whose accumulation over \(T/h\) steps is \(O(T\sqrt h)\).

\subsection{\texorpdfstring{A Gaussian fourth-\allowbreak{}moment integrability bound}{A Gaussian fourth-moment integrability bound}}


A polynomial of degree four can be absorbed by part of Gaussian decay: \(\|x\|^4e^{-a\|x\|^2}\le C e^{-(a/2)\|x\|^2}\) for \(a>0\). Gaussian integrability and measurable domination therefore prove the fourth weighted moment is finite.



The scalar function \(r\mapsto r^4e^{-ar^2/2}\) is bounded on \([0,\infty)\). Multiplying that bound by \(e^{-ar^2/2}\) provides an integrable Gaussian majorant. The same method works for other fixed polynomial degrees whenever a positive part of the quadratic decay is reserved.

\subsection{\texorpdfstring{The weighted Gibbs fourth moment}{The weighted Gibbs fourth moment}}


The quadratic strong-convexity lower bound gives \(e^{-U(x)}\le C e^{-(m/4)\|x\|^2}\). Multiply by \(\|x\|^4\) and invoke the fourth Gaussian weighted-moment theorem. This derives the unnormalized Gibbs fourth moment from the original potential assumptions.



A linear term in the strong-convexity lower model is absorbed by completing a square, leaving a positive quadratic coefficient and a finite multiplicative constant. Multiplication by \(\|x\|^4\) then falls under the Gaussian domination just established, with no stationarity assumption needed.

\subsection{\texorpdfstring{The normalized Gibbs fourth moment}{The normalized Gibbs fourth moment}}


Change measure to the normalized Gibbs density. Integrability of \(\|x\|^4\) under \(\pi_U\) reduces to the just-proved integrability of \(\|x\|^4e^{-U(x)}\) times a finite normalization factor. No fourth-moment premise is added.



Division by the strictly positive finite partition function preserves finiteness of the fourth moment. The result applies to the actual normalized target measure, so subsequent product-space estimates can use an ordinary Gibbs expectation rather than an unnormalized weighted integral.

\subsection{\texorpdfstring{A finite cubic state moment}{A finite cubic state moment}}


The fourth-moment theorem gives the identity variable in \(L^4(\pi_U)\). On a probability space, \(L^4\subset L^3\), so the identity belongs to \(L^3\). This supplies the cubic moment needed for the rejection envelope.



On a probability space, Hölder gives \(\mathbb E\|X\|^3\le(\mathbb E\|X\|^4)^{3/4}\). The mass-one assumption is what removes any additional measure factor. The cubic moment is sufficient even though the Gaussian-tail argument conveniently supplied the stronger fourth moment.

\subsection{\texorpdfstring{The Gibbs-\allowbreak{}Gaussian pair measure}{The Gibbs-Gaussian pair measure}}


Take \(\pi_U\otimes\mathcal N(0,I)\), representing a Gibbs state and an independent standard Gaussian proposal innovation. This is the joint law under which the local rejection error is averaged.



For a product test \(f(x)g(\xi)\), expectation factors into the Gibbs expectation of \(f\) and Gaussian expectation of \(g\), whenever integrable. This encodes fresh-noise independence explicitly and permits the separate score and Gaussian moment estimates to be combined.

\subsection{\texorpdfstring{The cubic moment of proposal amplitude}{The cubic moment of proposal amplitude}}


The Gibbs state is in \(L^3\); Lipschitz growth \(\|\nabla U(x)\|\le\|\nabla U(0)\|+L\|x\|\) puts its score in \(L^3\). Gaussian noise also belongs to \(L^3\). Lift both to their independent product and add their norm envelopes to obtain \(A\in L^3\).



The gradient-growth estimate follows from \(\|\nabla U(x)-\nabla U(0)\|\le L\|x\|\). Minkowski in \(L^3\) then bounds the sum of the two amplitude components. No pointwise boundedness of the score is required; its polynomial growth is enough under the Gibbs moments.

\subsection{\texorpdfstring{Integrability of the amplitude cube}{Integrability of the amplitude cube}}


For nonnegative \(A\), its third power is \(\|A\|^3\). The \(L^3\) characterization gives integrability of this function. This is the finite ordinary integral used to define the rejection constant.



Since \(A\ge0\), \(|A|^3=A^3\). Thus the cubic moment defining the rejection constant is exactly the integral furnished by \(L^3\) membership, rather than a new moment condition. Product independence ensures both amplitude components have the required moments on the same space.

\subsection{\texorpdfstring{The rejection-\allowbreak{}error constant}{The rejection-error constant}}


Set \(C_{U,L}=L(2+L)\int A(x,\xi)^3d(\pi_U\otimes\gamma)\). The preceding lemma proves the integral is finite under the potential assumptions. All dependence on the positive step size has been removed from this constant.



The constant depends on the potential and dimension through the score and Gaussian moments, but not on the mesh. It is used to prove stationarity in a limiting argument, where finiteness and uniformity are essential and a sharp numerical value is unnecessary.

\subsection{\texorpdfstring{Nonnegativity of that constant}{Nonnegativity of that constant}}


The cubic amplitude integrand is nonnegative and its integral is nonnegative. Since \(L\ge0\), the factor \(L(2+L)\) is nonnegative as well. Their product gives \(C_{U,L}\ge0\).



If \(L=0\), the prefactor and the resulting rejection-error constant vanish. For positive \(L\), its sign is still determined only by nonnegative factors. This permits accumulation bounds to multiply the constant by a nonnegative step count without changing inequality direction.

\subsection{\texorpdfstring{Integrability of rejected displacement}{Integrability of rejected displacement}}


Continuity of the potential and gradient makes the rejection-displacement function measurable. Its absolute value equals itself by nonnegativity, and its pointwise bound is the integrable function \(h\sqrt h\,L(2+L)A^3\). Apply measurable domination to prove ordinary integrability.



The dominating function is finite in integral because \(A\in L^3\). Therefore measurable rejection costs have finite first moments even though their pointwise acceptance ratios may be unbounded. This is the precise integrability needed for an expected-distance coupling bound.

\subsection{\texorpdfstring{The expected rejection-\allowbreak{}displacement bound}{The expected rejection-displacement bound}}


Integrate the pointwise rejection envelope using the established integrability of both sides. Pull the step-independent factors out of the integral and recognize \(C_{U,L}\). This gives \(\mathbb ED_h\le C_{U,L}h\sqrt h\) for \(0<h\le1\).



For a mesh \(h=T/n\), the cumulative bound \(n C_{U,L}h^{3/2}=C_{U,L}T\sqrt h\) tends to zero. This algebra explains why the local exponent \(3/2\), rather than merely \(1\), suffices to transfer exact MALA invariance to the limiting diffusion.

\subsection{\texorpdfstring{Existence of a finite rejection constant}{Existence of a finite rejection constant}}


Choose the explicitly defined the finite rejection constant \(C\). Its nonnegativity and the integrated envelope theorem verify the existential statement for every \(0<h\le1\) simultaneously. The same constant works throughout the small-step limit.



The quantifier order is important: one finite constant works for every sufficiently small positive step, instead of choosing a different uncontrolled constant for each mesh. The explicit cubic envelope establishes this uniformity and makes the stationary-limit argument valid.

\section{\texorpdfstring{Detailed balance through accepted and rejected flows}{Detailed balance through accepted and rejected flows}}
\label{proof-module-UnadjustedLangevin-Metropolis-lean}


The Metropolis correction preserves the initial marginal whenever the accepted proposal flow is invariant under swapping its two states. A positive proposal-pair density makes this symmetry explicit.



An accepted proposal moves mass from \(x\) to \(y\); rejection leaves that mass at \(x\). Their weights add to one, which preserves the first marginal's total mass. For a target-weighted proposal density \(w(x,y)\), the accepted density is \(\min\{w(x,y),w(y,x)\}\), symmetric under exchanging the coordinates. Its outgoing and incoming marginals therefore coincide, while the rejected mass already remains in place. Together these facts prove invariance of the original marginal.

\subsection{\texorpdfstring{The accepted proposal flow}{The accepted proposal flow}}


Weight the current/\allowbreak{}proposal joint measure \(\nu(dx,dy)\) by acceptance probability \(a(x,y)\). Thus \(A(dx,dy)=a(x,y)\nu(dx,dy)\), with a nonnegative scalar embedding its real nonnegative weight. The bounds \(0\le a\le1\) are supplied in the conservation theorems below.



For a measurable test \(f\), the accepted contribution is \(\int f(y)a(x,y)\,d\nu(x,y)\). Thus the auxiliary pair measure retains exactly the information needed both to compute the output law and to compare forward and reverse flow.

\subsection{\texorpdfstring{The rejected proposal flow}{The rejected proposal flow}}


The rejected pair measure is \(R(dx,dy)=(1-a(x,y))\nu(dx,dy)\). It retains the proposed state as an auxiliary coordinate; the eventual output maps it to the original state \(x\), preserving all rejection mass.



On rejection, the proposed \(y\) is discarded, but its acceptance weight still determines how much mass stays at \(x\). Integrating \(f(x)(1-a(x,y))\) over the pair measure implements this dependence without conditioning or dividing by a rejection probability.

\subsection{\texorpdfstring{The output law of the Metropolis step}{The output law of the Metropolis step}}


Define the output as \(\operatorname{snd}_*A+\operatorname{fst}_*R\). Accepted moves output \(y\), while rejected moves output \(x\). This is the complete transition measure and requires no renormalization of accepted moves.



For every bounded measurable \(f\), the transition expectation is \(\int[a(x,y)f(y)+(1-a(x,y))f(x)]d\nu\). Setting \(f=1\) exposes mass conservation. This is a convex mixture at each pair, rather than normalization of the accepted part alone.

\subsection{\texorpdfstring{Acceptance and rejection recover total mass}{Acceptance and rejection recover total mass}}


Since \(0\le a\le1\), the two nonnegative weights add to one after embedding in extended scalars. Additivity of weighted-density measures then gives \(A+R=\nu\). The proof checks both nonnegativity guards before combining a nonnegative scalar values.



The identity holds as an equality of measures: integrating any nonnegative test gives \(\int f a\,d\nu+\int f(1-a)\,d\nu=\int f\,d\nu\). Nonnegativity permits this argument even before integrability of signed tests is available.

\subsection{\texorpdfstring{Equal marginals for a symmetric pair measure}{Equal marginals for a symmetric pair measure}}


The identity \(\operatorname{fst}\circ\operatorname{swap}=\operatorname{snd}\) and pushforward composition give \(\operatorname{snd}_*\nu=\operatorname{fst}_*(\operatorname{swap}_*\nu)\). If the pair measure is swap invariant, this is \(\operatorname{fst}_*\nu\), proving equality of its marginals.



Equivalently, swap invariance gives \(\int f(y)d\nu=\int f(x)d\nu\) for every bounded measurable \(f\). The statement concerns the entire marginal distributions, so it can be applied to accepted flow even though that flow is generally a subprobability measure.

\subsection{\texorpdfstring{Preservation of the initial marginal}{Preservation of the initial marginal}}


Detailed balance makes the two marginals of the accepted flow equal. Replace its second marginal by its first, then combine the accepted and rejected first marginals using \(A+R=\nu\). The output is therefore \(\operatorname{fst}_*\nu\), exactly the original state law.



For a test \(f\), detailed balance replaces the accepted expectation \(\int f(y)dA\) by \(\int f(x)dA\). Adding the rejected expectation now gives \(\int f(x)d(A+R)=\int f(x)d\nu\). This proves stationarity at the level of all test expectations.

\subsection{\texorpdfstring{Mapping a density-\allowbreak{}weighted measure}{Mapping a density-weighted measure}}


For a measurable \(f\) and weight \(w\), prove \(f_*((w\circ f)\mu)=w(f_*\mu)\). On each measurable set, the pushforward definition and nonnegative change-of-variables formula turn both sides into \(\int\mathbf1_A(f(x))w(f(x))d\mu(x)\). This is the weight-transport identity used by the actual MALA proposal map.



This formula is the weighted version of change of variables. No injectivity is needed because the weight depends on the image. It applies first to nonnegative tests and then, when absolutely integrable, to ordinary real or vector tests by the corresponding norm bound.

\subsection{\texorpdfstring{Mapping through an invertible change of coordinates}{Mapping through an invertible change of coordinates}}


Specialize the preceding pushforward/\allowbreak{}weight identity to a measurable equivalence. Its measurability supplies the required change-of-variables hypothesis. The application below uses the coordinate swap equivalence to prove symmetry of an accepted density.



For coordinate swapping, the inverse map is the swap itself. Therefore transporting a weight through the equivalence simply replaces \(w(x,y)\) by \(w(y,x)\). This makes the forward/\allowbreak{}reverse symmetry calculation an exact measure identity.

\subsection{\texorpdfstring{The Metropolis acceptance probability}{The Metropolis acceptance probability}}


For positive joint density \(w(x,y)\), define \(a(x,y)=\min(1,w(y,x)/w(x,y))\). The ratio is reverse flow divided by forward flow; positivity ensures both its meaning and nonnegative acceptance.



If reverse flow exceeds forward flow, acceptance is one. If reverse flow is smaller, multiplication by its ratio reduces the accepted forward mass to exactly the reverse mass. These two cases explain why the minimum rule enforces balance locally for each ordered pair.

\subsection{\texorpdfstring{The symmetric accepted weight}{The symmetric accepted weight}}


Distribute the positive factor \(w(x,y)\) through the minimum. Cancellation of the nonzero denominator gives \(w(x,y)a(x,y)=\min(w(x,y),w(y,x))\). This algebra is the reason the accepted joint density is symmetric.



The right side is unchanged on exchanging \(x\) and \(y\). Strict positivity is used only to cancel the forward density in the ratio; it is not an assumption that either direction is more likely than the other.

\subsection{\texorpdfstring{The density of accepted flow}{The density of accepted flow}}


Two successive changes of density multiply their weights. Combine the original positive density \(w\) and acceptance \(a\), then apply the weight/\allowbreak{}minimum identity. The accepted flow becomes the base measure weighted by \(\min(w,w\circ\operatorname{swap})\), a visibly symmetric density.



The accepted measure has density \(\min(w(x,y),w(y,x))\) relative to the swap-invariant base measure. Symmetry of the density and of that base jointly imply symmetry of the accepted measure; symmetry of the density alone would not suffice with an asymmetric base.

\subsection{\texorpdfstring{Invariance from the symmetric accepted density}{Invariance from the symmetric accepted density}}


The base measure is swap invariant, and the accepted minimum density is unchanged by swapping its arguments. The weight-transport theorem therefore proves accepted-flow detailed balance. The general accepted/\allowbreak{}rejected conservation theorem then identifies the output with the initial marginal.



This closes the invariance argument without any limiting calculation. The later diffusion proof will use this exact discrete invariance for every mesh and pass the unchanged marginal to the continuous limit after controlling the rejection error.

\section{\texorpdfstring{Couple ULA to an invariant Metropolis chain}{Couple ULA to an invariant Metropolis chain}}
\label{proof-module-UnadjustedLangevin-MALACoupling-lean}


A common Gaussian proposal couples unadjusted ULA with MALA. Under a stable small step, their expected distance grows by at most the Gibbs-average rejected displacement at each step. This is an auxiliary approximation for stationarity.



Use the same Gaussian proposal noise for both coordinates and one uniform accept/\allowbreak{}reject variable for the corrected coordinate. If accepted, compare the two common-noise proposals; if rejected, compare the unadjusted proposal with the retained state. The discrepancy is bounded by the previous distance plus the rejected displacement. Integrating gives an affine cost estimate, and iterating accumulates the expected rejection correction. Every transition remains a probability law because acceptance and rejection weights are nonnegative and sum to one.

\subsection{\texorpdfstring{The Metropolis acceptance function}{The Metropolis acceptance function}}


Define \(a_h(x,y)=\min(1,e^{R_h(x,y)})\), with the proved target/\allowbreak{}proposal log ratio. The minimum caps acceptance at one and the exponential ensures nonnegativity; these bounds justify the accepted/\allowbreak{}rejected measure split.



Acceptance is a measurable weight in \([0,1]\), so the two alternatives define nonnegative masses. For a pair already favorable in the forward/\allowbreak{}reverse comparison, the exponential ratio is at least one and the move is accepted with full weight.

\subsection{\texorpdfstring{Acceptance of the Langevin proposal}{Acceptance of the Langevin proposal}}


Compose \(a_h\) with \((x,\xi)\mapsto(x,T_h(x,\xi))\). The resulting \(a_h(x,T_h(x,\xi))\) is the acceptance weight on the original state/\allowbreak{}noise product space, prior to pushing to state/\allowbreak{}proposal pairs.



Using the original innovation variable preserves the Gaussian product representation. This is useful when estimating rejection displacement: the state and Gaussian moments can be controlled before changing variables from innovations to proposed positions.

\subsection{\texorpdfstring{Nonnegativity of acceptance}{Nonnegativity of acceptance}}


Both entries of the minimum, one and the exponential, are nonnegative. Their minimum is nonnegative. This direct order proof is used when combining real weights through a nonnegative scalar.



The lower bound remains valid when the log ratio is very negative: the exponential stays strictly positive even though it may be small. In combination with the upper bound, it guarantees that acceptance and rejection weights are both legitimate probabilities.

\subsection{\texorpdfstring{Acceptance is at most one}{Acceptance is at most one}}


The minimum \(\min(1,e^{R_h})\) is at most its first entry. This proves \(a_h\le1\), so the rejection weight \(1-a_h\) is nonnegative.



Together the two bounds yield \(a_h+(1-a_h)=1\). Therefore any mixture formed with these weights has the total mass of its input, including when its accepted part alone has mass strictly below one.

\subsection{\texorpdfstring{Measurability of the Gaussian proposal}{Measurability of the Gaussian proposal}}


The gradient is continuous by Lipschitzness. Projections, scalar multiplication, subtraction, and addition preserve continuity, so the joint proposal map \((x,\xi)\mapsto x-h\nabla U(x)+\sqrt{2h}\xi\) is measurable.



Joint measurability, rather than measurability for each frozen state alone, is needed to push forward a state/\allowbreak{}innovation product measure. Continuity of the full formula supplies that joint property directly and avoids a separate kernel measurability assumption.

\subsection{\texorpdfstring{Measurability of proposal acceptance}{Measurability of proposal acceptance}}


Differentiability makes \(U\) continuous, and Lipschitzness makes its gradient continuous. The log-ratio formula is assembled from continuous norms, inner products, and arithmetic; exponentiation and minimum preserve measurability. Composing with the measurable proposal gives the state/\allowbreak{}noise acceptance weight.



The operation \(\min(1,\cdot)\) is continuous, as is the exponential. Their composition therefore preserves the regularity already obtained for the log ratio. No discontinuous indicator of an acceptance event is introduced in the weighted transition representation.

\subsection{\texorpdfstring{One Gaussian ULA transition}{One Gaussian ULA transition}}


Define \(P_h\mu\) as the pushforward of \(\mu\otimes\gamma\) by the Euler proposal \(T_h\). The product measure encodes that the fresh Gaussian is independent of the current state. This transition is the actual independent-noise ULA step on laws.



For a bounded test \(f\), \(\int f\,d(P_h\mu)=\iint f(T_h(x,\xi))d\mu(x)d\gamma(\xi)\). This formula identifies exactly the usual fresh-Gaussian transition and is independent of any auxiliary coupling used to analyze its error.

\subsection{\texorpdfstring{One Metropolis-\allowbreak{}adjusted transition}{One Metropolis-adjusted transition}}


Split \(\mu\otimes\gamma\) into acceptance- and rejection-weighted measures. Map the accepted part to \(T_h(x,\xi)\), and the rejected part to \(x\), then add them. This defines one full MALA transition with all rejection mass retained.



The expectation formula is \(\iint[a_h f(T_h)+(1-a_h)f(x)]d\mu d\gamma\). It retains the full rejected mass at the original state. Consequently invariance can be proved through weighted pair symmetry, without adding an acceptance random variable to the notation.

\subsection{\texorpdfstring{The common-\allowbreak{}noise coupled transition}{The common-noise coupled transition}}


For a pair \((x,y)\), draw one shared Gaussian \(\xi\). ULA always moves to \(T_h(x,\xi)\); MALA moves to \(T_h(y,\xi)\) with weight \(a_h(y,T_h(y,\xi))\), otherwise stays at \(y\). The two weighted pushforwards define the explicit coupling measure.



On acceptance, the distance between outputs is controlled by contraction of the common-noise proposal. On rejection, insert and subtract the MALA proposal to separate the same contraction term from the distance lost by rejection. This is the origin of the one-step affine cost bound.

\subsection{\texorpdfstring{Acceptance and rejection weights sum to one}{Acceptance and rejection weights sum to one}}


For a measurable \(a\in[0,1]\), combine the two weighted-density measures using \(a+(1-a)=1\). Nonnegativity of both terms justifies the a nonnegative scalar addition identity, so their sum is exactly the original measure.



This identity can be tested against any nonnegative function, not just the constant one. It will therefore be used both to prove total mass and to combine two weighted expectations into the unweighted proposal contribution in the cost calculation.

\subsection{\texorpdfstring{The first marginal of the coupled transition}{The first marginal of the coupled transition}}


The first-coordinate update is identical on acceptance and rejection. Combine those two weights to one before pushing forward. The remaining product marginal is the original first marginal times a fresh Gaussian, so the first output law is \(P_h(\operatorname{fst}_*\rho)\).



The first component moves by ULA in both acceptance branches. Their weights add to one, so its marginal is the unweighted ULA transition of the first input marginal. It does not depend on the second component's acceptance decisions.

\subsection{\texorpdfstring{The second marginal of the coupled transition}{The second marginal of the coupled transition}}


Project both branches of the explicit coupling to the second state. They are exactly the accepted proposal and rejected current-state branches of MALA. Product marginal identities replace \(\rho\) by its second marginal, giving \(\operatorname{snd}_*\rho'=M_h(\operatorname{snd}_*\rho)\).



The second component accepts or rejects exactly according to its own current state and common Gaussian proposal. Integrating out the first state therefore leaves the ordinary MALA transition of the second input marginal, regardless of their dependence under the input coupling.

\subsection{\texorpdfstring{Normalization of the coupled transition}{Normalization of the coupled transition}}


The first marginal is a measurable pushforward of a probability product by the first-marginal identity for the coupled transition. It has mass one, and a measure and its first-coordinate pushforward have the same total mass. Therefore the joint coupling also has total mass one.



The product input has mass one and the two nonnegative branch weights sum to one. Pushforward preserves each branch mass, so their sum is a probability. This legitimizes interpreting all later first-moment cost integrals as expectations under a coupling.

\subsection{\texorpdfstring{The iterated coupled chain}{The iterated coupled chain}}


Initialize on the diagonal \((X_0,Y_0)=(X,X)\) with \(X\sim\mu\), then recursively apply the explicit coupled step. Equal starting states give zero initial transport cost, while subsequent first and second marginals follow ULA and MALA respectively.



Each step applies the same coupled transition to the current pair distribution. This construction retains the dependence needed for the common-noise comparison while allowing each marginal to follow its own algorithm. It is a sequence of joint laws, rather than an independence assertion about the pair.

\subsection{\texorpdfstring{Normalization at every coupled iterate}{Normalization at every coupled iterate}}


The diagonal pushforward of the initial probability is a probability. Induct on the number of transitions, using probability preservation of the common-noise coupled transition at each successor. This gives the normalized law needed for product noise and finite-cost arguments at every iteration.



Induction uses normalization of one coupled transition at each successor. The initial diagonal coupling is already a probability, so every finite-time joint distribution is normalized before any marginal or integrability argument is applied.

\subsection{\texorpdfstring{The first iterated marginal}{The first iterated marginal}}


At zero, diagonal initialization has first marginal \(\mu\). In the successor step apply the first-marginal identity for the coupled transition and the induction hypothesis. The resulting recursion is precisely the iteration \(P_h^n\mu\).



The one-step first-marginal identity closes the induction: if the first marginal at step \(k\) is the ULA law, its successor is the next ULA law. Thus the constructed coupling represents the actual numerical chain at every iteration.

\subsection{\texorpdfstring{The second iterated marginal}{The second iterated marginal}}


At zero the second marginal is \(\pi\). The successor marginal theorem gives \(M_h\) applied to the current second marginal; substitute the induction hypothesis and the supplied one-step invariance \(M_h\pi=\pi\). Thus the second coordinate remains stationary through the coupling.



The same induction applies to the second marginal with the MALA transition. When initialized at Gibbs and combined with detailed-balance invariance, this marginal remains Gibbs at every step, even while its first component evolves by unadjusted Euler proposals.

\subsection{\texorpdfstring{Nonexpansion of common-\allowbreak{}noise proposals}{Nonexpansion of common-noise proposals}}


Shared noise cancels in \(T_h(x,\xi)-T_h(y,\xi)\). The sharp strong-convexity/\allowbreak{}smoothness contraction theorem for the deterministic gradient map bounds the remaining norm by \(\|x-y\|\) when \(0\le h\le2/(m+L)\). No Gaussian moment is used in this pointwise comparison.



Use the sharp gradient-step estimate before taking expectations. In the admissible small-step regime its factor is at most one, and the common additive Gaussian cancels. This yields the pointwise nonexpansiveness required to accumulate rejection errors additively.

\subsection{\texorpdfstring{The pointwise coupled-\allowbreak{}step discrepancy}{The pointwise coupled-step discrepancy}}


The accepted distance is at most \(\|x-y\|\) by nonexpansiveness. On rejection, insert \(T_h(y,\xi)\) and use the triangle inequality to bound the distance by \(\|x-y\|+\|T_h(y,\xi)-y\|\). Weight the two branches; their common old-distance coefficients add to one, and the remaining term is the rejection displacement.



On rejection, \(\|T_h(x,\xi)-y\|\le\|T_h(x,\xi)-T_h(y,\xi)\|+\|T_h(y,\xi)-y\|\). On acceptance only the first term remains. Weighting both alternatives gives the old distance plus the rejection-weighted displacement at the second state.

\subsection{\texorpdfstring{The mean separation of a coupled pair}{The mean separation of a coupled pair}}


Define \(\mathcal C(\rho)=\int\|x-y\|d\rho\) as a nonnegative extended integral. Keeping it extended avoids assuming integrability before the cost recurrence proves it. Subsequent theorems establish finiteness and then convert it to an ordinary expected distance.



The cost here is a first moment, \(\int\|x-y\|\), suited to comparing characteristic functions. It is different from the square-root second-moment cost defining \(W_2\). The stationarity argument needs only that this first-moment separation tends to zero.

\subsection{\texorpdfstring{Integrating a mapped weighted measure}{Integrating a mapped weighted measure}}


Pull an extended nonnegative integral through a measurable pushforward, then through a weighted-density measure. Multiplication of the two nonnegative real factors becomes the embedded product \(a(x)b(f(x))\). This gives an exact integration formula for each branch of the coupling.



For a nonnegative test, the weighted pushforward change-of-variables formula is valid without a prior finiteness assumption. This permits an exact calculation of coupled cost before proving that its rejection envelope has a finite ordinary expectation.

\subsection{\texorpdfstring{The exact coupled-\allowbreak{}step cost formula}{The exact coupled-step cost formula}}


Apply the change-of-variables formula for the weighted pushforward to accepted and rejected branches of the coupled measure. Their cost is the integral of the acceptance-weighted proposal distance plus rejection-weighted stay distance over \(\rho\otimes\gamma\). The formula exposes exactly the pointwise expression bounded by the pointwise common-noise discrepancy inequality.



Expand the two weighted pushforwards against the distance test. The accepted distance is between the two proposals; the rejected distance is between the ULA proposal and the old MALA state. Both use the same Gaussian variable, so the common-noise estimate applies inside the integral.

\subsection{\texorpdfstring{A one-\allowbreak{}step cost bound}{A one-step cost bound}}


Integrate the pointwise coupling bound. The old distance integrates to \(\mathcal C(\rho)\); the error depends only on the second state and innovation. Its second marginal is Gibbs by hypothesis, so that integral is bounded by \(C_{U,L}h\sqrt h\). Thus \(\mathcal C(\rho')\le\mathcal C(\rho)+C_{U,L}h\sqrt h\), valid initially as an extended inequality.



After integration, the proposal-contraction term contributes at most the old expected distance. The rejection term depends only on the second marginal and fresh Gaussian noise; if that marginal is Gibbs, the uniform local bound contributes at most \(C_{U,L}h^{3/2}\).

\subsection{\texorpdfstring{The accumulated coupled-\allowbreak{}chain cost}{The accumulated coupled-chain cost}}


Induct from zero diagonal cost. The second marginal remains Gibbs by the supplied one-step MALA invariance, so the one-step cost bound applies at every iteration. Summing the same local error gives \(\mathcal C(\rho_n)\le nC_{U,L}h\sqrt h\).



Starting from a diagonal pair, the initial cost is zero. Repeatedly applying \(c_{k+1}\le c_k+b_h\) gives \(c_n\le n b_h\), with \(b_h=C_{U,L}h^{3/2}\). On a fixed horizon this vanishes as the mesh tends to zero.

\subsection{\texorpdfstring{Integrability of the coupled-\allowbreak{}chain distance}{Integrability of the coupled-chain distance}}


The quantitative extended cost estimate bounds \(\int\|x-y\|d\rho_n\) by a finite embedded real number. The distance is measurable and nonnegative, so finite extended integral is equivalent to ordinary integrability. This discharges the moment condition rather than assuming it.



The nonnegative cost has a finite extended integral by the preceding explicit bound. Since the distance is measurable, this is precisely finiteness of its ordinary first moment. The signed characteristic-function calculation can therefore use ordinary integral subtraction legitimately.

\subsection{\texorpdfstring{The ordinary integral cost estimate}{The ordinary integral cost estimate}}


The coupled distance is nonnegative and integrable by the preceding finite-cost argument. Its nonnegative cost integral therefore equals its ordinary expectation. The accumulated bound becomes \(\mathbb E_{\rho_n}\|X-Y\|\le nC_{U,L}h\sqrt h\), with every quantity finite.



Identifying the ordinary expected distance with the finite nonnegative cost transfers the same bound to real expectations. This is the last bridge needed for the characteristic-function estimate that passes Gibbs invariance to the continuous solution.

\section{\texorpdfstring{Identify Gaussian proposals with their radial density}{Identify Gaussian proposals with their radial density}}
\label{proof-module-UnadjustedLangevin-GaussianProposalDensity-lean}


For \(h>0\), the affine standard Gaussian proposal \(c+\sqrt{2h}\,Z\) has density proportional to \(e^{-\|y-c\|^2/(4h)}\). The proof identifies probability measures by characteristic functions and does not require a preselected coordinate formula for the Gaussian normalizer.



For \(h>0\), the proposal centered at \(a\) has density \(q_h(a,y)=(4\pi h)^{-d/2}\exp(-\|y-a\|^2/(4h))\). Its normalization is obtained from the radial Gaussian integral in \(d\) dimensions. Translation changes its center without changing the integral. Setting \(a=x-h\nabla U(x)\) identifies the ULA transition. Fourier or characteristic-function calculations verify that this density describes the same affine Gaussian law as the independent-noise update.

\subsection{\texorpdfstring{The radial Gaussian integral}{The radial Gaussian integral}}


Set \(J(a)=\int_H e^{-a\|z\|^2}dz\) as an ordinary real integral. The positive-precision results below prove its integrability and positivity before using \(J(a)^{-1}\) as a normalizing constant.



The radial integral is evaluated in the Euclidean volume of the state space. A positive precision gives a product of finite scalar Gaussian integrals, retaining its dimension dependence.

\subsection{\texorpdfstring{The radial normalization constant}{The radial normalization constant}}


Define \(C(a)=J(a)^{-1}\). This is an actual integral normalizer on the given finite-dimensional Hilbert space, avoiding a coordinate-dependent choice of a formula for Lebesgue density.



The normalizer is the reciprocal of the positive finite radial integral. This choice is fixed before translations of the kernel are considered.

\subsection{\texorpdfstring{Positivity of the radial integral}{Positivity of the radial integral}}


For \(a>0\), Gaussian integrability was established in integrability of a radial Gaussian with positive precision. The integrand is strictly positive, so its ordinary integral is strictly positive by positivity of the integral of an everywhere positive continuous weight.



The radial integrand is strictly positive and continuous. Its integral over any nonempty ball is positive, establishing a nonzero denominator for normalization.

\subsection{\texorpdfstring{Positivity of the normalizer}{Positivity of the normalizer}}


The inverse of a strictly positive real number is strictly positive. Apply this to \(J(a)\), whose positivity was proved immediately above. This also proves the normalizer is nonzero for later cancellation.



The reciprocal of a positive finite scalar is positive. Thus normalization preserves strict positivity of every Gaussian density value.

\subsection{\texorpdfstring{The normalized radial density}{The normalized radial density}}


Define \(q_{a,c}(z)=C(a)e^{-a\|z-c\|^2}\), embedded in extended nonnegative scalars. The center changes the translated argument, while the normalizer depends only on precision and dimension.



The density is the radial exponential multiplied by its normalizing scalar. Translation to a nonzero mean later changes the exponent's center but leaves this scalar unchanged.

\subsection{\texorpdfstring{Integrability of the radial density}{Integrability of the radial density}}


Translation by \(c\) preserves Lebesgue measure, so \(e^{-a\|z-c\|^2}\) is integrable whenever \(a>0\). Multiply by the finite real constant \(C(a)\) to obtain integrability of the displayed real density.



The Gaussian exponential is integrable and the normalizer is finite. Scalar multiplication therefore preserves integrability of the normalized density.

\subsection{\texorpdfstring{Its unit integral}{Its unit integral}}


Translate \(z\mapsto z-c\) in the integral and pull out \(C(a)\). The result is \(C(a)J(a)=1\), justified by \(J(a)>0\). Thus normalization is independent of the center.



The chosen reciprocal cancels the radial integral, giving total integral one. Positivity and finiteness justify that cancellation in ordinary real arithmetic.

\subsection{\texorpdfstring{The radial Gaussian probability law}{The radial Gaussian probability law}}


The total mass of the weighted-density measure is the extended integral of its density. Nonnegativity and the proved real integrability identify it with a nonnegative scalar of the real integral, which is one. This constructs probability normalization, including finiteness.



A measurable nonnegative density of integral one defines a probability measure. This identifies the radial Gaussian density as an actual normalized law.

\subsection{\texorpdfstring{The complex Gaussian Fourier integral}{The complex Gaussian Fourier integral}}


For real \(a\), \(e^{-(a:\mathbb C)\|z\|^2}\) is the complex embedding of the real Gaussian exponential. Rewrite the integrand with a nonnegative scalar exp; the continuous real embedding commutes with integration, giving the complex embedding of \(J(a)\).



Completing the square in the Gaussian Fourier integral introduces the familiar quadratic frequency decay. The absolute integrability of the original Gaussian justifies its complex-valued integration.

\subsection{\texorpdfstring{The Gaussian character integral}{The Gaussian character integral}}


Apply the finite-dimensional complex Gaussian Fourier integral with imaginary linear coefficient. Combining the two exponentials and using \(i^2=-1\) gives \(\int e^{-a\|z\|^2}e^{i\langle z,t\rangle}dz=J(a)e^{-\|t\|^2/(4a)}\). The zero-frequency Gaussian formula identifies the normalization factor with the same real integral \(J(a)\).



The character is \(e^{i\langle v,x\rangle}\). Inserting it into the normalized radial integral gives the characteristic expression with the corresponding isotropic covariance.

\subsection{\texorpdfstring{The characteristic function of the radial law}{The characteristic function of the radial law}}


Change from the density measure to a weighted Lebesgue integral. Translate \(z=u+c\); the character splits into \(e^{i\langle u,t\rangle}e^{i\langle c,t\rangle}\). The Fourier formula gives a factor \(J(a)\), canceled by \(C(a)\), leaving \(e^{-\|t\|^2/(4a)}e^{i\langle c,t\rangle}\).



The computed characteristic function matches that of the designated centered Gaussian. Uniqueness of characteristic functions then identifies the radial density law with that Gaussian measure.

\subsection{\texorpdfstring{The proposal normalization constant}{The proposal normalization constant}}


Specialize the precision to \(a=1/(4h)\). The resulting constant normalizes a Gaussian with covariance \(2hI\). Its independence of the current state is essential when the same factor cancels in the forward and reverse Metropolis ratio.



For proposal variance \(2h\), the normalizer becomes \((4\pi h)^{-d/2}\). Positive \(h\) is necessary for both this power and the variance interpretation.

\subsection{\texorpdfstring{The Gaussian proposal density}{The Gaussian proposal density}}


Define \(q_h(c,y)=C(1/(4h))e^{-\|y-c\|^2/(4h)}\). For \(h>0\), this is the density of a covariance-\(2hI\) Gaussian centered at \(c\). The density identification is proved later rather than built into the definition.



The proposal exponent is the squared distance from its deterministic center divided by \(4h\). Its density is therefore a translation of the normalized radial Gaussian.

\subsection{\texorpdfstring{Positivity of the proposal normalizer}{Positivity of the proposal normalizer}}


For \(h>0\), the precision \(1/(4h)\) is positive. Apply positivity of the radial Gaussian normalizer at that precision. This nonzero constant is used in detailed-balance cancellation.



Positive time makes the Gaussian normalizer strictly positive. Combined with the positive exponential, this yields positivity of the whole proposal kernel.

\subsection{\texorpdfstring{Equality of radial and proposal density formulas}{Equality of radial and proposal density formulas}}


Unfold the two density definitions and the proposal normalizer. Their exponents agree by \(-(1/(4h))\|y-c\|^2=-\|y-c\|^2/(4h)\). Function extensionality gives equality of the density functions, with no probabilistic argument needed.



Substitution of the translated center identifies the two density formulas pointwise. The equality retains the same Euclidean norm and the same normalization constant.

\subsection{\texorpdfstring{The density of an affine Gaussian proposal}{The density of an affine Gaussian proposal}}


Compute characteristic functions on both sides. The affine image \(c+\sqrt{2h}Z\) has character \(e^{-h\|t\|^2}e^{i\langle c,t\rangle}\); the radial-density formula with \(a=1/(4h)\) gives the same expression. Both measures are probabilities, so uniqueness of a probability measure from its characteristic function identifies them. The square-root identity uses \(h>0\).



The affine noise map has Gaussian mean equal to the proposal center and covariance \(2hI_d\). Its characteristic function agrees with that of the explicit proposal density.

\subsection{\texorpdfstring{The radial representation of the proposal law}{The radial representation of the proposal law}}


Rewrite the proposal density using the equality between the translated radial Gaussian and proposal density formulas in the established affine-Gaussian law equality. This yields the equivalent precision-parameter form used in the detailed-balance proof.



Rewriting the proposal as a translated radial law separates the shape of the Gaussian from its center. This form is useful for uniform bounds independent of the starting state.

\subsection{\texorpdfstring{The radial representation of the ULA proposal}{The radial representation of the ULA proposal}}


Set the affine Gaussian center to \(c=x-h\nabla U(x)\). The fixed-state density theorem then gives \(x-h\nabla U(x)+\sqrt{2h}Z\) its exact radial density. No regularity of \(U\) is required because \(x\) is fixed throughout this statement.



For ULA the center is \(x-h\nabla U(x)\). Substituting this center into the radial representation gives the actual update law, with its drift dependence entirely in the translation.

\subsection{\texorpdfstring{The density of the ULA transition}{The density of the ULA transition}}


Expand the radial density at center \(x-h\nabla U(x)\). The exponent becomes \(-\|y-x+h\nabla U(x)\|^2/(4h)\), with the same positive normalizer. This is the explicit ULA transition density used in the Metropolis acceptance ratio.



The explicit density formula therefore describes the same transition as the independent Gaussian update. This law identification supports both mixture-density calculations and Metropolis proposal-ratio computations.

\section{\texorpdfstring{The joint initial-\allowbreak{}state and proposal density}{The joint initial-state and proposal density}}
\label{proof-module-UnadjustedLangevin-GaussianProposalPairDensity-lean}


For a measurable center \(c(x)\), independent \(X\sim\mu\) and \(Z\sim\mathcal N(0,I)\) produce the joint law \(\mu(dx)q_h(c(x),y)dy\).



Begin with an initial law \(\mu(dx)\) and an independent Gaussian noise. Conditional on \(x\), the affine change of variables sends the noise to the proposal density \(q_h(x,y)dy\). Fubini therefore gives \(\mathbb EF(X,Y)=\iint F(x,y)q_h(x,y)\mu(dx)dy\) for nonnegative measurable \(F\), and for integrable signed or vector-valued \(F\). This identifies the joint law before any posterior or score calculation uses it.

\subsection{\texorpdfstring{Disintegrate the Gaussian change of variables}{Disintegrate the Gaussian change of variables}}


For \(X\sim\mu\) and independent \(Z\sim\mathcal N(0,I)\), fix \(x\) and apply the proved proposal-density identity to \(c(x)+\sqrt{2h}Z\). It gives conditional integration against \(q_h(c(x),y)dy\).



To obtain equality of joint measures, evaluate a measurable set \(A\subset H\times H\). Tonelli writes both sides as integrals over \(x\) of their \(y\)-sections. The fixed-center density theorem identifies each section, hence the measures agree. Measurability of \(c\) supplies measurability of the pair update and joint density; no probability or moment condition on \(\mu\) is added.



For each starting state the conditional proposal has integral one. Integrating that kernel against \(\mu\) gives the joint pair measure with first marginal \(\mu\); integrating out the first coordinate gives the proposal mixture. The identity is first valid for nonnegative tests, then for absolutely integrable tests by positive-negative or norm domination.

\section{\texorpdfstring{Exact Gibbs invariance under the Metropolis-\allowbreak{}adjusted step}{Exact Gibbs invariance under the Metropolis-adjusted step}}
\label{proof-module-UnadjustedLangevin-MALAInvariance-lean}


The Gaussian proposal density identifies the MALA pair flow. Its reversed-to-forward density ratio is precisely the exponential of the implemented log ratio, so the Metropolis balance theorem proves exact Gibbs invariance.



The joint measure of a Gibbs starting state and its Langevin proposal has density \(Z^{-1}e^{-U(x)}q_h(x,y)\). Dividing its swapped density by itself gives the usual Metropolis ratio. Multiplying by acceptance leaves the symmetric density \(\min\{Z^{-1}e^{-U(x)}q_h(x,y),Z^{-1}e^{-U(y)}q_h(y,x)\}\). The generic accepted-flow calculation therefore applies to this actual proposal, proving the Gibbs marginal invariant for one MALA step.

\subsection{\texorpdfstring{The joint starting-\allowbreak{}proposal law}{The joint starting-proposal law}}


Push \(\mu\otimes\gamma\) through \((x,\xi)\mapsto(x,T_h(x,\xi))\). Its first coordinate retains the current state, and its second records the actual Gaussian proposal. This joint measure is the input to the accepted-flow detailed-balance theorem.



The first marginal is the starting law because the Gaussian conditional proposal has mass one. Its second marginal is the unadjusted proposal law. Keeping both coordinates makes forward and reverse density comparisons available on the same product space.

\subsection{\texorpdfstring{The Gibbs-\allowbreak{}weighted proposal density}{The Gibbs-weighted proposal density}}


Set \(w_h(x,y)=\exp(-U(x)-\|y-x+h\nabla U(x)\|^2/(4h))\). The Gibbs and Gaussian normalization constants are kept in a separate symmetric base measure, so \(w_h\) is strictly positive and its reverse/\allowbreak{}forward ratio is easy to compute.



With respect to Lebesgue product measure, the pair density is the Gibbs density at \(x\) multiplied by the proposal density from \(x\) to \(y\). Both factors are positive, so reverse/\allowbreak{}forward ratios can be formed pointwise and canceled without exceptional zero-density cases.

\subsection{\texorpdfstring{The forward-\allowbreak{}reverse density ratio}{The forward-reverse density ratio}}


Divide the two positive exponentials and use \(e^a/e^b=e^{a-b}\). Expanding the definition of \(w_h\) shows the exponent difference is exactly the logarithmic metropolis ratio. Thus \(w_h(y,x)/w_h(x,y)=e^{R_h(x,y)}\).



Taking logarithms turns multiplication of target and proposal factors into the potential and quadratic terms defining the log ratio. Exponentiating recovers the exact reverse/\allowbreak{}forward joint-density ratio, so the acceptance used in the coupling is precisely the detailed-balance acceptance.

\subsection{\texorpdfstring{The Metropolis representation of the actual transition}{The Metropolis representation of the actual transition}}


Transport acceptance and rejection weights through the measurable pair map using composition of a measurable map with a weighted pushforward. Projecting the accepted pair to its second coordinate gives the proposal; projecting rejected pairs to the first gives the unchanged state. These are precisely the two branches defining one metropolis-adjusted transition.



Changing variables from Gaussian innovations to proposed positions commutes with multiplication by the corresponding acceptance weight. Consequently the innovation-based transition and the pair-density Metropolis transition have identical accepted and rejected components, not merely the same total mass.

\subsection{\texorpdfstring{The joint proposal density at Gibbs}{The joint proposal density at Gibbs}}


Apply the Gaussian joint-density theorem with center \(x-h\nabla U(x)\), then multiply its density by the Gibbs density. Combine the exponential factors into \(w_h(x,y)\). Pull the two state-independent normalizers into the base measure \(Z_U^{-1}C_h(dx\,dy)\), which is invariant under coordinate swap.



Disintegrating the state/\allowbreak{}proposal law gives the displayed density against Lebesgue product measure. Tonelli is available for nonnegative tests; the Gaussian density normalization verifies the conditional mass and the target normalization verifies the total mass.

\subsection{\texorpdfstring{Invariance of Gibbs under the actual adjusted step}{Invariance of Gibbs under the actual adjusted step}}


The joint proposal measure has a positive density \(w_h\) against a symmetric base. Its acceptance rule is \(\min(1,w_h(y,x)/w_h(x,y))\) by the ratio between the reverse and forward target-weighted proposal densities. Hence its accepted density is the symmetric minimum \(\min(w_h(x,y),w_h(y,x))\).



Apply preservation of the target marginal by the symmetric accepted density and the exact step/\allowbreak{}output identity. The initial marginal of the proposal pair is Gibbs, since projecting \(\pi_U\otimes\gamma\) to its state returns \(\pi_U\). Thus \(M_h\pi_U=\pi_U\) is proved with no stationarity input.



The accepted density is \(\min(\pi(x)q_h(x,y),\pi(y)q_h(y,x))\), symmetric in \((x,y)\). Its two marginals coincide. Rejected mass stays at \(x\), and accepted plus rejected input mass is Gibbs, proving exact invariance for each positive admissible step.

\section{\texorpdfstring{Characteristic functions are continuous in L\ensuremath{{}^{2}}}{Characteristic functions are continuous in L²}}
\label{proof-module-UnadjustedLangevin-LpCharacteristic-lean}


At each fixed frequency, the characteristic function of the law of a real \(L^2\) variable depends continuously on that variable. Consequently centered Gaussian laws are closed under \(L^2\) convergence, with variance given by the squared norm.



For a scalar random variable \(X\), the characteristic test is \(e^{iX}\); the centered version subtracts its constant value at zero. The elementary estimate \(|e^{ia}-e^{ib}|\le|a-b|\), followed by Cauchy–Schwarz on a probability space, controls the difference of expectations by \(\|X-Y\|_2\). Thus convergence in \(L^2\) transfers characteristic functions. Gaussian characteristic functions also vary continuously with their mean and variance, allowing Gaussian laws to pass through the same limit.

\subsection{\texorpdfstring{The centered characteristic test}{The centered characteristic test}}


Set \(c_t(x)=e^{itx}-1\), represented by the unit-radius circle map centered at \(-1\). Subtracting one makes \(c_t(0)=0\), which is the condition needed by the Lipschitz composition construction on \(L^2\).



Subtracting the constant value at zero makes the test vanish at zero. This permits a norm estimate directly from the size of the random input without altering differences of characteristic expectations.

\subsection{\texorpdfstring{The Lipschitz bound for the centered character}{The Lipschitz bound for the centered character}}


Multiplication by \(t\) is \(|t|\)-Lipschitz, and the unit circle parametrization is \(1\)-Lipschitz. Composing gives \(|c_t(x)-c_t(y)|\le|t||x-y|\). The proof identifies the nonnegative Lipschitz constants before applying the composition theorem.



The scalar inequality \(|e^{ia}-e^{ib}|\le|a-b|\) integrates to an \(L^1\) bound. On a probability space, Cauchy–Schwarz bounds that first moment by the \(L^2\) norm.

\subsection{\texorpdfstring{The centered character at zero}{The centered character at zero}}


At zero the exponential is one, hence \(c_t(0)=e^0-1=0\). Unfolding the circle-map definition makes this a scalar identity; this vanishing proof permits composition with a Lipschitz map vanishing at zero to lift the map without an integrability assumption on a nonzero constant.



At a zero input the exponential is one, so the centered test is identically zero. Its expectation therefore vanishes exactly, agreeing with the Lipschitz normalization at the origin.

\subsection{\texorpdfstring{The characteristic random variable in the square-\allowbreak{}integrable space}{The characteristic random variable in the square-integrable space}}


Lift \(c_t\) to \(L^2\) by Lipschitz composition, then add the constant class one. Thus \(e^{itX}\) represents \(e^{itX}\). The probability assumption makes the constant function square integrable.



The expectation is well defined because the complex exponential has modulus one. Thus characteristic expectations remain meaningful even without moments; the \(L^2\) structure is used to quantify their continuity.

\subsection{\texorpdfstring{Its representative is the ordinary exponential}{Its representative is the ordinary exponential}}


Combine the almost-everywhere representative identities for lifted composition, \(L^2\) addition, and the constant class. Their intersection has full measure; there the representative is \((e^{itX}-1)+1=e^{itX}\). This connects quotient-space operations with the random function used in the characteristic integral.



Evaluating the expectation on a representative gives the usual characteristic integral. A change on a null set leaves the integral unchanged, so this expression depends only on the square-integrable class.

\subsection{\texorpdfstring{Recovering the characteristic function of a law}{Recovering the characteristic function of a law}}


The pushforward integral theorem gives \(\widehat{P_X}(t)=\int e^{itX}dP\). In complex \(L^2\), inner product with the constant one has exactly this integral because its conjugate is one. The representative identity above therefore yields \(\widehat{P_X}(t)=\langle1,e^{itX}\rangle\).



For a vector frequency \(v\), the scalar input is \(\langle v,X\rangle\). Pushing the law of \(X\) through this linear functional converts the vector characteristic function into the scalar characteristic expectation.

\subsection{\texorpdfstring{Continuity of characteristic functions under square-\allowbreak{}mean convergence}{Continuity of characteristic functions under square-mean convergence}}


Lipschitz composition is continuous on \(L^2\); adding the constant one preserves continuity. Pairing with one is a continuous linear functional. The preceding integral identity rewrites their composition as \(X\mapsto\widehat{P_X}(t)\), proving continuity of the actual-law characteristic function at every fixed frequency.



The characteristic difference is at most \(\|v\|\|X-Y\|_2\). Hence \(L^2\) convergence of random vectors yields convergence of their characteristic functions at every fixed frequency.

\subsection{\texorpdfstring{Continuity of Gaussian characteristic parameters}{Continuity of Gaussian characteristic parameters}}


The centered Gaussian characteristic function is \(e^{-t^2\|X\|_2^2/2}\). The norm, its square, and the scalar exponential are continuous, so this characteristic expression varies continuously with the square-integrable class \(X\). Its variance parameter is the actual squared norm, which is nonnegative.



The Gaussian characteristic function is \(e^{it m-vt^2/2}\). Ordinary continuity of the exponential transfers convergence of the mean and variance parameters into convergence of this expression.

\subsection{\texorpdfstring{Preservation of an eventual Gaussian law in the limit}{Preservation of an eventual Gaussian law in the limit}}


At each frequency, \(L^2\) convergence makes the actual characteristic functions converge. It also makes the centered Gaussian characteristic functions with variances \(\|X_i\|_2^2\) converge to variance \(\|Y\|_2^2\). Eventual Gaussianity identifies the two nets eventually, so uniqueness of limits identifies their limits. uniqueness of a probability measure from its characteristic function then gives the claimed Gaussian law of \(Y\); a nontrivial filter is required for uniqueness.



Along the tail on which the variables have the stated Gaussian laws, characteristic functions converge both by \(L^2\) continuity and by the Gaussian parameter formula. Uniqueness identifies the limiting law.

\subsection{\texorpdfstring{Preservation of a Gaussian law under square-\allowbreak{}mean convergence}{Preservation of a Gaussian law under square-mean convergence}}


If every approximant has the stated Gaussian law, then the law identity holds eventually for any filter. Feed the fact that a property holding at every index also holds eventually to the preceding closure theorem. This version is suited to finite orthogonal sums converging to an infinite Gaussian series.



If every approximating variable has the specified Gaussian law, the preceding eventual argument applies immediately. The conclusion concerns the limiting distribution, without requiring pointwise convergence of sample values.

\section{\texorpdfstring{A vanishing coupling distance determines the limiting law}{A vanishing coupling distance determines the limiting law}}
\label{proof-module-UnadjustedLangevin-CouplingLimit-lean}


If \(X_n\to X\) in \(L^2\) and each \(\operatorname{Law}(X_n)\) couples to a fixed probability \(\pi\) with mean distance tending to zero, then \(\operatorname{Law}(X)=\pi\).



For a frequency \(v\), the character \(x\mapsto e^{i\langle v,x\rangle}\) changes by at most \(\|v\|\|x-y\|\) between two points. Integrating under a coupling bounds the difference of characteristic functions by \(\|v\|\) times its first-moment cost. Thus couplings with costs tending to zero force agreement of all limiting characteristic functions. Uniqueness of characteristic functions identifies the probability measures, without requiring convergence of unbounded polynomial tests.

\subsection{\texorpdfstring{A Lipschitz bound for characteristic tests}{A Lipschitz bound for characteristic tests}}


The unit-circle map is \(1\)-Lipschitz, so the difference of the two characters is at most \(|\langle x-y,t\rangle|\). Cauchy–Schwarz then bounds this by \(\|t\|\|x-y\|\). Thus characteristic tests have an explicit Lipschitz constant at each fixed frequency.



For real numbers, \(|e^{iu}-e^{iv}|\le|u-v|\), obtained by integrating the derivative of the unit-circle map along the interval. Substitute \(u=\langle x,t\rangle\), \(v=\langle y,t\rangle\), then use Cauchy–Schwarz. Thus every fixed-frequency characteristic test is bounded and \(\|t\|\)-Lipschitz, linking a mean coupling distance to a difference of characteristic functions.

\subsection{\texorpdfstring{Characteristic differences controlled by coupling cost}{Characteristic differences controlled by coupling cost}}


Pull both marginal characteristic integrals back to the coupling space and subtract. Each character is bounded by one, hence integrable under the probability coupling. The integral norm inequality and the pointwise Lipschitz estimate give \(|\widehat\mu(t)-\widehat\pi(t)|\le\|t\|\mathbb E\|X-Y\|\), with the distance's integrability explicitly supplied.



The coupling's marginal identities let both characteristic integrals be evaluated on the same product-state space. Subtract under the integral, use its norm inequality, and apply the pointwise Lipschitz bound. The assumed integrability of distance makes the resulting right side finite. This gives a quantitative estimate for each frequency without requiring a second-moment transport bound on the coupling itself.

\subsection{\texorpdfstring{Continuity of the characteristic law in square mean}{Continuity of the characteristic law in square mean}}


Apply the continuous linear functional \(x\mapsto\langle t,x\rangle\) to an \(H\)-valued \(L^2\) variable. Its lifted map is continuous on \(L^2\). The Hilbert-valued characteristic function at \(t\) is the scalar projected characteristic function at frequency one, so the scalar continuity theorem proves the claim.



The scalar projection \(X\mapsto\langle t,X\rangle\) is continuous in \(L^2\) with norm at most \(\|t\|\). Also \(\mathbb E|\langle t,X-Y\rangle|\le\|t\|\|X-Y\|_2\). Applying the exponential Lipschitz inequality shows continuity of its characteristic expectation. This explains directly why \(L^2\) convergence of the approximating endpoints determines convergence of their characteristic functions.

\subsection{\texorpdfstring{Equality of limiting laws from vanishing coupling cost}{Equality of limiting laws from vanishing coupling cost}}


For each \(t\), the coupling inequality forces the approximating characteristic functions toward \(\widehat\pi(t)\), since the expected distance tends to zero. Independently, \(L^2\) convergence forces them toward \(\widehat{P_X}(t)\). Uniqueness of limits along a nontrivial filter identifies these values, and uniqueness of a probability measure from its characteristic function yields \(P_X=\pi\).



There are two limits of the same scalar characteristic sequence. The \(L^2\) convergence identifies its limit with the characteristic function of the continuous endpoint, while the coupling estimate identifies it with that of \(\pi\). Uniqueness of scalar limits makes them equal at every frequency. The characteristic-function uniqueness theorem for finite-dimensional probability measures then gives equality of the entire endpoint law and \(\pi\).

\section{\texorpdfstring{From scalar coordinate laws to independent Gaussian vectors}{From scalar coordinate laws to independent Gaussian vectors}}
\label{proof-module-UnadjustedLangevin-GaussianInnovationVectors-lean}


For finitely many spatial coordinates, independent standard normal scalar innovations assemble into independent standard Gaussian vectors. Scalar normal laws without joint independence would not suffice for this conclusion.



A vector innovation is assembled from its \(d\) scalar coordinates at a fixed index. Independence of the scalar array identifies its coordinate law with a product of standard one-dimensional Gaussians, which is the Euclidean standard Gaussian. Grouping the array by time index preserves independence of the resulting vectors. The conclusion therefore includes both the joint law within a vector and independence between successive vectors.

\subsection{\texorpdfstring{The law of the scalar innovation array}{The law of the scalar innovation array}}


Regroup the jointly independent scalar coordinates indexed by \((n,i)\) into vector blocks \(W_n\). A finite block with independent standard Gaussian coordinates has the standard Euclidean Gaussian law. Product-measure identification and reindexing then show that the full vector path the stated pushforward law \(\bigotimes_n\mathcal N(0,I)\).



The path-level product law specifies every finite array of innovation coordinates. It therefore determines the joint law of their sequences, including dependence across time.

\subsection{\texorpdfstring{The standard Gaussian vector law}{The standard Gaussian vector law}}


Project the established whole-path product law to its \(n\)-th vector coordinate. Product marginals give \(P\circ W_n^{-1}=\mathcal N(0,I)\). This derives the multivariate law from the scalar coordinate assumptions rather than adding it as an algorithm hypothesis.



At a fixed index, the independent scalar standard Gaussian coordinates form the Euclidean standard Gaussian vector. The product characterization identifies the actual vector law.

\subsection{\texorpdfstring{Independence of the vector innovations}{Independence of the vector innovations}}


The coordinate maps of the vector product law are independent. Pull that independence back through the actual innovation-path map using the product law of the complete scalar innovation array. Thus the vector-valued \(W_n\) are independent over time, not just uncorrelated.



Disjoint time-index groups of an independent scalar array remain independent after assembly into vectors. Hence the vector innovations form an independent family with the required common Gaussian law.

\section{\texorpdfstring{Prove Gibbs stationarity by refining invariant Metropolis chains}{Prove Gibbs stationarity by refining invariant Metropolis chains}}
\label{proof-module-UnadjustedLangevin-GibbsStationarity-lean}


A Brownian Langevin integral solution started from Gibbs stays Gibbs. The proof uses convergent Euler approximations and exactly invariant MALA chains, with no weak-PDE uniqueness assumption.



The invariant MALA marginal and the unadjusted Euler approximation are compared on a common noise construction. The rejection-displacement estimate controls their accumulated discrepancy on a fixed horizon as the mesh shrinks. The Euler approximation also converges in \(L^2\) to the integral solution. Bounded characteristic tests transfer these two vanishing errors to equality between the solution's marginal and Gibbs. This establishes stationarity through the actual approximation construction rather than through a separately assumed weak-equation uniqueness theorem.

\subsection{\texorpdfstring{The representative of the Hilbert-\allowbreak{}space chain}{The representative of the Hilbert-space chain}}


Induct over the chain recursion. The representative of the lifted gradient equals the pointwise gradient almost everywhere, and the representative formulas for addition and scalar multiplication give the same Euler successor. Intersect the finite set of null-set complements needed at each index to identify the \(L^2\) chain with the ordinary recursion.



At time zero both descriptions use the same initialization. For the successor, the composed score represents the pointwise gradient almost surely, and addition and scalar multiplication represent the displayed Euler arithmetic. Intersecting these full-measure identities with the induction event proves agreement at the next finite index. This establishes a concrete link between analytic chain estimates and ordinary state distributions.

\subsection{\texorpdfstring{Its law follows the Gaussian transition}{Its law follows the Gaussian transition}}


Use the almost-everywhere pointwise-chain identification, the independent-iid realization law theorem, and the canonical path/\allowbreak{}Markov-law identity. With innovation law \(\mathcal N(0,I)\), the generic independent-noise step is exactly the transition \(\mu\mapsto T_{h\#}(\mu\otimes\gamma)\). Thus the \(L^2\) construction the stated pushforward law \(P_h^n(P_{X_0})\).



The innovation assumptions identify the pointwise recursion law with the canonical iid input law, and the path-to-transition theorem identifies the latter with iterated Gaussian Euler transitions. Almost-sure equality transfers this law to the square-integrable chain. Its distribution is therefore the actual \(P_h^n\mu_0\), which allows transition-level invariant-chain comparisons to be applied to the constructed approximation.

\subsection{\texorpdfstring{Identification of the additive Euler chain}{Identification of the additive Euler chain}}


Induct on the mesh index. Both chains start at the same \(L^2\) variable. Their successor equations agree because the supplied additive increment equals \(\sqrt{2h}W_k\), making the noise terms identical as classes.



The additive Euler approximation uses increments of an input path, while the normalized-noise chain multiplies innovations by \(\sqrt{2h}\). When these two increments agree, their successor states agree whenever the preceding states agree. Induction thus identifies the two approximations exactly. This permits convergence results proved for additive-input Euler approximation to be combined with Gaussian transition-law theorems.

\subsection{\texorpdfstring{The grid increment of the Brownian input}{The grid increment of the Brownian input}}


Apply the Brownian input cocycle at \(a=kh\), \(t=h\). The elementary identity \((k+1)h=kh+h\) matches its endpoints to the grid notation. This supplies the exact additive increment used by the normalized Euler noise.



The input increment between grid points is \(N_0((k+1)h)-N_0(kh)=\sqrt2(B_{(k+1)h}-B_{kh})\). The cocycle identity rewrites it as \(N_{kh}(h)\). This is precisely the segment endpoint input that defined the normalized Euler noise. The algebra aligns the mesh convention and absolute Brownian times without introducing a discretization error.

\subsection{\texorpdfstring{The Brownian Euler chain equals the square-\allowbreak{}integrable chain}{The Brownian Euler chain equals the square-integrable chain}}


The Brownian grid-increment identity and the equality between the scaled innovation and Brownian mesh increment verify the increment premise of the generic chain equality. Therefore the deterministic matching-noise Euler approximation is exactly the already defined Brownian-normalized \(L^2\) Euler chain.



The normalized innovation satisfies \(\sqrt{2h}\,\xi_k=N_{kh}(h)\). Combining this endpoint identity with the grid-increment formula verifies the matching-increments condition of the chain comparison. Hence the Brownian additive Euler approximation is exactly the same square-integrable chain whose law has been identified as Gaussian Euler. Its convergence and its distributional interpretation now concern one common sequence.

\subsection{\texorpdfstring{The endpoint innovation has the Gaussian vector law}{The endpoint innovation has the Gaussian vector law}}


Each normalized coordinate is standard Gaussian. The supplied joint scalar-coordinate independence lets the standard Gaussian joint law of the independent scalar coordinates and independence after grouping scalar innovations into vectors assemble those coordinates into iid standard multivariate Gaussian vectors. These are the exact law hypotheses needed by the Markov-recursion theorem.



A Brownian increment over length \(h\) has variance \(h\), so division by \(\sqrt h\) gives a standard scalar Gaussian. Independence across Brownian coordinates and disjoint grid intervals gives joint independence over coordinate and step. Grouping each finite coordinate block therefore yields independent standard vector Gaussians. The statement verifies the complete input law required by the canonical Markov recursion.

\subsection{\texorpdfstring{Passing the invariant Gibbs marginal to the continuous limit}{Passing the invariant Gibbs marginal to the continuous limit}}


Fix \(T>0\) and choose meshes \(h_n=T/N_n\to0\) small enough for the nonexpansive MALA coupling. Matching-noise Euler approximations converge in \(L^2\) to \(Y_T\). Their laws are exactly \(P_{h_n}^{N_n}\pi_U\), by the proved iid Brownian innovation laws and initialization independence.



Couple each Euler law to \(N_n\) MALA steps started from Gibbs. Exact MALA invariance makes the second marginal \(\pi_U\). The rejection bound gives expected coupling distance at most \(N_nC h_n\sqrt{h_n}=TC\sqrt{h_n}\to0\); integrability of this distance is proved from the finite extended-cost estimate.



Apply identification of limiting laws by characteristic tests and vanishing coupling cost: characteristic functions of the Euler endpoints converge to those of \(Y_T\) by \(L^2\) convergence, and to those of \(\pi_U\) by the vanishing coupling cost. Equality of characteristic functions proves \(P_{Y_T}=\pi_U\). The zero-horizon case follows directly from the integral equation.



For fixed \(T\), choose mesh lengths tending to zero with an integral number of steps reaching \(T\). The Euler endpoints converge in \(L^2\) to the integral solution. Their coupling to invariant Metropolis-adjusted endpoints has mean distance at most \(TC\sqrt h\), tending to zero. Since the invariant endpoint law is exactly Gibbs at every mesh, characteristic-function continuity identifies the continuous endpoint law with Gibbs. Apply this argument at every fixed \(T\).

\section{\texorpdfstring{Identify actual Gaussian IPLA with the shared ULA construction}{Identify actual Gaussian IPLA with the shared ULA construction}}
\label{proof-module-IPLA-ULAAdapter-lean}


This identification makes the shared Gaussian Langevin theorem applicable to the actual particle recursion. The proof checks the complete innovation distribution and independence requirements, fixes the initial marginal on the auxiliary product construction, and then transfers the resulting output law to an arbitrary supplied realization.




This adapter translates the particle coordinates into shared ULA probability interfaces. It proves the needed Gaussian and law identities, while all Brownian construction and stationarity arguments remain in their shared modules.

\subsection{\texorpdfstring{Brownian increments satisfy the complete algorithm noise specification}{Brownian increments satisfy the complete algorithm noise specification}}




For \(h=\gamma/N>0\), the normalized increment \([B_a((k+1)h)-B_a(kh)]/\sqrt h\) has variance one. Disjoint Brownian time intervals give independence over \(k\), and independent scalar Brownian paths give independence over \(a\). The product construction separates the entire Brownian source from initialization. These three facts supply fresh standard state-space Gaussian innovations at every step.

\subsection{\texorpdfstring{Identify all realizations with the same Gaussian chain law}{Identify all realizations with the same Gaussian chain law}}


The Brownian product's initial random variable has the prescribed law \(P\circ Z_0^{-1}\), and its innovations satisfy the Gaussian structure just proved. The common initial and innovation laws therefore identify its IPLA law with any supplied Gaussian realization having the same initial law. The shared product-chain law identity then gives

\[\rho_k=\mu_0K_{\gamma/N}^{\,k}.\]

Finite coupling cost and the target second moment are used here to construct the product initial state in \(L^2\). Neither stationarity nor a reference diffusion is assumed for this law identity.






The auxiliary source may depend on the chosen initial coupling, but its first initial marginal is always \(\mu_0\). Independence makes the joint law of that initial state and its Gaussian path always \(\mu_0\otimes\mathcal G\). The measurable recursion therefore has the same pushforward for every coupling and agrees with the supplied realization. This common-output property is needed before taking an infimum over initial couplings.

\subsection{\texorpdfstring{The actual chain iterates the ULA Markov transition}{The actual chain iterates the ULA Markov transition}}


The innovation hypotheses imply independent vector-valued Gaussian increments by grouping the scalar coordinate family. Applying the shared Hilbert-space law theorem to \(V_N\), step \(\gamma/N\), and smoothness \(NL\) gives

\[\rho_k=\mu_0K_{\gamma/N}^{\,k}.\]

The continuous \(L^2\) representatives are measurable; the initial state is independent of the full noise path by hypothesis. This exact transition identity transfers both convergence and finiteness statements from ULA to the actual IPLA laws.





If \(K\) denotes the Gaussian Euler kernel, the law after one step is \(\rho K\), obtained by pushing \(\rho\otimes\mathcal N(0,I_D)\) through the update. Fresh-noise independence gives this formula conditionally on the current state. Iterating gives \(\rho_k=\mu_0K^k\). Equality of these laws transfers both a numerical convergence bound and its finite-moment conclusions to the explicitly defined particle process.

\section{\texorpdfstring{Construct an isonormal Gaussian family from product coordinates}{Construct an isonormal Gaussian family from product coordinates}}
\label{proof-module-UnadjustedLangevin-GaussianHilbert-lean}


An explicit product of standard normal coordinates realizes a linear isometry \(I:H\to L^2\) with \(I(x)\sim\mathcal N(0,\|x\|^2)\) and covariance \(\mathbb E[I(x)I(y)]=\langle x,y\rangle\).



Take independent scalar standard Gaussians \((Z_j)\) and an orthonormal basis \((e_j)\) of the real Hilbert space. For \(u\), the partial sums \(\sum_{j\le n}\langle u,e_j\rangle Z_j\) are Cauchy in \(L^2\), because the squared norm of a tail equals the sum of its squared coefficients. Their limit defines an isonormal linear map \(W\). It satisfies \(\mathbb EW(u)=0\) and \(\mathbb E[W(u)W(v)]=\langle u,v\rangle\); finite linear combinations are Gaussian by the same partial-sum limit.

\subsection{\texorpdfstring{The independent Gaussian source}{The independent Gaussian source}}


Use the product probability \(Q=\bigotimes_{i\in\iota}\mathcal N(0,1)\) on real coordinate families. The construction uses the countable product measure; independence and coordinate marginals are properties of this concrete source rather than extra assumptions on an unspecified Gaussian process.



The source is a countable product of standard Gaussian probabilities. Its coordinate projections provide independent Gaussian variables from which all later finite sums and \(L^2\) limits are synthesized.

\subsection{\texorpdfstring{The law of each scalar source coordinate}{The law of each scalar source coordinate}}


The coordinate projection \(\omega\mapsto\omega_i\) pushes the product measure to its \(i\)-th factor \(\mathcal N(0,1)\). The product marginal theorem proves both almost-everywhere measurability and the pushforward identity, packaged as the stated pushforward law.



A coordinate projection of a product measure has its designated marginal. Thus each scalar source variable has mean zero and variance one before any Hilbert coefficient is introduced.

\subsection{\texorpdfstring{The square-\allowbreak{}integrable coordinate class}{The square-integrable coordinate class}}


A scalar source coordinate has second moment one, so it represents an element \(Z_i\in L^2(Q)\). The associated equivalence class identifies variables agreeing almost everywhere. Its norm and law can therefore be computed from the original measurable coordinate projection.



Each scalar coordinate has finite second moment one, so it determines an \(L^2\) class. This places all source coordinates in a common Hilbert space where series convergence can be expressed by norm limits.

\subsection{\texorpdfstring{Its representative identity}{Its representative identity}}


The selected representative of \(Z_i\) agrees with the original source coordinate \(\omega_i\) almost everywhere. This equality preserves every integrable expectation and its pushforward law. Consequently the Gaussian marginal, centering, and independence calculations apply to the same square-integrable classes used for synthesis.



The class representative agrees almost everywhere with the original coordinate projection. Consequently law and moment calculations made for the coordinate projection transfer unchanged to the class.

\subsection{\texorpdfstring{Independence of the source coordinates}{Independence of the source coordinates}}


The coordinate projections of an infinite product measure are jointly independent. Invoke that product theorem for the scalar Gaussian factors. This is independence of every finite set of coordinates, which is precisely what finite Gaussian linear combinations require.



Independence follows from the product law on cylinder events. In particular, distinct coordinate products have expectation equal to the product of their centered means, which is zero.

\subsection{\texorpdfstring{The mean of a scalar coordinate}{The mean of a scalar coordinate}}


Change variables through the coordinate law and use the centered Gaussian mean: \(\int\omega_i\,dQ=\int x\,d\mathcal N(0,1)=0\). Gaussian integrability makes this an ordinary expectation identity.



The scalar standard Gaussian is centered. Integrating a coordinate against the product probability therefore gives zero, with all other coordinates integrating to total mass one.

\subsection{\texorpdfstring{Inner products of coordinate classes}{Inner products of coordinate classes}}


For \(i=j\), the \(L^2\) inner product is \(\mathbb EZ_i^2=1\). For \(i\ne j\), independence factors \(\mathbb EZ_iZ_j=(\mathbb EZ_i)(\mathbb EZ_j)=0\). Almost-everywhere representative identities move between \(Z_i\) and the product coordinates. Thus \(\langle Z_i,Z_j\rangle=\mathbf1_{i=j}\).



Equal indices give the scalar second moment one. Distinct indices give a product of two zero means by independence, yielding the Kronecker-delta inner product.

\subsection{\texorpdfstring{Orthonormality of the Gaussian coordinates}{Orthonormality of the Gaussian coordinates}}


The preceding inner-product formula is the orthonormality criterion for a family in a real Hilbert space. The identity converts the diagonal value one and off-diagonal value zero into orthonormality of the scalar coordinate classes, the input for the orthogonal-series extension.



The coordinate classes have unit norm and pairwise zero inner products. These are exactly the orthonormality conditions that turn coefficient-square summability into convergence of the synthesis series.

\subsection{\texorpdfstring{Normalization of the source probability}{Normalization of the source probability}}


Each scalar Gaussian factor is a probability measure. The product-probability hypothesis therefore gives \(Q(\Omega)=1\). This normalization permits constant functions and integrable functions to be treated in the usual probability-space \(L^2\) theory below.



A countable product of probability measures has total mass one. This permits the expectation and \(L^2\) inequalities used throughout the synthesis construction without additional finite-mass factors.

\subsection{\texorpdfstring{Synthesis from square-\allowbreak{}summable coefficients}{Synthesis from square-summable coefficients}}


Extend the orthonormal family \(Z_i\) to a linear isometry \(S:\ell^2(\iota)\to L^2(Q)\). The orthogonal-family extension theorem constructs the limit of finite sums and preserves squared norm, \(\|S(a)\|_2^2=\sum_i a_i^2\). It constructs an \(L^2\) limit rather than assuming pointwise convergence of a random series.



The coefficient sequence of a Hilbert vector is square summable by Parseval. Orthogonality of the Gaussian coordinates then makes the corresponding random series Cauchy in \(L^2\), defining its synthesis limit.

\subsection{\texorpdfstring{Convergence of the synthesis series}{Convergence of the synthesis series}}


The synthesis series \(\sum_i a_iZ_i\) converges in \(L^2\). Its convergence may be expressed by finite-subset partial sums: the squared norm of a tail is the sum of its squared coefficients, which tends to zero by square summability. No enumeration of the index set is required for this orthogonal-series limit.



The squared norm of a partial-sum difference is the sum of squared tail coefficients. Those tails tend to zero, proving that the prescribed series converges to the synthesized class.

\subsection{\texorpdfstring{The isonormal linear map}{The isonormal linear map}}


For a Hilbert basis \(b\), compose its coordinate isometry \(\mathcal C_b:H\to\ell^2(\iota)\) with gaussian synthesis from hilbert coefficients. The resulting linear isometry \(I_b\) sends \(x\) to the Gaussian series with coefficients \(\langle b_i,x\rangle\), and preserves all inner products.



The isonormal map packages synthesis as a linear map on the original Hilbert space. Its norm preservation is the same coefficient-square identity used to prove convergence.

\subsection{\texorpdfstring{The series representation of the isonormal map}{The series representation of the isonormal map}}


Apply the norm convergence of the Gaussian coordinate series to the coefficient family \(\mathcal C_b(x)\). Unfolding the composition gives \(I_b(x)=\sum_i\langle b_i,x\rangle Z_i\) in \(L^2\). This is the exact series represented by the isonormal linear isometry.



The series formula identifies the abstract linear map with the concrete Gaussian-coordinate expansion. This permits finite-sum moment calculations to pass to arbitrary Hilbert inputs by continuity.

\subsection{\texorpdfstring{Expectation as an inner product with one}{Expectation as an inner product with one}}


On a probability space, the constant one is in \(L^2\). The definition of the real \(L^2\) inner product gives \(\mathbb EX=\langle1,X\rangle\), after replacing the selected constant representative by one almost everywhere. This exhibits expectation as a continuous functional.



On a probability space the constant-one class has norm one. Pairing a scalar class with it is integration, so expectation is a continuous linear functional on \(L^2\).

\subsection{\texorpdfstring{Centering of each coordinate class}{Centering of each coordinate class}}


Use \(Z_i=\omega_i\) almost everywhere to replace the \(L^2\) representative inside the integral. The already established centered coordinate expectation then gives \(\mathbb EZ_i=0\).



The representative identity transfers the source coordinate's zero mean to its \(L^2\) class. This verifies centering in the same Hilbert-space language used for the infinite synthesis sum.

\subsection{\texorpdfstring{Centering of the synthesized series}{Centering of the synthesized series}}


Each finite partial sum \(\sum_{i\in F}a_iZ_i\) has mean zero by linearity and the coordinate means. Pair the \(L^2\)-convergent series with the constant one; continuity of the inner product passes the zero expectations to the limit, giving \(\mathbb ES(a)=0\).



Every finite synthesis sum has mean zero. Continuity of the expectation functional carries that equality to the \(L^2\) limit, proving the synthesized variable is centered.

\subsection{\texorpdfstring{Centering of every isonormal value}{Centering of every isonormal value}}


The isonormal variable is gaussian synthesis from hilbert coefficients applied to the Hilbert-basis coefficients of \(x\). Specialize the synthesis mean theorem to these coefficients to get \(\mathbb EI_b(x)=0\), without redoing the infinite-series limit.



The isonormal value is the synthesized class, so its zero mean follows from the same limit calculation. Linearity ensures that this centering holds for every index vector.

\subsection{\texorpdfstring{Variance from the Hilbert norm}{Variance from the Hilbert norm}}


For an \(L^2\) random variable with mean zero, the variance formula reduces to \(\operatorname{Var}(X)=\int X^2dQ\). The \(L^2\) norm identity identifies this with \(\|X\|_2^2\). The explicit mean premise is used to remove the centering term.



For a centered scalar class, variance equals its squared \(L^2\) norm. Norm preservation of synthesis therefore identifies the variance with the squared norm of the Hilbert index vector.

\subsection{\texorpdfstring{The representative of a finite partial sum}{The representative of a finite partial sum}}


Induct over the finite set. The zero sum has zero representative almost everywhere; addition and scalar multiplication of \(L^2\) classes agree almost everywhere with the corresponding pointwise operations. Intersect the finitely many full-measure sets to obtain \(\sum_{i\in F}a_iZ_i=\sum_{i\in F}a_i\omega_i\) almost everywhere.



The class of a finite linear combination agrees almost everywhere with the corresponding pointwise coordinate sum. Thus elementary finite-dimensional Gaussian formulas can be applied to that sum.

\subsection{\texorpdfstring{Gaussianity of finite partial sums}{Gaussianity of finite partial sums}}


Independent scalar standard Gaussians remain Gaussian after a finite linear combination. Apply that finite-sum Gaussian theorem to \(a_i\omega_i\), and transfer the law across the representative equality for the finite Gaussian sum. This proves Gaussianity of the actual \(L^2\) partial sum, not merely of a formally similar expression.



A finite linear combination of independent Gaussian coordinates is Gaussian. Its variance is the sum of squared coefficients because all cross covariances vanish.

\subsection{\texorpdfstring{The mean of each finite partial sum}{The mean of each finite partial sum}}


The integral of a finite \(L^2\) sum is the sum of its integrals, because every summand is integrable. Each coordinate mean is zero, hence \(\mathbb E\sum_{i\in F}a_iZ_i=0\). This supplies the centering parameter needed to identify its Gaussian law.



Each finite coefficient sum remains centered, since expectation is linear and every scalar coordinate is centered. This identifies its mean parameter before its exact scalar law is recorded.

\subsection{\texorpdfstring{The exact Gaussian law of a partial sum}{The exact Gaussian law of a partial sum}}


A scalar Gaussian law is determined by its mean and variance. The finite sum is Gaussian and centered; its variance is its squared \(L^2\) norm. Therefore its law is \(\mathcal N(0,\|\sum_{i\in F}a_iZ_i\|_2^2)\), with the variance converted to a nonnegative real.



The finite sum the stated pushforward law \(\mathcal N(0,\sum_{j\le n}\langle u,e_j\rangle^2)\). This records the full law, whose variance converges to \(\|u\|^2\).

\subsection{\texorpdfstring{The Gaussian law of the synthesis limit}{The Gaussian law of the synthesis limit}}


Finite partial sums converge in \(L^2\) to \(S(a)\), and each has the centered Gaussian law with its own squared norm as variance. Apply preservation of a Gaussian law under square-mean convergence; the isometry then rewrites the limiting variance as \(\|a\|_{\ell^2}^2\). This proves Gaussianity of the completed infinite series.



The finite sums converge in \(L^2\), and their variances converge by Parseval. Characteristic-function continuity therefore gives Gaussian law \(\mathcal N(0,\|u\|^2)\) for the synthesis limit.

\subsection{\texorpdfstring{The Gaussian law of an isonormal value}{The Gaussian law of an isonormal value}}


Apply the synthesis law theorem to \(\mathcal C_b(x)\). The Hilbert-basis coordinate map preserves norm, so the variance becomes \(\|x\|^2\). Hence \(I_b(x)\sim\mathcal N(0,\|x\|^2)\), including the degenerate law for \(x=0\).



Replacing synthesis by the equal isonormal map yields the same centered Gaussian law. This identifies each individual isonormal evaluation without selecting a particular expansion after the limit.

\subsection{\texorpdfstring{Joint Gaussianity of the isonormal family}{Joint Gaussianity of the isonormal family}}


Each isonormal evaluation has the exact law \(\mathcal N(0,\|x\|^2)\). By linearity, every finite linear combination of these evaluations is another evaluation, at the corresponding linear combination of indices. It is therefore Gaussian, which establishes joint Gaussianity of the finite-dimensional family.



Any finite linear combination of isonormal evaluations is one isonormal evaluation at the corresponding linear combination of indices. Its scalar Gaussian law verifies joint Gaussianity of the family.

\subsection{\texorpdfstring{The isonormal covariance identity}{The isonormal covariance identity}}


Both variables are centered, so their covariance is \(\mathbb E[I_b(x)I_b(y)]\), the \(L^2\) inner product. The linear isometry preserves that inner product, giving \(\operatorname{Cov}(I_b(x),I_b(y))=\langle x,y\rangle_H\). The moment bound verifies integrability through their \(L^2\) membership.



Polarization of the norm-preserving map gives \(\mathbb E[W(u)W(v)]=\langle u,v\rangle\). The means vanish, so this inner product is also the covariance.

\section{\texorpdfstring{Scalar Gaussian linear maps give jointly Gaussian processes}{Scalar Gaussian linear maps give jointly Gaussian processes}}
\label{proof-module-UnadjustedLangevin-GaussianLinearProcess-lean}


If a linear map \(T:H\to L^2(P;\mathbb R)\) sends every input to a Gaussian variable, all finite collections of its evaluations are jointly Gaussian.



A linear map from an index-vector space into \(L^2\) produces a jointly Gaussian family once every one-dimensional linear combination has a Gaussian law. For finitely many indices and coefficients, linearity rewrites their random linear combination as the image of a single index vector. Applying the scalar Gaussian property to that vector verifies the finite-dimensional Gaussian criterion. No independent-coordinate conclusion follows unless the corresponding covariances also vanish.

\subsection{\texorpdfstring{Finite-\allowbreak{}dimensional Gaussianity of a linear family}{Finite-dimensional Gaussianity of a linear family}}


Use the finite-dimensional Gaussian criterion: every scalar linear functional of \((T(f_i))_i\) must be Gaussian. A finite linear combination equals \(T(\sum_i a_if_i)\) as an \(L^2\) class by linearity, hence agrees almost everywhere as a random variable. The assumed Gaussianity of every \(T(x)\) proves the criterion.



For coefficients \(a_i\), linearity gives \(\sum_i a_iT(v_i)=T(\sum_i a_iv_i)\). The scalar Gaussian hypothesis applied to that one index verifies the Gaussian law of every finite linear combination.

\subsection{\texorpdfstring{A Gaussian process from a Gaussian linear map}{A Gaussian process from a Gaussian linear map}}


For each finite set of times, apply the preceding joint-Gaussian theorem to the corresponding finite family of vectors \(f(t)\). This verifies the finite-dimensional-law definition of joint Gaussianity of every finite collection. The index set may be uncountable; this statement alone asserts no sample-path continuity.



The finite-dimensional criterion just established applies to every finite choice of the process value at time \(t\)imes. Thus the family is a Gaussian process, with its joint distributions determined by its mean and covariance data.

\section{\texorpdfstring{An explicit Gaussian construction under pre-\allowbreak{}Brownian coordinates}{An explicit Gaussian construction under pre-Brownian coordinates}}
\label{proof-module-UnadjustedLangevin-CanonicalBrownian-lean}


The isonormal family on the time Hilbert space constructs a Gaussian process with covariance \(\min(s,t)\). Independent product copies supply the Brownian coordinates required by the canonical ULA theorem. The construction proves pre-Brownian laws; continuous sample paths are not asserted or needed, because continuity is established in \(L^2\).



In the time Hilbert space \(L^2([0,\infty))\), use the indicator \(1_{[0,t]}\) to index the isonormal map. Its inner product with \(1_{[0,s]}\) is the length of their overlap, \(\min(s,t)\). The resulting centered Gaussian family therefore has Brownian covariance. Independent copies give the spatial coordinates; zero covariance of disjoint time increments in a Gaussian family gives their independence. The construction supplies the marginal and increment laws needed by the later \(L^2\) reference argument.

\subsection{\texorpdfstring{The time-\allowbreak{}function Hilbert space}{The time-function Hilbert space}}


Take the deterministic Hilbert space \(L^2(\mathbb R,dt)\). Its vectors will encode time intervals before an isonormal map turns their inner products into covariances. It is distinct from the random-variable \(L^2\) space on the Gaussian source.



The time-index Hilbert space uses Lebesgue square integration on the nonnegative half-line. Indicators of finite intervals belong to it because their squared integral is the interval length.

\subsection{\texorpdfstring{An interval indicator as a time index}{An interval indicator as a time index}}


Let \(e_t=\mathbf1_{(0,t]}\) as a class in \(L^2(dt)\), for \(t\ge0\). Construct it with the square-integrable interval indicator and the finite Lebesgue measure of \((0,t]\). Endpoints do not affect this Lebesgue \(L^2\) class, but this half-open convention makes interval intersection identities direct.



The indicator is one before time \(t\) and zero afterward, up to irrelevant endpoint conventions. Its squared Hilbert norm is \(t\), so the associated isonormal variable has variance \(t\).

\subsection{\texorpdfstring{The overlap inner product of two indicators}{The overlap inner product of two indicators}}


The product of the two indicators is the indicator of \([0,\min(s,t)]\) almost everywhere. Integrating it gives \(\langle e_s,e_t\rangle=\min(s,t)\). The proof evaluates the Lebesgue measure of that finite interval, producing the Brownian covariance before introducing randomness.



The product of the two indicators is the indicator of their overlap. Its integral is \(\min(s,t)\), yielding the desired covariance kernel exactly.

\subsection{\texorpdfstring{A countable basis index}{A countable basis index}}


Choose an orthonormal basis of the time-function Hilbert space and use its index set to label the independent scalar Gaussian coordinates. Orthogonal-series synthesis depends only on square-summable coefficients, so its definition does not require fixing an enumeration.



A countable basis index permits synthesis from a countable Gaussian product source. The particular enumeration affects the construction's description but not its covariance or resulting Gaussian laws.

\subsection{\texorpdfstring{An orthonormal basis for time functions}{An orthonormal basis for time functions}}


Select a Hilbert basis for the hilbert space of time functions by Hilbert-space basis existence. This basis supplies the coordinate isometry for the explicit isonormal Gaussian family; all subsequent covariance identities follow from isometry and do not depend on which basis was selected.



An orthonormal basis expands each time indicator with square-summable coefficients. Parseval identifies their squared sum with the interval length, ensuring the Brownian synthesis has the correct variance.

\subsection{\texorpdfstring{The scalar Gaussian sample space}{The scalar Gaussian sample space}}


A sample is a real family indexed by the chosen time-Hilbert basis. This concrete product space stores the independent standard Gaussian coefficients used to synthesize all Brownian times.



The sample space is the Gaussian coordinate source underlying the isonormal map. Time is an index of synthesized variables on this fixed space, rather than an additional random coordinate.

\subsection{\texorpdfstring{The scalar source probability}{The scalar source probability}}


Take the product of standard scalar Gaussian measures indexed by the chosen time basis. This probability has independent coordinate projections, providing the source variables used to synthesize every time value on one common sample space.



The scalar source probability has mass one. All synthesized time values are therefore random variables under one probability law, permitting their finite-dimensional covariances to be compared.

\subsection{\texorpdfstring{Normalization of the scalar probability}{Normalization of the scalar probability}}


The scalar source is the Gaussian coordinate product probability. Reuse normalization of the independent Gaussian product probability to provide normalization for the chosen basis index. This normalization enables the \(L^2\) expectation and covariance theorems for the constructed Brownian process.



Normalization of the source is inherited from its Gaussian product construction. This supplies the common probability assumption for Gaussian process and moment statements that follow.

\subsection{\texorpdfstring{A time value in the square-\allowbreak{}integrable space}{A time value in the square-integrable space}}


Apply the isonormal isometry to \(e_t\), obtaining \(B_t=I(e_t)\) as an \(L^2\) random variable. The construction immediately gives a centered Gaussian with variance \(\|e_t\|^2=t\); joint Gaussianity follows from linearity.



Apply the isonormal map to the interval indicator to obtain a scalar \(L^2\) time value. Its centered Gaussian law and variance follow from the indicator norm calculation.

\subsection{\texorpdfstring{The scalar Gaussian process}{The scalar Gaussian process}}


At each time select the strongly measurable representative of a brownian value as a square-integrable variable. This defines an actual function \(B_t(\omega)\) for use in the stated pushforward law and finite-dimensional distributions. Equalities of \(L^2\) classes still translate to almost-everywhere equalities of these representatives.



Choose the prescribed representative of each \(L^2\) time value to define the scalar process. Fixed-time laws are unaffected by the representative choice on null sets.

\subsection{\texorpdfstring{Fixed-\allowbreak{}time measurability of the process}{Fixed-time measurability of the process}}


The representative selected for an \(L^2\) class is measurable in this setting. Apply the representative measurability theorem to the square-integrable value at time \(t\), obtaining measurability of the actual time-\(t\) random variable.



At each fixed time the selected representative is measurable. This is the measurability needed for its marginal pushforward law and for assembling finite collections of time values.

\subsection{\texorpdfstring{The mean of the process}{The mean of the process}}


The isonormal mean theorem gives \(\mathbb EI(e_t)=0\). Since the process value at time \(t\) is precisely the selected representative of this class, this is the expectation of the displayed random function \(B_t\).



The isonormal map is centered at every index vector. Applying it to a time indicator therefore gives zero mean for each process value.

\subsection{\texorpdfstring{The Brownian covariance formula}{The Brownian covariance formula}}


The isonormal covariance formula gives \(\operatorname{Cov}(I(e_s),I(e_t))=\langle e_s,e_t\rangle\). Substitute the overlap formula \(\langle1_{(0,s]},1_{(0,t]}\rangle=\min(s,t)\) to obtain \(\min(s,t)\). This calculation establishes the Brownian covariance for the concrete source and representatives.



The isonormal covariance equals the Hilbert inner product of the two indicators. Their overlap calculation gives \(\mathbb E[W_sW_t]=\min(s,t)\).

\subsection{\texorpdfstring{Joint Gaussianity of the time family}{Joint Gaussianity of the time family}}


The isonormal isometry is linear and every output is scalar Gaussian. Apply Gaussianity of linear combinations of the isonormal map to the family \(t\mapsto e_t\). Every finite collection of the selected random variables is therefore jointly Gaussian.



Finite linear combinations of time values are isonormal evaluations at finite linear combinations of indicators. Their Gaussian laws establish the process's joint Gaussian property.

\subsection{\texorpdfstring{The Brownian finite-\allowbreak{}dimensional and increment properties}{The Brownian finite-dimensional and increment properties}}


A centered jointly Gaussian family with covariance \(\min(s,t)\) has Brownian finite-dimensional distributions. Differences over disjoint intervals have zero covariance, so their joint Gaussianity makes them independent. These are the increment and distribution properties required by the square-mean construction; sample-path continuity is not asserted here.



Centered Gaussianity and covariance \(\min(s,t)\) determine the Brownian finite-dimensional laws. Orthogonality of disjoint indicator increments gives independence of disjoint time increments.

\subsection{\texorpdfstring{The product sample space for spatial coordinates}{The product sample space for spatial coordinates}}


For a coordinate index type \(I\), take one scalar Brownian sample for each \(i\in I\). This is the product sample type \(\Omega^I\), allowing independent entire coordinate processes rather than only independent values at one time.



The vector source is the finite product of scalar process sample spaces. Each spatial coordinate reads one separate factor, making the scalar paths independent by construction.

\subsection{\texorpdfstring{The vector source probability}{The vector source probability}}


Equip the vector sample space with the infinite product of the scalar Brownian source probabilities. The coordinate processes inherit independence from this actual product construction.



The finite product probability gives mass one and retains the scalar source law under each projection. It is the common probability space for the Euclidean Brownian vector.

\subsection{\texorpdfstring{A scalar spatial coordinate process}{A scalar spatial coordinate process}}


Define the \(i\)-th coordinate process by \(B_i(t,\omega)=B(t,\omega_i)\). It first projects to the \(i\)-th scalar sample, then evaluates the already constructed scalar process at \(t\).



A coordinate process is the scalar Brownian process composed with one product projection. This explicitly identifies which source factor supplies each spatial coordinate.

\subsection{\texorpdfstring{Brownian properties of each spatial coordinate}{Brownian properties of each spatial coordinate}}


The product coordinate projection has the scalar Brownian source as its law. Pull the scalar pre-Brownian process back along this projection; its finite-dimensional distributions, mean, and covariance are preserved. Thus each vector coordinate has the established scalar Brownian laws.



The coordinate projection preserves the scalar probability law. Hence all centered Gaussian increment and covariance properties transfer to the corresponding spatial coordinate.

\subsection{\texorpdfstring{Independence of the coordinate paths}{Independence of the coordinate paths}}


The product coordinate projections are independent as sample-valued maps. Compose each projection with the map sending a scalar sample to its entire time-indexed process. Independence is preserved under these coordinatewise functions, giving independence of the whole coordinate paths.



Independence is asserted for whole coordinate paths as observations of separate product factors. It therefore entails independence of every finite collection of observations drawn from distinct coordinates.

\subsection{\texorpdfstring{Independence after restriction to mesh times}{Independence after restriction to mesh times}}


Restrict each independent coordinate path to the deterministic grid \(nh\). This is a coordinatewise function of the path, so independence is preserved. The result supplies exactly the countable grid-path hypothesis used in the innovation theorems.



Restricting each path to mesh times is a measurable observation of that path. Independence is preserved by this restriction, so the grid paths retain spatial independence.

\subsection{\texorpdfstring{Normalization of the vector probability}{Normalization of the vector probability}}


Each scalar Brownian sample source has mass one. The infinite product over vector coordinates therefore also has mass one. This normalization concerns the actual vector sample measure, not an assumed normalization of an abstract process law.



Normalization of the vector product source supplies the probability setting for the later Euclidean Gaussian and Brownian moment calculations. No additional weighting of the coordinate paths is introduced.

\section{\texorpdfstring{Stability and uniqueness for additive-\allowbreak{}noise solutions}{Stability and uniqueness for additive-noise solutions}}
\label{proof-module-UnadjustedLangevin-LangevinUniqueness-lean}


Two solutions driven by the same continuous additive noise satisfy \(\|Y_t-Z_t\|\le e^{Lt}\|x-y\|\) on the prescribed interval. Equality of initial values gives uniqueness. Only global Lipschitzness is used here; the sharper contractive estimate uses strong convexity elsewhere.



For two solutions driven by the same additive input, that input disappears when their equations are subtracted. The remaining difference is bounded by its initial separation plus the time integral of \(L\) times its norm. Grönwall's inequality then yields the stability factor \(e^{Lt}\). With equal initial values the bound is zero. In a space of random variables, this proves equality of the corresponding \(L^2\) classes at each fixed time.

\subsection{\texorpdfstring{Stability for a common additive input}{Stability for a common additive input}}


Subtract the two integral equations so that their common additive input cancels. The integral norm inequality and Lipschitz drift give \(u(t)\le\|x-y\|+L\int_0^t u(s)ds\), where \(u(t)=\|Y(t)-Z(t)\|\) is continuous. Grönwall's inequality yields \(u(t)\le\|x-y\|e^{Lt}\). All time integrals are justified by continuity on the compact horizon.



After cancellation of the common input, the norm difference is bounded by initial separation plus \(L\) times its time integral. The exponential factor comes from this scalar Grönwall comparison, without differentiating the input.

\subsection{\texorpdfstring{Uniqueness with equal initialization}{Uniqueness with equal initialization}}


Set the two initial states equal in the stability estimate. The right side becomes \(0e^{Lt}=0\), forcing \(\|Y(t)-Z(t)\|=0\) for each \(t\in[0,h]\). Norm-zero equivalence gives equality of the vectors, packaged as equality at every point on the interval.



Equal initialization sets the multiplicative initial distance to zero. Nonnegativity of the norm then forces equality of the solutions at every time on the interval, giving uniqueness among the stated continuous integral solutions.

\subsection{\texorpdfstring{Stability of square-\allowbreak{}integrable solutions}{Stability of square-integrable solutions}}


Specialize Banach-space stability to the actual lifted gradient on \(L^2(P;K)\). Its Lipschitz theorem supplies the required constant, so \(\|Y(t)-Z(t)\|_2\le e^{Lt}\|Y(0)-Z(0)\|_2\). This compares solutions with a common \(L^2\)-noise input and possibly different initial random variables.



All operations now take place in \(L^2\), so the resulting distance is a root mean square difference. A different coupling of the initial random vectors changes that initial norm but not the stability constant.

\subsection{\texorpdfstring{Uniqueness of square-\allowbreak{}integrable solutions}{Uniqueness of square-integrable solutions}}


Apply the same-input uniqueness theorem in \(L^2(P;K)\). Equality here is equality of \(L^2\) classes for each time, hence almost-everywhere equality of representatives at each fixed time. The conclusion does not by itself select one exceptional null set valid simultaneously at every real time.



Equality of \(L^2\) classes at a fixed time means equality of representatives almost everywhere at that time. A simultaneous pathwise assertion would need an additional common-null-set argument and is not supplied by this fixed-time conclusion alone.

\section{\texorpdfstring{Identify chosen reference segments with one global solution}{Identify chosen reference segments with one global solution}}
\label{proof-module-UnadjustedLangevin-BrownianSegmentsGlobal-lean}


The segment construction is compatible with the unique global integral solution. This lets stationarity of a global Brownian Langevin solution prove stationarity of every chosen reference segment.



A global solution restricted to \([a,a+h]\), and reparametrized by \(t\mapsto a+t\), solves the same equation as the independently restarted segment. Subtracting the global equation at \(a\) cancels its earlier drift and noise history, leaving initial state \(Y_a\) and shifted input \(B_{a+t}-B_a\). Same-input uniqueness identifies the two solutions. Induction over knots then identifies every recursively constructed segment with the corresponding portion of the global trajectory.

\subsection{\texorpdfstring{The increment of the Langevin Brownian input}{The increment of the Langevin Brownian input}}


Expand the two increments from zero and subtract them. Their common \([B_0]\) term cancels, leaving \(\sqrt2([B_{a+t}]-[B_a])=N_a(t)\). This is an algebraic cocycle identity in \(L^2\), valid for the real-time extension as defined.



Both sides equal the same linear combination of three Brownian values. Expanding \(N_0(a+t)-N_0(a)\) cancels the two copies of \(B_0\), leaving \(\sqrt2(B_{a+t}-B_a)\). This cocycle property is the noise counterpart of additivity of time integrals and lets a global equation be restarted at any nonnegative absolute time.

\subsection{\texorpdfstring{Shifting the Langevin integral equation}{Shifting the Langevin integral equation}}


Subtract the global integral equations at \(a+t\) and \(a\). Integral additivity gives the drift integral over \([a,a+t]\), and time translation turns it into \(\int_0^t\nabla U(Y_{a+s})ds\). The noise cocycle gives \(N_a(t)\), yielding the local segment equation with initial state \(Y_a\).



Subtracting the equations removes the original initial state. Additivity gives \(\int_0^{a+t}-\int_0^a=\int_a^{a+t}\), and the substitution \(s=a+r\) puts the latter integral on \([0,t]\). Continuity of the composed score ensures these time integrals exist. The cocycle then produces exactly the local input used in the selected segment construction.

\subsection{\texorpdfstring{A restarted segment equals the global solution}{A restarted segment equals the global solution}}


The shifted global path and the selected local segment are continuous solutions of the same \(L^2\) integral equation with the same initial value and noise input. Apply \(L^2\) uniqueness to identify them on \([0,h]\). The horizon premise ensures the shifted interval lies inside the global solution interval.



Let \(D_t\) be the difference between the selected local segment and the shifted global path. Their identical inputs cancel, giving \(\|D_t\|_2\le L\int_0^t\|D_s\|_2ds\) with \(D_0=0\). Grönwall forces \(D_t=0\). Thus arbitrary choices made in selecting local solutions cannot change the path on the common interval.

\subsection{\texorpdfstring{The knots lie on the global solution}{The knots lie on the global solution}}


Induct on the knot index. At zero, the global equation and zero-starting noise give the prescribed initial state. If the current knot equals \(Y_{nh}\), segment/\allowbreak{}global uniqueness makes its endpoint equal \(Y_{(n+1)h}\), establishing the next knot identity.



The induction fixes the relationship between recursive endpoints and absolute time. If \(K_n=Y_{nh}\), local uniqueness identifies the whole next segment with \(t\mapsto Y_{nh+t}\). Its endpoint is therefore \(Y_{nh+h}=Y_{(n+1)h}\). Starting from the original initial equality proves the knot formula for every index whose time lies inside the chosen horizon.

\subsection{\texorpdfstring{Every reference segment lies on the global solution}{Every reference segment lies on the global solution}}


The knot/\allowbreak{}global theorem supplies the matching initial value for the \(n\)-th segment. Apply uniqueness of restarted and globally restricted integral solutions to conclude \(Y^{(n)}_t=Y_{nh+t}\) on \([0,h]\). This is a deterministic identification of the constructions, with no stationarity premise.



Once the initial knot has been identified, the segment equation and shifted global equation have the same data. Their equality therefore holds at every local time, including the endpoints. In particular any fixed-time law theorem for the global the process value at time \(t\)ransfers to every segment time \(nh+t\). Stationarity can thus be established once for the global solution instead of reproved separately for each recursively selected segment.

\section{\texorpdfstring{Dalalyan’s theorem for the actual Gaussian ULA transition}{Dalalyan’s theorem for the actual Gaussian ULA transition}}
\label{proof-module-UnadjustedLangevin-CanonicalULA-lean}


Let \(U:\mathbb R^d\to\mathbb R\) be differentiable and \(m\)-strongly convex with \(m>0\), with \(L\)-Lipschitz gradient. The canonical Brownian source derives all stochastic reference obligations and proves the Wasserstein bound for the iterated Gaussian transition \(P_h\), with finite initial \(W_2\).



The canonical output law is obtained by applying the independent standard-Gaussian transition repeatedly to the initial measure. The constructed product-Brownian realization has exactly these laws, because its normalized mesh increments are independent standard Gaussian vectors. Its reference estimates are uniform in the initial coupling. Taking the infimum consequently replaces the chosen initial \(L^2\) cost by \(W_2(\mu_0,\pi_U)\), while retaining the explicit step-dependent bias floor.

\subsection{\texorpdfstring{Identification of the canonical output law}{Identification of the canonical output law}}


The Brownian source provides jointly independent standard Gaussian innovations; the product extension makes initialization independent of the entire sequence. The realization/\allowbreak{}Markov-law theorem identifies its output with \(P_h^k\mu\). The product output-law identity removes the chosen finite coupling from this law equality.



The Brownian-generated innovations have the full iid standard vector Gaussian law. Their product construction makes the initial coordinate independent of the complete path. Thus the measurable recursion the stated pushforward law \(P_h^k\mu\). The product output definition removes the auxiliary reference coordinate, showing this identity is common to every finite initial coupling and depends only on the prescribed initial marginal.

\subsection{\texorpdfstring{A product reference from independent Brownian coordinates}{A product reference from independent Brownian coordinates}}


On each finite-coupling extension, pull back the independent Brownian coordinates and use the product independence theorems for the reference initialization. Construct a global \(L^2\) integral solution over the required finite horizon.



The proved MALA/\allowbreak{}Euler stationarity theorem gives Gibbs law for that global solution at every needed absolute time. Segment/\allowbreak{}global uniqueness identifies each recursively constructed reference segment with that solution. Thus the segment-stationarity premise of the stationary reference constructed on the independent initial-coupling and Brownian product space is discharged, filling every reference field from the stated source and potential assumptions.



For each finite coupling, solve the actual global Langevin integral equation from the target-distributed reference coordinate with independent Brownian forcing. The invariant-approximation argument proves its Gibbs law at every fixed time. Uniqueness then identifies the recursively chosen local segments with this global solution on their intervals. Hence all segment target-law premises are proved for the exact references used in the coupling family.

\subsection{\texorpdfstring{The canonical Euler innovations}{The canonical Euler innovations}}


Specialize the source Euler noise to the explicitly constructed independent coordinate Brownian source. This abbreviation fixes one concrete construction of the standard Gaussian innovations for the unconditional law theorem.



Fix the explicitly constructed independent Brownian coordinate source once and normalize its grid increments. For any \(h>0\), the resulting coordinates are standard Gaussian and independent over step and coordinate. This provides a concrete innovation sequence for the canonical algorithm rather than an existential Gaussian family chosen differently for each input law or coupling.

\subsection{\texorpdfstring{The canonical product reference}{The canonical product reference}}


The canonical Brownian construction supplies the coordinate Brownian laws and independence of their grid paths for every positive step. Apply the product reference with Brownian noise independent of the initial pair. No stochastic reference structure or stationary-law assertion remains an external premise.



The canonical source supplies its Brownian finite-dimensional laws and full-path coordinate independence. Product extension supplies independence from the initial coupling. The stationary integral solution and segment uniqueness then provide all reference target laws. Thus the analytic reference family is obtained entirely from the potential and initial transport assumptions, with no supplied stochastic convergence or invariant-law premise remaining.

\subsection{\texorpdfstring{The canonical output follows the Gaussian transition}{The canonical output follows the Gaussian transition}}


Finite initial transport provides a finite-cost coupling, and Gibbs has a finite second moment. Apply the general Brownian output-law identity on that extension with the canonical independent Brownian source. It proves the canonical product construction has exactly the ordinary ULA law \(P_h^k\mu\).



Finite initial transport supplies a finite-cost coupling, and target Gaussian tails supply the reference second moment needed on its extension. The canonical innovation law theorem identifies the numerical output with \(P_h^k\mu\). Since the output is common to all such coupling extensions, this identity does not retain any chosen coupling or auxiliary sample-space information.

\subsection{\texorpdfstring{Finiteness of the output transport distance}{Finiteness of the output transport distance}}


The canonical product reference supplies a finite-cost output coupling for every iteration whenever initial transport is finite. Its output law is \(P_h^k\mu\). Substituting this law proves \(W_2(P_h^k\mu,\pi_U)<\infty\), legitimizing the real Wasserstein notation used in the final estimates.



The constructed numerical and stationary reference states lie in \(L^2\) on any one finite-cost extension. Their joint law is therefore a finite-cost coupling at every iteration. The numerical marginal is \(P_h^k\mu\) and the reference marginal is \(\pi_U\). This proves output transport finiteness before passing to real-valued distances in the quantitative convergence bounds.

\subsection{\texorpdfstring{The small-\allowbreak{}step canonical Wasserstein bound}{The small-step canonical Wasserstein bound}}


Construct the canonical product reference from differentiability, \(m\)-strong convexity, and \(L\)-Lipschitz gradient. Nontriviality of the finite coordinate space proves \(m\le L\); the finite-coupling product theorem gives the small-step bound with bias \((33/20)(L/m)\sqrt{hd}\).



Finally replace the canonical product output by \(P_h^k\mu\) using the identification of the product-space chain law with the iterated Gaussian transition. The resulting inequality is about the actual Gaussian ULA laws, and output finiteness has already been proved. The independent Brownian construction and stationary reference are internal parts of this proof.



Apply the synchronous finite-coupling estimate to the fully constructed reference family, then take the initial cost infimum. The small-step factor is \(q=1-mh\), with gap \(mh>0\). The local forcing \((33/20)Lh\sqrt{hd}\) divided by that gap gives \((33/20)(L/m)\sqrt{hd}\). Rewriting the common output as the Gaussian transition iterate proves the stated actual-law bound.

\subsection{\texorpdfstring{The large-\allowbreak{}step canonical Wasserstein bound}{The large-step canonical Wasserstein bound}}


Use the same fully constructed canonical reference in the large-step product theorem. The step assumptions give \(0\le Lh-1<1\) and \(2-Lh>0\), so the bias is \((33/20)Lh\sqrt{hd}/(2-Lh)\). Rewrite its output law as the actual iterated Gaussian ULA law.



The same stationary reference and local defect estimate also apply in the second stable step regime. Here the deterministic factor is \(q=Lh-1\), so the positive gap is \(2-Lh\). Geometric summation gives bias \((33/20)Lh\sqrt{hd}/(2-Lh)\). The canonical output-law identity transfers this inequality to the actual Gaussian transition iterates without changing either the target or initial Wasserstein distance.

\section{\texorpdfstring{Dalalyan bounds for an implemented L\ensuremath{{}^{2}} Gaussian recursion}{Dalalyan bounds for an implemented L² Gaussian recursion}}
\label{proof-module-UnadjustedLangevin-GaussianULA-lean}


This endpoint applies the canonical-law theorem to an actual chain on a supplied probability space. Its inputs are the original potential assumptions, a square-integrable initialization, independent standard Gaussian scalar innovations and independence of the initial variable from the innovation sequence.



A chain driven by independent standard Gaussian vectors, independent of its initial state, has the same marginal recursion as the canonical law construction. Its finite initial transport distance supplies the required second-moment control relative to the Gibbs law. The canonical Wasserstein estimate therefore applies to this concrete realization at every finite index. The law identification is needed: coordinate marginal Gaussianity alone would not establish the independence required for the transition recursion.

\subsection{\texorpdfstring{Normalization of the initial law}{Normalization of the initial law}}


The representative of the \(L^2\) initialization is measurable. Its pushforward of a probability measure therefore has total mass one, giving the probability hypothesis required by the law-level canonical ULA theorem.



A square-integrable initialization has an almost-surely measurable representative. Its law is the pushforward of the supplied source probability and therefore has mass one. This normalization is the law-level prerequisite for applying the canonical transition theorem. It is independent of any Gaussian innovation assumption or finite-time convergence calculation.

\subsection{\texorpdfstring{Finiteness of the initial transport distance}{Finiteness of the initial transport distance}}


The actual initialization is in \(L^2\), so its law has a finite second moment. Strong convexity proves the same for Gibbs and proves its normalization. The independent product coupling of these two laws has finite quadratic cost, hence their Wasserstein distance is finite.



Changing variables shows that the initial law has the same finite second moment as its random-variable realization. The Gibbs target has a finite second moment by positive strong convexity. Their independent product coupling has finite mean squared separation, bounded by twice the sum of those moments. Thus the initial Wasserstein distance is finite and need not be supplied as an additional assumption.

\subsection{\texorpdfstring{The small-\allowbreak{}step bound for independent Gaussian innovations}{The small-step bound for independent Gaussian innovations}}


Assemble the jointly independent standard scalar Gaussian coordinates into iid multivariate standard innovations. Initialization independence and the \(L^2\) representative theorem identify the construction's law with \(P_h^k(P_{X_0})\). Derive finite initial transport from its moments, then apply the unconditional canonical small-step theorem and rewrite the law back to the construction.



Coordinate Gaussian laws together with full time-coordinate independence determine iid standard vector innovations. Initialization independence then identifies the supplied chain law with the canonical Gaussian transition iterate. Its initial transport is finite by the moment argument. Substituting this law equality into the constructed small-step theorem gives the realization's finite-time bound, with normalization and output finiteness already justified.

\subsection{\texorpdfstring{The large-\allowbreak{}step bound for independent Gaussian innovations}{The large-step bound for independent Gaussian innovations}}


Use the same construction/\allowbreak{}Markov-law identification and derived finite initial transport, then apply the canonical large-step theorem. The numerical construction therefore inherits contraction factor \(Lh-1\) and bias \((33/20)Lh\sqrt{hd}/(2-Lh)\) throughout the stated stable large-step regime.



The law-identification argument is unchanged in the large-step regime; only the deterministic contraction factor and geometric gap change. Apply the canonical bound with \(q=Lh-1\) and \(1-q=2-Lh>0\), then replace the transition iterate by the supplied realization law. The result retains the exact initial Wasserstein term and the explicit bias throughout the stated stable interval.

\section{\texorpdfstring{The Langevin generator annihilates Gibbs averages}{The Langevin generator annihilates Gibbs averages}}
\label{proof-module-UnadjustedLangevin-Generator-lean}


For compactly supported \(C^2\) test functions, \(\mathcal L\varphi=\Delta\varphi-\langle\nabla U,\nabla\varphi\rangle\) has zero Gibbs mean. These identities establish infinitesimal invariance; they alone do not assert stationarity of a stochastic process.



For diffusion coefficient \(\sqrt2\), the generator is \(\mathcal L\varphi=\Delta\varphi-\langle\nabla U,\nabla\varphi\rangle\). Multiplying its directional terms by \(e^{-U}\) makes each one a directional derivative of a weighted first derivative. Compact support removes the boundary integral at infinity and supplies absolute integrability. The conclusion \(\int\mathcal L\varphi\,d\pi_U=0\) is infinitesimal invariance on the stated test class; identifying a full time-dependent law needs additional evolution arguments.

\subsection{\texorpdfstring{The generator in one direction}{The generator in one direction}}


For a fixed direction \(v\), define \(L_v\varphi(x)=D(D\varphi(\cdot)[v])(x)[v]-\langle\nabla U(x),v\rangle D\varphi(x)[v]\). The first term is a directional second derivative of the test; the second is the directional gradient drift. No Hessian of \(U\) is part of this definition.



For a direction \(v\), the second derivative contributes \(D^2\varphi[v,v]\), while the drift contributes \(-\langle\nabla U,v\rangle D\varphi[v]\). Their common direction is what makes a weighted directional derivative possible.

\subsection{\texorpdfstring{The Langevin generator}{The Langevin generator}}


Choose the standard orthonormal basis and set \(L\varphi=\sum_iL_{e_i}\varphi\). The finite sum makes the construction meaningful in dimension zero and gives a direct path to directional integration by parts. A subsequent identity proves that this is the usual basis-independent \(\Delta\varphi-\langle\nabla U,\nabla\varphi\rangle\).



Summing over an orthonormal basis yields the Euclidean diffusion operator. The basis is finite, so all sums in the definition and subsequent integrations are finite sums.

\subsection{\texorpdfstring{The weighted Gibbs integral formula}{The weighted Gibbs integral formula}}


Under the positive finite partition hypothesis, changing measure gives, for every integrable test,

\[\int f\,d\pi_U=Z_U^{-1}\int e^{-U(x)}f(x)dx.\]

 Continuity of the potential makes the density measurable, and the same finite scalar \(Z_U^{-1}\) multiplies each weighted integral.



Changing to Gibbs inserts the same scalar normalization \(Z^{-1}\) into every weighted integral. Positivity and finiteness of \(Z\) make this the usual density-change formula for signed integrable tests.

\subsection{\texorpdfstring{Integrability of a compactly supported test}{Integrability of a compactly supported test}}


A continuous compactly supported \(f\) gives a continuous compactly supported product \(e^{-U}f\), hence a Lebesgue-integrable function. The weighted-density integrability criterion and multiplication by the real normalization factor transfer this fact to \(\pi_U\). This explicitly supplies the integrability needed by subsequent linearity arguments.



A continuous function with compact support is bounded on its support. Multiplication by the continuous Gibbs weight preserves compact support, and hence provides a Lebesgue-integrable weighted function.

\subsection{\texorpdfstring{The derivative of the Gibbs weight}{The derivative of the Gibbs weight}}


Apply the chain rule to \(x\mapsto e^{-U(x)}\): \(D(e^{-U})(x)[v]=-e^{-U(x)}DU(x)[v]\). The real Hilbert gradient identity rewrites \(DU(x)[v]\) as \(\langle\nabla U(x),v\rangle\). Only differentiability of \(U\) is used.



The minus sign arises from differentiating \(-U\), and the exponential factor remains unchanged. Identifying the directional derivative of \(U\) with its gradient pairing fixes the sign of the generator drift term.

\subsection{\texorpdfstring{The weighted directional generator integral}{The weighted directional generator integral}}


Differentiate \(e^{-U(x)}D\varphi(x)[v]\) in direction \(v\). The product rule gives exactly \(e^{-U(x)}L_v\varphi(x)\). Since derivatives of a compactly supported \(C^2\) test remain compactly supported, this weighted first derivative is integrable and has zero integral by integration by parts. The proof establishes continuity and compact support before invoking that theorem.



The product derivative is \(D_v(e^{-U}D_v\varphi)=e^{-U}(D_v^2\varphi-D_vU\,D_v\varphi)\). Its integral vanishes because the product has compact support, giving the weighted directional invariance identity.

\subsection{\texorpdfstring{The directional Gibbs generator integral}{The directional Gibbs generator integral}}


Apply the weighted-to-Gibbs integral formula to \(L_v\varphi\). Its weighted Lebesgue integral is zero, so multiplication by the normalization factor still gives zero. This is infinitesimal invariance for a single direction and does not yet identify the law of a time-evolving process.



Dividing the zero weighted integral by the positive partition leaves zero. This is the normalized directional invariance statement, retaining the same compactly supported test and requiring no new derivative calculation.

\subsection{\texorpdfstring{Continuity of the directional generator}{Continuity of the directional generator}}


The \(C^2\) hypothesis makes the first and second derivatives of \(\varphi\) continuous. The gradient of \(U\) is continuous by its Lipschitz hypothesis. Evaluating continuous linear maps in the fixed direction \(v\), multiplying the scalar factors, and subtracting therefore gives continuity of \(L_v\varphi\).



The second derivative of the test and the gradient of the potential are continuous. Evaluation in a fixed direction, multiplication, and subtraction therefore preserve continuity of the directional generator.

\subsection{\texorpdfstring{Compact support of the directional generator}{Compact support of the directional generator}}


Outside the compact support of \(\varphi\), the test is locally zero, so its first and second derivatives vanish there. Both terms of \(L_v\varphi\) therefore vanish. The support inclusion puts the support of the generator inside the same compact set; no support property of \(U\) is required.



Outside the support of a smooth compactly supported test, the test vanishes on a neighborhood, so all its derivatives vanish there. Both generator terms therefore have support inside that same compact set.

\subsection{\texorpdfstring{Integrability of the directional generator}{Integrability of the directional generator}}


Combine continuity of \(L_v\varphi\) with preservation of compact support. The established compact-support Gibbs integrability theorem then gives absolute integrability. This permits ordinary integral algebra, even though the ambient generator is not bounded uniformly for arbitrary tests.



Continuity alone would not ensure integrability on an unbounded space. Here compact support supplies the needed localization, and the Gibbs density is continuous and bounded on that compact set.

\subsection{\texorpdfstring{Integrability of the full generator}{Integrability of the full generator}}


The full generator is the finite sum of its directional pieces. Each piece is integrable under \(\pi_U\), so the finite-sum integrability theorem applies. Finiteness of the basis is the reason no exchange of an infinite series and integral is required.



Each directional summand has an absolutely integrable generator value. The finite orthonormal-basis sum is therefore integrable, allowing the full generator integral to be distributed term by term.

\subsection{\texorpdfstring{The weighted full generator integral}{The weighted full generator integral}}


Expand the finite basis sum defining \(L\varphi\) and distribute the Gibbs weight over it. Compact support supplies integrability of every summand, allowing \(\int\sum_i e^{-U}L_{e_i}\varphi=\sum_i\int e^{-U}L_{e_i}\varphi\). Every directional integral vanishes, hence so does their sum.



Distribute the Gibbs weight through the finite basis sum and integrate. Every directional integral has already been shown to be zero, so their sum vanishes without an infinite-series interchange.

\subsection{\texorpdfstring{Infinitesimal Gibbs invariance}{Infinitesimal Gibbs invariance}}


Normalization changes \(\int L\varphi\,d\pi_U\) into the real inverse partition times its weighted Lebesgue integral. The latter was proved to be zero. This establishes the full Gibbs infinitesimal invariance identity for compactly supported \(C^2\) tests.



The normalization formula identifies the Gibbs generator integral with \(Z^{-1}\) times its weighted Lebesgue counterpart. The latter is zero, proving infinitesimal invariance for the full generator.

\subsection{\texorpdfstring{The Laplacian-\allowbreak{}minus-\allowbreak{}drift formula}{The Laplacian-minus-drift formula}}


The sum of the directional second derivatives over an orthonormal basis is the Laplacian. For the drift sum, \(D\varphi(x)[e_i]=\langle\nabla\varphi(x),e_i\rangle\); Parseval turns the sum of products with \(\langle\nabla U(x),e_i\rangle\) into \(\langle\nabla U(x),\nabla\varphi(x)\rangle\). The proof rewrites the basis definition into this conventional formula pointwise.



Parseval's identity gives \(\sum_i\langle\nabla U,e_i\rangle\langle\nabla\varphi,e_i\rangle=\langle\nabla U,\nabla\varphi\rangle\). The analogous sum of second directional derivatives is the Laplacian.

\subsection{\texorpdfstring{The stationary identity in conventional notation}{The stationary identity in conventional notation}}


Replace the conventional expression \(\Delta\varphi-\langle\nabla U,\nabla\varphi\rangle\) by the equal basis-defined generator inside the integral. The established zero-generator integral then closes the result. The rewrite makes the displayed mathematical identity refer to the same function that the preceding calculation integrated.



Substitution of the pointwise generator formula yields the conventional stationary identity. It applies on the stated compactly supported test class; stationarity of an actual evolution is established later by identifying its laws.

\section{\texorpdfstring{Physical-\allowbreak{}time generator and Gibbs integration identity}{Physical-time generator and Gibbs integration identity}}
\label{proof-module-IPLA-Generator-lean}


In scaled coordinates the continuous physical-time model is \(dZ_t=-N^{-1}\nabla V_N(Z_t)dt+\sqrt{2/N}\,dB_t\). The generator computation and Gibbs integration identity are consistent with this time scaling. Infinitesimal invariance is a useful analytic identity, while full stationarity requires the separate construction of the Langevin evolution.




For scaled state \(z\), physical IPLA time is \(N\) times the ordinary lifted Langevin time. The generator is consequently \(N^{-1}(\Delta-\nabla V_N\cdot\nabla)\).

\subsection{\texorpdfstring{Identify the generator on a twice differentiable test function}{Identify the generator on a twice differentiable test function}}


Define \(\mathcal A_N\varphi=N^{-1}\mathcal L_{V_N}\varphi\), where \(\mathcal L_V\) is the shared Langevin generator. For \(C^2\) test functions the shared coordinate definition equals the Laplacian minus the gradient inner product, hence

\[\mathcal A_N\varphi(z)=\frac1N\left(\Delta\varphi(z)-\langle\nabla V_N(z),\nabla\varphi(z)\rangle\right).\]

This is a direct identity of operators on the scaled state; no invariant-measure conclusion is built into its definition.





The diffusion covariance is \((2/N)I\), so the second-order term in the generator is \(\Delta/N\). The drift contributes \(-\langle\nabla V_N,\nabla\varphi\rangle/N\). Adding these terms gives the displayed operator. The factor rescales time uniformly; it changes relaxation speed but leaves the Gibbs density associated with \(V_N\) unchanged.

\subsection{\texorpdfstring{Derive integrability and infinitesimal invariance}{Derive integrability and infinitesimal invariance}}


For compactly supported \(C^2\) \(\varphi\), the shared Gibbs integration-by-parts theorem proves that \(\mathcal L_{V_N}\varphi\) is integrable and has Gibbs integral zero. Its potential hypotheses follow from differentiability and Lipschitzness of the joint gradient by the lifted geometry. Multiplication by the constant \(N^{-1}\) preserves integrability and gives

\[\int\mathcal A_N\varphi\,d\Pi_N=\frac1N\int\mathcal L_{V_N}\varphi\,d\Pi_N=0.\]

The theorem also returns the probability property, derived from positive strong convexity. Explicit integrability rules out a zero caused by the default integral value of a nonintegrable function. Full time-evolution stationarity is supplied separately by the shared Langevin construction.




For a \(C^2\) compactly supported test function, the shared Gibbs integration identity gives \(\int[\Delta\varphi-\langle\nabla V_N,\nabla\varphi\rangle]d\Pi_N=0\). Physical IPLA time multiplies this generator by \(N^{-1}\), so linearity of the integral preserves zero. The theorem also proves integrability of the generator before asserting this identity, using the compact support and the proved lifted score regularity. This distinguishes an actual invariance calculation from a formal integral expression whose integrability has not been justified.





The identity comes from the divergence \(\nabla\!\cdot(e^{-V_N}\nabla\varphi)=e^{-V_N}(\Delta\varphi-\langle\nabla V_N,\nabla\varphi\rangle)\). Compact support eliminates the boundary term after integration. Local regularity of the score and the bounded support make the resulting terms integrable. Dividing by the positive Gibbs partition and by \(N\) then gives the zero-expectation statement for the normalized physical generator.

\section{\texorpdfstring{Zero-\allowbreak{}dimensional state spaces have zero transport cost}{Zero-dimensional state spaces have zero transport cost}}
\label{proof-module-UnadjustedLangevin-Degenerate-lean}


On a subsingleton state space every two states coincide, so every coupling has zero quadratic difference cost. Probability laws on such a space therefore have zero Wasserstein distance.



In a singleton state space, any two states agree and every displacement is zero. A joint law therefore has zero quadratic cost, and the product of the two probability measures is an admissible zero-cost coupling. The transport infimum is both nonnegative and at most zero, hence zero. This resolves dimension-zero or other degenerate cases directly, before invoking arguments that cancel a nonzero distance or select a unit direction.

\subsection{\texorpdfstring{Zero coupling cost in a singleton space}{Zero coupling cost in a singleton space}}


In a subsingleton state space every pair satisfies \(x=y\), so the displacement function is identically zero. Its extended \(L^2\) norm is zero under every pair measure, hence every coupling cost vanishes.



The conclusion uses only that the metric distance between any two states is zero. It therefore holds for every coupling, with no moment or regularity assumption on a potential. The quadratic cost integral is the integral of the identically zero function.

\subsection{\texorpdfstring{Zero transport distance in a singleton space}{Zero transport distance in a singleton space}}


The independent product of two probabilities is a coupling, so the family of couplings is nonempty. Every coupling has zero cost by the preceding result; taking their infimum gives \(W_2(\mu,\pi)=0\).



The product coupling demonstrates that the infimum is taken over a nonempty family. Every member of that family has value zero, so the infimum is zero as well. This avoids relying on any convention for an infimum over an empty set.

\subsection{\texorpdfstring{The ordinary zero-\allowbreak{}distance conclusion}{The ordinary zero-distance conclusion}}


Rewrite the extended Wasserstein distance as zero using the subsingleton theorem. Its real conversion is also zero. This covers dimension-zero spaces without invoking the nontrivial-space comparison \(m\le L\).



A zero-dimensional Euclidean space consists of the empty coordinate vector. Both the numerical update and every reference law are supported at that single state. Hence the transport conclusion is immediate and no smoothness-versus-curvature cancellation is necessary.

\section{\texorpdfstring{Zero-\allowbreak{}dimensional parameter and state spaces}{Zero-dimensional parameter and state spaces}}
\label{proof-module-IPLA-Degenerate-lean}


Zero-dimensional parameter and latent spaces are legitimate Euclidean spaces, each consisting of one point. Treating these cases explicitly removes artificial positive-dimension restrictions from the public convergence theorem. The arguments are elementary consequences of vanishing norms and uniqueness of a probability measure on a singleton.




Euclidean space over an empty coordinate type contains a single vector. These elementary cases allow the public theorem to retain independent dimensions that may both be zero.

\subsection{\texorpdfstring{The parameter error vanishes when d\ensuremath{\theta} = 0}{The parameter error vanishes when dθ = 0}}


If \(d_\theta=0\), every parameter vector equals every other one. In particular \(p_\theta z=\theta_*\), so \(\|p_\theta z-\theta_*\|^2=0\) pointwise. Its integral under any measure, and the square root of that integral, are zero. No probability, potential or noise assumption is required.





There is only one candidate parameter center when the parameter space has dimension zero. Its difference from every projected state is the zero vector, so the integrand vanishes before integration. In particular this conclusion holds even if the latent particle system remains nontrivial. The result separates the dimension of the statistical parameter from the dimension required for a nontrivial sampling process.

\subsection{\texorpdfstring{Every term of the full zero-\allowbreak{}dimensional bound vanishes}{Every term of the full zero-dimensional bound vanishes}}


If also \(d_x=0\), the whole state is a singleton. Any two probability measures on it are equal and their \(W_2\) distance is zero. The equilibrium radius \(\sqrt{2d_\theta/(Nm)}\) and discretization factor \(\sqrt{d_x+d_\theta/N}\) are also zero. Substituting these facts reduces the full claimed inequality to \(0\le0\), irrespective of the values of \(m,\gamma,L\). This handles the case in which the nonempty-coordinate analytic theorem cannot be instantiated.




When \(d_\theta=0\), the parameter space is a singleton, hence every parameter error vector and its square vanish. If also \(d_x=0\), the whole state space is a singleton. Any probability law there is the same Dirac law, so its Wasserstein distance to the Gibbs probability is zero. The dimension factors in the remaining terms are also zero, giving the public bound in this case directly, without invoking a theorem requiring a nonempty coordinate index.





A coupling of two singleton probability measures is necessarily concentrated on the diagonal. Its cost is zero, so their Wasserstein distance is zero directly from the definition. Both equilibrium and discretization radii also vanish through their dimension factors. These observations give the same quantitative formula as the nondegenerate theorem, without applying analytic estimates that presume a nonempty coordinate family.

\section{\texorpdfstring{Gibbs Dirichlet form, symmetry and dissipation}{Gibbs Dirichlet form, symmetry and dissipation}}
\label{proof-module-UnadjustedLangevin-GeneratorDirichlet-lean}


For compactly supported \(C^2\) tests \(\varphi,\psi\), the generator obeys \(\int\varphi\mathcal L\psi\,d\pi=-\int\langle\nabla\varphi,\nabla\psi\rangle d\pi\). This provides the symmetric Dirichlet form at the test-function level.



The identity \(\int\varphi\mathcal L\psi\,d\pi_U=-\int\langle\nabla\varphi,\nabla\psi\rangle d\pi_U\) is obtained by moving one derivative onto the weighted test. The derivative of the Gibbs weight supplies exactly the potential-drift term, so no term remains unaccounted for. Compact support makes the integration by parts global without a boundary contribution. Symmetry follows from symmetry of the inner product, and choosing \(\psi=\varphi\) identifies the nonnegative energy whose negative is the generator quadratic form.

\subsection{\texorpdfstring{Differentiating the Gibbs-\allowbreak{}weighted test}{Differentiating the Gibbs-weighted test}}


Differentiate the product \(e^{-U}\varphi\): \(D(e^{-U}\varphi)(x)[v]=e^{-U(x)}(D\varphi(x)[v]-\langle\nabla U(x),v\rangle\varphi(x))\). Combine the product rule with the already proved derivative of the Gibbs weight, then collects the common exponential factor algebraically.



Factoring out \(e^{-U}\) from the product rule leaves the test derivative minus its multiplication by the potential's directional derivative. This is the exact cancellation needed to pair with the generator.

\subsection{\texorpdfstring{Weighted directional integration by parts}{Weighted directional integration by parts}}


Integrate by parts the derivative of \(e^{-U}\varphi\) against \(D\psi[v]\). The weighted product rule separates a \(D\varphi[v]D\psi[v]\) term from the drift term, giving \(\int e^{-U}\varphi L_v\psi=-\int e^{-U}D\varphi[v]D\psi[v]\). Compact support and \(C^2\) regularity justify every derivative and integral involved.



Integrating \(D_v(e^{-U}\varphi D_v\psi)\) gives zero. Expanding this derivative and moving the term involving \(D_v\varphi\) to the other side yields the negative directional energy.

\subsection{\texorpdfstring{The directional Gibbs Dirichlet identity}{The directional Gibbs Dirichlet identity}}


Convert both sides of the weighted directional identity to Gibbs integrals. The same real normalization factor multiplies both integrals, producing \(\int\varphi L_v\psi\,d\pi_U=-\int D\varphi[v]D\psi[v]d\pi_U\). This is a change of measure, with no new integration-by-parts step.



Both weighted integrals carry the same inverse-partition factor. Applying the Gibbs change of density to each side therefore preserves the equality and its negative sign.

\subsection{\texorpdfstring{Integrability of the gradient pairing}{Integrability of the gradient pairing}}


The gradients are continuous; the gradient of the compactly supported test \(\varphi\) vanishes outside its compact support. Thus \(x\mapsto\langle\nabla\varphi(x),\nabla\psi(x)\rangle\) is continuous with compact support. The general compact-test Gibbs integrability theorem applies even though \(\psi\) itself need not have compact support in this statement.



Even if the second test is not compactly supported, the first gradient vanishes outside its compact support. The product of the two gradients is thus localized to a compact set where both are bounded.

\subsection{\texorpdfstring{Integrability of the test-\allowbreak{}generator product}{Integrability of the test-generator product}}


The generator of the \(C^2\) test \(\psi\) is continuous under the Lipschitz-gradient assumption on \(U\). Multiplication by the continuous compactly supported \(\varphi\) makes the product compactly supported, hence Gibbs integrable. This is the exact product needed on the left side of the Dirichlet identity.



The compactly supported multiplier localizes the continuous generator value. This makes the left side of the Dirichlet identity absolutely integrable before any integration-by-parts algebra is used.

\subsection{\texorpdfstring{The full Dirichlet identity}{The full Dirichlet identity}}


Sum the directional Dirichlet identities over the orthonormal basis. Finite-sum integrability allows the integral to pass through the sum, and Parseval identifies \(\sum_iD\varphi[e_i]D\psi[e_i]=\langle\nabla\varphi,\nabla\psi\rangle\). Therefore \(\int\varphi L\psi\,d\pi_U=-\int\langle\nabla\varphi,\nabla\psi\rangle d\pi_U\).



Summing the directional energies reconstructs the gradient inner product through orthonormal expansion. The finite sum and absolute integrability justify carrying this reconstruction through the integral.

\subsection{\texorpdfstring{Symmetry of the generator}{Symmetry of the generator}}


Apply the Dirichlet identity in both orders \((\varphi,\psi)\) and \((\psi,\varphi)\). Symmetry of the real inner product makes the resulting energy integrands equal pointwise. Hence \(\int\varphi L\psi\,d\pi_U=\int\psi L\varphi\,d\pi_U\).



Exchanging the two tests exchanges the two gradients in the energy pairing. Symmetry of the real inner product makes these energies identical, proving symmetry of the generator on the stated domain.

\subsection{\texorpdfstring{The generator self-\allowbreak{}pairing}{The generator self-pairing}}


Specialize the bilinear identity to \(\psi=\varphi\), then rewrite \(\langle\nabla\varphi,\nabla\varphi\rangle=\|\nabla\varphi\|^2\). This identifies the quadratic generator form exactly: \(\int\varphi L\varphi\,d\pi_U=-\int\|\nabla\varphi\|^2d\pi_U\).



With the two tests equal, their gradient pairing is the squared gradient norm. This identifies the generator quadratic form exactly, rather than merely bounding its sign.

\subsection{\texorpdfstring{Dissipativity of the generator}{Dissipativity of the generator}}


The self-pairing identity expresses the generator quadratic form as minus an integral of a nonnegative square. The integral is nonnegative, so its negative is at most zero. The proof uses the actual integrable energy established earlier, giving dissipativity on compactly supported \(C^2\) tests.



The squared gradient norm is nonnegative pointwise and integrable. Its integral is therefore nonnegative, so the self-pairing generator integral is nonpositive. This is the announced dissipativity statement.

\section{\texorpdfstring{A weak forward-\allowbreak{}equation interface and a conditional stationarity route}{A weak forward-equation interface and a conditional stationarity route}}
\label{proof-module-UnadjustedLangevin-FokkerPlanck-lean}


The weak Fokker–Planck equation is formulated by integrating \(\mathcal L\varphi\) against time-dependent probability measures. The final theorem in this module assumes uniqueness of this weak equation explicitly; the main ULA development proves stationarity by the separate Metropolis approximation route.



The weak equation is \(\int\varphi\,d\rho_t-\int\varphi\,d\rho_0=\int_0^t\int\mathcal L\varphi\,d\rho_sds\) for the stated compactly supported tests. For the constant Gibbs curve the left side is zero and the generator integral on the right is zero. Hence it is a weak solution. Comparing another curve to it yields stationarity only under the explicitly supplied uniqueness principle for this same solution class.

\subsection{\texorpdfstring{The integrated weak Fokker–Planck equation}{The integrated weak Fokker–Planck equation}}


A curve \(\rho_t\) is in this weak solution class when it is probability valued, has the stated moment and integrability properties, and satisfies the integrated generator equation against every compactly supported \(C^2\) test. The assumptions keep the analytic conditions visible, preventing a scalar expression from substituting for an unproved integrable quantity.



The integrated test identity avoids assuming a density or a pointwise time derivative of the law. Its moment and integrability conditions ensure that the generator expectation and its time integral are meaningful.

\subsection{\texorpdfstring{The constant Gibbs weak solution}{The constant Gibbs weak solution}}


Use the proved Gibbs normalization and second moment for the probability and moment fields. Compact-support generator integrability supplies the test-integral fields. Finally \(\int L\varphi\,d\pi_U=0\) makes the time integral vanish, while the constant curve has zero change in \(\int\varphi\,d\rho_t\). These verify each field of the weak equation for the stationary candidate.



For a constant law curve, every test expectation has zero increment. The Gibbs generator identity makes its integrated right side zero as well, verifying the actual weak equation for that candidate.

\subsection{\texorpdfstring{The displayed weak-\allowbreak{}solution uniqueness principle}{The displayed weak-solution uniqueness principle}}


The uniqueness assumption states that two probability-valued weak solutions in the displayed moment and test class, with the same initial law, agree throughout \([0,T]\). It is an explicit analytic premise of the comparison below. Its validity must be established for that solution class before the conditional stationarity argument is applied.



The uniqueness condition refers to exactly this test class and its accompanying moment assumptions. Stating it records the analytic premise required by the following comparison; it does not establish uniqueness independently.

\subsection{\texorpdfstring{Stationarity under that uniqueness principle}{Stationarity under that uniqueness principle}}


Both curves belong to the same weak-solution class and agree initially. The supplied uniqueness principle therefore identifies them throughout the interval. This conclusion remains conditional on that uniqueness premise.

\section{\texorpdfstring{The standalone strongly convex IPLA error bound}{The standalone strongly convex IPLA error bound}}
\label{proof-module-IPLA-Theorem-lean}


The final inequality separates three causes of parameter error: finite-particle equilibrium spread, incomplete relaxation from the initial law, and Euler discretization. Each is controlled by a different proved ingredient. The enlarged geometry and the exact algorithm-law identification allow these ingredients to be combined in the physical parameter and particle dimensions stated in the theorem.




For the fixed observed-data slice \(U_y\), assume differentiability, joint \(m\)-strong convexity with \(m>0\), \(L\)-Lipschitz gradient, \(N>0\), and \(0<\gamma\le2/(m+L)\). Initialization is square integrable, and innovations have the stated independent standard Gaussian coordinate laws and initialization independence. Let \(\theta_*\) maximize the actual latent marginal. The theorem applies to every pair of independent natural dimensions, including zero.

\subsection{\texorpdfstring{Combine lifted transport with proved equilibrium concentration}{Combine lifted transport with proved equilibrium concentration}}


Suppose the state has at least one coordinate. The lifted geometry gives \(Nm\)-strong convexity and \(NL\)-smoothness in dimension \(D=d_\theta+Nd_x\). The Gibbs theorem normalizes \(\Pi_N\), while square-integrable initialization and Gibbs moments imply finite initial transport. Apply the shared Gaussian ULA theorem at step \(h=\gamma/N\), using the supplied Gaussian hypotheses. The step restriction rescales exactly and gives

\[W_2(\rho_k,\Pi_N)\le(1-m\gamma)^kW_2(\mu_0,\Pi_N)+\frac{33}{20}\frac Lm\sqrt{d_x+d_\theta/N}\sqrt\gamma.\]

The ULA output-finiteness result and the proved law identity ensure this is finite actual transport, before any real conversion. The equilibrium concentration theorem gives \(\operatorname{RMS}_{\Pi_N}(p_\theta-\theta_*)\le\sqrt{2d_\theta/(Nm)}\). For any coupling of \(\rho_k\) and \(\Pi_N\), Minkowski and \(\|p_\theta\|\le1\) imply the sampled parameter RMS is at most coupling RMS plus that equilibrium radius. Taking the coupling infimum and using finiteness yields

\[\left(\int\|p_\theta z-\theta_*\|^2d\rho_k\right)^{1/2}\le\sqrt{\frac{2d_\theta}{Nm}}+(1-m\gamma)^kW_2(\mu_0,\Pi_N)+\frac{33}{20}\frac Lm\sqrt{d_x+\frac{d_\theta}{N}}\sqrt\gamma.\]

Stationarity, finite transport and concentration are derived here through the shared and particle theorems; they are not additional theorem hypotheses.




The proof first constructs target probability and initial finite transport from the lifted strong-convexity assumptions and \(L^2\) initialization. Gaussian ULA then gives the state-space transport rate, with step \(\gamma/N\), modulus \(Nm\), score constant \(NL\), and dimension \(d_\theta+Nd_x\). The law adapter identifies its output with the actual supplied IPLA recursion and proves the output transport distance is finite.





The projection step can be read directly from any coupling \((X,Y)\) with laws \(\rho_k,\Pi_N\): \(\|p_\theta X-\theta_*\|_2\le\|X-Y\|_2+\|p_\theta Y-\theta_*\|_2\). The second term is controlled by concentration, and taking the infimum of the first gives transport. The shared Gaussian theorem bounds that transport. Consequently the three terms are added after all laws and moments have been identified, rather than assuming a separate parameter convergence estimate.

\subsection{\texorpdfstring{Include every zero-\allowbreak{}dimensional case in the public theorem}{Include every zero-dimensional case in the public theorem}}




The full-state index is empty precisely when both dimensions vanish, provided \(N>0\). If parameter dimension is positive there is a parameter coordinate; if latent dimension is positive there is a coordinate in any one of the nonempty particle blocks. Thus every remaining case belongs to the nonempty-state analytic theorem. The singleton case has zero parameter error and zero transport, giving exactly the same public formula without an implicit dimension restriction.

\section{\texorpdfstring{Sharp parameter concentration and the IPLA RMS bound}{Sharp parameter concentration and the IPLA RMS bound}}
\label{proof-module-ProximalLangevin-SharpIPLA-lean}


The sharp continuous-potential concentration estimate improves the equilibrium contribution in the particle theorem from twice dimension over curvature to dimension over curvature. The transport and Gaussian discretization arguments are unchanged. Exact marginal identification ensures that the improved estimate uses parameter dimension and the \(N\)-amplified marginal curvature.




A differentiable \(m\)-strongly convex pair potential \(U\), \(m>0\), has a marginal potential whose \(N\)-replicate target is \(Nm\)-strongly convex. The sharp dilation theorem improves its parameter second moment to \(d_\theta/(Nm)\).

\subsection{\texorpdfstring{Concentration of the actual marginal target}{Concentration of the actual marginal target}}


For a continuous \(a\)-strongly convex potential \(f\) with minimizer \(x_\star\), the new dilation argument proves


\[
\int\|x-x_\star\|^2\,\pi(dx)\le\frac d a,
\qquad\pi(dx)\propto e^{-f(x)}dx.
\]


The proof uses a sharp radial growth inequality and a change of variables under dilation, rather than assuming an integration-by-parts identity for a nonsmooth density. It therefore applies directly to the continuous strongly convex Gibbs potentials needed here.



The existing Prékopa argument preserves strong convexity when latent variables are integrated out. Raising the marginal density to the \(N\)th power multiplies that modulus by \(N\). Applying the sharp estimate to its actual marginal potential gives


\[
W_2(\Pi_{\lambda,N}^{\theta},\delta_{\theta_{\lambda,\star}})
\le\sqrt{\frac{d_\theta}{Nm}}.
\]


This matches the constant in v3 equation (17) and Lemma A.14. The previous IPLA theorem's \(\sqrt{2d_\theta/(Nm)}\) estimate would have been weaker; the new result proves the sharper constant rather than labeling that factor as an error in the paper.




The exact parameter marginal of the lifted Gibbs measure is \(M(\theta)^N\) after normalization. Its negative log potential is \(NF(\theta)\), and the earlier Prékopa argument proves that it is \(Nm\)-strongly convex with minimizer \(\theta_*\). Sharp Gibbs concentration in dimension \(d_\theta\) gives \(\int\|\theta-\theta_*\|^2d\Pi_N\le d_\theta/(Nm)\). The pushforward identity transfers this same moment to the lifted state law before a square root is taken.





The parameter marginal density is proportional to \(e^{-N\Phi}\), so its mode is the actual statistical maximizer and its curvature is \(Nm\). Apply the sharp theorem in dimension \(d_\theta\), obtaining the second moment \(d_\theta/(Nm)\). The exact pushforward identity transfers this same moment to the enlarged target. Taking a nonnegative square root yields the improved equilibrium RMS radius.

\subsection{\texorpdfstring{Combine sharp concentration with the actual Gaussian-\allowbreak{}chain law}{Combine sharp concentration with the actual Gaussian-chain law}}


For \(N>0\), \(\gamma>0\), \(\gamma\le2/(m+L)\), Gaussian independent innovations and \(L^2\) initialization, use the scaled ULA transport theorem and the nonexpansive parameter projection. The projected triangle inequality gives

\[\operatorname{RMS}_n\le\sqrt{\frac{d_\theta}{Nm}}+(1-m\gamma)^nW_2(\rho_0,\Pi_N)+\frac{33}{20}\frac Lm\sqrt{d_x+d_\theta/N}\sqrt\gamma.\]

The private lemma handles nonempty coordinate sets. The public theorem splits off the empty-coordinate case, where the parameter norm and its error vanish, so no artificial positive-dimension assumption remains.






For any coupling of the current state and Gibbs reference, project the difference with the norm-one parameter map and add the equilibrium parameter error. Minkowski then bounds current RMS by transport plus the sharp radius. Insert the established state-space Gaussian Euler bound. The zero-dimensional parameter case has zero error, and the singleton-state case has zero transport, so the same public inequality covers all independent natural dimensions.

\section{\texorpdfstring{The normalized kinetic Gibbs candidate and sharp parameter concentration}{The normalized kinetic Gibbs candidate and sharp parameter concentration}}
\label{proof-module-KIPLMC-Target-lean}


The velocity-first phase target is the product of an explicit Gaussian thermal velocity measure and the already normalized lifted Gibbs position measure. Its parameter marginal is the powered likelihood, and its concentration radius has the paper's sharp dimension coefficient. These results concern a concrete probability measure; the section deliberately does not claim invariance of that measure under a kinetic diffusion.



The phase target combines explicit thermal Gaussian velocities with the normalized lifted Gibbs position law. Its powered-likelihood parameter marginal supplies the sharp equilibrium concentration around the unique marginal maximizer.

\subsection{\texorpdfstring{Construct thermal velocities as a Gaussian pushforward}{Construct thermal velocities as a Gaussian pushforward}}


Start with the standard Gaussian on the scaled particle Euclidean space and map \(v\mapsto v/\sqrt N\). thermal Gaussian velocity measure defines this measure, giving covariance \(I/N\) for positive \(N\). The probability proof uses measurability of the scalar map, evaluates its preimage of the whole space, and obtains mass one. The definition also makes sense at \(N=0\) under the convention \(0^{-1}=0\), but the physical target and subsequent normalization theorems require \(N>0\). No unspecified Gaussian velocity law is assumed.



The velocity pushforward is a concrete centered Gaussian with covariance \(I/N\) for positive particle count. Its finite moments follow from Gaussian scaling. This gives both the exact thermal initialization and the equilibrium velocity factor used in the particle coupling proof.

\subsection{\texorpdfstring{Normalize the product and compute its position marginal}{Normalize the product and compute its position marginal}}


the kinetic Gibbs target is \(\mathcal N(0,I/N)\otimes\pi_N\), with velocity first. Joint strong convexity and differentiability of the original pair potential imply that the shared lifted Gibbs target \(\pi_N\) is a probability measure. Both factors have total mass one, so their product does as well. Mapping the product by its second projection integrates out the first factor, whose total mass is one, so the result is exactly \(\pi_N\). The proof uses the second marginal of a product measure and simplifies the mass scalar. This computes the normalized marginal rather than merely comparing unnormalized density expressions.



The position density is proportional to \(e^{-V_N}\), while the kinetic velocity density is proportional to \(e^{-N\|v\|^2/2}\). Their product is the physical-temperature phase target. Product normalization immediately gives its independence structure and exact position marginal.

\subsection{\texorpdfstring{Identify the powered likelihood parameter marginal}{Identify the powered likelihood parameter marginal}}


Map the phase target to \(\theta\) by composing position projection with the parameter continuous linear map. The pushforward composition theorem identifies this with first taking the position marginal and then the parameter marginal. The preceding result gives the exact lifted Gibbs position law; the shared IPLA theorem identifies its parameter marginal as \(p_\theta(y)^N\) after normalization. Measurability of both maps is provided explicitly. The theorem therefore connects the concrete phase target with the same normalized marginal target used by the concentration proof.



Integrating latent particles factorizes conditionally on the common parameter: each contributes the same marginal likelihood factor. Their product is the \(N\)-th power of that likelihood. The normalized pushforward identity makes this statistical interpretation precise.

\subsection{\texorpdfstring{Construct the unique marginal maximizer and its sharp second moment}{Construct the unique marginal maximizer and its sharp second moment}}


Joint strong convexity makes the negative log marginal strongly convex after latent integration. The shared IPLA theorem supplies existence and uniqueness of its minimizer, equivalently of the marginal likelihood maximizer. For that \(\theta_*\), the sharp concentration theorem gives \(\int\|\theta-\theta_*\|^2d\pi_\Theta\le d_\theta/(Nm)\). These results apply to the concrete target for the target, not an abstract measure with a stated concentration field. The last theorem expresses the square-root bound in the existing extended-nonnegative transport notation as \(\left(\int\|P_\theta z-\theta_*\|^2\,d\pi_N(z)\right)^{1/2}\le\sqrt{d_\theta/(Nm)}\).



The parameter concentration scale is \(N^{-1/2}\), with the sharp dimension coefficient \(\sqrt{d_\theta/m}\). It describes equilibrium error relative to the likelihood maximizer. Numerical discretization and finite-time initialization errors must be added separately in an algorithmic bound.

\section{\texorpdfstring{Shared-\allowbreak{}noise cancellation and uniqueness of the integral solution}{Shared-noise cancellation and uniqueness of the integral solution}}
\label{proof-module-KineticLangevin-IntegralCoupling-lean}


Individual kinetic velocity paths contain continuous noise and need not be differentiable. Their difference under the same additive forcing is differentiable from the right, because the forcing cancels in the integral equations themselves. This section derives that assertion and proves uniqueness by Gronwall, including the initial time.



Cancellation is performed at the level of integral equations, where Brownian differentiability is unnecessary. The resulting deterministic difference equation supports both uniqueness and the transformed contraction theorem.

\subsection{\texorpdfstring{Subtract the two integral equations}{Subtract the two integral equations}}


Suppose \(Y_t-Y_0=-\int_0^t g(Y_s)ds+B_t\) and \(Z_t-Z_0=-\int_0^t g(Z_s)ds+B_t\), with continuous paths and a Lipschitz drift. Subtraction gives

\[Y_t-Z_t=Y_0-Z_0-\int_0^t[g(Y_s)-g(Z_s)]ds.\]

The two drift integrals exist on every compact interval because their integrands are continuous. First, the argument establishes interval integrability, distributes the difference through the integral, rewrites both supplied equations as expressions for the paths, and cancels the common forcing with the additive-group calculation. No derivative of the noise is taken.



The initial phase difference remains after subtraction, but the common forcing cancels identically at every time. The difference is an integral of the drift difference. This cancellation is valid for every continuous forcing path and does not use a distributional property of the noise.

\subsection{\texorpdfstring{Differentiate after cancellation, including the initial endpoint}{Differentiate after cancellation, including the initial endpoint}}


The integrand \(g(Y_s)-g(Z_s)\) is continuous. The fundamental theorem of calculus therefore differentiates the integral expression, producing the right derivative \(-(g(Y_t)-g(Z_t))\). The identity is known only for \(t\in[0,T]\); consequently the mathematically appropriate conclusion at zero is a derivative within \([t,\infty)\). The proof builds an eventual neighborhood lying inside the valid time interval and transfers the derivative from the integral formula using eventual equality. This avoids imposing an arbitrary differentiable extension of the stochastic path to negative time.



The integrand is continuous because the force and the two solution paths are continuous. The fundamental theorem of calculus therefore gives an ordinary derivative in the interior and a right derivative at zero. These are exactly the derivatives used by the transformed-energy calculation.

\subsection{\texorpdfstring{Uniqueness from the global Lipschitz bound}{Uniqueness from the global Lipschitz bound}}


For \(D=Y-Z\), the derived derivative obeys \(\|D\prime_t\|\le L\|D_t\|\). Equal initial conditions give \(\|D_0\|=0\). Gronwall then yields \(\|D_t\|\le0\) throughout the finite horizon. The proof passes the right-derivative theorem to its norm version of Gronwall, simplifies the zero initial bound, and uses definiteness of the norm to recover \(Y_t=Z_t\). This result supplies uniqueness for any continuous additive forcing, including the forcing used in the kinetic Picard construction.



Let \(D(t)\) be the norm of the difference of two solutions with the same initialization and input. Global Lipschitz continuity gives \(D(t)\le L\int_0^tD(s)ds\). Grönwall yields \(D(t)=0\), so every selected solution agrees with every other one on the horizon.

\section{\texorpdfstring{The transformed kinetic energy}{The transformed kinetic energy}}
\label{proof-module-KineticLangevin-Geometry-lean}


The phase convention is \((v,x)\), velocity first. The paper's transformation is \(T_{\gamma,a}(v,x)=(v+(\gamma-a)x,-v-ax)\). Its squared energy is the sum of the two squared Hilbert norms. We equip the phase pair with the Euclidean product norm \(\|(v,x)\|_2^2=\|v\|^2+\|x\|^2\), so the transformed energy is a Hilbert-space squared norm. The identities below also apply with the underlying Hilbert space equal to an L\ensuremath{{}^{2}} space of random vectors.



These are deterministic Hilbert-space identities. Their stochastic use requires only that the coupled position and velocity differences are square integrable; the same inner-product calculations then hold after expectation, with the corresponding root-mean-square norm.



The transformed norm is the single metric used to initialize, propagate, and project coupling errors. Its explicit inverse also explains which parameter choices are admissible and where prefactors become singular.

\subsection{\texorpdfstring{Define the change of variables and its nonnegative energy}{Define the change of variables and its nonnegative energy}}


Set \(y=v+(\gamma-a)x\), \(z=-v-ax\), and \(\mathcal E=\|y\|^2+\|z\|^2\). These are the two coordinates used to expose cancellation in the kinetic drift. The definition the linear transformation constructs the pair; the transformed energy explicitly sums their squares. Unfolding the latter leaves a sum of nonnegative terms, which nonnegativity proves without any sign restriction on \(\gamma\) or \(a\). Thus the energy construction is valid before selecting the admissible friction regime.



The state here may be a difference of two trajectories. In that application the same formula measures how far the trajectories are apart, rather than their absolute energy. Positivity is immediate, while strict definiteness additionally needs the transformation to be invertible.

\subsection{\texorpdfstring{Recover position and velocity}{Recover position and velocity}}


Adding the new coordinates gives \(y+z=(\gamma-2a)x\). The proof expands the scalar difference with distributivity of scalar multiplication and uses vector-space algebra to normalize the vector-space expression. If \(\gamma-2a\ne0\), multiply by its inverse to recover position. A second identity is \(ay+(\gamma-a)z=-(\gamma-2a)v\), giving \(v=-(\gamma-2a)^{-1}(ay+(\gamma-a)z)\). First, the argument proves this vector identity, then rewrites nested scalar multiplication and cancels the nonzero denominator. The hypothesis excludes the singular coordinate choice \(a=\gamma/2\).



The determinant of this two-by-two scalar transformation is \(\gamma-2a\). Thus the denominator in both inverse formulas has the same geometric origin. At the exceptional value both new coordinates are negatives of one another, and position and velocity can no longer be recovered separately.

\subsection{\texorpdfstring{Bound the position projection}{Bound the position projection}}


The triangle inequality and \((\|y\|-\|z\|)^2\ge0\) give \(\|y+z\|^2\le2(\|y\|^2+\|z\|^2)\). Substituting \(y+z=(\gamma-2a)x\) yields \((\gamma-2a)^2\|x\|^2\le2\mathcal E\). the triangle inequality supplies the triangle bound; elementary real inequalities combines it with nonnegative norms and the square of their difference. The final rewrite uses homogeneity of the norm and \(|c|^2=c^2\), so the squared result does not need a sign assumption on the denominator.



The factor two is the sharp constant for adding two Hilbert vectors: equality occurs when the two vectors coincide. Consequently the loss in position projection is a property of the chosen coordinates. It is independent of dimension and survives unchanged on a space of random variables.

\subsection{\texorpdfstring{Initialize a coupling with equal velocities}{Initialize a coupling with equal velocities}}


If the velocity difference is zero, the transformed pair is \(((\gamma-a)x,-ax)\). Its energy is therefore \(((\gamma-a)^2+a^2)\|x\|^2\). The proof unfolds both definitions, removes zero and negation, converts scalar multiplication norms to absolute scalar factors, and uses the identity \( |c|^2=c^2 \) before expansion of the scalar polynomials. This is the exact initialization coefficient in the synchronous-position contraction theorem.



For a coupling with a common initial velocity, the initial transformed cost therefore depends only on the position coupling. This is the bridge between a phase-space contraction estimate and an initial position Wasserstein distance. No initial velocity transport penalty is introduced by this choice.

\subsection{\texorpdfstring{Read off velocity control and phase energy}{Read off velocity control and phase energy}}


At \(a=0\), the second transformed coordinate is \(-v\), so \(\|v\|^2\le\|v+\gamma x\|^2+\|v\|^2=\mathcal E\). The argument simplifies the zero scalar terms and invokes nonnegativity of the other square. the ordinary phase energy separately records \(\|v\|^2+\|x\|^2\), the ordinary squared Euclidean phase norm. Both definitions avoid relying on the product's maximum norm; only the transformed energy contains the mixed position–velocity geometry needed for contraction.



Controlling both position and velocity follows by combining their projection inequalities. The transformed energy is therefore equivalent to the usual phase energy whenever the coordinate transformation is invertible. The constants of equivalence become large near the singular value, explaining the restriction on the shift.

\subsection{\texorpdfstring{Use the Euclidean transformed norm}{Use the Euclidean transformed norm}}


Define \(\|s\|_T\) as the norm of \(T_{\gamma,a}s\) in the Euclidean product Hilbert space. This the underlying measure theory type equips the two components with the Euclidean product norm. Its square is exactly the transformed energy, by the Euclidean squared-norm identity. Thus the norm, squared-energy calculations, and later triangle inequalities all describe the same geometry. Nonnegativity comes directly from nonnegativity of the norm. The construction applies to a Hilbert space of vectors or a Hilbert space of square-integrable random vectors.



Taking the square root is essential before combining two different errors. Squared energies do not obey the triangle inequality, whereas their square roots do. The Euclidean product construction supplies exactly the Hilbert norm needed for the subsequent insertion of an auxiliary exact trajectory.

\subsection{\texorpdfstring{Prove the triangle and velocity-\allowbreak{}projection bounds}{Prove the triangle and velocity-projection bounds}}


The transformation is additive: expand the two coordinates and let vector-space algebra normalize the vector identities. After this identity, linearity of the Euclidean identification turns the transformed sum into a sum in an normed space, so the triangle inequality gives its triangle inequality. Rewriting \(s-t=(s-r)+(r-t)\) gives the three-point bound needed to insert an auxiliary exact path. To pass from an energy bound \(\mathcal E(s)\le B^2\) to \(\|s\|_T\le B\), use nonnegativity and comparison of nonnegative squares, then replace the norm square by energy. Applying the same argument to the already proved velocity-square projection gives \(\|v\|\le\|(v,x)\|_T\) at \(a=0\).



In the coupling proof the intermediate state is the exact endpoint started from the numerical initial phase. The triangle inequality then separates propagation of the old discrepancy from the newly created local numerical error. Both terms remain in one common geometry until the final projection.

\subsection{\texorpdfstring{Project position only at the final endpoint}{Project position only at the final endpoint}}


The squared position bound is \(\gamma^2\|x\|^2\le2\|s\|_T^2\). Both sides are squares of nonnegative quantities, so comparison of nonnegative squares gives \(\gamma\|x\|\le\sqrt2\|s\|_T\). Division by positive friction yields the exact projection factor \(\sqrt2/\gamma\). If the velocity difference initially vanishes, the energy identity gives \(\|(0,x)\|_T^2=\gamma^2\|x\|^2\). Comparing squares with nonnegative \(\gamma\) proves \(\|(0,x)\|_T=\gamma\|x\|\). These two identities explain why the eventual position theorem loses only one factor \(\sqrt2\), while friction cancels out of its bias floor.



These identities must be applied in the correct order. Initialize the transformed discrepancy using the equal-velocity identity, propagate that norm through all iterations, and project position once at the endpoint. Projecting after every step would unnecessarily multiply the positional prefactor at every iteration.

\section{\texorpdfstring{Spectral cancellation and pathwise exponential decay}{Spectral cancellation and pathwise exponential decay}}
\label{proof-module-KineticLangevin-Contraction-lean}


This section proves a deterministic theorem for a kinetic difference equation \(v'=-\gamma v-H_tx\), \(x'=v\). The continuous linear operator \(H_t\) is allowed to vary with time but must be self-adjoint and satisfy quadratic-form bounds \(mI\preceq H_t\preceq MI\). The theorem derives energy decay from this differential equation. Constructing stochastic solutions whose difference has these properties and taking an infimum over probability couplings are separate layers.



The main calculation turns spectral information into a differential inequality with an exact coefficient. Keeping the shift explicit separates decay of the dynamics from conditioning of the coordinate transformation.

\subsection{\texorpdfstring{Write the derivative and prove exact spectral cancellation}{Write the derivative and prove exact spectral cancellation}}


The transformed derivatives are \(y'=-av-Hx\), \(z'=(\gamma-a)v+Hx\). Thus \(\mathcal E'=2\langle y,y'\rangle+2\langle z,z'\rangle\), which the energy derivative records. Multiplying by \(d=\gamma-2a\) gives

\[d\mathcal E'=2\bigl(a^2\|y\|^2-\langle y,Hy\rangle+\langle z,Hz\rangle-(\gamma-a)^2\|z\|^2\bigr).\]

The proof expands linear-map applications, scalar multiplication and inner products. Self-adjointness rewrites \(\langle v,Hx\rangle\) as \(\langle x,Hv\rangle\), so cross terms cancel. expansion of the scalar polynomials verifies the remaining scalar polynomial identity. Clearing the denominator keeps the identity valid even at singular parameter values.



To see why symmetry matters, expand \(\langle v+cx,Hx\rangle\) and its partner containing \(Hv\). The terms involving \(\langle v,Hx\rangle\) have equal and opposite coefficients after using \(\langle v,Hx\rangle=\langle Hv,x\rangle\). Without this identity the curvature bounds alone would not yield the displayed diagonal expression.

\subsection{\texorpdfstring{Use lower and upper curvature on different coordinates}{Use lower and upper curvature on different coordinates}}


Apply the lower bound \(\langle y,Hy\rangle\ge m\|y\|^2\) and upper bound \(\langle z,Hz\rangle\le M\|z\|^2\). Let \(k=\max\{a^2-m,M-(\gamma-a)^2\}\). Both coefficients are at most \(k\); multiply these scalar bounds by the nonnegative coordinate squares. The spectral identity then yields \(d\mathcal E'\le2k\mathcal E\). The explicit premise \(d>0\) permits division without reversing the inequality. Combining the quadratic-form and maximum inequalities and dividing by the positive denominator gives the result. This proves the dissipativity estimate from spectral bounds rather than assuming contraction.



The maximum is needed because the energy splits into two nonnegative coordinate contributions. Each contribution can be bounded by its own curvature coefficient, and a common coefficient must dominate both. One must divide only after checking \(\gamma-2a>0\); reversing that sign would reverse the dissipativity conclusion.

\subsection{\texorpdfstring{Specialize to the simple friction regime}{Specialize to the simple friction regime}}


At \(a=0\), the maximum becomes \(\max\{-m,M-\gamma^2\}\). The assumption \(m+M\le\gamma^2\) gives \(M-\gamma^2\le-m\), so the maximum equals \(-m\). With \(\gamma>0\), the differential bound becomes \(\mathcal E'\le-2m\mathcal E/\gamma\). The argument obtains the general bound, proves the scalar comparison by linear inequalities, and simplifies the maximum. This is the energy exponent; taking a square root later gives the norm exponent \(-m/\gamma\).



For \(m>0\) this yields strict exponential decay. Smoothness enters through the condition ensuring that the second transformed coordinate dissipates at least as fast as the first. The conclusion is uniform over all time-varying symmetric operators with the prescribed spectral interval.

\subsection{\texorpdfstring{Differentiate the transformed energy along a path}{Differentiate the transformed energy along a path}}


Assume ordinary differentiability witnesses the two kinetic equations at time \(t\). The linear combinations \(v+(\gamma-a)x\) and \(-v-ax\) inherit derivatives by addition, negation and constant scalar multiplication. The vector normalization calculation vector-space algebra reduces them to \(-av-Hx\) and \((\gamma-a)v+Hx\). The Hilbert-space squared-norm chain rule gives derivative \(2\langle y,y'\rangle\), and adding the two rules produces the energy derivative. No derivative of \(H_t\) is needed: only the value at the current time enters the first-order kinetic equation.



The calculation needs differentiability of the trajectories, not differentiability of the operator-valued path. This distinction is useful for averaged Hessians, which depend on both trajectories and may have limited temporal regularity. The kinetic equations already give every derivative used in the energy calculation.

\subsection{\texorpdfstring{Apply Gronwall with an explicitly proved derivative bound}{Apply Gronwall with an explicitly proved derivative bound}}


Define scalar functions \(F(t)=\mathcal E(v_t,x_t)\) and \(F'(t)=\mathcal D(H_t,v_t,x_t)\), with \(K=2k/(\gamma-2a)\). Continuous paths make \(F\) continuous on the closed horizon, and the preceding derivative theorem gives its derivative on the half-open interval. Spectral dissipativity proves \(F'\le KF+0\). The right-slope Grönwall theorem and its zero-forcing formula give \(F(t)\le F(0)e^{Kt}\). The zero additive term is significant: this is exact synchronous-flow contraction rather than a numerical defect estimate.



An integrating factor makes the argument transparent: the derivative of \(e^{-Kt}F(t)\) is at most zero. Its value at any later time is therefore bounded by its initial value. The proof uses a right-derivative version to accommodate the endpoint of the time interval.

\subsection{\texorpdfstring{Project the energy estimate to the position difference}{Project the energy estimate to the position difference}}


Suppose \(v_0=0\). Substitute the exact initial energy \(((\gamma-a)^2+a^2)\|x_0\|^2\) into the Gronwall estimate. The position projection theorem gives \((\gamma-2a)^2\|x_t\|^2\le2F(t)\). Combining the two inequalities yields the stated squared version of Proposition 1. The final elementary real inequalities only assembles already proved bounds. Its square root, under \(\gamma-2a>0\), is the paper's prefactor \(\sqrt{2((\gamma-a)^2+a^2)}/(\gamma-2a)\). This theorem concerns paths satisfying the explicit ODE; it is not yet a probability-law mixing result.



Every coefficient is retained when position is recovered. This is important because the change of variables can be poorly conditioned even when its energy decays rapidly. The final prefactor measures that conditioning, whereas the exponential describes contraction over time.

\subsection{\texorpdfstring{Include the initial time through right derivatives}{Include the initial time through right derivatives}}


Finite-horizon integral solutions naturally provide a derivative within \([t,\infty)\), including at \(t=0\). The coordinate differentiation proof works unchanged with right differentiability: add the two right derivatives, normalize their vector coefficients, apply the squared-norm derivative rule, and sum. The resulting derivative is still the energy derivative. The right-slope Grönwall theorem therefore gives the same exponential energy bound under the spectral hypotheses, without imposing a differentiable extension of the stochastic paths to negative time.



The stochastic forcing itself need not be differentiable. It has already disappeared from the difference of two common-input solutions. The resulting difference is an integral of a continuous drift, so its right derivative exists even though neither individual Brownian-driven velocity is differentiable.

\subsection{\texorpdfstring{Recover the full arbitrary-\allowbreak{}shift prefactor}{Recover the full arbitrary-shift prefactor}}


Write \(d=\gamma-2a>0\) and \(k=\max\{a^2-m,M-(\gamma-a)^2\}/d\). The energy estimate and equal initial velocities give \(d^2\|x_t\|^2\le2((\gamma-a)^2+a^2)e^{2kt}\|x_0\|^2\). One proves \((e^{kt})^2=e^{2kt}\) using the exponential addition formula, proves the square-root square identity from nonnegativity of its argument, and compares the squares of two nonnegative quantities. This first yields \(d\|x_t\|\le\sqrt{2((\gamma-a)^2+a^2)}e^{kt}\|x_0\|\). Dividing by the positive \(d\) gives precisely Proposition 1's prefactor and exponent.



If the maximum curvature coefficient is negative, the displayed exponent is a decay rate. If it is positive, theorem remains a valid finite-time stability estimate. Optimizing the shift is therefore a separate scalar problem whose admissible domain must preserve the positive denominator.

\section{\texorpdfstring{Nonlinear kinetic contraction directly from strong convexity and smoothness}{Nonlinear kinetic contraction directly from strong convexity and smoothness}}
\label{proof-module-KineticLangevin-PotentialContraction-lean}


The spectral proof can be replaced by a sector inequality for an gradient difference. This avoids selecting a time-varying averaged Hessian and only needs a differentiable strongly convex potential with Lipschitz gradient. The module proves the simple \(\sqrt2e^{-mt/\gamma}\) pathwise position estimate from the concrete noise-cancelled equations. Stochastic solution construction and Wasserstein law comparison are separate steps.



The two-point strong-smooth gradient inequality replaces an explicit Hessian representation. It yields the same transformed dissipativity under weaker regularity and applies directly to the potential assumptions of the sampler theorem.

\subsection{\texorpdfstring{Write the concrete energy derivative and the gradient sector inequality}{Write the concrete energy derivative and the gradient sector inequality}}


For velocity difference \(v\), position difference \(x\), and force difference \(g\), define \(D=2\langle v+\gamma x,-g\rangle+2\langle-v,\gamma v+g\rangle\). The strong-smooth interpolation inequality is \(mM\|x\|^2+\|g\|^2\le(m+M)\langle g,x\rangle\). Subtract the linear part \(mx\), expand its squared norm, and rearrange to obtain \(\|g-mx\|^2\le(M-m)\langle x,g-mx\rangle\). The proof performs this expansion with the Hilbert norm-square and inner-product identities, then uses elementary real inequalities. This sector estimate contains the joint lower–upper curvature information needed for contraction.



Write the displacement as \(q=x-y\) and the force difference as \(d=\nabla f(x)-\nabla f(y)\). Strong convexity and smoothness place \((q,d)\) in a sector narrower than a bare Lipschitz ball. This additional relation is what permits a dissipative transformed energy.

\subsection{\texorpdfstring{Complete a square to prove dissipativity}{Complete a square to prove dissipativity}}


Put \(k=g-mx\). Expanding the transformed energy yields

\[\gamma D+2m\mathcal E
=2\bigl[-(\gamma^2-2m)\|v\|^2-2\gamma\langle v,k\rangle-\gamma^2\langle x,k\rangle\bigr].\]

When \(M=m\), the sector inequality forces \(k=0\); friction gives \(\gamma^2\ge2m\), proving nonpositivity. When \(M>m\), combine the sector estimate with \(\|\gamma k+(M-m)v\|^2\ge0\) and \((M-m)(\gamma^2-m-M)\|v\|^2\ge0\). These imply \(\gamma D+2m\mathcal E\le0\). The argument handles the equal-curvature case separately, expands the square in the strict case, and combines the scalar inequalities. Positive \(\gamma\) then permits division, giving \(D\le-2m\mathcal E/\gamma\).



The potentially positive mixed term is paired with the negative position term and rewritten as a square plus a curvature contribution. The sector inequality controls the square. This preserves the strong-convexity constant rather than losing it through a separate triangle inequality on the gradient difference.

\subsection{\texorpdfstring{Obtain the sector condition from the potential}{Obtain the sector condition from the potential}}


For \(g=\nabla f(x)-\nabla f(y)\) and displacement \(x-y\), the shared ULA library has already proved the strong-smooth interpolation inequality from differentiability, strong convexity and gradient Lipschitzness. the strong-smooth force dissipativity inequality applies that theorem, converts it to the sector condition, and feeds it into the completed-square estimate. Consequently no Hessian, sector bound, or contraction conclusion is assumed in this potential theorem. The explicit hypotheses retain \(m\le M\), \(\gamma>0\), and \(\gamma^2\ge m+M\).



One may apply the strong-smooth interpolation inequality directly to \(f\). No continuous Hessian is required: differentiability, strong convexity, and a Lipschitz gradient already imply the needed two-point estimate. Thus this proof applies to the regularity assumptions of the first-order sampling theorem.

\subsection{\texorpdfstring{Differentiate the noise-\allowbreak{}cancelled path energy}{Differentiate the noise-cancelled path energy}}


Assume \(v'=-\gamma v-g\) and \(x'=v\). Then \((v+\gamma x)'=-g\), while \((-v)'=\gamma v+g\). The argument obtains these two derivative formulas by addition, negation and constant scalar multiplication, normalizes them with vector-space algebra, and applies the squared-norm derivative rule to each coordinate. Their sum is precisely the energy derivative. The theorem assumes differentiability of the difference paths; it does not incorrectly differentiate an individual Brownian velocity path.



For two paths driven by identical forcing, subtract their integral equations first. Their position difference has derivative equal to their velocity difference, and their velocity difference has derivative equal to the negative friction and force difference. The energy chain rule is then an ordinary deterministic calculation.

\subsection{\texorpdfstring{Integrate energy decay and take the position square root}{Integrate energy decay and take the position square root}}


For two continuous position paths \(x,y\), set \(F(t)=\mathcal E(v_t,x_t-y_t)\). The derivative theorem and potential dissipativity give \(F'\le-2mF/\gamma\). Closed-interval continuity and the explicit derivative witnesses allow the same Gronwall theorem as in the spectral proof, giving \(F(t)\le F(0)e^{-2mt/\gamma}\). If \(v_0=0\), then \(F(0)=\gamma^2\|x_0-y_0\|^2\). Combine with the position projection bound, cancel the positive \(\gamma^2\), and compare nonnegative squares. The exponential square identity \((e^{-mt/\gamma})^2=e^{-2mt/\gamma}\) yields \(\|x_t-y_t\|\le\sqrt2e^{-mt/\gamma}\|x_0-y_0\|\). This is a proved nonlinear pathwise estimate with the paper's correct prefactor.



After the differential inequality is established, integration gives exponential decay of the squared energy. Taking a square root halves its exponent. The equal-initial-velocity identity and the position projection then give the visible \(\sqrt2\) coefficient in the simple friction regime.

\subsection{\texorpdfstring{Use right derivatives at the initial time}{Use right derivatives at the initial time}}


A finite-horizon integral solution is constrained only for nonnegative time. Its extension to negative time need not have the same derivative. The correct initial derivative is therefore the derivative within the right half-line. The product and squared-norm derivative rules remain valid in this form, so differentiating the transformed energy gives the same force-energy expression. The argument applies those within-derivative rules to both transformed coordinates, cancels the opposing friction terms, and transfers the resulting derivative to the energy expression. The right-slope form of Gronwall requires exactly these hypotheses, and gives the same exponential energy decay throughout the closed interval. This is the version used by the integral-flow theorem.



At time zero all derivatives are interpreted from within the nonnegative time interval. This is sufficient for the integrating-factor argument on a closed finite horizon. It avoids extending either trajectory or its driving noise to negative times.

\section{\texorpdfstring{Arbitrary-\allowbreak{}friction contraction directly from the potential}{Arbitrary-friction contraction directly from the potential}}
\label{proof-module-KineticLangevin-GeneralContraction-lean}


The paper's full shifted contraction can be derived from differentiability, strong convexity and a Lipschitz gradient alone. This proof centers velocity at \(v+\gamma x/2\), uses the strong-smooth sector inequality, and completes a scalar square. It removes any need to assume that the Hessian is continuous when constructing a mean Hessian, while recovering the exact maximum of the two endpoint curvature exponents.



Centering the force around a scalar curvature isolates a residual controlled by the strong-smooth sector bound. Completing a square recovers the full shifted exponent without imposing continuous Hessians.

\subsection{\texorpdfstring{Center the phase energy and its derivative}{Center the phase energy and its derivative}}


The concrete force-difference derivative is \(2\langle v+(\gamma-a)x,-av-F\rangle+2\langle-v-ax,(\gamma-a)v+F\rangle\). Put \(y=v+\gamma x/2\), \(d=\gamma-2a\), and \(c=((\gamma-a)^2+a^2)/2\). Expanding the inner products gives the two exact identities

\[\mathcal E=2\|y\|^2+\tfrac12d^2\|x\|^2,\qquad \mathcal E'=-\gamma\mathcal E+4\langle y,cx-F\rangle.\]

The proof rewrites squared norms as self-inner-products, distributes inner products over sums, differences and scalar multiplication, uses symmetry, and finishes with expansion of the scalar polynomials. These identities expose a scalar residual rather than a matrix diagonalization.



Centering the curvature interval at a scalar value separates a uniformly controlled residual force from the linear central drift. The transformed energy can then be estimated by completing a square in this residual. This replaces the need to select an averaged Hessian explicitly.

\subsection{\texorpdfstring{Bound the residual around an arbitrary scalar curvature}{Bound the residual around an arbitrary scalar curvature}}


Strong-smooth interpolation places \(F\) in the ball centered at \(\mu x\), where \(\mu=(m+M)/2\), with radius \(R\|x\|\), where \(R=(M-m)/2\). Decompose \(cx-F=(c-\mu)x-(F-\mu x)\). The triangle inequality gives \(\|cx-F\|\le(|c-\mu|+R)\|x\|\). The coefficient is exactly \(\max\{c-m,M-c\}\): The proof splits on \(c\ge\mu\), simplifies the absolute value and the corresponding maximum, and checks the remaining arithmetic. No gradient or Hessian operator is invented in this argument.



The strong-smooth sector inequality bounds the residual by the half-width of the curvature interval. Its center and radius have different roles: the center determines the chosen linear comparison, whereas the radius controls the nonlinear deviation from it.

\subsection{\texorpdfstring{Complete the centered square to obtain the exact exponent}{Complete the centered square to obtain the exact exponent}}


Let \(K=\max\{c-m,M-c\}\ge0\). Cauchy–Schwarz and the residual bound give \(\langle y,cx-F\rangle\le K\|y\|\|x\|\). The nonnegative square \((2\|y\|-d\|x\|)^2\) implies \(4d\|y\|\|x\|\le2\mathcal E\). Multiplying by \(K\) and inserting the exact energy-derivative identity yields \(d\mathcal E'\le(-\gamma d+2K)\mathcal E\). Finally \(K=\max\{a^2-m,M-(\gamma-a)^2\}+\gamma d/2\), so the \(\gamma d\) terms cancel. The proof records the multiplied identities explicitly before elementary real inequalities, then divides by the positive \(d\). This proves the precise spectral differential estimate even in a friction regime with a positive growth exponent.



After completing the square, the two endpoint curvature contributions lead to a maximum. Dividing by the positive coordinate gap gives the same exponent as the symmetric-Hessian calculation. Thus both routes recover the paper’s shifted geometry with the same constants.

\subsection{\texorpdfstring{Supply the sector inequality from the potential}{Supply the sector inequality from the potential}}


For \(x-y\) and \(F=\nabla f(x)-\nabla f(y)\), the project's proved the strong-smooth gradient interpolation inequality gives the strong-smooth interpolation inequality from differentiability, strong convexity and global Lipschitz continuity. Its centered-ball consequence is exactly the residual theorem's premise. The proof composes these already proved results, yielding arbitrary-shift dissipativity for the nonlinear gradient. It assumes neither a Hessian representation nor the desired contraction conclusion.



The gradient of a differentiable strongly convex smooth potential obeys this sector inequality. Substituting it closes every deterministic premise of the general contraction estimate. No auxiliary vector field with an assumed sector bound is introduced in the final application.

\subsection{\texorpdfstring{Integrate the right-\allowbreak{}derivative inequality}{Integrate the right-derivative inequality}}


The two kinetic difference equations imply the transformed-coordinate right derivatives \(-av-F\) and \((\gamma-a)v+F\). Their squared-norm derivative rules add to the shifted energy derivative. For paths on \([0,T]\), define their transformed energy \(E(t)\). Continuity follows from continuity of the paths, while the preceding nonlinear result gives \(E'(t)\le2kE(t)\), with \(k=\max\{a^2-m,M-(\gamma-a)^2\}/(\gamma-2a)\). The right-slope form of Grönwall then proves \(E(t)\le E(0)e^{2kt}\), including the initial endpoint. The hypotheses are the force-difference differential equations, not an exponential bound supplied by the caller.



The difference of common-input integral solutions has the required right derivatives. Apply the scalar integrating-factor argument to its transformed energy and then use the inverse coordinate identities. This proves the general position estimate on every finite horizon.

\section{\texorpdfstring{Contraction of shared-\allowbreak{}input kinetic integral solutions}{Contraction of shared-input kinetic integral solutions}}
\label{proof-module-KineticLangevin-SynchronousFlow-lean}


This joins the Picard integral construction to the transformed-energy argument. Two solutions with the same continuous forcing have a noise-free difference equation as a consequence of their integral identities. The energy estimate and the paper's position decay then follow without assuming either a contraction result or a differentiable Brownian path.



The abstract energy inequalities are applied to solutions already constructed from their inputs. This closes the gap between a formal difference equation and a contraction statement for kinetic trajectories.

\subsection{\texorpdfstring{Derive the phase difference equation and contract its energy}{Derive the phase difference equation and contract its energy}}


Apply the shared-input derivative theorem to the concrete kinetic drift \(G(v,x)=(\gamma v+\nabla f(x),-v)\). Project its right derivative to velocity and position using the two continuous linear projection maps. The velocity derivative becomes \(-\gamma(V-V\prime)-[\nabla f(X)-\nabla f(X\prime)]\); the position derivative is \(V-V\prime\). Scalar multiplication and additive rearrangement justify both identities. The right-derivative energy theorem, which is valid at the initial endpoint, yields

\[\mathcal E_t\le e^{-2mt/\gamma}\mathcal E_0,\quad\mathcal E_t=\|V_t-V\prime_t+\gamma(X_t-X\prime_t)\|^2+\|V_t-V\prime_t\|^2.\]

In The proof the required Lipschitz constant for \(G\) is derived from the gradient bound; it is not an extra Hessian or solution assumption. The argument then supplies the original potential's strong convexity and smoothness hypotheses to the proved energy inequality.



The contraction theorem is now applied to integral solutions. Its differential hypotheses follow by cancellation and the fundamental theorem, so they are conclusions of the construction rather than extra assumptions about an arbitrary pair of random paths.

\subsection{\texorpdfstring{Equal initial velocities give the sharp position prefactor}{Equal initial velocities give the sharp position prefactor}}


If the initial velocities agree, \(\mathcal E_0=\gamma^2\|X_0-X\prime_0\|^2\). The projection estimate \(\gamma^2\|X_t-X\prime_t\|^2\le2\mathcal E_t\) cancels the factor \(\gamma^2\), giving the squared position bound. Positivity of \(\gamma\), the norm, the exponential and the square root permits taking square roots. The argument explicitly proves \(e^{-mt/\gamma\,\cdot2}=(e^{-mt/\gamma})^2\), rewrites \((\sqrt2)^2=2\), and applies the nonnegative square comparison theorem. The resulting prefactor is \(\sqrt2\), consistent with the explicit general prefactor in the supplied source.



The initial velocity difference is zero because the coupling shares one initial velocity. Thus its initial phase cost is exactly a constant times the position discrepancy. This special initialization is why the final theorem can be stated using position laws alone.

\subsection{\texorpdfstring{Contract the full transformed norm}{Contract the full transformed norm}}


The integral-flow energy theorem gives \(\mathcal E_t\le e^{-2mt/\gamma}\mathcal E_0\). Replacing energy by the square of the transformed norm and using \((e^{-mt/\gamma})^2=e^{-2mt/\gamma}\) gives \(\|Y_t-Z_t\|_T\le e^{-mt/\gamma}\|Y_0-Z_0\|_T\). One proves the exponential square identity by the exponential addition formula, notes that both norms and the exponential are nonnegative, and invokes the energy-to-norm lemma. This stronger phase estimate controls velocity differences and accumulates discretization errors without a separate projection loss at each step.



Unlike the equal-velocity projection, contraction in the full transformed norm permits any initial phase difference. This stronger form is needed at later numerical steps, where an accumulated velocity error is generally present as well as a position error.

\subsection{\texorpdfstring{Apply the full shifted theorem to integral flows}{Apply the full shifted theorem to integral flows}}


Take two kinetic integral solutions with the same additive input. Their noise-cancelled phase difference has a right derivative by the right derivative of the subtracted integral equations. Compose this derivative with the continuous linear velocity and position projections; vector-space algebra normalizes their drift formulas. The centered-sector theorem then supplies the full arbitrary-shift energy exponent from the original potential assumptions. Equal initial velocities identify the exact initial energy, and the position-square projection supplies the final prefactor. Squaring the desired nonnegative quantities, using the exponential-addition identity, and dividing by \(\gamma-2a>0\) proves

\[\|X_t-Y_t\|\le\frac{\sqrt{2((\gamma-a)^2+a^2)}}{\gamma-2a}\exp\!\left(\frac{\max\{a^2-m,M-(\gamma-a)^2\}}{\gamma-2a}t\right)\|X_0-Y_0\|.\]

No continuous-Hessian assumption or difference-ODE assumption remains; both the ODE and the coefficient are derived from the concrete integral paths.



The shift controls the relative weighting of the two phase components. The same integral-solution argument works for every admissible shift, so optimization changes only the scalar prefactor and exponent. The stochastic construction need not be repeated for each friction regime.

\section{\texorpdfstring{The nonlinear kinetic coupling in L\ensuremath{{}^{2}}}{The nonlinear kinetic coupling in L²}}
\label{proof-module-KineticLangevin-LpContraction-lean}


Random vectors form a Hilbert space under the L\ensuremath{{}^{2}} norm. The gradient acts on these vectors through the project's pointwise the square-integrable drift construction. This section lifts the centered strong-smooth inequality to that vector field, then proves contraction of kinetic integral paths in L\ensuremath{{}^{2}}. It does not introduce a hypothetical differentiable potential on the infinite-dimensional space of random variables.



Pointwise potential geometry lifts to the Hilbert space of square-integrable variables by integration. Consequently deterministic energy calculations become root-mean-square coupling bounds without changing the curvature constants.

\subsection{\texorpdfstring{Lift the centered gradient inequality pointwise}{Lift the centered gradient inequality pointwise}}


For deterministic points, strong convexity and Lipschitz gradient imply \(\|\nabla f(x)-\nabla f(y)-\frac{m+M}{2}(x-y)\|\le\frac{M-m}{2}\|x-y\|\). The lifted gradient agrees almost everywhere with pointwise application of the original gradient. The proof collects the almost-everywhere identities for subtraction, scalar multiplication and the pointwise representative of the lifted drift, rewrites the random-variable representatives, and applies the deterministic centered inequality at each sample. monotonicity of second-mean norm under a pointwise bound integrates this norm comparison, yielding the same constant in L\ensuremath{{}^{2}}. No independence or moment estimate beyond membership in L\ensuremath{{}^{2}} is needed for this step.



Apply the finite-dimensional sector inequality at almost every sample point, then integrate. All pairings are integrable by Cauchy–Schwarz. The integrated inequality is precisely the Hilbert-space sector bound for square-integrable random vectors, with the original curvature constants unchanged.

\subsection{\texorpdfstring{Derive the L\ensuremath{{}^{2}} energy inequality from integral paths}{Derive the L² energy inequality from integral paths}}


Two phase paths satisfy the concrete kinetic integral equations with the lifted gradient and common additive input. The noise cancels in those equations, so their difference has a right derivative. Continuous linear projections extract the velocity and position equations. The general centered dissipativity theorem applies in the Hilbert space L\ensuremath{{}^{2}} because the preceding pointwise lift supplies its sector premise. Under \(\gamma^2\ge m+M\), the maximum curvature coefficient equals \(-m\), giving \(E'(t)\le-2mE(t)/\gamma\). Continuity and the right-slope Grönwall theorem prove \(E(t)\le E(0)e^{-2mt/\gamma}\). Taking nonnegative square roots, using the exponential-addition identity, gives the full transformed L\ensuremath{{}^{2}} norm contraction.



The shared-input random paths are viewed as paths in a Hilbert space of square-integrable variables. Their integral equations give the same right derivatives and energy identity as in the deterministic argument. The resulting norm contraction is a root-mean-square coupling estimate.

\section{\texorpdfstring{Derive the frozen-\allowbreak{}force local defect}{Derive the frozen-force local defect}}
\label{proof-module-KineticLangevin-LocalError-lean}


The velocity defect is an explicit exponentially weighted integral of the changing force; the position defect is its primitive. Bounds follow from a Lipschitz force and a kinematic identity, not from a built-in defect hypothesis. The ambient norm may represent a deterministic Banach path or an L\ensuremath{{}^{2}} random-vector norm. A kinetic stationarity theorem supplying the required velocity-moment bound remains separate.



The key regularity gain is that position is a time integral of velocity. A linear-in-time force variation creates a quadratic velocity defect and a cubic position defect under a shared-noise comparison.

\subsection{\texorpdfstring{Define the force-\allowbreak{}freezing remainder as an integral}{Define the force-freezing remainder as an integral}}


Set \(D_v(t)=\int_0^te^{-\gamma(t-s)}(g(x_s)-g(x_0))ds\), and \(D_x(t)=\int_0^tD_v(s)ds\). These are exactly the terms left when one compares an exact kinetic equation with a frozen-force interpolation driven by the same additive noise. the velocity defect and the position defect define those vector-valued interval integrals. The definitions impose no desired error estimate. Interpreting them as a numerical-process difference additionally requires deriving the variation-of-constants identity for that process.



The comparison starts from the same phase and uses the same forcing. Every noise term therefore disappears from the error equation. What remains is the effect of replacing the exact moving force by its initial value, so the defect can be represented by an ordinary time integral.

\subsection{\texorpdfstring{Integrating velocity gives a linear position increment}{Integrating velocity gives a linear position increment}}


Assume \(x_t-x_0=\int_0^tv_sds\), with \(\|v_s\|\le C\) on the horizon. The norm-of-integral bound gives \(\|x_t-x_0\|\le C|t|=Ct\) for nonnegative time. The argument rewrites with the kinematic identity, applies the constant bound on the interval integral, and translates membership in the unoriented integration interval into \(0\le s\le t\le T\). This is the extra position regularity that kinetic Langevin gains from integrating velocity instead of receiving direct Brownian position noise.



If velocity is bounded by \(R\) on the step, then \(\|x_s-x_0\|\le Rs\). This uses only the kinematic relation \(x_s-x_0=\int_0^sv_rdr\). It explains why the force variation is linear in elapsed time even though velocity may be Brownian-driven.

\subsection{\texorpdfstring{Integrate a linear force increment to obtain the quadratic velocity defect}{Integrate a linear force increment to obtain the quadratic velocity defect}}


Lipschitzness gives \(\|g(x_s)-g(x_0)\|\le MCs\). With \(\gamma\ge0\) and \(s\le t\), the exponential kernel is between zero and one, so the norm of the weighted integrand is still at most \(MCs\). The proof derives this pointwise estimate from the preceding increment theorem, explicit exponential positivity and the norm of scalar multiplication. A continuous scalar majorant provides the required interval integrability. Integrating the bound gives \(\|D_v(t)\|\le\int_0^tMCsds=MCt^2/2\); The argument evaluates the integral and normalizes its coefficients with expansion of the scalar polynomials.



An \(M\)-Lipschitz force gives \(\|g(x_s)-g(x_0)\|\le MRs\). Dropping the exponential kernel only increases the integral when friction is nonnegative. Hence the velocity defect is bounded by \(MR\int_0^h s\,ds=MRh^2/2\).

\subsection{\texorpdfstring{Integrate once more to obtain the cubic position defect}{Integrate once more to obtain the cubic position defect}}


Apply the proved velocity-defect bound at each intermediate time \(s\): \(\|D_v(s)\|\le MCs^2/2\). Its continuous scalar majorant is integrable, so the norm of the position-defect integral is bounded by \(\int_0^tMCs^2/2\,ds=MCt^3/6\). First, the argument verifies that every intermediate time stays in the original horizon, rewrites the scalar coefficient into the majorant's form, applies norm-integral comparison, and uses the elementary integral of a power. The integration variable inside the bound is \(s\), preserving the correct cubic constant. No stationarity or Gaussian moment is assumed implicitly by this deterministic estimate.



For the position defect the first response is bounded by the remaining time: \(\psi_1(h-s)\le h-s\). Thus \(MR\int_0^h(h-s)s\,ds=MRh^3/6\). The extra integration is responsible for the one-higher power of the position error.

\subsection{\texorpdfstring{Recover the paper’s three-\allowbreak{}quarter transformed local constant}{Recover the paper’s three-quarter transformed local constant}}


Put \(B=MCh^2\). The component bounds give \(\|D_v\|\le B/2\) and \(\|D_x\|\le Bh/6\). If \(\gamma h\le1/4\), then \(\|D_v+\gamma D_x\|\le B/2+B/24=13B/24\). The transformed energy at parameter zero is \(\|D_v+\gamma D_x\|^2+\|D_v\|^2\), which is at most \(((13/24)^2+1/4)B^2\le(3B/4)^2\). The proof keeps all quantities nonnegative, squares each norm bound with comparison of nonnegative squares, unfolds the transformed energy and verifies the rational comparison by elementary real inequalities. This derives the constant for the explicit defect integrals; identifying them with an stationary-chain coupling still requires the kinetic process construction and moment estimates.



The transformed norm combines velocity and friction times position. Under the prescribed small-step restriction, the cubic positional contribution can be absorbed into the quadratic velocity budget. The resulting constant must be combined with exact-flow contraction before accumulating errors over many steps.

\section{\texorpdfstring{Identify the defined exponential-\allowbreak{}step error with its integral defect}{Identify the defined exponential-step error with its integral defect}}
\label{proof-module-KineticLangevin-MildError-lean}


The local errors are concrete integrals of changes in the force. This section proves that these integrals really are the errors of the defined exponential step when compared with a mild exact velocity and its integrated position. The variation-of-constants equation is an explicit input here; constructing its Brownian convolution and deriving it from the original integral equation remain distinct steps.



Identifying the exact mild response with the defined update is essential before estimating a numerical defect. The two paths share their complete filtered input, so their discrepancy is an ordinary nonlinear-force integral.

\subsection{\texorpdfstring{Separate the frozen force from its variation}{Separate the frozen force from its variation}}


Write \(D_v(t)=\int_0^t e^{-\gamma(t-s)}[g(x_s)-g(x_0)]ds\). Linearity and \(\int_0^t e^{-\gamma(t-s)}ds=\psi_1(t)\) give

\[D_v(t)=\int_0^t e^{-\gamma(t-s)}g(x_s)ds-\psi_1(t)g(x_0).\]

One proves interval integrability using continuity of the path, drift and exponential kernel; it then distributes scalar multiplication and the integral over subtraction. For a mild exact velocity \(V_t=\psi_0(t)V_0-\int_0^t e^{-\gamma(t-s)}g(x_s)ds+B_t\), subtracting it from the frozen velocity plus the same convolution \(B_t\) gives exactly \(D_v(t)\). The proof is algebra after the integral identity.



The decomposition \(g(x_s)=g(x_0)+(g(x_s)-g(x_0))\) separates the part integrated exactly by the method from the nonlinear remainder. The frozen term produces the response coefficients. The second term alone produces the local defect and can be bounded by gradient Lipschitzness.

\subsection{\texorpdfstring{Integrating the frozen velocity gives the position formula}{Integrating the frozen velocity gives the position formula}}


The frozen position is \(x+\psi_1(t)v-\psi_2(t)g\). Its derivative is the frozen velocity \(\psi_0(t)v-\psi_1(t)g\), and its value at zero is \(x\). The fundamental theorem of calculus gives \(\int_0^t V_{\rm frozen}(s)ds=X_{\rm frozen}(t)-x\). The argument invokes the already proved derivative theorem and proves integrability from continuity of the coefficients. The identity holds also at zero friction because the coefficients were defined by integrals.



The position equation is obtained by integrating the velocity formula. Interchanging the two deterministic integrals produces the first primitive response kernel. This explains why the position defect has an additional time factor and why the position noise is the primitive of the velocity response.

\subsection{\texorpdfstring{Both coordinates use the same noise convolution}{Both coordinates use the same noise convolution}}


The position innovation is \(\int_0^h B_sds\), the time integral of the same velocity convolution. If \(x_h-x_0=\int_0^hV_sds\), integrate the velocity error identity to obtain the position defect \(D_x(h)=\int_0^hD_v(s)ds\). Thus

\[\mathcal K_h(V_0,x_0;B_h,\int_0^hB_sds)-(V_h,x_h)=(D_v(h),D_x(h)).\]

One proves the velocity identity at every point in the interval, changes the integrand using that pointwise equality, distributes the integral over sums and differences, and uses the frozen-velocity integral and the kinematic identity. Extensionality then joins the two coordinate equalities into the phase-space equality. No stochastic error estimate is hidden inside the definitions.



The same convolution appears on the exact and numerical sides of the coupling. Thus the comparison involves no independent Gaussian variance penalty. Correctly retaining their common source is essential; merely matching separate marginal covariances would not justify pathwise cancellation.

\subsection{\texorpdfstring{Bound the defined step in the transformed norm}{Bound the defined step in the transformed norm}}


Apply the exact step-error identity to replace the numerical-minus-exact endpoint discrepancy by the explicit pair \((D_v,D_x)\). The local-error module has already integrated the kinematic velocity bound to prove \(\mathcal E(D_v,D_x)\le(3MC h^2/4)^2\) under \(\gamma h\le1/4\). The norm-square identity and nonnegativity allow comparison of the nonnegative transformed norm and its energy square to take square roots, giving \(\|\text{step}-\text{exact}\|_T\le3MC h^2/4\). Every input is visible: the mild velocity equation, continuous common convolution noise, the kinematic identity, and the uniform velocity norm bound. No stochastic moment or stationary law is asserted here.



Combine the explicit velocity and position remainder bounds before taking the transformed norm. The triangle inequality in that norm separates the two coordinates with their friction coefficient. The resulting one-step budget applies to the defined map because its frozen response was identified exactly.

\section{\texorpdfstring{Gaussian velocity moments and the stochastic local step}{Gaussian velocity moments and the stochastic local step}}
\label{proof-module-KineticLangevin-LpOneStep-lean}


This section proves the exact \(\sqrt p\) reference velocity norm from a standard Gaussian marginal, then combines the L\ensuremath{{}^{2}} nonlinear contraction with the exponential-integrator error. Its hypotheses display the Gaussian reference marginal and concrete exact integral paths; stationarity proving that marginal is a separate theorem.



The reference Gaussian marginal supplies the dimension-dependent velocity radius. Combining it with the auxiliary-flow comparison closes the concrete first-order one-step estimate on the coupling space.

\subsection{\texorpdfstring{Obtain the exact velocity moment from its Gaussian law}{Obtain the exact velocity moment from its Gaussian law}}


The standard Gaussian vector is the pushforward of independent scalar \(N(0,1)\) coordinates by the Euclidean identification. Compose that pushforward with coordinate evaluation; the finite-product marginal formula reduces its law to \(N(0,1)\), since every other coordinate measure has mass one. If an L\ensuremath{{}^{2}} velocity random variable has this standard Gaussian vector law, composition of laws gives the scalar coordinate laws. The project's proved Gaussian-coordinate moment identity then yields \(\|V\|_{L^2}=\sqrt{\#\iota}\). Thus the reference moment is derived from the explicit distribution rather than supplied as a numerical norm assumption.



A standard Gaussian vector in dimension \(p\) has squared Euclidean norm equal to a sum of \(p\) unit-variance coordinate squares. Hence its second moment is \(p\) and its root-mean-square radius is \(\sqrt p\). Stationarity supplies this same radius at every reference time.

\subsection{\texorpdfstring{Prove the concrete stochastic one-\allowbreak{}step inequality}{Prove the concrete stochastic one-step inequality}}


Let \(A\) be the auxiliary exact L\ensuremath{{}^{2}} path starting at the numerical phase point, and \(R\) a synchronously driven exact reference. Their transformed L\ensuremath{{}^{2}} distance is at most its initial value \(e\), because \(e^{-ms/\gamma}\le1\) for nonnegative time. Velocity projection and the exact Gaussian reference moment therefore give \(\|A_s^v\|_{L^2}\le\sqrt p+e\). Apply the mild-equation local-error theorem with this velocity bound: the numerical exponential update differs from \(A_h\) by at most \(3Mh^2(\sqrt p+e)/4\). Insert \(A_h\) with the transformed Hilbert triangle inequality and contract its distance to \(R_h\). The result is \(e_{\mathrm{next}}\le e^{-mh/\gamma}e+\frac34Mh^2(\sqrt p+e)\). Because the ambient space is L\ensuremath{{}^{2}}, the Bochner-integral and Hilbert triangle estimates are precisely the stochastic Minkowski estimates required by the paper.



The proof is factored through a shared theorem with an explicit reference velocity radius \(C\ge0\). A uniform bound \(\|V_t^{\rm ref}\|_{L^2}\le C\) supplies \(C+e_k\) in the local defect and the terminal bias \(Mh\sqrt2C/m\). The standard-Gaussian statement is a thin specialization with the radius \(\sqrt p\) proved from its marginal law. This same shared theorem applies to particle temperature \(1/N\), using the separately proved thermal radius \(\sqrt{(d_\theta+Nd_x)/N}\); there is no duplicate exponential-integrator argument.



At a given step, contraction controls the auxiliary-reference gap and the freezing estimate controls the numerical-auxiliary gap. The auxiliary velocity radius is at most \(\sqrt p\) plus the current transformed discrepancy. Combining the two inequalities gives the affine recurrence used in the global theorem.

\section{\texorpdfstring{Exact physical-\allowbreak{}temperature Gaussian velocity radius}{Exact physical-temperature Gaussian velocity radius}}
\label{proof-module-KIPLMC-Moments-lean}


In scaled particle coordinates every velocity coordinate has variance \(1/N\). The thermal target is explicitly the image of a standard Gaussian under multiplication by \(N^{-1/2}\). This section proves the reverse normalization as an law equality and computes the exact L\ensuremath{{}^{2}} radius \(\sqrt{(d_\theta+Nd_x)/N}\). No independence of coordinates is needed beyond what is already encoded by that Gaussian vector law; zero-dimensional spaces are included.



Gaussian scaling makes the equilibrium velocity radius intensive: the total coordinate dimension is divided by particle count. This exact radius determines the dimension dependence of the first particle method’s numerical bias.

\subsection{\texorpdfstring{Undo the thermal normalization as a pushforward equality}{Undo the thermal normalization as a pushforward equality}}


For \(N>0\), the scalar \(\sqrt N\) is nonzero. Composing the target's defining map \(v\mapsto N^{-1/2}v\) with \(v\mapsto\sqrt Nv\) gives the identity on the configuration space. The argument rewrites the two successive pushforwards with composition of pushforwards, proves the two scalar multiplications cancel using associativity of scalar multiplication and cancellation of a nonzero scalar, then applies the identity pushforward. Both maps are continuous and hence measurable. This is an exact distributional normalization, rather than a covariance-only assertion.



Multiplication by \(\sqrt N\) exactly reverses thermal velocity pushforward. The normalized velocity therefore has the standard Gaussian law on the full scaled coordinate space. This is an equality of measures and transfers all Gaussian moment identities.

\subsection{\texorpdfstring{Read off the L\ensuremath{{}^{2}} moment from the Gaussian marginal}{Read off the L² moment from the Gaussian marginal}}


Suppose the L\ensuremath{{}^{2}} random velocity \(V\) has thermal marginal. The preceding map equality and composition of distributional identities show that \(\sqrt NV\) has the standard Gaussian law. The almost-everywhere scalar-multiplication rule for Lp representatives transfers this statement to the Lp class. The shared Gaussian-moment theorem therefore gives \(\|\sqrt NV\|_{L^2}=\sqrt{d_\theta+Nd_x}\). Since \(\sqrt N>0\), the norm scaling rule and division yield

\[\|V\|_{L^2}=\sqrt{\frac{d_\theta+Nd_x}{N}}.\]

The argument simplifies the cardinality of the sum/\allowbreak{}product coordinate index, rewrites the square root of a quotient, and divides by the strictly positive \(\sqrt N\). The moment is exactly the constant needed by the shared kinetic local-error argument at temperature \(1/N\).



There are \(d_\theta+Nd_x\) independent scaled coordinates, each with variance \(1/N\). Their second moments sum to \(d_\theta/N+d_x\). Taking the square root gives the intensive thermal velocity radius used in the first-order particle bias.

\section{\texorpdfstring{The OBABO metric and a rigorous nonlinear average-\allowbreak{}force estimate}{The OBABO metric and a rigorous nonlinear average-force estimate}}
\label{proof-module-KineticLangevin-OBABOEnergy-lean}


The OBABO source metric is \(E(v,x)=\|x\|^2+2\langle x,v\rangle/\gamma+\|v\|^2/L\). Its continuous drift is contractive when \(\gamma^2\ge4L\). OBABO evaluates the force at two different position pairs; the proof below controls their average by an explicit centered-sector error. This avoids selecting a single Hessian point in a vector-valued mean-value theorem. The higher-order discrete remainder still has to be bounded separately.



The weighted metric exposes the dissipative part of the nonlinear two-kick update. The force average is derived from the endpoint evaluations, with its sector error retained explicitly.

\subsection{\texorpdfstring{Prove that the mixed energy defines the correct geometry}{Prove that the mixed energy defines the correct geometry}}


Write \(E_{a,b}(v,x)=\|x\|^2+2b\langle x,v\rangle+a\|v\|^2\), and let \(E_{a,0}\) be its diagonal part. Cauchy–Schwarz bounds the mixed inner product in both directions. The nonnegative square \((\|x\|-2b\|v\|)^2\), together with \(4b^2\le a\), gives \(|2b\langle x,v\rangle|\le E_{a,0}/2\). The proof records the two scalar inner-product comparisons, multiplies them by nonnegative \(2b\), and combines them with the squared difference using elementary real inequalities. Thus \(E_{a,0}/2\le E_{a,b}\le3E_{a,0}/2\), and energy is nonnegative. For \(a=1/L,b=1/\gamma\), the positive-denominator division rules show that the coefficient condition is exactly \(\gamma^2\ge4L\).



The mixed position–velocity energy is positive because it can be written as a sum of two squares after a linear change of variables. The inverse transformation yields quantitative coordinate bounds. Thus its square root is a true Hilbert norm in the admissible parameter regime.

\subsection{\texorpdfstring{Derive the source drift with an explicit sector error}{Derive the source drift with an explicit sector error}}


Along \(x'=v,v'=-\gamma v-F\), the energy derivative has common-denominator numerator \(-2L\langle x,F\rangle+(2L-2\gamma^2)\|v\|^2-2\gamma\langle v,F\rangle\). Suppose the force obeys \(\langle x,F\rangle\ge m\|x\|^2-e\|x\|\) and \(\|F\|^2\le L\langle x,F\rangle+Le\|x\|+e^2\). The nonnegative square \(\|\gamma v+F\|^2\) controls the velocity-force mixed term. Insert these inequalities, use the metric's upper equivalence, and use \(\gamma^2\ge4L\), \(0\le m\le L\) to discard the remaining nonpositive velocity term. Clearing the positive denominator \(\gamma L\) proves

\[E'\le-\frac{2m}{3\gamma}E+\frac{2e\|x\|}{\gamma}+\frac{e^2}{\gamma L}.\]

The error is explicit and vanishes when the average force satisfies the exact sector at the original position difference.



The ideal drift contribution is dissipative, but the two-kick method evaluates forces at different locations. The energy estimate isolates the resulting sector error rather than treating the whole force as arbitrary. This allows strong convexity to survive the numerical expansion.

\subsection{\texorpdfstring{Control the average of the two nonlinear force evaluations}{Control the average of the two nonlinear force evaluations}}


Let \(F,F'\) satisfy the strong-smooth centered bounds at differences \(x,x'\), and set \(S=(F+F')/2\). With \(\mu=(m+L)/2\) and \(R=(L-m)/2\), decompose \(S-\mu x\) into half of each centered force residual plus \(\mu(x'-x)/2\). The triangle inequality and \(\|x'\|\le\|x\|+\|x'-x\|\) yield \(\|S-\mu x\|\le R\|x\|+(L/2)\|x'-x\|\). Thus the sector error is precisely \(e=(L/2)\|x'-x\|\). A centered residual controls monotonicity by Cauchy–Schwarz. Moving its center from \(\mu x\) to \(Lx/2\) and squaring yields the approximate cocoercivity bound needed above. One proves both consequences from the explicit norm inequality.



The averaged force is formed from the endpoint evaluations of the nonlinear gradient. Its difference from a centered sector model is bounded using smoothness and the step displacement. This is necessary because an average of evaluations is not itself automatically a gradient at one intermediate point.

\subsection{\texorpdfstring{Apply the sector proof to the original potential}{Apply the sector proof to the original potential}}


For two arbitrary position pairs, use the project's proved strong-smooth interpolation and centered-ball theorem to obtain the two force residual bounds directly from differentiability, strong convexity and Lipschitz gradient. Their average is exactly the average of the two gradient differences used by the two OBABO kicks. The preceding average-sector lemma supplies the explicit position-increment error, and the perturbed dissipativity theorem supplies the source drift bound. No Hessian regularity, invented secant matrix, or assumed contraction inequality appears in this nonlinear step. The remaining discrete proof must bound the displayed error and the integrator remainder using its step-size restrictions.



The strong-smooth sector inequality holds for each gradient difference. Combining those inequalities and the displacement bound supplies the average-force premise. Consequently the contraction theorem applies to the nonlinear OBABO update, with no assumed averaged-Hessian representation.

\section{\texorpdfstring{OBABO: controlling every discrete remainder in the weighted metric}{OBABO: controlling every discrete remainder in the weighted metric}}
\label{proof-module-KineticLangevin-OBABORemainder-lean}


This section proves the numerical heart of the OBABO synchronous contraction. Starting from the explicit two-kick map, it bounds the averaged-force sector error, the Euler drift, and both components of the discrete remainder. These estimates combine to give the source factor \(1-mh/(3\gamma)\) under \(33\gamma^3h\le m\). The gradient and exponential damping are instantiated in the next module.



All positive terms in the discrete energy expansion are budgeted explicitly. The admissible step size guarantees that these remainders consume only part of the underlying curvature dissipation.

\subsection{\texorpdfstring{Realize the modified energy as an ordinary Hilbert norm}{Realize the modified energy as an ordinary Hilbert norm}}


Complete the square: \(E_{a,b}(v,x)=\|x+bv\|^2+(a-b^2)\|v\|^2\). The map \((v,x)\mapsto(x+bv,\sqrt{a-b^2}v)\) lands in the Euclidean product the Euclidean product geometry. This explicit choice matters: the ordinary The proof product norm is a maximum norm. the squared weighted phase-norm identity expands the two component norms, uses the square-root identity for nonnegative \(a-b^2\), and normalizes the scalar algebra. The associated bilinear pairing is exactly the Euclidean inner product of these transforms, so its upper bound follows from Cauchy–Schwarz without an extra geometric assumption.



A sum of squared transformed coordinates defines an inner product and hence a norm. In these coordinates Cauchy–Schwarz and the triangle inequality are available without estimating each mixed term separately. The inverse transform relates this norm to the physical phase variables.

\subsection{\texorpdfstring{Expand the exact energy of an Euler step plus a residual}{Expand the exact energy of an Euler step plus a residual}}


Expand \(E(s+t)=E(s)+2B(s,t)+E(t)\) and \(E(ct)=c^2E(t)\) directly from norms and inner products. Cauchy–Schwarz and \((\|s\|_E-\|t\|_E)^2\ge0\) give \(E(s+t)\le2E(s)+2E(t)\). For \(s_{+}=s+hA+r\), retain the first-order pairing \(2hB(s,A)\), bound the initial-residual pairing by \(2\|s\|_E\|r\|_E\), and apply the two-term estimate only to \(hA+r\). The resulting explicit bound is \(E(s_+)\le E(s)+2hB(s,A)+2\|s\|_E\|r\|_E+2h^2E(A)+2E(r)\). For the kinetic drift \(A=(-\gamma v-F,v)\), The argument checks that \(2B(s,A)\) is precisely the continuous derivative formula in the preceding module.



If an update is written as the old state plus a drift increment and a remainder, its energy expands into the old energy, twice a mixed pairing, and the squared increment. Dissipativity controls the leading pairing; all remaining terms must be included in the remainder budget.

\subsection{\texorpdfstring{Bound the drift and residual using their components}{Bound the drift and residual using their components}}


The averaged force obeys the explicit envelope \(\|F\|\le17L\|x\|/16+\gamma\|v\|/8\). Expanding the drift energy and using the nonnegative square \((\gamma\|x\|-2\|v\|)^2\) bounds its mixed term. The friction condition \(4L\le\gamma^2\) and metric equivalence then give \(E(A)\le13\gamma^2E(s)/4\). For the residual, assume only the directly verifiable component bounds \(\|r_x\|\le h^2(\gamma\|v\|+L\|x\|)/2\) and \(\|r_v\|\le h^2(9\gamma L\|x\|/16+5\gamma^2\|v\|/8)\). Squaring each envelope uses \((a-b)^2\ge0\); multiplying by positive coefficients preserves the inequalities. Metric equivalence and \(4L\le\gamma^2\) yield \(E(r)\le3h^4\gamma^4E(s)\). No norm error estimate is supplied as a convergence hypothesis.



The two force evaluations and thermostat damping determine the remainder. Smoothness bounds force changes by the position displacement. Coordinate bounds in the weighted norm then express the whole residual in terms of the original phase discrepancy.

\subsection{\texorpdfstring{Pay for the approximate centered sector}{Pay for the approximate centered sector}}


The change in position gives \(e\le hL(\|v\|/2+\gamma\|x\|/16)\). The continuous dissipativity therefore contributes the positive correction \(h(2e\|x\|/\gamma+e^2/(\gamma L))\). The proof bounds \(e^2\) with the square of the difference of its two summands, and bounds the position-velocity term using \((\gamma\|x\|-2\|v\|)^2\). The restrictions \(\gamma h\le1\), \(4L\le\gamma^2\), together with the lower metric equivalence, give the explicit budget \(\gamma^2h^2E(s)/2\). Denominators are cleared only with positive factors; the remaining scalar inequalities are by elementary real inequalities.



The averaged nonlinear force only approximately satisfies the desired centered sector relation. Its error contributes a positive term in the energy balance. Bounding that contribution explicitly prevents an unjustified use of the exact continuous-time sector identity for the discrete update.

\subsection{\texorpdfstring{Recover the source step-\allowbreak{}size constant from the complete budget}{Recover the source step-size constant from the complete budget}}


The residual energy bound implies \(\|r\|_E\le7h^2\gamma^2\|s\|_E/4\), since \(\sqrt3\le7/4\). Thus the initial-residual cross term costs at most \(7\gamma^2h^2E(s)/2\). The drift term costs \(13\gamma^2h^2E(s)/2\), the sector correction costs \(\gamma^2h^2E(s)/2\), and \(12\gamma^2h^2\le1\) makes twice the residual energy cost at most another \(\gamma^2h^2E(s)/2\). Their sum is exactly \(11\gamma^2h^2E(s)\). Consequently the map contracts by \(1-2mh/(3\gamma)+11\gamma^2h^2\), which is at most \(1-mh/(3\gamma)\) when \(33\gamma^3h\le m\).



The final algebra theorem constructs the residual explicitly: \(r_x=h(\delta-1)v-h^2F_0/2\), \(r_v=(\delta^2-1+\gamma h)v+h(1-\delta)F\). Norm triangle inequalities and the damping envelopes prove the two component bounds internally. The identity between \(s+hA+r\) and the difference update is by vector-space algebra, so no update-error contract is assumed.



The step restriction ensures that the sum of every positive remainder contribution is smaller than a fixed portion of the dissipative term. After that scalar comparison the weighted energy decreases by a factor below one. Taking square roots gives the one-step norm contraction.

\section{\texorpdfstring{The first-\allowbreak{}order scalar accumulation and its exact scope}{The first-order scalar accumulation and its exact scope}}
\label{proof-module-KineticLangevin-Recurrence-lean}


The numerical accumulation part of the KLMC proof begins with an explicit one-step coupling inequality. This section derives the stable contraction coefficient and sums its affine recurrence, retaining the coupling inequality and reference velocity norm as visible hypotheses. It therefore checks the numerical implication without claiming the still-open Gibbs stationarity and Gaussian-chain Wasserstein endpoint.



The scalar proof explains the transition from a local quadratic error to a global linear bias. Its hypotheses retain the coupling-dependent part of the local error, which must be absorbed into contraction.

\subsection{\texorpdfstring{Prove the one-\allowbreak{}sided exponential quadratic estimate}{Prove the one-sided exponential quadratic estimate}}


For \(x\ge0\), let \(F(x)=1-x+x^2/2-e^{-x}\). Its derivative is \(-1+x+e^{-x}\), nonnegative because \(e^{-x}\ge1-x\). Therefore \(F\) is nondecreasing and \(F(x)\ge F(0)=0\), proving \(e^{-x}\le1-x+x^2/2\). The argument constructs the derivative through the polynomial and exponential derivative rules, proves its sign with the inequality \(1+s\le e^s\), applies the derivative monotonicity criterion, and evaluates the zero endpoint. The estimate is proved analytically rather than introduced as an unchecked Taylor remainder.



The scalar estimate is applied only for a nonnegative step parameter. Its quadratic remainder quantifies how much the exact exponential contraction differs from its linear approximation. That remainder is later paid for using the small-step restriction, rather than silently discarded.

\subsection{\texorpdfstring{Keep most of the exact-\allowbreak{}flow contraction after the local correction}{Keep most of the exact-flow contraction after the local correction}}


Set \(r=m/\gamma\). Suppose \(r,M,h\ge0\), \(2r^2\le M\), and \(4Mh\le r\). The quadratic exponential estimate gives \(e^{-rh}+3Mh^2/4\le1-rh+r^2h^2/2+3Mh^2/4\). Since \(r^2/2\le M/4\), the positive correction is at most \(Mh^2\le rh/4\). Thus the total is at most \(1-3rh/4\). The proof multiplies the curvature and step inequalities by the nonnegative time powers and combines them with elementary real inequalities. The next theorem derives both scalar conditions from \(m>0\), \(m\le M\), \(\gamma^2\ge m+M\) and \(4\gamma Mh\le m\), dividing only by positive \(\gamma\) or \(\gamma^2\).



The local error contains a term proportional to the old coupling discrepancy as well as a stationary forcing term. The step restriction keeps the proportional correction small enough to retain a fixed fraction of the continuous contraction. Only the remaining stationary term contributes to the persistent bias.

\subsection{\texorpdfstring{Sum the affine recurrence with the visible local-\allowbreak{}coupling premise}{Sum the affine recurrence with the visible local-coupling premise}}


Assume nonnegative errors satisfy \(e_{n+1}\le e^{-rh}e_n+3Mh^2(C+e_n)/4\), with \(r,h>0\), \(C\ge0\), and the scalar restrictions above. Put \(q=1-3rh/4\). The coefficient bound and positivity of the exponential prove \(0\le q<1\), and multiplying that bound by \(e_n\ge0\) yields \(e_{n+1}\le qe_n+3Mh^2C/4\). The shared affine-recurrence theorem gives \(e_n\le q^ne_0+(3Mh^2C/4)/(1-q)\). Positive \(r,h\) permit denominator cancellation, leaving \(e_n\le q^ne_0+MhC/r\). The argument uses cancellation of nonzero denominators and expansion of the scalar polynomials only for that final quotient. The one-step coupling inequality is an explicit assumption, so this result must not be presented as the full actual-law KLMC theorem.



If \(A_{n+1}\le qA_n+b\) and \(0\le q<1\), induction gives \(A_n\le q^nA_0+b\sum_{j=0}^{n-1}q^j\). The sum is at most \((1-q)^{-1}\). A local defect of order \(h^2\) therefore produces a global floor of order \(h\).

\subsection{\texorpdfstring{Restore the paper’s physical parameters and position coefficient}{Restore the paper’s physical parameters and position coefficient}}


Substitute \(r=m/\gamma\) into the scalar recurrence theorem. The derived scalar conditions justify the call, and rational simplification gives \(e_n\le(1-3mh/(4\gamma))^ne_0+\gamma MhC/m\). Suppose in addition that \(d_n\le\sqrt2e_n/\gamma\) and \(e_0\le\gamma d_0\). The nonnegative projection factor and nonnegative power of the proved stable contraction coefficient allow multiplication of these inequalities. Canceling \(\gamma\) yields

\[d_n\le\sqrt2(1-3mh/(4\gamma))^nd_0+Mh\sqrt2C/m.\]

The argument explicitly proves the projection factor and contraction coefficient nonnegative before transporting inequalities, then uses cancellation of nonzero denominators for the scalar cancellations. Setting \(C=\sqrt p\) matches the paper’s numerical expression. The local coupling, projection and initialization bounds remain explicit hypotheses on sequences; neither theorem identifies \(d_n\) with Wasserstein error or constructs a stationary reference.



The physical substitution identifies the dimension-dependent stationary velocity radius and the curvature-dependent contraction gap. The transformed initial error contributes a factor \(\gamma\), while endpoint position projection contributes \(1/\gamma\). Their cancellation leaves the stated first-order position bias.

\section{\texorpdfstring{The nonlinear OBABO map contracts with the source constants}{The nonlinear OBABO map contracts with the source constants}}
\label{proof-module-KineticLangevin-OBABOContraction-lean}


For a differentiable, m-strongly convex potential with M-Lipschitz gradient, the OBABO step driven by common thermostat vectors contracts the modified energy whenever \(\gamma^2\ge4M\) and \(33\gamma^3h\le m\). The thermostat amplitude is arbitrary: it disappears through synchronous cancellation. This proves a property of the defined map and its iterates, rather than postulating a contractive transition.



Synchronous thermostat noises cancel from the discrepancy, but the two nonlinear force evaluations remain. Their full weighted remainder estimate proves contraction for the defined map and all its iterates.

\subsection{\texorpdfstring{Derive the small-\allowbreak{}step and exponential damping bounds}{Derive the small-step and exponential damping bounds}}


The source restrictions imply \(\gamma h\le1/132\): multiply by the positive \(\gamma^2\) and use \(33\gamma^3h\le m\le M\le\gamma^2/4\). Squaring gives the weaker estimates \(\gamma h\le1\), \(12\gamma^2h^2\le1\) used by the residual budget. With \(\delta=e^{-\gamma h/2}\), positivity and monotonicity of the exponential give \(0\le\delta\le1\). The global tangent inequality \(e^{-u}\ge1-u\) gives \(1-\delta\le\gamma h/2\) and \(\delta^2-1+\gamma h\ge0\). The proved quadratic exponential bound gives its upper bound \(\gamma^2h^2/2\). These are scalar facts about the exact damping used in the code.



The exponential damping coefficient lies between zero and one for nonnegative friction and step. Elementary exponential inequalities convert the source restriction into usable bounds on its distance from one. These estimates make the subsequent scalar remainder budget explicit.

\subsection{\texorpdfstring{Apply the complete remainder estimate to the two-\allowbreak{}kick step}{Apply the complete remainder estimate to the two-kick step}}


For two states driven by the same noises, set \(x=x_0-y_0\), \(v=v_0-w_0\), \(F_0=\nabla f(x_0)-\nabla f(y_0)\), and let \(x_1,y_1\) be the drift-stage positions. The average second-kick force is \(S=(F_0+\nabla f(x_1)-\nabla f(y_1))/2\). Its sector error is exactly \(e=M\|(x_1-y_1)-x\|/2\). Lipschitzness bounds \(F_0\), while the position update bounds the change in difference by \(h\|v\|+h^2M\|x\|/2\). The centered interpolation inequality for the original potential at both pairs yields the average-force sector estimate and then \(\|S\|\le M\|x\|+e\).



The source-metric drift estimate and all damping bounds now satisfy the preceding module’s explicit algebra theorem. Finally vector-space algebra checks that the exact difference of the two defined O-B-A-B-O steps equals that algebraic map, including both noisy intermediate positions. The result is \(E(z_+-w_+)\le(1-mh/(3\gamma))E(z-w)\). No Gaussian assumption is needed for this pathwise theorem, because the same additive vectors cancel exactly.



Insert the two force kicks and the displaced position into the weighted-energy expansion. The remainder estimate accounts for both nonlinear force variation and damping. The final inequality therefore bounds the defined OBABO map under synchronous thermostat noises.

\subsection{\texorpdfstring{Iterate and project to the physical position}{Iterate and project to the physical position}}


The small-step restrictions ensure the factor \(q=1-mh/(3\gamma)\) is nonnegative. Induction then applies the one-step theorem and multiplies the preceding energy inequality by q, giving \(E(Z_n-W_n)\le q^nE(Z_0-W_0)\) for any two recursions with the same two input vectors at every step. When the initial velocities agree, the initial mixed energy is exactly the squared initial position distance. The lower metric equivalence gives \(\|X_n-Y_n\|^2\le2q^n\|X_0-Y_0\|^2\). Both sides are nonnegative, so The argument applies the square comparison and the nonnegative square-root identity to obtain \(\|X_n-Y_n\|\le\sqrt{2q^n}\|X_0-Y_0\|\). This is a direct statement about numerical iterates; identifying an invariant law and quantifying discretization bias are separate analytic steps.



Apply the same one-step inequality to the discrepancy at each iteration, using common noises at both thermostat stages. Induction gives the geometric phase decay. The position projection is then applied at the final step, keeping its geometric constant outside the iteration.

\section{\texorpdfstring{The kinetic solution depends continuously on its initial state and noise path}{The kinetic solution depends continuously on its initial state and noise path}}
\label{proof-module-KineticLangevin-Flow-lean}


Picard existence alone does not ensure a selected solution can be used as a measurable random process. This section proves stability when both the initial condition and continuous forcing vary, then constructs the kinetic flow by selection from the existing existence theorem. The stability estimate proves that the selected endpoint, and indeed the whole compact solution path, is jointly Lipschitz. No measurability, output-law identity or contraction conclusion is postulated.



Continuity of the solution operator turns pathwise existence into a usable probabilistic construction. It ensures that random initial states and random continuous inputs have measurable endpoint outputs.

\subsection{\texorpdfstring{Subtract the rough forcing before differentiating}{Subtract the rough forcing before differentiating}}


Let \(Y_t-x=-\int_0^tg(Y_s)ds+B_t\) and \(Z_t-y=-\int_0^tg(Z_s)ds+C_t\). The inputs \(B,C\) are continuous and may be nowhere differentiable. Set \(F_t=Y_t-Z_t-(B_t-C_t)\). Subtract the integral equations to obtain \(F_t=x-y-\int_0^t(g(Y_s)-g(Z_s))ds\). This expression is differentiable from the right, including at the initial endpoint. If \(\|B_t-C_t\|\le D\), then \(\|Y_t-Z_t\|\le\|F_t\|+D\), and an \(L\)-Lipschitz drift gives \(\|F'_t\|\le(L+1)\|F_t\|+(L+1)D\). Gronwall and the initial equality \(F_0=x-y\) yield \(\|Y_t-Z_t\|\le(\|x-y\|+D)e^{(L+1)t}\). One proves equality with the differentiable integral expression on a right neighborhood and uses the within-derivative Gronwall theorem. The harmless extra one makes the coefficient strictly positive and avoids a special zero-Lipschitz case.



Subtracting the additive input replaces a rough forced path by a differentiable integral path. The resulting equation still contains the input through its state argument, but its time derivative is now an ordinary drift evaluation. This is the appropriate route to pathwise stability.

\subsection{\texorpdfstring{Use compact paths with a uniform norm}{Use compact paths with a uniform norm}}


the space of continuous forcing paths is the Banach space of continuous velocity-forcing paths on \([0,T]\), equipped with the supremum norm. The extension \(\widehat B_t=B(\operatorname{clip}_{[0,T]}t)-B(0)\) is continuous on the whole real line and starts at zero. This permits direct use of the proved finite-horizon Picard construction. Two normalized extensions satisfy \(\|\widehat B_t-\widehat C_t\|\le2\|B-C\|_\infty\): each of the two evaluation differences is bounded by the uniform norm. The argument uses continuity of the interval projection, evaluates at zero, and applies the norm triangle inequality to the two differences. No derivative of a forcing path is taken.



Continuity on a compact interval implies a finite uniform bound. Consequently the drift is bounded along each path and its integral is well defined. The supremum norm also controls every endpoint evaluation, linking path estimates to continuity of the terminal map.

\subsection{\texorpdfstring{Select an solution from proved existence}{Select an solution from proved existence}}


the kinetic solution map uses classical choice on the kinetic integral-solution existence theorem, supplying the normalized forcing \((\widehat B_t,0)\). Its specification theorem retains all conclusions of that construction: the path is globally continuous, its initial phase is exactly the supplied phase, and it satisfies the concrete kinetic integral equation at every time on the horizon. Classical selection is justified because existence was proved analytically. It does not add a solution axiom. Subsequent stability between any two solutions makes every selected representative agree on the horizon and provides continuous parameter dependence.



uniqueness of the selected kinetic solution applies the proved Lipschitz integral-equation uniqueness theorem to identify the selection with any other continuous solution driven by the same normalized forcing and initial phase. Thus a concrete Brownian realization can be identified with this canonical flow without assuming equality of endpoint laws.



Existence and uniqueness together turn the integral equation into a well-defined solution operator. Any choice made while constructing a solution has the same output. This matters when the operator is later pushed forward under a probability law on initial states and continuous inputs.

\subsection{\texorpdfstring{Prove joint endpoint continuity quantitatively}{Prove joint endpoint continuity quantitatively}}


The kinetic negative drift has Lipschitz constant \(K_0=\gamma+uL+1\) in the product maximum norm. Apply the two-input estimate with forcing difference \(D=2\|B-C\|_\infty\). This gives

\[\|Y_t^{z,B}-Y_t^{w,C}\|\le(\|z-w\|+2\|B-C\|_\infty)e^{(K_0+1)t}.\]

In the product parameter norm both input differences are bounded by the same distance, so the endpoint is Lipschitz with constant \(3e^{(K_0+1)t}\). The argument derives the estimate from the selected paths and their specifications, then bounds each coordinate by the product norm. continuous dependence of the kinetic endpoint follows from the Lipschitz theorem. Thus the eventual random endpoint map is measurable for a proved reason, even though its original existence definition used classical choice.



For two initial phases and inputs, the difference is bounded by the initial discrepancy plus the uniform input discrepancy plus a Lipschitz drift integral. Grönwall gives a bound proportional to \(e^{LT}\). Hence small perturbations of both arguments produce small changes at every fixed endpoint.

\subsection{\texorpdfstring{Control the entire solution as a compact continuous path}{Control the entire solution as a compact continuous path}}


the continuous kinetic path restricts the selected solution to \([0,T]\), retaining its proved continuity. At each time its difference has the endpoint bound. Since \(K_0+1\ge0\), \(e^{(K_0+1)t}\le e^{(K_0+1)T}\) on the horizon. The supremum-norm characterization then bounds the whole path difference by \(3e^{(K_0+1)T}\) times the initial-state/\allowbreak{}forcing-path distance. The argument applies the supremum norm bound, proves the pointwise exponential comparison, and uses the endpoint Lipschitz estimate. The resulting path-valued Lipschitz and continuity theorems support pushforwards into a space of continuous trajectories, not merely separate endpoint choices.



The same estimate holds uniformly in terminal time on the compact horizon. Thus the solution map is continuous as a map into the continuous-path space, not merely pointwise in time. In particular, composing it with a random input gives a measurable path-valued object.

\section{\texorpdfstring{Actual independent-\allowbreak{}source kinetic laws and their Wasserstein infimum}{Actual independent-source kinetic laws and their Wasserstein infimum}}
\label{proof-module-KineticLangevin-FlowLaw-lean}


The position law is an explicit pushforward of an independent product of the initial position law and a fixed source law for common initial velocity and continuous forcing. Every initial coupling is extended by that same source. The proof computes the two output marginals, preserves the initial cost, derives a coupling estimate and takes the infimum. The kinetic pointwise contraction is separately proved from the integral equation and the potential hypotheses; it is not an assumed convergence field.



The proof works over every initial transport plan extended by an independent common source. Exact output marginals and cost control allow the plan cost to be replaced by the true Wasserstein infimum.

\subsection{\texorpdfstring{Define the independent random-\allowbreak{}map law and compute both output marginals}{Define the independent random-map law and compute both output marginals}}


For a measurable endpoint \(F(x,s)\), set \(\mathcal T_Q\mu=F_\#(\mu\otimes Q)\). Given any coupling \(c\) of \(\mu,\nu\), extend it to \(c\otimes Q\). Projecting to \((y,s)\) produces \(\nu\otimes Q\), because the second marginal of \(c\) is \(\nu\). Pushing further by \(F\) therefore gives \(\mathcal T_Q\nu\). One proves this with the product-map identity, the coupling's marginal equality, and pushforward composition. The corresponding first marginal uses the existing the marginal identity for the independent coupling extension. Both output-law identities are derived from the concrete product construction and are independent of which coupling was chosen.



Given an initial position coupling, add one independent common source containing velocity and forcing. Evaluating the same endpoint map at the two coupled positions gives a coupling of the two desired output laws. Product integration verifies each marginal without assuming the initial coupling is optimal.

\subsection{\texorpdfstring{Lift pathwise contraction through L\ensuremath{{}^{2}} and take all initial couplings}{Lift pathwise contraction through L² and take all initial couplings}}


Suppose the endpoint obeys \(\|F(x,s)-F(y,s)\|\le q\|x-y\|\) for each common source. On \(c\otimes Q\), the two outputs form a measurable coupling of \(\mathcal T_Q\mu,\mathcal T_Q\nu\). Pointwise contraction passes to the extended L\ensuremath{{}^{2}} norm. The initial difference depends only on \(c\), so its norm on the product is exactly the original coupling cost. Thus \(W_2(\mathcal T_Q\mu,\mathcal T_Q\nu)\le q\,\operatorname{cost}(c)\) for every \(c\). The argument then applies the infimum theorem, using the independent product \(\mu\otimes\nu\) as a witness that couplings exist. No optimal-coupling attainment and no supplied output-law independence are required.



Square the pathwise Lipschitz inequality and integrate over the extended coupling. Taking square roots bounds the output transport cost by the same coefficient times the initial cost. Since the construction works for every initial coupling, taking the infimum gives the Wasserstein inequality.

\subsection{\texorpdfstring{Retain probability and justify the real-\allowbreak{}valued distance}{Retain probability and justify the real-valued distance}}


A measurable pushforward of \(\mu\otimes Q\) has total mass one when both inputs are probability measures; the probability property of the random-map pushforward records this. The primary distance bound lives in the extended nonnegative reals, so infinite distances are preserved. If the initial Wasserstein distance is finite, the finite coefficient \(q\) makes its upper bound finite, and monotonicity of the finite real value gives the real-valued estimate. One proves that finiteness before converting the inequality, then rewrites the real value of the product. This avoids accidentally converting an infinite output distance to zero.



A measurable map sends a probability measure to a probability measure. Finiteness of the second-moment transport cost is separately before interpreting the extended transport distance as a real number. This prevents real-valued inequalities from concealing an infinite-distance case.

\subsection{\texorpdfstring{Make the concrete kinetic endpoint a measurable random map}{Make the concrete kinetic endpoint a measurable random map}}


The compact forcing space receives its Borel measurable structure. The source \(s=(v,B)\) contains initial velocity and a continuous forcing path; the separate input is initial position \(x\). the position endpoint map evaluates the selected flow at phase \((v,x)\) and takes its position coordinate. The coordinate rearrangement \((x,(v,B))\mapsto((v,x),B)\) is continuous. Composing it with the proved joint endpoint continuity and projecting position gives continuity, hence measurability, of the map. Independence of source and position is defined later by the explicit product measure rather than encoded as an unspecified output-law premise.



Continuity of the solution with respect to initialization and forcing implies measurability of its endpoint. Thus the law being compared is the pushforward of an explicitly defined random map. Its probability property follows from the input probability measure, not from a formal density calculation.

\subsection{\texorpdfstring{Derive pointwise contraction from integral solutions}{Derive pointwise contraction from integral solutions}}


Take two values of the selected kinetic flow with phases \((v,x)\) and \((v,y)\) and the same forcing path. Their specification gives the integral equations and equal initial velocities. The proved integral-synchronous theorem cancels their common forcing, derives the right-derivative difference equation, and applies the nonlinear potential energy estimate. Under differentiability, \(m\)-strong convexity, \(M\)-Lipschitz gradient and \(\gamma^2\ge m+M\), this yields \(\|F_t(x,s)-F_t(y,s)\|\le\sqrt2e^{-mt/\gamma}\|x-y\|\). The argument obtains the integral premises by rewriting the flow's initial-value specification, then simplifies the initial position and velocity projections. No differential equation or contraction conclusion is assumed of the chosen flow.



Both endpoints are evaluated with a common initial velocity and common additive path. Their difference satisfies the already proved deterministic kinetic equation. This ties the pointwise Lipschitz coefficient to curvature of the potential, rather than postulating an abstract contractive transition.

\subsection{\texorpdfstring{Conclude the law-\allowbreak{}to-\allowbreak{}law kinetic theorem}{Conclude the law-to-law kinetic theorem}}


the kinetic position law is the explicit independent-product pushforward for this concrete endpoint. Its probability theorem follows from endpoint measurability. Substituting the derived coefficient \(q=\sqrt2e^{-mt/\gamma}\) in the general random-map transport theorem gives

\[W_2(\mu_t,\nu_t)\le\sqrt2e^{-mt/\gamma}W_2(\mu_0,\nu_0).\]

The source law may be any fixed probability law of initial velocity and continuous forcing, independent of initial position. The statement therefore applies to any such stochastic realization; a stationary Gibbs law is unnecessary for comparing two evolving laws. The extended-valued theorem retains infinite initial costs, and the real-valued theorem explicitly requires finite initial Wasserstein distance. The full all-friction spectral prefactor and a canonical Brownian-source identification are separate refinements beyond this simple-friction law theorem.



The conclusion concerns the independently defined output laws for each initial distribution. Dependence on the auxiliary initial coupling has disappeared through the marginal identities. The coupling argument therefore proves a property of the kinetic transition family.

\subsection{\texorpdfstring{The full shifted prefactor and exponent pass to output laws}{The full shifted prefactor and exponent pass to output laws}}


The actual-integral general contraction theorem uses the centered gradient sector estimate, so it requires differentiability, strong convexity and a Lipschitz gradient without adding Hessian continuity. For \(d=\gamma-2a>0\), its coefficient is

\[q_a(t)=\frac{\sqrt{2((\gamma-a)^2+a^2)}}{\gamma-2a}
\exp\!\left(\frac{\max\{a^2-m,M-(\gamma-a)^2\}}{\gamma-2a}t\right).\]

The proof instantiates that actual-integral theorem with the same two selected flows and their proved specifications. The coefficient is nonnegative, so it can be bundled as an \(\mathbb R_{\ge0}\) scalar in the already proved random-map transport theorem. This yields \(W_2(\mu_t,\nu_t)\le q_a(t)W_2(\mu_0,\nu_0)\), with the coupling infimum and internally derived output marginals. The real-valued version uses the same finite-initial-cost conversion. The theorem holds for every independent continuous-forcing source law; tying a constructed Brownian source and its required independence to the paper’s named diffusion remains separate, and no Gibbs invariant law is assumed.



The positivity of the transformed-coordinate denominator makes the full prefactor nonnegative. That coefficient can therefore be pulled through the infimum over coupling costs. This yields the entire shifted bound, including its conditioning factor, rather than only the simple zero-shift specialization.

\section{\texorpdfstring{Actual independent-\allowbreak{}noise OBABO laws contract in Wasserstein distance}{Actual independent-noise OBABO laws contract in Wasserstein distance}}
\label{proof-module-KineticLangevin-OBABOLaw-lean}


The preceding pathwise energy contraction now becomes a theorem about the Markov laws. An explicit invertible change of coordinates turns the modified energy into Euclidean squared distance. Measurability, the commuting law identity, independent source extension, and passage to the Wasserstein infimum are all proved. The result holds for every probability law of the two thermostat inputs, including the independent Gaussian law used by the sampler.



Changing to contracting coordinates transports the complete transition measure as well as its deterministic update. This makes the pathwise OBABO contraction a theorem about the iterated probability laws.

\subsection{\texorpdfstring{Build and verify the weighted-\allowbreak{}coordinate inverse}{Build and verify the weighted-coordinate inverse}}


The weighted Hilbert coordinates are \(T(v,x)=(x+bv,c v)\), where \(c=\sqrt{a-b^2}\). For \(b^2<a\), the inverse is \(T^{-1}(u,w)=(w/c,u-bw/c)\). First, the argument proves \(c\ne0\), then uses the inverse cancellation identities for real scalars to check both compositions. The transform commutes with subtraction by direct vector-space algebra. Its continuity follows from continuity of scalar multiplication, coordinate projections, addition, and the Euclidean-product identification. The local measurable structure on the Euclidean product geometry is explicitly its Borel structure, so these continuous maps supply the measurable changes of variables required for the measure arguments.



The weighted coordinate map is linear and has an explicitly computed inverse. Applying both compositions verifies the identification of physical and contracting coordinates. This ensures that transformed laws can be returned to their original physical meaning.

\subsection{\texorpdfstring{Conjugate the noisy step and prove its Lipschitz bound}{Conjugate the noisy step and prove its Lipschitz bound}}


Define the transformed update by the conjugation \(\widetilde P_\xi=T\circ P_\xi\circ T^{-1}\), retaining the original two noise arguments and their damping and amplitude. Continuity is a composition of the verified inverse, the O-B-A-B-O formulas, and the verified transform. The condition \(\gamma^2\ge4M>0\) makes \(b^2<a\). Apply the nonlinear energy contraction to \(T^{-1}y,T^{-1}z\); replace each energy by the square of its transformed norm and use the two inverse identities. The nonnegative square comparison yields \(\|\widetilde P_\xi y-\widetilde P_\xi z\|\le\sqrt{1-mh/(3\gamma)}\|y-z\|\). Every geometric fact needed for this estimate is derived explicitly.



Conjugating the noisy map means changing coordinates before and after its update. Under common noises its Lipschitz coefficient is the weighted pathwise contraction factor. The innovation still enters through its pushforward law.

\subsection{\texorpdfstring{Prove that Markov iteration commutes with the coordinate change}{Prove that Markov iteration commutes with the coordinate change}}


the OBABO laws is the repeated independent-noise product-and-pushforward transition of the original defined step; the transformed OBABO laws use its transformed realization. A generic induction proves the commuting diagram for any measurable maps satisfying \(T(F(x,\xi))=G(Tx,\xi)\). At each step, the pushforward formula for product measures proves that transforming the current state commutes with adjoining an independent innovation. Two applications of composition of pushforwards reduce the law identity to the pointwise conjugacy. Instantiating with the verified inverse gives \(T_\#\mathcal L_n(\mu)=\widetilde{\mathcal L}_n(T_\#\mu)\), so the contracting measures are the transformed sampler laws.



A measurable coordinate change commutes with a single transition pushforward after conjugation. Induction extends this identity to every Markov iterate. Thus contraction in transformed coordinates describes the same probability chain as the physical OBABO recursion.

\subsection{\texorpdfstring{Take the true coupling infimum and iterate the laws}{Take the true coupling infimum and iterate the laws}}


Let \(r=\sqrt{1-mh/(3\gamma)}\). For every coupling of two current laws, append one common innovation independently. Its two output marginals are the Markov outputs, and the pointwise bound gives output transport cost at most r times the input coupling cost. The previously proved random-map theorem takes the infimum over all input couplings; it assumes no optimal coupling exists. Induction, with probability of each iterated law proved internally, yields \(W_2(\widetilde{\mathcal L}_n\mu,\widetilde{\mathcal L}_n\nu)\le r^nW_2(\mu,\nu)\). The commuting law identity transfers this to the original laws measured in the modified geometry. The extended-real theorem is unconditional on transport finiteness. If the initial weighted transport distance is finite, the proved upper bound first establishes finiteness of the output and then permits conversion to the real-valued Wasserstein estimate. Stationarity and numerical target bias are separate claims, not hidden inputs to this law contraction.



For every input coupling, use the same independent fresh noise in both updates and bound the output cost by the pathwise contraction. Taking the coupling infimum yields one-step Wasserstein contraction. Iterating applies this to the transformed Markov laws.

\section{\texorpdfstring{Exact KIPLMC1 clock change for correlated Gaussian innovations}{Exact KIPLMC1 clock change for correlated Gaussian innovations}}
\label{proof-module-KIPLMC-ExponentialRescaling-lean}


The same positions and a velocity multiplication by \ensuremath{\surd}N convert physical-time KIPLMC1 into the shared unit-temperature exponential integrator for the sum potential. This file proves the covariance, Gaussian measure, update, initialization, and complete iterated-law identities. Both correlated noise components and their mutual covariance are transported explicitly.



The two Gaussian response coordinates scale by different clock powers. Their full covariance, spatial product law, initialization, and iterated transition identities establish an exact probabilistic conjugacy.

\subsection{\texorpdfstring{The Brownian clock changes each correlated block diagonally}{The Brownian clock changes each correlated block diagonally}}


Write \(c>0\), \(\gamma^\prime=c\gamma\), and \(h^\prime=h/c\). A time change contributes a Brownian factor \(c^{-1/2}\); the position response contributes another factor \(c^{-1}\). The resulting map on the two-component Gaussian block is \[D_c(v,x)=(c^{-1/2}v,c^{-3/2}x).\] The argument constructs this diagonal map as an continuous linear map on Euclidean two-space. The coordinate theorem verifies the factors directly. The kernel accessor theorem identifies the two entries of the algorithmic kernel vector with \(\psi_0\) and \(\psi_1\), so the subsequent proof can use the already arbitrary-index response scaling.



The velocity response contains one Brownian integration factor, whereas the position response has an additional time primitive. Their different powers of the clock scale are therefore unavoidable. A single common scalar multiplication would not transform the two-component innovation correctly.

\subsection{\texorpdfstring{Prove the covariance transformation and the Gaussian measure identity}{Prove the covariance transformation and the Gaussian measure identity}}


The covariance is the Gram integral \(C_{ij}(\gamma,h)=\int_0^h\psi_i(\gamma,t)\psi_j(\gamma,t)\,dt\), with \(i,j\in\{0,1\}\). The exact response identity \(\psi_i(c\gamma,t)=c^{-i}\psi_i(\gamma,ct)\), followed by the substitution \(s=ct\), gives \[C_{ij}(c\gamma,h/c)=c^{-(i+j+1)}C_{ij}(\gamma,h).\] These are precisely the entries of \(D_c C D_c^\mathsf T\), including the off-diagonal velocity-position covariance. The argument rewrites both response factors, moves their constants outside the integral, performs the interval-integral change of variables, and reduces the four finite index cases using \((\sqrt c)^2=c\).



Both measures are centered Gaussian measures. The mean of their linear pushforward is zero by commutation of the Bochner integral with the continuous linear map. Their covariance bilinear forms agree on the standard orthonormal basis by the displayed Gram identity, so Gaussian uniqueness proves \[(D_c)_\#\mathcal G_{\gamma,h}=\mathcal G_{c\gamma,h/c}.\] No Gaussian covariance or distributional scaling is introduced as a premise.



The off-diagonal covariance acquires the product of the two different diagonal scale factors. Checking that entry is essential: it verifies that the common Brownian forcing is preserved. Gaussian uniqueness then upgrades the covariance calculation to an equality of measures.

\subsection{\texorpdfstring{Transport the independent spatial blocks without losing correlation}{Transport the independent spatial blocks without losing correlation}}


The particle noise measure is the finite product of these two-component Gaussian blocks over the parameter and latent coordinate set. The argument applies the proved scalar Gaussian pushforward to each independent block, uses the exact product-measure mapping identity, and checks that assembling blocks into velocity and position vectors commutes with \(D_c\). Thus the full particle innovation obeys the same clock law. Independence between distinct spatial blocks is retained, while each block retains its exact velocity-position correlation.



The diagonal block transformation is applied independently at each spatial coordinate. Product pushforward preserves independence between blocks. Because each whole block is transformed together, its internal response correlation is retained exactly.

\subsection{\texorpdfstring{Conjugate the exponential update}{Conjugate the exponential update}}


For \(c=\sqrt N\), multiply velocity by \(c\) and leave position fixed. Under the changed clock, the force becomes \(\nabla V_N\), the friction is \(c\gamma\), and the step is \(h/c\). The exact deterministic exponential response conjugacy has already been proved. What remains is the additive innovation: \[\sqrt{2c\gamma}\,D_c\xi=\bigl(c\sqrt{2\gamma/N}\,\xi_v,\sqrt{2\gamma/N}\,\xi_x\bigr).\] One proves these two scalar identities from the positive square roots and the particle amplitude identity, then rewrites the shared exponential map. The resulting equality is an identity of the defined updates for every phase state and every supplied correlated innovation.



The transformed force is the gradient of the sum potential and the transformed velocity has unit temperature. The response scaling and thermal amplitude identities together identify both coordinates of the update. Position remains the original scaled particle position.

\subsection{\texorpdfstring{Pass the clock identity through the Markov laws}{Pass the clock identity through the Markov laws}}


The generic innovation-mapping lemma proves that replacing the noise law by \(R_\#Q\) is exactly substituting \(R\xi\) in each transition. Its proof is an induction on the true iterated Markov measure: at each step, the product law is pushed through \((z,\xi)\mapsto(z,R\xi)\), then through the measurable update. The phase conjugacy theorem applies this fact together with the proved update identity and the exact Gaussian block scaling. Consequently the entire KIPLMC1 Gaussian Markov law, pushed through \((v,x)\mapsto(\sqrt N v,x)\), is exactly the shared unit-temperature KLMC Markov law for \(V_N\) at \((\sqrt N\gamma,h/\sqrt N)\). The equality includes the noise measure and every iteration, rather than merely comparing update formulas.



One-step transition equality extends by induction to all iterated laws. At each successor step, product pushforward transports the fresh innovation measure and the update together. Thus the law identity includes the full source distribution, not only a deterministic formula.

\subsection{\texorpdfstring{Identify the thermal initialization and the unchanged position marginal}{Identify the thermal initialization and the unchanged position marginal}}


The physical velocity target is the explicit Gaussian pushforward \(N^{-1/2}\mathcal N(0,I)\). Multiplication by \(\sqrt N\) returns the standard Gaussian law. Since the position is left fixed, the independent thermal initialization becomes \(\mathcal N(0,I)\otimes\nu\) for every initial position law \(\nu\). One proves this using the product-measure pushforward and the already thermal velocity normalization. Finally, applying the position projection to the full iterated-law identity proves that the KIPLMC1 position law equals the unit-temperature KLMC position law at the changed clock. This supplies the exact law transfer needed to import the shared kinetic error estimate without duplicating its analysis.



The independent thermal initialization is transported to the independent standard-Gaussian initialization required by the shared theorem. Projecting the full chain-law identity onto position removes the velocity rescaling. This allows the unit-temperature sampling estimate to be imported without changing the observable.

\section{\texorpdfstring{Construct physical-\allowbreak{}time particle paths and derive their synchronous contraction}{Construct physical-time particle paths and derive their synchronous contraction}}
\label{proof-module-KIPLMC-Dynamics-lean}


The particle-average potential has the original pair curvature and smoothness constants, so the physical-time kinetic system can reuse the shared existence and nonlinear contraction theorems directly. The temperature changes the additive noise amplitude; the synchronous contraction does not depend on that amplitude because the common forcing cancels. The results here are explicit path statements, not stationary-law endpoints.



The average force meets the shared kinetic existence and contraction hypotheses. Common-input cancellation explains why physical-time deterministic contraction does not deteriorate with particle number or total state dimension.

\subsection{\texorpdfstring{Verify differentiability and smoothness of the average}{Verify differentiability and smoothness of the average}}


The lifted potential is differentiable whenever the pair potential is differentiable. Multiplying its value by the constant \(1/N\) preserves differentiability, yielding differentiability of the average potential. The gradient has already been identified pointwise with the particle force. Function extensionality turns those equalities into equality of the vector fields, allowing the average-gradient smoothness bound to rewrite the force's proved Lipschitz estimate as an estimate on the gradient. The result uses the original pair constant \(L\), independent of \(N\).



The gradient identity is essential for applying a potential-based contraction theorem: it ensures the defined force is the derivative of the stated average potential. Smoothness is then inherited with the original pair constant, while strong convexity is inherited by averaging the lifted inequality.

\subsection{\texorpdfstring{Apply the kinetic integral-\allowbreak{}equation construction to the real particle force}{Apply the kinetic integral-equation construction to the real particle force}}


With \(N>0\), the force is globally Lipschitz. The shared kinetic existence theorem therefore constructs a continuous phase path satisfying \(Y_t-Y_0=-\int_0^tD_F(Y_s)ds+(B_t,0)\) on any finite nonnegative horizon, for every continuous velocity forcing starting at zero. The call uses inverse mass one because the force already equals \(\nabla\bar U_N\). Choosing physical-time Brownian amplitude \(\sqrt{2\gamma/N}\) gives the intended scaled particle forcing, but a complete probabilistic strong-solution result also identifies the driving Brownian law and independence conditions.



The forcing is any continuous velocity input in the deterministic existence statement. The canonical Brownian choice supplies the desired thermal process afterward. This separates solving the integral equation from identifying the probability law of its driving source.

\subsection{\texorpdfstring{The nonlinear contraction rate is independent of particle dimension}{The nonlinear contraction rate is independent of particle dimension}}


Apply the shared potential contraction theorem to \(\bar U_N\). Its required hypotheses are derived by the average-potential differentiability, strong-convexity and smoothness theorems. Rewrite each gradient in the difference equation using \(\nabla\bar U_N=F\). With equal initial velocities, \(\gamma>0\), and \(\gamma^2\ge m+L\), the resulting bound is \(\|x_t-y_t\|\le\sqrt2e^{-mt/\gamma}\|x_0-y_0\|\). Neither \(N\) nor \(d_\theta+Nd_x\) appears in its rate. The theorem still explicitly requires continuous paths satisfying the noise-cancelled derivative equations; it does not replace construction of a stationary kinetic reference or taking the Wasserstein infimum.



The common additive noise cancels even though its amplitude depends on particle count. The remaining difference equation sees only the average-potential geometry. The dimension and temperature therefore affect equilibrium and discretization moments, while leaving the synchronous physical-time rate unchanged.

\section{\texorpdfstring{Iterating concrete stochastic segments to a Wasserstein bound}{Iterating concrete stochastic segments to a Wasserstein bound}}
\label{proof-module-KineticLangevin-LpConvergence-lean}


The numerical phase knots, auxiliary exact segments, reference segments, common raw inputs and convolution noises are explicit. The local coupling inequality is proved internally from these paths. The remaining law inputs state the reference Gaussian velocity and stationary position marginals. They do not assume a convergence conclusion.



The numerical and reference states are constructed together throughout the iteration. The scalar recurrence bounds their root-mean-square discrepancy, which then bounds the transport distance of their endpoint laws.

\subsection{\texorpdfstring{Accumulate update errors and project to laws}{Accumulate update errors and project to laws}}


For each step, the auxiliary path begins at the numerical knot, while the reference segment joins consecutive reference knots. The numerical knot update is exactly the exponential update with innovation \((B(h),\int_0^hB(s)\,ds)\). Rewriting these endpoint identities in the proved L\ensuremath{{}^{2}} one-step theorem produces the recurrence \(e_{k+1}\le e^{-mh/\gamma}e_k+3Mh^2(\sqrt p+e_k)/4\). The source step restriction \(4\gamma Mh\le m\), together with \(m\le M\), implies \(\gamma h\le1/4\), as needed for the local constant. The scalar recurrence theorem then gives \(e_n\le(1-3mh/(4\gamma))^ne_0+\gamma Mh\sqrt p/m\).



The numerical and reference position random variables provide a finite-cost Wasserstein coupling. The reference position law is the stated target \(\pi\), and position projection gives \(W_2\le\sqrt2e_n/\gamma\). Equal initial velocities identify \(e_0=\gamma\|X_0-Y_0\|_{L^2}\). The proof cancels friction and combines \(\sqrt2\sqrt p=\sqrt{2p}\), proving

\[W_2(\mathcal L(X_n),\pi)\le\sqrt2\left(1-\frac{3mh}{4\gamma}\right)^n\|X_0-Y_0\|_{L^2}+\frac{Mh\sqrt{2p}}m.\]

The initial term is the cost of the supplied coupling. Obtaining the infimum over all initial couplings, and discharging reference stationarity from the Brownian kinetic process, remain separate probabilistic requirements.



The proof is factored through a shared theorem with an explicit reference velocity radius \(C\ge0\). A uniform bound \(\|V_t^{\rm ref}\|_{L^2}\le C\) supplies \(C+e_k\) in the local defect and the terminal bias \(Mh\sqrt2C/m\). The standard-Gaussian statement is a thin specialization with the radius \(\sqrt p\) proved from its marginal law. This same shared theorem applies to particle temperature \(1/N\), using the separately proved thermal radius \(\sqrt{(d_\theta+Nd_x)/N}\); there is no duplicate exponential-integrator argument.



The recurrence is applied on a single common probability space throughout the iteration. Its initial norm is the cost of the prescribed initial position coupling times the transformation coefficient. Endpoint projection and the coupling-cost upper bound then convert the accumulated root-mean-square error into an output Wasserstein bound.

\section{\texorpdfstring{The shared exponential proof at particle temperature}{The shared exponential proof at particle temperature}}
\label{proof-module-KIPLMC-ReferenceBound-lean}


The first particle algorithm uses exactly the shared exponential step. Its analytical comparison should therefore reuse the same proof. This section applies the shared L\ensuremath{{}^{2}} kinetic theorem to the averaged lifted potential, uses the thermal Gaussian radius, and projects the resulting transport estimate to the MMLE parameter. The result has all three constants in the supplied KIPLMC1 theorem, with the initial supplied coupling cost. Its concrete integral-segment and reference-marginal hypotheses are explicit: construction of a canonical stationary reference and identification of the independent Gaussian chain are still needed for the unconditional sampler theorem.



The conditional comparison displays how the shared local kinetic proof yields initialization, discretization, and concentration terms. Later canonical constructions discharge its reference-path and stationary-law requirements.

\subsection{\texorpdfstring{Lift the physical particle force once}{Lift the physical particle force once}}


The averaged lifted potential \(\bar U_N=V_N/N\) has gradient \(F_N\), the physical-time force. The earlier geometry proves its original smoothness constant \(L\) and strong convexity constant \(m\), independently of particle count. the square-integrable force applies the existing pointwise the lifted gradient field construction to this gradient. It is therefore the particle force on almost-everywhere representatives, with a well-defined square-integrable output. No new exponential update or substitute gradient algorithm is introduced.



The lifted force is still the pointwise gradient of the average potential. Its square-integrable action inherits global Lipschitzness and the same curvature geometry. This permits the shared kinetic local-error proof to be reused without defining a different particle algorithm.

\subsection{\texorpdfstring{Inherit the local comparison, recurrence, and sharp equilibrium term}{Inherit the local comparison, recurrence, and sharp equilibrium term}}


Let \(S_n\) be numerical phase states and \(R_n\) reference phase states in L\ensuremath{{}^{2}}. On each interval, the auxiliary exact path starts at \(S_n\), the reference exact path starts at \(R_n\) and ends at \(R_{n+1}\), and their integral equations have the same additive input. The numerical endpoint is the shared exponential step with velocity noise \(B_h\) and position noise \(\int_0^hB_sds\). Its mild velocity and kinematic identities are explicit inputs. The reference velocity marginal is thermal Gaussian and its position marginal is the normalized lifted Gibbs target.



The thermal moment theorem supplies \(C=\sqrt{(d_\theta+Nd_x)/N}\). Applying the shared kinetic theorem gives the full-state transport estimate with initial term \(\sqrt2(1-3mh/(4\gamma))^n\|S_0^X-R_0^X\|_{L^2}\) and bias \(Lh\sqrt2C/m\), under \(\gamma^2\ge m+L\) and \(4\gamma Lh\le m\). Every local-error and recurrence step is inherited, not assumed as a transport inequality.



The particular L\ensuremath{{}^{2}} pair \((S_n^X,R_n^X)\) supplies a finite transport coupling, so real-valued conversion cannot conceal infinity. The parameter projection has operator norm at most one. The earlier sharp Gibbs concentration theorem contributes \(\sqrt{d_\theta/(Nm)}\). The argument invokes the proved transport-to-projected-RMS triangle bound, proves the projected moment is finite, rewrites it as the square root of the parameter second moment, and adds the three terms. The resulting theorem remains explicitly a reference comparison; it does not assert that the unresolved reference construction has already been discharged.



The three contributions have distinct origins: the initial coupling contracts, force freezing produces a linear-step bias, and the powered marginal contributes finite-particle concentration. The reference hypotheses identify exactly which path and law facts must be supplied before this conditional comparison becomes an unconditional sampler theorem.

\section{\texorpdfstring{The L\ensuremath{{}^{2}} recursion is the original defined KLMC chain}{The L² recursion is the original defined KLMC chain}}
\label{proof-module-KineticLangevin-LpDynamics-lean}


Analysis on equivalence classes of random vectors is useful only if it represents the intended algorithm. The lemmas here prove that the lifted exponential update, KLMC noise scaling and every recursively defined iterate agree almost everywhere with the original pointwise vector formulas. No Gaussian or independence assumptions are required for this algebraic correspondence.



The functional-analytic formulation is tied to pointwise random vectors by representative identities. Thus the sequence analyzed in second mean is exactly the original exponential recursion with its prescribed innovations.

\subsection{\texorpdfstring{Read L\ensuremath{{}^{2}} phase points and identify the exponential update}{Read L² phase points and identify the exponential update}}


the phase representative reads the velocity and position representatives of a pair of L\ensuremath{{}^{2}} classes at the same sample. A private affine-operation lemma combines the almost-everywhere identities for scalar multiplication, subtraction and addition. Apply it to the velocity formula \(\psi_0v-\psi_1G(x)+\eta_v\). For position, rewrite its affine expression as \(\psi_1v-\psi_2G(x)+(x+\eta_x)\), then use the same lemma. Finally the pointwise representative of the lifted drift identifies \(G(x)(\omega)\) with \(\nabla f(x(\omega))\). The two component equalities give equality of the phase update and the original pointwise the exponential update.



A square-integrable phase is represented almost everywhere by a pair of ordinary random vectors. Evaluating the exponential update on these representatives gives the original pointwise algorithm. The lifted force acts by applying the original gradient separately at each sample point.

\subsection{\texorpdfstring{Expand affine operations on L\ensuremath{{}^{2}} representatives}{Expand affine operations on L² representatives}}


The representative of a sum or scalar multiple of Lp classes agrees with the corresponding operation almost everywhere. The private helper gathers the four rules needed for \(aX-bY+Z\) into one common full-measure set. On that set, it rewrites the representative of the outer sum, the inner difference, and the two scalar multiples. The resulting equality is precisely the pointwise affine formula. This step accounts for the fact that Lp elements are equivalence classes; pointwise equalities of their arbitrarily selected representatives cannot be assumed everywhere.



Addition, scalar multiplication, and the bounded coordinate projections commute with passage to square-integrable equivalence classes. Consequently the affine update identities hold almost everywhere simultaneously for the finitely many operations in each step. No alternative numerical map is introduced by the functional-analytic formulation.

\subsection{\texorpdfstring{Retain the exact KLMC noise scaling through every iterate}{Retain the exact KLMC noise scaling through every iterate}}


The two phase components of a scaled L\ensuremath{{}^{2}} innovation agree almost everywhere with the same scalar times their representatives. Combining this with the exponential-step correspondence proves the KLMC correspondence, including its \(\sqrt{2\gamma}\) fluctuation-dissipation factor. Induction then proves the entire chain identity: at zero both chains use the supplied initial point; at a successor, intersect the almost-everywhere previous-chain identity with the current-step identity, and rewrite the recursive definitions. This shows the L\ensuremath{{}^{2}} error estimates concern the original numerical sampler for the supplied innovation stream, rather than a different abstract sequence.



Induct on the iteration number, using the same representative innovation at each successor step. The physical amplitude is retained in the induction identity. Thus the sequence whose error is analyzed in second mean has exactly the same random outputs as the defined KLMC recursion.

\section{\texorpdfstring{Actual KLMC output laws on concrete initial-\allowbreak{}coupling extensions}{Actual KLMC output laws on concrete initial-coupling extensions}}
\label{proof-module-KineticLangevin-ProductLaw-lean}


An estimate for one coupling must eventually be optimized over all initial couplings. Its numerical output law must therefore be independent of that choice. This section proves that independence for the KLMC chain. For every coupling \(c\) of the initial position law \(\mu\) and reference law \(\pi\), the sample measure is explicitly \(c\otimes Q\), where \(Q\) supplies the common initial velocity and all phase innovations. The velocity and position components of each innovation stay together, so their correlation is preserved. No convergence or stationary-reference assumption enters these law identities.



The numerical law is fixed independently of an initial transport plan. Exact source projections and representative identities then show that every coupling construction used in analysis has that same output marginal.

\subsection{\texorpdfstring{Build velocity and innovations on each product extension}{Build velocity and innovations on each product extension}}


The initial coupling has a probability first marginal, hence its joint measure is a probability measure. Its product with the fixed probability source \(Q\) is therefore a probability measure; the local instance records this derived fact. Pull the source velocity \(V\) back through the second projection and take its canonical L\ensuremath{{}^{2}} class. Do the same separately for the two components of each supplied phase innovation \(W_n\). Both use the same source point. This gives the pair of L\ensuremath{{}^{2}} classes needed by the kinetic analysis without assuming that velocity and position noise are independent. The realization reuses the existing product-noise construction.



The initial position coupling and the canonical Gaussian/\allowbreak{}Brownian source are independent product factors. Pulling source variables onto this extension preserves their moments and joint laws. Both numerical and reference paths can therefore use the same source while retaining their prescribed initial marginals.

\subsection{\texorpdfstring{Identify the chosen L\ensuremath{{}^{2}} representatives almost everywhere}{Identify the chosen L² representatives almost everywhere}}


The canonical Lp representative equals the original pulled-back random variable almost everywhere. For velocity this is the existing product-noise representative theorem. For a phase innovation, intersect the two full-measure sets for its components and apply pair extensionality. Thus the L\ensuremath{{}^{2}} pair represents precisely \(W_n(s)\), including all its correlation. This is an almost-everywhere equality of random variables, not a comparison of moments or covariances.



A square-integrable equivalence class has many pointwise representatives, but any two agree outside a null set. The coordinate and innovation identities are therefore stated almost everywhere. Their finite intersections suffice at each fixed iteration to identify the random update.

\subsection{\texorpdfstring{Define one numerical output law from initialization and fixed source}{Define one numerical output law from initialization and fixed source}}


Given a position \(x\) and source point \(s\), start the original pointwise KLMC recursion at \((V(s),x)\), use its phase innovations \(W_n(s)\), and project the final phase state to position. This is the numerical endpoint map. Its law is the pushforward of \(\mu\otimes Q\), recorded by the numerical position law. The definition contains no reference coupling. In particular, it is the same candidate output law for every coupling used to estimate its transport error.



Define the numerical endpoint first on the product of its own initial position distribution and the fixed independent source. This gives one output law without referring to a transport plan. Couplings are introduced afterward solely to estimate that law.

\subsection{\texorpdfstring{Prove the L\ensuremath{{}^{2}} numerical chain is exactly the pointwise chain}{Prove the L² numerical chain is exactly the pointwise chain}}


On \(c\otimes Q\), the L\ensuremath{{}^{2}} numerical start consists of the pulled-back common velocity and first initial-position coordinate. The reference position's finite second moment together with finite coupling cost supplies a square-integrable first coordinate. Apply the shared L\ensuremath{{}^{2}}-to-pointwise KLMC chain theorem. A countable intersection supplies agreement of every phase innovation on one full-measure set; intersect it with the initial velocity and position agreements. Pair extensionality identifies the exact initial phase state, function extensionality identifies the entire innovation sequence, and projection to the position coordinate finishes the equality. Equal initial velocities arise by using the same constructed L\ensuremath{{}^{2}} variable; no separate velocity coupling cost is inserted.



The initial phases agree almost everywhere with the pointwise construction. If the iterates agree at step \(n\), applying the same measurable update and the same innovation proves agreement at step \(n+1\). Induction identifies the entire analyzed recursion at every fixed time.

\subsection{\texorpdfstring{Derive law independence with almost-\allowbreak{}everywhere measurable data}{Derive law independence with almost-everywhere measurable data}}


The projection \(((x,y),s)\mapsto(x,s)\) sends \(c\otimes Q\) exactly to \(\mu\otimes Q\). Consequently every almost-everywhere measurable function of \((x,s)\) has the same law on both spaces. One proves this using the projection marginal and the almost-everywhere version of the composition rule for pushforwards. This avoids an unnecessary assumption that arbitrarily selected Lp representatives are everywhere measurable.



The output map is indeed almost-everywhere measurable: the source L\ensuremath{{}^{2}} assumptions give that property for velocity and both innovation components. The force is continuous by Lipschitz regularity, so the KLMC update is measurable. Induction propagates measurability through every phase iterate, and projection gives the position output. The final output-law theorem rewrites the L\ensuremath{{}^{2}} chain with its proved representative equality, then applies the projection-law identity. Both its initial marginal and its numerical law are derived rather than supplied as fields of a reference interface.



Projecting an extended coupling onto its first position coordinate and source gives the original numerical input law. Composing with the pointwise endpoint map proves that its output marginal is unchanged. This is the precise reason the output does not depend on which transport coupling was selected.

\subsection{\texorpdfstring{Check that the common output has total mass one}{Check that the common output has total mass one}}


The independent product of the initial probability law and source law is a probability measure. The preceding almost-everywhere measurability theorem ensures its pushforward remains a probability measure. The argument applies the probability property of a measurable pushforward. This establishes an numerical output law before any Wasserstein estimate is attached to it. Canonical Gaussian source construction and the stationary continuous reference are additional layers needed for the final unconditional sampler theorem.



The input product is a probability measure, and the endpoint map is almost everywhere measurable. Its pushforward therefore has total mass one. This probability verification is needed before transport comparisons are made with the target law.

\section{\texorpdfstring{Exact Hamiltonian energy and volume-\allowbreak{}preserving Verlet geometry}{Exact Hamiltonian energy and volume-preserving Verlet geometry}}
\label{proof-module-KineticLangevin-Hamiltonian-lean}


The Hamiltonian subflow is the kinetic equation with zero friction and zero noise. Its total energy is conserved by an exact cancellation between force work and kinetic-energy change. This section proves conservation for the Picard-constructed integral flow and proves that the explicit kick/\allowbreak{}drift/\allowbreak{}Verlet maps preserve product Haar volume. These are prerequisites for Gibbs stationarity; energy conservation alone and volume preservation of a numerical map do not yet establish invariance of the full kinetic diffusion.



Exact energy conservation belongs to the continuous Hamiltonian flow. Verlet instead supplies reversible volume-preserving geometry with a nonzero energy defect, making it suitable for a corrected invariant comparison.

\subsection{\texorpdfstring{Differentiate total energy along the Hamiltonian equations}{Differentiate total energy along the Hamiltonian equations}}


Define \(H(v,x)=f(x)+\|v\|^2/2\). If \(x'=v\) and \(v'=-\nabla f(x)\), the chain rule gives \((f(x))'=\langle\nabla f(x),v\rangle\), while the squared-norm derivative gives \((\|v\|^2/2)'=-\langle v,\nabla f(x)\rangle\). Symmetry of the real inner product makes their sum zero. The proof composes the Fréchet derivative of \(f\) with the position derivative, uses the gradient representation of that derivative, and combines it with the Hilbert squared-norm derivative. The statement uses right derivatives within \([t,\infty)\), allowing the initial endpoint of an integral solution.



For \(H(v,x)=\|v\|^2/2+f(x)\), the two derivative contributions are \(\langle v,-\nabla f(x)\rangle\) and \(\langle\nabla f(x),v\rangle\). They cancel exactly. This calculation concerns the true Hamiltonian differential equation, before any numerical discretization is introduced.

\subsection{\texorpdfstring{Turn vanishing right derivatives into exact conservation}{Turn vanishing right derivatives into exact conservation}}


A continuous scalar function with zero right derivative on \([0,T)\) is constant on \([0,T]\), including the final endpoint. The proof applies the scalar Gronwall comparison with coefficient and forcing both zero, first to \(F\) and then to \(-F\). These give \(F(t)\le F(0)\) and \(-F(t)\le-F(0)\). Their combination proves equality. This avoids assuming a derivative of the chosen extension beyond the terminal time or on negative time.



A continuous scalar function with zero right derivative on the interval is constant. Applying this to the composed Hamiltonian proves energy conservation up to the endpoints. The argument does not require differentiability of the flow at negative times.

\subsection{\texorpdfstring{Conserve energy for the constructed Hamiltonian integral flow}{Conserve energy for the constructed Hamiltonian integral flow}}


the Hamiltonian flow is the already constructed the kinetic solution map with friction zero, inverse mass one, and the zero compact forcing path. Its specification reduces to \(Y_t=Y_0-\int_0^tD(Y_s)ds\), with \(D(v,x)=(\nabla f(x),-v)\). Continuity of the Lipschitz drift gives the integral derivative. The argument identifies the selected path with its differentiable primitive on a right neighborhood, projects its derivative onto velocity and position, and applies the energy derivative cancellation. The scalar constancy theorem then yields \(H(Y_t)=H(Y_0)\). This is an exact statement about an existing solution, not an energy-preservation hypothesis on an abstract flow.



The selected Hamiltonian solution satisfies the integral equations with zero friction and zero forcing. Differentiating those equations supplies the hypotheses of the conservation calculation. Consequently conservation belongs to the constructed flow, rather than to an unspecified trajectory satisfying a formal equation.

\subsection{\texorpdfstring{Build the reversible proposal and preserve the Hamiltonian under reversal}{Build the reversible proposal and preserve the Hamiltonian under reversal}}


Velocity reversal \(R(v,x)=(-v,x)\) conjugates every kick and drift to its negative-time inverse. Applying this to the palindromic kick–drift–kick composition gives \(R\,V_h\,R=V_{-h}\). Since \(V_{-h}V_h=\mathrm{id}\), the proposal \(J_h=R\,V_h\) satisfies \(J_h^2=\mathrm{id}\). The argument rewrites the three component conjugacies and then uses the proved inverse Verlet identity. Negating velocity leaves its squared norm unchanged, so \(H(Rz)=H(z)\). The exact involution will give symmetric old/\allowbreak{}new proposal graph measures in the Metropolis argument.



Velocity reversal preserves the kinetic square and leaves position fixed. Composing it with a reversible Verlet map creates an involutive proposal. An involution is useful for Metropolis balance because its forward and backward proposal graphs are paired by the same transformation.

\subsection{\texorpdfstring{Expose the exact numerical energy defect}{Expose the exact numerical energy defect}}


Let \(w=v-hg(x)/2\), \(y=x+hw\), and \(v^+=w-hg(y)/2\). Expanding both kinetic-energy squares gives

\[H(v^+,y)-H(v,x)=f(y)-f(x)-\frac h2\langle w,g(x)+g(y)\rangle
+\frac{h^2}{8}\bigl(\|g(y)\|^2-\|g(x)\|^2\bigr).\]

The first part is a trapezoidal force-work remainder. The second is an explicit force-norm correction. First, the argument proves the two-half-kick norm identity for arbitrary vectors using squared-norm and inner-product expansions, then substitutes the intermediate Verlet velocity and position. This identity preserves the error terms that must later be bounded to control rejection; no second-order error estimate is assumed by definition.



The numerical energy change is obtained by evaluating the true Hamiltonian before and after one Verlet step. It generally does not vanish. Keeping it explicit permits both a Metropolis acceptance formula and quantitative rejection estimates derived from force smoothness.

\subsection{\texorpdfstring{Prove the elementary subflows preserve Haar volume}{Prove the elementary subflows preserve Haar volume}}


A drift \((v,x)\mapsto(v,x+hv)\) is a skew product whose second component is a translation for each fixed velocity. Translation invariance of Haar measure therefore proves preservation of the product measure. A kick \((v,x)\mapsto(v-hg(x),x)\) is the same construction after swapping the two phase coordinates; measurability of \(g\) suffices. The argument applies the measure-preserving skew-product theorem to translations, then conjugates the kick by the measure-preserving swap. Composition proves the complete Verlet map preserves product Haar volume. This derives the Jacobian-one property without presupposing differentiability of the force.



Each kick is a translation of velocity depending only on position; each drift is a translation of position depending only on velocity. Fubini and translation invariance show that these shears preserve product volume. Their compositions and velocity reversal preserve volume as well.

\subsection{\texorpdfstring{The reversed Verlet proposal is a measure-\allowbreak{}preserving involution}{The reversed Verlet proposal is a measure-preserving involution}}


For a Haar factor also invariant under negation, velocity reversal preserves product measure: apply the negation-preserving map in the velocity coordinate and identity in the position coordinate. Composing this with the proved volume-preserving Verlet map gives volume preservation of the explicit proposal \(J_h\). Together with \(J_h^2=\mathrm{id}\), this supplies the two exact properties used by deterministic Metropolis detailed balance. Neither property asserts that the uncorrected finite-step Verlet map itself preserves the Gibbs weight.



The reversed Verlet proposal combines volume preservation with involutivity. These are exactly the geometric ingredients needed for the accepted-flow symmetry argument. Approximate energy conservation is not needed to prove corrected invariance, although it controls how often proposals are rejected.

\section{\texorpdfstring{Prove Verlet converges to the Hamiltonian flow}{Prove Verlet converges to the Hamiltonian flow}}
\label{proof-module-KineticLangevin-HamiltonianConvergence-lean}


Exact reversibility and volume preservation of the Verlet map do not by themselves identify the continuous Hamiltonian flow. This section supplies that analytic limit. For a globally Lipschitz force, the Verlet iterates on a mesh of size \(T/(n+1)\) converge to every continuous Hamiltonian integral solution at time \(T\). Specializing to the gradient force identifies the limit with the Picard-selected Hamiltonian flow. The proof uses a local state-error calculation, uniform growth of an Euler comparison chain, a discrete Grönwall estimate, and the existing Euler convergence theorem. It needs no Hessian assumption.



The argument combines local consistency with finite-horizon stability. Comparing Verlet to Euler and Euler to the integral solution identifies the numerical limit with the constructed Hamiltonian dynamics.

\subsection{\texorpdfstring{Control the elementary shears and their symmetric composition}{Control the elementary shears and their symmetric composition}}


Use the standard maximum norm on phase pairs \((v,x)\). A kick of duration \(h\ge0\) is \((1+Mh)\)-Lipschitz when the force is \(M\)-Lipschitz, and a drift is \((1+h)\)-Lipschitz. The two half kicks around one drift therefore have factor \((1+Mh/2)^2(1+h)\). For \(0\le h\le1\), expansion and \(h^2,h^3\le h\) bound it by

\[\|\operatorname{Verlet}_h(z)-\operatorname{Verlet}_h(w)\|\le\bigl(1+(M+1)^2h\bigr)\|z-w\|.\]

First, the argument bounds the velocity and position projections by the phase norm. It rewrites each difference so the shared force difference is explicit, applies the triangle inequality and the force Lipschitz estimate, and handles the product norm using comparison with a maximum. The Verlet estimate composes the three bounds; a scalar polynomial inequality controls the resulting product. This is a stability estimate for the symmetric map, not an assumed numerical property.



Each shear is Lipschitz with a coefficient controlled by its step size and the force Lipschitz constant. Multiplying the stage bounds gives a stable one-step Verlet coefficient. A product of many such coefficients remains bounded on a fixed total time interval.

\subsection{\texorpdfstring{Compute and bound the local Verlet state defect}{Compute and bound the local Verlet state defect}}


Let \(v_{1/2}=v-hg(x)/2\) and \(x_+=x+hv_{1/2}\). Compare Verlet with the explicit Hamiltonian Euler step \((v-hg(x),x+hv)\). Their exact coordinate differences are

\[\Delta v=-\frac h2\bigl(g(x_+)-g(x)\bigr),\qquad\Delta x=-\frac{h^2}{2}g(x).\]

The Lipschitz bound gives \(\|g(x)\|\le(M+1)(\|z\|+\|g(0)\|)\) and \(\|v_{1/2}\|\le(M+2)(\|z\|+\|g(0)\|)\) for \(h\le1\). Consequently,

\[\|\operatorname{Verlet}_h(z)-(z-hD(z))\|\le h^2(M+1)(M+2)(\|z\|+\|g(0)\|),\quad D(v,x)=(g(x),-v).\]

One proves the two displayed coordinate identities by vector-space algebra. It then applies the force Lipschitz estimate to \(x_+-x=hv_{1/2}\), rewrites scalar-multiple norms, and bounds each coordinate before taking their maximum. This deliberately sufficient quadratic local estimate proves convergence under only a Lipschitz force; it does not claim a second-order global rate under those minimal assumptions.



Expand the two kicks and the intervening drift relative to one explicit Euler step. Lipschitz continuity controls the force evaluated at the displaced position. The difference is bounded by a quadratic step-size factor times the current phase envelope.

\subsection{\texorpdfstring{Bound the entire Euler comparison chain on a fixed horizon}{Bound the entire Euler comparison chain on a fixed horizon}}


Write \(X_{k+1}=X_k-hd(X_k)\), with \(d\) globally \(L\)-Lipschitz and \(nh=T\). The distance from the initial point satisfies

\[e_{k+1}\le(1+Lh)e_k+h\|d(X_0)\|,\qquad e_0=0.\]

The finite recurrence estimate and \((1+Lh)^k\le e^{Lkh}\le e^{LT}\) yield, for every \(k\le n\),

\[\|X_k\|\le\|X_0\|+T\|d(X_0)\|e^{LT}.\]

The bound is uniform over the number of mesh steps. The argument derives the recurrence by comparing one Euler update with the update of the fixed initial point, invoking the previously proved Lipschitz bound for explicit Euler. the finite discrete Grönwall inequality sums the recurrence. The inequality \(1+r\le e^r\), the exponential power identity, and the partition equation \(nh=T\) control the amplification. No boundedness of the numerical trajectory is added as a hypothesis.



The Euler recursion satisfies a discrete Grönwall inequality. On a finite horizon this gives a uniform bound on all intermediate states in terms of initialization and the force at the origin. Such a bound prevents the local defect coefficient from changing uncontrolledly with iteration.

\subsection{\texorpdfstring{Propagate the local difference between Verlet and Euler}{Propagate the local difference between Verlet and Euler}}


Let \(V_k\) be the Verlet iterates and \(X_k\) the Euler iterates from the same initial phase. Split their next difference into a Verlet stability term and the local Verlet-versus-Euler defect evaluated at \(X_k\). The proved Euler growth bound makes the local coefficient uniform in \(k\). With \(A=(M+1)^2\) and

\[C=(M+1)(M+2)\bigl(\|z\|+T\|D(z)\|e^{(M+1)T}+\|g(0)\|\bigr),\]

the error satisfies \(e_{k+1}\le(1+Ah)e_k+h^2C\), hence

\[\|V_n-X_n\|\le ThC e^{AT}.\]

The argument constructs the two concrete chains and proves their recursive update formulas. The triangle inequality combines the established stability and local-defect theorems. The finite recurrence bound gives \(nh^2C(1+Ah)^n\); the partition identity converts \(nh^2\) into \(Th\), and the exponential bound removes dependence on \(n\). This is the full finite-horizon numerical comparison required by the limit theorem.



The error at the next step is at most the propagated previous error plus the new local defect. Iterating this inequality gives a geometric sum with one term per time interval. With total time fixed, the quadratic local discrepancy yields a vanishing global discrepancy.

\subsection{\texorpdfstring{Send the mesh size to zero}{Send the mesh size to zero}}


Choose \(h_n=T/(n+1)\). This mesh tends to zero, and eventually \(h_n\le1\), so the finite comparison applies. Its right side is the vanishing mesh times a fixed, finite coefficient depending only on the horizon, force, and initial phase. Therefore the norm of the difference between the two endpoints tends to zero.



One proves the mesh limit from convergence of a constant divided by the mesh count, obtains the eventual small-step condition from the neighborhood of zero, and applies the squeeze theorem to the nonnegative error norm. norm convergence of the difference upgrades this scalar estimate to convergence of the phase-vector difference itself.



As the mesh count tends to infinity, step length tends to zero and the accumulated defect tends to zero. The stability factors remain bounded by an exponential depending on the fixed horizon. The comparison therefore proves convergence, rather than merely formal consistency of one step.

\subsection{\texorpdfstring{Identify the selected Hamiltonian flow as the numerical limit}{Identify the selected Hamiltonian flow as the numerical limit}}


The existing Euler theorem states that a globally Lipschitz drift's matching-input Euler approximations converge to every continuous solution of its integral equation. Here the additive input is zero and the negative Hamiltonian drift is \(D(v,x)=(g(x),-v)\), which has Lipschitz constant at most \(M+1\). Adding the vanishing Verlet-minus-Euler difference to that convergent Euler endpoint proves convergence of Verlet itself.



For \(g=\nabla f\), the selected Hamiltonian flow is already constructed by Picard from this precise integral equation. the selected solution’s integral equation supplies its continuity, initialization, and integral identity. The final theorem therefore proves

\[\operatorname{Verlet}_{T/(n+1)}^{\,n+1}(z)\longrightarrow\Phi_T^f(z).\]

The argument invokes convergence of additive-input Euler approximation with zero input, combines the two limits using addition of convergent terms, and simplifies the exact vector identity \((V_n-X_n)+X_n=V_n\). The specialization unfolds only the zero-friction, unit-inverse-mass flow definition and its normalized zero forcing. Neither the desired numerical convergence nor a stationarity conclusion occurs among the hypotheses.



The Euler scheme already converges to the unique integral solution. Combining that convergence with the Verlet–Euler comparison identifies the Verlet limit with the same solution. This makes the later stationary-limit argument about the constructed Hamiltonian flow.

\section{\texorpdfstring{Concrete Verlet–OU splitting and local defects}{Concrete Verlet–OU splitting and local defects}}
\label{proof-module-KineticLangevin-KineticSplitting-lean}


The splitting is an explicit numerical map: velocity Verlet for the Hamiltonian drift, then exponential velocity damping and a supplied filtered additive increment. Its stability and local kinetic Euler defect are proved from the map algebra. The filtered increment is written by deterministic integration by parts, matching the canonical Brownian segment identity proved separately.



The deterministic splitting is driven by filtered increments of a prescribed continuous input. Its exact consistency estimates prepare a convergence theorem that later specializes to canonical Brownian forcing.

\subsection{\texorpdfstring{Damp velocity while retaining position}{Damp velocity while retaining position}}


The map is \(D_\delta(v,x)=(\delta v,x)\). For \(0\le\delta\le1\), its velocity difference has norm \(\delta\|v-w\|\) and its position difference is unchanged. Each is bounded by the original maximum product norm, hence the full map is nonexpansive. One proves the two coordinate bounds and combines them using comparison with a maximum. This uses the usual product norm rather than an undeclared Euclidean phase norm.



The thermostat modifies velocity but leaves position unchanged. Its deterministic damping is contractive, and its additive response is supplied separately. This separation permits a pathwise convergence proof for a prescribed continuous input before identifying that input probabilistically.

\subsection{\texorpdfstring{Construct and stabilize the Verlet–OU step}{Construct and stabilize the Verlet–OU step}}


The explicit step is \(S_h(z,N)=O_{e^{-\gamma h},N}(V_hz)\). Two synchronously driven steps have exactly the same additive increment, which cancels. Exponential damping lies in \((0,1]\) for nonnegative friction and step, so it is nonexpansive. The already proved Verlet stability bound gives \[\|S_h(z,N)-S_h(w,N)\|\le[1+(M+1)^2h]\|z-w\|\] for \(h\le1\). First, the argument checks the exact noise-cancellation identity and then composes these two bounds.



The numerical map combines a Hamiltonian Verlet update with the damped velocity response. Stagewise Lipschitz estimates give stability of the composition. The forcing increment is the filtered increment associated with the same input interval.

\subsection{\texorpdfstring{Derive exponential damping errors without dividing by friction}{Derive exponential damping errors without dividing by friction}}


The response balances \(\gamma\psi_1=1-e^{-\gamma h}\) and \(\gamma\psi_2=h-\psi_1\) imply \(e^{-\gamma h}-1+\gamma h=\gamma^2\psi_2\). Since \(0\le\psi_2\le h^2/2\), the absolute Euler remainder is bounded by \(\gamma^2h^2/2\). Also the exponential tangent inequality gives \(0\le1-e^{-\gamma h}\le\gamma h\). The argument obtains both estimates from exact identities and checks all sign conditions. The formulas remain valid at zero friction.



The bounds follow from integral representations of \(1-e^{-s}\) and its linear remainder. They remain valid at zero friction because no division by friction is used. This is useful when comparing the splitting with the deterministic Hamiltonian limit.

\subsection{\texorpdfstring{Compare the damped Hamiltonian Euler stage with full kinetic Euler}{Compare the damped Hamiltonian Euler stage with full kinetic Euler}}


Damp the Hamiltonian Euler stage and subtract full kinetic Euler. Their position components agree exactly, while their velocity difference is \((e^{-\gamma h}-1+\gamma h)v+(1-e^{-\gamma h})h\,g(x)\). The preceding scalar estimates bound this by a quadratic expression. Lipschitz force growth \(\|g(x)\|\le M\|z\|+\|g(0)\|\) yields the displayed \(h^2[\gamma^2/2+\gamma(M+1)](\|z\|+\|g(0)\|)\) bound. One proves the exact vector identity, bounds its two terms, and combines nonnegative scalar factors.



An Euler comparison exposes the kinetic drift as the sum of Hamiltonian transport and friction. The splitting’s damping introduces higher-order terms because these pieces are applied in succession. Those terms are kept explicitly in the local consistency estimate.

\subsection{\texorpdfstring{Bound the complete deterministic kinetic Euler defect}{Bound the complete deterministic kinetic Euler defect}}


Insert the damped Hamiltonian Euler stage between the split step and full kinetic Euler. Nonexpansiveness of damping preserves the proved Verlet-versus-Hamiltonian-Euler defect. Add the damped-stage defect just established. The noise increment is retained exactly on both sides and cancels, giving \[\|S_h(z,N)-[z-hd(z)+(N,0)]\|\le h^2C(\|z\|+\|g(0)\|),\quad C=(M+1)(M+2)+\gamma^2/2+\gamma(M+1).\] The argument applies the triangle inequality through this explicit intermediate state and normalizes the coefficient algebra.



The full defect includes the Verlet approximation error, the damping remainder, and their interaction. Each is controlled by the phase envelope and a small-step factor. Summing these bounds gives a local state comparison with the kinetic Euler increment.

\subsection{\texorpdfstring{Use the filtered additive increment}{Use the filtered additive increment}}


For a continuous additive velocity input \(B\), define \[N_{a,h}=B_{a+h}-B_a-\gamma\int_0^h e^{-\gamma(h-s)}(B_{a+s}-B_a)\,ds.\] This is the deterministic integration-by-parts representation of the exact OU response. If the input oscillation over the interval is at most \(\eta\), exponential damping is at most one, so \(\|N_{a,h}-(B_{a+h}-B_a)\|\le\gamma h\eta\). The proof subtracts the matching increment exactly, moves the scalar outside the norm, and bounds the integral by its interval length and pointwise integrand bound. The Brownian segment module independently proves precisely this identity for the constructed Gaussian source.



The additive increment must be the filtered response to the supplied input, rather than an unrelated Gaussian with the same variance. Its exact relation to the input gives the integral equation in the limit. This is what identifies the splitting limit with the desired kinetic process.

\section{\texorpdfstring{The Verlet–OU chain converges to the kinetic integral solution}{The Verlet–OU chain converges to the kinetic integral solution}}
\label{proof-module-KineticLangevin-KineticSplittingConvergence-lean}


This section proves the missing numerical limit directly in a Banach space, including the intended space of square-integrable random variables. The scheme uses explicit Verlet stages, exact exponential damping, and the filtered increments of the same continuous additive input as the kinetic equation. Compact continuity supplies every local modulus used in the proof. No stationary solution, local-error contract, or numerical-convergence hypothesis is assumed.



The filtered-input splitting converges to the selected kinetic integral solution as the mesh is refined. Uniform continuity controls forcing variation, while stability prevents local errors from accumulating uncontrollably.

\subsection{\texorpdfstring{Construct the splitting on the grid input}{Construct the splitting on the grid input}}


Starting from \(X_0=z_0\), the next state is \(X_{k+1}=S_h(X_k,N_{kh,h})\), where \(S_h\) is the explicit Verlet–OU step and \(N_{a,h}\) is the filtered increment of the supplied continuous input \(B\). The proof defines this recursion directly. Each stage uses the correct grid time \(kh\); the forcing is not replaced by an unrelated Gaussian with the same marginal variance.



The recursion uses increments of the same continuous path on a uniform time grid. Its initialization is fixed independently of the mesh. Thus refinement changes only the discretization, making convergence to one candidate continuous solution a meaningful statement.

\subsection{\texorpdfstring{Compare one numerical increment with the exact kinetic path}{Compare one numerical increment with the exact kinetic path}}


Let the path satisfy \(Y_t-z_0=-\int_0^t d(Y_s)\,ds+(B_t,0)\), where \(d(v,x)=(\gamma v+g(x),-v)\). Subtract its equations at \(a\) and \(a+h\). Freezing its drift gives an exact Euler identity with remainder \(D_{a,h}=\int_0^h[d(Y_{a+s})-d(Y_a)]\,ds\). A drift oscillation bound \(\eta\) gives \(\|D_{a,h}\|\le h\eta\). The matching OU input differs from \(B_{a+h}-B_a\) by at most \(\gamma h\eta\) when the input oscillation is also at most \(\eta\). Finally, the proved explicit splitting defect contributes at most \(h^2C(R+\|g(0)\|)\), where \(\|Y_t\|\le R\) and \(C=(M+1)(M+2)+\gamma^2/2+\gamma(M+1)\). The argument constructs the grid and interval-membership facts, derives the exact increment equation through adjacent-interval integration, proves the local-defect identity, and bounds each of these three terms.



For \(e_k=\|X_k-Y_{kh}\|\), synchronous numerical stability then gives \[e_{k+1}\le[1+(M+1)^2h]e_k+h\{hC(R+\|g(0)\|)+(1+\gamma)\eta\}.\] The initial equation and \(B_0=0\) prove \(e_0=0\). The finite recurrence estimate and \((1+ah)^n\le e^{aT}\), with \(nh=T\), yield the displayed finite-horizon bound. Every modulus premise belongs to this intermediate estimate and is derived from continuity in the next theorem.



Subtract the exact integral equation over one grid interval from the numerical increment. Separate the local deterministic consistency defect from the variation of the exact drift within that interval. Continuity of the exact path controls the latter by a modulus tending to zero.

\subsection{\texorpdfstring{Derive all moduli and prove convergence}{Derive all moduli and prove convergence}}


The continuous exact path is bounded on \([0,T]\). Its drift path is continuous by global Lipschitz continuity of \(g\), and the additive input \(B\) is continuous. Compactness makes both uniformly continuous. Given any target error \(\varepsilon>0\), choose a common oscillation tolerance \(\eta=\varepsilon/(D+1)\), where the explicit finite horizon constant \(D\) contains the state bound, the splitting defect coefficient, friction, and the stability exponential. Since \(h_n=T/(n+1)\to0\), the mesh is eventually at most one, at most \(\eta\), and below both continuity radii. The finite error estimate is then at most \(\eta D<\varepsilon\). The argument obtains the two compact uniform-continuity radii, combines the four eventual mesh bounds, instantiates the finite estimate, and checks this final scalar comparison. Thus the concrete scheme converges to every continuous kinetic integral solution driven by the same input.



Uniform continuity on a compact horizon gives moduli for both the input and the exact solution. Stability propagates local defects without uncontrolled amplification. A discrete Grönwall estimate then shows that the maximum grid-point error tends to zero as the mesh is refined.

\subsection{\texorpdfstring{Identify the limit with the constructed kinetic flow}{Identify the limit with the constructed kinetic flow}}


The earlier Picard construction supplies a continuous solution for every initial state and every compact continuous forcing path, with its kinetic integral identity already proved. Its normalized forcing is globally continuous and vanishes at zero. Substituting these facts into the generic Banach convergence theorem proves that the splitting limit is exactly the kinetic solution map at the requested horizon. This ties the approximation to the existing jointly continuous solution map, not to an unspecified solution chosen by a new contract.



The limit satisfies the same integral equation and starting condition as the selected kinetic flow. Uniqueness identifies the two. This conclusion is pathwise for the prescribed input and can subsequently be applied to the canonical Brownian construction in second mean.

\section{\texorpdfstring{The L\ensuremath{{}^{2}} splitting is the vector algorithm almost everywhere}{The L² splitting is the vector algorithm almost everywhere}}
\label{proof-module-KineticLangevin-KineticSplittingLp-lean}


The deterministic convergence theorem applies to a Banach space of square-integrable random variables. To connect that analysis to the intended stochastic method, this section checks every elementary stage on representatives. The lifted force is the original force almost everywhere, and the supplied noise is retained unchanged. Induction then identifies the entire L\ensuremath{{}^{2}} recursion with the vector recursion.



Every splitting stage agrees almost everywhere with its original vector operation. The Hilbert-space convergence proof therefore concerns the vector chain driven by the stated filtered innovations.

\subsection{\texorpdfstring{Expose the chain driven by explicit filtered innovations}{Expose the chain driven by explicit filtered innovations}}


The innovation chain starts at its supplied phase state and recursively applies the explicit Verlet–OU step with innovation \(N_k\). The earlier continuous-input chain is exactly this chain with \(N_k=N_{kh,h}(B)\). One proves this equality by induction, unfolding the two recursions at the successor step. This representation lets the probability layer identify the innovation array law independently of the deterministic input analysis.



The numerical splitting is first written as a deterministic recursion with explicit filtered inputs. This makes the step index and the additive noise visible. Later identification with a pointwise chain can then be proved directly by induction.

\subsection{\texorpdfstring{Check force kicks on L\ensuremath{{}^{2}} representatives}{Check force kicks on L² representatives}}


The lifted force \(G\) is constructed from the Lipschitz force \(g\), and the shared lifting theorem gives \(G(X)(\omega)=g(X(\omega))\) almost everywhere. The representative identities for subtraction and scalar multiplication therefore give \([V-hG(X)](\omega)=V(\omega)-hg(X(\omega))\). Position is unchanged. The proof intersects the three almost-everywhere events and proves equality of the two phase coordinates. No new vector-field or output-law assumption is introduced.



The lifted force is evaluated by applying the original gradient to each representative position. A half kick subtracts the prescribed scalar multiple of that force. The representative identity verifies that the lifted kick is exactly the physical kick almost everywhere.

\subsection{\texorpdfstring{Check free transport and the full supplied thermostat noise}{Check free transport and the full supplied thermostat noise}}


The representative of \(X+hV\) is \(X(\omega)+hV(\omega)\), proving the free drift identity. Likewise the representative of \(\delta V+\sigma\Xi\) is \(\delta V(\omega)+\sigma\Xi(\omega)\), proving the OU stage identity. The argument explicitly intersects the addition and scalar-multiplication representative events and retains the exact supplied innovation \(\Xi\). These are pathwise stage identities and do not presume any Gaussian distribution.



Transport adds the current velocity times the step to position, while thermostat damps velocity and adds its filtered innovation. All these affine operations commute with representatives. Their almost-everywhere identities therefore preserve both phase components.

\subsection{\texorpdfstring{Compose the stages into the Verlet–OU map}{Compose the stages into the Verlet–OU map}}


The Verlet step is the kick–drift–kick composition. Apply the representative identity at each intermediate L\ensuremath{{}^{2}} state, intersect the resulting almost-everywhere events, and substitute them in order. The complete splitting then composes this Verlet map with the OU stage, using exactly the damping \(e^{-\gamma h}\) and supplied filtered innovation. The proof performs these successive substitutions, so the state at which the second force is evaluated remains the updated position.



Intersect the finitely many full-measure sets on which the stage identities hold. On that intersection their composition is exactly the original vector splitting step. This verifies the complete map rather than only its separate projected distributions.

\subsection{\texorpdfstring{Identify every numerical iterate with the original vector chain}{Identify every numerical iterate with the original vector chain}}


The initial phase representatives agree by definition. Assuming agreement at step \(n\), the full-step identity and the induction hypothesis give agreement at step \(n+1\), driven by the same innovation representative \(N_n(\omega)\). The proof intersects these two almost-everywhere events and rewrites the recursive definitions. Hence every L\ensuremath{{}^{2}} iterate analyzed by the convergence theorem is the vector splitting algorithm on the same probability source.



Initialization and innovation representatives agree with the pointwise construction. Successive application of the verified step identity proves agreement at every fixed iteration. Consequently distributional conclusions for the pointwise sampler apply to the sequence analyzed in the Hilbert space.

\section{\texorpdfstring{Actual phase laws and their L\ensuremath{{}^{2}} limits}{Actual phase laws and their L² limits}}
\label{proof-module-KineticLangevin-LpPhaseLaw-lean}


The phase variable is a pair of L\ensuremath{{}^{2}} classes, and its law is the pushforward of their measurable representatives. Dual characteristic functions depend continuously on this pair. This identifies a limiting stationary phase law from bounded-test estimates without incorrectly treating the usual maximum product norm as a Hilbert norm.



Characteristic functions bridge second-mean convergence and equality of probability laws. This permits a stationary law to be identified from numerical weak-error bounds without assuming almost-sure convergence of the full sequence.

\subsection{\texorpdfstring{Verify measurability of phase representatives}{Verify measurability of phase representatives}}


Each L\ensuremath{{}^{2}} class has a fixed strongly measurable representative. Taking the pair of its velocity and position representatives gives a measurable phase map. The argument uses the two strong-measurability results and the measurable product constructor, making the subsequent pushforward an measure.



The Euclidean phase representative is formed from its two measurable coordinate representatives. Its finite second moment follows from the sum of their squared norms. Hence it has a well-defined probability law and may be used as a transport coupling variable.

\subsection{\texorpdfstring{Reduce phase characteristic functions to one real L\ensuremath{{}^{2}} variable}{Reduce phase characteristic functions to one real L² variable}}


For a continuous real linear functional \(L\) on phase space, let \(L_v(v)=L(v,0)\) and \(L_x(x)=L(0,x)\). The real L\ensuremath{{}^{2}} variable \(R=L_v(V)+L_x(X)\) represents \(L(V(\omega),X(\omega))\) almost everywhere. Its law therefore has scalar characteristic function at frequency one equal to the dual characteristic function of the phase law. The map \((V,X)\mapsto R\) is continuous, since the two continuous linear maps lift continuously to L\ensuremath{{}^{2}} and addition is continuous. The existing scalar characteristic-function continuity theorem proves the desired phase continuity. The argument explicitly intersects the representative events for both lifted linear maps and their sum, proves the coordinate decomposition by linearity, rewrites the pushforward through \(L\), and composes the two continuous maps.



Pairing the phase with a fixed frequency vector gives a real square-integrable variable. Its characteristic exponential is bounded, and convergence in second mean implies convergence of these scalar pairings in first mean. This gives continuity of characteristic functions under the phase limit.

\subsection{\texorpdfstring{Identify the law of an L\ensuremath{{}^{2}} limit from vanishing weak errors}{Identify the law of an L² limit from vanishing weak errors}}


Suppose phase pairs \(S_n\) converge in the product of L\ensuremath{{}^{2}} spaces to \(S\), and their bounded complex-test errors relative to a fixed finite measure \(\pi\) are eventually bounded by \(\delta_n\to0\). For each dual functional use the norm-one test \(F(z)=e^{iLz}\). The error squeeze gives convergence of the approximating characteristic functions to that of \(\pi\); the continuity theorem gives convergence of the same sequence to the characteristic function of the limiting phase law. Uniqueness of limits equates them, and uniqueness of finite measures from their dual characteristic functions gives \(P\circ S^{-1}=\pi\). The argument rewrites the phase expectation by the measurable pushforward formula and invokes the scalar squeeze and measure-uniqueness theorems. The weak estimates must be proved by the concrete numerical and probability layers; stationary output is not assumed.



If approximating phase laws differ from the target by a vanishing bounded-test error, their characteristic functions converge to those of the target. The second-mean limit has the same characteristic functions. Uniqueness then identifies its phase law with the target.

\section{\texorpdfstring{The splitting chain has its independent-\allowbreak{}noise Markov law}{The splitting chain has its independent-noise Markov law}}
\label{proof-module-KineticLangevin-SplittingLaw-lean}


A law identity must be established for the random variables actually used by the construction. The input here is the complete joint law of the initial phase and the entire innovation stream. Composing this law with the measurable defined recursion proves the endpoint law. The same identity is then established for the L\ensuremath{{}^{2}} realization using its proved almost-everywhere representatives.



The lifted and pointwise descriptions are proved to have identical outputs under the complete input law. Consequently the splitting analyzed for convergence is the same independent-noise Markov chain used in stationarity transfer.

\subsection{\texorpdfstring{Push the full input law through the recursion}{Push the full input law through the recursion}}


Let \(C_n(z,w)\) be the explicit recursion run for \(n\) steps. Its joint measurability has already been proved by induction. If \((X,(N_k)_k)\) has the product law \(\pi\otimes Q^{\otimes\mathbb N}\), the pushforward composition rule says its output has law \((\pi\otimes Q^{\otimes\mathbb N})\circ C_n^{-1}\). The existing finite-prefix proof identifies this with the iterated Markov law \(\pi P^n\). On an explicit initial-law/\allowbreak{}source product, the initial phase and a source stream are independent by the product pushforward identity. One proves both statements directly using composition of distributional identities and the pushforward formula for product measures.



The complete input consists of the initial phase and the full innovation stream. The recursion is a measurable function of that input. Pushing the product input law through the recursion determines its endpoint law uniquely.

\subsection{\texorpdfstring{Match the two concrete realizations step for step}{Match the two concrete realizations step for step}}


the splitting sequence uses Verlet followed by the exact OU update. Induction shows it is precisely the additive-noise sequence for that step. Continuity of the force makes every kick, drift and thermostat continuous, so their composition is measurable. The law theorem can therefore be applied to the splitting map, with no additional transition-measurability premise.



The lifted and pointwise maps agree at initialization and at every successor step under the proved representative identities. Induction gives equality of their numerical endpoints almost everywhere. Their pushforward laws are consequently equal.

\subsection{\texorpdfstring{Eliminate the choice of L\ensuremath{{}^{2}} representatives}{Eliminate the choice of L² representatives}}


The lifted chain has already been proved equal almost everywhere to the pointwise vector chain with the same initial state and innovations. Replace its pushforward using this equality, then apply the full-input-law theorem. This proves the L\ensuremath{{}^{2}} endpoint law is the independent-noise Markov iterate. Quotient representatives cannot create an extra source of randomness or change the chain law.



Changing a representative on a null set changes neither its law nor the law of a measurable finite-step recursion. The construction therefore defines a sampler law intrinsically. This is needed when the analytic convergence proof is carried out in equivalence classes of random variables.

\section{\texorpdfstring{The defined KLMC law is the independent-\allowbreak{}noise Markov law}{The defined KLMC law is the independent-noise Markov law}}
\label{proof-module-KineticLangevin-ProductMarkovLaw-lean}


This section identifies the complete law of the concrete exponential-integrator recursion. It proves the independent phase/\allowbreak{}source product by coordinate reassociation, carries the whole innovation-stream law to the initialization product, and identifies the resulting position law with the iterated Markov transition.



The canonical source’s entire innovation stream is independent of the true initial phase. Pushing that joint law through the original recursion identifies the analytic output with the Gaussian Markov iterate.

\subsection{\texorpdfstring{Prove the true initial phase and source product law}{Prove the true initial phase and source product law}}


The concrete sampler begins on \(\nu\otimes(\kappa\otimes Q)\): the first coordinate is initial position, the second is initial velocity, and the third is the full future random source. The coordinate map \((x,(v,\omega))\mapsto((v,x),\omega)\) has law \((\kappa\otimes\nu)\otimes Q\). One proves this by the measure-preserving product associator followed by the measure-preserving swap of position and velocity. If the source function \(W\) has the entire iid innovation-stream law \(\eta^{\otimes\mathbb N}\), composing that map gives \[\operatorname{Law}((V_0,X_0),(\xi_k)_{k\ge0})=(\kappa\otimes\nu)\otimes\eta^{\otimes\mathbb N}.\] This proves joint independence of initialization and the whole future stream from the product source; there is no separate independence hypothesis at each iteration.



The independent standard Gaussian velocity and initial position law form the initial phase product. The canonical innovation stream is independent of this whole phase. Reassociating the product factors proves the joint input law needed for the explicit Markov recursion.

\subsection{\texorpdfstring{Identify the recursion without changing its innovations}{Identify the recursion without changing its innovations}}


The shared KLMC recursion is exactly the ordinary recursion \(z_{k+1}=\mathcal K_h(z_k,\xi_k)\). One proves equality with the generic noise-driven chain by induction: both initial values are the supplied phase state, and the successor equations apply the same update to the same preceding state and the same correlated phase innovation. This identity allows the generic independent-noise law theorem to describe the defined exponential chain.



The original pointwise KLMC update is applied with the same canonical innovation at each index. Its endpoint agrees with the fixed-source numerical map by induction. Consequently the output-law identification does not rely on replacing the source by a different but similar noise sequence.

\subsection{\texorpdfstring{Prove the full Markov position-\allowbreak{}law identity}{Prove the full Markov position-law identity}}


The pointwise sampler uses \((V_0,X_0)=(v,x)\) and innovations \(W(\omega)_k\) on the explicit initialization/\allowbreak{}source product. The gradient Lipschitz assumption proves joint measurability of the defined update. The joint initial/\allowbreak{}stream law then identifies the true phase pushforward of the noise-driven recursion with the \(n\)-fold Markov transition. The argument applies the position projection to that equality, combines the two successive measurable pushforwards, unfolds the concrete product-sampler map, and rewrites the proved recursion identity. The conclusion identifies the numerical position law with the position marginal of the Markov iterate initialized at \(\kappa\otimes\nu\). The only source-law premise concerns the input stream; canonical Gaussian modules prove it from their constructed Brownian source. No output-law agreement is assumed.



Push the verified joint input law through the measurable recursion and then project position. This gives equality with the iterated independent-Gaussian KLMC transition law. The coupling-independent output appearing in the analytic theorem is therefore the Markov sampler output.

\section{\texorpdfstring{Construct the correlated innovations from the same Brownian source}{Construct the correlated innovations from the same Brownian source}}
\label{proof-module-KineticLangevin-BrownianConvolution-lean}


The exponential schemes use correlated Brownian integrals, for example \(\xi_v=\int_0^t e^{-\gamma(t-s)}\,dB_s\) and \(\xi_x=\int_0^t\psi_1(t-s)\,dB_s\). Drawing two independent Gaussians would give the wrong algorithm. This section constructs these random variables on the explicit Gaussian probability space already used for the canonical Brownian motion. Their joint Gaussian law, all covariances, and continuity in L\ensuremath{{}^{2}} are proved from the deterministic time kernels. The construction uses the existing isonormal isometry, which is the deterministic-kernel Wiener integral, rather than adding a hypothesis asserting that an unspecified noise has the required law.



One isonormal source simultaneously realizes all response kernels. This gives exact covariance, temporal continuity, primitive identities, and the common-noise equations needed by both continuous and numerical kinetic constructions.

\subsection{\texorpdfstring{Reuse the existing Brownian probability space}{Reuse the existing Brownian probability space}}


The time Hilbert space is \(H=L^2(\mathbb R,ds)\). The existing canonical Brownian construction selects a Hilbert basis of this space and places independent standard Gaussian coordinates on an explicit product probability space. Its isonormal map \(I:H\to L^2(\Omega,\mathbb P)\) is linear and preserves inner products. Brownian motion is represented there by \(B_t=I(1_{(0,t]})\). The three abbreviations reuse exactly this Hilbert space, sample space, and probability measure. No second independent Brownian source is introduced.



An isonormal Gaussian map assigns a centered Gaussian random variable to every deterministic square-integrable time function. Its covariance is the time-function inner product. This single source realizes all response integrals together and makes their cross covariances available automatically.

\subsection{\texorpdfstring{Turn a continuous response into a L\ensuremath{{}^{2}} time kernel}{Turn a continuous response into a L² time kernel}}


For a continuous scalar response \(k\), define \(K_{k,t}(s)=1_{(0,t]}(s)k(t-s)\). The time reversal is essential: a Brownian fluctuation arriving at time \(s\) has \(t-s\) units of damping before the endpoint. Continuity of \(k\) gives a finite bound on the compact interval \([0,t]\). Since \((0,t]\) has finite Lebesgue measure, this bound proves square integrability.



First, the argument uses square integrability on the kernel support to work on the finite interval measure. boundedness of a continuous function on a compact interval provides an bound, and square integrability of bounded variables proves membership in L\ensuremath{{}^{2}}. the square-integrable time kernel then uses the square-integrable equivalence class to form the equivalence class. Its selected representative agrees almost everywhere with the explicit indicator function, recorded by the pointwise time-kernel representative. Subsequent pointwise calculations therefore apply to the L\ensuremath{{}^{2}} object.



The response kernel is supported on the finite interval preceding the endpoint. Its boundedness and continuity give a finite square integral there. Extending it by zero creates a element of the time Hilbert space to which the Gaussian map can be applied.

\subsection{\texorpdfstring{Compute the time-\allowbreak{}kernel inner product exactly}{Compute the time-kernel inner product exactly}}


The fundamental identity is

\[\langle K_{k,t},K_{\ell,t}\rangle_H=\int_0^t k(s)\ell(s)\,ds.\]

Start with the definition of the L\ensuremath{{}^{2}} inner product as an integral. Replace both chosen representatives using their almost-everywhere formulas. Their product is zero outside \((0,t]\), and inside it equals \(k(t-s)\ell(t-s)\). Integrating the indicator restricts the integral to that interval. The substitution \(u=t-s\) reverses the endpoints and gives the displayed expression. These steps follow with the integral formula for the square-integrable inner product, equality of integrals under almost-everywhere equality, integration over the support, and the reflected-time substitution. This identity is the source of the covariance matrix, including its off-diagonal entries.



The covariance of two responses is the overlap integral of their kernels. Reflection about the endpoint changes this into the response Gram integral already used to define the sampler covariance. Thus the probabilistic construction and the numerical covariance agree entry by entry.

\subsection{\texorpdfstring{Prove time continuity in L\ensuremath{{}^{2}}, including the moving endpoint}{Prove time continuity in L², including the moving endpoint}}


A moving indicator has a discontinuity at its endpoint, so pointwise continuity at every time coordinate is unavailable. Fix an endpoint \(t\). For every \(r\ne t\), the membership of \(r\) in \((0,s]\) is constant for all \(s\) sufficiently close to \(t\). If \(r<t\), ordinary continuity of \(k(s-r)\) applies; if \(r>t\), the kernel stays zero. The single exceptional time coordinate has Lebesgue measure zero.



The exact norm formula is

\[\|K_{k,s}-K_{k,t}\|_H^2=\int_{\mathbb R}|K_{k,s}(r)-K_{k,t}(r)|^2\,dr.\]

For \(s\) near \(t\), both kernels are supported in the fixed finite interval \((0,t+1]\). Compactness gives a common bound \(C\) for their values. The squared difference is therefore bounded by the integrable function \(4C^2 1_{(0,t+1]}\). Dominated convergence makes the squared L\ensuremath{{}^{2}} distance tend to zero; continuity of the square root gives convergence of the distance itself.



One proves the pointwise limit by splitting the cases \(r<t\), \(r>t\), and \(r\le0\), with eventual inequalities supplied by the neighborhood filter. It obtains the norm formula from the L\ensuremath{{}^{2}} inner product and the almost-everywhere subtraction formula. The continuity proof then applies dominated convergence, explicitly supplies the compact bound and the null singleton exception, and finishes with continuity of the square root. This proves continuity as an L\ensuremath{{}^{2}}-valued map; it does not assert that the selected representatives have continuous sample paths.



Changing the endpoint changes both the kernel values and their support. On a common bounded interval the responses admit an integrable dominating bound. Pointwise convergence away from the moving endpoint and dominated convergence then yield convergence in the time Hilbert norm.

\subsection{\texorpdfstring{Apply the isonormal map and derive Gaussian moments}{Apply the isonormal map and derive Gaussian moments}}


Define \(Z_k(t)=I(K_{k,t})\). This is the Wiener integral \(\int_0^t k(t-s)\,dB_s\) for a deterministic continuous kernel. The existing construction proves that every output of \(I\) is centered Gaussian and that \(\operatorname{Cov}(Iu,Iv)=\langle u,v\rangle_H\). Consequently,

\[\mathbb E Z_k(t)=0,\qquad \operatorname{Cov}(Z_k(t),Z_\ell(t))=\int_0^t k(s)\ell(s)\,ds.\]

The same linear map proves joint Gaussianity of every finite collection. This last statement is stronger than proving Gaussianity of each coordinate separately: all linear combinations must be Gaussian. The argument invokes the previously proved the joint Gaussian law of linear responses, whose proof uses finite linear combinations and the linearity of the isonormal map. Covariance follows by rewriting the isonormal covariance theorem and the time-kernel inner-product identity. Time continuity follows by composing the continuous isometry with continuity of the response in the time Hilbert norm.



The Gaussian map is an isometry in second mean. It therefore sends continuity in the deterministic time-kernel norm to continuity in the random-variable norm. Its covariance identity gives the exact second moments, while linearity establishes joint Gaussianity of any finite response family.

\subsection{\texorpdfstring{Identify the exact multivariate innovation law}{Identify the exact multivariate innovation law}}


For finitely many kernels \((k_i)\), assemble \(Z=(Z_{k_i}(t))_i\) as a Euclidean random vector. Its law is proved equal to \(\mathcal N(0,C)\), where \(C_{ij}=\int_0^t k_i(s)k_j(s)\,ds\). The conclusion compares the pushforward probability measure with the centered Gaussian response law; equality of laws is not a premise.



First, the argument transports the jointly Gaussian coordinate family through the continuous linear equivalence from a finite function space to Euclidean space. Coordinate L\ensuremath{{}^{2}} membership gives vector L\ensuremath{{}^{2}} membership, hence integrability. To identify the mean, the continuous coordinate projections commute with the Bochner integral and reduce every coordinate to the already proved zero mean. To identify covariance, the coordinate covariance formula expands the covariance bilinear form as a finite double sum of coordinate covariances. The Brownian covariance identity turns that double sum into the matrix expression for the Gram covariance. uniqueness of a Gaussian law from its mean and covariance then concludes equality of the two Gaussian probability measures from equality of their means and covariance bilinear forms.



For a finite collection of responses, every real linear combination is another isonormal image. Hence the collection is jointly Gaussian. Its mean vector is zero and its covariance is the Gram matrix, which identifies its entire distribution rather than only separate coordinate moments.

\subsection{\texorpdfstring{Specialize to both kinetic schemes and check zero friction}{Specialize to both kinetic schemes and check zero friction}}


The first-order velocity and position noises use \(k=\psi_0\) and \(k=\psi_1\), respectively. Their covariance is exactly the velocity–position response covariance, which generally is nonzero. The two-dimensional law theorem supplies the full first-order KLMC innovation, and the four-dimensional theorem uses \((\psi_0,\psi_1,\varphi_2,\varphi_3)\) to supply every correlated innovation needed by the Hessian-corrected second-order scheme. Both are direct specializations of the finite-family joint-law theorem, so each uses one common Brownian motion.



At zero friction, \(\psi_0\equiv1\), so the velocity convolution must reduce to the original Brownian variable \(B_t\). The argument verifies this as an equality in L\ensuremath{{}^{2}}: the time kernel is precisely \(1_{(0,t]}\), the same indicator used in the canonical Brownian construction. The proof compares the almost-everywhere representatives with equality almost everywhere and the indicator-kernel representative identity. The position and velocity L\ensuremath{{}^{2}}-continuity statements follow from the general continuous-kernel theorem.



The first-order and second-order schemes are obtained by choosing their respective kernel families. When friction is zero the kernels become polynomials and the same covariance construction still applies. No singular quotient formula is needed to identify the limiting innovations.

\subsection{\texorpdfstring{Justify stochastic Fubini in the deterministic Hilbert space}{Justify stochastic Fubini in the deterministic Hilbert space}}


The essential primitive identity is

\[K_{\psi_1,t}=\int_0^t K_{\psi_0,s}\,ds\quad\text{in }L^2(\mathbb R,dr).\]

Point evaluation on L\ensuremath{{}^{2}} is not continuous, so one cannot simply evaluate a Bochner integral at a time coordinate and interchange the two integrals. The proof instead tests both sides against an arbitrary L\ensuremath{{}^{2}} function \(f\). The inner product with a kernel is the ordinary integral \(\int_0^s f(r)e^{-\gamma(s-r)}\,dr\), a fact obtained from the almost-everywhere kernel representatives.



On the finite square \((0,t]^2\), consider \(1_{\{r\le s\}}e^{-\gamma(s-r)}f(r)\). For nonnegative friction the triangular exponential weight is bounded by one. The restriction of an L\ensuremath{{}^{2}} function to a finite-measure interval is integrable, so the product is integrable on the square. This establishes the hypothesis of Fubini's theorem; the interchange is proved rather than assumed. After swapping, the inner integral is

\[\int_r^t e^{-\gamma(s-r)}\,ds=\psi_1(t-r).\]

Thus the two inner products agree for every \(f\), and Hilbert-space extensionality gives the vector identity.



One proves the two triangular-domain reductions by converting indicator integrals to restricted measures and explicitly identifying the intersections \((0,t]\cap(-\infty,s]=(0,s]\) and \((0,t]\cap[r,\infty)=[r,t]\). It removes the endpoint distinction using the null singleton property of Lebesgue measure. integrability on an independent product extension and integrability under bounded multiplication establish product integrability. Fubini’s theorem applies Fubini, and equality of all inner products finishes after commuting the continuous inner-product functional with the Bochner integral.



The triangular integration region has finite measure and the response kernels satisfy the required norm bound there. Fubini may therefore interchange the time primitive and the isonormal response. This is the rigorous origin of the positional innovation as an integrated velocity innovation.

\subsection{\texorpdfstring{Obtain the kinetic Brownian integral equation}{Obtain the kinetic Brownian integral equation}}


The scalar coefficient identity \(\psi_0(u)+\gamma\psi_1(u)=1\) holds pointwise, including at zero friction. Multiplying it by the time indicator and applying linearity of the isonormal map gives the exact relation

\[V_t^{\mathrm{noise}}+\gamma X_t^{\mathrm{noise}}=B_t.\]

The deterministic Hilbert primitive identity just proved transfers through the continuous linear isonormal map. Consequently,

\[X_t^{\mathrm{noise}}=\int_0^tV_s^{\mathrm{noise}}\,ds,\qquad V_t^{\mathrm{noise}}=B_t-\gamma\int_0^tV_s^{\mathrm{noise}}\,ds.\]

These are equations for the L\ensuremath{{}^{2}} random variables whose joint Gaussian innovation law has already been identified. They connect the numerical exponential-integrator noise to the same additive Brownian forcing used by the kinetic diffusion.



For the first identity The proof compares representatives of the sum and scalar multiple of the time kernels, then applies the first friction-response balance. For the primitive identity, commutation of a bounded linear map with integration transports the deterministic Bochner integral through the isonormal map; its integrability is supplied by the proved L\ensuremath{{}^{2}} time continuity. Finally, ordinary vector-space rearrangement gives the damped Brownian integral equation. No matching-noise law or integral-solution equation is introduced as a premise.



The exponential response \(B_t\) satisfies \(B_t=W_t-\gamma\int_0^tB_sds\). Multiplying by the physical amplitude gives the filtered forcing in the velocity equation. Its time primitive supplies position, so the exact and numerical equations use one common Brownian source.

\subsection{\texorpdfstring{Derive the explicit Brownian convolution by an integrating factor}{Derive the explicit Brownian convolution by an integrating factor}}


The proved primitive identity implies \(P'_t=V_t\) for positive time. Combined with \(V_t+\gamma P_t=B_t\), this is the ordinary L\ensuremath{{}^{2}}-valued equation \(P'_t+\gamma P_t=B_t\), with \(P_0=0\). Its integrating factor gives

\[\frac{d}{dt}\bigl(e^{\gamma t}P_t\bigr)=e^{\gamma t}B_t,\qquad P_t=\int_0^t e^{-\gamma(t-s)}B_s\,ds.\]

Substitution into the exact Brownian balance then gives the integration-by-parts representation

\[V_t=B_t-\gamma\int_0^t e^{-\gamma(t-s)}B_s\,ds.\]

Thus the isonormal Wiener integral and the continuous additive-Brownian forcing formulation are proved equal as L\ensuremath{{}^{2}} vectors.



First, the argument differentiates the indefinite integral defining the primitive. On a neighborhood of a positive time, coercing through the nonnegative real value does not change the endpoint, so locality of differentiation transfers the derivative to the square-integrable position. The product rule differentiates the scalar exponential times the L\ensuremath{{}^{2}} vector. The Brownian balance cancels the damping term. The fundamental theorem of calculus on \([0,t]\), requiring derivatives only in its interior, then integrates this derivative. Multiplication by \(e^{-\gamma t}\) cancels the integrating factor. Finally, scalar multiplication commutes with the interval integral, and the exponential addition law converts \(e^{-\gamma t}e^{\gamma s}\) into \(e^{-\gamma(t-s)}\). This proof uses no sample-path differentiability of Brownian motion.



Integration by parts writes the convolution as a Brownian increment minus an ordinary convolution of the continuous Brownian path. This representation avoids differentiating Brownian paths and gives a pathwise expression compatible with the continuous-input solution construction.

\section{\texorpdfstring{Construct independent correlated innovations on Brownian time intervals}{Construct independent correlated innovations on Brownian time intervals}}
\label{proof-module-KineticLangevin-BrownianSegments-lean}


The one-step Gaussian covariance is only part of the algorithm: innovations at different iterations must be independent and must come from increments of the same Brownian source used by the exact reference. This section translates deterministic time kernels onto successive intervals of that source. Translation preserves the within-step Gram covariance. Disjoint time supports give zero cross covariances, and joint Gaussianity turns this orthogonality into independence of the entire innovation blocks. Correlations among coordinates within a block are retained.



Translated response kernels realize successive numerical steps on one Brownian source. Disjoint support gives zero covariance, and joint Gaussianity upgrades that orthogonality to independence of entire response blocks.

\subsection{\texorpdfstring{Translate a time kernel without changing its L\ensuremath{{}^{2}} geometry}{Translate a time kernel without changing its L² geometry}}


For an interval beginning at \(a\), translate the original response kernel by \((S_aK)(r)=K(r-a)\). Lebesgue measure is invariant under translation, so \(S_a\) is a real linear isometry of the deterministic time Hilbert space. The translated response is

\[K^{a,t}_k(r)=1_{(a,a+t]}(r)k(a+t-r).\]

Its inner product with another response on the same interval remains \(\int_0^t k(s)\ell(s)\,ds\).



The argument constructs the isometry with the isometric change of probability source and translation invariance of time measure. The representative theorem combines the almost-everywhere composition formula with the previously proved kernel representative. Elementary inequalities identify \(r-a\in(0,t]\) with \(r\in(a,a+t]\); scalar algebra identifies the two elapsed-time expressions. Inner-product preservation of the linear isometry then reduces the covariance geometry to the unshifted kernel identity.



Translation of a deterministic time kernel preserves its square integral. The same response covariance is therefore obtained on every interval of the same length. This identifies numerical time steps with translated pieces of one underlying Brownian source.

\subsection{\texorpdfstring{Prove that disjoint time segments are orthogonal}{Prove that disjoint time segments are orthogonal}}


If \(a+t\le b\), the intervals \((a,a+t]\) and \((b,b+u]\) are disjoint. At every time coordinate, at least one of their indicators is zero. Their L\ensuremath{{}^{2}} inner product is therefore zero, regardless of the two continuous response functions.



The proof unfolds the L\ensuremath{{}^{2}} inner product as an integral, substitutes the explicit almost-everywhere segment representatives, and splits on membership in the first interval. Membership implies \(r\le a+t\le b\), which excludes the strict lower bound \(b<r\) of the second interval. Thus the integrand vanishes almost everywhere, and the vanishing integral of an almost-everywhere zero function concludes. The half-open convention makes adjacent segments disjoint exactly, including their common endpoint.



Kernels supported on disjoint time intervals have zero product almost everywhere. Their Hilbert-space inner product vanishes, including when intervals meet at a single endpoint. That endpoint has zero time measure and contributes no covariance.

\subsection{\texorpdfstring{Apply the same isonormal map to every time segment}{Apply the same isonormal map to every time segment}}


Define \(Z_k^{a,t}=I(K_k^{a,t})\) using the same isonormal map and Gaussian product sample space as the original Brownian motion. Each response is centered, and every finite family of responses on an interval is jointly Gaussian. The covariance on one interval is the same Gram integral as before. Responses on separated intervals have covariance zero by the deterministic support calculation.



One proves these assertions by composing the translation's linear map with the existing Gaussian isonormal map. The finite-family Gaussian theorem applies to that composite linear map. For covariances, the isonormal covariance identity reduces the probability calculation to the deterministic inner product, where the within-interval and disjoint-interval theorems apply. Independence is not assumed at this stage; zero covariance alone would not imply independence without the jointly Gaussian property proved next.



All translated kernels are fed into the same isonormal map. This makes them jointly Gaussian, not merely separately Gaussian. Orthogonality of the time kernels can therefore be converted into independence of the response blocks.

\subsection{\texorpdfstring{Identify the full Gaussian block law on every interval}{Identify the full Gaussian block law on every interval}}


For finitely many response functions \((k_i)\), the Euclidean vector \((Z_{k_i}^{a,t})_i\) has exactly the centered Gaussian law with covariance \(C_{ij}=\int_0^t k_i(s)k_j(s)\,ds\). The starting time \(a\) does not appear in this law, because deterministic translation preserves all inner products.



The proof assembles the L\ensuremath{{}^{2}} representatives into a Euclidean random vector and proves its Gaussianity via the continuous linear equivalence from finite function space. Coordinate projections commute with the Bochner integral, giving the zero vector mean. The covariance bilinear form expands into a finite double sum; each entry is rewritten by the proved segment covariance formula. Equality with the Gram Gaussian's bilinear form, together with equality of means, lets uniqueness of a Gaussian law from its mean and covariance identify the two probability measures. No distributional equality is supplied as a hypothesis.



Within one interval, kernel inner products reproduce the full correlated response Gram matrix. Across spatial coordinates the construction uses independent Brownian coordinates. The resulting vector block has precisely the noise law specified by the exponential integrator.

\subsection{\texorpdfstring{Derive independence of all successive Brownian innovation blocks}{Derive independence of all successive Brownian innovation blocks}}


On the uniform grid, the \(n\)-th block uses the interval \((nh,(n+1)h]\). the grid convolution places every response function on that interval. The family indexed by both iteration and response coordinate is a Gaussian process because every member is obtained by applying one Gaussian linear map to a deterministic Hilbert vector.



For \(n<m\), the elementary inequality \((n+1)h\le mh\) proves disjoint support and therefore zero cross covariance between every coordinate of the two blocks. The Gaussian-process independence theorem converts these coordinate covariance equalities into independence of the entire block-valued random variables, over all iteration numbers simultaneously. Packing a finite coordinate family into Euclidean space is a measurable map, so time independence survives that packing.



The argument constructs the jointly Gaussian process with Gaussianity of linear response families. It uses independence of uncorrelated jointly Gaussian blocks, proves the disjoint-segment inequality from the successor-index inequality, and handles the reversed ordering through symmetry of covariance. Finally, preservation of independence under separate measurable maps transports independence through the Euclidean identification. The per-block law is the segment joint-law theorem, so the Brownian-derived Euclidean blocks are identically distributed with the required correlated Gram-Gaussian law.



For any finite collection of disjoint intervals, the jointly Gaussian covariance matrix is block diagonal. Its characteristic function factors, proving independence of those blocks. Consistency over finite collections gives the independent innovation stream needed for the full numerical recursion.

\section{\texorpdfstring{The exact Brownian equations on every numerical time interval}{The exact Brownian equations on every numerical time interval}}
\label{proof-module-KineticLangevin-BrownianSegmentEquations-lean}


The numerical method uses fresh correlated noise on each interval \((a,a+t]\). The preceding module constructs those noises on the same canonical Brownian source and proves independence across disjoint intervals. Here the noises also satisfy the integrated kinetic equations, so their relationship to the continuous diffusion is an equality in L\ensuremath{{}^{2}} rather than a matching-law assumption.



The local filtered responses are exact translated pieces of the global source. Their damped and primitive equations ensure every numerical interval uses the same forcing as its corresponding exact segment.

\subsection{\texorpdfstring{Identify the Brownian increment on a translated interval}{Identify the Brownian increment on a translated interval}}


Define \(W^a_t=B_{a+t}-B_a\). The unit response kernel on \((a,a+t]\) is exactly the difference \(1_{(0,a+t]}-1_{(0,a]}\). One proves the indicator identity almost everywhere using equality almost everywhere, the explicit representatives, and real order arithmetic. Linearity of the isonormal Wiener-integral map then identifies its image with the Brownian increment. This fixes the driving noise as a constructed object on the original probability space.



The increment over a translated interval is represented by the indicator of that interval in the deterministic time Hilbert space. Linearity identifies it with the difference of the two Brownian endpoint values. This connects the local response to the original global process.

\subsection{\texorpdfstring{Integrate the shifted velocity response}{Integrate the shifted velocity response}}


Write \(V^a_t=\int_a^{a+t}e^{-\gamma(a+t-r)}\,dB_r\) and \(P^a_t=\int_a^{a+t}\psi_1(a+t-r)\,dB_r\). Translation of deterministic time kernels is a continuous linear isometry. Therefore continuity of the unshifted kernels gives continuity of each translated stochastic response in L\ensuremath{{}^{2}}. The already proved deterministic primitive identity \(K_{\psi_1}(t)=\int_0^tK_{\psi_0}(s)ds\) survives translation and application of the isonormal map. The proof composes their continuous linear maps and invokes commutation of a bounded linear map with integration, proving \(P^a_t=\int_0^tV^a_sds\) without pointwise evaluation of an L\ensuremath{{}^{2}} class.



Integrating the shifted velocity response creates the corresponding position response because both use the same translated time support. The equality is an identity in the square-integrable Hilbert space. It retains the correlation needed for exact local noise cancellation.

\subsection{\texorpdfstring{Prove the damped Brownian integral equation}{Prove the damped Brownian integral equation}}


The exact coefficient identity \(\psi_0+\gamma\psi_1=1\) gives \(V^a_t+\gamma P^a_t=W^a_t\). First, the argument uses injectivity of the isonormal linear isometry to recover the deterministic kernel identity from the unshifted balance, then applies time translation and the same map. Substituting the primitive gives \(V^a_t=W^a_t-\gamma\int_0^tV^a_sds\). This equality supplies the additive forcing needed when a continuous reference trajectory and a numerical step share one Brownian interval.



The translated filtered response satisfies the same damped integral equation as the response starting at zero. Only the global time origin has changed. This allows every numerical interval to reuse the same kinetic local-error argument with its own fresh Brownian increment.

\subsection{\texorpdfstring{Recover the ordinary Brownian convolution formula}{Recover the ordinary Brownian convolution formula}}


The position response also equals \(P^a_t=\int_0^te^{-\gamma(t-s)}W^a_sds\), hence \(V^a_t=W^a_t-\gamma\int_0^te^{-\gamma(t-s)}W^a_sds\). To prove this for the exact translated construction, The proof pulls the established unshifted formula back through the injective isonormal isometry, then sends the deterministic identity through the translated continuous linear map. Continuity ensures Bochner integrability. Commuting that map with the integral and its scalar factors yields the formula; the velocity version follows by the balance identity. These are exact integration-by-parts identities on the constructed Brownian source.



Integration by parts converts the filtered response into an ordinary integral of the Brownian increment path. The formula uses only path continuity. It therefore interfaces directly with the deterministic continuous-input integral equation used to construct exact kinetic segments.

\section{\texorpdfstring{The canonical vector Brownian innovations and initial Gaussian velocity}{The canonical vector Brownian innovations and initial Gaussian velocity}}
\label{proof-module-KineticLangevin-CanonicalVectorSegments-lean}


The paper draws correlated velocity and position innovations in each spatial coordinate, independently across coordinates and time steps. Here every such draw is realized from the existing product of scalar canonical Brownian sources. The entire constructed array is proved to have exactly the previously defined Gaussian product law; equality of a few moments alone is not used as a substitute for equality of laws.



The scalar Brownian responses are assembled into finite-dimensional vector innovations with their complete joint stream law. Independent Gaussian initialization is adjoined once, independently of the entire future source.

\subsection{\texorpdfstring{Transport independence along an law}{Transport independence along an law}}


A measurable map \(Y\) with law \(Q\) can pull back any independent measurable family defined on \(Q\). One proves this by the characterization of mutual independence through equality with the infinite product of marginal laws. It applies the pushforward composition theorem to the joint family and to each marginal. This elementary tool is used repeatedly to transfer established independence to the concrete Brownian source.



If a random vector has a product law, its coordinate projections are independent with the stated marginals. Pushing this law through separate measurable coordinate maps preserves independence. This law-based route allows independence to be transferred to the canonical response variables.

\subsection{\texorpdfstring{Build correlated scalar and spatial blocks}{Build correlated scalar and spatial blocks}}


A scalar block consists of all deterministic response kernels applied to the same Brownian interval. Packing their values with the Euclidean identification gives a finite Euclidean vector whose law is the Gram Gaussian already proved in the translated Brownian interval construction. A vector block evaluates that scalar block on each spatial coordinate source. Every map is measurable because each L\ensuremath{{}^{2}} representative is strongly measurable, coordinate evaluation is measurable, and finite or infinite coordinatewise packing preserves measurability.



At each time interval and spatial coordinate, retain the entire correlated response block. Assembling like response coordinates into Euclidean vectors gives the velocity and position innovations. Separate spatial coordinates remain independent, while the within-coordinate covariance is unchanged.

\subsection{\texorpdfstring{Prove the law of the entire array}{Prove the law of the entire array}}


For one spatial source, disjoint Brownian intervals give independent scalar blocks. Different spatial sources are independent by the explicit product construction. The theorem regrouping independent coordinates combines these two facts into independence over iteration-coordinate pairs; reindexing swaps their order. Each marginal has the established Gram-Gaussian law, so the law of the entire independent stream identifies the full array distribution. Currying the array is a measurable equivalence, and the product-measure curry theorem identifies the distribution as an independent stream of complete spatial blocks. Pullback along this path law proves independence across time steps.



Finite-dimensional cylinder distributions determine the law of the countable innovation array. Each cylinder is the product of its independent block laws. Consequently the whole array has the independent-product stream law required by Markov iteration.

\subsection{\texorpdfstring{Identify the existing KLMC Gaussian stream exactly}{Identify the existing KLMC Gaussian stream exactly}}


For \(d\) finite spatial coordinates, the pair-indexed array has exactly the existing the filtered noise path law. Specializing the response family to \((\psi_0,\psi_1)\) and packing each response into a vector produces the unscaled velocity-position innovation. Composition of the stream law with the existing the first-order Gaussian innovation map proves its exact distribution. Independence of the full vector blocks survives this measurable packing. Each vector coordinate lies in L\ensuremath{{}^{2}} by the Brownian isometry; pulling it back along spatial evaluation and using square integrability of a coordinate proves both phase components have finite second moments.



The Brownian response pair and the numerical Gaussian block have the same centered joint Gaussian covariance. Equality of their laws identifies them as two realizations of the same sampler innovation. The construction thereby connects the source process to the prescribed transition measure.

\subsection{\texorpdfstring{Adjoin the independent initial Gaussian velocity}{Adjoin the independent initial Gaussian velocity}}


The source for a sampler is \(Q=\mathcal N(0,I_d)\otimes P_{\mathrm{Brownian}}\). Its first projection is the standard Gaussian initial velocity. Its second projection carries all constructed innovations. Product-measure projection laws prove the velocity law and preserve the innovation law. The product independence theorem proves that initial velocity is independent of the entire future innovation stream, while measure-preserving pullback gives all required L\ensuremath{{}^{2}} bounds. An arbitrary initial-position coupling can subsequently be extended by this same fixed source, independently of that coupling.



Pullback along the second projection also proves mutual independence of complete innovations across all iterations on this enlarged source; the new Gaussian velocity does not alter the noise law.



An additional independent standard Gaussian vector supplies the initial velocity. Taking its product with the entire Brownian source proves independence from every future innovation simultaneously. This is stronger than checking independence from only the first numerical interval.

\section{\texorpdfstring{Continuous vector Brownian responses in L\ensuremath{{}^{2}}}{Continuous vector Brownian responses in L²}}
\label{proof-module-KineticLangevin-VectorSegmentEquations-lean}


The scalar Brownian equations must also hold for Euclidean random vectors so they can drive the L\ensuremath{{}^{2}} kinetic diffusion used by the contraction proof. A continuous linear spatial assembly map achieves this directly: it preserves sums, scalar multiples, continuity, and Bochner time integrals.



Bounded linear spatial assembly transfers scalar response identities to vector-valued second-mean paths. The resulting equations identify raw inputs, filtered inputs, and numerical endpoints on the same source.

\subsection{\texorpdfstring{Assemble scalar L\ensuremath{{}^{2}} variables through a continuous linear map}{Assemble scalar L² variables through a continuous linear map}}


For a scalar random variable \(Z\), spatial assembly produces the vector \((Z(\omega_i))_i\), evaluating it on independent spatial copies of the canonical source. The proof definition is a finite sum. Each summand first pulls back along a measure-preserving coordinate evaluation, then inserts the scalar into one Euclidean basis coordinate. Both operations are continuous linear maps on L\ensuremath{{}^{2}}. The representative theorem is proved by the finite-sum almost-everywhere identity and the explicit representatives of these two maps; Euclidean extensionality verifies each coordinate. This construction avoids interchanging pointwise evaluation with an L\ensuremath{{}^{2}} Bochner integral.



A finite family of scalar square-integrable variables can be assembled into one Euclidean square-integrable vector. The assembly map is bounded and linear. It therefore commutes with limits, time integrals, and the affine operations used by the kinetic equations.

\subsection{\texorpdfstring{Define vector responses and prove continuity}{Define vector responses and prove continuity}}


The vector response is the spatial assembly of the scalar response on \((a,a+t]\). The vector Brownian input is the assembly of the scalar Brownian increment. Their definitions remain on the existing canonical product source. The response representative theorem gives exactly the expected independent-coordinate vector. Continuity in time follows by composing the continuous scalar response with the continuous linear assembly map.



Each spatial response is continuous in the scalar second-mean norm. Finite-dimensional assembly transfers that continuity to the vector response. Its squared norm is the sum of the coordinate second moments, so the required bounds remain explicit in dimension.

\subsection{\texorpdfstring{Lift the exact primitive and damped integral equations}{Lift the exact primitive and damped integral equations}}


The position response satisfies \(P^a_t=\int_0^tV^a_sds\), and the velocity response satisfies \(V^a_t=W^a_t-\gamma\int_0^tV^a_sds\). These equalities follow from the already proved scalar equations by applying the spatial assembly map. The proof commutes that continuous linear map with each Bochner integral using commutation of a bounded linear map with integration. Thus the exact filtered Brownian input required by the generic mild-equation theorem is constructed for the vector L\ensuremath{{}^{2}} diffusion.



Apply the assembly map to the scalar primitive and damped integral identities. Bounded linearity permits it to pass through the Bochner integrals. The vector equations follow componentwise with the exact same friction and response coefficients.

\subsection{\texorpdfstring{Identify the input with canonical vector Brownian increments}{Identify the input with canonical vector Brownian increments}}


The lifted input is proved almost everywhere equal to the vector whose coordinates are \(B_i(a+t)-B_i(a)\) for the existing canonical coordinate process. The scalar subtraction identity is pulled back through each measure-preserving spatial projection and combined with the spatial representative theorem. Continuity follows by identifying the increment with the unit-kernel vector response. This connects the L\ensuremath{{}^{2}} forcing to the canonical vector Brownian construction.



The zero-response kernel is the indicator of the local Brownian interval. Its assembled isonormal image is therefore the vector Brownian increment. This verifies that the local input used in the kinetic equation comes from the canonical vector process.

\subsection{\texorpdfstring{Identify the OU increment and numerical innovation endpoints}{Identify the OU increment and numerical innovation endpoints}}


The vector velocity response is exactly \(V^a_t=W^a_t-\gamma\int_0^te^{-\gamma(t-s)}W^a_sds\). The argument applies spatial assembly to the established scalar integration-by-parts formula, uses its linearity, and commutes it with the continuous Bochner integrand. This is precisely the OU input increment used by the kinetic splitting limit. At each grid interval, the paired L\ensuremath{{}^{2}} velocity and position response endpoints are proved almost everywhere equal to the pointwise the canonical first-order innovation already equipped with the exact Gaussian law. The endpoint proof combines the two representative equalities with equality of the two phase coordinates; it does not replace the constructed inputs by new independent Gaussian variables.



At the step endpoint the vector responses equal the numerical innovations. The filtered velocity response is also the Ornstein–Uhlenbeck increment used by the splitting. These identities keep the continuous and numerical couplings tied to one underlying source.

\section{\texorpdfstring{Splitting convergence on the constructed vector Brownian source}{Splitting convergence on the constructed vector Brownian source}}
\label{proof-module-KineticLangevin-CanonicalSplitting-lean}


The deterministic input in the splitting theorem is now identified with the canonical vector Brownian motion in L\ensuremath{{}^{2}}. Its filtered increments are exactly the constructed Gaussian segment convolutions, including the fluctuation-dissipation scaling. Consequently the innovation-driven numerical chain and the continuous-input chain coincide on the same source, and the L\ensuremath{{}^{2}} chain converges to its kinetic integral solution.



The continuous Brownian input, filtered increments, and full thermal stream are identified on one canonical probability source. Applying deterministic splitting convergence there gives an second-mean stochastic limit.

\subsection{\texorpdfstring{Keep scaling and source lifts exact}{Keep scaling and source lifts exact}}


The OU input construction is linear in the additive input. Multiplying \(B\) by a real coefficient \(c\) multiplies both its matching increment and its exponential convolution by \(c\). More generally, a continuous linear map \(L\) commutes with this construction: it commutes with subtraction, scalar multiplication, and the Bochner interval integral of the continuous integrand. One proves these exact identities by integral linearity and continuous-linear integral commutation. The latter applies, in particular, to measure-preserving pullbacks that extend the Brownian source by an independent initial state.



The Brownian source lift includes its physical amplitude before any numerical comparison is made. Since source pullback is linear, it preserves every response identity with that amplitude. This keeps thermal normalization consistent across the continuous equation and splitting increments.

\subsection{\texorpdfstring{Construct the continuous L\ensuremath{{}^{2}} Brownian input}{Construct the continuous L² Brownian input}}


The input is the already constructed vector Brownian increment from time zero, evaluated at \(t_+=\max\{t,0\}\). This supplies a globally defined real-time function, anchored at zero. The proved continuity of vector Brownian increments in L\ensuremath{{}^{2}} and continuity of \(t\mapsto t_+\) prove continuity of this input. At zero the underlying increment is the difference of a Brownian value with itself, so it is exactly zero. The proof unfolds the source definitions to check this anchor.



Assemble the canonical scalar Brownian coordinates into a vector and lift them onto the product initialization space. Bounded linearity transfers second-mean continuity and the zero initial value. This produces the precise continuous Hilbert-space input to which splitting convergence applies.

\subsection{\texorpdfstring{Identify input increments and OU segment responses}{Identify input increments and OU segment responses}}


The difference between the anchored values at \(a+t\) and \(a\) cancels their common time-zero value, leaving the canonical Brownian increment on \((a,a+t]\). The independently proved segment convolution identity gives \[Z_{a,t}=\Delta B_{a,t}-\gamma\int_0^t e^{-\gamma(t-s)}\Delta B_{a,s}\,ds.\] This is exactly the deterministic the filtered Ornstein–Uhlenbeck increment used in the splitting theorem. The argument rewrites the endpoint increment, compares the interval integrands, and uses nonnegativity of the integration time to remove its real-to-nonnegative-time clamp. Scaling by any \(c\), including \(\sqrt{2\gamma}\), follows from the linearity identity. The equality is of L\ensuremath{{}^{2}} variables, stronger than merely matching marginal Gaussian laws.



The increment of the full Brownian input over a grid interval equals the corresponding segment input. Its filtered response equals the splitting thermostat innovation. These identities are proved from the source kernels, so the discretization uses the same forcing as the limiting solution.

\subsection{\texorpdfstring{Identify the entire canonical innovation chain}{Identify the entire canonical innovation chain}}


On a mesh of nonnegative size \(h\), the continuous-input chain uses the increment \(N_{kh,h}\). The preceding identity replaces this by \(cZ_{kh,h}\) for every grid index. The explicit-innovation recursion therefore agrees exactly with the continuous-input recursion. First, the argument uses the proved equality between the two chain representations and then substitutes the OU increment identity at each nonnegative grid time.



The entire stream has the product of its explicit thermal Gaussian laws. Equality is established for the source variables, including their dependence across response coordinates within an interval. This permits the deterministic splitting sequence to be identified with the intended random chain.

\subsection{\texorpdfstring{Conclude L\ensuremath{{}^{2}} convergence on the canonical source}{Conclude L² convergence on the canonical source}}


Apply the generic Banach splitting theorem to the continuous canonical input \(cB\), its proved zero initial value, and the continuous kinetic integral solution \(Y\) driven by that same input. Rewrite the resulting continuous-input chain by the proved canonical innovation-chain identity. This yields convergence of the explicit chain with Gaussian convolution innovations on the mesh \(T/(n+1)\) to \(Y_T\). The argument checks the nonnegative-time casts in the mesh and grid endpoints. The integral-solution premise is the equation constructed elsewhere by Picard existence; no stationarity or numerical convergence is supplied as a contract.



Apply the pathwise splitting convergence theorem in the Hilbert space of square-integrable variables. All continuity and increment premises have already been verified on the canonical source. The conclusion is second-mean convergence of the sampler splitting to the constructed kinetic solution.

\section{\texorpdfstring{Exact Gaussian invariance of the velocity thermostat}{Exact Gaussian invariance of the velocity thermostat}}
\label{proof-module-KineticLangevin-Thermostat-lean}


The O substep removes velocity variance through damping and restores exactly that variance through independent Gaussian noise. This section proves an equality of probability measures, not only a matching second moment. Invariance of this substep is one component of kinetic stationarity; the Hamiltonian evolution and the full diffusion require separate arguments.



The invariant Gaussian temperature is determined by balancing the variance lost through damping against the variance supplied by fresh independent noise. This exact identity is one of the invariant building blocks for kinetic stationarity.

\subsection{\texorpdfstring{Independent Gaussian linear combinations}{Independent Gaussian linear combinations}}


Let \(G,G\prime\) be independent standard Gaussian vectors in a finite-dimensional real Hilbert space, and suppose \(a^2+b^2=1\). For a continuous linear functional \(L\), the characteristic function of \(aG+bG\prime\) is

\[\exp(-a^2\|L\|^2/2)\exp(-b^2\|L\|^2/2)=\exp(-\|L\|^2/2).\]

Independence is represented by the product of two standard Gaussian measures. The argument constructs the continuous linear map \(A(v,w)=av+bw\), uses the characteristic function of a linear pushforward to pull the test functional back through \(A\), and uses factorization of characteristic functions on independent products to factor its characteristic function. The restrictions of \(L\circ A\) to the two factors are \(aL\) and \(bL\), proved by extensionality. The norm of a scalar multiple supplies the squared coefficients. The relation \(a^2+b^2=1\) then simplifies the exponent, and uniqueness of characteristic functions proves the distributional identity.



For independent centered Gaussian vectors \(V\) and \(G\), the characteristic function of \(aV+bG\) factors into the product of the two Gaussian characteristic functions. Its covariance is \((a^2+b^2)I\). Independence is essential: an unspecified cross covariance would contribute an additional term.

\subsection{\texorpdfstring{Apply the fluctuation-\allowbreak{}dissipation balance}{Apply the fluctuation-dissipation balance}}


For nonnegative friction and time step, put \(\delta=e^{-\gamma h/2}\) and \(s=\sqrt{1-\delta^2}\). The variance under the square root is nonnegative, so \(s^2=1-\delta^2\). Hence \(\delta^2+s^2=1\), and the preceding theorem gives \(\delta G+sG\prime\sim\mathcal N(0,I)\). The argument applies the general linear-combination theorem, rewrites the square root squared using its proved nonnegativity, unfolds the variance, and finishes with a polynomial identity. This is precisely the standard-temperature O half-step used in OBABO.



Take \(a=e^{-\gamma t}\) and \(b=\sqrt{1-e^{-2\gamma t}}\). Their squared sum is one, so a standard Gaussian velocity remains standard Gaussian after thermostat. Position is unchanged, hence the full Gaussian-times-position target is preserved whenever the refresh noise is independent.

\section{\texorpdfstring{The OU innovation law and independent time stream}{The OU innovation law and independent time stream}}
\label{proof-module-KineticLangevin-OUMarginal-lean}


The canonical convolution is an random variable on the constructed Brownian source. This section proves that its fluctuation-dissipation scaling has exactly thermal Gaussian law, and that full spatial vectors on distinct grid intervals are independent. The complete stream therefore has the independent Gaussian product law required by the concrete Markov transition.



The thermal response law is computed from its exponential kernel and assembled across independent spatial coordinates. Disjoint intervals yield the complete independent Gaussian stream required by the splitting chain.

\subsection{\texorpdfstring{Identify independent spatial copies through their scalar law}{Identify independent spatial copies through their scalar law}}


A scalar L\ensuremath{{}^{2}} variable \(Z\) is evaluated on each independent coordinate Brownian source. Each copy has the same proved law \(\nu\), and the coordinate product construction proves their independence. The infinite-product law theorem, specialized to finitely many spatial coordinates, gives the finite product measure \(\nu^{\otimes\iota}\). Transporting this law through the Euclidean coordinate map identifies the law of the coordinate assembly map; its representative theorem supplies the almost-everywhere equality. For \(\nu\) equal to a scaled real standard Gaussian, the product pushforward identity and the definition-independent standard Gaussian theorem identify this law as the same scalar scaling of the Euclidean standard Gaussian.



The scalar response at each spatial index has the computed centered Gaussian law. Independence of the coordinate sources gives their product law. Assembling coordinates then yields a Euclidean Gaussian with a scalar covariance multiple of the identity.

\subsection{\texorpdfstring{Compute the scalar thermal variance}{Compute the scalar thermal variance}}


The scalar segment convolution is the constructed isonormal map applied to its time kernel \(K\), so it has centered Gaussian variance \(\|K\|^2\). With \(c=\sqrt{2\gamma}\), its scaled variance is \(2\gamma\|K\|^2\). The proved kernel inner-product identity gives \(\|K\|^2=\int_0^h e^{-2\gamma s}\,ds\), and the exact fluctuation-decay balance gives \(2\gamma\|K\|^2=1-e^{-2\gamma h}\). This equals \(\sigma_h^2\) for \(\sigma_h=\sqrt{1-e^{-2\gamma h}}\). The argument derives nonnegativity from the damping bound, uses the isonormal Gaussian law, checks the equality of nonnegative-real variances, and rewrites the scaled real Gaussian pushforward. Zero friction is included without division.



The filtered velocity response variance is the square integral of the exponential kernel. After multiplying by \(2\gamma\), it is \(1-e^{-2\gamma h}\). This is the missing variance of a velocity damped for time \(h\).

\subsection{\texorpdfstring{Identify the entire spatial OU Gaussian law}{Identify the entire spatial OU Gaussian law}}


Apply the independent-spatial-copy result to the proved scalar thermal law. Linearity of the coordinate assembly map identifies its scaled scalar input with the scaled vector segment convolution. Thus \(\sqrt{2\gamma}\,Z_{a,h}\) has exactly the law of \(\sqrt{1-e^{-2\gamma h}}\,\xi\), where \(\xi\) is the Euclidean standard Gaussian. This proves the full vector distribution rather than only its covariance.



A centered Euclidean Gaussian is determined by its covariance bilinear form. The independent coordinates give covariance equal to the scalar thermal variance times the identity. Thus the assembled innovation is exactly thermal Gaussian required by thermostat.

\subsection{\texorpdfstring{Prove independence of full-\allowbreak{}vector innovations across time}{Prove independence of full-vector innovations across time}}


The canonical KLMC innovation pairs are already proved independent over disjoint time-grid intervals by their Brownian kernel supports. Project each pair to its velocity component and apply any fixed scalar scaling. The L\ensuremath{{}^{2}} vector-segment representatives agree almost everywhere with these projected canonical innovations. Independence is preserved under these measurable projections and almost-everywhere changes. The argument explicitly combines the representative scalar-multiplication event with the canonical endpoint identity.



Entire vector blocks from different intervals are measurable functions of disjoint collections of Gaussian response variables. Independence of those collections therefore gives independence of the vectors, rather than only pairwise independence of their individual coordinates.

\subsection{\texorpdfstring{Identify the full independent Gaussian stream}{Identify the full independent Gaussian stream}}


Every scaled OU grid innovation has the exact thermal Gaussian law, and the preceding theorem proves independence of all grid vectors. The infinite-product law theorem therefore identifies the entire innovation stream with the product of this Gaussian law over all iterations. Measurability follows from the strongly measurable L\ensuremath{{}^{2}} representative at each countable grid index. This stream identity is the source-level fact used to transfer the concrete Brownian-driven splitting chain to its iterated Markov laws.



The finite-dimensional distributions of the full stream are products of thermal Gaussian vector law. This identifies the stream as independent and identically distributed. It supplies the exact fresh-noise law required when the splitting transition is iterated.

\section{\texorpdfstring{Canonical kinetic innovations on arbitrary finite coordinate sets}{Canonical kinetic innovations on arbitrary finite coordinate sets}}
\label{proof-module-KineticLangevin-GenericCanonicalSource-lean}


Particle states are naturally indexed by a finite disjoint union of parameter and particle coordinates. These constructions work directly for any finite index type, while retaining the same canonical Brownian source, correlated Gaussian innovations and independent Gaussian velocity.



All constructions are indexed by an arbitrary finite coordinate set. This permits direct reuse on particle spaces while preserving Gaussian covariance, response identities, and complete time-stream independence.

\subsection{\texorpdfstring{Construct the independent velocity/\allowbreak{}Brownian source}{Construct the independent velocity/Brownian source}}


The source is the explicit product of a standard Euclidean Gaussian velocity and the existing independent-coordinate Brownian source. Product projection laws prove the velocity distribution and its second moment. The phase innovation packs the two exponential-response coordinates from every spatial Brownian interval. Measurability, temporal independence and L\ensuremath{{}^{2}} bounds follow from the generic spatial-block construction. Pullback along the source product projection gives the corresponding L\ensuremath{{}^{2}} innovation components on the enlarged source.



A finite coordinate set specifies a Euclidean space without requiring a particular indexing convention. The source is the product of its standard Gaussian velocity and coordinate Brownian family. The construction works unchanged under any relabeling of the coordinates.

\subsection{\texorpdfstring{Identify the correlated Gaussian phase law and its complete stream}{Identify the correlated Gaussian phase law and its complete stream}}


The scalar block on each coordinate has the two-dimensional Gram-Gaussian law. The already proved spatial/\allowbreak{}time array law identifies each complete spatial block as the product of these scalar laws. Packing its first and second response coordinates gives exactly the correlated phase innovation used by the algorithm. Mutual independence across time identifies the entire defined innovation stream with the infinite product of that same phase law. The proof retains the correlation between velocity and position in each coordinate.



Each interval response is a centered Gaussian phase pair with the prescribed Gram covariance. Distinct spatial coordinates and time intervals are independent. The complete stream law follows from its finite cylinder laws, retaining within-pair correlations exactly.

\subsection{\texorpdfstring{Connect L\ensuremath{{}^{2}} endpoints and the exact thermal OU stream}{Connect L² endpoints and the exact thermal OU stream}}


The two vector-response representative equalities prove that the L\ensuremath{{}^{2}} endpoint pair is the pointwise canonical innovation. Taking and scaling its velocity component preserves temporal independence. The exact generic vector OU marginal theorem identifies its scale as \(\sqrt{1-\psi_0(h)^2}\). The independent-family law theorem then identifies the entire OU stream with the iid thermal Gaussian product required by the kinetic Gibbs stationarity proof.



Endpoint representative identities link the Hilbert-space source variables to the pointwise innovations. The filtered velocity endpoints also have the exact thermal stream law. This single generic source therefore supports both the KLMC coupling and the splitting-based stationarity proof.

\section{\texorpdfstring{From the integral equation to the exponential mild equation}{From the integral equation to the exponential mild equation}}
\label{proof-module-KineticLangevin-MildEquation-lean}


The KLMC update freezes the force inside an exponentially weighted integral. To compare it with an exact path, that weighted formula must follow from the integral equation driven by the same Brownian input. The theorem below proves this bridge for any complete real normed space, including the L\ensuremath{{}^{2}} space of random Euclidean vectors. No time derivative of Brownian motion is used.



The integrating-factor derivation removes the filtered forcing before differentiating. It gives an exact exponential response to the moving force while preserving the stochastic convolution.

\subsection{\texorpdfstring{Subtract the filtered noise and apply an integrating factor}{Subtract the filtered noise and apply an integrating factor}}


Assume continuous velocity \(V\), continuous force path \(F\), and a continuous filtered input \(B\), satisfying \(V_t-V_0=-\int_0^t(\gamma V_s+F_s)ds+W_t\) and \(B_t=W_t-\gamma\int_0^tB_sds\). The conclusion is the exact identity \(V_T=e^{-\gamma T}V_0-\int_0^Te^{-\gamma(T-s)}F_sds+B_T\). In the proof, evaluation at zero first gives \(B_0=0\). Subtracting the two equations shows \(D=V-B\) solves \(D_t=V_0-\int_0^t(\gamma D_s+F_s)ds\). Continuous integrands allow the fundamental theorem of calculus on the interval interior. The argument derives \(D\!\prime=-\gamma D-F\), differentiates \(e^{\gamma t}D_t\), and obtains derivative \(-e^{\gamma t}F_t\). A second application of the fundamental theorem integrates this identity over the closed interval; the proof handles the endpoints by continuity. Finally it multiplies by \(e^{-\gamma T}\), commutes scalar multiplication with the Bochner integral, and simplifies the exponent. The hypotheses are the integral equations, and the resulting mild equation is proved rather than postulated. The argument works for arbitrary real friction and includes \(T=0\).



After subtracting the filtered noise, the remaining velocity solves a linear inhomogeneous equation with the moving force as forcing. Multiplication by \(e^{\gamma t}\) gives its integral formula. Adding the filtered noise back yields the exponential mild identity needed to compare the exact path with a frozen-force step.

\section{\texorpdfstring{Constructing the reference and auxiliary coupling segments}{Constructing the reference and auxiliary coupling segments}}
\label{proof-module-KineticLangevin-CoupledSegments-lean}


The local KLMC analysis compares a numerical state with an exact reference state. On each step it inserts an auxiliary exact path started at the numerical state and driven by precisely the same additive input as the reference. This section constructs those paths from the proved kinetic flow. The differential, kinematic and mild identities used in the local-error theorem follow from their integral equations.



The exact reference, numerical chain, and auxiliary exact path have different initial phases but compatible local inputs. Their proved integral identities provide the complete structural foundation of the coupling argument.

\subsection{\texorpdfstring{Restart an integral solution at any time}{Restart an integral solution at any time}}


If \(Y_t-Y_0=-\int_0^tg(Y_s)ds+W_t\), then \(Y_{a+t}-Y_a=-\int_0^tg(Y_{a+s})ds+(W_{a+t}-W_a)\) whenever the shifted interval lies in the original horizon. The proof subtracts the two equations at \(a+t\) and \(a\), splits the continuous drift integral at \(a\), and changes variables in the remaining integral. Additive group algebra cancels the shared starting point. This supplies the exact time-shift identity required when segmenting a continuous diffusion.



Subtract the global integral equation at a restart time from its value later in the interval. The new initial phase is the global phase at the restart, and the new forcing is the corresponding input increment. Thus the shifted path solves the same kinetic equation locally.

\subsection{\texorpdfstring{Use a raw continuous input with the selected kinetic flow}{Use a raw continuous input with the selected kinetic flow}}


A continuous input \(W\) is restricted to the compact solution horizon, producing the existing the space of continuous forcing paths object. If \(W_0=0\), the flow's normalized clipped extension equals \(W_t\) at every time in that horizon. The argument uses the exact projection identity the identity of the interval projection on its interval to prove this. The selected segment is then the existing Picard-constructed the kinetic solution map with that compact input. Its specification gives continuity, its true initial condition, and the exact integral equation with the original raw input.



The selected flow only needs a continuous additive velocity input starting at zero. Subtracting the input’s initial value supplies that normalization for a raw continuous path. This permits local segments to be defined uniformly without choosing a new stochastic solution for each interval.

\subsection{\texorpdfstring{Derive velocity, kinematic and mild equations by projection}{Derive velocity, kinematic and mild equations by projection}}


Applying the velocity continuous linear projection to the kinetic integral equation gives \(V_t-V_0=-\int_0^t(\gamma V_s+u g(X_s))ds+W_t\). Applying the position projection gives \(X_t-X_0=\int_0^tV_sds\), because the positional input is exactly zero. In both cases The proof commutes the projection with the Bochner integral. If the filtered input satisfies \(B_t=W_t-\gamma\int_0^tB_sds\), the proved Banach-space integrating-factor theorem gives \(V_T=\psi_0(T)V_0-\int_0^Te^{-\gamma(T-s)}u g(X_s)ds+B_T\). Thus the exponential mild equation is a consequence of the kinetic equation, rather than a new stochastic hypothesis.



Projection of the phase equation gives the velocity and position integral identities. The integrating-factor argument then gives the mild velocity formula, while the position equation is already the kinematic primitive. These are the exact analytic identities required by the local-error estimate.

\subsection{\texorpdfstring{Construct exact reference segments and the defined numerical chain}{Construct exact reference segments and the defined numerical chain}}


The reference chain starts from its supplied phase vector. At each iteration, its next state is the endpoint of the selected exact kinetic segment driven by that iteration's input. The corresponding reference segment is the same selected path, so its endpoint matches the next reference state by definition. The numerical chain instead performs the frozen-force exponential update with innovation \((B_n(h),\int_0^hB_n(s)ds)\). This uses the constructed filtered input and its positional primitive.



The reference uses exact selected segments; the numerical sequence uses the defined exponential update with the same interval responses. Their endpoint matching is built into the recursions. The comparison is therefore between trajectories rather than arbitrary sequences satisfying a desired bound.

\subsection{\texorpdfstring{Construct the auxiliary exact path from the current numerical state}{Construct the auxiliary exact path from the current numerical state}}


The auxiliary segment is an selected exact kinetic solution whose initial condition is the current exponential-integrator state. It receives exactly the same raw Brownian input as the reference segment on that iteration. The two specification theorems prove continuity, exact starting values, reference endpoint matching, and their shared-input integral equations. Every field comes from the kinetic flow construction; no contraction or local-error inequality is inserted into the definitions.



At step \(n\), start an auxiliary exact segment at the numerical phase and drive it by the same input as the reference segment. It agrees initially with the numerical state, enabling the freezing-error estimate. Its discrepancy from the reference is controlled by exact-flow contraction.

\subsection{\texorpdfstring{Discharge the local-\allowbreak{}error proof's kinematic and mild requirements}{Discharge the local-error proof's kinematic and mild requirements}}


The constructed auxiliary path has its positional primitive identity by the general kinetic projection theorem. Its exponentially weighted force formula follows by restricting its integral equation to any smaller time \(t\le h\), restricting the filtered-input balance to the same interval, and applying the generic mild theorem. The argument explicitly checks the interval inclusions and simplifies the inverse mass \(u=1\). The remaining stochastic work concerns the law and stationarity of the reference; the coupling segments and their analytic identities are concrete constructions.



The local-error argument needs continuity, the mild equation, and the integrated-velocity position equation. Each follows from the selected solution’s specification. These hypotheses are therefore discharged internally when the constructed segments are inserted into the one-step coupling theorem.

\section{\texorpdfstring{Stitched exact segments are the global kinetic solution}{Stitched exact segments are the global kinetic solution}}
\label{proof-module-KineticLangevin-SegmentGlobal-lean}


A reference defined one interval at a time must be identified with the continuous kinetic process whose stationarity is proved separately. The bridge here uses the true shifted-input equation and uniqueness. It does not assume a Markov transition formula or independence of arbitrarily chosen reference segments.



Uniqueness identifies local restarts with shifts of one global kinetic solution. Therefore grid states and all intermediate segment times inherit the global reference laws used by local-error integrals.

\subsection{\texorpdfstring{Identify a selected local solution with a shifted global path}{Identify a selected local solution with a shifted global path}}


Start with an continuous global solution \(Y\), and suppose the local input is the increment \(W_t=B_{a+t}-B_a\). The previously proved time-shift lemma shows \(Y_{a+t}\) satisfies the same kinetic integral equation on the local horizon. The selected local solution starts at \(Y_a\). The argument checks the two initial conditions, the shared additive input, and continuity, then invokes the proved uniqueness theorem for the globally Lipschitz kinetic drift. Their equality follows at every time in the local interval.



The shifted global path and the selected local path have the same starting phase and the same input increment. Subtract their integral equations and apply uniqueness. Their equality follows on the whole local interval, not only at its endpoint.

\subsection{\texorpdfstring{Identify every stitched grid point with the global solution}{Identify every stitched grid point with the global solution}}


The reference chain starts at \(Y_0\). Inductively assume its state at step \(n\) equals \(Y_{nh}\). The next selected segment begins at that same point, so the local-global uniqueness result identifies its endpoint with \(Y_{nh+h}\). The proof performs induction over the chain recursion, checks the finite-horizon inequalities using nonnegative step size, and simplifies the real cast of \(n+1\). The result is the exact identity \(R_n=Y_{nh}\), including zero step size.



Induct on the grid index. At zero the initial states coincide; equality of the selected local solution with the shifted global solution identifies the next stitched endpoint. Thus the recursively constructed exact reference sequence is exactly the global flow sampled on the grid.

\subsection{\texorpdfstring{Identify all intermediate times within a reference segment}{Identify all intermediate times within a reference segment}}


Once the starting grid point is identified, the same uniqueness result proves \(Z_n(t)=Y_{nh+t}\) throughout the whole segment. This lets a stationary global phase law supply both the stationary reference position at grid points and the Gaussian reference velocity at intermediate times required by the KLMC local-error proof. No stationary law is assumed for a newly invented process: the stitched reference is proved to be the global solution.



The same uniqueness argument applies before the local endpoint. Hence every intermediate reference time is a value of the global solution. Global stationarity and moment statements can therefore be transferred to all local times needed in integral error bounds.

\section{\texorpdfstring{Exact Gibbs invariance of an explicit corrected Verlet proposal}{Exact Gibbs invariance of an explicit corrected Verlet proposal}}
\label{proof-module-KineticLangevin-HamiltonianMetropolis-lean}


A volume-preserving involution has a symmetric old/\allowbreak{}new-state graph measure. Weighting that graph by a positive target density and applying the ordinary Metropolis ratio makes its accepted flow symmetric. The rejected mass stays at the old state. This section proves exact invariance for the explicit reversed-Verlet proposal. Passing from these corrected proposals to the unadjusted Hamiltonian flow and then to the full kinetic diffusion still requires quantitative limits.



Involutivity and volume preservation produce a symmetric accepted-flow measure. The rejected mass completes a probability transition, while the energy defect controls how closely it resembles the uncorrected Verlet step.

\subsection{\texorpdfstring{Construct the proposal graph and derive its swap symmetry}{Construct the proposal graph and derive its swap symmetry}}


For \(J^2=\mathrm{id}\), the graph map is \(G(z)=(z,Jz)\). Swapping the old and new state gives \(\operatorname{swap}\circ G=G\circ J\). If \(J\) preserves a base measure \(\lambda\), pushforward composition yields \(\operatorname{swap}_\#G_\#\lambda=G_\#J_\#\lambda=G_\#\lambda\). One proves the pointwise graph identity using the involution equation, verifies graph measurability, and rewrites the two pushforward compositions. Thus symmetry of the graph measure is derived from the proposal rather than introduced as a detailed-balance field.



The proposal graph associates each state with its deterministic proposal. Swapping the two graph coordinates corresponds to applying the proposal again. Involutivity and volume preservation therefore identify the swapped graph measure with the original one after the appropriate density change.

\subsection{\texorpdfstring{Apply the accepted-\allowbreak{}flow argument to an arbitrary positive density}{Apply the accepted-flow argument to an arbitrary positive density}}


Let the target be \(w\lambda\), with measurable \(w>0\). Its proposal joint law is \(G_\#(w\lambda)\). The pushforward-with-density theorem expresses this as the symmetric base graph weighted by \(w\) at its first coordinate. The Metropolis ratio is \(\min\{1,w(y)/w(x)\}\), so accepted joint density is \(\min\{w(x),w(y)\}\), symmetric under swapping. The shared Metropolis theorem proves that accepted new-state mass plus rejected old-state mass has the initial marginal. The proof computes that marginal by projecting the graph back to its first coordinate, giving the exact identity with \(w\lambda\). No invariant measure or acceptance-balance conclusion is assumed.



The accepted mass is the smaller of the initial and proposed density weights. This minimum is symmetric under swapping the graph coordinates. Rejected mass remains at the initial state, so the accepted and rejected contributions together have the required invariant marginal.

\subsection{\texorpdfstring{Bound rejection by an energy increase}{Bound rejection by an energy increase}}


For energy increment \(d\), rejection is \(1-\min\{1,e^{-d}\}\). If \(d\le0\), the proposed energy decreases and rejection is zero. If \(d\ge0\), the exponential tangent inequality \(e^{-d}\ge1-d\) gives rejection at most \(d\). Thus rejection is at most \(\max\{d,0\}\le|d|\). The proof splits the sign cases, rewrites the minimum and maximum with the corresponding exponential bounds, and applies the tangent inequality. This is a scalar probability bound; converting it to an L\ensuremath{{}^{2}} trajectory error also requires the size of the rejected jump and must not silently assume that jump is of order \(h\).



For a nonnegative energy increase \(u\), the rejection probability is \(1-e^{-u}\le u\); for a nonpositive increase it is zero. Thus rejection is bounded by the positive part of the energy defect. The inequality turns an analytic energy estimate into a probabilistic rejection bound.

\subsection{\texorpdfstring{Instantiate the exact corrected Verlet Gibbs theorem}{Instantiate the exact corrected Verlet Gibbs theorem}}


Take \(\lambda\) to be product Haar phase measure, \(J_h=R\,V_h\) to be the constructed reversed Verlet proposal, and \(w(z)=e^{-H(z)}\). The Hamiltonian module proves that \(J_h\) is both an involution and measure preserving. Measurability of the potential and gradient proves measurability of the weight and proposal, and the exponential weight is strictly positive. The generic theorem therefore proves exact preservation of the Gibbs weighted measure by this Metropolis step. Its accepted state is velocity-reversed Verlet and its rejected state is the old state. Applying a further velocity reversal changes those to ordinary Verlet on acceptance and a velocity flip on rejection; any limiting kinetic argument must account for this distinction explicitly.



For the Gibbs density the ratio of proposal and initial weights is the exponential of the negative energy defect. Substitution in the general balance calculation gives the corrected Verlet invariant measure. This exact statement does not imply invariance of uncorrected Verlet.

\section{\texorpdfstring{Explicit Verlet energy and rejection estimates}{Explicit Verlet energy and rejection estimates}}
\label{proof-module-KineticLangevin-HamiltonianError-lean}


The exact Hamiltonian defect proved earlier becomes quantitative under a globally Lipschitz gradient. The proof uses only first derivatives: two-sided Taylor control bounds the trapezoidal force work, while a norm factorization bounds the force-square correction. These estimates concern the algebraic Verlet map and can control the rejection mass of the invariant corrected approximation.



Symmetric Taylor expansions match the two half kicks of Verlet and expose cancellation of the leading energy terms. The remaining bound is expressed through velocity and force moments suitable for Gibbs integration.

\subsection{\texorpdfstring{Obtain the two-\allowbreak{}sided Taylor remainder}{Obtain the two-sided Taylor remainder}}


The shared smooth upper bound gives \(f(y)-f(x)-\langle\nabla f(x),y-x\rangle\le M\|y-x\|^2/2\). Apply the same theorem to \(-f\). Its gradient is \(-\nabla f\), proved by rewriting the Fréchet derivative and the linear Riesz inverse, and its Lipschitz constant remains \(M\). The second application gives the reverse inequality, so the absolute remainder has the same bound. The argument explicitly rewrites the negative gradient and combines the upper and lower scalar inequalities with the two-sided absolute-value bound.



The fundamental theorem along the segment gives \(f(y)-f(x)-\langle\nabla f(x),y-x\rangle\). Gradient Lipschitzness bounds its absolute value by \(M\|y-x\|^2/2\). The two-sided form is needed because numerical energy can increase or decrease.

\subsection{\texorpdfstring{Average Taylor errors from both endpoints}{Average Taylor errors from both endpoints}}


Write \(d=y-x\), \(r_x=f(y)-f(x)-\langle\nabla f(x),d\rangle\), and \(r_y=f(x)-f(y)+\langle\nabla f(y),d\rangle\). The trapezoidal force-work error is exactly \((r_x-r_y)/2\). Both remainders have absolute value at most \(M\|d\|^2/2\). The triangle inequality therefore gives \[\left|f(y)-f(x)-\tfrac12\langle\nabla f(x)+\nabla f(y),d\rangle\right|\le\tfrac M2\|d\|^2.\] One proves the algebraic average identity after expanding the inner product, reverses the displacement in the second Taylor bound, and applies the absolute-value triangle inequality.



Writing the potential difference once from each endpoint and averaging produces a trapezoidal force term. Its remainder is controlled by the same segment displacement. This symmetric arrangement matches the two force evaluations of velocity Verlet.

\subsection{\texorpdfstring{Control the squared-\allowbreak{}force correction}{Control the squared-force correction}}


Factor \(\|g(y)\|^2-\|g(x)\|^2=(\|g(y)\|-\|g(x)\|)(\|g(y)\|+\|g(x)\|)\). The reverse triangle inequality and Lipschitz continuity bound the first factor by \(M\|y-x\|\). The ordinary triangle inequality bounds the second by \(2\|g(x)\|+M\|y-x\|\). Multiplying these nonnegative estimates proves the displayed force-square bound. The argument derives the endpoint-force growth estimate before applying multiplication of nonnegative inequalities, so every sign condition is checked.



The kinetic-square expansion contains differences of squared force norms. Factor each difference as an inner product of the force difference and force sum. Lipschitz continuity controls the first factor; the second is bounded by the endpoint force magnitudes.

\subsection{\texorpdfstring{Bound the complete Verlet energy defect}{Bound the complete Verlet energy defect}}


Set \(w=v-h\nabla f(x)/2\), \(y=x+hw\). The exact identity from the Hamiltonian module writes the defect as a trapezoidal work error plus \(h^2(\|\nabla f(y)\|^2-\|\nabla f(x)\|^2)/8\). Since \(\|y-x\|=|h|\|w\|\), the preceding two lemmas give \[|H(V_hz)-H(z)|\le\frac M2h^2\|w\|^2+\frac{h^2}{8}M|h|\|w\|\bigl(2\|\nabla f(x)\|+M|h|\|w\|\bigr).\] The proof substitutes the intermediate stage, changes the inner product to the displacement form, applies the triangle inequality to the two exact defect terms, and simplifies \(|h|^2=h^2\). The leading term is quadratic in the step under only gradient Lipschitz regularity.



Combine the potential and kinetic expansions before estimating their absolute values. Their leading terms cancel because the Verlet stages use matching half kicks. The remaining estimate is expressed in terms of velocity and force magnitudes so it can later be integrated under a Gibbs moment bound.

\section{\texorpdfstring{The accepted-\allowbreak{}Verlet and rejected-\allowbreak{}flip invariant measure}{The accepted-Verlet and rejected-flip invariant measure}}
\label{proof-module-KineticLangevin-HamiltonianWeak-lean}


The reversed-Verlet involution supplied exact Metropolis invariance. Applying one final velocity reversal gives ordinary Verlet on acceptance and a velocity flip on rejection. This section computes the resulting two weighted pushforward measures explicitly and proves that their sum equals the Hamiltonian Gibbs measure. This identity is the concrete starting point for bounded-test estimates and a numerical limit.



The corrected transition is described by its accepted and rejected measures. Exact invariance supplies the stationary comparison used to prove weak error bounds for uncorrected numerical Hamiltonian evolution.

\subsection{\texorpdfstring{Compute the two Metropolis graph marginals}{Compute the two Metropolis graph marginals}}


For a deterministic proposal \(J\), the old/\allowbreak{}new graph is \(G(z)=(z,Jz)\). Weighting its pushforward by a measurable acceptance function \(a(z,Jz)\) commutes with mapping the correspondingly weighted initial measure through \(G\). Projecting onto the second coordinate yields \(J_\#(a\pi)\). Projecting rejected mass onto the first coordinate yields \((1-a)\pi\). The argument applies the previously proved pushforward-with-density formula twice and composes each graph with its coordinate projection. This proves an equality of measures without an integrability premise.



The proposal graph records both the old and proposed states. Its accepted mass has symmetric first and second marginals, while the rejected part remains at the old state after the prescribed reversal. Computing both marginals explicitly proves invariance without omitting rejected probability mass.

\subsection{\texorpdfstring{Construct the Hamiltonian Gibbs measure and prove velocity symmetry}{Construct the Hamiltonian Gibbs measure and prove velocity symmetry}}


The measure is \(e^{-H(v,x)}(\mu\otimes\mu)(d(v,x))\), where \(\mu\) is a translation-invariant symmetric Haar factor and \(H(v,x)=f(x)+\|v\|^2/2\). It is initially unnormalized; finiteness and normalization are independent questions. Velocity reversal preserves product Haar measure and preserves \(H\) exactly. The weighted pushforward identity therefore proves that it preserves this Gibbs measure. The argument verifies measurability of the exponential density and rewrites the two exact preservation identities.



The Hamiltonian density factors into a velocity Gaussian weight and a position Gibbs weight. Reversing velocity leaves that density and volume unchanged. This symmetry permits the rejection flip to be included without changing the target measure.

\subsection{\texorpdfstring{Use the Verlet energy change as acceptance}{Use the Verlet energy change as acceptance}}


The acceptance probability is \(a_h(z)=\min\{1,\exp[-(H(V_hz)-H(z))]\}\). Its measurability follows by expanding the explicit kick-drift-kick stages and composing the measurable potential and gradient. Positivity of the exponential gives \(a_h\ge0\), and the minimum gives \(a_h\le1\). These scalar bounds later justify finite weighted measures and nonnegative rejection estimates.



The acceptance ratio must use the difference of the true Hamiltonian evaluated at the proposed and current states. Replacing it by a truncated local approximation would define a different transition. The exact algebraic defect from the Verlet map therefore enters the formula unchanged.

\subsection{\texorpdfstring{Derive exact corrected Verlet invariance}{Derive exact corrected Verlet invariance}}


The proved Metropolis result uses the involution \(J_h=R\circ V_h\). Since \(H\circ R=H\), its ratio is exactly the acceptance \(a_h\) just defined. Compute its accepted and rejected graph marginals, then push the resulting equality through \(R\). The accepted state becomes \(R(J_hz)=V_hz\), whereas the rejected state becomes \(Rz\). Gibbs velocity symmetry fixes the right-hand side, giving \[(V_h)_\#(a_h\pi)+R_\#((1-a_h)\pi)=\pi.\] The proof performs precisely these graph, density, and pushforward rewrites and uses the identity \(R^2=\mathrm{id}\). This theorem does not identify rejection with a small spatial jump; the rejected velocity flip is retained explicitly.



Integrating any bounded measurable test against the corrected transition gives its original target integral. This is an equality of probability measures after normalization. It will be used as a comparator when bounding weak error of the uncorrected Hamiltonian step.

\section{\texorpdfstring{Bounded-\allowbreak{}test estimates from the exact corrected measure}{Bounded-test estimates from the exact corrected measure}}
\label{proof-module-KineticLangevin-HamiltonianTests-lean}


The rejected move can reverse a velocity of order one. A weak approximation therefore uses bounded complex-valued tests rather than assuming the rejected jump is proportional to the step. The exact corrected measure identity implies that one uncorrected proposal has test bias at most twice the test bound times its mean rejection probability.



Bounded tests permit rejection errors to be controlled by probability mass alone. This yields a robust route from an exact corrected invariant measure to vanishing weak error under numerical refinement.

\subsection{\texorpdfstring{Keep all acceptance and rejection mass}{Keep all acceptance and rejection mass}}


For measurable \(0\le a\le1\), the two weighted measures \(a\pi\) and \((1-a)\pi\) sum exactly to \(\pi\). The proof combines the two densities, uses nonnegativity to rewrite the sum of the nonnegative extended-real embedding values as the value of their sum, and reduces the resulting density to one. In particular both weighted measures are bounded above by the original finite measure.



For each state, the accepted and rejected coefficients sum to one. Keeping both terms makes the corrected transition a probability kernel and its target identity exact. Dropping the rejected term would produce a subprobability expression and an incorrect invariance statement.

\subsection{\texorpdfstring{Turn the explicit invariant measure into an integral identity}{Turn the explicit invariant measure into an integral identity}}


Suppose the already constructed corrected measure satisfies \(V_\#(a\pi)+R_\#((1-a)\pi)=\pi\). This identity is proved for the Verlet and reversal maps in the preceding module. For bounded measurable complex \(F\), integration gives \[\int F\,d\pi=\int a(z)F(Vz)\,d\pi+\int(1-a(z))F(Rz)\,d\pi.\] First, the argument proves finiteness of each weighted measure from the exact mass split. The bound on \(F\) supplies integrability for both pushforward measures. It then applies the Bochner integral formula for a pushforward and the formula for a density-weighted integral; the acceptance and rejection factors are finite nonnegative real numbers.



An invariant measure equality implies equality of integrals for every bounded measurable test. This formulation is convenient for passing to numerical limits, because the error between corrected and uncorrected tests can be controlled pointwise by the rejection probability.

\subsection{\texorpdfstring{Bound the entire rejected move through the test bound}{Bound the entire rejected move through the test bound}}


Subtract the invariant-test identity from the ordinary proposal expectation. The exact difference is \[\int F(Vz)\,d\pi-\int F\,d\pi=\int(1-a(z))\,[F(Vz)-F(Rz)]\,d\pi.\] If \(\|F\|\le C\), the norm of the bracket is at most \(2C\), irrespective of the size of the rejected physical jump. Since \(1-a\ge0\), the norm of the integral is at most \(2C\int(1-a)\,d\pi\). The argument explicitly checks bounded integrability of the proposal test, accepted test, rejected test, and rejection factor, rewrites the subtraction under the integral, applies the norm-of-integral inequality, and integrates the pointwise triangle bound. Thus the velocity flip is fully accounted for.



If a test is bounded in absolute value by \(B\), changing its argument changes its value by at most \(2B\). Multiplying this bound by the rejection probability controls the entire rejected move. No Lipschitz regularity of the test is required for this step.

\section{\texorpdfstring{Linear accumulation of weak error and exact limiting laws}{Linear accumulation of weak error and exact limiting laws}}
\label{proof-module-KineticLangevin-HamiltonianIteration-lean}


Corrected invariance can be used entirely through bounded tests. Composing a test with more deterministic proposal steps preserves its bound. Each telescoping increment is therefore controlled under the same fixed measure, giving a linear accumulation of rejection mass. Pointwise convergence and vanishing accumulated error then identify the limiting law exactly by dual characteristic functions.



Telescoping bounded-test errors produces linear accumulation over a fixed horizon. Once the numerical maps converge, characteristic-function uniqueness identifies the invariant law of the limiting Hamiltonian flow.

\subsection{\texorpdfstring{Telescope tests composed with deterministic iterates}{Telescope tests composed with deterministic iterates}}


Let \(A_k=\int F(V^kz)\,d\pi\). Apply the one-step weak estimate to \(G=F\circ V^k\). Measurability of iterates proves measurability of \(G\), and the same bound \(\|G\|\le C\) holds. Thus \(\|A_{k+1}-A_k\|\le2C\int(1-a)\,d\pi\). The triangle inequality and induction give \[\|A_n-A_0\|\le2Cn\int(1-a)\,d\pi.\] The argument explicitly rewrites the successor iterate as \(V^k\circ V\), applies the one-step result, and normalizes the successor natural-number cast. Every increment is integrated against the fixed measure \(\pi\); no law of an unadjusted intermediate chain is assumed invariant.



Write the difference between the \(n\)-step test integral and its invariant value as a sum of \(n\) one-step differences. Composition with deterministic iterates preserves the test’s uniform bound. Each term therefore costs the same local rejection estimate, giving linear rather than exponential accumulation.

\subsection{\texorpdfstring{Identify the limiting pushforward through characteristic tests}{Identify the limiting pushforward through characteristic tests}}


Suppose measurable maps \(Y_n\) converge pointwise to a measurable map \(Y\), and the bounded-test error is eventually at most a real sequence \(\delta_n\to0\). For every continuous real linear functional \(L\), use \(F(z)=\exp(iLz)\), whose norm is one. Dominated convergence under the finite measure gives \(\int F(Y_nz)\,d\pi\to\int F(Yz)\,d\pi\). The error bound gives convergence of the same sequence to \(\int F\,d\pi\). Uniqueness of limits equates the dual characteristic functions of \(Y_\#\pi\) and \(\pi\); the uniqueness theorem for finite measures gives \(Y_\#\pi=\pi\). The argument uses dual characteristic functions, so the result applies to Banach phase space with its usual product norm and does not silently require that norm to arise from an inner product.



Characteristic functions determine probability laws in finite-dimensional Euclidean space. If the numerical maps converge to the Hamiltonian flow and the accumulated test error vanishes, bounded convergence identifies the limiting pushforward with the target. This gives invariance of the flow.

\section{\texorpdfstring{Quadratic rejection mass under a concrete phase moment}{Quadratic rejection mass under a concrete phase moment}}
\label{proof-module-KineticLangevin-HamiltonianRejection-lean}


The exact Verlet energy defect controls its Metropolis rejection probability. For steps between zero and one, the state-dependent bound simplifies to a single quadratic phase envelope. Integrating this bound requires precisely the velocity and force second moments displayed in theorem, and supplies the vanishing error needed for a fixed-horizon weak limit.



The numerical energy defect is reduced to an integrable quadratic phase envelope. Its expectation gives a quantitative rejection bound, enabling exact corrected invariance to approximate uncorrected evolution.

\subsection{\texorpdfstring{Apply the scalar rejection inequality to the energy defect}{Apply the scalar rejection inequality to the energy defect}}


The acceptance is defined by the defect \(d=H(V_hz)-H(z)\). The already proved scalar inequality \(1-\min\{1,e^{-d}\}\le |d|\) therefore applies by substitution, with no acceptance approximation.



The scalar acceptance bound is first applied pointwise to the exact numerical energy change. An upper bound on its absolute value also bounds its positive part. The resulting rejection envelope is measurable and can therefore be integrated under the target law.

\subsection{\texorpdfstring{Reduce all intermediate stages to a quadratic phase envelope}{Reduce all intermediate stages to a quadratic phase envelope}}


Put \(a=\|v\|\), \(b=\|\nabla f(x)\|\), \(w=v-h\nabla f(x)/2\), and \(q=a^2+b^2\). For \(0\le h\le1\), the triangle inequality gives \(\|w\|\le a+b\), hence \(\|w\|^2\le2q\). Also \(2\|w\|b\le\|w\|^2+b^2\), so \(\|w\|(2b+Mh\|w\|)\le(2M+3)q\). Substitution in the earlier exact defect estimate yields \[|H(V_hz)-H(z)|\le h^2\frac{M(2M+11)}8\bigl(\|v\|^2+\|\nabla f(x)\|^2\bigr).\] The proof retains the intermediate nonnegative scalar bounds, checks every multiplication by \(h\) and \(M\), and combines them by polynomial arithmetic. The following theorem composes this estimate with the scalar rejection inequality.



Intermediate kick velocities and drift positions are bounded by the original velocity and force using the triangle inequality and force Lipschitzness. Squaring introduces cross terms, controlled by \(2ab\le a^2+b^2\). This reduces the envelope to finitely many second moments.

\subsection{\texorpdfstring{Integrate the proved rejection envelope}{Integrate the proved rejection envelope}}


For any finite phase measure with integrable \(q(z)=\|v\|^2+\|\nabla f(x)\|^2\), the rejection factor is measurable and lies between zero and one, so it is integrable. The pointwise bound integrates to \[\int(1-a_h)\,d\pi\le h^2\frac{M(2M+11)}8\int q\,d\pi.\] One proves bounded integrability of rejection directly from the acceptance bounds and applies monotonicity of the ordinary integral against the explicitly assumed finite quadratic moment. No one-step trajectory or stationary-flow conclusion is supplied as a premise.



Under a finite quadratic phase envelope, the pointwise rejection estimate has a finite integral. The step-size factors remain explicit while the phase dependence becomes a target moment constant. This separation is what makes the later weak-error bound uniform over a fixed time horizon.

\section{\texorpdfstring{Gibbs finiteness and the exact phase moment from strong convexity}{Gibbs finiteness and the exact phase moment from strong convexity}}
\label{proof-module-KineticLangevin-HamiltonianMoments-lean}


The weak Hamiltonian limit needs finite mass and an integrable velocity-plus-force quadratic envelope. This section derives both for the exponential Hamiltonian density in finite dimensions. It factors that density into a radial Gaussian velocity weight and the potential Gibbs weight, uses existing direct Gaussian-domination estimates, and transfers position moments to force moments by Lipschitz growth.



Strong convexity makes the Gibbs phase measure finite and gives the moments required by the rejection envelope. Smoothness converts position control into force control, closing the integrability assumptions of the weak limit.

\subsection{\texorpdfstring{Factor the Hamiltonian density}{Factor the Hamiltonian density}}


Because \(H(v,x)=f(x)+\|v\|^2/2\), the positive exponential weight factors as \(e^{-H(v,x)}=e^{-\|v\|^2/2}e^{-f(x)}\). Therefore its product-Haar weighted measure is the product of the separately weighted velocity and position measures. The argument applies theorem for products of measures with densities, combines the two nonnegative real weights through multiplicativity of the nonnegative embedding, and checks the exponent identity. Normalization has not been assumed or imposed at this step.



The kinetic and potential parts of the Hamiltonian depend on different coordinates, so their exponential weights factor. Product integration separates velocity Gaussian integrals from position Gibbs integrals. This factorization is valid for the reference measure before normalization.

\subsection{\texorpdfstring{Derive finite mass of both factors and their product}{Derive finite mass of both factors and their product}}


The velocity weight is the radial Gaussian with precision \(1/2\), whose ordinary integral is finite. The position weight has finite mass by the previously proved strong-convexity Gaussian domination: a supporting quadratic lower bound for \(f\) makes \(e^{-f}\) integrable. The proof converts these finite ordinary integrals to finite extended density integrals, instantiates finite weighted measures, and uses the exact factorization to prove finiteness of the full phase measure.



Strong convexity gives a quadratic lower envelope for the potential. Its exponential is dominated by a shifted Gaussian density, yielding finite positive partition mass. The velocity factor is an ordinary Gaussian integral, so the phase measure also has finite positive mass.

\subsection{\texorpdfstring{Derive velocity and position second moments before normalization}{Derive velocity and position second moments before normalization}}


A quadratic factor is absorbed by radial Gaussian decay, so \(\int\|v\|^2e^{-\|v\|^2/2}\,dv<\infty\). The same strong-convexity tail estimate gives \(\int\|x\|^2e^{-f(x)}\,dx<\infty\). The argument uses the ordinary-integral characterization of integration under a density, checks that each exponential density is finite and measurable, and invokes the existing Gaussian-tail estimates. These are moments of the unnormalized factors, not moment contracts attached to a stationary process.



Polynomial factors remain integrable under the Gaussian lower envelope. Thus both velocity and position second moments exist before taking ratios by the partition function. Establishing this first ensures that later normalized moment expressions represent finite expectations.

\subsection{\texorpdfstring{Transfer position moments to force moments}{Transfer position moments to force moments}}


Gradient Lipschitz continuity implies \(\|\nabla f(x)\|\le M\|x\|+\|\nabla f(0)\|\). Squaring and using \((a+b)^2\le2a^2+2b^2\) gives \(\|\nabla f(x)\|^2\le2M^2\|x\|^2+2\|\nabla f(0)\|^2\). The right-hand side is integrable under the finite position Gibbs weight by its second moment. The argument derives the force growth inequality from the Lipschitz difference bound and triangle inequality, then applies domination by this explicit integrable envelope.



Force Lipschitzness gives \(\|\nabla f(x)\|\le\|\nabla f(x_0)\|+M\|x-x_0\|\). Squaring and using \((a+b)^2\le2a^2+2b^2\) reduces force moments to the finite position second moment. Choosing a minimizer removes the constant force term when available.

\subsection{\texorpdfstring{Assemble the precise phase envelope needed by the weak limit}{Assemble the precise phase envelope needed by the weak limit}}


On the factored finite phase measure, \(\|v\|^2\) is integrable by pulling the velocity moment through the first coordinate, and \(\|\nabla f(x)\|^2\) is integrable by pulling the force moment through the second coordinate. Their sum is exactly the envelope appearing in the rejection theorem. The argument rewrites the phase measure by its proved factorization, supplies finiteness of both factors, and adds these two coordinate integrability results. Thus every tail premise in the later Hamiltonian stationarity limit follows from positive strong convexity and a Lipschitz gradient.



The rejection envelope is a fixed quadratic combination of velocity and force norms. The preceding moment bounds therefore make its target integral finite. This is the exact integrability required to make the accumulated rejection error vanish in the Hamiltonian-flow limit.

\section{\texorpdfstring{Exact Gibbs invariance of the Hamiltonian flow}{Exact Gibbs invariance of the Hamiltonian flow}}
\label{proof-module-KineticLangevin-HamiltonianStationarity-lean}


This section closes the deterministic Hamiltonian stationarity argument. Explicit corrected Verlet proposals preserve the Gibbs measure exactly. Their rejection probability is quadratically small, while the uncorrected Verlet iterates converge to the selected Hamiltonian integral solution. Bounded tests transfer invariance through this limit. The finite-dimensional strong-convexity theorem derives every required tail premise, so its statement assumes no invariant flow, numerical convergence, or error contract.



Exact corrected invariance, a vanishing rejection discrepancy, and Verlet convergence are combined in one limit. This proves invariance for the constructed Hamiltonian solution rather than assuming a stationary comparator.

\subsection{\texorpdfstring{Make the endpoint a measurable function of its initial state}{Make the endpoint a measurable function of its initial state}}


The constructed kinetic endpoint is jointly continuous in initialization and continuous forcing. Fix zero forcing, friction zero, inverse mass one, and endpoint time \(T\). Composing this joint map with \(z\mapsto(z,0)\) proves continuity of the selected Hamiltonian endpoint. The argument supplies the compact forcing-space zero explicitly and simplifies the composed definition. The endpoint pushforward is therefore a concrete measurable law.



Continuous dependence on the initial phase makes the terminal Hamiltonian map measurable. Its pushforward under the normalized Gibbs phase law is therefore well defined. The selected solution’s uniqueness ensures this map does not depend on a choice in the existence construction.

\subsection{\texorpdfstring{Pass exact corrected invariance to the integral solution}{Pass exact corrected invariance to the integral solution}}


Let \(\pi\) be the Hamiltonian weighted measure, \(B=\int(\|v\|^2+\|\nabla f(x)\|^2)\,d\pi\), and \(K=M(2M+11)/8\). Use \(N=n+1\) Verlet steps of size \(h=T/N\). Exact corrected invariance and the iterated bounded-test estimate imply, for \(\|F\|\le1\), \[\left\|\int F(V_h^Nz)\,d\pi-\int F\,d\pi\right\|\le 2N\int(1-a_h)\,d\pi\le\frac{2T^2KB}{N}.\] The mesh is eventually at most one, so the proved quadratic rejection estimate applies; the displayed bound tends to zero. The independently proved numerical theorem gives \(V_{T/N}^Nz\to\Phi_Tz\), where \(\Phi\) is the Picard integral solution. The bounded-test limit theorem identifies \((\Phi_T)_\#\pi=\pi\). The proof assembles these results, proves the mesh limit, checks nonnegative multipliers, and simplifies the exact scalar bound. It assumes only the finite measure and precisely displayed quadratic moment required by this intermediate general theorem.



Corrected Verlet preserves the target exactly, while its difference from uncorrected Verlet is bounded by integrated rejection. Over a fixed horizon the total rejection error tends to zero as the mesh is refined. Convergence to the flow then transfers invariance to that flow.

\subsection{\texorpdfstring{Discharge every tail premise from the potential assumptions}{Discharge every tail premise from the potential assumptions}}


In finite dimensions, positive strong convexity and differentiability make the Hamiltonian Gibbs measure finite by the factorization and Gaussian domination proved in the moments module. The globally Lipschitz gradient supplies the finite force second moment, and thermal velocity factor supplies the velocity second moment. Substituting these results into the preceding limit theorem gives exact Hamiltonian Gibbs invariance from the original first-order potential assumptions. The theorem applies to every nonnegative finite horizon. It proves the Hamiltonian subflow; the full friction-and-noise kinetic diffusion additionally needs its thermostat splitting limit.



The necessary integrable envelope follows from strong convexity and smoothness through the phase moment estimates. These hypotheses also supply finite normalization. Consequently the stationarity theorem closes its own tail and moment requirements rather than leaving an assumed invariant reference measure.

\section{\texorpdfstring{The Gaussian-\allowbreak{}times-\allowbreak{}Gibbs phase target}{The Gaussian-times-Gibbs phase target}}
\label{proof-module-KineticLangevin-GibbsTarget-lean}


The normalized kinetic target is the independent product of a standard Gaussian velocity and the intended Gibbs position measure. This section identifies that product with an explicit scalar normalization of the Hamiltonian weighted measure, transfers the proved Hamiltonian invariance, and proves exact Gaussian thermostat invariance on the correct independent-noise product. Its position and velocity marginals are exact pushforward identities. These two invariant substeps are ingredients for kinetic Gibbs stationarity; passage to the continuous kinetic limit is a separate theorem.



The canonical target is identified by normalized measure equalities. Its exact product structure supplies both invariant building blocks, the desired position marginal, and the Gaussian moments used in numerical error analysis.

\subsection{\texorpdfstring{Identify the Gaussian thermal density by its law}{Identify the Gaussian thermal density by its law}}


The already Gaussian proposal density at step \(h=1/2\) and center zero is exactly the standard Gaussian. Its scalar precision is \(1/(4h)=1/2\), and its affine map is the identity. Consequently its density is \(c_Ge^{-\|v\|^2/2}\), where \(c_G\) is the inverse of the finite Gaussian integral. First, the argument normalizes the precision and the identity pushforward, then factors the positive normalizer out of the density using scaling of a density-weighted measure and multiplicativity of the nonnegative embedding. The remaining real exponential identity follows by scalar algebra. This is an equality with the existing probability measure the standard Gaussian measure, not an informal naming of a radial weight.



The standard Gaussian velocity law has a density proportional to \(e^{-\|v\|^2/2}\) with respect to Euclidean volume. This identifies the explicit Gaussian measure with the kinetic factor in the Hamiltonian Gibbs weight. The normalization is therefore known, not left implicit.

\subsection{\texorpdfstring{Normalize the Hamiltonian weighted measure explicitly}{Normalize the Hamiltonian weighted measure explicitly}}


Define \(\Pi=\varphi\otimes\pi\), velocity first, with \(\varphi\) the standard Gaussian and \(\pi\) the normalized Gibbs measure of \(f\). The Gaussian density identity and the existing definition of \(\pi\) give

\[\Pi=(c_G/Z_f)\,e^{-f(x)-\|v\|^2/2}\,dv\,dx.\]

The proof distributes the two scalar normalizers through the product of measures, then rewrites the previously proved factorization of the Hamiltonian Gibbs weight. Strong convexity with \(m>0\) makes the position normalization positive and finite, so the product is a probability measure. Its probability proof invokes the already derived Gibbs probability theorem; normalization is not an assumption about a stationary process.



Product integration gives the phase partition function as the product of the velocity and position partition functions. Dividing the factored density by that positive finite mass yields the independent Gaussian-times-position probability measure used in the sampler theorem.

\subsection{\texorpdfstring{Transfer the proved Hamiltonian invariance to the canonical target}{Transfer the proved Hamiltonian invariance to the canonical target}}


The selected Hamiltonian flow already preserves the unnormalized Hamiltonian weighted measure, by the corrected-Verlet approximation and its proved numerical limit. Multiplying a measure by a fixed scalar commutes with a measurable pushforward. The explicit normalization identity therefore gives \((H_t)_\#\Pi=\Pi\). The argument supplies measurability from continuous dependence of the endpoint on its initial phase state, uses compatibility of pushforward and measure scaling, and applies the proved Hamiltonian invariance theorem. This statement has only the potential assumptions and finite-dimensional setting; no invariant-target hypothesis remains.



The Hamiltonian flow preserves its normalized energy-weighted measure. The density factorization identifies that measure exactly with the canonical Gibbs target. Hence the proved Hamiltonian invariance transfers to the same concrete target used by the kinetic analysis.

\subsection{\texorpdfstring{Prove the thermostat identity with fresh independent Gaussian noise}{Prove the thermostat identity with fresh independent Gaussian noise}}


Let \(X\) have any probability position law, independently of a standard Gaussian velocity \(V\), and let \(\xi\) be a fresh independent standard Gaussian. If \(a^2+b^2=1\), then \((aV+b\xi,X)\) again has the original independent velocity-position law. The input is explicitly \((\varphi\otimes\pi_X)\otimes\varphi\). A measure-preserving reassociation and swap reorder its three coordinates to \((V,\xi),X\). The previously proved Gaussian linear-combination law acts on the first pair; the position coordinate is unchanged. The argument verifies both the coordinate permutation and its product-measure pushforward before applying the pushforward formula for product measures. Thus independence and the output marginal are derived, not supplied.



Specializing the arbitrary position law to the normalized Gibbs target proves thermostat invariance of \(\Pi\). Exact OU evolution chooses \(a=e^{-\gamma h}\) and \(b=\sqrt{1-e^{-2\gamma h}}\); an OBABO half-step uses the corresponding half-interval coefficients. The theorem admits both choices through the explicit variance balance.



The thermostat preserves a standard Gaussian velocity by the fluctuation–dissipation identity and leaves position unchanged. Independence of the fresh refresh noise from the full initial phase ensures that the resulting velocity remains independent of the position.

\subsection{\texorpdfstring{Recover the two equilibrium marginals exactly}{Recover the two equilibrium marginals exactly}}


Projecting the independent product \(\varphi\otimes\pi\) to position integrates out a unit-mass Gaussian and leaves \(\pi\). Projecting to velocity integrates out the normalized position measure and leaves \(\varphi\). The target has total mass one, so applying the two product-projection identities removes the unit total-mass factors. These equalities supply the target position marginal and Gaussian velocity marginal used by the later L\ensuremath{{}^{2}} coupling theorem once stationarity of the full kinetic dynamics is proved.



The first marginal of the product is the Gaussian velocity law and the second is the normalized position Gibbs law. These exact marginals supply the stationary velocity moment and the desired sampling reference simultaneously. They are consequences of the concrete target construction.

\section{\texorpdfstring{Bounded tests of the independent-\allowbreak{}noise Markov chain}{Bounded tests of the independent-noise Markov chain}}
\label{proof-module-KineticLangevin-MarkovTests-lean}


To prove that a continuous dynamics preserves its Gibbs distribution, approximate it by concrete random steps. A small one-step weak error at the Gibbs measure must be propagated through those steps. The dual transition integrates the next test value against the independent fresh noise. Its measurability, uniform bound and exact relation to the forward law are proved here.



Fresh-noise independence identifies explicit random recursions with iterated transition operators. Probability averaging preserves test bounds, making one-step weak errors accumulate linearly through the chain.

\subsection{\texorpdfstring{Integrate the next test against the fresh noise}{Integrate the next test against the fresh noise}}


For a measurable step map \(S(x,\xi)\) and noise probability \(Q\), define \(PF(x)=\int F(S(x,\xi))Q(d\xi)\). The measurable parameter-integral theorem proves that \(PF\) is measurable. If \(\lVert F\rVert\le C\) everywhere, the norm-of-integral inequality gives \(\lVert PF(x)\rVert\le\int C\,dQ=C\). Independence is encoded by the product measure used by the Markov update.



Conditioning on the present state, integrate a test of the next updated state against the independent fresh-noise law. The resulting function is the transition operator applied to that test. Its uniform bound is no larger than the original test bound.

\subsection{\texorpdfstring{Prove the duality identity for the defined recursion}{Prove the duality identity for the defined recursion}}


The law update is the pushforward of \(\mu\otimes Q\) by \(S\). The pushforward integration formula and Fubini give \(\int F\,d(\mu P)=\int PF\,d\mu\). Boundedness on a probability space supplies integrability. Induction proves that every \(P^nF\) remains measurable and bounded by the original \(C\), and that \(\int F\,d(\mu P^n)=\int P^nF\,d\mu\). Thus the dual operator describes the law recursion rather than an assumed sequence of measures.



Independence of the next innovation from the previous history gives the one-step conditional expectation identity. Iterating it identifies expectations of the explicit recursion with iterated transition operators. This is the dual relation between the pathwise chain and its Markov laws.

\subsection{\texorpdfstring{Accumulate one-\allowbreak{}step weak errors without losing the test bound}{Accumulate one-step weak errors without losing the test bound}}


Suppose every measurable test bounded by \(C\) has one-step defect at \(\pi\) at most \(\delta\). Apply this estimate to \(P^nF\), whose bound is still \(C\). Insert and subtract \(\int P^nF\,d\pi\), apply the triangle inequality, and use the induction hypothesis. The resulting estimate is \(\left|\int F\,d(\pi P^n)-\int F\,d\pi\right|\le n\delta\). No exact invariance is assumed. In the splitting application the rejection bound is quadratic in the time step, so \(n\delta\to0\) on a fixed time interval.



Telescope the difference between numerical and reference test expectations over the number of steps. Each propagated test retains its uniform bound because transition operators average probability measures. Consequently a uniform one-step weak error accumulates at most linearly in the step count.

\section{\texorpdfstring{Quadratic stationary weak error for the kinetic splitting}{Quadratic stationary weak error for the kinetic splitting}}
\label{proof-module-KineticLangevin-SplittingTarget-lean}


This section applies the exact corrected Verlet identity and the exact Gaussian thermostat identity to the normalized Gaussian-times-Gibbs target. Averaging a bounded test over thermostat preserves its bound. The uncorrected splitting therefore has a proved quadratic one-step stationary weak defect, and its Markov iterate accumulates at most a linear number of these defects.



Corrected Hamiltonian invariance and exact thermostat invariance yield a vanishing weak stationary error for the splitting. Its normalized moment envelope keeps every local error integrable.

\subsection{\texorpdfstring{Use the thermal Gaussian innovation measure}{Use the thermal Gaussian innovation measure}}


The additive innovation measure is the pushforward of the Euclidean standard Gaussian by the coefficient \(\sigma_h=\sqrt{1-e^{-2\gamma h}}\). It is a probability measure because it is a measurable pushforward of a probability measure. The preceding OU marginal theorem identifies this exact law with the canonical convolution innovation; here it is the concrete measure used by the Markov transition.



The innovation law is the exact Gaussian with thermostat’s thermal variance. It is independent of the initial phase. This independence and the fluctuation–dissipation relation are what preserve the Gaussian velocity factor, rather than a covariance calculation alone.

\subsection{\texorpdfstring{Normalize the exact corrected Verlet identity}{Normalize the exact corrected Verlet identity}}


The canonical target is an explicitly proved finite scalar normalization of the Hamiltonian weighted measure. A density-weighted measure, its measurable pushforward, and their sum all commute with this scalar multiplication. Scaling the earlier exact identity therefore gives \((V_h)_\#(a_h\pi)+R_\#((1-a_h)\pi)=\pi\) for the normalized target. The argument rewrites the explicit normalization, both density scalings, both pushforward scalings, and the corrected Hamiltonian Gibbs identity.



The corrected Verlet identity initially concerns an unnormalized Hamiltonian density. Its finite positive partition mass permits division to obtain an identity for the probability target. Every subsequent weak-test comparison therefore uses normalized measures.

\subsection{\texorpdfstring{Derive the normalized target phase moment}{Derive the normalized target phase moment}}


The unnormalized Hamiltonian measure already has integrable \(\|v\|^2+\|\nabla f(x)\|^2\) from first-order strong convexity and Lipschitz force growth. The explicit normalizing scalar is finite: its Gaussian factor is an ordinary finite real number and its inverse potential partition is finite because the partition is strictly positive. The argument checks these extended-real facts and applies integrability under finite scalar multiplication. Thus the rejection envelope remains an finite target moment.



The normalized quadratic moment is computed from the independent Gaussian velocity and Gibbs position factors. Smoothness transfers the position moment to a force moment. This yields a finite constant for integrating the numerical rejection envelope.

\subsection{\texorpdfstring{Preserve the target with the additive OU law}{Preserve the target with the additive OU law}}


The damping \(\delta_h=e^{-\gamma h}\) and innovation scale satisfy \(\delta_h^2+\sigma_h^2=1\). The already proved thermostat result preserves the target when driven by a fresh standard Gaussian. Mapping that Gaussian through \(\sigma_h\) changes its law to the explicit additive innovation measure. One proves the coefficient identity, maps the independent product source through this noise scaling, composes the two pushforwards, and obtains exact preservation by the additive OU transition.



The thermostat acts only on velocity and adds the independent filtered Gaussian response. The exact variance identity preserves the velocity Gaussian, while position is unchanged. Hence the concrete product target is preserved by this transition.

\subsection{\texorpdfstring{Apply corrected invariance to the thermostat-\allowbreak{}averaged test}{Apply corrected invariance to the thermostat-averaged test}}


For \(\|F\|\le C\), define \(G(z)=\mathbb E_N[F(O_h(z,N))]\). The noise integral makes \(G\) measurable with \(\|G\|\le C\). Exact thermostat invariance gives \(\int G\,d\pi=\int F\,d\pi\), while the split transition has expectation \(\int G(V_hz)\,d\pi\). Apply the proved corrected-Verlet weak estimate to \(G\), then integrate the quadratic rejection envelope. This gives \[\left\|\int F\,d(P_h\pi)-\int F\,d\pi\right\|\le2Ch^2\frac{M(2M+11)}8\int(\|v\|^2+\|\nabla f(x)\|^2)\,d\pi.\] The argument derives both expectation identities from the Markov-step integral formula, checks measurability and boundedness of the averaged test, and composes the rejection and target moment estimates.



Average a bounded test over thermostat noise before applying the corrected Hamiltonian identity. Probability averaging preserves its uniform bound. This permits the two invariant building blocks to be combined while retaining a controlled weak error for the uncorrected splitting.

\subsection{\texorpdfstring{Accumulate only a linear number of local defects}{Accumulate only a linear number of local defects}}


The generic Markov test theorem applies the one-step estimate to each iterated dual test, whose bound remains \(C\), and telescopes their expectations under the fixed target. Consequently the \(n\)-step splitting law has weak error at most \(n\) times the displayed quadratic defect. The argument verifies nonnegativity of the target moment and all scalar coefficients, supplies the concrete measurable transition, and instantiates the proved generic theorem. With \(h=T/n\), this error tends to zero.



Over a fixed horizon the number of steps is proportional to the inverse step length. A local weak discrepancy with a higher step power therefore has a vanishing accumulated total. The telescoping argument keeps this accumulation linear and avoids a spurious growth of the test bound.

\section{\texorpdfstring{Exact kinetic Gibbs stationarity from the initial/\allowbreak{}noise laws}{Exact kinetic Gibbs stationarity from the initial/noise laws}}
\label{proof-module-KineticLangevin-StationarityTransfer-lean}


The stationarity argument starts with a continuous L\ensuremath{{}^{2}} kinetic integral solution whose joint input law is the product of the target initial phase law and the complete canonical Ornstein–Uhlenbeck innovation stream. The independent product construction supplies this joint law. Convergence of splitting approximations, vanishing stationary weak defects and equality of the limiting endpoint laws then prove invariance.



The splitting approximants have a vanishing target weak discrepancy and converge to the constructed kinetic solution. Their common limit identifies the solution’s full Gibbs phase law at every time.

\subsection{\texorpdfstring{Derive exact stationarity through the splitting limit}{Derive exact stationarity through the splitting limit}}


Let \(\pi\) be the normalized Gaussian-times-Gibbs target, and let \(Y\) solve the kinetic integral equation driven by the continuous L\ensuremath{{}^{2}} input \(B\). On a mesh \(h_n=T/(n+1)\), form the concrete Verlet–OU approximation \(S_n\) using the filtered increments of this same \(B\). The proved Banach splitting theorem gives \(S_n\to Y_T\) in the product of L\ensuremath{{}^{2}} spaces.



The full initial/\allowbreak{}innovation joint law is the explicit product \(\pi\otimes Q_{h_n}^{\otimes\mathbb N}\). The source-law transfer theorem, together with the almost-everywhere vector-chain representation, identifies the law of \(S_n\) with the defined Markov iterate. Its bounded-test error is at most \[2(n+1)h_n^2 K A=\frac{2T^2KA}{n+1},\qquad K=\frac{M(2M+11)}8,\quad A=\int(\|v\|^2+\|\nabla f(x)\|^2)\,d\pi.\] The moment is finite by the proved target tail theorem. The mesh is eventually at most one, so every quadratic rejection estimate applies, and the displayed scalar error tends to zero.



The phase-law continuity theorem now identifies the L\ensuremath{{}^{2}} limiting phase law with \(\pi\). The argument constructs the true approximation sequence, applies the proved convergence theorem, rewrites its chain law using the explicit initial/\allowbreak{}noise product distribution, rewrites the bounded test expectation through the pushforward, invokes the proved Markov weak-error bound, and checks the scalar \(T^2/(n+1)\) decay. It then applies dual-characteristic uniqueness for finite measures. No stationary solution, local-error conclusion, output-law independence field, or numerical-limit conclusion is assumed. The remaining input-law hypothesis is a concrete source probability identity, supplied by the canonical independent-product construction.



The constructed splitting chain has the exact independent thermal-noise law and its endpoint converges in second mean to the kinetic solution. Its bounded-test target discrepancy vanishes along the same refinement. Characteristic functions therefore identify the limiting phase law with the Gibbs target, proving stationarity of the continuous solution.

\section{\texorpdfstring{The complete canonical KLMC theorem on arbitrary finite coordinate sets}{The complete canonical KLMC theorem on arbitrary finite coordinate sets}}
\label{proof-module-KineticLangevin-GenericCanonicalKLMC-lean}


The entire actual-source KLMC proof is stated for a general finite coordinate type \(I\). This includes the particle index \(\{\theta\}\sqcup(\{1,\ldots,N\}\times\{x\})\) without a separate coordinate-isometry transport. All initial couplings, canonical Brownian inputs, selected reference flows, Gibbs stationary laws and Wasserstein infima are the same concrete constructions as in the finite-dimension presentation; the dimension in the final bias is exactly \(|I|\).



The canonical proof is carried out on a finite coordinate set without choosing a special enumeration. It constructs the sampler, stationary reference, local comparisons, and true transport infimum in one consistent argument.

\subsection{\texorpdfstring{Specify the concrete probability space and source projection}{Specify the concrete probability space and source projection}}


The phase variable is velocity first. The sample space contains the two coupled initial positions, one shared Gaussian velocity and one vector Brownian source. The measure is the initial-coupling product extension; probability follows from its first marginal and from the Gaussian/\allowbreak{}Brownian probability factors. The Brownian-source projection preserves its measure. Pullback along that projection gives a continuous linear map on L\ensuremath{{}^{2}}, and its representative theorem identifies it almost everywhere with evaluation on the product sample.



Let the finite coordinate set be \(I\), so the physical dimension is \(|I|\). Extending the initial plan by one independent common source preserves both initial position marginals. The same Gaussian velocity and Brownian source can then drive the two coupled phase paths.

\subsection{\texorpdfstring{Construct the common Brownian input and filtered innovation}{Construct the common Brownian input and filtered innovation}}


The scaled source lift includes the physical amplitude \(\sqrt{2\gamma}\). At iteration \(n\), the input is the vector Brownian increment on \((nh,nh+t]\), and the filtered input is the exponential response to that same increment. The unscaled innovation pair consists of the exact velocity and position Wiener responses at the step endpoint. Both input paths are continuous in L\ensuremath{{}^{2}}, and the additive input starts at zero. These follow from the source lift's continuity and linearity and the proved vector-response properties.



The physical amplitude multiplies both the raw Brownian input and its filtered responses. Keeping it outside the linear response construction preserves all covariance identities. The velocity and position innovations therefore have exactly thermal normalization of the prescribed KLMC method.

\subsection{\texorpdfstring{Identify the exact noise equations and defined numerical step}{Identify the exact noise equations and defined numerical step}}


The constructed filtered input satisfies \(B_n(t)=W_n(t)-\gamma\int_0^tB_n(s)ds\). Its endpoint is the scaled velocity innovation, and its time primitive is the scaled positional innovation. The argument transports the vector Brownian equations through the continuous linear source lift and commutes that lift with Bochner integration. Induction on the recursions then proves that the exponential-input chain is exactly the KLMC sequence, with the canonical correlated innovation pair and the prescribed fluctuation-dissipation scaling.



The velocity response solves a damped integral equation and its time integral is the position response. At the endpoint these identities recover both innovation terms of the exponential step. Thus the constructed recursion is identified through equations, not only through matching marginal noise laws.

\subsection{\texorpdfstring{Connect Brownian L\ensuremath{{}^{2}} innovations with the generic output-\allowbreak{}law construction}{Connect Brownian L² innovations with the generic output-law construction}}


The L\ensuremath{{}^{2}} innovation pair has the same representatives as the pointwise canonical Brownian innovation. The proof pulls back the previously established vector-response representative theorem to the coupling product and combines the two component identities. Comparing with the generic product-law layer's representatives gives equality of the L\ensuremath{{}^{2}} pairs. Thus the numerical chain analyzed here has the coupling-independent output law already proved for the defined KLMC chain.



Equality almost everywhere of response representatives implies equality of their pushforward laws. This transfers the fixed-source output identification to the coupling extension. In particular, the analytic random variables do not define a new sampler depending on the chosen initial plan.

\subsection{\texorpdfstring{Build the numerical and stationary-\allowbreak{}reference initial states}{Build the numerical and stationary-reference initial states}}


The two initial phases share exactly one Gaussian velocity variable. The numerical phase uses the first coupled position; the reference phase uses the second. The reference representative is proved to be the explicit pair \((V,X_\mathrm{ref})\). Product independence and the two marginal laws prove that its law is \(\mathcal N(0,I_{|I|})\otimes\pi\). The finite target second moment and finite coupling cost supply the required L\ensuremath{{}^{2}} initial positions.



The shared initial velocity makes the initial phase discrepancy purely positional. Its standard Gaussian law and independence from the reference position give the reference product target. Finite coupling cost ensures the numerical initial position also has the required second moment.

\subsection{\texorpdfstring{Construct all three chains and segment families}{Construct all three chains and segment families}}


The numerical chain applies the lifted original gradient and the constructed correlated innovations. The reference chain stitches exact selected kinetic segments, each driven by the canonical Brownian increment for that iteration. The corresponding reference segments start from these reference states. The auxiliary exact segment starts instead from the numerical state and receives the identical additive input. The drift is the pointwise gradient lifted to L\ensuremath{{}^{2}}; no substitute potential on the L\ensuremath{{}^{2}} space is assumed.



The reference exact path is restarted from the reference phase, while the auxiliary exact path restarts from the numerical phase. Both use the same interval input. This arrangement separates contraction of the old error from the force-freezing defect of the new numerical step.

\subsection{\texorpdfstring{Prove every analytic premise of the KLMC coupling theorem}{Prove every analytic premise of the KLMC coupling theorem}}


The bundled specification supplies continuity of auxiliary paths, reference paths and filtered inputs; exact initial states; reference endpoint matching; the defined exponential update; both common-input integral equations; the exponential mild identity; the kinematic primitive identity; and equal initial velocities. Each field follows from the concrete definitions and earlier proved equations. The mild and kinematic fields are derived internally. The only remaining probabilistic input for a final stationary-error theorem is the stationary law of the global reference diffusion, proved separately and connected to these stitched segments by uniqueness.



Continuity, endpoint matching, the mild formula, and the position primitive are all consequences of the solution specification. The one-step theorem therefore receives verified path identities. Its eventual stationary velocity radius will come from the separate global-reference stationarity argument.

\subsection{\texorpdfstring{Construct the full canonical input and identify its increments}{Construct the full canonical input and identify its increments}}


The full input is the scalar canonical Brownian value assembled spatially and pulled onto the initial-coupling extension with amplitude \(\sqrt{2\gamma}\). It is continuous in L\ensuremath{{}^{2}} because each assembly and pullback is continuous linear. The Brownian value at zero vanishes: The argument verifies that its deterministic time indicator has empty support, then uses linearity of the isonormal map. For nonnegative local time, the constructed segment input equals the full input difference at \(nh+t\) and \(nh\). For nonnegative local times, linearity of spatial assembly and source pullback preserves this increment identity.



The full canonical input and each interval input are two descriptions of the same Brownian source. Linearity gives their increment identity exactly. This is the compatibility condition needed to identify stitched local reference segments with one continuous global reference.

\subsection{\texorpdfstring{Construct the full reference solution}{Construct the full reference solution}}


The global reference is the selected kinetic solution with the lifted gradient, the canonical reference initial phase and the full canonical Brownian input. Its specification proves continuity, its exact starting phase, and the kinetic integral equation over the whole finite horizon. There is no stationary-law requirement in its definition.



Solve the kinetic equation on the complete horizon from the reference initial phase. The global equation uses the same lifted gradient as every local segment. Existence and uniqueness are established before using any claimed invariant law.

\subsection{\texorpdfstring{Identify every stitched reference point by uniqueness}{Identify every stitched reference point by uniqueness}}


The generic stitching theorem now applies because the local inputs are proved to be increments of the full input. Consequently \(R_n=Y_{nh}\), and \(Z_n(t)=Y_{nh+t}\) throughout each step whose endpoint lies in the global horizon. The argument instantiates the existing exact uniqueness results with the L\ensuremath{{}^{2}} gradient and checks the horizon arithmetic. Thus a stationary law for the full constructed kinetic solution supplies the law of all reference states used by the quantitative KLMC proof.



At the first grid point both reference constructions start identically. Induction and local uniqueness then identify each later endpoint; uniqueness on each entire interval also identifies intermediate times. This transfers global laws to every local time in the error integrals.

\subsection{\texorpdfstring{Project the coupling source onto an independent phase/\allowbreak{}Brownian product}{Project the coupling source onto an independent phase/Brownian product}}


The projection sends the full sample to \(((V,X_\mathrm{ref}),\omega_\mathrm{Brownian})\). Its law is \((\mathcal N(0,I_{|I|})\otimes\pi)\otimes P_\mathrm{Brownian}\). First, the argument projects the initial coupling onto its second marginal, retaining the independent Gaussian/\allowbreak{}Brownian source. It then applies measure-preserving product associativity and swaps the position and velocity factors. Thus independence of the complete reference initial phase and the complete Brownian source is proved for the product construction.



Only the reference position marginal of the initial transport plan enters this projection. That marginal is the target, while the adjoined source remains independent. Reordering the product factors therefore yields the required joint Gibbs-phase and Brownian law.

\subsection{\texorpdfstring{Identify the full-\allowbreak{}input OU increment}{Identify the full-input OU increment}}


The deterministic integration-by-parts input used in the kinetic splitting proof is \(\Delta W-\gamma\int_0^he^{-\gamma(h-s)}\Delta W_sds\). The local-input increment identity replaces each difference of the full canonical input by its constructed segment input. Applying the vector exponential-response formula through the continuous linear source lift proves that this OU increment is exactly the filtered velocity response. Its representative is consequently the scaled canonical vector response on the original Brownian source, pulled onto the coupling extension.



The filtered increment can be written entirely using the continuous full input and a deterministic exponential convolution. Substituting the proved local increment identity turns it into the canonical segment response. Hence the stationarity splitting uses the source of the reference equation.

\subsection{\texorpdfstring{Derive the full independent thermal Gaussian joint law}{Derive the full independent thermal Gaussian joint law}}


The vector OU stream already has the exact iid Gaussian law with scale \(\sqrt{1-\psi_0(h)^2}\). The product-source joint-law theorem combines this stream with the initial phase target. Composition with the proved phase/\allowbreak{}Brownian projection law transfers that product law to the coupling extension. Finally, the representative identities identify its initial phase with the constructed reference initial L\ensuremath{{}^{2}} phase and its stream with the full-input OU increments. The resulting joint law is precisely the source obligation of the proved kinetic Gibbs stationarity transfer theorem.



The reference initial phase is independent of the entire filtered stream, not just each block separately. The stream’s finite cylinder laws are products of thermal Gaussian law. This full joint law is precisely what repeated independent thermostat updates require.

\subsection{\texorpdfstring{Fix the numerical output law}{Fix the numerical output law}}


The definition uses the original pointwise KLMC method, the standard Gaussian initial velocity and the exact canonical correlated Brownian innovations. The numerical L\ensuremath{{}^{2}} chain is identified with the generic product-law construction by equality of the innovation variables. The established output-law theorem then proves that its pushforward is independent of which initial position coupling was chosen.



The numerical endpoint map uses only its own initial position and the fixed source. Projection of any extended initial plan recovers this same input law. Thus every coupling proof estimates the same output measure \(\nu_n\).

\subsection{\texorpdfstring{Transfer one global stationary phase law to the reference segments}{Transfer one global stationary phase law to the reference segments}}


The named stationary property concerns the full constructed kinetic flow: at every time its phase law is \(\mathcal N(0,I_{|I|})\otimes\pi\). It contains no numerical convergence conclusion. Exact chain-to-global identification transfers its position marginal to every reference grid point. Exact segment-to-global identification transfers its Gaussian velocity marginal to every intermediate time in a reference segment. The argument obtains these marginals by composition with the product-measure projections.



Global stationarity gives both position and velocity marginals at all times. The former identifies the reference endpoint as the sampling target; the latter supplies \(\sqrt{|I|}\) throughout the local-error integral. These roles use different projections of one proved phase-law identity.

\subsection{\texorpdfstring{Apply the proved coupling estimate to one initial plan}{Apply the proved coupling estimate to one initial plan}}


The initial position L\ensuremath{{}^{2}} discrepancy equals the true cost of the supplied transport coupling; this follows from the concrete extension's representative and norm identities. The segment specification furnishes all analytic local-error premises internally. Together with the derived stationary position and Gaussian velocity marginals, the quantitative L\ensuremath{{}^{2}} KLMC theorem gives \(W_2(\nu_n,\pi)\le\sqrt2(1-3mh/(4\gamma))^n\operatorname{cost}(c)+Mh\sqrt{2|I|}/m\). Output-law equality and initial-cost equality rewrite this into a bound for the numerical law.



The supplied initial plan gives an coupling of the numerical endpoint with the stationary reference endpoint. Its root-mean-square discrepancy bounds Wasserstein distance. At this stage the right-hand side still contains the cost of that particular initial plan.

\subsection{\texorpdfstring{Prove output finiteness and take the transport infimum}{Prove output finiteness and take the transport infimum}}


Finite initial Wasserstein distance supplies an finite-cost initial coupling. Its numerical and reference states are L\ensuremath{{}^{2}} variables, so their distance gives a finite transport bound on the output law. This proves output finiteness before real-valued transport arithmetic is used. The scalar recurrence estimate proves the contraction factor is nonnegative under the paper's step restrictions. The argument then applies the proved infimum theorem over all finite-cost couplings, obtaining the same estimate with \(W_2(\mu,\pi)\) in place of the cost of a chosen plan. No existence of an optimal coupling is assumed.



At least one finite-cost plan exists because the initial Wasserstein distance is finite. It proves finiteness of the output distance. Since every finite-cost plan obeys the same affine bound with a nonnegative coefficient, the infimum gives the exact initial Wasserstein term.

\subsection{\texorpdfstring{Prove Gibbs stationarity of the global coupling reference}{Prove Gibbs stationarity of the global coupling reference}}


The reference begins with standard Gaussian velocity independent of the target Gibbs position. Its full Brownian forcing and its entire filtered innovation stream are constructed objects on the initial-coupling extension. Their exact initial/\allowbreak{}noise product law is already proved. The kinetic stationarity transfer theorem therefore applies directly to the integral solution. For a time \(t\) inside a larger horizon, The proof restricts the integral equation to \([0,t]\), supplies the continuous input and its zero starting value, and applies the proved splitting-limit stationarity result. This yields the full phase law \(\mathcal N(0,I_{|I|})\otimes\pi\) at every reference time, without a remaining stationarity or Gaussian-marginal hypothesis.



The full initial/\allowbreak{}noise product law and canonical splitting convergence discharge the stationarity-transfer hypotheses. Restricting a longer solution to a smaller horizon preserves its integral equation. Thus every reference time has the target phase law, with no remaining assumed stationary process.

\subsection{\texorpdfstring{Conclude the source-\allowbreak{}paper KLMC Wasserstein estimate}{Conclude the source-paper KLMC Wasserstein estimate}}


Let \(f\) be differentiable, \(m\)-strongly convex with \(m>0\), and have an \(M\)-Lipschitz gradient. If \(\gamma^2\ge m+M\), \(\gamma>0\), and \(0<h\le m/(4\gamma M)\), the KLMC position law satisfies \[W_2(\nu_n,\pi)\le\sqrt2\left(1-\frac{3mh}{4\gamma}\right)^nW_2(\mu,\pi)+\frac{Mh\sqrt{2|I|}}m.\] The initial position distribution may be arbitrary with finite initial Wasserstein distance; its velocity is the explicit independent standard Gaussian. First, the argument derives the finite Gibbs second moment from strong convexity. It then supplies the just-proved global-reference stationarity for every initial transport coupling to the complete canonical coupling/\allowbreak{}infimum theorem. The conclusion has no conditional local-error, reference-process, Gaussian-law, stationary-law or optimal-coupling premise. The smoothness assumptions are expressed by the gradient and strong convexity, so theorem also covers the paper's Hessian-bounded setting through the existing deterministic geometry results.



The transient factor is geometric and the persistent floor is linear in the step. The latter arises from a quadratic local freezing defect divided by a first-order contraction gap. All coefficients retain the dimension \(|I|\), so theorem applies directly to particle-scaled Euclidean spaces.

\section{\texorpdfstring{KIPLMC1: the unconditional parameter RMS bound for the Gaussian sampler}{KIPLMC1: the unconditional parameter RMS bound for the Gaussian sampler}}
\label{proof-module-KIPLMC-CanonicalExponential-lean}


The KIPLMC1 Gaussian Markov laws are identified with the shared canonical KLMC laws at the changed clock. The shared kinetic theorem constructs its Brownian reference, proves exact Gibbs stationarity, controls the real local error, and takes the true coupling infimum. The particle geometry and sharp parameter concentration then give the paper’s three explicit error terms.



The canonical shared theorem applies to the sum potential and correlated-Gaussian chain. Exact clock cancellation and sharp parameter concentration yield the three explicit physical-time KIPLMC1 error contributions.

\subsection{\texorpdfstring{Identify the particle noise source with the Gaussian Markov law}{Identify the particle noise source with the Gaussian Markov law}}


The particle coordinate set is \(\mathrm{Fin}(d_\theta)\sqcup(\mathrm{Fin}(N)\times\mathrm{Fin}(d_x))\). The shared kinetic construction now works directly on this arbitrary finite set. Its complete canonical innovation stream has the proved product law of independent spatial blocks, each with the exact correlated two-component Gaussian distribution. The argument instantiates the product-source Markov-law theorem with this Brownian source, its whole-stream law, and an independent standard Gaussian initial velocity. This identifies the concrete canonical kinetic position law with the position marginal of the Gaussian Markov iterate. The previously particle clock theorem then gives the exact identity between the physical-time KIPLMC1 output and this same canonical output for \(V_N\) at friction \(\sqrt N\gamma\) and step \(h/\sqrt N\). The position target is unchanged.



The finite particle coordinate set has cardinality \(d_\theta+Nd_x\), even when the parameter and latent dimensions differ. The generic shared construction applies directly to that set. Its Brownian response law therefore matches the particle Gaussian transition without an artificial dimensional identification.

\subsection{\texorpdfstring{Check the transformed assumptions and all particle-\allowbreak{}count cancellations}{Check the transformed assumptions and all particle-count cancellations}}


In scaled particle coordinates, the sum potential \(V_N\) has strong convexity \(Nm\) and gradient Lipschitz constant \(NL\). Put \(c=\sqrt N\), \(\gamma^\prime=c\gamma\), and \(h^\prime=h/c\). The original restrictions \(m+L\le\gamma^2\) and \(4\gamma Lh\le m\) become exactly \[Nm+NL\le(\gamma^\prime)^2,\qquad4\gamma^\prime(NL)h^\prime\le Nm.\] One proves positivity, curvature compatibility, friction, and step conditions together using \(c^2=N\), then proves the exact rate identity \[1-\frac{3(Nm)h^\prime}{4\gamma^\prime}=1-\frac{3mh}{4\gamma}.\] For \(D=d_\theta+Nd_x\), the shared bias simplifies exactly to \[\frac{NLh^\prime\sqrt{2D}}{Nm}=\frac{Lh}{m}\sqrt{\frac{2D}{N}}.\] The scalar proof checks the square-root quotient with nonnegative dimensions and the positive particle count. This cancellation explains why the first algorithm’s error contains \(d_x+d_\theta/N\), with no additional growing power of particle count.



Both curvature constants scale by \(N\), while friction and step scale by reciprocal square-root particle factors. The transformed rate cancels these factors exactly. The bias retains only thermal radius \(\sqrt{d_x+d_\theta/N}\), which explains the favorable particle dependence.

\subsection{\texorpdfstring{Derive finite transport moments for the output}{Derive finite transport moments for the output}}


The initial position law has the stated finite second moment. Strong convexity proves the corresponding moment of the normalized Gibbs target; their independent product therefore supplies a finite transport coupling. The shared canonical construction places this coupling on the same Gaussian/\allowbreak{}Brownian source as the chain. Its numerical and exact reference positions are L\ensuremath{{}^{2}} random variables, and exact reference stationarity has already been proved from the splitting limit. Their L\ensuremath{{}^{2}} discrepancy supplies a finite output transport coupling. The proof transfers this finiteness through the exact particle law identity. Thus the later real-valued Wasserstein and parameter RMS expressions are used only after their finiteness has been proved.



Finite initialization and target moments construct a finite-cost coupling. The shared canonical chain and stationary reference then give finite output transport on the source. This validates every later real-valued transport and parameter-error integral.

\subsection{\texorpdfstring{Apply the shared kinetic theorem to the particle sampler}{Apply the shared kinetic theorem to the particle sampler}}


The shared theorem is applied to the sum potential, using its derived \(Nm\)-strong convexity and \(NL\)-smoothness. It proves, rather than assumes, the stationary Brownian reference, source independence, mild identities, local exponential-integrator error, scalar recurrence, and transport infimum. The argument supplies the transformed assumptions proved above, rewrites the output through the exact clock/\allowbreak{}Gaussian-law identity, counts the finite particle coordinates, and simplifies the rate and bias. The resulting position-law estimate is \[W_2(\nu_n,\pi_N)\le\sqrt2\left(1-\frac{3mh}{4\gamma}\right)^nW_2(\nu,\pi_N)+\frac{Lh\sqrt2}{m}\sqrt{\frac{d_\theta+Nd_x}{N}}.\] The initialization in this theorem is explicitly the independent thermal velocity law times the supplied position law; all algorithmic Gaussian correlations are those of the constructed noise measure.



The shared theorem is applied only after the exact noise law, whole-chain conjugacy, and transformed assumptions have been established. Rewriting its conclusion gives the physical-time particle sampling bound. Every term refers to the independent correlated-Gaussian algorithm.

\subsection{\texorpdfstring{Project to the parameter and add the sharp equilibrium concentration}{Project to the parameter and add the sharp equilibrium concentration}}


The parameter projection is a continuous linear map with norm at most one. The powered Gibbs parameter marginal satisfies the proved sharp concentration bound \(\mathbb E_{\pi_N}\|\theta-\theta_\star\|^2\le d_\theta/(Nm)\). The transport-to-RMS inequality therefore gives \[\bigl(\mathbb E\|\theta_n-\theta_\star\|^2\bigr)^{1/2}\le\sqrt2\left(1-\frac{3mh}{4\gamma}\right)^nW_2(\nu,\pi_N)+\frac{Lh\sqrt2}{m}\sqrt{\frac{d_\theta+Nd_x}{N}}+\sqrt{\frac{d_\theta}{Nm}}.\] The argument invokes the output finiteness theorem, the sharp equilibrium concentration theorem, and the nonexpansive projection theorem. It rewrites the finite RMS as the displayed square root of the Bochner integral and substitutes the proved Gaussian position-law estimate. These are the initialization, discretization, and finite-particle concentration errors in the paper’s theorem.



The parameter projection is nonexpansive, so it cannot increase the transport discrepancy. Minkowski’s inequality adds the equilibrium parameter radius around the maximizer. This decomposes the parameter root-mean-square error into transient, discretization, and statistical contributions.

\subsection{\texorpdfstring{Recover the paper’s independent-\allowbreak{}initialization constant}{Recover the paper’s independent-initialization constant}}


The true initial Wasserstein distance is bounded by every initial coupling cost. In particular, independent \(Z_0\sim\nu\) and \(\bar Z_\star\sim\pi_N\) give \[W_2(\nu,\pi_N)\le\bigl(\mathbb E\|Z_0-\bar Z_\star\|^2\bigr)^{1/2}.\] The argument constructs this independent product coupling, derives the integrability of its discrepancy from the two proved second moments, and rewrites its finite L\ensuremath{{}^{2}} cost as the square-root integral. It also proves that the stated step restrictions make the geometric rate nonnegative. Multiplying the initial coupling inequality by this nonnegative rate gives exactly the paper’s printed constant \(C_1\); the other two constants are unchanged. Thus the true-transport statement and the printed independent-initialization statement are both consequences for the same sampler.



The independent initial product coupling is one admissible plan, so its root-mean-square cost bounds the infimum defining Wasserstein distance. A nonnegative transient coefficient preserves that comparison. The printed initialization constant is therefore a corollary of the stronger true-transport estimate.

\section{\texorpdfstring{Actual kinetic flows under measure-\allowbreak{}preserving source changes}{Actual kinetic flows under measure-preserving source changes}}
\label{proof-module-KineticLangevin-FlowPullback-lean}


Changing a probability source by a measure-preserving projection preserves the random variables and their output laws. This section proves that the lifted nonlinear drift and the selected kinetic integral solution commute with that change of source. It supplies coupling-independent continuous-flow laws without assuming samplewise continuity of an arbitrary L\ensuremath{{}^{2}} representative.



Measure-preserving source changes preserve the kinetic integral equation because pullback is bounded linear and commutes with the pointwise drift. Uniqueness then identifies solutions constructed on different coupling spaces.

\subsection{\texorpdfstring{Define continuous linear source and phase pullbacks}{Define continuous linear source and phase pullbacks}}


Measure-preserving composition is a linear isometry on L\ensuremath{{}^{2}}. The source pullback is this isometry viewed as a continuous linear map, and the phase pullback applies it separately to velocity and position. Their representative theorems identify the resulting random variables with composition. Pushforward composition proves preservation of every output law. A further representative argument proves that two successive pullbacks equal pullback along the composed projection.



Changing probability source by a measure-preserving map defines an isometry on square-integrable variables. Apply it separately to velocity and position to obtain the phase pullback. Bounded linearity preserves continuity and commutes with the time integrals in the kinetic equation.

\subsection{\texorpdfstring{Commute the nonlinear drift with pullback}{Commute the nonlinear drift with pullback}}


The lifted vector field is pointwise application of the original Lipschitz field. The proof compares its representatives before and after source projection, obtaining exact equality in L\ensuremath{{}^{2}}. Combining this with linearity proves the kinetic drift commutes with the phase pullback, including both friction and force terms and the position equation.



The lifted drift is defined pointwise from the original gradient. Composing a representative with the source map therefore commutes with evaluating that gradient. Together with the affine phase terms, this proves that pullback preserves the nonlinear kinetic drift.

\subsection{\texorpdfstring{Identify the selected integral solutions}{Identify the selected integral solutions}}


Apply the phase pullback to the original kinetic integral equation. Continuity and linearity commute it with the Bochner time integral, and the drift commutation identity produces the target-space kinetic equation. The transformed initial condition and additive input are exact. The proved integral-solution uniqueness theorem therefore identifies the selected target-space solution with the pulled original solution at every time in the horizon. This equality transfers laws between an independent initial-state/\allowbreak{}noise product and any larger initial-coupling extension.



The reusable integral-equation theorem proves this transformation for any continuous kinetic solution before specializing to selected flows.



Pulling back the solution equation gives a solution with the pulled-back initial phase and forcing. Uniqueness identifies it with the independently selected solution on the new source. This validates the coupling-space comparison of flows defined on different product spaces.

\section{\texorpdfstring{The arbitrary-\allowbreak{}shift kinetic contraction for L\ensuremath{{}^{2}} random vectors}{The arbitrary-shift kinetic contraction for L² random vectors}}
\label{proof-module-KineticLangevin-LpGeneralContraction-lean}


The continuous mixing theorem permits an arbitrary shifted kinetic metric, not only the convenient zero-shift case used by KLMC. The pointwise strong-smooth gradient sector bound already lifts to L\ensuremath{{}^{2}}. Here it gives the complete shifted energy and position estimates for shared-input L\ensuremath{{}^{2}} kinetic solutions.



The full arbitrary-shift transformed geometry extends to second-mean random paths. Its exact initialization and endpoint projection recover the same prefactor and exponent needed for general kinetic law contraction.

\subsection{\texorpdfstring{Prove the full shifted L\ensuremath{{}^{2}} energy decay}{Prove the full shifted L² energy decay}}


Let \(d=\gamma-2a>0\). The exact common-noise integral equations make the difference differentiable, with the original gradient difference lifted pointwise. The proof differentiates the explicit shifted energy and applies the general dissipativity inequality using the lifted sector bound. The resulting rate is \(2\max(a^2-m,M-(\gamma-a)^2)/d\). A right-derivative Grönwall argument integrates this differential inequality, including the starting endpoint. No potential on the L\ensuremath{{}^{2}} space is introduced.



The strong-smooth sector bound lifts by integration to the square-integrable Hilbert space. Applying the shifted energy calculation there yields the same exponent as for deterministic paths. Thus no dimension-dependent norm conversion is introduced by taking root-mean-square errors.

\subsection{\texorpdfstring{Project the shifted energy to the position error}{Project the shifted energy to the position error}}


Equal initial velocities make the initial energy equal to \(((\gamma-a)^2+a^2)\|X_0-Y_0\|_2^2\). The invertible transformed coordinates bound the position error by the energy divided by \((\gamma-2a)^2\). The proof combines these identities with the energy decay, proves the exact squared-exponential identity, and takes the nonnegative square-root inequality. The conclusion is the full source-paper prefactor \(\sqrt{2((\gamma-a)^2+a^2)}/(\gamma-2a)\) multiplied by the explicit exponential rate and the initial L\ensuremath{{}^{2}} position cost.



Equal initial velocities reduce the initial transformed energy to its position coefficient. The endpoint position projection then supplies the full shifted prefactor. The resulting estimate applies to common-input random solutions and can be pushed forward to their probability laws.

\section{\texorpdfstring{Actual continuous kinetic laws and the true Wasserstein infimum}{Actual continuous kinetic laws and the true Wasserstein infimum}}
\label{proof-module-KineticLangevin-ProductContinuousFlow-lean}


The initial position is independent of a common source carrying the initial velocity and continuous forcing. This construction makes the output law independent of any chosen coupling. We then couple two initial position laws, pull their kinetic solutions onto the same source, and apply the proved arbitrary-shift kinetic energy contraction. Taking the infimum over all initial couplings gives the Wasserstein estimate.



The continuous output laws are fixed independently of initial couplings. Source pullback and common-input contraction bound every coupling cost, producing the full arbitrary-shift diffusion Wasserstein theorem.

\subsection{\texorpdfstring{Construct the independent-\allowbreak{}source diffusion}{Construct the independent-source diffusion}}


The probability space is the product of the initial position law and the common source. The first L\ensuremath{{}^{2}} coordinate is the pulled-back velocity; the second is the original position coordinate. Pulling back the continuous additive forcing gives the kinetic input. The selected Picard solution is then pushed forward at the terminal time to define the position law. The solution specification proves continuity, the exact initial condition and both integral equations. Measurability makes the output a probability measure, and its L\ensuremath{{}^{2}} construction proves a finite second moment.



The output law is defined separately for each initial distribution using one fixed independent velocity-and-forcing source. The endpoint is a measurable function of those inputs. This makes the subsequent coupling argument a comparison of already defined physical output laws.

\subsection{\texorpdfstring{Put any initial coupling on the same source}{Put any initial coupling on the same source}}


Adjoin the same independent source to an arbitrary coupling of the two initial position distributions. The two projections recover the original independent-source spaces exactly, so they preserve measure. Composition of L\ensuremath{{}^{2}} pullbacks shows that both projected solutions receive exactly the same velocity and forcing. The representative formula identifies their initial positions with the two original coupling coordinates.



The initial plan may correlate the two starting positions arbitrarily. Adding one common independent source preserves both position marginals and supplies equal initial velocities and equal forcings. Its two natural projections recover the separate input spaces exactly.

\subsection{\texorpdfstring{Prove the bound before taking an infimum}{Prove the bound before taking an infimum}}


The source-pullback theorem transforms each solution equation to the coupling extension. Their common velocity and noise cancel, so the arbitrary-shift L\ensuremath{{}^{2}} contraction theorem applies. The initial position L\ensuremath{{}^{2}} norm equals the root-mean-square transport cost of the supplied coupling: the independent extra source contributes a factor one. At the endpoint, measure-preserving projection recovers each independently defined output law. The L\ensuremath{{}^{2}} transport inequality therefore gives the stated bound for this particular initial coupling.



Pullback of the solution equations puts both paths on the common extension. The pathwise contraction coefficient then bounds their second-mean endpoint difference. Product integration identifies their initial second-mean difference with the original transport-plan cost.

\subsection{\texorpdfstring{Pass to the Wasserstein distance}{Pass to the Wasserstein distance}}


Both initial laws and both output laws have finite second moments. The prefactor is nonnegative because the shifted damping denominator is positive. The proved affine bound over every finite-cost initial coupling therefore passes to the infimum defining Wasserstein distance. The result has the paper’s full arbitrary-shift exponential factor and holds for any common square-integrable velocity and any continuous common input, so Brownian forcing is a direct specialization.



The affine all-couplings bound has a nonnegative coefficient. Taking its infimum yields the Wasserstein contraction, without an optimizer assumption. The argument accommodates every square-integrable common velocity law and continuous square-integrable forcing source satisfying the construction hypotheses.

\section{\texorpdfstring{Theorem 1 for Brownian kinetic diffusions}{Theorem 1 for Brownian kinetic diffusions}}
\label{proof-module-KineticLangevin-CanonicalContinuousMixing-lean}


The continuous mixing theorem is instantiated with the existing vector Brownian source. Its initial velocity distribution is arbitrary, subject only to a finite second moment, and is independent of the initial positions and Brownian noise. The output measures are terminal pushforwards of the selected kinetic integral solution, on any finite coordinate set.



The independent-source law theorem is specialized to canonical Brownian forcing. It gives explicit position mixing for the constructed kinetic diffusion with any admissible square-integrable common velocity law.

\subsection{\texorpdfstring{Construct arbitrary initial velocity and Brownian forcing}{Construct arbitrary initial velocity and Brownian forcing}}


The source is the independent product of the prescribed velocity distribution and the canonical vector Brownian measure. Its first projection gives the L\ensuremath{{}^{2}} initial velocity. The second projection pulls back the canonical Brownian process, multiplied by the fluctuation-dissipation factor \(\sqrt{2\gamma}\). Continuity and the zero initial value follow from the corresponding proved Brownian facts and linear source pullback.



The Brownian factor supplies the physical forcing independently of the prescribed velocity law. Gaussianity of that velocity is unnecessary for the common-input difference estimate. Its finite second moment is enough to construct the phase path in the relevant Hilbert space.

\subsection{\texorpdfstring{Define the position laws and prove arbitrary-\allowbreak{}shift mixing}{Define the position laws and prove arbitrary-shift mixing}}


For each initial position distribution, the independent-source construction solves the kinetic equation with the original gradient and Brownian input. Its terminal position pushforward defines the law. Applying the general product-source coupling theorem proves the complete arbitrary-shift exponential coefficient. The common velocity law needs no Gaussian assumption. The initial transport term is the true Wasserstein distance between the initial position distributions, obtained by infimizing over couplings rather than assuming an optimizer or a prescribed coupling cost.



The general source-law theorem now specializes to the canonical Brownian input. The complete shifted prefactor remains visible, including any deterioration near the singular coordinate choice. Choosing the simple zero shift yields the familiar exponential kinetic position-mixing bound.

\section{\texorpdfstring{The continuous Gibbs target and explicit exponential mixing}{The continuous Gibbs target and explicit exponential mixing}}
\label{proof-module-KineticLangevin-ContinuousGibbs-lean}


The canonical Brownian diffusion starts on an independent position–velocity–noise product. This section proves its exact Gibbs stationarity on that same source, rather than assuming a reference law. Combining it with the arbitrary-shift continuous coupling theorem gives unconditional transport to the true Gibbs target. The simple friction regime yields the explicit factor \(\sqrt{2}e^{-mt/\gamma}\).



The constructed Brownian kinetic flow preserves the normalized Gibbs phase law. Using its position marginal as a stationary comparator turns law-to-law contraction into explicit convergence to the sampling target.

\subsection{\texorpdfstring{Reassociate the source and identify each thermal increment}{Reassociate the source and identify each thermal increment}}


The measurable source map sends \((x,(v,B))\) to \(((v,x),B)\). Product associativity and exchange of independent position and velocity prove that this map preserves the exact phase-times-Brownian law. The continuous linear source pullbacks commute with the OU input integral, and the proved Brownian convolution formula identifies each filtered increment with the translated exponential kernel on the independent Brownian coordinate.



Reassociate the initial-phase and Brownian product factors without altering their joint law. Each filtered increment is a measurable function of the corresponding Brownian interval. The established stream law therefore supplies the precise thermal innovations of the stationarity argument.

\subsection{\texorpdfstring{Prove the complete initial/\allowbreak{}noise law}{Prove the complete initial/noise law}}


Disjoint Brownian intervals give an iid stream of thermal Gaussian innovations. On the phase-times-Brownian product this stream is independent of the initial phase. Pulling that law back along the source reassociation and using the L\ensuremath{{}^{2}} representative identities proves the full initial-phase/\allowbreak{}entire-stream product law on the position–velocity–noise source.



The initial Gibbs phase is independent of the entire future source. Combining this product identity with the innovation-stream law proves the full initialization-and-noise input law. This supports repeated transition identities without a hidden independence assumption.

\subsection{\texorpdfstring{Identify the exact stationary position law}{Identify the exact stationary position law}}


The selected continuous solution satisfies the kinetic integral equation by construction. The initial/\allowbreak{}noise product law discharges every source premise of the proved splitting-limit stationarity theorem. This first proves the full phase has exactly the independent Gaussian-times-Gibbs target law. Projecting that phase law then gives the Gibbs position distribution at the terminal time. Strong convexity supplies the finite target second moment used in the construction.



Phase stationarity of the kinetic solution implies position stationarity by projection. The projected target is exactly the normalized density \(e^{-f}\). Thus the reference endpoint in the mixing theorem has the desired sampling law.

\subsection{\texorpdfstring{Derive the full and simple explicit mixing bounds}{Derive the full and simple explicit mixing bounds}}


In the arbitrary-shift two-law theorem, start the reference position from the Gibbs distribution and use its newly proved stationarity. This produces the full explicit shifted coefficient with the true initial Wasserstein distance. For shift zero and \(m+M\le\gamma^2\), the maximum in the exponent equals \(-m\); the prefactor simplifies to \(\sqrt2\). Finite L\ensuremath{{}^{2}} output and Gibbs moments also prove the transport distance is finite, ensuring the real-valued endpoint never conceals an infinite distance.



Apply the arbitrary-shift law contraction with one initial law equal to the stationary target. Its output remains the target, so the result is an explicit convergence-to-equilibrium bound. The simple friction regime follows by setting the shift to zero.

\section{\texorpdfstring{Push Wasserstein bounds through coordinate projections}{Push Wasserstein bounds through coordinate projections}}
\label{proof-module-KineticLangevin-TransportMap-lean}


These lemmas prove the Wasserstein effect of any measurable Lipschitz map between two normed spaces. They support coordinate projections and metric changes without assuming an optimal coupling.



Projection and randomized initialization are treated by pushing couplings through Lipschitz maps. The proof takes the full coupling infimum and therefore applies without assuming an optimal plan exists.

\subsection{\texorpdfstring{Push each coupling through the map and then take the infimum}{Push each coupling through the map and then take the infimum}}


For any coupling c of the two input laws, apply the same measurable map T to both coordinates. The two marginal identities are proved by composing pushforwards with the original coupling marginals. The pointwise Lipschitz inequality and monotonicity of L\ensuremath{{}^{2}} norms bound the output coupling cost by q times the input cost, even though the input and output spaces may differ. Taking the infimum over every original coupling gives \(W_2(T_\#\mu,T_\#\nu)\le qW_2(\mu,\nu)\). A product coupling establishes that the input coupling type is nonempty. For the real-valued statement, the extended upper bound first implies finite output cost from finite initial cost, then the nonnegative-real coefficient and the finite real value identities give the claimed real inequality. This prevents an infinite transport cost from disappearing through conversion to a real number.



A Lipschitz map sends any coupling to a coupling of the pushforward laws. Its pointwise distance bound controls the squared coupling cost. Taking square roots and then the infimum yields the Wasserstein Lipschitz bound, including for projections onto a coordinate block.

\subsection{\texorpdfstring{Adjoin an independent common source before changing the output space}{Adjoin an independent common source before changing the output space}}


The same argument applies to a measurable map from an input state and a common independent innovation into a different normed output space. Extend every original coupling by the source law, prove both output marginals, and retain the unchanged input coupling cost. The pointwise contraction and L\ensuremath{{}^{2}} monotonicity bound each output cost. Taking the true infimum yields the cross-space random-map transport bound. In particular, adjoining one shared Gaussian velocity to two position laws has no cost of its own.



Adjoining a common independent source does not change an initial coupling’s transport cost when the relevant discrepancy depends only on the original coordinates. This permits randomized initialization and output maps to be treated by the same coupling-pushforward argument.

\section{\texorpdfstring{The Brownian KIPLD flow and its stationary mixing theorem}{The Brownian KIPLD flow and its stationary mixing theorem}}
\label{proof-module-KIPLMC-Continuous-lean}


This section connects the canonical Brownian kinetic diffusion to the physical particle system and proves Proposition 3 for its output law. The initial position is an arbitrary probability law with finite second moment. Standard Gaussian velocity and the canonical Brownian source are independent of that position; velocity scaling gives the required physical thermal law. The Gibbs law is a proved stationary target, and the Wasserstein distance is the infimum over all couplings.



The canonical physical KIPLD process has the intended force, thermal initialization, and stationary target. Its law-level contraction and parameter projection separate dynamical mixing from finite-particle statistical error.

\subsection{\texorpdfstring{Construct the source, the physical flow, and its endpoint law}{Construct the source, the physical flow, and its endpoint law}}


Set \(c=\sqrt N\). The canonical source contains a standard Gaussian velocity and an independent vector Brownian source on the coordinate set \(\mathrm{Fin}(d_\theta)\sqcup(\mathrm{Fin}(N)\times\mathrm{Fin}(d_x))\). Use the proved selected kinetic flow with sum potential \(V_N\), smoothness \(NL\), and friction \(c\gamma\). The physical process is \(D_{c^{-1}}Y_{t/c}\). Position is unchanged, so its endpoint pushforward is exactly the canonical position law at time \(T/c\). This exact equality is the bridge used by the convergence proof.



The source combines the prescribed position initialization, independent thermal Gaussian velocity, and canonical Brownian process. The selected physical-time solution is a measurable function of these inputs. Projecting its endpoint defines the continuous particle position law.

\subsection{\texorpdfstring{Verify the independent thermal initialization}{Verify the independent thermal initialization}}


The complete source law theorem supplies the initial phase law \(\mathcal N(0,I)\otimes\nu\). Push it through velocity multiplication by \(N^{-1/2}\). The result is precisely the independent thermal-velocity and position law; the velocity and position remain independent. The selected flow’s time-zero identity and the almost-everywhere scalar-multiplication rule identify the L\ensuremath{{}^{2}} representatives with this pushforward. No initialization-law hypothesis is imposed on the final theorem.



The thermal initialization is the explicit inverse-square-root-particle Gaussian pushforward. Its independence from position and the entire Brownian source is proved by the product construction. This is precisely the phase/\allowbreak{}source law needed to transfer stationarity.

\subsection{\texorpdfstring{Cancel the clock factors and prove stationary mixing}{Cancel the clock factors and prove stationary mixing}}


The lifted potential has curvature \(Nm\) and smoothness \(NL\). Multiplication by \(N\) transforms \(m+L\le\gamma^2\) into \(Nm+NL\le(\sqrt N\gamma)^2\), so the canonical Gibbs mixing theorem applies. Its exponent is \(-(Nm)/(\sqrt N\gamma)\,(T/\sqrt N)=-mT/\gamma\). The resulting estimate is \(W_2(\nu_T,\pi_Z)\le\sqrt2e^{-mT/\gamma}W_2(\nu,\pi_Z)\). The canonical theorem proves Gibbs stationarity internally. Separate moment and probability theorems certify that all real Wasserstein quantities are finite.



The clock transform turns the process into the shared unit-temperature diffusion for the sum potential. Curvature and friction scaling cancel in the physical-time mixing exponent. Since the target position is unchanged, the resulting convergence bound is directly a particle position estimate.

\subsection{\texorpdfstring{Identify the physical force in the integral equation}{Identify the physical force in the integral equation}}


The drift lifted to L\ensuremath{{}^{2}} represents pointwise application of the true sum gradient. The almost-everywhere scalar multiplication and composition identities show that \(N^{-1}\mathcal G(\nabla V_N)=\mathcal G(F)\), where \(F=\nabla V_N/N\) is the previously verified particle force. Apply the exact kinetic clock theorem to the selected canonical solution. The result is the full physical-time integral equation with friction \(\gamma\), the force \(F\), and input \(N^{-1/2}B_{t/\sqrt N}\). The shared Brownian scaling theorem identifies its thermal amplitude as \(\sqrt{2\gamma/N}\).



The average-potential gradient identity rewrites the transformed drift as the parameter-average and latent-pair force. Consequently the constructed physical path solves the intended KIPLD equation. This identification links the law-level mixing theorem to the stated particle dynamics.

\subsection{\texorpdfstring{Project to the parameter and add the statistical error}{Project to the parameter and add the statistical error}}


The parameter projection has operator norm at most one, so pushing any coupling through it cannot increase its quadratic cost. Taking the true coupling infimum bounds the parameter Wasserstein distance by the full position distance. The initial independent coupling \(\nu\otimes\pi_Z\) bounds that infimum by \((\mathbb E\|Z_0-\bar Z_*\|^2)^{1/2}\), giving precisely the printed proposition. A second endpoint adds the already proved sharp equilibrium concentration radius \(\sqrt{d_\theta/(Nm)}\) by the L\ensuremath{{}^{2}} triangle inequality, giving a direct parameter RMS bound about the unique marginal-likelihood maximizer.



Nonexpansive parameter projection transfers the sampling transport error. The powered Gibbs marginal contributes the sharp statistical radius around the likelihood maximizer. Their sum yields a continuous-time parameter RMS bound with separate dynamical and finite-particle contributions.

\section{\texorpdfstring{Exact unit-\allowbreak{}temperature transfer of the KIPLMC2 law}{Exact unit-temperature transfer of the KIPLMC2 law}}
\label{proof-module-KIPLMC-OBABORescaling-lean}


Put \(c=\sqrt N\). Multiply the scaled physical velocity by \(c\), keep position unchanged, replace the clock step by \(h/c\), and replace friction by \(c\gamma\). The averaged force then becomes the gradient of the sum potential. Every subflow, the whole defined trajectory, its independent Gaussian Markov law and its Gibbs target satisfy this exact change of variables. This is the probability-level bridge needed to reuse the shared OBABO sampling theorem.



The complete five-stage particle update and its weighted metric are transported through the clock change. Exact target and whole-chain identities enable reuse of the unit-temperature OBABO sampling theorem.

\subsection{\texorpdfstring{Check thermostat, drift and kick under the clock change}{Check thermostat, drift and kick under the clock change}}


The thermostat sends \(v\) to \(\delta v+\sigma\xi\); multiplying velocity by \(c\) changes its noise amplitude to \(c\sigma\). The drift satisfies \((h/c)(cv)=hv\), so position is unchanged. For the kick, \((h/c)g=c h(c^{-2}g)\). One proves each identity by unfolding the map, distributing scalar multiplication and cancelling the nonzero clock factor. Combining these equalities in their original O–B–A–B–O order proves the full five-subflow conjugacy, including both distinct fresh innovations.



Under the velocity-and-clock change, the damping product is unchanged, while the drift and kick coefficients compensate for the velocity scale. The thermal amplitude is transformed at the same time. Verifying all five stages identifies the complete physical and unit-temperature steps.

\subsection{\texorpdfstring{Keep the contracting metric exactly invariant}{Keep the contracting metric exactly invariant}}


The transformed potential has smoothness \(c^2L\) and friction \(c\gamma\). In its phase energy, the velocity term is \(\lVert cv\rVert^2/(c^2L)=\lVert v\rVert^2/L\), and the mixed term is \(2\langle x,cv\rangle/(c\gamma)=2\langle x,v\rangle/\gamma\). The position term is unchanged. Thus the entire weighted energy is invariant. The proof expands its definition, applies the norm and inner-product scalar rules, and cancels the positive \(c\).



The weighted metric is designed from friction and phase coordinates. Scaling friction and velocity together preserves its value after the stated clock normalization. Thus transporting the contraction theorem introduces no hidden particle-dependent norm-equivalence factor.

\subsection{\texorpdfstring{Identify the defined particle step and entire trajectory}{Identify the defined particle step and entire trajectory}}


Take \(c=\sqrt N\). The force \(N^{-1}\nabla V_N\) is precisely \(c^{-2}\nabla V_N\), damping depends only on \(\gamma h\), and \(\sqrt N\sqrt{(1-\delta^2)/N}=\sqrt{1-\delta^2}\). These established identities specialize the generic conjugacy to Algorithm 2. Induction through its concrete successor recurrence proves the entire transformed trajectory is the standard-temperature OBABO trajectory for \(V_N\), driven by exactly the same paired Gaussian vectors.



Apply the step conjugacy inductively with the transformed noises. The entire physical trajectory is mapped to the shared unit-temperature trajectory. Since position is unchanged by the velocity map, their position sequences agree under this coupling.

\subsection{\texorpdfstring{Transport the Gibbs target and Markov laws}{Transport the Gibbs target and Markov laws}}


The unit-temperature phase target is \(\mathcal N(0,I)\otimes\pi_{V_N}\). Multiplying its velocity by \(N^{-1/2}\) gives exactly the defined KIPLMC thermal target; position remains \(\pi_{V_N}\). The product pushforward identity proves this equality after strong convexity supplies probability normalization. For numerical laws, the proved step conjugacy and joint measurability feed the general Markov-law conjugacy theorem. It establishes the equality for the defined independent-Gaussian law at every iteration, so transferring a sampling estimate requires no assumed correspondence between two algorithms.



The thermal Gaussian velocity becomes standard Gaussian, and the position Gibbs target is unchanged. The exact transition conjugacy then transports both the invariant comparison target and all numerical Markov iterates. This is the probabilistic bridge used by the particle OBABO error theorem.

\section{\texorpdfstring{A local state comparison with exact Hamiltonian dynamics}{A local state comparison with exact Hamiltonian dynamics}}
\label{proof-module-KineticLangevin-VerletLocal-lean}


To quantify sampling bias, compare the defined O–Verlet–O step with a reference that uses exact Hamiltonian motion between the same two thermostat updates. Exact Hamiltonian motion and thermostat preserve the Gibbs target separately. This section derives the local state error from the Hamiltonian integral equations. It works in any real Hilbert space, including L\ensuremath{{}^{2}}, where uniform velocity and force norm bounds are second-moment statements.



The defect is derived by integrating the true Hamiltonian equations. A common-noise comparison then inserts this deterministic estimate into the same weighted geometry used for OBABO contraction.

\subsection{\texorpdfstring{Integrate the Hamiltonian equations instead of assuming a local error}{Integrate the Hamiltonian equations instead of assuming a local error}}


The Hamiltonian equations are \(v(t)-v(0)=-\int_0^t g(x(s))ds\) and \(x(t)-x(0)=\int_0^t v(s)ds\). If the force norm is at most G, the integral norm inequality gives \(\|v(t)-v(0)\|\le Gt\). Subtract \(tv(0)\) from the position equation; the result is the integral of the velocity increment. Integrating the derived envelope \(Gs\) gives the Euler-position defect \(\|x(t)-x(0)-tv(0)\|\le Gt^2/2\). One proves the exact integral identities using integrability from continuity before applying either inequality.



The exact Hamiltonian solution satisfies integrated velocity and position equations. Substitute the position primitive into the force difference and use smoothness to obtain explicit integral remainders. This derives a local state error from the differential equation itself.

\subsection{\texorpdfstring{Compare both Verlet components to the true endpoint}{Compare both Verlet components to the true endpoint}}


The defined drift-stage position is \(x_V=x_0+hv_0-h^2g(x_0)/2\). Its difference from the Euler position has norm at most \(Gh^2/2\); adding the exact position defect proves \(\|x_V-x(h)\|\le Gh^2\). For velocity, subtract the frozen-force Euler value \(v_0-hg(x_0)\) from both trajectories. The exact difference is the integral of \(g(x(s))-g(x_0)\), bounded by \(MC h^2/2\) using Lipschitzness and the kinematics with velocity norm C. The Verlet difference is exactly \(-h(g(x_V)-g(x_0))/2\), whose norm is at most \(Mh^2C/2+MGh^3/4\). Combining the two gives \(\|v_V-v(h)\|\le MC h^2+MGh^3/4\). Every residual is computed from the update and the exact path; none is supplied as an opaque accuracy assumption.



The first half kick, drift, and final half kick are compared componentwise with those exact integrals. Terms sharing the same leading Taylor contribution cancel. The remaining bounds retain the initial velocity and force magnitudes needed for subsequent moment integration.

\subsection{\texorpdfstring{Express the state defect in the contracting OBABO metric}{Express the state defect in the contracting OBABO metric}}


The upper metric equivalence and completing-square Hilbert representation imply the convenient component estimate \(\|(v,x)\|_E\le2(\|x\|+\|v\|/\sqrt M)\). The proof squares both nonnegative sides, expands the scalar square, and uses \(\sqrt M^2=M\). Inserting the two component defects yields

\[\|\mathrm{Verlet}_h(v_0,x_0)-(v(h),x(h))\|_E\le2h^2\bigl(G+\sqrt M C+\sqrt M G h/4\bigr).\]

This estimate uses only a Lipschitz force and the explicit moment envelopes; no Hessian regularity is required.



The physical component errors are inserted into the verified weighted transformation. The triangle inequality yields a phase defect in the same norm used for contraction. Keeping this common metric avoids an extra norm-conversion loss at each later iteration.

\subsection{\texorpdfstring{Couple the defined OBABO step with the exact target-\allowbreak{}preserving split}{Couple the defined OBABO step with the exact target-preserving split}}


Start the Hamiltonian path at the first Ornstein–Uhlenbeck half-step. Apply the same final thermostat vector to its endpoint and to the Verlet endpoint. The final additive noise cancels exactly, while the velocity defect is multiplied by \(\delta\). Since \(|\delta|\le1\), its norm is no larger than the Verlet velocity defect, and the position defect is unchanged. Applying the component metric bound therefore gives the same explicit local estimate for the O–Verlet–O and O–exact-Hamiltonian–O outputs. In L\ensuremath{{}^{2}}, an invariant Gibbs Hamiltonian reference supplies \(C=\sqrt d\) and \(G\le\sqrt{Md}\) from the Gaussian velocity marginal and Gibbs score moment theorem. The local bound itself keeps the exact path equations and norm envelopes visible so the probabilistic reference construction must discharge them.



The defined OBABO step and its exact comparator use identical first and second thermostat noises. Their only difference is replacing the central exact Hamiltonian evolution by Verlet. Thus the deterministic Hamiltonian local defect directly controls the coupled stochastic step discrepancy.

\section{\texorpdfstring{The L\ensuremath{{}^{2}} Hamiltonian comparator has the exact Gibbs law}{The L² Hamiltonian comparator has the exact Gibbs law}}
\label{proof-module-KineticLangevin-LpHamiltonianLaw-lean}


This section closes the probabilistic reference obligations for the Hamiltonian part of the OBABO bias argument. It constructs the L\ensuremath{{}^{2}} Picard flow, identifies its finite-dimensional output law, derives exact Gibbs stationarity and velocity/\allowbreak{}score moments, and applies the proved local O–Verlet–O comparison. No stationary reference-law or local-error bound is postulated.



The second-mean Hamiltonian comparator is identified with the pointwise deterministic flow. Its exact target law then supplies the finite stationary moments needed to integrate the OBABO local state defect.

\subsection{\texorpdfstring{Identify an L\ensuremath{{}^{2}} limit with a pointwise numerical limit}{Identify an L² limit with a pointwise numerical limit}}


Suppose L\ensuremath{{}^{2}} phase variables \(S_n\) converge to S, and their representatives agree almost everywhere with measurable vector approximants \(f_n\) converging almost everywhere to Y. For any continuous linear functional \(\ell\), the characteristic test \(e^{i\ell(z)}\) has norm one. Bounded dominated convergence identifies the characteristic-function limit of the pointwise approximants, while the proved L\ensuremath{{}^{2}} characteristic continuity identifies the limit of the L\ensuremath{{}^{2}} variables. Their equality yields \(\mathcal L(S)=\mathcal L(Y)\) through characteristic uniqueness. The next lemma proves the required almost-everywhere correspondence for every Verlet iterate by induction from the lifted kick and drift identities. Thus the two convergence arguments refer to precisely the same defined approximants.



A second-mean convergent sequence has a subsequence converging almost everywhere. If the same numerical maps converge pointwise to the deterministic Hamiltonian endpoint, the two limits coincide almost everywhere. This identifies the Hilbert-space limit with the intended pointwise flow.

\subsection{\texorpdfstring{Identify the coupled representatives almost everywhere}{Identify the coupled representatives almost everywhere}}


Marginal equality alone is insufficient when the local numerical error must be coupled with the same initial state and thermostat inputs. This stronger result identifies the representatives almost everywhere. Each component of an L\ensuremath{{}^{2}}-convergent phase sequence converges in measure. The known almost-everywhere equality transfers this convergence to the same pointwise Verlet approximants. Their almost-everywhere convergence also implies convergence in measure, because the source is a probability measure. Uniqueness of the limit in measure therefore gives equality almost everywhere of each component, hence of the phase pair.



Apply this argument to the L\ensuremath{{}^{2}} Hamiltonian flow and its Verlet approximants. The result states that its representative at time t is almost everywhere the finite-dimensional Hamiltonian flow evaluated at the same representative of the initial state. Thus adjoining any initial state or thermostat variable preserves the coupling used in the local comparison; no equality of marginal laws is substituted for a joint-law argument.



The initial and numerical representatives are chosen on the same probability source. The verified stage identities identify them with the original vector algorithm. The resulting law statement concerns random variables, independently of any representative chosen for their equivalence classes.

\subsection{\texorpdfstring{Construct the true L\ensuremath{{}^{2}} Hamiltonian integral flow and identify its law}{Construct the true L² Hamiltonian integral flow and identify its law}}


The L\ensuremath{{}^{2}} force is the Nemytskii lift the square-integrable drift; its global Lipschitz bound permits the existing Picard construction with zero friction and zero noise. Its specification gives a continuous global path, the prescribed initial value, and the Hamiltonian integral equation on the finite horizon. Verlet convergence holds in this Banach space. The same vector Verlet maps converge pointwise to the selected finite-dimensional Hamiltonian flow. The proved almost-everywhere pointwise identification and the composition rule for pushforwards give the exact law identity \(\mathcal L(A_t)=(H_t)_\#\mathcal L(A_0)\). This does not infer a pathwise identification from an unproved interchange of representatives and time integrals.



Solve the Hamiltonian integral equation in second mean and identify its representative with the pointwise selected flow. Pushing forward the known initial law then gives its exact endpoint law. For Gibbs initialization, the previously proved Hamiltonian invariance makes that law stationary.

\subsection{\texorpdfstring{Derive stationarity and both moment envelopes from the target}{Derive stationarity and both moment envelopes from the target}}


If the initial phase law is the canonical independent Gaussian-times-Gibbs target, the exact Hamiltonian invariance theorem and the derived output-law identity prove that the L\ensuremath{{}^{2}} flow has that target law at every time. Projecting the integral equation with the continuous linear coordinate maps gives the velocity and kinematic component equations; commuting those maps with Bochner integrals is justified by continuity and interval integrability. Projecting the target law gives standard Gaussian velocity and normalized Gibbs position. The exact Gaussian coordinate calculation therefore yields \(\|V_t\|_{L^2}=\sqrt d\), and the already proved Gibbs score estimate gives \(\|\nabla f(X_t)\|_{L^2}\le\sqrt{Md}\). These are conclusions from the law, not supplied moment hypotheses.



The target factorization supplies the Gaussian velocity marginal and the Gibbs position marginal. Their finite moments control the velocity and force envelopes along the stationary Hamiltonian reference. These are the quantities entering the Verlet local-error estimate.

\subsection{\texorpdfstring{Apply the local OBABO comparison to this constructed stationary reference}{Apply the local OBABO comparison to this constructed stationary reference}}


Start the L\ensuremath{{}^{2}} Hamiltonian path at the first thermostat half-step and assume its phase law is the canonical target, as supplied by thermostat invariance. Its continuity, initial value, velocity equation, kinematic equation, and uniform moments have all been proved above. The local O–Verlet–O versus O–Hamiltonian–O theorem therefore applies with \(C=\sqrt d\) and \(G=\sqrt{Md}\). It yields the weighted error bound

\[2h^2\left(\sqrt{Md}+\sqrt M\sqrt d+\frac{h}{4}\sqrt M\sqrt{Md}\right).\]

The two thermostat noises are retained, their common final contribution cancels, and the damping needs only \(|\delta|\le1\). This provides the numerical local budget for the global sampling-bias estimate.



Apply the deterministic common-noise OBABO comparison to the constructed stationary reference variables. Integrating its phase envelope gives the stated local root-mean-square cost. Thus the local kernel discrepancy uses Hamiltonian and thermostat laws.

\section{\texorpdfstring{Identify weighted phase L\ensuremath{{}^{2}} errors with finite-\allowbreak{}dimensional costs}{Identify weighted phase L² errors with finite-dimensional costs}}
\label{proof-module-KineticLangevin-LpWeightedGeometry-lean}


A local estimate in a pair of L\ensuremath{{}^{2}} spaces must refer to the same Euclidean phase geometry as the vector sampler. This section constructs the explicit L\ensuremath{{}^{2}} Euclidean pair and proves its exact norm and representative identities. It then proves that the weighted representative cost is precisely the weighted phase norm used by the local analysis.



The weighted Hilbert-space norm is proved equal to the second-mean finite-dimensional phase cost. This equality converts analytic local-error bounds directly into transport costs of coupled outputs.

\subsection{\texorpdfstring{Construct a Euclidean phase variable in L\ensuremath{{}^{2}}}{Construct a Euclidean phase variable in L²}}


The coordinate injections \(u\mapsto(u,0)\), \(w\mapsto(0,w)\) are continuous linear maps into the Euclidean product Hilbert space. Lift each through the bounded pullback map and add their outputs. Almost-everywhere identities for lifted linear maps and L\ensuremath{{}^{2}} addition show that the resulting variable represents exactly \((A(\omega),B(\omega))\). Expand its L\ensuremath{{}^{2}} inner product as an integral, use the Euclidean product inner product pointwise, and split the two integrable self-inner-products. This proves \(\|\operatorname{pair}(A,B)\|_{L^2}^2=\|A\|_{L^2}^2+\|B\|_{L^2}^2\). Thus the construction uses the Euclidean geometry, not the default maximum norm of an ordinary product.



Package velocity and position as one Euclidean phase-valued random variable. Its second-mean squared norm is the sum of the two component second moments. This identifies the functional-analytic phase cost with the finite-dimensional cost used in transport.

\subsection{\texorpdfstring{Lift the weighted transform without changing its norm}{Lift the weighted transform without changing its norm}}


Apply the Euclidean pair construction to \(X+bV\) and \(\sqrt{a-b^2}V\). The representative is almost everywhere the original vector weighted transform \(T(V(\omega),X(\omega))\), with all addition and scalar-multiplication representatives justified explicitly. The squared norm identity for the pair is exactly the defining squared norm of the weighted phase transform on the L\ensuremath{{}^{2}} Hilbert space. Both norms are nonnegative, so equality of their squares gives equality of the norms. Finally the standard L\ensuremath{{}^{2}} norm definition and finiteness of an L\ensuremath{{}^{2}} variable give

\[\|T(V(\cdot),X(\cdot))\|_{L^2}=\|(V,X)\|_E.\]

The formal statement retains extended-real L\ensuremath{{}^{2}} norms on the left and converts the finite norm with the nonnegative extended-real embedding. This is the needed bridge from local Hilbert estimates to true coupling costs.



The weighted transformation is bounded linear, so applying it before or after passage to square-integrable classes gives the same element. Its norm identity therefore transfers exactly to random phase errors. No extra coordinatewise inequality is needed.

\subsection{\texorpdfstring{Keep differences and the entire OBABO step tied to their representatives}{Keep differences and the entire OBABO step tied to their representatives}}


The first lemma records the almost-everywhere representative of a phase difference, using the two component subtraction identities. The second expands OBABO into first thermostat, Verlet, and final thermostat, then composes the already proved representative identities for all three. Consequently the lifted nonlinear force update represents exactly the original finite-dimensional O-B-A-B-O sampler with the same two supplied noise variables. These equalities preserve the joint coupling; a marginal-law equality is never substituted for the common-noise realization.



The representative identities hold for differences and for each affine stage of the OBABO update. Thus the weighted local-error norm is the root-mean-square distance between the coupled vector outputs. This connects the analytic budget with a transport coupling.

\section{\texorpdfstring{The local OBABO kernel discrepancy has the proved L\ensuremath{{}^{2}} budget}{The local OBABO kernel discrepancy has the proved L² budget}}
\label{proof-module-KineticLangevin-OBABOLocalCost-lean}


This section turns the L\ensuremath{{}^{2}} Hamiltonian local comparison into the integrated cost of the original finite-dimensional maps, retaining the same initial state and both noise variables on one probability space.



The local kernel discrepancy is a common-noise transport cost. It is obtained by integrating the deterministic Verlet error against an stationary Hamiltonian comparator in the contracting phase metric.

\subsection{\texorpdfstring{Record the explicit local budget}{Record the explicit local budget}}


The local budget is \(b_h=2h^2(\sqrt{Md}+\sqrt M\sqrt d+h\sqrt M\sqrt{Md}/4)\). Its nonnegativity follows directly from \(h\ge0\) and square-root nonnegativity. It is O(h\ensuremath{{}^{2}}\ensuremath{\surd}d) with coefficients depending only on M and the allowed step-size range. No numerical error is defined to equal this budget; the next theorem proves it bounds the error.



The local budget retains the explicit step size, friction, smoothness, and dimension factors obtained from the Hamiltonian comparator. Its order in the step is quadratic. This order is needed because geometric accumulation will later divide by a contraction gap of first order.

\subsection{\texorpdfstring{Identify the full common-\allowbreak{}noise defect before taking its cost}{Identify the full common-noise defect before taking its cost}}


Construct the L\ensuremath{{}^{2}} Hamiltonian path at the first thermostat output. The strengthened pointwise Hamiltonian lifting theorem identifies its endpoint almost everywhere with the finite-dimensional exact Hamiltonian flow at that same initial state. The lifted OBABO theorem and final-thermostat representative theorem therefore identify the two outputs jointly, including both original noise variables. Subtracting the phase representatives and applying the weighted-transform L\ensuremath{{}^{2}} identity gives the exact integrated weighted cost of the original vector map discrepancy. The stationary-reference local estimate then proves that this cost is at most \(b_h\). The only reference-law input is the law of the first thermostat output; canonical target initialization and exact thermostat invariance discharge that input in the full sampler theorem.



Both updates use exactly the same thermostat noises. The weighted random-variable discrepancy is therefore a coupling cost for their output measures. Integrating the previously derived deterministic local defect under the stationary reference yields the stated kernel error.

\section{\texorpdfstring{Convert the local h\ensuremath{{}^{2}} cost into a global h\ensuremath{\surd}d bias}{Convert the local h² cost into a global h√d bias}}
\label{proof-module-KineticLangevin-OBABOBiasScalar-lean}


The contracting metric factor is \(r=\sqrt{1-mh/(3\gamma)}\). These scalar lemmas prove its strict contraction and a lower bound on its gap, then simplify the genuinely derived local budget into an explicit first-order bias coefficient.



The contraction gap is of first order in the step, while local kernel error is of second order. Their ratio produces the explicit linear-step global bias, with all friction and dimension factors visible.

\subsection{\texorpdfstring{Quantify the contraction gap}{Quantify the contraction gap}}


For positive m, gamma and h, the squared rate is strictly below one, so the nonnegative square root r is also below one. When its square is nonnegative, \(1-r^2=mh/(3\gamma)\). The nonnegative square \((1-r)^2\) gives \(2(1-r)\ge1-r^2\), hence \(1-r\ge mh/(6\gamma)\). All division steps retain their positive denominator hypotheses.



For a contraction coefficient \(r=\sqrt{1-s}\), rationalization gives \(1-r=s/(1+\sqrt{1-s})\ge s/2\) when \(0\le s\le1\). This lower bound turns the local error into an explicit global stationary-bias coefficient.

\subsection{\texorpdfstring{Simplify the local budget and divide by the proved gap}{Simplify the local budget and divide by the proved gap}}


The nonnegative square-root product identity yields \(b_h=h^2\sqrt d(4\sqrt M+Mh/2)\). Dividing the local budget by the proved lower bound on \(1-r\) gives

\[\frac{b_h}{1-r}\le\frac{6\gamma}{m}h\sqrt d(4\sqrt M+Mh/2).\]

Thus the accumulated bias is O(h\ensuremath{\surd}d), and restricting h to its source upper bound makes its coefficient depend only on m, M and gamma. This is arithmetic applied to an independently proved local map error and contraction rate.



The local cost is proportional to \(h^2\sqrt d\), while the contraction gap is bounded below by a constant times \(mh/\gamma\). Dividing cancels one power of \(h\). The remaining terms give the stated linear-step bias with all physical constants retained.

\section{\texorpdfstring{A concrete independent initial phase and canonical Brownian source}{A concrete independent initial phase and canonical Brownian source}}
\label{proof-module-KineticLangevin-CanonicalProduct-lean}


The stationarity proof takes place on one explicit probability space: the product of an initial phase law and the constructed vector Brownian source. This section builds the initial L\ensuremath{{}^{2}} variables, the continuous Brownian forcing and its OU-filtered innovations, then proves the complete initial/\allowbreak{}noise joint law. The resulting splitting law identity has no assumed independence or Gaussian marginal fields.



Independent initialization and the canonical Brownian source give the exact input product law. This identifies the additive-input and explicit thermal-noise splitting endpoints without changing their probability distributions.

\subsection{\texorpdfstring{Represent the original initial phase on the product source}{Represent the original initial phase on the product source}}


the response Hilbert space is the finite-dimensional Euclidean state, and the product sample space pairs a phase with a vector Brownian sample. The probability measure is exactly \(\pi\otimes P_B\). Pull the velocity and position projections through the first product projection. Their square integrability follows from that of the original phase law. The two the square-integrable representative representative identities prove that the resulting phase agrees almost everywhere with the original initial coordinate.



The initial phase is realized as the first coordinate of an independent product with the canonical Brownian source. Its law is exactly the prescribed initial law. The second coordinate supplies the full source independently of every initial phase component.

\subsection{\texorpdfstring{Lift the continuous Brownian input}{Lift the continuous Brownian input}}


The second product projection preserves the Brownian source law. Its pullback is a linear isometry on L\ensuremath{{}^{2}}, hence a continuous linear map. Apply it to \(\sqrt{2\gamma}B_t\), where the constructed Brownian input is continuous in L\ensuremath{{}^{2}} and zero at time zero. Linearity and continuity give the same properties after adjoining the independent initial phase. Thus the forcing used by the deterministic integral solver is an Brownian forcing on this very source.



Pullback of a square-integrable variable along a measure-preserving projection is an isometry. It preserves the Brownian response identities and second-mean continuity. Thus adjoining initial data does not alter the source equations or moments.

\subsection{\texorpdfstring{Use the exact OU-\allowbreak{}filtered Brownian intervals}{Use the exact OU-filtered Brownian intervals}}


The innovation in interval \([nh,(n+1)h]\) is the pulled-back isonormal convolution
\[\sqrt{2\gamma}\int_{nh}^{(n+1)h}e^{-\gamma((n+1)h-s)}dB_s\]
. Its representative is exactly the original source variable composed with the Brownian projection. Continuous linear maps commute with the deterministic OU filter. Combined with the proved canonical integration-by-parts identity, this shows the continuous-input filter is exactly this innovation, not merely a random variable with a matching covariance.



The splitting uses the exponentially filtered increment on each numerical interval. The source calculation identifies that increment exactly with the corresponding canonical Gaussian response. The physical amplitude is retained, so its law has the correct thermal variance.

\subsection{\texorpdfstring{Derive the complete independent Gaussian stream law}{Derive the complete independent Gaussian stream law}}


The canonical OU marginal theorem gives covariance \((1-e^{-2\gamma h})I\); the disjoint-interval theorem gives independence of the full vector innovations across time. Their whole stream therefore has the infinite product of thermal Gaussian law. The explicit initial-law/\allowbreak{}source product makes the initial phase independent of that entire stream. Finally use the almost-everywhere initial and innovation identities to obtain the complete joint law of the L\ensuremath{{}^{2}} variables.



Independence of disjoint Brownian intervals gives independence of the filtered vector responses. Their common covariance identifies their marginal law. Together with the independent initialization this proves the full input product law required by the splitting Markov chain.

\subsection{\texorpdfstring{Identify both splitting endpoint laws}{Identify both splitting endpoint laws}}


Apply the proved full-input-law theorem to the lifted Verlet–OU chain. Its phase pushforward is exactly the Markov iterate with fresh thermal Gaussian OU innovations and initial law \(\pi\). The exact OU-filter identity also identifies the continuous-input chain with the same noise-driven chain, so its endpoint has that law as well. These identities provide the missing connection between the pathwise convergence theorem and the measure-level weak stationarity calculation.



The deterministic additive-input splitting and the explicit independent-noise splitting have equal endpoints under these identities. Their laws therefore coincide. This lets the pathwise convergence theorem and the stationary weak-test theorem refer to the same numerical process.

\section{\texorpdfstring{The canonical Gibbs-\allowbreak{}started kinetic Langevin solution}{The canonical Gibbs-started kinetic Langevin solution}}
\label{proof-module-KineticLangevin-CanonicalStationarity-lean}


The first-order potential assumptions construct a square-integrable initial phase and an independent canonical Brownian source. The selected Picard solution satisfies the kinetic integral equation. The previously proved splitting limit establishes its exact Gaussian-times-Gibbs law at every time; stationarity is a conclusion of this concrete construction.



The initial Gibbs phase and canonical Brownian source are realized on a concrete independent product. Splitting convergence and the verified joint source law establish stationarity of the selected kinetic integral solution.

\subsection{\texorpdfstring{Derive the two target projection moments}{Derive the two target projection moments}}


The target is the product \(\pi=\mathcal N(0,I)\otimes\pi_f\), where \(\pi_f(dx)=Z_f^{-1}e^{-f(x)}dx\). Its velocity projection lies in L\ensuremath{{}^{2}} because a finite-dimensional Gaussian has a finite second moment. Its position projection lies in L\ensuremath{{}^{2}} because strong convexity with \(m>0\) gives Gaussian tail domination for the normalized Gibbs measure. First, use that \(\pi_f\) has total mass one. Factor the phase product and compose the existing Gaussian or Gibbs L\ensuremath{{}^{2}} identity with the corresponding measure-preserving product projection. These lemmas prove the exact integrability required by the random initial state; they do not introduce moment hypotheses.



The target’s product structure gives the position second moment and the exact Gaussian velocity second moment. These ensure that the phase initialization belongs to the square-integrable Hilbert space. They are derived from the target rather than imposed as unexplained moment assumptions.

\subsection{\texorpdfstring{Build the independent Gibbs and Brownian probability space}{Build the independent Gibbs and Brownian probability space}}


The state space is \(E=\mathbb R^d\), defined as Euclidean space on \(\mathrm{Fin}\,d\). The probability space is the product of the normalized phase law with the already constructed canonical Brownian source. Its initial phase is the two L\ensuremath{{}^{2}} coordinate projections, whose integrability was proved above. Thus the starting velocity is genuinely standard Gaussian, the starting position has the normalized Gibbs law, and the whole Brownian source is independent of both. The argument uses the derived projection \(square-integrability\) statements to construct the L\ensuremath{{}^{2}} elements through the initial phase variable. The definitions name this concrete source and initial state so subsequent theorems cannot silently substitute a source with an assumed stationary output law.



Take the product of the normalized Gibbs phase measure and the canonical Brownian probability space. The coordinate projections have exactly those laws and are independent. This concrete probability space supports both the initial phase and all future filtered responses.

\subsection{\texorpdfstring{Construct the fluctuation-\allowbreak{}dissipation input}{Construct the fluctuation-dissipation input}}


On the independent product source, the forcing is \(B_t=\sqrt{2\gamma}\,W_t\), where the vector Brownian motion has the already covariance and independent increment laws. It is a continuous function of time taking values in the L\ensuremath{{}^{2}} space, and \(B_0=0\). The proof pulls back the canonical Brownian coordinates along the product projection, carries over their proved L\ensuremath{{}^{2}} continuity, and rewrites their value at zero. These are properties of the constructed input, rather than regularity assumptions on an unspecified process. The drift is the pointwise potential gradient lifted to this same L\ensuremath{{}^{2}} space; its global Lipschitz constant is inherited from the stated gradient bound.



The additive velocity input is the vector Brownian process multiplied by \(\sqrt{2\gamma}\). Its filtered increments have thermal Gaussian law required by the splitting. The amplitude is fixed by the fluctuation–dissipation relation, ensuring the target velocity temperature is one.

\subsection{\texorpdfstring{Select a continuous solution and record its equation}{Select a continuous solution and record its equation}}


For a horizon \(T\ge0\), Picard existence for a globally Lipschitz additive equation constructs the phase path \(Y_t=(V_t,X_t)\) in L\ensuremath{{}^{2}}. It satisfies \[V_t-V_0=-\int_0^t(\gamma V_s+\nabla f(X_s))\,ds+\sqrt{2\gamma}\,W_t,\qquad X_t-X_0=\int_0^t V_s\,ds.\] The product notation in The proof combines these two equations into one integral identity using the negative drift. The flow definition supplies the proved L\ensuremath{{}^{2}} Lipschitz gradient, the constructed initial state, and the continuous input to the restarted kinetic solution. Its specification theorem then proves continuity, the exact initial value, and the displayed integral equation on the whole interval. No stationary process or solution identity is postulated.



Existence and uniqueness give a continuous square-integrable phase path with the prescribed initial state and input. Recording its integral equation links this object to the physical kinetic dynamics. Stationarity is then a property to prove about that specific path.

\subsection{\texorpdfstring{Prove the exact target law at every time}{Prove the exact target law at every time}}


Fix \(t\in[0,T]\). On every mesh size \(h\ge0\), the initial phase and the entire filtered Brownian innovation stream have joint law \[\pi\otimes Q_h^{\otimes\mathbb N},\qquad Q_h=\mathcal N\!\left(0,(1-e^{-2\gamma h})I\right).\] The argument obtains this from the proved product-source joint-law theorem and the exact identity between filtered Brownian increments and OU innovations. It restricts the already solution equation from \([0,T]\) to \([0,t]\), then applies the splitting stationarity theorem. That theorem has already proved convergence of the concrete splitting chain, an order-\(h^2\) stationary weak defect from Metropolis rejection, Markov telescoping, and continuity of the limiting law. The conclusion is the exact pushforward identity \[\operatorname{Law}(V_t,X_t)=\mathcal N(0,I)\otimes\pi_f.\] Every source, initial moment, solution, innovation law, and limit used here has been constructed or proved; there is no invariant-target or output-law assumption remaining in this endpoint.



Apply the splitting-limit stationarity transfer at an arbitrary time in the horizon. The verified joint input law and the exact splitting convergence discharge its hypotheses. The resulting phase law is the full product target, yielding both stationary marginals simultaneously.

\section{\texorpdfstring{An Gibbs-\allowbreak{}preserving O–Hamiltonian–O reference step}{An Gibbs-preserving O–Hamiltonian–O reference step}}
\label{proof-module-KineticLangevin-ExactOBABO-lean}


To measure OBABO discretization bias, keep both thermostat noises exactly the same and replace its Verlet middle portion by the proved exact Hamiltonian flow. This comparison step preserves the Gibbs target. Its invariance and its common-source L\ensuremath{{}^{2}} witnesses are constructed here, so the numerical error bound can compare concrete random variables.



Replacing only the central Verlet stage by exact Hamiltonian evolution creates a target-preserving comparison transition. Retaining the same two thermostat noises yields the natural coupling for its local numerical discrepancy.

\subsection{\texorpdfstring{Define the exact comparison map with the same two noises}{Define the exact comparison map with the same two noises}}


Starting from \(z=(v,x)\), apply \(O_{\delta,\sigma,\xi}\), evolve the resulting phase under the selected Hamiltonian flow for time \(h\), and then apply \(O_{\delta,\sigma,\xi\prime}\). The map uses the same first and second innovations as the OBABO update. Only the deterministic Hamiltonian evolution changes. Continuity of the selected Hamiltonian endpoint and of the two affine thermostat maps proves joint continuity in the input phase and the paired noise.



The comparison replaces the central Verlet stage by exact Hamiltonian evolution while retaining the same thermostat half steps. Since all three exact stages preserve the Gibbs target, their composition is a natural target-preserving comparator for the numerical OBABO map.

\subsection{\texorpdfstring{Preserve the first updated phase and its independence of fresh noise}{Preserve the first updated phase and its independence of fresh noise}}


The source law is explicitly \(\pi\otimes(\mathcal N(0,I)\otimes\mathcal N(0,I))\), where \(\pi\) is the normalized Gaussian-times-Gibbs phase target. Reassociate it as \((\pi\otimes\mathcal N(0,I))\otimes\mathcal N(0,I)\). The first thermostat preserves \(\pi\) because \(\delta^2+\sigma^2=1\), while the untouched second Gaussian remains independent. The proof implements this using the measure-preserving product associator and the product pushforward identity. The conclusion is the joint law of the updated phase and the second fresh noise, not just a marginal invariance statement.



After the first thermostat update the phase still has the Gibbs law. The second thermostat innovation remains independent because it was chosen independently of the initial phase and the first innovation. This justifies applying its invariance property after the Hamiltonian stage.

\subsection{\texorpdfstring{Compose the three proved invariant transitions}{Compose the three proved invariant transitions}}


After the first O step the joint law is \(\pi\otimes\mathcal N(0,I)\). The exact Hamiltonian map preserves \(\pi\) under the original strong-convexity and Lipschitz-gradient assumptions; applying it only to the phase therefore leaves the same joint law with the second noise. The last O step preserves \(\pi\) again. Three explicit pushforward composition identities prove that the complete O–Hamiltonian–O random map preserves the normalized target. No invariance hypothesis is supplied to this theorem.



Push the target successively through the first thermostat, Hamiltonian flow, and final thermostat. Each pushforward equals the target, so their composition does as well. This proves invariance for the concrete comparison kernel, including both noises.

\subsection{\texorpdfstring{Construct the L\ensuremath{{}^{2}} variables used by the local comparison}{Construct the L² variables used by the local comparison}}


Pull the initial velocity and position through the first product projection, and the two Gaussian innovations through their separate projections. The target projection moment theorems and Gaussian finite moments prove all four coordinates belong to L\ensuremath{{}^{2}}. Their the square-integrable representative representatives agree almost everywhere with the original coordinates. The proved pointwise representation of a lifted thermostat identifies the first updated L\ensuremath{{}^{2}} phase with the first O map. Projecting the established phase/\allowbreak{}fresh-noise joint law now gives its exact target marginal. The resulting witnesses simultaneously retain the common noise coupling and supply the stationarity input needed by the local Hamiltonian–Verlet error theorem.



The initial phase and both innovations are represented jointly on their product space. The exact Hamiltonian comparator is square integrable under its stationary law. These variables provide the common-noise coupling needed to integrate the Verlet local state error.

\section{\texorpdfstring{Pass explicit local kernel errors through the true coupling infimum}{Pass explicit local kernel errors through the true coupling infimum}}
\label{proof-module-KineticLangevin-RandomMapPerturbation-lean}


These transport lemmas connect measurable random maps to their iterated independent-noise laws. They preserve every marginal and take the infimum over initial couplings. Their inputs are a pathwise contraction and the measured L\ensuremath{{}^{2}} discrepancy between two explicit maps; these analytic inputs are proved separately for OBABO and its Hamiltonian comparator.



A common independent innovation extends any input coupling into an output coupling. Minkowski’s inequality separates contraction from local kernel discrepancy, producing an affine transport recurrence against an invariant comparator.

\subsection{\texorpdfstring{Extend every coupling by a common independent innovation}{Extend every coupling by a common independent innovation}}


Given any input coupling \(c\) of \(\mu,\nu\), append one independent innovation with law Q. Let the numerical output be \(F(x,s)\) and the comparator output be \(G(y,s)\). Their difference decomposes exactly as

\[F(x,s)-G(y,s)=[F(x,s)-F(y,s)]-[G(y,s)-F(y,s)].\]

Minkowski and the pathwise contraction bound the first term by q times the original coupling cost. The law of \((y,s)\) is proved to be exactly \(\nu\otimes Q\); the pushforward formula for L\ensuremath{{}^{2}} therefore identifies the second term with the local discrepancy under that fixed product law. Both output marginals are the random-map laws. The resulting bound holds for every c, so the established affine-infimum theorem gives \(W_2(F_\#(\mu\otimes Q),G_\#(\nu\otimes Q))\le qW_2(\mu,\nu)+b\). No optimal-coupling existence is assumed.



Take any coupling of the two input measures and adjoin one independent innovation. Updating both coupled states produces the desired output marginals. Minkowski’s inequality separates the contractive shared-map discrepancy from the local discrepancy between the two maps.

\subsection{\texorpdfstring{Iterate the Markov laws against an invariant comparator}{Iterate the Markov laws against an invariant comparator}}


Suppose the explicit comparator map preserves \(\pi\) under the independent source Q and its local L\ensuremath{{}^{2}} error is at most b. Applying the preceding transport theorem at every numerical iteration gives \(e_{n+1}\le qe_n+b\). The proof first establishes that each iterated law is a probability and that every transport distance is finite, using the extended-real recurrence and finite initial distance. Only then does it convert to real distances. The proved geometric recurrence yields

\[W_2(\mu_n,\pi)\le q^nW_2(\mu_0,\pi)+\frac{b}{1-q},\qquad0\le q<1.\]

This general lemma keeps the invariant-map and integrated-local-error obligations explicit. The Hamiltonian and local-state modules supply them for the OBABO comparison; the convergence theorem itself is not used as an input.



If the comparator transition preserves the target, each step compares the numerical law directly with the same target. Iterating the one-step affine Wasserstein inequality yields a geometric transient and a sum of local errors. This is the law-level analogue of the pathwise recurrence.

\subsection{\texorpdfstring{Expose finite output transport before taking real values}{Expose finite output transport before taking real values}}


The extended affine one-step bound also proves finite output distance at every Markov iteration. Start with finite initial transport. At each step the finite nonnegative coefficient times the preceding finite distance, plus the finite local budget, is finite. The distance lies below that finite quantity. This induction is exposed separately so coordinate projections and mean-square consequences can use a real Wasserstein estimate without silently converting infinity to zero.



A finite-cost coupling and a finite local root-mean-square defect give a finite output cost. Establishing this before converting the extended transport value to a real number makes the affine recurrence legitimate. The final infimum is taken over couplings.

\section{\texorpdfstring{The independent-\allowbreak{}Gaussian OBABO sampler has a proved global bias bound}{The independent-Gaussian OBABO sampler has a proved global bias bound}}
\label{proof-module-KineticLangevin-OBABOBias-lean}


This is the unconditional numerical endpoint for OBABO under the original strong-convexity and smoothness assumptions. The Markov laws converge toward the canonical Gaussian-times-Gibbs target up to an explicit O(h\ensuremath{\surd}d) bias. All stochastic reference laws, Hamiltonian invariance, moment envelopes, common-noise local errors, weighted-coordinate law identities, and the initial Wasserstein infimum are proved internally.



The numerical transition is compared with an exact invariant O–Hamiltonian–O transition in the same weighted coordinates. Local cost and contraction combine into the complete global transport-bias theorem.

\subsection{\texorpdfstring{Construct the exact reference in the same contracting coordinates}{Construct the exact reference in the same contracting coordinates}}


The reference takes the first thermostat output, evolves it by the Hamiltonian integral flow, and applies the same final thermostat input. Conjugate this exact map by the explicit invertible weighted transform. Its continuity follows from the verified inverse, the exact Hamiltonian endpoint continuity, the two thermostat formulas, and the final transform. The local Borel structure is made explicit for the Euclidean weighted phase space. The numerical and reference maps now share both the same innovations and the same contracting geometry.



The exact O–Hamiltonian–O transition is conjugated by the same weighted transformation as the numerical OBABO transition. Its invariant law is the transformed Gibbs target. This places contraction and local kernel discrepancy in one common transport geometry.

\subsection{\texorpdfstring{Prove the canonical source cost and exact reference stationarity}{Prove the canonical source cost and exact reference stationarity}}


The source is precisely the canonical phase target times a fresh pair of independent standard Gaussian vectors. The canonical source theorem supplies L\ensuremath{{}^{2}} coordinate representatives and proves that the first thermostat output has the target law. The constructed Hamiltonian L\ensuremath{{}^{2}} path therefore has the derived Gaussian velocity and Gibbs score moments. Its almost-everywhere identification with the vector flow keeps the whole local coupling intact. The integrated local-cost theorem yields the true cost \(b_h\), with no analytic error premise.



Exact O–Hamiltonian–O invariance on the original source is already proved. The general Markov conjugacy identity transports this invariance to the weighted target. For the weighted local cost, transforming the initial state commutes with adjoining the independent Gaussian pair. The L\ensuremath{{}^{2}} pushforward formula, inverse coordinate identity, and invariance of the norm under reversing a difference identify its cost with the previously bounded canonical cost. Thus the exact target invariance and the local discrepancy needed by the global transport theorem are both established for the explicit maps.



The source consists of the stationary phase and two independent thermostat noises. The verified exact transition preserves its target law. Coupling it with the numerical map using those same noises gives the finite local root-mean-square cost used by the perturbation theorem.

\subsection{\texorpdfstring{Derive the full Markov bound and simplify the accumulated bias}{Derive the full Markov bound and simplify the accumulated bias}}


Let \(r=\sqrt{1-mh/(3\gamma)}\). Strong convexity, gradient Lipschitzness, \(\gamma^2\ge4M\), and \(33\gamma^3h\le m\) give the numerical-map contraction. The preceding proofs give exact reference invariance and local cost \(b_h\). The all-couplings Markov perturbation theorem therefore applies to the independent-Gaussian laws. The coordinate-law identity proves these are exactly the original defined sampler laws pushed into the weighted geometry. The contraction-gap and budget calculations yield

\[W_{2,E}(\mu_n,\pi)\le r^nW_{2,E}(\mu_0,\pi)+\frac{6\gamma}{m}h\sqrt d\left(4\sqrt M+\frac{Mh}{2}\right).\]

The only initial transport premise is finite weighted Wasserstein distance, as required to use a real-valued distance. The thermostat parameters satisfy their variance identity \(\delta^2+\sigma^2=1\), with \(\delta=e^{-\gamma h/2}\); choosing \(\sigma=\sqrt{1-\delta^2}\) discharges it directly. The source step-size upper bound makes the displayed O(h\ensuremath{\surd}d) coefficient uniform in h. There is no externally supplied stationary process, Gaussian reference marginal, local coupling inequality, or convergence theorem among the hypotheses.



The one-step transport inequality has the form \(D_{n+1}\le rD_n+\delta_h\). Summing its geometric series gives \(D_n\le r^nD_0+\delta_h/(1-r)\). The scalar gap estimate converts this into the explicit dimension-dependent global bias.

\subsection{\texorpdfstring{Expose finite output distance for projection and mean-\allowbreak{}square use}{Expose finite output distance for projection and mean-square use}}


The same fully constructed exact comparator and integrated local cost establish that every weighted OBABO output has finite Wasserstein distance to the target whenever its initial weighted distance is finite. This theorem exposes the extended-real induction before conversion to the real endpoint. It allows subsequent coordinate projections and RMS bounds to use the result without any possibility of infinity being hidden by a real-valued distance.



A finite initial discrepancy and the finite local cost imply finite distance at every numerical step. This is proved before the distance is used as a real-valued bound. Subsequent coordinate projection and root-mean-square consequences therefore apply to finite transport quantities.

\section{\texorpdfstring{The original OBABO position laws satisfy the full sampling bound}{The original OBABO position laws satisfy the full sampling bound}}
\label{proof-module-KineticLangevin-OBABOPosition-lean}


The initial phase law is standard Gaussian velocity independent of the specified position law. Coupling these velocities identically costs nothing. The explicit inverse weighted coordinate map then projects the proved full Markov theorem to the position laws, with factor \ensuremath{\surd}2 and finite output transport established separately.



The weighted phase estimate is projected only at the final endpoint. A common initial Gaussian velocity then expresses its initial cost using position transport alone, yielding the position sampling theorem.

\subsection{\texorpdfstring{Project the contracting geometry to the physical position}{Project the contracting geometry to the physical position}}


The projection is the position coordinate of the explicitly verified inverse transform. The metric lower bound gives \(\|x-y\|^2\le2\|T(v,x)-T(w,y)\|^2\). Both sides are nonnegative, so the square comparison yields the global \ensuremath{\surd}2 Lipschitz coefficient. Composing this projection with the original weighted transform is exactly the ordinary position projection; this is proved from the inverse identity, not supplied as a law equivalence.



The inverse weighted map recovers position by a fixed linear combination of its coordinates. Cauchy–Schwarz gives its global Lipschitz constant. Applying this map only after the phase-law estimate introduces one positional factor, rather than a new factor at every numerical iteration.

\subsection{\texorpdfstring{Adjoin the same Gaussian velocity to every initial position coupling}{Adjoin the same Gaussian velocity to every initial position coupling}}


For every coupling of two position laws, append one common independent Gaussian velocity. Their phase difference has zero velocity, so its weighted energy is exactly the squared original position distance. The cross-space random-map theorem takes the coupling infimum and yields \(W_{2,E}(\mathcal N\otimes\mu,\mathcal N\otimes\nu)\le W_2(\mu,\nu)\). Product-measure swapping proves that the pushforward initialization is the stated independent Gaussian-times-position law. No optimal position coupling is assumed.



Using a common initial velocity leaves only the original position difference in the weighted cost. The added Gaussian source integrates to mass one. This produces a valid phase coupling for every initial position plan and therefore passes directly to the position Wasserstein infimum.

\subsection{\texorpdfstring{Obtain the unconditional position estimate and expose finite transport}{Obtain the unconditional position estimate and expose finite transport}}


Start from \(\mathcal N(0,I_d)\otimes\mu\), use the independent Gaussian-pair OBABO transition at every step, and let \(\nu_n\) be its position marginal. Under the original strong-convexity, smoothness, friction and step-size restrictions, theorem proves both finite \(W_2(\nu_n,\pi)\) and

\[W_2(\nu_n,\pi)\le\sqrt2\,r^nW_2(\mu,\pi)+\frac{6\sqrt2\gamma}{m}h\sqrt d\left(4\sqrt M+\frac{Mh}{2}\right),\qquad r=\sqrt{1-\frac{mh}{3\gamma}}.\]

The initial weighted distance is bounded by the initial position distance using the preceding shared-velocity construction. The full weighted theorem and its finite-output counterpart apply to the defined Markov laws. The \ensuremath{\surd}2 Lipschitz transport theorem projects both output measures; the target projection is the already proved normalized Gibbs marginal. Output finiteness is proved before the real-valued inequality is used, so this endpoint can safely imply RMS consequences. The finite-dimensional index type is arbitrary, allowing direct reuse for the scaled particle configuration.



The exact target’s projected law is \(\pi\), while the projected numerical law is the OBABO position output. Apply the Lipschitz transport bound to both. The geometric transient and accumulated local bias then have their physical position interpretation, with finite transport already established.

\section{\texorpdfstring{The KIPLMC2 sampling and parameter RMS bounds}{The KIPLMC2 sampling and parameter RMS bounds}}
\label{proof-module-KIPLMC-OBABOBound-lean}


The entire KIPLMC2 error estimate follows from the shared OBABO sampler, its exact particle clock change and the sharp parameter concentration theorem. Every initial law, Gaussian innovation, target transport and finite output moment is established internally. The explicit bound below has the intensive dimension \(\sqrt{(d_\theta+Nd_x)/N}\), and consequently also implies the supplied manuscript’s coarser particle dependence.



The KIPLMC2 sampling estimate follows from whole-law conjugacy and the shared invariant-comparator argument. Its intensive thermal dimension factor then combines with sharp equilibrium parameter concentration.

\subsection{\texorpdfstring{Transfer the independent thermal initialization and position law}{Transfer the independent thermal initialization and position law}}


The physical scaled velocity law is \(N^{-1/2}\) times a standard Gaussian, independently of the initial position \(\nu\). Multiplication by \(\sqrt N\) restores the exact independent standard-Gaussian initialization. The product pushforward identity proves this equality. Applying the already proved whole-chain Markov conjugacy and then projecting to positions gives equality of the KIPLMC2 position law and the standard-temperature OBABO position law at \(\gamma\prime=\sqrt N\gamma\), \(h\prime=h/\sqrt N\), for the sum potential \(V_N\).



The velocity scaling changes only thermal coordinate, so the position output law is unchanged. Whole-chain conjugacy identifies that output with the shared sampler under its transformed curvature, friction, and step. This permits direct transfer of the numerical theorem.

\subsection{\texorpdfstring{Cancel all clock factors in rate and admissibility}{Cancel all clock factors in rate and admissibility}}


With a general positive clock factor \(c\), curvature becomes \(c^2m,c^2L\), friction becomes \(c\gamma\), and the step becomes \(h/c\). Hence \(m\prime h\prime/(3\gamma\prime)=mh/(3\gamma)\). The friction threshold \(4L\le\gamma^2\) and the step threshold \(33\gamma^3h\le m\) are preserved by multiplying through by \(c^2\). One proves these cancellations directly, keeping the nonzero clock factor explicit.



The contraction ratio contains one curvature factor, one step factor, and one inverse friction factor. Their clock powers cancel. The friction and stability assumptions scale by a common positive factor, so their validity is preserved by the change of units.

\subsection{\texorpdfstring{Recover the intensive particle bias}{Recover the intensive particle bias}}


The shared OBABO position bias is \(6\sqrt2\gamma\prime h\prime\sqrt{d_z}(4\sqrt{L\prime}+L\prime h\prime/2)/m\prime\). Substitute the exact clock and curvature values. Since \(\sqrt{c^2L}=c\sqrt L\), cancellation gives

\[\frac{6\sqrt2\gamma}{m}h\sqrt{\frac{d_z}{N}}\left(4\sqrt L+\frac{Lh}{2}\right),\qquad d_z=d_\theta+Nd_x.\]

The algebraic theorem records the exact equality of the two expressions, not just an asymptotic comparison.



After substitution, the only dimensional fluctuation factor is the physical thermal radius \(\sqrt{d_z/N}\). This is intensive in the number of latent particles. The remaining friction, smoothness, and step factors are exactly those of the shared OBABO local-error accumulation.

\subsection{\texorpdfstring{Apply the OBABO sampling theorem}{Apply the OBABO sampling theorem}}


The sum potential inherits strong convexity \(Nm\) and smoothness \(NL\) from the original joint potential. Its Gibbs target is exactly the defined lifted particle target. The theorem applies the proved unit-temperature sampler bound, transports its position law through the exact conjugacy, and simplifies the rate and bias. It proves both finiteness of the output Wasserstein distance and

\[W_2(\nu_n,\pi_N)\le\sqrt2\left(1-\frac{mh}{3\gamma}\right)^{n/2}W_2(\nu,\pi_N)+\frac{6\sqrt2\gamma}{m}h\sqrt{\frac{d_z}{N}}\left(4\sqrt L+\frac{Lh}{2}\right).\]

The two-thermostat Gaussian laws are part of the constructed Markov recursion. No stationarity, reference-path or local-error estimate is assumed.



The numerical law, target, and all analytic assumptions have been identified before invoking the shared result. Its finite-output transport conclusion therefore transfers along with its inequality. The resulting physical-time bound applies to the two-thermostat Gaussian particle recursion.

\subsection{\texorpdfstring{Project to the parameter and add the sharp concentration term}{Project to the parameter and add the sharp concentration term}}


The initial finite second moment and strong convexity give a finite initial transport distance. The preceding theorem gives finite output transport. The parameter projection has norm at most one, so its RMS error around the unique MMLE is bounded by the position transport distance plus the Gibbs parameter RMS error. The latter is at most \(\sqrt{d_\theta/(Nm)}\). The final theorem therefore adds this third term to the KIPLMC2 sampling estimate and expresses the result as the square root of the parameter-error integral.



Project an output coupling to the common parameter and use its nonexpansiveness. The equilibrium powered marginal lies within the sharp statistical radius of the unique maximizer. Adding this radius to the numerical transport bound yields the displayed parameter RMS estimate.

\section{\texorpdfstring{The supplied KIPLMC2 rate with an explicit dimension-\allowbreak{}independent constant}{The supplied KIPLMC2 rate with an explicit dimension-independent constant}}
\label{proof-module-KIPLMC-OBABOPaperRate-lean}


The preceding sampler theorem gives a stronger intensive bias. This section derives the precise coarser form printed in the supplied manuscript, including its additional transient term. The implicit constant is replaced by an explicit expression that depends only on \(m,L,\gamma\), and the final result concerns the parameter-error integral.



The sharper proved physical bound implies the manuscript’s coarser envelope term by term. An explicit constant reveals its dependence on curvature and friction while remaining uniform in particle count and dimension.

\subsection{\texorpdfstring{Turn the square-\allowbreak{}root geometric factor into the printed exponential}{Turn the square-root geometric factor into the printed exponential}}


Let \(r=\sqrt{1-mh/(3\gamma)}\). The scalar bound \(1-x\le e^{-x}\), followed by taking square roots, gives \(r\le e^{-mh/(6\gamma)}\). Raising to the integer iteration index proves \(r^n\le e^{-mhn/(6\gamma)}\). The proof also handles a nonpositive square-root argument directly, so the scalar envelope itself does not need an extra admissibility hypothesis.



The square-root geometric factor is converted to an exponential with half its linear exponent. This gives the manuscript’s continuous-looking decay in iteration count. Raising the scalar bound to an integer power preserves it because both sides are nonnegative.

\subsection{\texorpdfstring{Make the hidden constant explicit}{Make the hidden constant explicit}}


The allowed step satisfies \(h\le m/(33\gamma^3)\). Consequently the physical bias coefficient is bounded by

\[B=\frac{6\sqrt2\gamma}{m}\left(4\sqrt L+\frac{Lm}{66\gamma^3}\right).\]

Choose \(C=\sqrt2+B+1/\sqrt m\). All three summands are nonnegative, so \(C\) dominates the transient coefficient, bias coefficient and equilibrium concentration coefficient separately. One proves these comparisons from positivity, without permitting dependence on the particle count, dimensions or step.



The step-dependent coefficient is bounded using the explicit stability cap. The resulting constant depends only on curvature and friction. Taking one larger constant dominating all three coefficients makes its independence from particle count, dimension, and requested accuracy explicit.

\subsection{\texorpdfstring{Compare each term with the supplied coarse particle envelope}{Compare each term with the supplied coarse particle envelope}}


For \(N\ge1\), \(D\le\sqrt N D+hN\sqrt{d_z}\) and \(\sqrt{d_z/N}\le\sqrt{Nd_z}\). The equilibrium term satisfies \(\sqrt{d_\theta/(Nm)}=(1/\sqrt m)\sqrt{d_\theta/N}\). Apply the exponential rate envelope, dominate each coefficient by \(C\), and add the three inequalities. This proves the exact printed shape, including the harmless nonnegative extra transient contribution \(hN\sqrt{d_z}\), from the sharper estimate.



The sharper intensive estimate can be enlarged term by term into the printed particle envelope. Every enlargement adds a nonnegative quantity or uses \(N\ge1\). This derives the coarser displayed statement without asserting that its particle dependence is optimal.

\subsection{\texorpdfstring{Apply the envelope to the particle parameter law}{Apply the envelope to the particle parameter law}}


The final theorem starts from the unconditional KIPLMC2 RMS estimate and supplies every scalar hypothesis using the original curvature, friction and step restrictions. With \(D=W_2(\nu,\pi_N)\), it proves

\[\sqrt{\mathbb E\lVert\theta_n-\theta_\star\rVert^2}\le C\left[(\sqrt ND+hN\sqrt{d_z})e^{-mhn/(6\gamma)}+h\sqrt{Nd_z}+\sqrt{d_\theta/N}\right].\]

This is the supplied manuscript theorem with an explicit constant and the independent Gaussian algorithm law. The proof does not assume its desired RMS estimate as an input.



The scalar comparison is applied to the unconditional parameter theorem, with its true initial transport discrepancy. Thus the coarser rate still describes the independent-Gaussian sampler. Its explicit constant makes clear which parameters are suppressed in asymptotic notation.

\section{\texorpdfstring{Logistic regression: gradient, Hessian, and convexity scope}{Logistic regression: gradient, Hessian, and convexity scope}}
\label{proof-module-KIPLMC-LogisticExample-lean}


The Bayesian logistic example has a Gaussian location prior and finitely many Bernoulli observations. The model is \(U(\theta,x)=\frac{\tau}{2}\|x-\theta\|^2+\sum_j[\log(1+e^{\langle v_j,x\rangle})-y_j\langle v_j,x\rangle]\), where \(\tau=\sigma^{-2}\). Additive normalization constants do not affect any derivative. This chapter derives the joint gradient and its Fréchet derivative, proves the exact global smoothness constant printed in the paper, and proves convexity and the full-span strict-convexity statement. It also proves that no finite unregularized model in a nonzero latent dimension has a positive global strong-convexity constant. Consequently the experimental model illustrates the algorithms but does not automatically satisfy the hypotheses of the quantitative convergence theorems.



The full joint derivatives distinguish global smoothness, pointwise strict convexity, and uniform strong convexity. Common-translation directions prove that finite unregularized logistic models need additional regularization for the global quantitative convergence assumptions.

\subsection{\texorpdfstring{Differentiate a single observation}{Differentiate a single observation}}


For a single response, write \(\ell_y(u)=\log(1+e^u)-yu\). The logarithm is evaluated at a positive argument, so the ordinary chain rule gives \(\ell_y'(u)=e^u/(1+e^u)-y=s(u)-y\). First, the argument identifies the sigmoid with this exponential ratio, then composes the exponential and logarithm derivative rules. The curvature is \(q(u)=s(u)(1-s(u))\). Since \(0<s(u)<1\), it is positive. The identity \(q(u)=1/4-(s(u)-1/2)^2\) gives \(q(u)\le1/4\). The mean-value theorem therefore supplies a global \(1/4\)-Lipschitz theorem for the sigmoid. These arguments do not require the response to be binary; the source's Bernoulli choices are included.



The response label enters the first derivative but cancels from the second derivative. Logistic curvature therefore depends only on the predictor. Its pointwise positivity and global upper bound distinguish strict convexity along observed directions from a uniform lower curvature bound.

\subsection{\texorpdfstring{Derive the first and second variations from the model}{Derive the first and second variations from the model}}


Take any affine joint line \((\theta,x)+t(u,w)\). The prior contributes \(\tau\langle x-\theta+t(w-u),w-u\rangle\) to the score. Observation \(j\) contributes \([s(\langle v_j,x+tw\rangle)-y_j]\langle v_j,w\rangle\). Differentiating this scalar function again gives \(\tau\|w-u\|^2+\sum_jq(\langle v_j,x+tw\rangle)\langle v_j,w\rangle^2\). The proof constructs the derivative of each affine coordinate, applies the squared-norm rule and the single-observation rule, and adds the finitely many derivatives. The second proof differentiates the resulting score. Thus the displayed quadratic form is justified by calculus on the original potential.



The common-translation direction \(u=w\) removes the Gaussian prior contribution entirely. Curvature in that direction must come from observations. The displayed second variation therefore reveals both the prior’s nullspace and the role of feature span in identifying the full joint state.

\subsection{\texorpdfstring{Find the exact Hessian nullspace}{Find the exact Hessian nullspace}}


Every term in the second variation is nonnegative when \(\tau\ge0\). If \(\tau>0\), a zero total forces \(w-u=0\) and \(\langle v_j,w\rangle=0\) for every feature. Conversely those equations make every term vanish. The proof isolates the nonnegative prior term, bounds each observation term by the whole sum, and uses strict positivity of every logistic weight to cancel it. If the features span the latent space, orthogonality to every feature extends to their linear span by submodule induction. Taking the inner product with \(w\) itself then gives \(w=0\), and hence \(u=0\). Every nonzero joint direction therefore has strictly positive second variation at each finite state.



Because every finite predictor has a positive logistic weight, zero observation curvature means orthogonality to every feature. Full span eliminates that nullspace at each finite state. This is a pointwise positive-definiteness statement; it does not provide a positive lower eigenvalue uniformly over all states.

\subsection{\texorpdfstring{Prove global convexity and full-\allowbreak{}span strict convexity}{Prove global convexity and full-span strict convexity}}


To compare two joint states \(z,z'\), restrict the potential to \(F(t)=U(z+t(z'-z))\). The two derivative theorems identify \(F''\) with the nonnegative quadratic form. The real-variable second-derivative criterion proves convexity of \(F\). Applying its Jensen inequality at \(t=0,1\) yields the full joint convexity inequality. Under the span assumption and \(z\ne z'\), the direction is nonzero, so the second derivative is positive at every finite \(t\). The strict second-derivative criterion yields strict Jensen inequality. The argument explicitly rewrites \(z+b(z'-z)=az+bz'\) using \(a+b=1\), so the conclusion concerns the original joint potential on the whole space.



Convexity is proved on arbitrary line segments, so it is a global property of the joint potential. Strict positivity of the second derivative on each nonconstant segment yields strict convexity. The argument does not turn the varying second-derivative lower bound into a uniform constant.

\subsection{\texorpdfstring{Use the Euclidean norm for the joint gradient}{Use the Euclidean norm for the joint gradient}}


The source norm is the Euclidean product norm \(\|(\theta,x)\|^2=\|\theta\|^2+\|x\|^2\). the joint state is equipped with this Euclidean Hilbert norm for the gradient claims. The coordinate equivalence is continuous and linear. Differentiating the prior and all finite observations proves Fréchet differentiability. The explicit gradient is \((-\tau(x-\theta),\;\tau(x-\theta)+\sum_j[s(\langle v_j,x\rangle)-y_j]v_j)\). To identify it with the Hilbert gradient, The proof compares the Fréchet derivative applied to an arbitrary direction with the already proved directional score, expands the product inner product, and uses uniqueness of derivatives. The resulting equality is with the the gradient, not a separately named surrogate force.



The Euclidean joint norm is needed because the gradient is defined through its inner product with arbitrary joint directions. A maximum product norm would not give that Hilbert gradient interpretation. The explicit formula therefore uses the same geometry as the smoothness and Hessian statements.

\subsection{\texorpdfstring{Recover the exact printed global smoothness constant}{Recover the exact printed global smoothness constant}}


The prior block maps a joint difference \((a,b)\) to \((a-b,b-a)\). Its squared norm is \(2\|b-a\|^2\le4(\|a\|^2+\|b\|^2)\), hence its operator norm is at most two. For observation \(j\), the sigmoid Lipschitz bound and Cauchy–Schwarz give \(\|(s(\langle v_j,x\rangle)-s(\langle v_j,x'\rangle))v_j\|\le\frac14\|v_j\|^2\|x-x'\|\). The latent coordinate projection has norm at most one. Adding these bounds proves \(\|\nabla U(z)-\nabla U(z')\|\le(2\tau+\frac14\sum_j\|v_j\|^2)\|z-z'\|\). The proof first certifies the squared-norm estimate, then bounds each summand, and finally decomposes the gradient difference into its prior and observation parts. With \(\tau=\sigma^{-2}\), this is exactly the source constant.



The prior couples parameter and latent directions through a block operator with norm two. Each observation contributes a rank-one latent operator with norm bounded by one quarter of its squared feature norm. Summing these operator bounds yields the printed global Lipschitz constant.

\subsection{\texorpdfstring{Identify the full block Hessian as a Fréchet derivative}{Identify the full block Hessian as a Fréchet derivative}}


The operator is \(\nabla^2U(\theta,x)=\tau\begin{pmatrix}I&-I\\-I&I\end{pmatrix}+\begin{pmatrix}0&0\\0&\sum_jq(\langle v_j,x\rangle)v_jv_j^{\mathsf T}\end{pmatrix}\). Each rank-one observation block sends a direction \((u,w)\) to \((0,q_j\langle v_j,w\rangle v_j)\). The argument constructs these maps as continuous linear maps, forms their finite sum, and transports the block operator to the Euclidean joint space. Differentiating both gradient coordinates with the Fréchet chain rule gives exactly this operator. Expanding its product inner product recovers the earlier curvature formula. This closes the connection between the matrix printed in the experiment and the derivative used by the formal calculus.



The block Hessian is obtained by differentiating the joint gradient. Its quadratic form recovers exactly the previously derived line second variation. This closes the consistency between the coordinate formula, directional calculus, and the operator used in theoretical assumptions.

\subsection{\texorpdfstring{An explicit full-\allowbreak{}rank counterexample separates strict from strong convexity}{An explicit full-rank counterexample separates strict from strong convexity}}


Already one scalar feature \(v_1=1\) spans the latent space. With response one and common location \(\theta=x=t\), the Gaussian prior vanishes and \(U(t,t)=\log(1+e^{-t})\). For \(t\ge0\), this lies between zero and \(\log2\). A positive strong-convexity constant \(m\) would force the midpoint inequality between \((0,0)\) and \((2t,2t)\) to contain the negative correction \(-mt^2\). Choosing \(t^2=(\log2+1)/m\) makes that right-hand side negative, contradicting nonnegativity. One proves the logarithm identity, both bounds, the exact Euclidean distance, and this contradiction. The full-span theorem applies to this example, so strict convexity is compatible with failure of every positive uniform lower curvature bound.



Along common positive translation in the one-feature example, the potential remains bounded while a strongly convex function would grow quadratically away from its midpoint. The contradiction is explicit despite full feature rank. Thus strict convexity alone cannot justify the sampler’s strong-convexity convergence theorem.

\subsection{\texorpdfstring{Prove the strong-\allowbreak{}convexity caveat for every finite unregularized model}{Prove the strong-convexity caveat for every finite unregularized model}}


The obstruction is general. For every real response and predictor, \(\ell_y(u)\le\log2+(1+|y|)|u|\), because \(1+e^u\le2e^{|u|}\). Along any nonzero common translation \((\theta,x)=(tw,tw)\), the prior stays zero, while the finitely many observation terms grow at most linearly: \(U(tw,tw)\le|J|\log2+|t|C_w\), where \(C_w=\sum_j(1+|y_j|)|\langle v_j,w\rangle|\). Applying a hypothetical strong-convexity midpoint inequality to the opposite states at \(\pm t\) forces \(m\|w\|^2t^2\le C_wt\). Choose \(t=(C_w+1)/(m\|w\|^2)>0\) to contradict it. This proves the source caveat in every nonzero latent dimension, regardless of feature rank. It explains why further regularization is needed before applying the paper's global strong-convexity convergence bounds to these experiments.



For any finite observation set, common translation eliminates the prior and leaves at most linear growth of the observation losses. Uniform strong convexity would require quadratic growth in that same nonzero direction. Additional regularization is therefore necessary to meet the global quantitative hypotheses in these unregularized examples.

\section{\texorpdfstring{Actual Brownian paths under a kinetic clock change}{Actual Brownian paths under a kinetic clock change}}
\label{proof-module-KineticLangevin-BrownianScaling-lean}


Time change is a probability statement as well as a deterministic substitution in the drift. If \(c>0\), the process \(\widehat B_t=\sqrt c B_{t/c}\) has the same Brownian finite-dimensional laws, because its covariance is \(c\min(s/c,t/c)=\min(s,t)\). This section proves that assertion on the already constructed canonical source, retains independence of entire coordinate paths, and calculates the physical noise amplitude.



The clock change is proved at the level of the Brownian coordinate law. Applying it together with velocity scaling gives the correct normalized thermal input for the kinetic equation.

\subsection{\texorpdfstring{Change the coordinate process and prove its Brownian law}{Change the coordinate process and prove its Brownian law}}


the rescaled Brownian motion applies the clock directly to each coordinate of the canonical source. standard analysis’s proved Brownian scaling theorem uses \(a^{-1/2}B_{at}\); setting \(a=c^{-1}\) gives the desired formula. The argument simplifies the inverse square root and the double inverse to identify this process exactly. the Brownian finite-dimensional laws supplies all finite-dimensional Gaussian distributions and independent increments, rather than a single matching variance.



For a positive clock factor, scale time and Brownian amplitude together so that increment variance remains equal to elapsed new time. Centering, covariance, Gaussianity, and independent increments then identify the rescaled coordinate process as standard Brownian motion.

\subsection{\texorpdfstring{Preserve independence and recover the kinetic temperature}{Preserve independence and recover the kinetic temperature}}


Each entire changed-clock coordinate path is a measurable function of its own original coordinate path. Independence is therefore preserved by measurable composition. Finally, velocity scaling by \(c^{-1}\) changes the standard thermal input to \(c^{-1}\sqrt{2c\gamma}\,B_{t/c}=\sqrt{2\gamma/c^2}\,\widehat B_t\). The scalar proof uses the square-root product and quotient identities with their nonnegativity hypotheses. For \(c=\sqrt N\), this is exactly the physical particle temperature \(1/N\).



The rescaled Brownian path is a measurable function of the original source, so independence from initialization is preserved. Applying the same clock change to the kinetic noise amplitude restores the normalized thermal coefficient, provided velocity and friction are transformed consistently.

\section{\texorpdfstring{The canonical KLMC coupling on each initial transport plan}{The canonical KLMC coupling on each initial transport plan}}
\label{proof-module-KineticLangevin-CanonicalCoupledSegments-lean}


Fix any finite-cost coupling of the initial position law and the target position law. The probability space is its explicit product with an independent standard Gaussian velocity and the canonical vector Brownian source. This section constructs the numerical chain, the exact reference segments and the auxiliary exact segments on that same space, and proves every analytic segment requirement used by the quantitative KLMC theorem.



Each initial transport plan is realized together with a shared Gaussian velocity and canonical Brownian source. Exact response and solution identities connect this construction to the fixed numerical output law.

\subsection{\texorpdfstring{Specify the concrete probability space and source projection}{Specify the concrete probability space and source projection}}


The phase variable is velocity first. The sample space contains the two coupled initial positions, one shared Gaussian velocity and one vector Brownian source. The measure is the initial-coupling product extension; probability follows from its first marginal and from the Gaussian/\allowbreak{}Brownian probability factors. The Brownian-source projection preserves its measure. Pullback along that projection gives a continuous linear map on L\ensuremath{{}^{2}}, and its representative theorem identifies it almost everywhere with evaluation on the product sample.



An initial transport plan is extended by an independent Gaussian velocity and Brownian source. The source projection preserves its probability law. This one extension realizes the numerical and reference processes together without altering either prescribed input marginal.

\subsection{\texorpdfstring{Construct the common Brownian input and filtered innovation}{Construct the common Brownian input and filtered innovation}}


The scaled source lift includes the physical amplitude \(\sqrt{2\gamma}\). At iteration \(n\), the input is the vector Brownian increment on \((nh,nh+t]\), and the filtered input is the exponential response to that same increment. The unscaled innovation pair consists of the exact velocity and position Wiener responses at the step endpoint. Both input paths are continuous in L\ensuremath{{}^{2}}, and the additive input starts at zero. These follow from the source lift's continuity and linearity and the proved vector-response properties.



The common additive input is the Brownian increment multiplied by the physical thermal amplitude. The filtered response uses the exponential kernel on that same increment. Their second-mean continuity permits the deterministic integral-solution machinery to be applied directly.

\subsection{\texorpdfstring{Identify the exact noise equations and defined numerical step}{Identify the exact noise equations and defined numerical step}}


The constructed filtered input satisfies \(B_n(t)=W_n(t)-\gamma\int_0^tB_n(s)ds\). Its endpoint is the scaled velocity innovation, and its time primitive is the scaled positional innovation. The argument transports the vector Brownian equations through the continuous linear source lift and commutes that lift with Bochner integration. Induction on the recursions then proves that the exponential-input chain is exactly the KLMC sequence, with the canonical correlated innovation pair and the prescribed fluctuation-dissipation scaling.



The filtered response obeys its damped integral equation and has exactly the numerical innovation at its endpoint. Its primitive gives the position innovation. These identities verify the defined exponential step using Brownian responses rather than a merely matching covariance matrix.

\subsection{\texorpdfstring{Connect Brownian L\ensuremath{{}^{2}} innovations with the generic output-\allowbreak{}law construction}{Connect Brownian L² innovations with the generic output-law construction}}


The L\ensuremath{{}^{2}} innovation pair has the same representatives as the pointwise canonical Brownian innovation. The proof pulls back the previously established vector-response representative theorem to the coupling product and combines the two component identities. Comparing with the generic product-law layer's representatives gives equality of the L\ensuremath{{}^{2}} pairs. Thus the numerical chain analyzed here has the coupling-independent output law already proved for the defined KLMC chain.



The L\ensuremath{{}^{2}} responses agree almost everywhere with the source’s pointwise innovations. Their joint law is therefore the canonical correlated Gaussian law already used to define the numerical output. The abstract coupling analysis and the fixed sampler law are now connected.

\subsection{\texorpdfstring{Build the numerical and stationary-\allowbreak{}reference initial states}{Build the numerical and stationary-reference initial states}}


The two initial phases share exactly one Gaussian velocity variable. The numerical phase uses the first coupled position; the reference phase uses the second. The reference representative is proved to be the explicit pair \((V,X_\mathrm{ref})\). Product independence and the two marginal laws prove that its law is \(\mathcal N(0,I_d)\otimes\pi\). The finite target second moment and finite coupling cost supply the required L\ensuremath{{}^{2}} initial positions.



The two initial phases share velocity and use the two positions of the supplied plan. The reference position has the Gibbs marginal, and the velocity is an independent standard Gaussian. Their phase law is therefore the concrete stationary candidate.

\subsection{\texorpdfstring{Construct all three chains and segment families}{Construct all three chains and segment families}}


The numerical chain applies the lifted original gradient and the constructed correlated innovations. The reference chain stitches exact selected kinetic segments, each driven by the canonical Brownian increment for that iteration. The corresponding reference segments start from these reference states. The auxiliary exact segment starts instead from the numerical state and receives the identical additive input. The drift is the pointwise gradient lifted to L\ensuremath{{}^{2}}; no substitute potential on the L\ensuremath{{}^{2}} space is assumed.



The numerical recursion, exact reference recursion, and auxiliary exact segments all use the same interval source. Their distinct roles are fixed by their initial states: numerical update, stationary exact continuation, and exact continuation from the numerical phase.

\subsection{\texorpdfstring{Prove every analytic premise of the KLMC coupling theorem}{Prove every analytic premise of the KLMC coupling theorem}}


The bundled specification supplies continuity of auxiliary paths, reference paths and filtered inputs; exact initial states; reference endpoint matching; the defined exponential update; both common-input integral equations; the exponential mild identity; the kinematic primitive identity; and equal initial velocities. Each field follows from the concrete definitions and earlier proved equations. The mild and kinematic fields are derived internally. The only remaining probabilistic input for a final stationary-error theorem is the stationary law of the global reference diffusion, proved separately and connected to these stitched segments by uniqueness.



Every local analytic identity follows from these definitions and the selected-flow specification. The only distributional step still needed for the final bound is stationarity of the global reference. Uniqueness transfers that stationarity to all the constructed local segments.

\section{\texorpdfstring{One canonical global kinetic reference behind all coupling segments}{One canonical global kinetic reference behind all coupling segments}}
\label{proof-module-KineticLangevin-CanonicalGlobalSegments-lean}


The exact reference segments are constructed on one canonical Brownian source. This section identifies their increment input with the increments of one full L\ensuremath{{}^{2}} Brownian input and proves that every stitched reference state and intermediate segment point belongs to the full kinetic solution.



The canonical local inputs are increments of one global Brownian input. A global kinetic solution and uniqueness then identify every stitched reference state and segment used by the sampler coupling.

\subsection{\texorpdfstring{Construct the full canonical input and identify its increments}{Construct the full canonical input and identify its increments}}


The full input is the scalar canonical Brownian value assembled spatially and pulled onto the initial-coupling extension with amplitude \(\sqrt{2\gamma}\). It is continuous in L\ensuremath{{}^{2}} because each assembly and pullback is continuous linear. The Brownian value at zero vanishes: The argument verifies that its deterministic time indicator has empty support, then uses linearity of the isonormal map. For nonnegative local time, the constructed segment input equals the full input difference at \(nh+t\) and \(nh\). For nonnegative local times, linearity of spatial assembly and source pullback preserves this increment identity.



The full canonical input is a continuous Hilbert-space Brownian path with zero initial value. Its local increments agree with the already constructed interval inputs. This consistency is what permits independently defined local solution segments to be compared with one global solution.

\subsection{\texorpdfstring{Construct the full reference solution}{Construct the full reference solution}}


The global reference is the selected kinetic solution with the lifted gradient, the canonical reference initial phase and the full canonical Brownian input. Its specification proves continuity, its exact starting phase, and the kinetic integral equation over the whole finite horizon. There is no stationary-law requirement in its definition.



The global reference is the selected kinetic solution over the complete horizon, initialized at the Gibbs phase. Its definition requires only the drift and forcing, while stationarity is proved afterward. The construction therefore avoids using invariance as an existence assumption.

\subsection{\texorpdfstring{Identify every stitched reference point by uniqueness}{Identify every stitched reference point by uniqueness}}


The generic stitching theorem now applies because the local inputs are proved to be increments of the full input. Consequently \(R_n=Y_{nh}\), and \(Z_n(t)=Y_{nh+t}\) throughout each step whose endpoint lies in the global horizon. The argument instantiates the existing exact uniqueness results with the L\ensuremath{{}^{2}} gradient and checks the horizon arithmetic. Thus a stationary law for the full constructed kinetic solution supplies the law of all reference states used by the quantitative KLMC proof.



Uniqueness identifies each shifted global path with the corresponding selected local segment. Induction identifies all grid endpoints, and the same argument identifies intermediate times. Hence the local coupling proof can use global stationary laws throughout each step.

\section{\texorpdfstring{The KLMC law bound with the initial Wasserstein distance}{The KLMC law bound with the initial Wasserstein distance}}
\label{proof-module-KineticLangevin-CanonicalKLMCBound-lean}


The numerical law is the pushforward of the defined KLMC chain on the independent Gaussian-velocity/\allowbreak{}Brownian product source. Every initial transport coupling yields this same output law. The complete constructed coupling proof gives the paper's first-order bound, and taking the transport infimum replaces the supplied coupling cost by the true initial Wasserstein distance. This intermediate module makes the sole stationary global-reference law property explicit; the final Gibbs theorem discharges it from the Brownian source.



The stationary global reference supplies all local marginals, and the numerical output is independent of the initial plan. The final bound is obtained by estimating every coupling and then taking its infimum.

\subsection{\texorpdfstring{Fix the numerical output law}{Fix the numerical output law}}


The definition uses the original pointwise KLMC method, the standard Gaussian initial velocity and the exact canonical correlated Brownian innovations. The numerical L\ensuremath{{}^{2}} chain is identified with the generic product-law construction by equality of the innovation variables. The established output-law theorem then proves that its pushforward is independent of which initial position coupling was chosen.



The numerical law is fixed by the initial position law, the independent standard Gaussian velocity, and the canonical innovation stream. It is defined without choosing a reference coupling. Every later coupling estimate therefore concerns one and the same sampler output.

\subsection{\texorpdfstring{Transfer one global stationary phase law to the reference segments}{Transfer one global stationary phase law to the reference segments}}


The named stationary property concerns the full constructed kinetic flow: at every time its phase law is \(\mathcal N(0,I_d)\otimes\pi\). It contains no numerical convergence conclusion. Exact chain-to-global identification transfers its position marginal to every reference grid point. Exact segment-to-global identification transfers its Gaussian velocity marginal to every intermediate time in a reference segment. The argument obtains these marginals by composition with the product-measure projections.



Exact uniqueness identifies the stitched reference segments with the global kinetic solution. Its phase stationarity supplies the Gibbs position at grid endpoints and the standard Gaussian velocity at every local time. These are the distributional facts used in the local-error integral.

\subsection{\texorpdfstring{Apply the proved coupling estimate to one initial plan}{Apply the proved coupling estimate to one initial plan}}


The initial position L\ensuremath{{}^{2}} discrepancy equals the true cost of the supplied transport coupling; this follows from the concrete extension's representative and norm identities. The segment specification furnishes all analytic local-error premises internally. Together with the derived stationary position and Gaussian velocity marginals, the quantitative L\ensuremath{{}^{2}} KLMC theorem gives \(W_2(\nu_n,\pi)\le\sqrt2(1-3mh/(4\gamma))^n\operatorname{cost}(c)+Mh\sqrt{2d}/m\). Output-law equality and initial-cost equality rewrite this into a bound for the numerical law.



For one finite-cost initial plan, the numerical and reference variables form an output coupling. The analytic recurrence bounds its cost, and endpoint projection gives a position-law transport estimate. The plan’s own root-mean-square cost remains visible at this stage.

\subsection{\texorpdfstring{Prove output finiteness and take the transport infimum}{Prove output finiteness and take the transport infimum}}


Finite initial Wasserstein distance supplies an finite-cost initial coupling. Its numerical and reference states are L\ensuremath{{}^{2}} variables, so their distance gives a finite transport bound on the output law. This proves output finiteness before real-valued transport arithmetic is used. The scalar recurrence estimate proves the contraction factor is nonnegative under the paper's step restrictions. The argument then applies the proved infimum theorem over all finite-cost couplings, obtaining the same estimate with \(W_2(\mu,\pi)\) in place of the cost of a chosen plan. No existence of an optimal coupling is assumed.



Finite initial transport provides at least one finite-cost plan, proving output finiteness. Since the same affine estimate holds for every such plan and its coefficient is nonnegative, infimizing replaces the plan cost by the initial Wasserstein distance.

\section{\texorpdfstring{The initial phase and Brownian innovation joint law}{The initial phase and Brownian innovation joint law}}
\label{proof-module-KineticLangevin-CanonicalReferenceNoise-lean}


The Gibbs stationarity proof needs the real joint law of the initial phase and all future filtered Brownian increments. This section derives that law on each arbitrary initial-position coupling extension. It combines the explicit product-source projection, the exact integration-by-parts response and the established independent thermal Gaussian stream.



The reference initial phase and the entire Brownian source have the required independent product law on the coupling extension. Exact filtered-increment identities supply thermal stream used by stationarity transfer.

\subsection{\texorpdfstring{Project the coupling source onto an independent phase/\allowbreak{}Brownian product}{Project the coupling source onto an independent phase/Brownian product}}


The projection sends the full sample to \(((V,X_\mathrm{ref}),\omega_\mathrm{Brownian})\). Its law is \((\mathcal N(0,I_d)\otimes\pi)\otimes P_\mathrm{Brownian}\). First, the argument projects the initial coupling onto its second marginal, retaining the independent Gaussian/\allowbreak{}Brownian source. It then applies measure-preserving product associativity and swaps the position and velocity factors. Thus independence of the complete reference initial phase and the complete Brownian source is proved for the product construction.



Project the coupling source onto the reference initial phase and the Brownian coordinate. Product associativity and the reference marginal identify its law as the product of the Gibbs phase and the entire Brownian source. This proves the required independence jointly, rather than interval by interval.

\subsection{\texorpdfstring{Identify the full-\allowbreak{}input OU increment}{Identify the full-input OU increment}}


The deterministic integration-by-parts input used in the kinetic splitting proof is \(\Delta W-\gamma\int_0^he^{-\gamma(h-s)}\Delta W_sds\). The local-input increment identity replaces each difference of the full canonical input by its constructed segment input. Applying the vector exponential-response formula through the continuous linear source lift proves that this OU increment is exactly the filtered velocity response. Its representative is consequently the scaled canonical vector response on the original Brownian source, pulled onto the coupling extension.



The filtered increment expressed through the full Brownian input equals the local canonical velocity response. This identity follows from the increment formula and bounded linearity of source pullback. Consequently the global-flow stationarity argument uses the same innovations as the local construction.

\subsection{\texorpdfstring{Derive the full independent thermal Gaussian joint law}{Derive the full independent thermal Gaussian joint law}}


The vector OU stream already has the exact iid Gaussian law with scale \(\sqrt{1-\psi_0(h)^2}\). The product-source joint-law theorem combines this stream with the initial phase target. Composition with the proved phase/\allowbreak{}Brownian projection law transfers that product law to the coupling extension. Finally, the representative identities identify its initial phase with the constructed reference initial L\ensuremath{{}^{2}} phase and its stream with the full-input OU increments. The resulting joint law is precisely the source obligation of the proved kinetic Gibbs stationarity transfer theorem.



The vector response stream has the independent thermal Gaussian law. Combining it with the phase/\allowbreak{}source projection gives the joint product law of reference initialization and all thermostat innovations. This supplies the precise source hypothesis of the splitting-limit stationarity theorem.

\section{\texorpdfstring{The complete first-\allowbreak{}order KLMC theorem for the sampler}{The complete first-order KLMC theorem for the sampler}}
\label{proof-module-KineticLangevin-CanonicalKLMC-lean}


This final theorem concerns the defined KLMC chain on the explicit independent Gaussian-velocity and canonical Brownian source. It proves the first-order Wasserstein bound of Dalalyan and Riou-Durand with the true initial Wasserstein distance. All continuous reference laws, common-noise identities, local-error estimates, output-law identities and transport-infimum steps are discharged inside the proof chain.



All reference, source-law, and local-error premises have been constructed internally. Combining them proves the explicit first-order position theorem for the independent correlated-Gaussian sampler.

\subsection{\texorpdfstring{Prove Gibbs stationarity of the global coupling reference}{Prove Gibbs stationarity of the global coupling reference}}


The reference begins with standard Gaussian velocity independent of the target Gibbs position. Its full Brownian forcing and its entire filtered innovation stream are constructed objects on the initial-coupling extension. Their exact initial/\allowbreak{}noise product law is already proved. The kinetic stationarity transfer theorem therefore applies directly to the integral solution. For a time \(t\) inside a larger horizon, The proof restricts the integral equation to \([0,t]\), supplies the continuous input and its zero starting value, and applies the proved splitting-limit stationarity result. This yields the full phase law \(\mathcal N(0,I_d)\otimes\pi\) at every reference time, without a remaining stationarity or Gaussian-marginal hypothesis.



The reference initialization and its entire Brownian innovation stream have the required independent product law. The splitting-limit stationarity theorem applies to the integral solution. Uniqueness then identifies its stationary samples with all reference segments used by the coupling argument.

\subsection{\texorpdfstring{Conclude the source-\allowbreak{}paper KLMC Wasserstein estimate}{Conclude the source-paper KLMC Wasserstein estimate}}


Let \(f\) be differentiable, \(m\)-strongly convex with \(m>0\), and have an \(M\)-Lipschitz gradient. If \(\gamma^2\ge m+M\), \(\gamma>0\), and \(0<h\le m/(4\gamma M)\), the KLMC position law satisfies \[W_2(\nu_n,\pi)\le\sqrt2\left(1-\frac{3mh}{4\gamma}\right)^nW_2(\mu,\pi)+\frac{Mh\sqrt{2d}}m.\] The initial position distribution may be arbitrary with finite initial Wasserstein distance; its velocity is the explicit independent standard Gaussian. First, the argument derives the finite Gibbs second moment from strong convexity. It then supplies the just-proved global-reference stationarity for every initial transport coupling to the complete canonical coupling/\allowbreak{}infimum theorem. The conclusion has no conditional local-error, reference-process, Gaussian-law, stationary-law or optimal-coupling premise. The smoothness assumptions are expressed by the gradient and strong convexity, so theorem also covers the paper's Hessian-bounded setting through the existing deterministic geometry results.



Combine the verified local comparison, scalar recurrence, endpoint projection, and all-couplings infimum. The result concerns the explicit independent-Gaussian KLMC sampler. Its linear step-size bias is the accumulated effect of force freezing, while the geometric term transports the initial position discrepancy.

\section{\texorpdfstring{Explicit accuracy budgets and integer iteration counts}{Explicit accuracy budgets and integer iteration counts}}
\label{proof-module-KineticLangevin-Complexity-lean}


The paper's step-size and iteration formulas are scalar consequences of its numerical error estimate. This section verifies that implication, retains a stability-compatible bias budget, and rounds the logarithmic threshold to a natural number. It handles zero initial error separately and clips negative logarithmic thresholds at zero. These accuracy corollaries have an explicit error inequality as input; they do not replace the outstanding unconditional stochastic theorem.



The error theorem becomes an executable prescription only after selecting a positive step and an integer count. The logarithmic threshold handles initialization, while an explicit allocation controls the persistent bias.

\subsection{\texorpdfstring{Replace geometric decay by its exponential envelope}{Replace geometric decay by its exponential envelope}}


For \(\alpha\le1\), the base \(1-\alpha\) is nonnegative. The elementary exponential inequality gives \(1-\alpha\le e^{-\alpha}\). Raising this comparison to a natural power preserves order, so \((1-\alpha)^n\le e^{-\alpha n}\). The argument uses the inequality \(1+s\le e^s\), monotonicity of nonnegative powers, and the exponential power identity. The lemma itself needs no positive lower bound on \(\alpha\); that requirement enters when dividing the logarithmic threshold.



For \(0\le\alpha\le1\), \((1-\alpha)^n\le e^{-\alpha n}\). The exponential form can be solved explicitly for the required iteration number. Its use must retain the admissible range of \(\alpha\), which comes from the sampler’s step restrictions.

\subsection{\texorpdfstring{Round the logarithmic threshold and include small initial errors}{Round the logarithmic threshold and include small initial errors}}


For initial coupling radius \(D\), define

\[n_0=\left\lceil\max\left\{0,\frac{\log(24D/\varepsilon)}\alpha\right\}\right\rceil.\]

This is a natural number even when the logarithm is negative. If \(\alpha>0\) and \(n\ge n_0\), then \(\alpha n\ge\log(24D/\varepsilon)\). The proof combines the ceiling's lower bound, the right-hand inequality for a maximum, and the cast of the natural comparison, before multiplying by the positive \(\alpha\). The explicit maximum and ceiling handle the initialization-dependent small-error regime and integer rounding.



The count is a natural-number ceiling of a nonnegative threshold. Taking the maximum with zero ensures that an already small initial error requires no negative number of steps. The ceiling adds fewer than one extra iteration beyond the real-valued requirement.

\subsection{\texorpdfstring{Verify the transient budget only when logarithmic cancellation is valid}{Verify the transient budget only when logarithmic cancellation is valid}}


For \(D>0\) and \(\varepsilon>0\), exponentiating the iteration inequality gives \(e^{-\alpha n}\le e^{-\log(24D/\varepsilon)}=\varepsilon/(24D)\). The geometric envelope and multiplication by \(D\) imply \((1-\alpha)^nD\le\varepsilon/24\). One proves positivity of the logarithm argument before using the exponential–logarithm identity, rewrites the negative exponential as an inverse, and cancels the nonzero factors. Zero initial error is handled by the next theorem without a logarithmic cancellation.



When the logarithm is positive, multiplying the chosen threshold by the contraction rate cancels its denominator. Exponentiating then reduces the initial error to the allocated fraction of accuracy. When the logarithm is nonpositive, the initial error already satisfies that allocation.

\subsection{\texorpdfstring{Combine the manuscript's 0.\allowbreak{}94 bias allocation with the transient}{Combine the manuscript's 0.94 bias allocation with the transient}}


Suppose an explicit estimate gives \(E_n\le\sqrt2(1-\alpha)^nD+b\), with \(0<\alpha\le1\), \(D\ge0\), and \(b\le0.94\varepsilon\). When \(D=0\), the bias already gives the conclusion and there is no iteration requirement. Otherwise use the rounded threshold. The exact rational inequality \(\sqrt2\le36/25\) bounds the remaining contribution by \(0.06\varepsilon\), so the total error is at most \(\varepsilon\).



One proves the square-root comparison by squaring, splits on zero versus positive initial radius, derives the logarithmic condition in the positive case, and combines the inequalities with nonnegative scalar multiplication. For KLMC, substitute \(\alpha=3mh/(4\gamma)\) and bias \(Mh\sqrt{2p}/m\), with a step size that is the minimum of the stability threshold and bias allowance. The same scalar argument applies to KIPLMC1 with its thermal radius and a separately budgeted concentration term.



The persistent bias and the decaying transient consume separate portions of the requested accuracy. The specific numerical allocation leaves enough slack for the \(\sqrt2\) positional prefactor and the rounded count. Adding the two verified contributions gives the final epsilon guarantee.

\section{\texorpdfstring{An explicit epsilon guarantee for the KLMC sampler}{An explicit epsilon guarantee for the KLMC sampler}}
\label{proof-module-KineticLangevin-Accuracy-lean}


The convergence theorem becomes an executable mathematical budget: choose a positive step, an integer iteration count, and obtain a certified error at most \(\varepsilon\). The result concerns the constructed sampler. Neither a recurrence nor a local-error estimate is left as a premise.



The first-order sampling theorem is converted into a positive-step, integer-iteration prescription. Explicit bias allocation and a rounded logarithmic transient bound handle finite initialization and requested accuracy.

\subsection{\texorpdfstring{Verify that the geometric factor is admissible}{Verify that the geometric factor is admissible}}


Set \(\alpha=3mh/(4\gamma)\). The curvature, friction and step assumptions imply \(0<\alpha\le1\). The proof reuses the exponential-coefficient estimate: \(e^{-mh/\gamma}+(3/4)Mh^2\le1-\alpha\). Its left side is positive, so the right side is nonnegative. This justifies every later geometric and logarithmic estimate.



A positive contraction rate ensures the logarithmic threshold is meaningful. The upper bound on the rate ensures the geometric base is nonnegative. Both properties are consequences of the friction and step restrictions, and should be before solving the error inequality for the iteration number.

\subsection{\texorpdfstring{Apply the scalar error allocation to the law}{Apply the scalar error allocation to the law}}


The KLMC theorem gives \(W_2(\nu_n,\pi)\le\sqrt2(1-\alpha)^nD+b\). The scalar budget allocates \(b\le0.94\varepsilon\) and \((1-\alpha)^nD\le\varepsilon/24\). The sum is at most \(\varepsilon\). The threshold is the natural ceiling of \(\max\{0,\log(24D/\varepsilon)/\alpha\}\); if \(D=0\), no initial-error iterations are needed. The finite-transport hypothesis prevents the real-valued Wasserstein convention from hiding an infinite distance.



When the initial discrepancy is already below its transient allocation, the maximum-logarithm prescription selects zero iterations. Otherwise the exponential envelope of the geometric factor reaches that allocation after the selected ceiling count. The bias allocation then completes the total-error bound.

\subsection{\texorpdfstring{Choose the paper step and round the count}{Choose the paper step and round the count}}


For positive dimension choose \(h=\min\{m/(4\gamma M),\,0.94\varepsilon m/(M\sqrt{2p})\}\). Both terms are positive. The first enforces stability and the second enforces the bias budget by multiplication through positive denominators. Construct the corresponding nonnegative-real The step, use the rounded logarithmic count, and apply the accuracy theorem. At \(\gamma=\sqrt{m+M}\), this is precisely the manuscript’s minimum-step prescription. The explicit reciprocal dependence displays the \(\sqrt p/\varepsilon\) iteration scaling up to the logarithm.



The chosen step is the minimum of a stability limit and an accuracy limit. Taking that minimum guarantees both requirements simultaneously. Its reciprocal makes explicit the tradeoff between small requested error, dimension-dependent stationary velocity fluctuations, and the curvature-dependent contraction rate.

\section{\texorpdfstring{Explicit particle, step, and iteration accuracy budgets}{Explicit particle, step, and iteration accuracy budgets}}
\label{proof-module-KIPLMC-Accuracy-lean}


These corollaries turn both particle-sampler theorems into finite epsilon guarantees. Statistical concentration, discretization, and initialization receive separate budgets. Integer ceilings and positive minimum steps make the prescriptions valid without informal asymptotic rounding.



Particle count, step size, and iteration count control separate statistical, numerical, and transient budgets. Positive minimum steps and natural-number ceilings turn the two sampler bounds into finite accuracy prescriptions.

\subsection{\texorpdfstring{Control KIPLMC1 with three concrete budgets}{Control KIPLMC1 with three concrete budgets}}


Write its bound as \(\sqrt2(1-\alpha)^nD+b+s\). Allocate \(b,s\le\varepsilon/3\), and choose the logarithmic count so the geometric factor times \(D\) is at most \(\varepsilon/24\). These fit within the shared \(0.94\varepsilon\) persistent-error allocation. The algorithm theorem supplies the bound internally, with its Gaussian law and MMLE target. Only choices of \(N,h,n\) remain in the hypotheses. Since \(b=(\sqrt2 L/m)h\sqrt{d_\theta/N+d_x}\) and \(s=\sqrt{d_\theta/(Nm)}\), the formula exposes the first method’s dimension and accuracy dependence directly.



The statistical and numerical floors each receive an explicit fraction of the requested error. The remaining geometric transient is reduced by the rounded logarithmic iteration count. The sum of these verified allocations proves the requested parameter accuracy for the first sampler.

\subsection{\texorpdfstring{Round the exponential and statistical requirements}{Round the exponential and statistical requirements}}


The exponential budget uses \(n\ge\lceil\max\{0,\log(24D/\varepsilon)/\alpha\}\rceil\), with a separate zero-\(D\) case. Monotonicity of the exponential and \(e^{-\log x}=x^{-1}\) give \(De^{-\alpha n}\le\varepsilon/24\). For a square-root statistical floor \(C\sqrt{d/N}\), choose \(N=\max\{1,\lceil9C^2d/\varepsilon^2\rceil\}\). The ceiling inequality, multiplication through \(N>0\), and the square-root identity prove the floor is at most \(\varepsilon/3\). Zero parameter dimension is included.



The particle ceiling ensures a positive integer count and the desired concentration budget, including zero parameter dimension. The iteration ceiling similarly converts a real decay threshold into a valid integer execution count. The zero-initial-error case avoids taking an unnecessary positive logarithm.

\subsection{\texorpdfstring{Choose the step and inspect the dimension algebra}{Choose the step and inspect the dimension algebra}}


For a positive stability cap \(A\) and positive linear-bias coefficient \(B\), the explicit choice \(h=\min\{A,\varepsilon/(3B)\}\) is positive, respects stability, and has \(hB\le\varepsilon/3\). This applies to each method’s displayed bias. The exact source identity \(\sqrt{N(d_\theta+Nd_x)}=(d_\theta/\varepsilon)\sqrt{1+d_x/\varepsilon^2}\) at the idealized choice \(N=d_\theta/\varepsilon^2\) explains the supplied KIPLMC2 complexity. Integer ceilings affect constants and are handled by the previous theorem. One proves the identity by equality of nonnegative squares.



The minimum step enforces stability and a linear bias budget simultaneously. Substituting the particle prescription into the dimension factor explains the source’s small-error scaling. Integer rounding is controlled separately rather than hidden inside the idealized algebra.

\subsection{\texorpdfstring{Apply the supplied KIPLMC2 rate with its uniform constant}{Apply the supplied KIPLMC2 rate with its uniform constant}}


Let \(C\) be the explicit constant depending only on curvature and friction, and set \(A=\sqrt N D_0+hN\sqrt{d_z}\). Apply the rounded exponential budget to \(D=CA\), with \(\alpha=mh/(6\gamma)\). Allocate \(Ch\sqrt{Nd_z}\le\varepsilon/3\) and \(C\sqrt{d_\theta/N}\le\varepsilon/3\). The Gaussian-chain theorem then gives RMS error at most \(\varepsilon\). For fixed curvature constants, the particle choice has order \(d_\theta/\varepsilon^2\), the small-error step has order \(\varepsilon^3/(d_\theta\sqrt{d_x})\), and the displayed reciprocal-step logarithmic threshold yields the source’s iteration order. Polynomial growth of initial distance is the additional condition under which that logarithm is suppressed in the source’s tilde notation.



The second method’s printed envelope has both a persistent particle-dependent bias and an additional transient coefficient. Both enter the rounded prescription. Suppressing the initialization logarithm in asymptotic notation requires the stated growth condition on the initial discrepancy.


\end{document}
